\documentclass[11pt,a4paper]{article}
\usepackage[margin=28mm]{geometry}
\usepackage[T1]{fontenc}
\usepackage[utf8]{inputenc}
\usepackage{lmodern}
\usepackage{amsmath,amssymb,amsthm,mathtools}
\usepackage{booktabs,longtable,array,enumitem,microtype,needspace}
\usepackage[hidelinks]{hyperref}
\hypersetup{pdftitle={Grothendieck's collars and cones, and Elliott's integer-valued polynomial conjectures},pdfauthor={}}
\usepackage{bookmark}
\setlength{\emergencystretch}{3em}
\setlength{\parskip}{2pt}
\newtheorem{theorem}{Theorem}[section]
\newtheorem{lemma}[theorem]{Lemma}
\newtheorem{proposition}[theorem]{Proposition}
\newtheorem{corollary}[theorem]{Corollary}
\theoremstyle{definition}
\newtheorem{definition}[theorem]{Definition}
\newtheorem{example}[theorem]{Example}
\BeforeBeginEnvironment{theorem}{\Needspace{6\baselineskip}}
\BeforeBeginEnvironment{lemma}{\Needspace{6\baselineskip}}
\newcommand{\R}{\mathbb R}
\newcommand{\Z}{\mathbb Z}
\newcommand{\Q}{\mathbb Q}
\newcommand{\C}{\mathbb C}
\newcommand{\F}{\mathbb F}
\DeclareMathOperator{\Homeo}{Homeo}
\DeclareMathOperator{\id}{id}
\DeclareMathOperator{\im}{im}
\DeclareMathOperator{\coker}{coker}
\DeclareMathOperator{\Spec}{Spec}
\DeclareMathOperator{\Hom}{Hom}
\DeclareMathOperator{\End}{End}
\newcommand{\workbenchurl}{https://github.com/KokunoYumeto/homotopy-and-categories-workbench}
\newcommand{\recordurl}{https://github.com/KokunoYumeto/mathematics-workbenches}
\title{Grothendieck's collars and cones, and Elliott's integer-valued polynomial conjectures}
\author{}
\date{23 November 2026}
\begin{document}
\maketitle
\begin{abstract}
Grothendieck's Additif proposes that collar and cone choices are
contractible and uses this expectation to compare categories of spaces
with boundary and spaces with marked points. We identify the obstruction
to forgetting these choices as a concordance space. The six-dimensional
torus disproves the unparametrized assertions; strict parametrized collars
and fixed-cone rigidifications are contractible. Of the eleven attachment
and forgetting functors, five are equivalences and six fail.
We construct a rational semialgebraic torus loop with no finite
semialgebraic filling carrying its specified face witnesses over any
real closed field. A countable cycle presentation computes stable
concordance components and unstable components from dimension six;
stabilization on components is surjective in dimension five.
We also compare ordinary TOP concordances with joint PL families in
low dimensions. Galvin's relative construction and the surgery formula
of Hambleton, Milgram, Taylor and Williams give a surjective
Casson--Sullivan homomorphism for oriented three-manifolds with finite
fundamental group. This obstruction vanishes on the standard three-torus.
For real projective three-space, joint odd germs
identify the full component group as an abelian kernel of forgetting
antipodal symmetry. An intrinsic PL partition with saturated singular
loci need not have conical charts along its strata. Coarsest embedded
smooth partitions also exist for arbitrary semialgebraic sets,
but depend on the embedding, as the compact cusp shows.

The algebraic part answers Elliott's integer-valued polynomial
conjectures. For every one-dimensional analytically irreducible
Noetherian local domain, its ring of integer-valued polynomials is a
free module of countably infinite rank. This contradicts the three
published nonflatness conjectures for this class. Elliott had already
announced freeness for the cusp example in 2013; we prove the full
local-domain statement and give polynomial witnesses for both singular
examples. Valuation orderings, finite-conductor approximation, Milnor
patching and Kaplansky's theorem carry the proof. All multivariable
tensor maps are isomorphisms and the compatible plethory is unique.
The one-branch hypothesis distinguishes this result from the
two-branch counterexamples of Ma and Jiang, Li, Sun, Wang, Xiao and Yu.
Over a Noetherian base and any coefficient extension, injectivity of
the binary polynomial tensor map forces isomorphisms in every arity;
we give the exact criterion with zero divisors and coefficient torsion.
Outside this setting, all tensor isomorphisms can hold without flatness.
Finally, compact parameter resolutions and exact relative replacement
prove that jointly PL and jointly definable concordance families have
the same weak homotopy type over every o-minimal real closed field.
The mathematics and text were produced and checked by GPT-6 Astra
(ChatGPT Pro mode), with a separate editorial reading of the freeness
conclusion by Claude Opus 5.5 at high effort; no human has checked them.
\end{abstract}

\section{Grothendieck's questions and the conclusions}
Collars and cones encode different kinds of local choices.  Choosing a
map from a product cylinder, choosing its image, and choosing radial
coordinates on a fixed cone produce different spaces of choices.  In
the Additif, cote 119, archive sheets 147--149 (PDF pp.~148--150), Grothendieck asks whether
these choices can be forgotten without changing the homotopy theory.
Beside the diagram on sheet~149 he writes: ``Syst\`eme des 11 foncteurs
qui sont des \'equivalences $\infty$-isotopiques'' \cite{Groth119}.
The archive sheet numbers printed below the scans are one less than the
PDF page numbers; both are supplied at the main citations.
The distinction matters for a programme that constructs global
stratified spaces by gluing tubular or conical neighbourhoods: a
noncontractible choice leaves automorphisms in the gluing problem.

The geometric and algebraic conclusions use different data. We state
their complete hypotheses first and give the proofs in the subsequent
sections. The topological arguments combine finite collar sliding,
Alexander contraction, and classical concordance theory. The all-field
comparison uses compact parameter resolutions, finite joins and marked
product prisms; the algebraic classification uses coefficient reduction,
local completion and explicit tensor obstructions.

\subsection{Collar and cone data}
We work in a fixed o-minimal expansion of a real closed field $K$.
Throughout the geometric statements, compact means closed and bounded and every map and family
is definable. A family over a finite parameter polyhedron $P$ is continuous
as a map of total spaces and has a continuous inverse on its image.
Its collar interior is required to be a neighborhood of $B\times P$ in
$X\times P$. These requirements are part of the meaning of a family;
pointwise collarhood alone is not used to substitute for them.

A \emph{strict collar embedding} is an embedding
$c:B\times[0,1]\hookrightarrow X$ with $c(b,0)=b$, whose image minus
$c(B\times\{1\})$ is open in $X$, and which extends to a collar on
$B\times[0,2]$. Equivalently, its outer boundary is collared from the
exterior, after matching the two boundary parametrizations. An
\emph{unparametrized strict collar} is $(C,\dot C)$ with such a
parametrization existing but not chosen. A rigidification chooses $c$.
 A \emph{bordant pair}
will mean a closed definable pair $(X,B)$ with a finite definable local
collar atlas. No local-manifold hypothesis is part of this definition.
The categorical statements about general pairs use compact $X,B$.
The local collar-gluing theorem also treats a finite atlas on a
noncompact definable pair, producing an open collar there.

The simplicial space of a choice has as its $q$-simplices the indicated
families over $\Delta^q_K$. We call it contractible when its geometric
realization is contractible. In particular, this is not the assertion
that the set of all definable maps is a definable finite-dimensional set.


A \emph{strict cone neighbourhood} of a marked point $o$ is a closed
subset $C\subset X$ with outer face $Y$, a homeomorphism
$CY=(Y\times[0,1])/(Y\times\{0\})\to C$ carrying the vertex to $o$
and fixing every point of the base $Y$, with $C\setminus Y$ open in
$X$, and with an exterior collar of $Y$ in
$X\setminus(C\setminus Y)$. The parametrization need only exist for an
unrigidified cone; a rigidification chooses such a base-fixed map.
A parametrization carrying the base merely onto $Y$ can always be
normalized by precomposition with the cone of the inverse of its
base restriction. We set
$C\varnothing=\{o\}$. For families these conditions hold on total spaces,
with every map over its parameter polyhedron.

Write $PB$ for the simplicial category of compact bordant pairs and
definable homeomorphisms of pairs.  Write $PBC$ for pairs equipped with
an unparametrized strict collar and $PBCR$ for pairs with a chosen collar
parametrization.  A morphism in the latter category is the product of its
boundary map and the identity of the interval on that collar.  Similarly,
$EP$ denotes compact spaces with a marked point, $EPC$ those with a strict
conical neighbourhood, and $EPCR$ those with its radial parametrization.
Morphisms preserve all indicated data; a rigid cone map is the cone of
its boundary map.  All mapping simplices are families over a standard
simplex.  An equivalence means a Dwyer--Kan equivalence: essential
surjectivity on the homotopy category and weak equivalences on mapping
spaces.  For every compact definable $B$, put
\[
 P_{\mathrm{def}}(B)=\operatorname{Homeo}_{\mathrm{def}}
 (B\times I;B\times\{0\}),\qquad I=[0,1].
\]
The bottom is fixed pointwise and the top is preserved setwise.  When $B$ is a finite PL polyhedron, the notation $P_{\mathrm{PL}}(B)$
imposes joint piecewise linearity on every finite parameter family.
A closed PL manifold here is a finite polyhedron locally PL homeomorphic
to a Euclidean space, without boundary. The simplicial operators are
pullbacks along the usual affine maps of standard simplices; weak
homotopy equivalence refers to geometric realizations.

\begin{theorem}[Fibrewise collar automorphisms and cone automorphisms]\label{result-1}
For every compact definable $B$, the simplicial group of automorphisms
\[
 (b,t)\longmapsto(b,h(b,t))\quad\text{of }B\times[0,1]
\]
is contractible when $h(b,0)=0$ and $h(b,1)=1$. For every compact
definable $Y$, the simplicial group of homeomorphisms of the closed
cone $CY$ fixing its base $Y$ pointwise and its vertex $o$ is contractible.
These prove the fibrewise form of cote~119, sheet~147 (PDF p.~148) I c) and cote~119, sheet~148 (PDF p.~149) II c),
respectively; they do not prove I c) for arbitrary automorphisms.
\end{theorem}

\begin{theorem}[Parametrized strict collars are contractible]\label{result-2}
Suppose a compact definable pair $(X,B)$ admits a strict collar.
The simplicial space $\operatorname{Col}^{\mathrm{rig}}(X,B)$ of all
strict collar embeddings is contractible. The statement concerns the family convention just defined.
\end{theorem}

\begin{theorem}[Unparametrized collars retain the classifying space of concordances]\label{result-7}
Use parametrizable families for the unparametrized collar simplicial set:
a simplex is a strict collar family admitting a parametrization, with no
parametrization retained. If $(X,B)$ admits a collar, then
\[
 \operatorname{Col}(X,B)\simeq BP_{\mathrm{def}}(B).
\]
In particular cote~119, sheet~147 (PDF p.~148) I b) is false for $B=(S^1)^6$:
$\pi_1\operatorname{Col}(X,B)\cong\pi_0P_{\mathrm{def}}(B)\ne0$.
The rigidified assertion Theorem~\ref{result-2} and the fibrewise assertion Theorem~\ref{result-1} are
the precise different assertions that remain true.
\end{theorem}

\begin{theorem}[A fully explicit nonzero concordance group]\label{result-3}
First let $B=(S^1)^d$, $d\geq6$, be the ordinary real torus with
its finite semialgebraic PL model. For $\mathrm{CAT}=\mathrm{PL}$ or
$\mathrm{TOP}$, put
\[
 P_{\mathrm{CAT}}(B)=
 \Homeo_{\mathrm{CAT}}(B\times[0,1]\,;B\times\{0\}).
\]
Then
\[
 \pi_0P_{\mathrm{CAT}}(B)
 \cong \F_2[t_1^{\pm1},\ldots,t_d^{\pm1}]/\F_2\{1\}
 =\bigoplus_{v\in\Z^d\setminus\{0\}}\F_2.
\]
The same component computation holds for definable concordances of this
model in every o-minimal expansion of a real closed field. In particular
we may take $B=(S^1)^6\subset K^{12}$, given by the six equations
$x_i^2+y_i^2=1$.
\end{theorem}

\begin{theorem}[The strict-cone choice obstruction in the stipulated categorical model]\label{result-12}
For a simplicial category of homeomorphisms, define its derived space
of choices over an object to be the homotopy fibre of the map on
classifying spaces of its maximal simplicial groupoids. For
$X=C B$ with $B=(S^1)^6$, the component of the standard strict cone
in the derived space of strict conical neighborhoods of the vertex
has homotopy type
\[
 BP_{\mathrm{def}}(B).
\]
Consequently II b) is false with this categorical interpretation of
choice, and remains false for rigidified strict-cone choices. The
space of base-fixed rigidifications of a \emph{fixed} cone is contractible, as
II c) asserts. These three assertions concern different spaces.
\end{theorem}

There is also an ordinary topological statement about cone images.
Let $B$ be a closed connected PL manifold of dimension at least four,
let $X=C_2B$ be its cone of radius two, and consider strict cone
neighborhoods of the vertex whose outer faces lie in $B\times(0,2)$.
Give their faces the quotient topology from the compact-open space
of two-sided locally flat embeddings of $B$, modulo
$\operatorname{Homeo}(B)$.  The component of the standard radius-one
cone has weak homotopy type $BP_{\mathrm{TOP}}(B)$.
If a parameterization of its outer face is chosen, the corresponding
component instead has the classifying weak type of homeomorphisms of
$B\times[1,2]$ fixing both ends.  For $B=T^6$, a rational
semialgebraic loop does not bound a real definable disk family whose
outer faces admit local jointly continuous parameterizations with
jointly continuous inverses.  Over every real closed field, it has
no filling on a finite triangulated disk by semialgebraic strict cones
with joint face parameterizations on each closed simplex and matching
images on common faces.  Theorem~\ref{thm:locally-parameterized-outer-faces}
proves these assertions using the completely regular mapping theorem
of Chapman--Ferry and the compact-family ambient extension theorem
of Chernavskii.  The quotient topology is essential to the statement;
the example following that theorem explains the difference from
Hausdorff convergence of images.

The attachment and collapse functors in the table below have the
following object maps. Collar attachment sends $(X,B)$ to
$X\cup_B(B\times[-1,0])$ with new boundary $B\times\{-1\}$ and
the displayed product collar. Cone attachment sends it to
$X\cup_B CB$ with the cone vertex marked and the displayed rigid cone.
Boundary collapse sends it to the pushout $X\cup_B\{*\}$ with $*$
marked; for $B=\varnothing$ this adjoins an isolated point. On a rigid
collar, collapse gives its corresponding radially parametrized cone.
The other arrows forget exactly the data named in the table. All maps
on morphisms are induced by the indicated products, cones and pushouts.

\begin{theorem}[The eleven arrows on cote~119, sheet~149 (PDF p.~150)]\label{result-13}
For the strict collar and strict cone conventions above, five of the
eleven arrows are equivalences of the definable simplicial categories.
The other six fail in the torus example of Theorem~\ref{result-3}. The following table
distinguishes the two copies of $PB$ in the drawing by position.
\begin{center}
\begin{tabular}{p{0.58\linewidth}p{0.30\linewidth}}
\toprule
Arrow & Outcome\\
\midrule
Top $PB\to PBCR$, attach a collar & equivalence (Theorem~\ref{result-10})\\
Top $PB\to EPCR$, attach a cone & equivalence (Theorem~\ref{result-10})\\
Top $PB\to$ bottom $PB$, identity & equivalence\\
Top $PB\to EP$, collapse & fails (Theorem~\ref{result-5})\\
$PBCR\to EPCR$, collapse & equivalence (Theorem~\ref{result-6})\\
$PBCR\to PBC$, forget product coordinates & fails\\
$EPCR\to EPC$, forget cone coordinates & equivalence (Theorem~\ref{result-6})\\
$PBC\to EPC$, collapse & fails\\
$PBC\to PB$, forget the collar & fails\\
$EPC\to EP$, forget the cone & fails (Theorem~\ref{result-12})\\
Bottom $PB\to EP$, collapse & fails (Theorem~\ref{result-5})\\
\bottomrule
\end{tabular}
\end{center}
Here ``fails'' means that the functor is not an equivalence in general;
we do not assert failure at every bordant pair with that boundary.
\end{theorem}

The torus component calculation is an application of the established
pseudoisotopy theorems of Hatcher--Wagoner, Igusa and
Burghelea--Lashof, with the formulation of Courte--Porcelli used in
its proof. The eleven-arrow conclusion is about the stipulated
simplicial categories and parametrizable choice families. In particular,
an ordinary set of neighbourhoods with an unspecified topology is not
a derived space of choices. The collar and cone absorption maps require
chosen extended collars and product-preserving maps to be strict natural
isomorphisms; their precise naturality statement is
Theorem~\ref{result-11}.

The explicit countable cycle presentation $E(K)$ in
\eqref{eq:countable-concordance-presentation} describes the stable
concordance component group, with its composition law, of every compact
smooth or PL manifold with finite simplicial homotopy model $K$.
Bottom and side boundary are fixed, and the upper face is free.
For every compact PL manifold whose components have dimension at
least six, it also gives the unstable component group. These
statements extend to the stipulated definable models over every real
closed field and fixed o-minimal expansion. No smoothing hypothesis
is needed; no terminating word problem is asserted. Theorems
\ref{thm:countable-stable-smooth-components} and
\ref{thm:stable-pl-component-presentation} give the stable presentation.
Theorem~\ref{thm:smoothable-concordance-six} first proves the unstable
range for smoothable models from Igusa's corrected geometric
obstruction sequence. Theorem~\ref{thm:component-followup-all-PL-six}
then treats every PL model: extension from a smoothable three-handle
core is surjective on components, and its natural stable comparison
is an isomorphism. For any specified triangulation $K$, one may use
its three-skeleton in the complete presentation, by
Corollary~\ref{cor:component-followup-chosen-three-skeleton}.

For a compact connected PL five-manifold, with the same bottom and
side fixed and upper face free, positive suspension is surjective
on concordance components. Thus the stable presentation $E(K)$ is
a quotient of its component group for every finite pointed homotopy
model $K$. No condition on $\pi_2$ is required. This statement holds
also smoothly, with boundary allowed, and for the stipulated
definable models over every real closed field and fixed o-minimal
expansion. Finite disjoint unions give finite products.
Proposition~\ref{prop:five-dimensional-stable-surjectivity} proves
this from Hatcher--Wagoner's realization theorem and the corrected
six-dimensional sequence. The kernel in dimension five remains
to be determined.



The bottom-fixed, top-free concordance component group is zero
for every PL ball, every compact manifold of dimension zero or one,
every closed PL surface, every compact orientable PL surface with
boundary, every standard PL sphere $S^n$ with $n\geq2$, and every
closed connected simply connected PL manifold of dimension three
or five. Finite disjoint unions of these cases are included. No condition on
$\pi_2M$ is required in the five-dimensional case. These assertions
include the standard spheres of dimensions three and four. The same
component statements hold for the stipulated definable models over
every real closed field, by the relative comparison.
Theorems~\ref{thm:low-concordance-components},
\ref{thm:standard-sphere-components},
\ref{thm:projective-plane-components}, and
\ref{thm:simply-connected-five-components}, together with
Corollaries~\ref{cor:all-closed-surface-components}
and~\ref{cor:simply-connected-three-components}, prove these
claims from the cited classical inputs. They assert no vanishing of
higher homotopy groups.

There is a different conclusion for every closed connected orientable
PL three-manifold $B$ with infinite fundamental group. Its concordance
component group contains a subgroup isomorphic to
$\bigoplus_{j\geq1}\mathbb Z$, and over the real field the inclusion
of this subgroup remains injective after all boundary conditions are
forgotten in $\pi_0\operatorname{Homeo}(B\times I)$.
The subgroup assertion holds for the stipulated definable models over
every real closed field as well. Theorem~\ref{thm:infinite-rank-three-components}
deduces this from the detection theorems of Lin, Xie and Zhang
\cite[Theorems~C and~D, pp.~13 and~28]{LinXieZhang2026} and the
relative PL--PD triangulation comparison. This describes an abstract subgroup, without identifying the entire component group.

The zero component group of the three-sphere does not extend to
vanishing of all higher concordance groups. With bottom fixed and top
free, Theorem~\ref{thm:three-sphere-higher-top-discrepancy} identifies
the actual real TOP concordance space with the based loop space of
\(\operatorname{hofib}(BO(4)\to BTOP(4))\), where
\(TOP(4)=\operatorname{Homeo}(\mathbb R^4;0)\) has its compact-open
topology. Both the TOP and jointly PL models have zero component and
first homotopy groups, but their second homotopy groups are respectively
$\mathbb Z/2$ and zero. The proof derives the germ comparison with
exact parameter-face lifting, then uses Hatcher's three-sphere theorem,
the four-connected part of Freedman--Quinn stabilization, and the
classical stable smoothing calculation. Thus an unrestricted
PL-to-TOP weak comparison already fails on this base.

For every standard real sphere $S^n$, $n\geq2$, define
$Y_k=\operatorname{hofib}_*(B\mathrm O(k)\to
B\operatorname{Homeo}(\mathbb R^k;0))$, with the usual topology
on the orthogonal group and the compact-open topology on the
origin-fixing homeomorphism group. Theorem~\ref{thm:top-sphere-finite-rank-stabilization}
identifies the ordinary bottom-fixed, top-free TOP concordance
space of $S^n$ with $\operatorname{hofib}_*(Y_n\to Y_{n+1})$,
where the map is product with a line and all basepoints are standard.
The proof combines Kuiper--Lashof's parametric germ extension
\cite[Lemma~3, p.~11; proof, pp.~12--13]{KuiperLashof1966I}
with local ambient extension and an explicit identification of
the cone map. Corollary~\ref{cor:four-sphere-low-concordance-groups}
then gives $\pi_iP_{\mathrm{TOP}}(S^4)=
\pi_iP_{\mathrm{PL}}(S^4)=0$ for $0\leq i\leq3$, using the
classical finite-rank stabilization calculations and the corrected
Freedman--Quinn theorem. In degree zero each space has one component.
The argument leaves degree four and general four-manifold bases
unevaluated.

\Needspace{5\baselineskip}
For every finite ordered simplicial model of a closed PL manifold,
Theorem~\ref{thm:effective-countable-concordance-models} constructs
effective countable PL and semialgebraic concordance groups. Their
finite codes have decidable validity, equality and degeneracy; their
group and simplicial operations are computable. Each model separately
has the same weak homotopy type as its geometric group over every
real closed field. Complete countable presentations of all homotopy
groups can be enumerated, with integer relation matrices in positive
degrees. The PL component presentation also applies to the stipulated
definable model, and in component dimensions at least five the PL
model gives the real TOP concordance space. The input is promised to
be a manifold; no recognition algorithm or terminating word problem
is asserted.

For closed real PL manifolds whose components have dimension at most
two, Theorem~\ref{thm:low-dimensional-real-concordance-comparison}
gives the natural weak equivalence from joint PL concordances to the
singular compact-open TOP concordance group, with bottom fixed and top
free. In dimensions zero and one both groups are weakly contractible.
The same theorem compares the full automorphism groups of every closed
surface, including nonorientable surfaces and all components. It follows
from the classical three-manifold family theorems and annulus calculation.
Corollary~\ref{cor:effective-low-dimensional-top-models} combines it with
the finite coding theorem to give effective countable TOP weak models
and complete enumerable homotopy-group presentations in these dimensions.

For a closed connected oriented real PL three-manifold $B$ with
finite cyclic fundamental group of even order,
Theorem~\ref{thm:cyclic-three-concordance-obstruction} constructs a
homeomorphism of $B\times I$ fixing both ends whose ordinary
bottom-fixed concordance component contains no PL or definable
representative. This includes $\mathbb {RP}^3$. The definable
assertion holds in every fixed o-minimal expansion of $\mathbb R$.
Corollary~\ref{cor:cyclic-three-absolute-obstruction} proves the
stronger conclusion that its absolute ordinary homeomorphism
component contains no such representative, even when the boundary
may move during an isotopy. Galvin's relative Casson--Sullivan
realization and earlier block comparison
\cite[Section~4.2 and p.~2528]{Galvin2026CassonSullivan}
\cite[Chapter~9, pp.~109--117]{GalvinThesis2024}
provide the obstruction; the argument below gives the ordinary
component conclusion. The absolute conclusion additionally uses
Galvin--Ladu's boundary-forgetting lemma
\cite[Lemma~5.3, pp.~17--18]{GalvinLadu2025Boundary}
and the uniqueness of compatible smoothings in this dimension.

The obstruction is natural on every closed connected real PL
three-manifold, including nonorientable ones.
Theorem~\ref{thm:natural-three-concordance-obstruction} constructs
a homomorphism under composition
\[
 c_B:\pi_0P_{\mathrm{TOP}}(B)\longrightarrow H^2(B;\mathbb Z/2)
\]
which vanishes on the PL and definable images. For a PL homeomorphism
$f:B\to B'$, it satisfies
$c_{B'}([(f\times1)H(f^{-1}\times1)])=(f^{-1})^*c_B([H])$.
If $B$ is oriented and has finite fundamental group,
Theorem~\ref{thm:finite-three-concordance-realization} proves that
$c_B$ is onto and realizes each class by a homeomorphism fixing
both ends pointwise. Thus the quotient by its kernel is
$H^2(B;\mathbb Z/2)$; for even cyclic fundamental group this is
$\mathbb Z/2$. Hambleton--Milgram--Taylor--Williams's relative
surgery formula and Oliver's Sylow-group calculation provide
the arithmetic input, and Galvin's marked relative construction
provides the realization. The descent to ordinary free-top
components normalizes the upper map and proves independence
of that choice. Determining which zero-class components have
PL representatives remains a further question.

The same construction realizes classes with prescribed local support.
Let $B$ be any closed connected real PL three-manifold, possibly
nonorientable, and let $W_1,\ldots,W_s\subset B$ be pairwise
disjoint compact codimension-zero PL submanifolds. Suppose
$W_i\cong C_i\setminus\operatorname{int}D^3$, where each $C_i$
is closed, connected, oriented and has finite fundamental group.
For $j_i:W_i\hookrightarrow B$, extension by zero satisfies
\[
\begin{split}
\operatorname{im}c_B
&\supseteq\sum_i(j_i)_!H^2(W_i,\partial W_i;\mathbb Z/2)\\
&=\operatorname{PD}_B^{-1}
       \left(\sum_i(j_i)_*H_1(W_i;\mathbb Z/2)\right).
\end{split}
\]
Theorem~\ref{thm:supported-three-concordance-realization}
realizes each local class by a concordance supported in
$\operatorname{int}W_i\times(0,1)$ and equal to the identity
near its entire boundary. For
$B=C_1\#\cdots\#C_s\#D$, with any closed connected PL
three-manifold $D$, the composite of $c_B$, mod-two duality
and projection onto $\bigoplus_iH_1(C_i;\mathbb Z/2)$ is onto;
its kernel contains the PL and each fixed-real-o-minimal
definable comparison image. Thus local finite-group pieces
give ordinary concordance obstructions without a restriction
on the ambient fundamental group. The proof combines the
relative realization with PL collars and the naturality of
relative duality.

For the standard real PL three-torus the same obstruction vanishes:
Theorem~\ref{thm:three-torus-concordance-obstruction} proves
\[
 c_{T^3}:\pi_0P_{\mathrm{TOP}}(T^3)
              \longrightarrow H^2(T^3;\mathbb Z/2)
 \quad\text{is zero}.
\]
The convention is again bottom fixed and top free. One proof uses
Ranicki's relative five-dimensional surgery sequence and the precise
quadratic assembly isomorphism for the free abelian group. A second
proof shows that, for each concordance, its bottom-fixed lifts through
all sufficiently large covers $[x]\mapsto[Nx]$ lie in the ordinary
identity component. It combines bounded lift displacement, time
compression and Edwards--Kirby's relative local deformation theorem.
The obstruction pulls back under these covers after a simultaneous
PL correction of their upper deck translations; odd covers detect
all of $H^2(T^3;\mathbb Z/2)$. Neither proof classifies the original
concordance components or asserts surjectivity of PL comparison.

For $D=D^4$ and $\iota(x)=-x$, let
$H=\operatorname{Homeo}_{\iota}(D)$ and
$A=\operatorname{Homeo}_{\iota}(S^3)$ have the compact-open
topology. Let $\mathcal G_{\iota}$ be the simplicial group
whose $q$-simplices are odd, parameter-preserving homeomorphisms
between neighborhoods of $\Delta^q\times\{0\}$, fixing that
zero section, modulo equality near the entire zero section.
The maps and their inverses are jointly continuous.
Let $\mathcal G_4$ be the same group without the oddness condition.
Theorem~\ref{thm:projective-concordance-joint-germ} proves that
restriction $\operatorname{Sing}H\to\mathcal G_{\iota}$ is
a vertex-surjective Kan fibration with identity fibre weakly
equivalent to $\operatorname{Sing}P_{\mathrm{TOP}}(\mathbb{RP}^3)$.
Radial coning and the linear inclusion give
\[
\begin{split}
P_{\mathrm{TOP}}(\mathbb{RP}^3)
&\simeq_w\operatorname{hofib}_1(A\to|\mathcal G_{\iota}|)\\
&\simeq_w\operatorname{hofib}_1(\mathrm O(4)\to|\mathcal G_{\iota}|)\\
&\simeq_w\Omega|\mathcal G_{\iota}/\operatorname{Sing}\mathrm O(4)|,
\end{split}
\]
including every orientation component. Arrows from topological
groups to realizations are represented by singular evaluation
and the realized simplicial homomorphism. Forgetting oddness gives
\[
\pi_0P_{\mathrm{TOP}}(\mathbb{RP}^3)
\cong\ker\bigl(\pi_1|\mathcal G_{\iota}|
                   \to\pi_1|\mathcal G_4|\bigr).
\]
In particular, this full component group is abelian. The proof
combines the projective annulus theorem of Kwasik--Schultz in
Boyle--Chen--Conway's formulation, Edwards--Kirby's local
extension theorem, and Cerf's comparison with Bamler--Kleiner's
Smale theorem. Exact parameter lifting is the additional step
that relates these results to joint germs. The displayed kernel
and the PL comparison subgroup are left unevaluated.

For a closed compact real PL four-manifold $B$, possibly
disconnected or nonorientable, let $F_A$ and $F_P$ be the identity
homotopy fibres of PL-to-TOP comparison for automorphisms of $B$ and
bottom-fixed, top-free concordances on $B\times I$. Let $F_H$ be the
corresponding fibre with both ends fixed. Theorem~\ref{thm:four-dimensional-comparison-fibres}
identifies
\[
 \operatorname{Map}(B,K(\mathbb Z/2,1))\simeq F_H
 \longrightarrow F_P\longrightarrow F_A
\]
as a homotopy fibre sequence. The mapping space is unbased, with
constant map as basepoint. Thus $\pi_jF_P\to\pi_jF_A$ is an
isomorphism for $j\geq3$, and the lower groups are related by the
exact sequence in that theorem.

Theorem~\ref{thm:four-dimensional-stabilization-fibre} specifies more:
$F_P$ is the zero homotopy fibre of
$\beta_B:F_A\to\operatorname{Map}(B,K(\mathbb Z/2,2))$.
The map comes from product with a line on sliced PL structures,
using the given PL structure as origin in the open five-dimensional
classification. On a pair $(a,\gamma)$, with $a$ PL and $\gamma$
a TOP path from the identity to $a$, its component value is the
relative PL-structure difference of $(x,t)\mapsto(\gamma_t(x),t)$
in $H^3(B\times I,B\times\partial I;\mathbb Z/2)
\cong H^2(B;\mathbb Z/2)$.

For the full component image, let $I_B$ consist of TOP concordance
components whose top map is TOP-isotopic to a PL map, and set
$J_B=\operatorname{im}(\pi_0\beta_B)$. Theorem~\ref{thm:four-dimensional-component-image}
proves that $J_B$ is a subgroup and constructs a homomorphism
$\lambda_B:I_B\to H^2(B;\mathbb Z/2)/J_B$ whose kernel is exactly
the PL concordance image. Its value is the relative structure class
after the top has been corrected to a PL map. The quotient of $I_B$
by the PL image is therefore a finite elementary abelian $2$-group
of order at most $2^{\dim H^2(B;\mathbb Z/2)}$. No surjectivity of
$\lambda_B$ is assumed. These statements follow from the classical
relative PL-structure theory with the full boundary conditions made
explicit; the endpoint condition and $\beta_B$ remain to be evaluated
for general four-manifolds.

The realization theorem of Orson--Powell--Randal-Williams
\cite[Theorem~2.3(D), p.~6; Definition~5.8, p.~21]
{OrsonPowellRandalWilliams2026} gives a conditional range calculation.
Let $B$ be connected, closed, oriented and real PL of dimension four,
and let $I_2:H_2(\pi_1B;\mathbb Z_{(2)})\to
L^q_6(\mathbb Z[\pi_1B])_{(2)}$ be their quadratic assembly summand,
specified in \eqref{eq:four-dimensional-assembly-summand}.
Let $L_B\subset H^2(B;\mathbb Z/2)$ be the Poincar\'e dual of the
kernel of the coefficient Bockstein
$H_2(B;\mathbb Z/2)\to H_1(B;\mathbb Z_{(2)})$.
If $I_2=0$, Theorem~\ref{thm:four-dimensional-obstruction-realization}
proves that $\operatorname{im}\lambda_B$ contains $(L_B+J_B)/J_B$.
If also $H_1(B;\mathbb Z_{(2)})$ has no $2$-torsion, then
$I_B/\operatorname{im}\pi_0c_P\cong H^2(B;\mathbb Z/2)/J_B$.
The proof identifies the relative obstruction classes and realizes
them by cylinders fixing both ends. For simply connected $B$,
Corollary~\ref{cor:simply-connected-four-dimensional-trace} gives
$\pi_0P_{\mathrm{TOP}}(B)=1$ and $J_B=H^2(B;\mathbb Z/2)$.
Every class in that cohomology group is the obstruction of an ordinary
automorphism loop, whereas jointly PL loops have zero obstruction.
Thus PL-to-TOP comparison on $\pi_1$ of automorphisms fails to be
surjective whenever $H^2(B;\mathbb Z/2)\ne0$. This last passage uses
the ordinary pseudo-isotopy theorem of
\cite[Theorem~1.5, p.~3]{GabaiGayHartmanKrushkalPowell2026}.
For $S^2\times S^2$ the obstruction target is $(\mathbb Z/2)^2$.

The same obstruction can be realized with prescribed support.
Let $j:W\hookrightarrow B$ be a compact, connected, simply connected,
codimension-zero PL submanifold with collared boundary; allow $W=B$.
Theorem~\ref{thm:supported-four-dimensional-traces} proves that ordinary
based TOP automorphism loops supported in $\operatorname{int}W$ realize
every class in
\[
 j_!H^2(W,\partial W;\mathbb Z/2)
 =\operatorname{PD}_B^{-1}\operatorname{im}
       \bigl(H_2(W;\mathbb Z/2)\longrightarrow H_2(B;\mathbb Z/2)\bigr).
\]
The complete loop fixes one common collar of $\partial W$.
If the displayed homology map is onto, then
$J_B=H^2(B;\mathbb Z/2)$ and $\operatorname{im}\pi_0c_P=I_B$.
This applies to $B=C\#D$ with $C$ closed, connected, simply connected
and oriented of dimension four, and $D$ closed, connected and oriented
with $H_2(D;\mathbb Z/2)=0$, including $C\#k(S^1\times S^3)$.
Thus a simply connected part can detect failure of comparison on
automorphism loops even when the ambient fundamental group is nontrivial.
The proof combines the compact relative realization theorem just cited,
the ordinary pseudo-isotopy theorem, and extension of the relative
PL difference class across a common identity collar.

For aspherical four-manifolds, surgery assembly gives a further
relative calculation. Let $B$ be closed, connected, oriented and real
PL of dimension four, with contractible universal cover. Put
$\pi=\pi_1B$, $N=B\times I$, and
$H_C=\operatorname{Aut}_C(N\,\mathrm{rel}\,\partial N)$ for
$C\in\{\mathrm{PL},\mathrm{TOP}\}$, with joint PL families and
ordinary TOP singular complexes. Assume $\operatorname{Wh}(\pi)=0$,
and let
\[
A_j:H_j(B;\mathbb L\langle1\rangle)\longrightarrow L_j(\mathbb Z\pi)
\]
be quadratic surgery assembly with trivial orientation character
and the geometric surgery decoration. Whitehead-group vanishing
identifies the simple and homotopy decorations here.
Theorem~\ref{thm:aspherical-four-relative-comparison} proves that
$A_7$-injectivity identifies the component kernel, as a group, with
$H^1(B;\mathbb Z/2)$ by the natural relative structure boundary.
If also $A_8$ is injective, then
\[
0\longrightarrow H^0(B;\mathbb Z/2)
\longrightarrow\pi_1H_{\mathrm{PL}}
\longrightarrow\pi_1H_{\mathrm{TOP}}\longrightarrow0
\]
is exact, and comparison is an isomorphism in every degree at least two.
If $A_6$ and $A_7$ are injective and every central element of $\pi$
is evaluated by a based jointly PL loop of automorphisms of $B$,
then component comparison is onto as well, giving
\[
0\longrightarrow H^1(B;\mathbb Z/2)
\longrightarrow\pi_0H_{\mathrm{PL}}
\longrightarrow\pi_0H_{\mathrm{TOP}}\longrightarrow0.
\]
No splitting is asserted. Under $A_6$-injectivity alone, the
ordinary-loop contribution $\eta_B$ to $\pi_0\beta_B$ vanishes
on any loop whose evaluation class has a jointly PL representative.
Thus the central realization condition gives $\eta_B=0$; it is
automatic when the centre is trivial.

The standard $T^4$ satisfies all these hypotheses. Its component
kernel is $(\mathbb Z/2)^4$, its fundamental-group kernel is
$\mathbb Z/2$, both comparisons are onto, and comparison is an
isomorphism in higher degrees. More generally, the assembly and
Whitehead hypotheses hold if $\pi$ is word-hyperbolic or acts
properly, cocompactly and isometrically on a finite-dimensional
CAT(0) space. Component surjectivity still requires the central
realization condition. Bartels--L\"uck's assembly theorem and
Ranicki's relative normal-invariant sequence carry the argument.
Orson--Powell--Randal-Williams use the same central realization
condition in a distinct stabilized smooth setting
\cite[Section~8.1, Remark~8.2 and Example~8.3, pp.~25--26;
Section~8.2, pp.~26--27]{OrsonPowellRandalWilliams2026}.
Here it corrects the relative homotopy class before the ordinary
PL-to-TOP comparison. The full free-top fibre, the remaining values
of $\beta_B$ and the endpoint-image condition remain separate questions.

Related four-dimensional smoothing questions are studied by Galvin
\cite[Theorem~1.1, p.~2497]{Galvin2026CassonSullivan}, who realizes every
Casson--Sullivan class for compact connected orientable smooth
four-manifolds after one $S^2\times S^2$ stabilization. The vanishing
criterion in \cite[Proposition~2.27, pp.~2510--2511]{Galvin2026CassonSullivan}
gives pseudo-isotopy to a diffeomorphism after finitely many such
stabilizations, following Freedman--Quinn. The endpoint condition in
our component criterion concerns ordinary isotopy on the original
four-manifold, while its relative cylinder obstruction is taken
modulo the classes of endpoint traces.

For an elementary extension of o-minimal structures expanding real
closed fields, Theorem~\ref{thm:elementary-extension-marked-automorphisms}
proves invariance of the joint definable automorphism group of any
closed bounded definable set with finitely many setwise preserved or
pointwise fixed definable subsets. The marks may be nonclosed. The
invariance is a weak equivalence on all components and in all degrees;
it includes the bottom-fixed, top-free concordance groups. For a
countable language, Corollary~\ref{cor:countable-definable-automorphism-model}
also gives a countable weak model. This is a cardinality conclusion:
effective recognition and enumerable presentations require the
additional constructions of the preceding theorem.

There is also a local statement for any finite collection of graph
formulas for homeomorphisms of a closed bounded definable set $K$
over a fixed o-minimal real closed field. Impose any finite list of
definable marked-set or boundary conditions and identify parameters
which define the same whole map. Give this quotient the sum of the
uniform distances between the maps and between their inverses.
Theorem~\ref{thm:finite-format-homeomorphism-contractions} proves that
it is definably homeomorphic to a Euclidean definable set. Around
each map and inside every prescribed positive metric ball, there is
an open definable neighborhood containing a positive metric ball
which strongly contracts to that map within itself. Evaluation and
inverse evaluation throughout the contraction are jointly definable
and continuous. The conclusion requires neither a group of formats
nor closed marked subsets; an empty underlying set gives an empty
or singleton quotient. The proof verifies the discrete-subspace
hypothesis of Walsberg's affineness theorem
\cite[Theorem~7.1, p.~16]{Walsberg2015}, then uses compatible triangulation.
The neighborhood may depend on the chosen formulas.

Proposition~\ref{prop:infinitesimal-pl-approximation-obstruction} gives
an elementary limitation on approximation: a continuous function on
$[0,1]$ which is affine on at most an ordinary finite number $m\geq1$
of intervals has uniform distance at least $1/(8m^2)$ from $x^2$,
even with arbitrary field coefficients and breakpoints. Thus a
positive infinitesimal tolerance cannot be met by a finite PL map.
The family $(1-s)x^2+sx$ nevertheless gives a semialgebraic isotopy
to the identity fixing the endpoints. Approximation at that scale
and existence of an isotopy are distinct assertions.

\Needspace{12\baselineskip}
Two relative consequences specify the permissible changes of parameter.
Let $M$ be any o-minimal real closed field, let $Z\subset M^n$ be
closed, bounded and definable, and let $F:[0,1]\times Z\to M$ be
continuous and definable with spatial Lipschitz constant $L$,
uniformly in time. Corollary~\ref{cor:increasing-definable-clock}
gives an increasing definable Lipschitz homeomorphism of $[0,1]$
fixing any prescribed finite set of times, after which
$|F(h(s),z)-F(h(s'),z')|\le K|s-s'|+L\|z-z'\|_2$ for some
field-valued $K$. A bounded Hausdorff-continuous definable family
of nonempty closed bounded sets has the same clock conclusion
in Hausdorff distance. The spatial constant stays $L$.
Over the same field, for $k\ge1$ and any basepoint on
$\partial[-1,1]^{k+1}$, Corollary~\ref{cor:based-sphere-hausdorff}
gives a based definable Lipschitz self-map, based homotopic to
the identity, making any bounded Hausdorff-continuous definable
family of nonempty closed bounded definable sets Hausdorff-Lipschitz. On a parameter
cube the corresponding map fixes its whole boundary whenever
the boundary family is already Hausdorff-Lipschitz. A family of
homeomorphisms of a closed bounded definable set which is the identity
on that entire boundary consequently admits such a based homotopy
with Lipschitz graphs, preserving all common spatial conditions.
The full family and its inverse are jointly continuous throughout.

\subsection{Comparison and intrinsic partitions}
\begin{theorem}[All-field weak equivalence]\label{result-44}
Let $M$ be any fixed o-minimal expansion of a real closed field,
and let $B$ be a closed compact PL manifold over $M$. Then
\begin{equation}
 P_{\mathrm{PL}}(B)\longrightarrow P_{\mathrm{def}}(B)
 \label{equation-115}
\end{equation}
is a weak equivalence. The statement compares finite jointly
PL parameter families with finite jointly definable parameter
families. It does not identify either space with the whole
ordinary TOP concordance space in exceptional dimensions,
nor does it by itself compute their homotopy type.
\end{theorem}

The comparison also makes the effective model apply to every fixed
o-minimal expansion. Let $R$ be its underlying real closed field,
let $B$ be a closed definably compact definable manifold, and choose
a definable triangulation $\tau:|K|_R\to B$ by a finite abstract
PL manifold. Corollary~\ref{cor:effective-whole-definable-concordances}
gives a specified weak equivalence
\[
             G^{\mathrm{PL}}(K)\longrightarrow P_{\mathrm{def}}(B),
\]
obtained by extending real-algebraic coefficients, including the
full PL families, and conjugating by $\tau\times\operatorname{id}_I$.
From $K$ one can enumerate complete presentations of every homotopy
group, including the full component group under composition;
the positive groups use the identity basepoint and translation.
The same output describes ordinary TOP concordances of
$|K|_{\mathbb R}$ when every component has dimension at most two
or at least five. The algorithm has the manifold condition as
an input promise and uses no decision procedure for the chosen
o-minimal theory. These recursively enumerable presentations
provide explicit descriptions; solving their word problems
is a separate question. Baro--Otero's homotopy comparison for
finite-dimensional semialgebraic sets is the relevant classical
background \cite[Theorems~3.1 and 4.1, preprint pp.~5 and 9]
{BaroOtero2008v1}. The corollary additionally uses the full
homeomorphism-family comparison just stated.

There is also a direct mapping-cylinder description of small
definable neighbourhoods. Let $X$ be closed and bounded in a fixed
o-minimal real closed field, let $f:X\to K_{\ge0}$ be continuous
and definable, and put $Y=f^{-1}(0)$. For every sufficiently small
$\epsilon>0$, the sublevel $V=f^{-1}([0,\epsilon])$ is a mapping
cylinder of a continuous definable map $B\to Y$, where
$B=f^{-1}(\epsilon)$. The identification fixes the entire core $Y$
and the specified outer face $B$, extends to an exterior collar of
$B$, and identifies the pointed quotient $V/Y$ with the rigid cone
on $B$. The same conclusion holds for a prescribed closed strict
tube $V\subset X$ with that core, outer face, exterior collar and
specified quotient-cone product off $Y$. For fixed $(V,B,Y)$,
any two such attaching maps are definably homotopic
(Theorem~\ref{thm:zero-locus-mapping-cylinder}). Empty cores and faces
are included. Delfs--Knebusch's semialgebraic neighbourhood theorem
is the classical antecedent
\cite[Theorem~2, p.~48]{DelfsKnebusch1984Retractions}; the proof here
applies scalar triangulation to the given function and preserves
the specified outer labels. The conclusion concerns the attaching
map's homotopy class, without asserting uniqueness of its homotopy.

The corresponding boundary-bundle collapse is not full for unrestricted
definable homeomorphisms, already when the core is one point, and it
also fails to be a Dwyer--Kan equivalence in the stated simplicial
categories. There is a precise product-preserving equivalence
(Proposition~\ref{prop:boundary-bundle-retained-equivalence}). Its
source objects are closed bounded bordant pairs $(E,B)$, with a
specified collared subset $B$, a closed bounded core $Y$, and a
locally trivial definable bundle $\varphi:B\to Y$. Its target objects
are closed bounded $(X,Y,V,B)$ with a specified identification
$M_\varphi\cong V$ fixing the core and the outer face, with
$V\setminus B$ open in $X$ and an exterior collar of $B$.
Target morphisms preserve the marked sets and have the exact product
form $[b,t]\mapsto[h_B(b),t]$, with
$h_Y\varphi=\varphi'h_B$. Attachment of $M_\varphi$ is an
equivalence for these objects and morphisms and an isomorphism on
their definable simplicial mapping spaces. The elementary failure
of fullness loses boundary directions; the enriched failure uses
the six-torus concordance component of Theorem~\ref{result-3} and
the cone contraction of Theorem~\ref{result-1}. The positive assertion
follows by cutting the tube and extending maps in its specified
coordinates. A scalar radius alone does not determine that extension.

The dimensions of the attachment fibres do not by themselves
characterize these neighborhoods. A compact semialgebraic disk
with a transverse segment attached at its centre in $S^3$ has
local codimension one there, yet every strict mapping-cylinder
tube has a positive-dimensional attachment fibre at that point
(Theorem~\ref{thm:tube-fibre-dimension-obstruction}). For a locally
trivial conical bundle over a pure-dimensional manifold, by
contrast, the link fibre does have dimension
$\dim_yX-\dim_yY-1$, with the empty-link convention $-1$.
That numerical equality also holds throughout the noncontractible
derived strict-cone component for $T^6$.

There is a stronger conclusion when the core is a boundary.
Let $V$ be a closed bounded definable manifold with boundary,
$Y$ a union of its boundary components, and $B$ a closed bounded
definable set. If $q:B\times[0,1]\to V$ is continuous and definable,
maps the bottom to $Y$, and is a homeomorphism from
$B\times(0,1]$ onto $V\setminus Y$, its bottom attachment is onto
and all its fibres are definably contractible. Over the reals,
with compact topological $V$ and compact metric ANR $B$, the same
map hypotheses imply cell-like fibres
(Theorem~\ref{thm:boundary-attachment-fibres}). The real conclusion
is a special case of the classical mapping-cylinder and neighborhood
theory of Lacher; the finite triangulation retraction proves the
internal definable contraction.

Consequently, for a closed bounded definable manifold $Y$ without
boundary, the jointly definable simplicial space of parametrizations
of its fixed tube $Y\times[0,1]$ with fixed top and zero-dimensional
attachment fibres is exactly $P_{\mathrm{def}}(Y)$ after interval
reversal (Theorem~\ref{thm:numerical-tube-concordances}). This holds
over every o-minimal real closed field. For the six-torus its
component group is the direct sum in Theorem~\ref{result-3}.

Allowing arbitrary attachment fibres still leaves these components
distinct. For $V=T^6\times[0,1]$, let $\mathcal R$ consist of all
continuous self-maps fixing the top, mapping the bottom into itself,
and homeomorphic off the bottom, with its compact-open topology.
The inclusion $G=\operatorname{Homeo}(V\;\mathrm{rel}\;
T^6\times\{1\})\subset\mathcal R$ induces an isomorphism of
component monoids, and every $q\in\mathcal R$ is the time-zero
limit of a path lying in $G$ at all positive times
(Theorem~\ref{thm:unrestricted-tube-component-bijection}).
Edwards--Kirby local contractibility and Siebenmann's approximation
theorem and published limiting-isotopy Complement carry this
ordinary real assertion. Two fixed finite PL representatives of
distinct components cannot be joined by a jointly definable family
over the reals or a semialgebraic family over any real closed field
(Theorem~\ref{thm:unrestricted-tube-components}). The component
bijection concerns ordinary real paths; its countable isotopy
construction gives no definable path-surjectivity conclusion.

For the next independent conclusion, PL regularity is intrinsic to a
compatible finite definable triangulation. A non-PL-manifold point is
one with no neighbourhood PL homeomorphic to a Euclidean space of any
dimension. A set is \emph{saturated} for a partition when it is a union
of its strata. Coarser means that every stratum of the finer partition
is contained in one stratum of the coarser partition.
All strata below are definable. For the displayed dimension indices
we use $\dim\varnothing=-1$; the empty space has the empty partition
and the filtration consisting only of $\varnothing$.
\begin{theorem}[A canonical intrinsic PL partition and its dimension filtration]\label{result-14}
Let $X$ be a definable set in a fixed o-minimal expansion
of a real closed field, and let $\Phi$ be a finite family of closed
definable subsets of $X$. There is a unique coarsest finite
partition $\mathcal P(X;\Phi)$ satisfying the following requirements:
its strata are locally closed PL manifolds, each of a single dimension;
each stratum closure is a union of strata; the sets in $\Phi$ are
unions of strata; and for every closed union $Y$ of strata, the
non-PL-manifold locus $\operatorname{Sing}_{\rm PL}(Y)$ is a union
of strata. Dimension may differ between strata. One may instead
require connected strata (definably connected over a general real
closed field); this class also has a unique coarsest
partition.

Both versions are preserved by definable homeomorphisms of the
marked data $(X;\Phi)$. Their dimension filtration
\[
 X_{\le k}=\bigcup_{\substack{S\in\mathcal P(X;\Phi)\\\dim S\le k}}S
 \quad(-1\le k\le\dim X)
\]
is closed and canonical, with pure-dimensional PL-manifold layers
after empty layers are omitted. If a closed filtration
$\varnothing=Y_{-1}\subset Y_0\subset\cdots\subset Y_s=X$ is
given, take $\Phi=\{Y_0,\ldots,Y_s\}$ and insert, in lexicographic
order, the closed sets
\[
 E_{j,k}=Y_{j-1}\cup(Y_j\cap X_{\le k}),
 \qquad 0\le j\le s,\quad -1\le k\le\dim X.
\]
Omitting repetitions gives a canonical regular equidimensional
filtration compatible with the preassigned one. The coarseness
assertion concerns the stated class of partitions; it is not the
false least-fine assertion among all weak filtrations.
\end{theorem}

The extra conditions are necessary. A weak regular equidimensional
filtration means a finite closed chain whose nonempty layers are
manifolds without boundary, each of one dimension; layers need not
be connected, and coarsening refers to their layer partitions.
Proposition~\ref{prop:weak-filtration-r3} constructs a compact
semialgebraic subset of $\mathbb R^3$ with two such filtrations,
with embedded smooth layers of dimensions $0,1,2$, whose only
common coarsening is the whole nonmanifold. The same formulas
give closed and bounded examples over every real closed field.
Proposition~\ref{prop:weak-filtration-pure-r4} gives a compact
semialgebraic example in $\mathbb R^4$ with local dimension two
at every point, again with smooth $0,1,2$ layers and only the
indiscrete common coarsening. Its local homology excludes a
topological or PL manifold. Both examples therefore rule out
a common weak regular coarsening for smooth or intrinsic PL
regularity, including the more restrictive subsequence convention.

Embedded smooth regularity gives a related construction over
$\mathbb R$. Let $X\subset\mathbb R^N$ be any semialgebraic set
and let $\Phi$ be any finite family of semialgebraic subsets of $X$.
For $A\subset X$, write $\Sigma(A)$ for the points at which $A$
is not locally an embedded $C^\infty$ manifold without boundary,
allowing all local dimensions. There is a unique coarsest finite
semialgebraic partition adapted to $\Phi$ whose strata are locally
closed, whose stratum closures are unions of strata, and for which
$\Sigma(Y)$ is a union of strata for every union $Y$ that is closed
in $X$. Every stratum is smooth. The same existence and uniqueness
hold separately with connected strata, with strata of one local
dimension, and with both requirements
(Theorems~\ref{thm:coarsest-embedded-smooth-partition}
and~\ref{thm:coarsest-smooth-partition-variants}).
Neither compactness nor local closedness of $X$ is required.
The classical inputs are the semialgebraicity of smooth loci in
Stasica and Koike--Shiota and the finite Nash frontier partitions
described by Coste; the proof below gives the finite construction
and the coarsest property, including all closed saturated unions.

A semialgebraic homeomorphism of the marked sets preserves these
partitions if it carries $\Sigma(Y)$ to $\Sigma(h(Y))$ for every
relatively closed semialgebraic $Y\subset X$. Restrictions of
ambient Nash diffeomorphisms satisfy this condition. Arbitrary
semialgebraic homeomorphisms do not: the interval $[-1,1]$ and
the cusp $\{(t^2,t^3):-1\le t\le1\}$ have finite singular
strata of two and three points respectively. Their coarsest
connected versions have three and five strata
(Proposition~\ref{prop:embedded-smooth-partition-naturality}).
Thus this embedded smooth construction records information that
the intrinsic PL partition does not record. It imposes no local
triviality or equisingularity condition.

\subsection{The polynomial operations and their exact tests}
For a unital inclusion of commutative rings $R\subset S$, define
\[
 \Omega_n(R,S)=\{f\in S[X_1,\ldots,X_n]:f(R^n)\subset R\},
 \qquad A=\Omega(R,S)=\Omega_1(R,S).
\]
Multiplication in disjoint variables gives
$\theta_n:A^{\otimes_R n}\to\Omega_n(R,S)$.
\emph{Unique polynomial coaddition and comultiplication} means that,
for every $f\in A$, both $f(X+Y)$ and $f(XY)$ have exactly one preimage
under $\theta_2$. This includes $f=0$, so uniqueness requires
injectivity of $\theta_2$. A ring epimorphism means an epimorphism of
commutative unital rings; equivalently, its multiplication map on the
tensor square is an isomorphism.

For $d\ge1$, let $D_d(R)$ be the ideal of $R$ generated by the
Vandermonde determinants $\prod_{0\le i<j\le d}(r_j-r_i)$, with
$r_0,\ldots,r_d\in R$. Write
$(R:_S D_d(R))=\{c\in S:cD_d(R)\subset R\}$.
The union $C$ used in the next theorem is a subring of $S$; it is
exactly the set of coefficients, or leading coefficients, that occur
in integer-valued polynomials. This coefficient assertion is proved in
Theorem~\ref{result-19}.
The notation $\operatorname{App}(R,C/R)$ denotes the $R$-module of all
functions from the underlying set of $R$ to $C/R$.
\begin{theorem}[A universal coefficient-syzygy criterion for the two polynomial tensors]\label{result-20}
Let $R\subset S$ be any unital inclusion of commutative rings. The
following construction uses only ideals, coefficient functions and
finite linear relations associated with this inclusion. It imposes
no finiteness, reducedness, flatness or domain hypothesis.

First form
\begin{equation}
\begin{gathered}
 C=\bigcup_{d\ge1}(R:_S D_d(R)),\qquad \mathcal F=C[X],\\
 E=\operatorname{im}\left(\mathcal F\xrightarrow{\varepsilon}
                       \operatorname{App}(R,C/R)\right),\\
 \varepsilon\left(\sum c_iX^i\right)(r)=\sum c_ir^i+R.
\end{gathered}
                                                        \label{equation-50}
\end{equation}
The first arithmetic test is that every $c\in C$ have a finite
matrix presentation
\begin{equation}
\begin{gathered}
 c=Y M Z,\qquad M\in\operatorname{Mat}_{n}(R),\\
 Y\in C^{1\times n},\quad Z\in C^{n\times1},\quad
 YM\in R^{1\times n},\quad MZ\in R^{n\times1}.
\end{gathered}                                           \label{equation-51}
\end{equation}
This test is equivalent to epicness of $R\to C$.
If it fails, the requested tensors are not uniquely determined.
When it holds, define the following finite-depth syzygy invariants.
Set
\begin{equation}
\begin{aligned}
 U_q&=\operatorname{Tor}_q^R(\mathcal F,\mathcal F),\\
 V_q&=\operatorname{Tor}_q^R(\mathcal F,E)
                    \oplus\operatorname{Tor}_q^R(E,\mathcal F),\\
 W_q&=\operatorname{Tor}_q^R(E,E),\qquad
 U_q\xrightarrow{a_q}V_q\xrightarrow{b_q}W_q,
\end{aligned}                                           \label{equation-52}
\end{equation}
where
\begin{equation}
\begin{aligned}
 a_q&=(\operatorname{Tor}_q(1,\varepsilon),
                         \operatorname{Tor}_q(\varepsilon,1)),\\
 b_q&=\operatorname{Tor}_q(\varepsilon,1)
                        -\operatorname{Tor}_q(1,\varepsilon).
\end{aligned}
                                                        \label{equation-53}
\end{equation}
The two terms of $b_q$ act on the respective summands of $V_q$;
thus $b_qa_q=0$. Put
\begin{equation}
 Z_q=\ker a_q,\quad T_q=\operatorname{coker}b_q,\quad
 K_1=\ker b_1/\operatorname{im}a_1.                       \label{equation-54}
\end{equation}
The canonical maps specified below are
\begin{equation}
\begin{gathered}
 \delta_0:C[X,Y]=Z_0\longrightarrow T_1,
 \qquad \delta_1:Z_1\longrightarrow T_2,\\
 K_2=T_2/\operatorname{im}\delta_1.
\end{gathered}                                           \label{equation-55}
\end{equation}
To specify these maps completely, choose the functorial free
resolution $P_\bullet(N)$ which is free on the elements of $N$ in
degree zero and free on the elements of each successive kernel in the
next degree. Put $P=P_\bullet(\mathcal F)$, $Q=P_\bullet(E)$ and
$u=P_\bullet(\varepsilon)$. Form the three vertical chain complexes
\[
\begin{split}
 \mathcal T_0&=\operatorname{Tot}(P\otimes_R P),\\
 \mathcal T_1&=\operatorname{Tot}(P\otimes_R Q)
                    \oplus\operatorname{Tot}(Q\otimes_R P),\\
 \mathcal T_2&=\operatorname{Tot}(Q\otimes_R Q),
\end{split}
\qquad
 h_0=(1\otimes u,u\otimes1),\quad h_1=(u\otimes1,-1\otimes u).
\]
On a summand of bidegree $(a,b)$ the vertical differential is
$d=d_1\otimes1+(-1)^a1\otimes d_2$. For $q=0,1$, represent a class
$[z]\in Z_q$ by a vertical cycle $z\in\mathcal T_{0,q}$, choose
$w\in\mathcal T_{1,q+1}$ with $dw=(-1)^{q+1}h_0z$, and define
\[
 \delta_q[z]=[(-1)^{q+1}h_1w]
       \pmod{\operatorname{im}b_{q+1}}\in T_{q+1}.
\]
The proof verifies that this is independent of the choices.

The two polynomial tensors exist uniquely if and only if
\eqref{equation-51} holds and
\begin{equation}
\begin{gathered}
                         K_1=K_2=0,\\
 \delta_0\left(\sum_{i,j\ge0}\binom{i+j}{i}
                        a_{i+j}X^iY^j\right)=0,\\
 \delta_0\left(\sum_{i\ge0}a_iX^iY^i\right)=0            \label{equation-56}
\end{gathered}
\end{equation}
for every finite coefficient vector $(a_i)$ in $C$ satisfying
$\sum_i a_ir^i\in R$ for every $r\in R$. Missing coefficients are zero.
All quantities in this test are defined before forming any tensor
power of $\Omega(R,S)$; in particular it is not a test that presupposes
the image or kernel of that tensor map. Its free-module implementation
below uses only finite coefficient relations in homological degrees
at most three. The criterion does not assert a uniformly terminating
algorithm for arbitrary presentations of infinite rings, and it
does not replace the separate higher-coherence conditions for a
plethory by a binary test.
\end{theorem}

\subsection{Noetherian bases and arbitrary coefficients}
For a Noetherian base, binary tensor injectivity already determines
all polynomial tensor maps. Here $\Omega_n(R,S)$ consists of the
polynomials over $S$ taking $R^n$ into $R$, and $\theta_n$ multiplies
one-variable factors in distinct variables. A compatible plethory
has the standard polynomial coaddition, comultiplication, scalar
operations and associative substitution, with identity $X$.
\begin{theorem}[Binary injectivity over a Noetherian base]
\label{thm:arbitrary-noetherian-coefficients}
Let $R$ be a nonzero Noetherian commutative ring and $R\subset S$ a
unital commutative inclusion. For every maximal ideal $\mathfrak m$
with finite residue field put
\[
 C_{\mathfrak m}=\bigcup_{j\ge0}
 (R_{\mathfrak m}:_{S_{\mathfrak m}}(\mathfrak mR_{\mathfrak m})^j).
\]
The following conditions are equivalent.
\begin{enumerate}[label=\textup{(\roman*)}]
\item The binary map $\theta_2$ is injective.
\item Every map $\theta_n$ is an isomorphism, for $n\ge1$.
\item For every $F\in\Omega_1(R,S)$, each of $F(X+Y)$ and $F(XY)$
      has exactly one preimage under $\theta_2$, including when $F=0$.
\item At every finite-residue maximal ideal $\mathfrak m$, either
      $C_{\mathfrak m}=R_{\mathfrak m}$, or
      $R_{\mathfrak m}$ is a one-dimensional domain whose completion
      is a domain and
      \[
       C_{\mathfrak m}=\operatorname{Frac}(R_{\mathfrak m})
                        \subseteq S_{\mathfrak m},
      \]
      with its canonical fraction-field inclusion.
\end{enumerate}
No condition is imposed at maximal ideals with infinite residue
field. In these equivalent cases the compatible plethory is unique.
\end{theorem}
Thus requiring each bivariate polynomial to have at most one tensor
preimage forces existence and all higher tensor identities. Only coefficients
supported at finite-residue closed points can affect this answer.
Completion reduces the question to a saturated ring epimorphism.
Ferrand's complete-local epimorphism theorem separates its closed
fiber, and explicit tensor detectors exclude additional components.
The nilpotent detector has unbounded parameter-torsion order and
works in every characteristic, including when the residue
characteristic is a zero divisor. The proof appears in
Section~\ref{sec:binary-injective-classification}.

\subsection{The rational-overring specialization}
For the structural specialization, a torsion-free extension
$R\subset S$ of a domain embeds in $K\otimes_R S$, where
$K=\operatorname{Frac}(R)$. Identify its rational coefficient overring
\[
 S_{\mathrm{rat}}=S\cap(K\cdot1)\subset K
\]
inside that algebra. For a nonfield Noetherian local domain
$(R_{\mathfrak m},\mathfrak mR_{\mathfrak m})$, set
\[
 R_{\mathfrak m}^{\sharp}
 =\{x\in K:(\mathfrak mR_{\mathfrak m})^N x
                       \subset R_{\mathfrak m}\text{ for some }N\ge0\}.
\]
The hat denotes completion in the maximal-ideal topology. The phrase
\emph{finite local normalization with one branch} means that the
integral closure in $K$ is a finite module and is local. A compatible
plethory on $A$ means a biring whose coaddition and comultiplication
realize the displayed polynomial substitutions, together with
associative polynomial composition and identity $X$, respecting these
ring operations. Its uniqueness concerns this specified compatibility.
\begin{theorem}[Binary injectivity and the complete Noetherian-domain criterion]\label{result-39}
Let $R$ be a nonfield Noetherian domain with
fraction field $K$, and let $R\subset S$ be an
$R$-torsion-free unital extension. Put
$L=S_{\mathrm{rat}}$ and $\Omega=\Omega_1(R,S)$.
The following conditions are equivalent.
\begin{enumerate}[label=(\roman*)]
\item The binary map
$\theta_2:\Omega\otimes_R\Omega\to\Omega_2(R,S)$ is injective.
\item Every map $\theta_n:\Omega^{\otimes_R n}\to\Omega_n(R,S)$
is an isomorphism, for $n\ge1$.
\item For every $f\in\Omega$, each of $f(X+Y)$ and $f(XY)$
has exactly one preimage under $\theta_2$.
\item At every maximal ideal $\mathfrak m$ with finite residue field,
the following local conditions hold:
\begin{equation}
\begin{gathered}
 \dim R_{\mathfrak m}=1:\quad
 L_{\mathfrak m}=R_{\mathfrak m}\quad\hbox{or}\\
 \bigl(L_{\mathfrak m}=K\text{ and }
        \widehat{R_{\mathfrak m}}\text{ is a domain}\bigr);\\[2mm]
 \dim R_{\mathfrak m}\ge2:\quad
 L_{\mathfrak m}\cap R_{\mathfrak m}^{\sharp}=R_{\mathfrak m}.
\end{gathered}                                                   \label{equation-79}
\end{equation}
\end{enumerate}
No hypothesis of finite normalization or uniform bound on
Krull dimension is imposed, and there is no condition at
infinite-residue-field maximal ideals. In these equivalent
cases the compatible plethory exists and is unique.
For $S=K$, the dimension-one clause of (iv) says finite local
normalization with one branch; its higher-dimensional clause
is exactly $\operatorname{depth}R_{\mathfrak m}\ge2$.
\end{theorem}
Thus binary injectivity alone supplies existence of the two
operations and all higher tensor isomorphisms under the stated
hypotheses. These conditions express the preceding theorem in terms of a rational
overring and the punctured-spectrum coefficient ring.

\subsection{Freeness over analytically irreducible local domains}
For a domain $R$ with fraction field $K$, the notation
$\operatorname{Int}(R)$ means the subring of $K[X]$ consisting of
polynomials taking $R$ into $R$. A Noetherian local domain is
\emph{analytically irreducible} when its completion in the maximal
ideal is a domain.
\begin{theorem}[Freeness and polynomial operations in one branch]
\label{thm:int-analytic-free}
Let $(R,\mathfrak m)$ be a one-dimensional analytically irreducible
Noetherian local domain, with fraction field $K$. Then
$\operatorname{Int}(R)$ is a free $R$-module of countably infinite
rank. For every integer $n\geq1$, multiplication in disjoint variables
induces an isomorphism
\[
 \operatorname{Int}(R)^{\otimes_R n}
 \xrightarrow{\ \sim\ }
 \{F\in K[X_1,\ldots,X_n]:F(R^n)\subset R\}.
\]
Polynomial addition, multiplication, and composition give
$\operatorname{Int}(R)$ a unique compatible $R$-plethory structure.
In particular it is flat over $R$. If $R/\mathfrak m$ is infinite,
then $\operatorname{Int}(R)=R[X]$.
\end{theorem}
Elliott conjectured that some ring in precisely this class has a
nonflat integer-valued polynomial module
\cite[Conjecture~2.10, author version p.~11]{Elliott2011IntegerTClosure}.
The theorem disproves that conjecture and the two proposed instances
\cite[Conjectures~2.15 and 2.17, author version p.~12]{Elliott2011IntegerTClosure}:
$\mathbb F_2+t\mathbb F_4[[t]]$ and $\mathbb F_2[[t^2,t^3]]$.
Elliott subsequently announced that the cusp example is free in his
26 December 2013 addition to \cite{Elliott2013CuspAnnouncement}.
That announcement states the cusp conclusion without presenting its
proof; the general theorem above and its proof below also encompass
the residue-extension example. Proposition~\ref{prop:elliott-explicit-witnesses} gives explicit
polynomials exhibiting $X^2+X\in\mathfrak m\operatorname{Int}(R)$
for both rings. The freeness here concerns arbitrary module bases;
Lemma~\ref{lem:cusp-no-regular-basis} shows that the cusp ring has
no basis with one polynomial of each degree.

Problems~19--21 of Cahen, Fontana, Frisch and Glaz ask about flatness
and freeness, multivariable tensor maps, and compatible birings
\cite[authors' version pp.~9--10]{CFFG2014OpenProblems}.
Theorem~\ref{thm:int-analytic-free} gives positive answers throughout
the stated class. The passage from freeness to tensor isomorphisms
and a biring is also part of Elliott's earlier work
\cite[Theorem~4 and Corollary~5, p.~54]{Elliott2010}; the full argument
used here appears in Theorem~\ref{result-24}.
For general one-dimensional local domains the answer is different.
Ma constructs
$R=\mathbb F_2+t(t+1)T$, where $T$ is the localization of
$\mathbb F_2[t]$ away from $(t)\cup(t+1)$, and proves that
$\operatorname{Int}(R)$ is not flat
\cite[Lemma~3.1 and Proposition~4.1, pp.~2--5]{Ma2026}.
For this same ring Jiang, Li, Sun, Wang, Xiao and Yu prove that
the binary tensor map is not surjective and that no compatible
biring exists \cite[Theorem~3.4, pp.~7--8]{Rethlas2026}.
Its normalization has two maximal ideals. Our theorem treats
finite local normalization with one branch, which is equivalent
to analytic irreducibility in dimension one.

The mechanism is described and proved in
Section~\ref{sec:integer-free}. Finite normalization supplies a
conductor ideal. A valuation-ordering basis handles polynomials
valued in the normalization, while approximation of residue-class
indicators makes the reduction modulo the conductor free.
Milnor patching then gives projectivity, and Kaplansky's theorem
turns projectivity over the local ring into freeness. The
normalization, approximation, patching, and projective-module
inputs are classical and are credited separately there.

The algebraic question comes from cote~153, sheet~4.
Elliott's sufficient tensor criteria \cite[Theorem~12, pp.~56--57]{Elliott2010}
and the recent domain counterexamples \cite[\S3.4]{Rethlas2026}
\cite[\S\S2--3]{Ma2026} provide part of its context. The coefficient-and-Tor criterion
also allows nilpotents and torsion in the base. It is a binary existence
and uniqueness test; the higher compatibility in the Noetherian theorem
follows separately from all the multivariable tensor isomorphisms.
The intrinsic-partition theorem is compared below with the classical
intrinsic CS construction and Friedman's marked version, with their
different hypotheses and the resulting difference in local structure.

\begingroup
\small
\setlength{\parskip}{0pt}
\tableofcontents
\endgroup
\clearpage
\part{Collars, cones, and intrinsic partitions}
\section{Collars and the concordance obstruction}
We use the strict collar, cone and family conventions fixed in Section~1.
\begin{proof}[Proof of Theorem~\ref{result-1}]
Every $h(b,-)$ is strictly increasing. The formula
$h_s(b,t)=(1-s)h(b,t)+st$ preserves strict increase, endpoints,
continuity and definability, and is stationary at the identity. It
contracts each family simultaneously. The inverse is continuous:
the total map over the compact parameter and base is a continuous
definable bijection between definably compact sets, hence a homeomorphism.

Realize $CY$ as the cone of a bounded embedding of $Y$, with radius
$r\in[0,1]$. For $g\in\Homeo(CY;Y,o)$ put, for $s>0$,
\[
 A_s(g)(y,r)=
 \begin{cases}
 s\,g(y,r/s),&0\leq r\leq s,\\
 (y,r),&s\leq r\leq1,
 \end{cases}
 \qquad A_0(g)=\id.
\]
Multiplication by $s$ here is radial dilation on the embedded cone.
The two formulas agree at $r=s$ because $g$ fixes the base. The inverse
is $A_s(g^{-1})$. At $s=0$ both image and domain of the nonidentity
part lie in the cone of radius $s$, proving continuity, uniformly in
every compact parameter family. The vertex is fixed for every $s$.
This is the Alexander contraction. For $Y=\varnothing$ we define
$C\varnothing=\{o\}$, and the assertion is immediate.
\end{proof}

\paragraph{Parametrizations of one cone.}\label{result-cone-parametrizations}
The base condition can be described exactly. Let
$\operatorname{Par}_{\rm set}(C,Y,o)$ be the simplicial space of all
parametrizations of a fixed cone that carry the base onto $Y$ and the
vertex to $o$, and fix one base-pointwise parametrization
$\psi_0:CY\to C$. Put
$G_Y=\operatorname{Homeo}_{\rm def}(CY;Y\text{ pointwise},o)$.
There is an isomorphism of simplicial sets
\[
 \operatorname{Par}_{\rm set}(C,Y,o)
 \cong G_Y\times\operatorname{Homeo}_{\rm def}(Y).
\]
Indeed, for a parametrization $\psi$, let
$h=\psi_0^{-1}\psi$ and $a=h|_Y$. Then
$h=k\,Ca$ for the unique $k=h\,C(a^{-1})\in G_Y$.
Conversely $(k,a)$ gives $\psi=\psi_0k\,Ca$.
These formulas are inverse on total parameter families and commute
with faces and degeneracies. The contraction just proved contracts
$G_Y$ while leaving $a$ fixed. Thus arbitrary base-setwise
parametrizations have the homotopy type of
$\operatorname{Homeo}_{\rm def}(Y)$, whereas those fixing the base
pointwise form a contractible space.
For example, take a smaller centered interval inside a larger interval
as a strict cone on its two endpoints $Y=S^0$. Identity and reflection
belong to different components of its base-setwise parametrization
space, since their restrictions permute the two base points
 differently. The base-fixed convention is the one used throughout
the rigid cone category. Its morphisms are not required to fix the
outer faces pointwise.

\begin{proof}[Proof of Theorem~\ref{result-2}]
We give the parameter-dependent argument, including the map that makes
the comparison. Fix one strict collar $c$. Let $d_p$, $p\in P$, be any
family, where $P$ is a finite compact parameter polyhedron. Uniform
definable continuity at $B\times\{0\}\times P$ gives
$\epsilon\in K$, $0<\epsilon<1$, such that
\[
 d_p(B\times[0,\epsilon])\subset c(B\times[0,1))
 \quad\text{for every }p.
\]
Indeed the inverse image of the right-hand open set contains the
definably compact bottom face; its complement has positive distance
from that face. Put $k_p=c^{-1}d_{p,\epsilon}$, where
$d_{p,\epsilon}(b,t)=d_p(b,\epsilon t)$. Thus $k_p$ is a collar
embedding into $B\times[0,1]$, with image disjoint from the top.

Here is the finite formula behind Armstrong's collar comparison
\cite[Theorem 2, pp.~123--124]{Armstrong1970}. For $0\leq s\leq1$,
let $F_s:B\times[0,1]\to B\times[-s,1]$ send
$(b,t)$ to $(b,(1+s)t-s)$. Define $G_{p,s}$ on the latter cylinder by
\[
\begin{aligned}
 G_{p,s}(b,u)&=k_p\left(b,\frac{u+s}{1+s}\right)
 &&(-s\leq u\leq0),\\
 G_{p,s}(k_p(b,v))&=k_p\left(b,\frac{v+s}{1+s}\right)
 &&(0\leq v\leq1),
\end{aligned}
\]
and as the identity outside $k_p(B\times[0,1])$ in the nonnegative
cylinder. The formulas agree on their common boundary. Their images
partition the small collar into radii $[0,s/(1+s)]$ and
$[s/(1+s),1]$, so $G_{p,s}$ is a homeomorphism. Consequently
$K_{p,s}=G_{p,s}F_s$ is a cylinder homeomorphism fixing both ends,
$K_{p,0}=\id$, and
\[
 K_{p,1}(b,t)=k_p(b,t)\quad(0\leq t\leq1/2).
\]
All these maps and inverses are definable continuous families. Glue
$cK_{p,s}c^{-1}$ to the identity on $X\setminus c(B\times[0,1))$.
The top boundary is fixed, so this is an ambient family of homeomorphisms.
It carries $c_{1/2}$ to $d_{p,\epsilon/2}$. Concatenate the shrinking
path from $c$ to $c_{1/2}$, this ambient path, and the expansion from
$d_{p,\epsilon/2}$ to $d_p$. Every intermediate embedding is still a
strict collar, by its extension or by the ambient homeomorphism.
Thus every finite parameter family is homotopic to a constant one.

A finite simplicial sphere of choices is such a family, by gluing its
finitely many defining formulas. Hence all homotopy groups of the
realization vanish and it is nonempty. A simplicial realization has
CW homotopy type, so Whitehead's theorem makes it contractible.
The proof also covers $B=\varnothing$, when the choice is unique.
\end{proof}

\begin{proof}[Proof of Theorem~\ref{result-3}]
We use the corrected smooth pseudoisotopy sequence
\[
 \operatorname{Wh}_3(\pi)\longrightarrow
 \operatorname{Wh}_1^+(\pi;\Z/2)\longrightarrow
 \pi_0P_{\rm Diff}(B)\longrightarrow
 \operatorname{Wh}_2(\pi)\longrightarrow0,
 \qquad \pi=\pi_1B,
\]
when $\dim B\geq6$ and $\pi_2B=0$
\cite[Theorem~8.a.2, pp.~149--150]{Igusa1984}.
The aspherical form and the reduced Whitehead terms are also stated in
\cite[Theorem~7.1, p.~23; Remark~7.3, p.~26]{Igusa2025}.
In the definition on p.~2 of the latter source,
$\operatorname{Wh}_3(\pi)$ is $K_3(\Z[\pi])$ modulo the constant
$K_3(\Z)$ and the framed assembly image from
$\Omega^{\rm fr}_3(B\pi)$. We compute the third Whitehead term for the torus.

The Laurent fundamental theorem for regular rings, including its
splitting by multiplication by the unit classes $[t_i]$, gives
\[
\begin{split}
K_3(\Z[\Z^d])\cong{}&K_3(\Z)
 \oplus\bigoplus_iK_2(\Z)[t_i]
 \oplus\bigoplus_{i<j}K_1(\Z)[t_i][t_j]\\
 &\oplus\bigoplus_{i<j<k}K_0(\Z)[t_i][t_j][t_k].
\end{split}
\]
There are no Nil terms, by regularity
\cite[Theorem~6.3, Corollary~6.3.1, p.~40; Theorem~8.2, p.~61]{WeibelV}.
The unit from stable homotopy to algebraic $K$-theory is
multiplicative and sends $\eta$ to $[-1]$
\cite[Exercise~1.12, pp.~15--16; Corollary~1.20, p.~13]{WeibelIV}.
Hence $\eta^2$ maps to $\{-1,-1\}$, the generator of
$K_2(\Z)=\Z/2$ \cite[Example~5.2.2, p.~36]{WeibelIII}.
The classes $\eta^2$ times each coordinate circle, $\eta$ times
each coordinate two-torus, and the framed coordinate three-tori
therefore hit every nonconstant summand. The constant summand
is quotiented out separately. Thus
$\operatorname{Wh}_3(\Z^d)=0$. In degree two the identical
Laurent decomposition has only the constant $K_2$, the unit
$K_1[t_i]$ summands and $K_0[t_i][t_j]$ summands; their
monomial generators are precisely the defining relations in
$\operatorname{Wh}_2$, so that group is also zero.
Because $\Z^d$ acts trivially by conjugation, the middle
Whitehead term is $\F_2[\Z^d]/\F_2[0]$. Exactness proves the
smooth formula without any splitting assertion.

For the category comparison, Burghelea--Lashof's
\cite[Theorem~6.1, p.~41; proof pp.~42--43]{BurgheleaLashof1974II}
states directly that smooth-to-PL concordances induce an
isomorphism on components for a compact connected smooth
manifold. Its proof identifies the relative smoothing space
with a section space and proves its components and fundamental
group vanish. This is an unstable comparison, exactly the map
we need for the smooth torus with its differentiable
triangulation. Theorem~6.2 of the same source supplies the
PL--TOP comparison in the safe range $d\geq6$ adopted here.
Both fix the bottom and the side boundary; the latter is empty
for the torus. Thus no stability inference is needed.
The torus formula is also explicitly recorded in
\cite[\S6.2, p.~34, equation~(6.12), Lemmas~6.16--6.17]{CourtePorcelli2025};
we do not use those stable statements as a substitute for the
unstable comparison.
The nonconstant monomial $t_1$ is nonzero in the displayed
additive quotient, proving nontriviality.

For the definable comparison use the relative homeomorphism and isotopy
arguments in Shiota \cite[Theorem 2.1 and its complement; Lemma 3.3,
proof of injectivity, pp.~44--46]{Shiota2010}. They give two precise
facts for a compact polyhedral pair: every definable self-homeomorphism
fixed on a subpolyhedron is definably isotopic, relative to it, to a PL
one; and two such PL homeomorphisms joined by a definable isotopy are
joined by a PL isotopy. Apply them to
$(B\times[0,1],B\times\{0\})$, keeping the other boundary component
as a distinguished subpolyhedron. Shiota's Lemma 1.2 identifies the
PL isotopy classes over different ordered fields. Thus inclusion of PL
concordances induces a bijection on components, compatible with
composition. We are using a component theorem, not asserting from
these locators alone a comparison of all higher homotopy groups.
\end{proof}

\begin{theorem}[The general component comparison and the PL--TOP comparison]\label{result-4}
Let $B$ be any closed definable manifold and choose its finite PL
model, with the closed-manifold convention including definable
compactness. Inclusion gives a natural isomorphism
\[
 \pi_0P_{\rm PL}(B)\xrightarrow{\sim}\pi_0P_{\rm def}(B).
\]
These component groups are independent of the ordered field used
to realize the finite PL model. Over the real field, if every
component of $B$ has dimension at least five, the map of
simplicial groups $P_{\rm PL}(B)\to P_{\rm TOP}(B)$ is a weak
equivalence. Thus for general $K$ the definable component group
also agrees with the TOP component group of a chosen real
realization of the finite PL model in that range. We do not claim
a canonical map from a non-Archimedean realization to the real
TOP space.
The last statement does not assert a weak equivalence with the
whole definable concordance group.
\end{theorem}
\begin{proof}
The two relative facts applied in the proof of Theorem~\ref{result-3} did not use
the torus. Apply the same relative homeomorphism straightening
and relative isotopy theorem of \cite[Theorem 2.1 with complement;
Lemma 3.3, pp.~44--46]{Shiota2010} to the compact polyhedron
$B\times I$, fixing $B\times\{0\}$ pointwise and preserving the
other boundary component. The first theorem replaces every
definable concordance by a PL concordance through such maps,
proving surjectivity. The relative injectivity argument takes an
isotopy between two PL maps, triangulates the images of finitely
many simplex cylinders while preserving the single isotopy
coordinate, and uses the level-preserving Alexander extension
on each simplex cylinder, fixed on its already treated boundary.
It gives a PL isotopy with the same two endpoints and the same
fixed subpolyhedron. This proves injectivity. Composition is
unchanged by the inclusion, so the bijection is a group
isomorphism. Shiota's Lemma 1.2 transports these PL isotopy
classes between real closed fields.

Burghelea--Lashof's Theorem 6.2, printed pp.~41--43, proves
this for compact PL manifolds of dimension at least six
\cite{BurgheleaLashof1974II}; its original range is $n>5$.
Its concordances fix $\partial B\times I\cup B\times\{0\}$,
exactly our condition because $\partial B$ is empty.

For the dimension-five endpoint, we use the relative sliced
classification, retaining the free upper boundary. We state the
classical inputs precisely. For a compact PL manifold $N$ whose
interior and boundary dimensions both differ from four, the simplicial
space $\mathcal S_{\rm PL}(N\operatorname{rel} C)$ of sliced PL
structures equal to the given structure near a closed subset $C$
is weakly equivalent to the space of compatible lifts of its interior
and boundary tangent classifiers to
$B\mathrm{PL}(\dim N)$ and $B\mathrm{PL}(\dim N-1)$, relative to $C$.
The compatibility uses stabilization along a boundary collar.
Moreover the orbit map of TOP automorphisms has PL automorphisms as
fibre and identifies their simplicial quotient with the component of
the given structure. These are the relative fibration (1.5), printed
p.~229, and the relative boundary classification (4.3)--(4.5), printed
p.~244, of \cite[Essay V]{KirbySiebenmann1977}. Their sliced hypothesis
means that projection to each parameter simplex is a PL submersion;
the relative bundle theorem, or Essay II, Theorem 2.1, printed
pp.~73--74, gives the required horn lifting. For compact $N$, TOP
simplices are precisely continuous compact-open families, and PL
simplices are jointly PL families, so these are our models. Finally
the canonical
stabilization maps
\[
 \mathrm{TOP}_j/\mathrm{PL}_j\longrightarrow\mathrm{TOP}/\mathrm{PL}
       \simeq K(\mathbb Z/2,3)\qquad(j\ge5)
\]
are weak equivalences: Essay V, \S5.1--5.2 and Theorem 5.5(I), printed
pp.~246--247 and 251--252. Thus these inputs concern whole parameter
spaces, not merely concordance classes of individual structures.

Apply them to $N=B\times I$ and $C=B\times\{0\}$, where
$\dim B=5$. The interior dimension is six and both boundary components
have dimension five, so every displayed dimension restriction holds.
The tangent bundle of $N$ is the pullback of $TB\oplus\mathbb R$.
The space of its $B\mathrm{PL}(6)$ lifts relative to the fixed bottom
germ is weakly contractible: restriction to that bottom germ is a
fibration and a weak equivalence, since the germ inclusion is a
homotopy equivalence. Equivalently, contract the cylinder coordinate
to zero in the associated lifting problem. Over the upper boundary
one must also choose a $B\mathrm{PL}(5)$ lift compatible with the
six-dimensional lift. The homotopy fibre of this additional choice is
contractible, because
$\mathrm{TOP}_5/\mathrm{PL}_5\to\mathrm{TOP}_6/\mathrm{PL}_6$
is a weak equivalence. The relative version of the fibre-product
square (4.3) therefore makes
$\mathcal S_{\rm PL}(N\operatorname{rel} C)$ weakly contractible.
There is no prescribed PL structure at the upper end. Prescribing
one would change this lifting problem and would invalidate the
cylinder contraction used here. The orbit fibration now proves that
the inclusion of PL into TOP automorphisms fixing a neighborhood of
the bottom is a weak equivalence, including components.

To pass from a fixed bottom germ to the pointwise condition in our
concordance groups, no collar approximation theorem is needed. For
$0\le s\le1/4$, put
\[
 \lambda_s(t)=\min\{t+s,(1+t)/2\}:I\xrightarrow{\sim}[s,1].
\]
For a bottom-fixed homeomorphism $h$ of $B\times I$, let $h_s$ be
the identity below height $s$ and, above it, the conjugate of $h$ by
$\operatorname{id}_B\times\lambda_s$. The two formulas agree at height
$s$, because $h$ fixes the bottom. This construction is jointly
continuous, starts at $h$, and at $s=1/4$ fixes a whole bottom collar.
The map $(s,t)\mapsto(s,\lambda_s(t))$ and its inverse are PL on finite
triangulations, so the same construction preserves jointly PL
parameter families. It preserves the identity family and any already
fixed bottom collar, since $\lambda_s(t)\ge t$. Consequently inclusion
of the bottom-germ subgroup into the pointwise bottom-fixed group
is a weak equivalence in both categories, with the same assertion
relative to parameter faces. This proves the dimension-five endpoint.
Finite disjoint unions give finite products, because fixing the bottom
prevents permutation of components. This finishes the comparison.

\end{proof}

\begin{theorem}[I c) and the proposed collapse equivalence]\label{result-5}
For $B=(S^1)^6$, the automorphism space in cote~119, sheet~147 (PDF p.~148) I c) is not
contractible. Moreover the boundary-collapse functor from bordant pairs
to pointed spaces is not a Dwyer--Kan equivalence in the stipulated
definable simplicial categories. Its map on automorphism spaces at
$(B\times[0,1],B\times\{0\})$ has homotopy fibre equivalent to
$P_{\mathrm{def}}(B)$. Thus collar rigidification alone does not make boundary collapse an equivalence.
\end{theorem}
\begin{proof}
The space in I c) is precisely $P_{\mathrm{def}}(B)$: only its lower
boundary is fixed pointwise, and its upper boundary is free.
Theorem~\ref{result-3} proves that its component group is nonzero.

Write $I=[0,1]$ and put
\[
 G=\Homeo_{\mathrm{def}}(B\times I;B\times\{0\}\text{ setwise}),
 \qquad H=\Homeo_{\mathrm{def}}(CB;o).
\]
The upper boundary of the cylinder and the base of the cone are
intrinsic boundaries of their manifold loci. The cone vertex is the
specified point. Hence upper-boundary restriction
$e:G\to\Homeo_{\mathrm{def}}(B)$ and base restriction
$r:H\to\Homeo_{\mathrm{def}}(B)$ are defined. Both have sections,
by constant extension and coning, respectively.
The kernel of $e$ is the group of cylinder homeomorphisms fixed on the
upper boundary; interval reversal identifies it with $P_{\mathrm{def}}(B)$.
Multiplying by the constant extension of $e(g)^{-1}$ gives an actual
product decomposition of the underlying simplicial spaces of $G$.
Similarly the kernel of $r$ consists of cone homeomorphisms fixed on
base and vertex and contracts by Theorem~\ref{result-1}. The same decomposition makes
$r$ a homotopy equivalence.

The quotient $q:B\times I\to CB$ collapses the lower boundary.
Every $g\in G$ descends uniquely to $\bar g\in H$, and
\[
 r(\bar g)=e(g).
\]
Thus the homotopy fibre of $G\to H$ is that of $e$, namely
$P_{\mathrm{def}}(B)$. By Theorem~\ref{result-3} it is not contractible. The induced
map on this mapping space is not a weak equivalence, which disproves
the required enriched equivalence.
Equivalently choose any concordance representing $t_1$ in Theorem~\ref{result-3} and
reverse its interval. It is nontrivial in $\pi_0G$, by the split
upper-boundary restriction, whereas its descended cone automorphism
lies in the Alexander-contractible kernel of $r$.
We do not assert that every bottom-fixed concordance descends to that
kernel: its upper endpoint need not be the identity.
\end{proof}

\subsection{Countable presentations of concordance component groups}

For a compact manifold $M$, let $P_{\mathrm{CAT}}(M)$ denote
concordances of $M\times I$ fixing
$M\times\{0\}\cup\partial M\times I$ pointwise, with the upper
face free.  Smooth corners are treated by product collars or compatible
rounding.  PL components use jointly PL families and jointly PL inverses.
The group law is composition, and stabilization is positive suspension.

We first specify the presentation independently of any manifold.  Let
$K$ be a connected finite simplicial complex.  Collapse a maximal tree
of its one-skeleton to obtain a reduced finite simplicial set $X$, let
$G(X)$ be its Kan loop group, and put $R=\Sigma^\infty_+G(X)$.
Define
\[
 H_n=\operatorname{hAut}_R(R^{\vee n}),\qquad
 W_n=G(X)^n\rtimes\Sigma_n.
\]
Here $H_n$ is the derived simplicial monoid of homotopy automorphisms
of a free right module over the spherical group ring.  In particular,
it is not a discrete general linear group over $\mathbb Z[\pi_1K]$.
Left multiplication and permutation of free summands define
$W_n\to H_n$.  Adjoining an identity summand gives countable simplicial
telescope models
\[
 C=\operatorname*{hocolim}_n BH_n,\qquad
 D=\operatorname*{hocolim}_n BW_n,\qquad f:D\longrightarrow C.
\]
The existence of fully countable models is verified below.

For a connected simplicial model $Y$, write $Y^{\mathrm{ab}}$ for the
connected cover corresponding to
$\ker(\pi_1Y\to H_1(Y;\mathbb Z))$.  This cover has an explicit
simplicial description.  In a one-vertex model put
$L_Y=C_1(Y;\mathbb Z)/\partial C_2(Y;\mathbb Z)$.  Its simplices
are pairs $(a,\sigma)$ with $a\in L_Y$ and $\sigma\in Y_q$, with
\[
 d_0(a,\sigma)=(a+[\sigma|[0,1]],d_0\sigma),
 \qquad d_i(a,\sigma)=(a,d_i\sigma)\quad(i>0).
\]
Degeneracies leave $a$ unchanged.  The edge relation
$[01]+[12]=[02]$ verifies the first-face identities; the remaining
identities are those of $Y$.  If a telescope has several vertices,
choose paths in a maximal tree from a base vertex and use the
corresponding edge labels.  This gives the same cover without needing
to collapse the tree functorially.  The based map $f$ has an induced
based map $f^{\mathrm{ab}}:D^{\mathrm{ab}}\to C^{\mathrm{ab}}$.

All chains in the following definition are normalized integral
simplicial chains with finite support.  Set
\begin{equation}\label{eq:countable-concordance-presentation}
 E(K)=
 \frac{Z_2(C^{\mathrm{ab}})}
 {\partial C_3(C^{\mathrm{ab}})
        +(f^{\mathrm{ab}})_*Z_2(D^{\mathrm{ab}})}.
\end{equation}
Equivalently, take one generator $[z]$ for every two-cycle in
$C^{\mathrm{ab}}$, impose $[z+z']=[z]+[z']$, kill every three-boundary,
and kill the image of every two-cycle from $D^{\mathrm{ab}}$.
These are all the relations.  For a finite disjoint union of complexes,
define $E$ to be the product of the component presentations.

\begin{theorem}\label{thm:countable-stable-smooth-components}
For every compact smooth manifold $M$ with finite simplicial homotopy
model $K$, the group $E(K)$ presents the stable smooth concordance
component group, including its composition law:
\[
 \operatorname*{colim}_j
 \pi_0P_{\mathrm{Diff}}(M\times I^j)\cong E(K).
\]
If each component of $M$ has dimension at least nine, it also presents
$\pi_0P_{\mathrm{Diff}}(M)$.  The same unstable conclusion holds for
a smoothable compact PL manifold in that range, and for the
stipulated definable concordance group of a compact definable manifold
with such a PL model, over every real closed field and every fixed
o-minimal expansion.  When a boundary is present, concordances fix the
bottom and the entire side boundary.

The presentation is countable and depends, up to group isomorphism,
only on the finite homotopy model.  No terminating word problem or
finite invariant-factor decomposition is asserted.
\end{theorem}

\begin{proof}
\emph{Countable derived models.}
The loop group has
\[
 G(X)_q=F(X_{q+1})/\langle[s_0x]=1\rangle,
\]
with faces and degeneracies on generators
\[
 d_0[x]=[d_1x][d_0x]^{-1},\qquad
 d_i[x]=[d_{i+1}x]\ (i>0),\qquad
 s_i[x]=[s_{i+1}x].
\]
These are the convention in
\cite[V.5, formula~(5.1), p.~276]{GoerssJardine1999}.
Theorem~5.10 and Corollary~5.11 there, p.~285, give the loop-space
comparison.  Each degree consists of finite words in countably many
generators, so $G(X)$ and $R$ have countable simplicial models.

Work in right $R$-modules in simplicial symmetric spectra.  Retain
finite cell presentations built from the free cells
\[
 R\wedge F_j\partial\Delta^q_+
       \longrightarrow R\wedge F_j\Delta^q_+.
\]
Maps between finite cell modules are determined by the images of
finitely many cell generators in countable simplicial sets.  Thus
both the objects with these presentations and their morphisms form
countable sets.  For acyclic cell extensions use the free $R$-module
images of the generating stable acyclic cofibrations
\[
 J=FI_\Lambda\cup\{c_j\mathbin\square i\},
\]
where $c_j:F_{j+1}S^1\to C\lambda_j$ is the mapping-cylinder
cofibration of $F_{j+1}S^1\to F_jS^0$, and $i$ runs through boundary
inclusions.  The domains and codomains are finite cell spectra;
see \cite[Theorem~3.4.4 and Definition~3.4.9, pp.~180--182]
{HoveyShipleySmith2000}.  The transferred module model structure,
applied to $R^{\mathrm{op}}$, is
\cite[Theorem~4.1(1), p.~499; proof p.~505]{SchwedeShipley2000}.
Include finite cylinders $P\wedge\Delta^1_+$ as well.

Here finite cell extensions suffice to encode the weak equivalences
and all derived maps that are needed.  Construct a fibrant replacement
of a finite cell module by attaching the stated acyclic cells in
countably many stages.  A map from a finite cell module, including its
tensor with any finite simplicial parameter set, uses finitely many
cells at a finite stage.  Each of their attaching maps uses finitely
many cells at earlier stages.  Induction on these finitely many stages
produces a finite set closed under all attaching dependencies.
The map therefore factors through a finite relative acyclic cell
extension.  The same argument applies simultaneously to a finite
collection of maps and homotopies.

If a map between finite cell modules is a stable equivalence, choose
its inverse and inverse homotopies in fibrant replacements.  The
preceding factorization realizes this inverse and its two homotopies
using only finite acyclic extensions and finite cylinders.  Conversely,
inverting such extensions and identifying cylinder homotopies inverts
only stable equivalences.  Thus this is a finite-certificate
description of the stable equivalences; it does not decide when no
such certificate exists.

Take the hammock localization of this countable category at its
stable equivalences.  Its simplices are finite diagrams, hence remain
countable.  The function-complex comparison of
\cite[Proposition~4.4 and Application~4.6, p.~432]
{DwyerKan1980Functions}, together with the simultaneous finite
factorization just proved, identifies its mapping spaces between
finite cell modules with the derived module mapping spaces.  In
particular, finite simplicial representatives of maps, homotopies
and nullhomotopies all factor through finite extensions; this checks
both surjectivity and injectivity on the homotopy groups of the
comparison.  Take the invertible components of the endomorphism
monoid of $R^{\vee n}$ to obtain $H_n$.  Invertibility has a finite
inverse diagram with paths to the two identities.  Direct sums give
the required coherent operations; a bar rectification uses finite
strings and preserves countability.  Classifying spaces, telescopes,
the abelian covers above, and their finite-support chain groups are
therefore countable models determined by $K$.

\emph{What the two-cycle quotient computes.}
Direct sum makes
$\coprod_nBH_n$ and $\coprod_nBW_n$ homotopy commutative monoids,
with commutativity given by block permutations.  Their component
monoids are $\mathbb N$, and stabilization by one identity summand
is cofinal.  Their group completions, and the map between them, are
the standard monomial-matrix models
\[
 Q(|K|_+)\longrightarrow A(|K|);
\]
see \cite[Section~2, p.~247]{Jahren2010}.
Only the rank-zero components of these group completions are used
below; free ranks are not being identified with all projective
$K_0$ classes.

The comparison from a stabilization telescope to its group completion
is acyclic here.  Indeed, the component monoid is central in homology
and has the cofinal stabilization sequence required in
\cite[Proposition~1.1 and proof, pp.~1--7]{Gritschacher2026Acyclicity}.
Consequently the maps from $C$ and $D$ to the rank-zero group-completion
components induce homology isomorphisms with all local coefficient
systems.  The targets have abelian fundamental group because they
are group-like homotopy commutative spaces.  The ordinary homology
isomorphism in degree one therefore shows that the induced
fundamental-group map is abelianization.  Pulling back their simply
connected covers gives exactly $C^{\mathrm{ab}}$ and
$D^{\mathrm{ab}}$, and the lifted maps are integral homology
isomorphisms.  Hurewicz on the simply connected \emph{targets} gives
\[
 H_2(C^{\mathrm{ab}};\mathbb Z)\cong\pi_2A(|K|),\qquad
 H_2(D^{\mathrm{ab}};\mathbb Z)\cong\pi_2Q(|K|_+).
\]
These identifications commute with the lifted monomial map.
There is no claim that the abelian covers themselves are simply
connected.  Taking the cokernel yields
\[
 E(K)\cong
 \operatorname{coker}\bigl(\pi_2Q(|K|_+)\to\pi_2A(|K|)\bigr)
 \cong\pi_2\operatorname{Wh}^{\mathrm{Diff}}(|K|).
\]
The last identification uses the split unit
$A(Y)\simeq Q(Y_+)\times\operatorname{Wh}^{\mathrm{Diff}}(Y)$,
as in \cite[pp.~247--248]{Jahren2010} and
\cite[Definition~1.3.4 and (1.3.5), p.~18]
{WaldhausenJahrenRognes2013}.

\emph{The composition law and the geometric range.}
The stable parametrized h-cobordism theorem identifies
$H^{\mathrm{st}}_{\mathrm{Diff}}(M)$ with
$\Omega\operatorname{Wh}^{\mathrm{Diff}}(M)$
\cite[Theorem~0.1, p.~2]{WaldhausenJahrenRognes2013}.
The trivial h-cobordism component is $BP_{\mathrm{Diff}}(M)$,
with the side-fixed concordance group included by a group
homomorphism and a weak equivalence into the group of automorphisms
fixed on the bottom
\cite[Remark~1.1.11 and (1.1.12), p.~12]
{WaldhausenJahrenRognes2013}.
For a simplicial group $P$, the bar two-simplex $(g,h)$ identifies
composition in $\pi_0P$ with concatenation in $\pi_1BP$.
Stabilizing h-cobordisms sends an automorphism to its product with
the identity, so it preserves this group law.  The stabilization maps
are simplicial cofibrations
\cite[Definition~1.1.3, p.~8]{WaldhausenJahrenRognes2013}, whence
\[
 \operatorname*{colim}_j\pi_0P_{\mathrm{Diff}}(M\times I^j)
 \cong\pi_1H^{\mathrm{st}}_{\mathrm{Diff}}(M)_0
 \cong\pi_2\operatorname{Wh}^{\mathrm{Diff}}(M).
\]
Every arrow is a group homomorphism.  In particular, the additive
cycle presentation records the full extension law, without choosing
a splitting of a pseudo-isotopy obstruction sequence.

Igusa's smooth stability bound is
$n\geq\max\{2k+7,3k+4\}$ for $k$-connected stabilization
\cite[Theorem~1.4.1, p.~21]{WaldhausenJahrenRognes2013}.
For $k=1$ it gives a component-group isomorphism at every stage from
dimension nine onward.  Burghelea and Lashof's smooth-to-PL component
isomorphism applies to compact manifolds with the side boundary fixed
\cite[Theorem~6.1 and proof, pp.~41--43]{BurgheleaLashof1974II}.
For a definable manifold, the proofs of Theorems~\ref{result-4}
and~\ref{result-40} apply relative to
$M\times\{0\}\cup\partial M\times I$: include this entire
subpolyhedron among the fixed cells.  Their straightening, coning,
and coefficient-transfer constructions retain these fixed cells at
each step.  They identify the PL and definable component groups and
transport the finite PL model between real closed fields.  This
proves the asserted unstable and definable conclusions.  The bottom
is fixed, so disconnected components cannot be permuted and their
groups form the stated finite product.
\end{proof}

\begin{theorem}\label{thm:stable-pl-component-presentation}
Let $M$ be any compact PL manifold and let $K$ be a finite simplicial
homotopy model.  Then, with the side boundary fixed,
\[
 \operatorname*{colim}_j\pi_0P_{\mathrm{PL}}(M\times I^j)
       \cong E(K).
\]
No smoothing or dimension hypothesis is required.  The stable
definable component group of the same finite PL model over any real
closed field has the same presentation.  This does not identify an
unstable nonsmoothable concordance group with its stable group.
\end{theorem}

\begin{proof}
In this proof use Whitehead \emph{spectra}.  The defining assembly
maps are
\[
 \mathbb S\wedge K_+\longrightarrow A(*)\wedge K_+
                \longrightarrow A(K).
\]
Taking the cofibers of this composable pair gives a cofiber sequence
\[
 \operatorname{Wh}^{\mathrm{Diff}}(*)\wedge K_+
 \longrightarrow\operatorname{Wh}^{\mathrm{Diff}}(K)
 \longrightarrow\operatorname{Wh}^{\mathrm{PL}}(K);
\]
the definitions and the factorization of the unit through assembly
are \cite[Definitions~1.3.2 and~1.3.4, pp.~16--18]
{WaldhausenJahrenRognes2013}.  This comparison does not require a
smoothing of $M$.

The connective spectrum $\operatorname{Wh}^{\mathrm{Diff}}(*)$
has zero homotopy groups through degree two.  Its degree-zero group
vanishes because the unit is an isomorphism on rank, and its
degree-one group is $\operatorname{Wh}_1(1)=0$
\cite[Theorem~0.3 and the following paragraph, p.~4]
{WaldhausenJahrenRognes2013}.  To check degree two, use a smooth
disk of dimension at least nine.  The preceding theorem identifies
that group with its stable smooth concordance component group.
Every stage is a disk after rounding corners; its smooth component
group equals its PL component group by
\cite[Theorem~6.1, pp.~41--43]{BurgheleaLashof1974II}.
The PL disk concordance space contracts by the Alexander trick,
as also used in
\cite[proof of Corollary~1.4.2, p.~23]
{WaldhausenJahrenRognes2013}.  Hence degree two vanishes as well.

The finite cell filtration of $K_+$ builds
$\operatorname{Wh}^{\mathrm{Diff}}(*)\wedge K_+$ from
nonnegative suspensions of this spectrum.  Its homotopy groups in
degrees at most two therefore vanish.  The comparison cofiber
sequence gives
\[
 \pi_2\operatorname{Wh}^{\mathrm{Diff}}(K)
       \xrightarrow{\sim}
 \pi_2\operatorname{Wh}^{\mathrm{PL}}(K).
\]
The left side is $E(K)$; the stable PL h-cobordism theorem and the
bar-space group-law argument above identify the right side with the
stable PL concordance component group.  Finally, the relative PL and
definable comparisons commute with positive suspension chosen using
fixed rational polyhedral bends.  They apply at every finite stage
and give the stable definable assertion.
\end{proof}


\subsection{Concordance components from dimension six}

The stable presentation already computes the unstable group for
smoothable manifolds from dimension six. The reason is geometric:
positive suspension preserves both the first obstruction and the
complete framed second obstruction. The corrected geometric exact
sequence then identifies the whole group, including its extension
law. This argument uses Igusa's correction to the
Hatcher--Wagoner sequence; it does not discard the framing term or
choose a splitting.

\begin{theorem}\label{thm:smoothable-concordance-six}
Let $M$ be a compact connected smoothable PL manifold of dimension
$n\geq6$, possibly with boundary, and let $K$ be a finite simplicial
complex homotopy equivalent to $M$. With the bottom and the entire
side boundary fixed pointwise and the upper face free,
\[
                     \pi_0P_{\mathrm{PL}}(M)\cong E(K).
\]
PL components mean jointly PL isotopies. The same group presents
the stipulated definable concordance components of any chosen
finite model over any real closed field and in any fixed o-minimal
expansion. Finite disjoint unions give finite products, provided
every component is smoothable and has dimension at least six.

For a compact smooth manifold of dimension at least six, positive
suspension induces a group isomorphism
\[
 \sigma^+_*:\pi_0P_{\mathrm{Diff}}(M)
       \xrightarrow{\ \cong\ }\pi_0P_{\mathrm{Diff}}(M\times I).
\]
No smoothing or isotopy is prescribed on the boundary in advance;
the concordances themselves fix the side boundary. No assertion
about higher homotopy groups is made here.
\end{theorem}

\begin{proof}
For a connected space $Z$, set
\[
 A_Z=\operatorname{Wh}_1^+
       (\pi_1Z;\mathbb Z/2\oplus\pi_2Z),\qquad
 I_Z=\operatorname{im}(\chi_Z:\operatorname{Wh}_3(\pi_1Z)\to A_Z).
\]
Here $\chi_Z$ is the full corrected obstruction homomorphism.  Its
image is unchanged if its domain is replaced by the surjecting group
$K_3(\mathbb Z[\pi_1Z])$.
For a compact smooth $n$-manifold, $n\geq6$, the corrected geometric
sequence gives
\begin{equation}\label{eq:full-geometric-concordance-sequence}
 0\longrightarrow A_M/I_M\xrightarrow{\overline\theta_M}
 \pi_0P_{\mathrm{Diff}}(M)\xrightarrow{\Sigma_M}
 \operatorname{Wh}_2(\pi_1M)\longrightarrow0.
\end{equation}
This is Igusa's corrected \emph{unstable} geometric sequence
\cite[Section~8.a and Theorem~8.a.2, pp.~149--150]{Igusa1984}.
Chapter~8 assumes $\dim M\geq6$. The map $\chi_M$ includes both
the Postnikov contribution and the Grassmann, or framing,
contribution. In particular, we take the cokernel of the full map
\[
 K_3(\mathbb Z[\pi_1M])\longrightarrow
 \operatorname{Wh}_1^+(\pi_1M;\pi_2M)
 \oplus\operatorname{Wh}_1^+(\pi_1M;\mathbb Z/2).
\]
The fixed-boundary convention is specified on printed p.~104.
The sequence is also stated in
\cite[Theorem~3.1, p.~6]{Igusa2021Dimension3}; its functorial stable
form is \cite[Corollary~2.7, p.~250]{Jahren2010}. We use the
geometric sequence, rather than infer an unstable dimension range
from the stable formula.

Projection $M\times I\to M$ identifies the fundamental groups and their
modules $\pi_2$.  Functoriality of the full $\chi$ identifies both $A$
and $I$ for these two spaces.  To apply these identifications to the
geometric sequence, we check its two obstruction maps under suspension.

Normalize concordances to be the identity near the bottom and side
boundary and products near the top.  Collar straightening induces the
same components and works for paths of concordances
\cite[Part~I, Chapter~I, Section~5, p.~55]{HatcherWagoner1973}.
In the positive-suspension construction one applies a normalized
concordance on each of a fixed collection of polar rays and uses the
identity outside them.  This is the formula in
\cite[Section~4.1, p.~8]{Igusa2021Dimension3}.  Since the ray coordinates
are fixed independently of the concordance, it gives
\[
                    \sigma^+(gh)=\sigma^+(g)\sigma^+(h).
\]
In particular it induces a homomorphism on components, with the required
boundary restrictions.

For the second obstruction choose marked lens models in indices $3,4$
in every dimension $n\geq6$.  Hatcher and Wagoner's two-index theorem
permits $2\leq i\leq n-2$ for the one-parameter representatives
\cite[Part~I, Chapter~V, Theorem~3.1(a), p.~184;
Part~II, Section~2, p.~244]{HatcherWagoner1973}.
The marked-lens calculation identifies its group with
\[
 \frac{(\mathbb Z/2\oplus\pi_2M)[\pi_1M]}
      {(\mathbb Z/2\oplus\pi_2M)[1]}.
\]
For this specialization keep $k=3$. Then
\[
                         n-k\geq k=3
\]
for every $n\geq6$, exactly the inequality used in the proof of
\cite[Lemma~1.10, p.~9]{Igusa2021SecondObstructionI}.
The theorem there is stated with $k=\lfloor n/2\rfloor$; the
fixed-index specialization follows from its proof, the preceding
two-index representatives, and the general fixed-index suspension
spaces of \cite[Section~4.2, pp.~8--9]{Igusa2021Dimension3}.
Changing markings passes to $A_M$, and forgetting the lens data gives
$\theta_M$.  Use the stable sign convention
$\lambda^s_0=(-1)^k\lambda_0$.
Positive suspension keeps $k$ fixed and preserves the complete framed
invariant, including its $\mathbb Z/2$ coordinate
\cite[Section~4.2, pp.~8--9]{Igusa2021Dimension3}.  Thus, under the
identification of the groups $A$,
\[
                  \sigma^+_*\theta_M=\theta_{M\times I}.
\]
The first obstruction is also unchanged: this is the suspension remark
in \cite[Part~I, Chapter~IV, Section~1, p.~130]{HatcherWagoner1973},
based on the trajectory formula (5.2) on p.~60.  Therefore
\[
                    \Sigma_{M\times I}\sigma^+_*=\Sigma_M.
\]

We have obtained a commutative diagram of
\eqref{eq:full-geometric-concordance-sequence} for $M$ and $M\times I$
whose outside vertical arrows are isomorphisms and whose middle arrow
is $\sigma^+_*$.  Exactness proves that middle arrow is an isomorphism.
For injectivity, a class killed by suspension has zero first obstruction,
so lies in $A_M/I_M$; injectivity of the lower left map then kills it.
For surjectivity, first match a target class's first obstruction and then
correct the residual element using $A_M/I_M$.  This uses no splitting
of the extension.

All these statements concern compact manifolds with side boundary
fixed.  Igusa explicitly uses this convention on p.~1 and fixes the
projection near the full boundary in the setup on p.~3 of
\cite{Igusa2021SecondObstructionI}.  The handle and lens modifications
occur in the interior.  Thus the suspension argument may be iterated
on $M\times I^j$, with corners rounded, and identifies the unstable
smooth group with the stable group.  The countable presentation theorem
and the classical smooth-to-PL component comparison therefore give
$E(K)$ whenever $M$ is smoothable and $n\geq6$.

For clarity, the smooth--PL component comparison preserves
\emph{composition}. Let $u$ and $v$ be the forgetful homomorphisms
from the smooth and PL concordance component groups to the TOP
concordance component group. Theorems~6.1 and~6.2 of
\cite[pp.~41--43]{BurgheleaLashof1974II} give the required
component bijections and commute with these forgetful maps.
Consequently $v^{-1}u$ is the comparison group isomorphism.
This does not require a group of piecewise differentiable maps.

The relative PL--definable component comparison applies with
$M\times\{0\}\cup\partial M\times I$ as fixed subpolyhedron.
It preserves composition and transfers the presentation to every
stipulated real closed field model. Fixing the bottom prevents
permutation of connected components, proving the finite-product
assertion.
\end{proof}

\begin{corollary}\label{cor:pl-six-seven-components}
Every compact PL manifold whose connected components have
dimension six or seven satisfies the component presentation in
Theorem~\ref{thm:smoothable-concordance-six}, without an additional
smoothability assumption. The same is true of its stipulated
definable models.
\end{corollary}

\begin{proof}
The smoothing classification uses the fibration
$BO\to BPL$ with fibre $PL/O$
\cite[Part~II, Theorem~4.1, p.~108, and Theorem~5.3,
pp.~122--123]{HirschMazur1974}. The fibre is $6$-connected
\cite[Part~II, p.~123]{HirschMazur1974}. An obstruction to extending
a lift across a $j$-cell lies in $\pi_{j-1}(PL/O)$.
All these groups vanish for $j\leq7$, so the stable PL tangent
classifying map of a manifold of dimension at most seven admits
a lift. The smoothing theorem turns this lift into a compatible
smoothing.

For a manifold with boundary, first apply the boundaryless
classification to its interior, which has the same homotopy type.
A PL collar gives a compact interior copy of the original manifold.
Its locally flat boundary may be moved to a smooth hypersurface by
\cite[Part~I, Theorem~7.4, pp.~37--38]{HirschMazur1974}.
The enclosed compact smooth domain is PL homeomorphic to the
original manifold and supplies a compatible smoothing, including
the boundary. There was no prescribed boundary smoothing to extend.
The preceding theorem now applies.
\end{proof}

\begin{lemma}\label{lem:three-skeleton-cycle-presentation}
For every connected finite simplicial complex $K$, inclusion of
the three-skeleton induces a group isomorphism
\[
                              E(K^{(3)})\cong E(K).
\]
\end{lemma}

\begin{proof}
The functorial corrected stable sequence is
\[
 0\longrightarrow A_Z/I_Z\longrightarrow
 \pi_2\operatorname{Wh}^{\mathrm{Diff}}(Z)\longrightarrow
 \operatorname{Wh}_2(\pi_1Z)\longrightarrow0,
\]
with $A_Z$ and $I_Z$ as in the proof above
\cite[Corollary~2.7, p.~250]{Jahren2010}.
The inclusion $K^{(3)}\to K$ induces isomorphisms on $\pi_1$
and on its module $\pi_2$, since the remaining cells have
dimension at least four. It therefore induces isomorphisms on
$\operatorname{Wh}_3$, $A$, and $\operatorname{Wh}_2$.
Naturality of the complete $\chi$ identifies its image subgroups
$I$; no assertion that these images are determined by abstract
groups without their Postnikov data is needed. The diagram of
short exact sequences then identifies the middle terms.
The cycle-presentation identification
$E(Z)\cong\pi_2\operatorname{Wh}^{\mathrm{Diff}}(Z)$ gives the
claimed isomorphism.
\end{proof}


\subsubsection{The extension to arbitrary PL manifolds}

We retain the convention that a concordance of a compact manifold $M$
fixes $M\times\{0\}\cup\partial M\times I$ pointwise, with the top face
free.  PL components are defined by jointly PL families with jointly PL
inverses, and their group law is composition.  For a compact manifold
$X$, write $E(X)$ for the previously defined countable presentation
formed from a finite homotopy model of $X$.

\begin{lemma}\label{lem:component-followup-E-naturality}
If $a:X\to Y$ is a based three-connected map of connected finite CW
homotopy types, the induced homomorphism
\[
                         E(a):E(X)\longrightarrow E(Y)
\]
is an isomorphism.  In particular, $E$ is invariant under homotopy
equivalence and $E(K^{(3)})\to E(K)$ is an isomorphism for every connected
finite simplicial complex $K$.
\end{lemma}

\begin{proof}
The loop-group map is two-connected.  Its map of connective spherical
group rings is therefore two-connected.  At each finite rank, extension
of scalars induces a two-connected map of derived endomorphism spaces
and an isomorphism of their degree-zero matrix rings.  It identifies
the invertible components, so the homotopy-automorphism monoids receive
a two-connected map and their bar constructions receive a
three-connected map.  The same conclusion holds for the bars of the
monomial groups: on each permutation component, the underlying map is
a finite product of the loop-group map.

These maps commute with adding an identity summand and with the
monomial inclusion.  The corresponding maps of the two telescopes are
three-connected.  They identify fundamental groups and their
abelianization kernels, and hence lift to three-connected maps of the
specified abelian covers.  In particular, the maps on integral second
homology are isomorphisms.  The presentation $E$ is the cokernel of the
monomial map on these second homology groups, so its induced map is an
isomorphism.  The argument uses no simple-connectivity assumption on
the abelian covers.

Homotopy equivalences are three-connected.  The inclusion
$K^{(3)}\to K$ is three-connected because its omitted cells have
dimensions at least four.  This proves the last assertions.
\end{proof}

\begin{lemma}\label{lem:component-followup-handle-core}
Let $M$ be a compact connected PL manifold of dimension $n\geq6$.
There is a smoothable compact codimension-zero PL submanifold
$j:N\hookrightarrow\operatorname{int}M$ having the homotopy type of a
finite CW complex of dimension at most three, such that
\[
                         \pi_i(M,N)=0\quad(i\leq3),
\]
and $\partial N\to N$ induces isomorphisms on $\pi_0$ and $\pi_1$.
The manifold $N$ is connected.
\end{lemma}

\begin{proof}
Take an absolute finite handle decomposition of $M$, relative to the
empty set, with indices ordered increasingly, and let $N_0$ be the
stage through index three.  The triangulation construction supplies
such a decomposition, with one handle of index $\dim A$ for each
simplex $A$; its associated core complex has one cell for each handle
\cite[Proposition~6.9 and the following discussion,
pp.~82--83]{RourkeSanderson1982Component}.

If $M$ has boundary, push a compact PL copy of $M$ into its interior
using a boundary collar and let $N$ be the corresponding copy of
$N_0$.  This collar embedding is homotopic to the identity as a map to
$M$.  For closed $M$ take $N=N_0$.  Adding the usual collars to handle
stages does not change their indices or their homotopy types.
Thus $N$ has the homotopy type of a finite three-complex.  The remaining
handles have indices at least four, which proves
$\pi_i(M,N)=0$ for $i\leq3$.  They cannot join components, so $N$ is
connected.

The dual handle decomposition of $N$ relative to its entire boundary
has indices at least $n-3\geq3$; see the duality construction on
page~82 of the same source.  These attachments preserve components
and fundamental groups.  Hence $\partial N\to N$ is an isomorphism on
$\pi_0,\pi_1$.

Finally, the fiber $PL/O$ of the stable smoothing lift is
6-connected.  All lifting obstructions vanish for a PL manifold with
the homotopy type of a three-dimensional CW complex.  The smoothing
theorem therefore gives a compatible smoothing
\cite[Part~II, Theorem~4.1, p.~108, and \S5.3,
pp.~122--123]{HirschMazur1974}.
For boundary, apply the lifting theorem to the interior, then use a
collar to take a compact PL copy inside that interior and move its
locally flat boundary to a smooth hypersurface
\cite[Part~I, Theorem~7.4, pp.~37--38]{HirschMazur1974}.
The resulting compact domain gives a compatible smoothing of $N$.
No smoothing of $M$ or preassigned smoothing of $\partial N$ is required.
\end{proof}

\begin{lemma}\label{lem:component-followup-stable-square}
For a codimension-zero inclusion $j:N\hookrightarrow\operatorname{int}M$,
extension by the identity induces a homomorphism
$\gamma_0:\pi_0P_{\mathrm{PL}}(N)\to\pi_0P_{\mathrm{PL}}(M)$.
If $s_X:\pi_0P_{\mathrm{PL}}(X)\to E(X)$ denotes stabilization followed
by the established stable comparison, then
\[
\begin{array}{ccc}
\pi_0P_{\mathrm{PL}}(N)&\xrightarrow{\ \gamma_0\ }&
       \pi_0P_{\mathrm{PL}}(M)\\
\downarrow s_N&&\downarrow s_M\\
E(N)&\xrightarrow{\ E(j)\ }&E(M)
\end{array}
\tag{$*$}\label{eq:component-followup-stable-square}
\]
commutes.
\end{lemma}

\begin{proof}
On a parameter simplex, extend the given map on
$\Delta^q\times N\times I$ by the identity.  The two definitions agree
on $\Delta^q\times\partial N\times I$.  Common finite subdivisions
therefore give a jointly PL map and its jointly PL inverse on the
whole product.  The parameter projection, bottom fixation, and side
fixation are retained.  Extension commutes with composition and
inversion, which proves the first assertion.

For naturality, use the collared h-cobordism model of the stable
comparison.  A codimension-zero embedding induces the covariant
pushout
\[
                    W\longmapsto M\times I\cup_{N\times I}W
\]
on collared h-cobordisms
\cite[Definition~1.1.3(b), pp.~8--9]{WaldhausenJahrenRognes2013}.
On automorphisms of the product h-cobordism, this is exactly extension
by the identity.  The product component is the classifying space of
the concordance group with the bottom-and-side convention used here;
thus its fundamental group identifies with the concordance component
group with composition
\cite[Remark~1.1.11 and (1.1.12),
pp.~12--13]{WaldhausenJahrenRognes2013}.

The stable h-cobordism comparison is natural for these maps
\cite[Theorem~0.1, p.~1, and Theorem~1.1.8,
p.~11]{WaldhausenJahrenRognes2013}.
In the retractive-space model its map is pushout along $j$; in the
spherical-group-ring model it is extension of scalars.  This preserves
the monomial maps and rank stabilization and induces the maps of the
based abelian covers and their second homology cokernels.  Consequently
the induced map on the specified presentation is $E(j)$.
For the PL comparison, the assembly cofiber sequence in the proof of
Theorem~\ref{thm:stable-pl-component-presentation} is natural in the
base. Its isomorphism from the degree-two smooth Whitehead group to
the degree-two PL Whitehead group, and therefore its inverse, commute
with these maps. This verifies the same identification for a possibly
nonsmoothable $M$; no arbitrary choice of a group isomorphism is used.

Positive suspension commutes with extension by the identity on
components: choose collar-normalized representatives and the same
positive bending in the interval coordinates.  Both composites agree
on the $N$-part and are the identity outside it.  The same argument
applies to compact families.  Passing to stabilization colimits and
then to the natural stable comparison gives the displayed square.
\end{proof}

\begin{theorem}\label{thm:component-followup-all-PL-six}
For every compact connected PL manifold $M$ of dimension at least six,
stabilization induces an isomorphism
\[
                      \pi_0P_{\mathrm{PL}}(M)\xrightarrow{\ \cong\ }E(M).
\]
Consequently, for every connected finite simplicial complex $K\simeq M$,
\[
                             \pi_0P_{\mathrm{PL}}(M)\cong E(K).
\]
No compatible smoothing of $M$ is assumed.
\end{theorem}

\begin{proof}
Choose $N$ by Lemma~\ref{lem:component-followup-handle-core}.
We need only that extension by the identity is surjective on
components.  The original component reduction to an absolute
three-handle stage is explicitly recorded in Pedersen's appendix
\cite[\S2, pp.~148--149]{Pedersen1975ComponentAppendix}.
For the precise extension map and a direct check of the required
weak consequence, use the PL statement in
\cite[Theorem~3.3 and the preceding paragraph,
p.~17]{GalvinNonino2025v1Component}.
That statement assumes
\[
 \begin{gathered}
 \pi_i(\partial N)\cong\pi_i(N)\ (i=0,1),\qquad
 \pi_i(M,N)=0\ (i\leq r),\\
 N\text{ is }k\text{-connected},\qquad 0\leq k\leq r\leq n-4.
 \end{gathered}
\]
Its PL connectivity bound for the extension map is
\[
                              \min(2r-3,r+k-2,r+2).
\]
Take $r=2$ and $k=0$.  All the hypotheses hold, and the bound is
$\min(1,0,4)=0$, since $2\leq n-4$.  Hence $\gamma_0$ is surjective.
This use requires no hypothesis on $\partial M\to M$ and does not use
an $r=3$ application in dimension six.

The smoothable dimension-six result established above applies to
$N$ and gives that $s_N$ is an isomorphism.
Lemma~\ref{lem:component-followup-E-naturality} gives that $E(j)$ is
an isomorphism because $j$ is three-connected.
The natural square in Lemma~\ref{lem:component-followup-stable-square}
therefore gives
\[
                              s_M\gamma_0=E(j)s_N
\]
as an isomorphism.  If $\gamma_0(a)=\gamma_0(b)$, applying $s_M$
and the injectivity of the right-hand side gives $a=b$.
Thus $\gamma_0$ is also injective, and $s_M$ is an isomorphism.
Homotopy invariance of $E$ proves the assertion for $K$.

\end{proof}

\begin{corollary}\label{cor:component-followup-chosen-three-skeleton}
Let $M$ be a compact connected PL manifold of dimension at least six,
possibly with boundary, and let $K$ be a specified finite
combinatorial triangulation of $M$.  Then
\[
                              \pi_0P_{\mathrm{PL}}(M)\cong E(K^{(3)}).
\]
\end{corollary}

\begin{proof}
The theorem gives $\pi_0P_{\mathrm{PL}}(M)\cong E(K)$.
The inclusion $K^{(3)}\to K$ is three-connected, so
Lemma~\ref{lem:component-followup-E-naturality} gives
$E(K^{(3)})\cong E(K)$.
This proves the conclusion for the specified skeleton.  It does not
require identifying the particular handle core used in the proof of
the theorem with that skeleton. 
\end{proof}

\begin{corollary}\label{cor:component-followup-fields-products}
The same presentations compute the stipulated definable concordance
component groups for finite PL models over every real closed field
and in each fixed o-minimal expansion.  For a compact manifold with
finitely many components, use the finite product of the connected
presentations, with each component in the corresponding dimension
range.
\end{corollary}

\begin{proof}
Use the established relative finite-model PL/definable component
comparison and its field independence, which preserve composition
and the entire fixed bottom-and-side subpolyhedron.
Bottom fixation forces each concordance and each permitted family
to preserve every connected component of the manifold.  Consequently
the groups and their component groups decompose as the asserted
finite products.
\end{proof}


\subsection{The topology of locally parameterized outer faces}

Pseudoisotopy obstructions to families of cone structures also appear
in Ayala--Francis--Tanaka's comparison between naive smooth families
and conically smooth stratified spaces.  In their category, the
automorphism space of a compact link is equivalent to that of its
open cone \cite[arXiv version~6, pp.~3--4; Theorem~4.3.1,
p.~46]{AyalaFrancisTanaka2017v6}.
The family category and the topology on the space of faces are
therefore part of the question.
Nariman's model for a three-manifold homeomorphism-group classifying space uses a
quotient of locally flat embeddings into an infinite-dimensional
half-space \cite[Section~7, pp.~13--14, especially Definition~7.1]
{Nariman2025}.  Here a fixed finite-dimensional cone and its shell
determine the stabilizer.
Salvatore--Turchin's smoothing-theory deloopings concern spaces of
disc embeddings and diffeomorphisms.  For $1\le m\le n$, their
Theorem~B excludes ambient dimension four and restricts to
invertible components in codimension two
\cite[pp.~2--3]{SalvatoreTurchin2026}.
The argument below uses compact-open locally flat embeddings,
supported isotopy extension, and a contraction of the ambient
cone homeomorphism group.

Let $B$ be a closed connected PL manifold of dimension $d\geq4$.
Write $X=C_2B$ for the cone with radial interval $[0,2]$, let $o$ be
its vertex, and put $C=C_1B$ and $Y_*=B\times\{1\}$.
A strict cone neighborhood $(D,Y)$ has a closed cone image $D\subset X$,
an outer face $Y\subset B\times(0,2)$, a homeomorphism $CY\to D$
fixing the base and taking the vertex to $o$, an open deleted-face
part $D\setminus Y$, and an exterior collar of $Y$.
The cone chart and exterior collar give an interior and an exterior
collar, so $Y$ is two-sided and locally flat.
The cone chart is required to exist but is not retained in the object.

The face determines its cone image: $D\setminus Y$ is the connected
component of $o$ in $X\setminus Y$, since it is connected and both
open and closed there, and $D=\overline{D\setminus Y}$.
We give outer faces the subspace topology in
\[
 \operatorname{Emb}_{\mathrm{lf}}(B,B\times(0,2))
       /\operatorname{Homeo}(B),
\]
where embeddings have the compact-open topology and the action is by
precomposition.  This topology is specified before any choice of
parameterization is forgotten.

A local joint face witness over a parameter neighborhood $U$ is a
homeomorphism
\[
 B\times U\longrightarrow
       \{(p,y):p\in U,\ y\in Y_p\}
\]
over $U$, with jointly continuous forward and inverse maps.
In a semialgebraic family these maps are semialgebraic.
A finite semialgebraic simplex filling consists of a finite
triangulated parameter disk, semialgebraic cone and face families
on its closed simplices, and a joint semialgebraic face witness
on each simplex.  Cone and face images agree on common parameter
faces and equal the prescribed loop on the boundary of the disk.
The witness parameterizations need not agree and are not retained.
Over a real closed field the cone charts, collar witnesses, and
their inverses are also required to be semialgebraic; ``compact''
means closed and bounded.  This finite convention imposes no
retained radial ruling and invokes no topological finite-subcover
property.

Put
\[
 P_{\mathrm{TOP}}(B)=
 \operatorname{Homeo}(B\times[1,2]\ \mathrm{rel}\ B\times\{2\}).
\]
Interval reversal identifies this with the usual pseudoisotopy group.

\begin{theorem}\label{thm:locally-parameterized-outer-faces}
The component $\mathcal W$ of the standard cone in this outer-face
space has weak homotopy type $BP_{\mathrm{TOP}}(B)$.
For $B=T^6$ there is a semialgebraic loop of strict cones that does
not bound a real definable family with local joint face witnesses.
The loop can be defined over $\mathbb Q$, and over every real closed
field it has no finite semialgebraic simplex filling with the
specified witnesses.

If a parameterization $B\to Y$ is retained, the component of the
standard parameterized face instead has weak homotopy type
\[
 B\operatorname{Homeo}
       (B\times[1,2]\ \mathrm{rel}\ B\times\{1,2\}).
\]
\end{theorem}

\begin{proof}
\emph{The quotient metric and finite-parameter face bundles.}
Let $E$ be the space of embeddings occurring as strict outer faces,
and let $U=E/\operatorname{Homeo}(B)$.
Fix a compatible metric $d_X$ on the compact cone.
The compact-open topology on $E$ is the topology of
$d_\infty(e,f)=\sup_{b\in B}d_X(e(b),f(b))$.
Precomposition acts by isometries, and
\[
 \delta([e],[f])=
 \inf_{a\in\operatorname{Homeo}(B)}
       \sup_{b\in B}d_X(e(b),f(a(b)))
\]
induces the quotient topology.
Indeed, inverse images of $\delta$-balls are unions of uniform balls;
conversely, an open saturated set containing $e$ contains a uniform
ball about $e$, whose quotient contains the corresponding
$\delta$-ball.  This is a metric: the Hausdorff distance between
$e(B)$ and $f(B)$ is bounded by $\delta$, and equal images differ by
a homeomorphism of $B$.

Let $P$ be a compact finite-dimensional metric space and let
$p\mapsto Y_p$ be a continuous map $P\to U$.  Its total face space
\[
 Z=\{(p,x)\in P\times X:x\in Y_p\}
\]
is closed in $P\times X$.  If $(p_i,x_i)\to(p,x)$ with
$x_i\in Y_{p_i}$, quotient-metric continuity implies Hausdorff
convergence to $Y_p$, and hence $x\in Y_p$.
Thus $Z$ is compact and its projection to $P$ is proper.
For any $p$ and $\varepsilon>0$, sufficiently close $q$ satisfy
$d_P(p,q)<\varepsilon/2$ and $\delta(Y_p,Y_q)<\varepsilon/2$.
The latter inequality supplies a homeomorphism between the two
fibres moving every point by less than $\varepsilon/2$ in $X$,
and therefore by less than $\varepsilon$ in $Z$ with its product
metric.  The projection is completely regular.

The completely regular mapping theorem
\cite[Proposition~7.1 and preceding definition, p.~605]{ChapmanFerry1979}
now applies.  Its base is locally finite-dimensional, its projection
is proper, and the homeomorphism group of the closed manifold fibre
is locally contractible by
\cite[Corollary~1.1, p.~63]{EdwardsKirby1971}.
Consequently $Z\to P$ is a locally trivial $B$-bundle.
In particular, a quotient family over a disk has a joint face
parameterization.  Over $D^k\times I$ a trivialization can be made
to extend any prescribed bottom parameterization: its discrepancy
at $t=0$ is a continuous map $D^k\to\operatorname{Homeo}(B)$,
and precomposition by that map, held constant in $t$, gives the
required correction.  Thus $E\to U$ has the Serre lifting property.
Only finite-dimensional compact parameter bases have been used.

\emph{The ambient action.}
Consider
\[
 G=\operatorname{Homeo}
       (X\ \mathrm{rel}\ \{o\}\cup B\times\{2\}).
\]
This group is contractible.  If $s_t:X\to C_{2t}B$ denotes radial
scaling, set, for $t>0$,
\[
 A_t(g)(x)=
 \begin{cases}
  s_tg s_t^{-1}(x),&x\in C_{2t}B,\\
  x,&x\notin C_{2t}B,
 \end{cases}
 \qquad A_0(g)=1.
\]
The formulas agree on the joining face because $g$ fixes the outer
base.  The inverse is $A_t(g^{-1})$.
At $t=0$ the displacement of both maps is bounded, uniformly in
$g$, by the diameter of $C_{2t}B$, which tends to zero in a conical
metric.  This proves continuity of the Alexander contraction,
including its dependence on $g$.

We use the compact-family extension theorem of
\cite[Theorem~1, pp.~1247--1248; covering-isotopy theorem,
p.~1250]{Chernavskii1969}.
In that theorem a compact family is jointly continuous and each
fibre embedding is locally flat; a locally flat total parameter
track is not an additional hypothesis.
Here the source is compact, the ambient manifold $B\times(0,2)$
has dimension $d+1\geq5$, and the codimension is one, hence differs
from two.  The local complementary simple-connectivity condition
in the small-extension theorem holds in a local flat chart.
Its displacement control function can vanish off a neighborhood
of the embedded source.  The extension of a compact parameter
family can therefore be supported in a compact annulus containing
its trace, and extends by the identity to an element of $G$.

The small-extension part also shows that every ambient orbit in
$E$ and in $U$ is open, and that its orbit map from $G$ is open.
For a prescribed identity neighborhood in $G$, choose the ambient
push small enough in both the forward and inverse uniform metrics.
Uniformly close parameterized embeddings then lie in its orbit
image; after reparameterization the same holds in the quotient.
Every orbit is connected, since $G$ is contractible, and so the
orbits are exactly the connected components.
Write $E_0$ and $U_0=\mathcal W$ for the standard components.

The map $G\to E_0$, $g\mapsto g|_{Y_*}$, is a Serre fibration.
Indeed, a family over $D^k\times I$ extends by the compact-family
theorem to a continuous family of ambient maps.  If $g(z,t)$ is
one such lift and $g_0(z)$ is the prescribed bottom lift, then
\[
 \widetilde g(z,t)=g(z,t)g(z,0)^{-1}g_0(z)
\]
has the prescribed initial value and still induces the same face
parameterization, since the right-hand correction fixes $Y_*$
pointwise.  The lifting property of $E\to U$ restricts to
$E_0\to U_0$, since lifted paths remain in their component.
It follows that $G\to U_0$ is a Serre fibration as well.

\emph{The two stabilizers.}
Let $H\subset G$ preserve $Y_*$ setwise, and let $H^{\mathrm{pt}}$
fix it pointwise.  The open orbit maps identify
$U_0=G/H$ and $E_0=G/H^{\mathrm{pt}}$ as topological quotients.
Put
\[
 A=\operatorname{Homeo}(C\ \mathrm{rel}\ \{o\}\cup Y_*),
 \qquad
 Q=\operatorname{Homeo}
       (B\times[1,2]\ \mathrm{rel}\ B\times\{1,2\}).
\]
Restriction to the outer shell gives a split group homomorphism
$H\to P_{\mathrm{TOP}}(B)$.  Its section extends the inner-end
boundary map conically over $C$, and its kernel is $A$.
Thus $H\cong A\rtimes P_{\mathrm{TOP}}(B)$, with a product
decomposition of the underlying spaces.  The group $A$ has the
same Alexander contraction.  For the pointwise stabilizer,
independent restrictions to the cone and shell give
$H^{\mathrm{pt}}\cong A\times Q$.

The principal Serre fibrations with contractible total group now
give
\[
 U_0\simeq_w BH\simeq_w BP_{\mathrm{TOP}}(B),
 \qquad
 E_0\simeq_w BH^{\mathrm{pt}}\simeq_w BQ.
\]
One may obtain the classifying-space conclusion directly on
singular simplicial sets: every simplex in the quotient lifts,
and two lifts differ by a unique continuous simplex in the
stabilizer.  The quotient of the contractible singular total
group therefore has the stabilizer's classifying weak type.
This uses no global section over the infinite-dimensional
quotient.

\emph{A rational semialgebraic essential loop.}
For $B=T^6$, the classical torus component calculation gives
\[
 \pi_0P_{\mathrm{PL}}(B)
       \cong\mathbb F_2[\mathbb Z^6]/\mathbb F_2[0]\ne0,
\]
and these classes survive in $P_{\mathrm{TOP}}(B)$, by
Theorem~\ref{result-3} and its unstable category comparison
\cite[Theorems~6.1--6.2, pp.~41--43]{BurgheleaLashof1974II}.
Choose a nonzero PL representative $p$ on the shell, fixed at its
outer end.  Use a rational finite PL model of the torus and a
rational PL product identification of the shell.

The representative may be chosen with rational PL coefficients.
To see this, choose finite source and target triangulations on
which $p$ is simplicial, refining a fixed rational triangulation
of the shell and agreeing on the fixed end.  Perturb their
vertices to rational points in their carriers, identically on
the fixed end.  Sufficiently small carrier-preserving perturbations
of this finite triangulation preserve its simplex incidences,
nondegeneracy, and coverage; the straight vertex motions therefore
give continuous families of PL homeomorphisms
$\alpha_s,\beta_s$ of the shell.
They agree on the fixed end, so $\beta_s p\alpha_s^{-1}$ is a
path in the topological pseudoisotopy group and its final member
has rational affine pieces.  Its TOP component is therefore the
same nonzero component as that of $p$.  Replace $p$ by that member.

Cone the free-end map of $p$ to obtain a rational PL element
$h\in H$, and set
\[
 g_t=A_t(h),\qquad C_t=g_t(C),\qquad Y_t=g_t(Y_*).
\]
The images form a loop because $g_0=1$ and $g_1=h\in H$.
The family is semialgebraic over $\mathbb Q$: on an affine piece
of $h$, the scaled expression $t\,h(x/t)$ is affine in $(x,t)$,
and its piece conditions are homogeneous linear inequalities.
The inverse has the same finite form using $h^{-1}$.
Transporting the standard cone chart and collar supplies joint
semialgebraic witnesses on the parameter interval.
If $k=h|_{Y_*}$, precompose the terminal face parameterizations
by $k^{-1}$ near $t=1$.  Their limiting value then equals the
initial parameterization, giving a joint local witness across
the identified endpoints of the circle.

The boundary homomorphism
$\pi_1(U_0)\to\pi_0(H)$ sends this loop to $[h]$, up to inverse.
Shell restriction sends it to the chosen nonzero class $[p]$.
Hence the loop is nonzero in $\pi_1(U_0)$.

\emph{Real fillings.}
A real family with local joint face witnesses gives a continuous
map to $U$: compactness of $B$ turns each total-space witness into
a continuous map to $E$, and equality of images glues the quotient
maps on overlaps.  A disk filling would therefore nullhomotope the
essential loop.  This excludes the asserted real definable fillings;
the argument applies to every continuous filling with these
witnesses.  No conclusion about Hausdorff-only image families is
being used.

\emph{Finite simplex fillings over real closed fields.}
Suppose such a filling exists over a real closed field $F$.
Its finitely many simplex families already supply finite graph
formulas for the joint face witnesses and their inverses.
To retain strictness under transfer, we first produce finitely
many graph formulas for cone and collar witnesses.

On each closed parameter simplex $\sigma$, apply the finite-subset
form of semialgebraic triviality to $X\times\sigma\to\sigma$,
compatible with the total cone and face families,
$\{o\}\times\sigma$, and $(B\times\{2\})\times\sigma$
\cite{Hardt1980}\cite[Theorem~5.46]{BasuPollackRoy2006}.
There is a finite semialgebraic partition into pieces $T_j$ with
ambient-pair trivializations.  Choose one parameter in each
nonempty piece and one semialgebraic cone chart and exterior
collar for its fibre.  Transport these along the trivialization.
If $a_z$ is its ambient map and $\phi_j:CY_j\to C_j$ the chosen
cone chart, the transported chart is
\[
 [y,r]\longmapsto a_z\phi_j([a_z^{-1}(y),r]).
\]
It fixes its base and vertex; the same conjugation transports
the collar.  These joint maps and their inverses are semialgebraic
on their indicated pieces.  They need not agree across pieces,
since they witness strictness and are not part of the object.
The construction uses a finite partition theorem, not a finite
subcover of a non-Archimedean topological space.

Fix finite polynomial formulas for all these sets and maps,
and collect their coefficients into one finite tuple $c$.
The parameter complex may be replaced by the rational realization
of its finite abstract triangulation, retaining the graph of its
boundary identification as part of this finite datum.
For these fixed formulas, being graphs of mutually inverse maps,
continuity, preservation of base coordinates, closedness and
boundedness of cone fibres, openness of their deleted-face parts,
the cone and collar identities, and equality of images on common
faces and on the boundary are first-order conditions.
Continuity is the epsilon--delta formula with the domain
predicates included.  The cone domain can be represented by
\[
 \{(ry,r):y\in Y,\ 0<r\leq1\}\cup\{0\},
\]
so no quotient sort is required.  In particular, no unrestricted
quantification over all semialgebraic chart formulas occurs.

All coefficients of the fixed torus, cone, and loop are rational.
Thus the existence of this finite-format filling is a sentence
$\exists c\,\Phi(c)$ over $\mathbb Q$.
Transfer for real closed fields makes it true over $\mathbb R$
\cite[Theorems~2.80--2.81]{BasuPollackRoy2006}.
The resulting real datum gives a continuous quotient map on each
closed parameter simplex.  Their image agreement and the finite
closed-simplex pasting lemma give a continuous disk in $U$ with
the prescribed boundary loop, contrary to the preceding paragraph.
This proves the field-independent non-filling assertion.
\end{proof}


\subsection{Hausdorff continuity does not supply face parameterizations}

The topology on unparameterized outer faces cannot be replaced silently
by the Hausdorff topology.  The following example separates these two
notions even for semialgebraic polygonal disks.

\begin{example}[A collapsing hairpin]\label{ex:hausdorff-hairpin}
In the square $X=[-3,3]^2$, let $C_0=[-2,2]^2$ and, for
$0<\varepsilon\leq 1/4$, put
\[
 C_\varepsilon=[-2,2]^2
 \cup\bigl([1,2]\times[2,2+\varepsilon]\bigr)
 \cup\bigl([-1,2]\times[2+\varepsilon,2+2\varepsilon]\bigr).
\]
Every $C_\varepsilon$ is a PL disk containing the origin in its interior,
so it is a strict cone neighborhood with a collared outer face.
Both $C_\varepsilon$ and $\partial C_\varepsilon$ converge in Hausdorff
distance to $C_0$ and $\partial C_0$, respectively.  Nevertheless, no
neighborhood of $0$ in the parameter interval admits a jointly
continuous family of homeomorphisms
\[
 e_\varepsilon:S^1\xrightarrow{\sim}\partial C_\varepsilon.
\]
Over an arbitrary real closed field the same conclusion holds for
jointly definable continuous parameterizations.
\end{example}

\begin{proof}
The two added rectangles attach successively along boundary intervals.
Their union with the square is therefore a PL disk.  Equivalently, its
boundary is the simple polygon with the following cyclic list of
vertices:
\[
 \begin{gathered}
 (-2,2),\ (1,2),\ (1,2+\varepsilon),\ (-1,2+\varepsilon),\\
 (-1,2+2\varepsilon),\ (2,2+2\varepsilon),\ (2,-2),\ (-2,-2).
 \end{gathered}
\]
A PL disk with a specified interior point is a PL cone on its boundary
with that point as vertex.  Its polygonal boundary also has an exterior
collar inside $X$.  All changed points lie in the strip
$2\leq y\leq2+2\varepsilon$.  Vertical projection, together with the
unchanged sides, shows that both Hausdorff distances are at most
$2\varepsilon$.  For positive parameters, a piecewise linear increasing
change of the $y$ coordinate identifies the disks and their boundaries,
so the only issue is the limit at zero.

Suppose a joint parameterization $e$ exists on $[0,\delta]\times S^1$.
Set
\[
 P_\varepsilon=(0,2),\qquad
 Q_\varepsilon=(0,2+\varepsilon),\qquad P=(0,2).
\]
Both marked points belong to $\partial C_\varepsilon$ and converge to
$P$.  If $a=e_0^{-1}(P)$, compactness of the circle and injectivity of
$e_0$ imply
\[
 e_\varepsilon^{-1}(P_\varepsilon)\longrightarrow a,
 \qquad
 e_\varepsilon^{-1}(Q_\varepsilon)\longrightarrow a.
\]
Choose a small closed arc $I$ about $a$ whose image under $e_0$ lies in
the ball of radius $1/4$ about $P$.  Joint continuity implies that, for
small positive $\varepsilon$, the whole arc $e_\varepsilon(I)$ lies in
the ball of radius $1/2$ about $P$ and contains both marked points.
The subarc joining them is one of the two arcs of the polygonal
boundary.  The first passes through $(1,2)$, while the other passes
through $(-1,2+\varepsilon)$.  Each contains a point at distance at
least one from $P$, a contradiction.

For the definable version, the map
\[
 [0,\delta]\times S^1\longrightarrow
 \{(\varepsilon,z):z\in\partial C_\varepsilon\},
 \qquad(\varepsilon,x)\longmapsto(\varepsilon,e_\varepsilon(x))
\]
would be a continuous definable bijection from a definably compact set
to a definable Hausdorff set.  Its inverse is continuous.  This gives
the same convergence of the two preimages, and definable compactness
gives the uniform control on $e_\varepsilon(I)$.  The polygonal arc
contradiction uses only the ordered-field inequalities above.
\end{proof}

Thus a result about outer-face families with local joint
parameterizations does not by itself answer a question about
Hausdorff-continuous families.  In particular, this example makes no
claim about whether a prescribed six-torus loop can be filled in the
Hausdorff-only space.


\subsection{Concordances of the sphere and projective plane}

We calculate components for the sphere and projective plane.  The arguments
use classical results of Cerf, Epstein, Laudenbach, and
Kalliongis--McCullough.  Throughout, a concordance fixes the lower face
pointwise and leaves the upper face free.  A PL isotopy means a jointly
PL total family, with a jointly PL inverse.

\begin{theorem}[The two-sphere]\label{thm:two-sphere-components}
For the standard PL sphere,
\[
                    \pi_0P_{\mathrm{PL}}(S^2)=0.
\]
The same holds for its stipulated definable models over every real
closed field.
\end{theorem}

\begin{proof}
We first work smoothly over the real numbers.  Let $D$ be the closed
unit three-ball, let $D_\epsilon$ be a smaller concentric ball, and put
\[
 G=\operatorname{Diff}^{+}(D),\qquad
 E=\operatorname{Emb}^{+}(D_\epsilon,\operatorname{int}D),\qquad
 H=\{g\in G:g|_{D_\epsilon}=\operatorname{id}\}.
\]
The mapping spaces carry their $C^\infty$ topologies.  Restriction
$G\to E$ is locally trivial over its image, by Cerf's restriction
theorem and its boundary-incidence version
\cite[Theorem~1, p.~8-01, and Section~2.3, p.~8-05]{Cerf1962Fibration}.
The local extensions may be supported away from $\partial D$.

Here are the two properties of $E$ needed below.  Given a compact
family $e_z$ of embeddings, write $p_z=e_z(0)$ and $A_z=De_z(0)$.
After replacing $e_z(x)$ by $e_z(ux)$, for one sufficiently small
$u>0$, compactness and uniform derivative continuity give
\[
 \sup_{z,x\in D_\epsilon}\|De_z(ux)-A_z\|
       <\tfrac12\inf_z\sigma_{\min}(A_z).
\]
The straight interpolation to $p_z+uA_zx$ is consequently embedded:
its error relative to the invertible linear map $uA_z$ has Lipschitz
constant less than half the least singular value of that map.
All images remain in $D$, by convexity and the uniform positive
distance of the centers from its boundary.  Decreasing $u$ if
necessary, move the centers to zero and use the polar decomposition
$A_z=Q_zP_z$ to interpolate $P_z$ to the identity.  The resulting
family is $x\mapsto uQ_zx$, with $Q_z\in SO(3)$.  These operations
are continuous in the smooth mapping-space topology.

For one embedding, connectedness of $SO(3)$ and expansion of the
small standard ball give a path to the inclusion $j$.  Thus $E$ is
path connected, and local path lifting shows that the restriction
map has image all of $E$.  For a loop based at $j$, the preceding
deformation moves the basepoint along precisely the radial path
$j\leadsto uj$ and then keeps it fixed.  The final loop $uQ_zj$
comes from the actual rotations $Q_z$ of $D$.  Shrinking their
restrictions uses the same radial basepoint path.  The basepoint
transports therefore agree, and
\[
                        \pi_1G\longrightarrow\pi_1E
                  \quad\text{is surjective}.
\]
This is the ball version of Cerf's frame-space argument
\cite[Proposition~3, p.~8-10]{Cerf1962Fibration}.

Cerf's three-disk theorem gives $\pi_0G=0$ in the formulation
recorded by Laudenbach
\cite[Chapter~I, Theorem~1.2, p.~8]{Laudenbach1974}.
The exact sequence of the restriction fibration, including its
action on fiber components, now gives $\pi_0H=0$.

We identify the correct boundary condition on the annulus
$A=D\setminus\operatorname{int}D_\epsilon\cong S^2\times I$.
Restriction identifies $H$ with those annulus diffeomorphisms whose
infinite jets on the inner boundary equal those of the identity.
Equality of all jets permits smooth extension by the identity over
$D_\epsilon$, for the map and its inverse.

Let $h$ be a smooth concordance of $S^2$, and let $a$ be its upper
restriction.  This is orientation preserving.  The connectedness
of $\operatorname{Diff}^{+}(S^2)$
\cite[Theorem~4 and Corollary~2, p.~8-13]{Cerf1962Fibration}
gives a smooth isotopy $a_s$ from the identity to $a$.  If
$\eta:I\to I$ is zero near zero and one near one, postcomposition by
\[
                  (x,t)\longmapsto(a_{s\eta(t)}^{-1}(x),t)
\]
joins $h$ to a map $h'$ fixing both boundary components.
Set $M=S^2\times\{0\}$ and
\[
 D_\infty=\operatorname{Diff}(A;J^\infty_{\partial A}),\qquad
 D_M=\operatorname{Diff}(A;\partial A;J^\infty_M),\qquad
 D_0=\operatorname{Diff}(A;\partial A).
\]
Cerf's boundary-jet proposition and its displayed special case
\cite[Proposition~1, p.~8-06]{Cerf1962Fibration}
say that $D_\infty\to D_M$ and $D_\infty\to D_0$ induce
isomorphisms on $\pi_i$ for $i\ge0$.  Hence $D_M\to D_0$ induces
an isomorphism on components.  A path relative to the entire
boundary therefore joins $h'$ to some $h''\in D_M$.
Extending $h''$ over the inner ball gives an element of $H$.
A path in $H$ to the identity restricts to a bottom-fixed path on
$A$.  We have proved
$\pi_0P_{\mathrm{Diff}}(S^2)=0$.

For the category change, use the three-dimensional comparison in
\cite[Chapter~I, Section~1, pp.~7--9]{Laudenbach1974}.
A PL concordance $h$ is TOP isotopic to a smooth annulus
diffeomorphism $g$.  Every map in this isotopy preserves the two
boundary components.  If $b_s$ is its bottom restriction,
postcompose by $b_s^{-1}\times\operatorname{id}_I$.  This gives a
bottom-fixed TOP isotopy ending at a smooth concordance, because
the final bottom restriction is smooth.  The smooth calculation
above then joins that endpoint to the identity.

Laudenbach's injectivity from PL to TOP isotopy classes gives a
jointly PL isotopy $L_s$ from the original $h$ to the identity.
Write $c_s=L_s|_{S^2\times\{0\}}$.  The corrected family
\[
                    (c_s^{-1}\times\operatorname{id}_I)L_s
\]
is jointly PL and fixes the bottom.  Indeed, the total boundary
map is a PL homeomorphism and its inverse is PL.  Since
$c_0=c_1=\operatorname{id}$, the endpoints are unchanged.
This proves the real PL assertion.  The relative component
comparison and coefficient transfer of
Theorems~\ref{result-4} and~\ref{result-40} give the stated
definable conclusion.
\end{proof}

\begin{theorem}[The projective plane]\label{thm:projective-plane-components}
For the standard PL projective plane,
\[
                  \pi_0P_{\mathrm{PL}}(\mathbb{RP}^2)=0.
\]
The same holds for its stipulated definable models over every real
closed field.
\end{theorem}

\begin{proof}
Put $B=\mathbb{RP}^2$ and let $h\in P_{\mathrm{PL}}(B)$.
Its upper restriction $\phi$ is homotopic to the identity, through
$(x,t)\mapsto\operatorname{pr}_B h(x,t)$.
Epstein's surface theorem
\cite[Theorem~6.4(a), p.~101]{Epstein1966}
provides a jointly PL isotopy $a_u$ from the identity to $\phi$.
The total inverse isotopy is PL.  With
$r(s,t)=\max(0,s+t-1)$, define
\[
                     F_s(x,t)=(a_{r(s,t)}^{-1}(x),t).
\]
This is a jointly PL family, with inverse obtained by replacing
$a^{-1}$ by $a$.  It starts at the identity and fixes the bottom.
At $s=1$ it has upper restriction $\phi^{-1}$, so
$h'=F_1h$ fixes both ends.

Kalliongis--McCullough prove that the relative homeotopy classes of
$B\times I$, with both ends fixed, are generated by rotation loops
\cite[Section~4.1 and Lemma~4.1.1, p.~110]{KalliongisMcCullough1992}.
Thus $h'$ is TOP isotopic relative to both ends to a map
\[
                  R(x,t)=(\rho(t)x,t),
\]
where $\rho:I\to SO(3)$ is a based loop acting on projective lines.
The generating conclusion suffices here.  Once the upper end is
free, the formula
\[
                  R_s(x,t)=(\rho((1-s)t)x,t)
\]
contracts $R$ through bottom-fixed TOP homeomorphisms.
This formula is used in TOP.

The three-dimensional PL--TOP comparison
\cite[Chapter~I, Section~1, pp.~7--9]{Laudenbach1974}
therefore gives a jointly PL isotopy $H_s$ from $h'$ to the identity.
It preserves each boundary component.  If $b_s$ is its restriction
to the bottom, then
\[
                  \widehat H_s=(b_s^{-1}\times\operatorname{id}_I)H_s
\]
is jointly PL and fixes that bottom pointwise.  To verify the
parameter assertion, invert the total PL boundary homeomorphism,
take its product with the remaining interval, and compose with
the total $H$-map.  Both endpoints have bottom restriction equal
to the identity, so they remain $h'$ and the identity.
Concatenating this path with $F_sh$, using a piecewise affine
reparameterization, proves the PL statement.
Theorems~\ref{result-4} and~\ref{result-40} transfer the same
fixed-bottom component calculation to the stipulated definable
models over every real closed field.
\end{proof}



\subsection{Closed simply connected five-manifolds}

\begin{theorem}[Cerf's theorem and the PL component comparison]
\label{thm:simply-connected-five-components}
Let $M$ be a closed connected PL five-manifold with $\pi_1M=0$.
Then
\[
                         \pi_0P_{\mathrm{PL}}(M)=0.
\]
The corresponding component group of every stipulated definable model
over a real closed field is also zero.  There is no restriction on
$\pi_2M$.
\end{theorem}

\begin{proof}
First choose a compatible smoothing of $M$.  The smoothing
obstruction for a PL handle of index $k$ lies in the twisted-sphere
group $\Gamma_{k-1}$
\cite[Essay~IV, Appendix~B, p.~211]{KirbySiebenmann1977}.
The low-dimensional compatible-smoothing results and
$\Gamma_4=0$ remove every such obstruction for $k\le5$
\cite[Essay~I, p.~39; Essay~IV, p.~200]{KirbySiebenmann1977}.
There is no prescribed boundary smoothing to extend, since $M$ is
closed.  Denote the resulting smooth manifold by $M^{\mathrm{sm}}$.
It is still simply connected.

Cerf's full pseudoisotopy theorem says that the group of
diffeomorphisms of $V\times I$ fixing $V\times\{0\}$ is connected
when $V$ is closed, smooth, simply connected, and
$\dim V\ge5$
\cite[Introduction, Definition~2 and Theorem~0, pp.~7--8]{Cerf1970}.
The theorem concerns the whole pseudoisotopy group.
That group is locally path connected: in a product metric near
the boundary, a sufficiently small local-addition chart consists
of small smooth vector fields tangent to the boundary and zero
on the bottom.  Scaling such a field to zero stays in this chart.
Consequently
\[
                        \pi_0P_{\mathrm{Diff}}(M^{\mathrm{sm}})=0.
\]

For a compatible differentiable triangulation, the
Burghelea--Lashof component comparison
\cite[Theorem~6.1, pp.~41--43]{BurgheleaLashof1974II} gives
\[
 \pi_0P_{\mathrm{Diff}}(M^{\mathrm{sm}})
       \xrightarrow{\ \cong\ }\pi_0P_{\mathrm{PL}}(M).
\]
The five-dimensional endpoint is included in their proof:
Theorem~5.1(a), p.~37, has the bound $k\ge5$.
The concordance groups defined immediately before Theorem~6.1
fix the bottom and side boundary, as here; the side is empty in
the present closed case.  This proves the PL assertion.

Finally, the relative component comparison and finite PL
coefficient transfer in
Theorems~\ref{result-4} and~\ref{result-40} preserve the fixed
bottom and the group law.  They give the asserted result for
the stipulated definable models over every real closed field.
\end{proof}


\subsection{Stable quotients in dimension five}

The complete component presentation in dimension six also gives a
quotient of the component group in dimension five. The next proposition
uses only the realization part of the classical obstruction theory.

\begin{proposition}\label{prop:five-dimensional-stable-surjectivity}
Let $M$ be a compact connected smooth five-manifold, possibly with
boundary. With the bottom and entire side boundary fixed and the
upper face free, positive suspension is surjective on components:
\[
 \sigma^+_*:\pi_0P_{\mathrm{Diff}}(M)
       \longrightarrow\pi_0P_{\mathrm{Diff}}(M\times I).
\]
Consequently every stable smooth concordance component is represented
on $M$. The corresponding assertions hold for compact PL
five-manifolds and their stipulated definable models over every real
closed field and in every fixed o-minimal expansion. For a finite
pointed homotopy model $K$ of the manifold, the
canonical stabilization map onto the complete presentation $E(K)$ is
a surjective homomorphism. Finite disjoint unions give finite products.

No hypothesis on $\pi_2M$ is required. Injectivity and a description
of the kernel are not asserted.
\end{proposition}

\begin{proof}
Write $C_5=\pi_0P_{\mathrm{Diff}}(M)$ and
$C_6=\pi_0P_{\mathrm{Diff}}(M\times I)$, rounding corners when
needed. As above, collar normalization and the fixed positive bending
give the homomorphism $s=\sigma^+_*:C_5\to C_6$. The first
obstruction is a surjective homomorphism
\[
 \Sigma_5:C_5\longrightarrow\operatorname{Wh}_2(\pi_1M).
\]
Its dimension-five surjectivity is exactly
\cite[Part~I, Chapter~VI, Theorem~2 and its surjectivity proof,
p.~213]{HatcherWagoner1973}. The proof constructs the elementary
handle-addition word and then cancels its identity intersection
matrix in dimension at least five. The first obstruction is
unchanged under positive suspension, by the remark on p.~130 of
the same source, so $\Sigma_6s=\Sigma_5$.

Only in dimension six do we use the complete corrected sequence
\eqref{eq:full-geometric-concordance-sequence}. Put
\[
 A=\operatorname{Wh}_1^+
       (\pi_1M;\mathbb Z/2\oplus\pi_2M),\qquad
 \mathcal R_\chi=\operatorname{im}\chi_{M\times I}.
\]
Thus $\ker\Sigma_6$ is the image of the injective map
$\overline\theta_6:A/\mathcal R_\chi\to C_6$. This includes the entire
Postnikov and framing correction. There is no dimension-five
exact sequence being assumed.

We claim that every element of $\overline\theta_6(A/\mathcal R_\chi)$ belongs
to $s(C_5)$. For this claim it suffices to realize any specified
finite coefficient sum before passing to the quotient $A/\mathcal R_\chi$.
Hatcher and Wagoner construct a marked one-parameter lens family
with any prescribed intersection coefficient in
$(\mathbb Z/2\oplus\pi_2M)[\pi_1M]$ whenever
$3\leq i\leq n-2$
\cite[Part~I, Chapter~VII, Proposition~5 and its complete
construction, pp.~227--228]{HatcherWagoner1973}.
Choose $n=5$ and $i=3$. This is the allowed endpoint $i=n-2$.

At that endpoint the normal bundle used in the local Whitney
construction has only one framed direction. The cited proof treats
this case separately: changing the inward pushing direction is
classified by $\pi_1S^1=\mathbb Z$, whose stabilization realizes
either value in $\pi_1O=\mathbb Z/2$. The choice of Whitney disk
realizes the specified $\pi_2M$ coefficient, and iteration gives
every finite sum with its fundamental-group labels. Thus the
endpoint has not been obtained by applying a higher-codimension
framing argument outside its range.

Positive suspension leaves the lens indices $3,4$ unchanged.
In the six-dimensional target these are in the stable marked-lens
range. The intersection curves, their maps to the loop space, and
their labels remain unchanged. The added positive coordinate is
a trivial line and preserves the stable framed bordism class. The stable convention
$\lambda^s_0=(-1)^3\lambda_0$ is compatible with this positive
suspension
\cite[Sections~4.1--4.2, pp.~8--9]{Igusa2021Dimension3}.
Since the source construction realizes every coefficient, the fixed
sign convention does not restrict its image. Forgetting the marked
lens data in the six-dimensional target gives exactly its geometric
$\theta_6$ map. Therefore each required coefficient class, and each
finite sum of them, has image $s(g)$ for a product of
five-dimensional lens concordances $g$. Passing to $A/\mathcal R_\chi$ can only
identify such images. This proves the claim without constructing a
five-dimensional homomorphism from $A/\mathcal R_\chi$ or classifying
five-dimensional lenses with equal invariants.

Now let $y\in C_6$. By surjectivity of $\Sigma_5$, choose
$x\in C_5$ with $\Sigma_5(x)=\Sigma_6(y)$. Then
$y\,s(x)^{-1}$ belongs to $\ker\Sigma_6$, hence by the claim
equals $s(g)$ for some $g\in C_5$. Thus $y=s(gx)$.
This proves surjectivity as a group statement and uses no splitting
of the corrected sequence.

The handle additions, disks, and lens modifications in the cited
realization occur in the interior. Collar normalization and positive
bending fix the entire side boundary, so the argument applies when
$M$ has boundary. The six-dimensional component theorem makes every
later stabilization an isomorphism. Hence the same source maps
surjectively onto the stable group.

A compact PL five-manifold admits a compatible smoothing, including
the boundary existence statement proved above. Let $N=M\times I$.
The smooth/PL component isomorphisms $\kappa_M$ and $\kappa_N$,
implemented through the common piecewise differentiable model,
respect composition and are compatible with forgetting to TOP.
The smooth and PL positive bends give the same map on relative
TOP components. Indeed, their marked parameter-disk straightenings
differ by an orientation-preserving homeomorphism preserving the
marked boundary arcs. Isotope its boundary-interval parameterizations
to the identity, extend over the disk, and apply the disk Alexander
isotopy to the remaining boundary-fixed map. Conjugation by this
isotopy preserves the pointwise bottom-and-side condition.
Only a TOP isotopy of the bends is used here.

The forgetful map
\[
 j_N:\pi_0P_{\mathrm{PL}}(N)\longrightarrow
          \pi_0P_{\mathrm{TOP}}(N)
\]
is injective because $\dim N=6$
\cite[Theorem~6.2, p.~41, and proof pp.~42--43]{BurgheleaLashof1974II}.
Consequently the equality
\[
 j_N\kappa_N\sigma^+_{\mathrm{Diff}}
   =j_N\sigma^+_{\mathrm{PL}}\kappa_M
\]
on relative components implies
$\kappa_N\sigma^+_{\mathrm{Diff}}
 =\sigma^+_{\mathrm{PL}}\kappa_M$.
Thus smooth surjectivity transfers to PL.
The relative finite-model PL/definable comparisons commute with
the fixed rational polyhedral suspension and with field extension,
so they transfer the same conclusion to each stipulated definable
model. The stable presentation identifies the target with $E(K)$
as a group under composition. Bottom fixation prevents permutation
of connected components, giving the finite-product assertion.
The empty case is immediate.
\end{proof}

The construction is a consequence of the cited realization theorems
and the corrected six-dimensional sequence. It does not resolve the
dimension-five injectivity question: equality of stabilized framed
coefficients does not by itself give a deformation of the original
lens families.


% Authored continuation for original A.128--A.129.
% Targeted source/proof audit. Independent/main acceptance remains a separate gate.
\subsection{Standard spheres and simply connected three-manifolds}

The standard-sphere calculation is a consequence of the classical
unstable PL bundle comparison.  It concerns the jointly PL concordance
group, with the bottom fixed and the top free.

\begin{theorem}\label{thm:standard-sphere-components}
For every standard PL sphere of dimension $n\geq2$,
\[
                    \pi_0P_{\mathrm{PL}}(S^n)=0.
\]
The same holds for its stipulated definable models over every real
closed field.  In particular, this includes the standard three- and
four-spheres.
\end{theorem}

\begin{proof}
Put $X_j=\mathrm{PL}_j/\mathrm O_j$, where $\mathrm{PL}_j$ is the
simplicial group of germs at the origin.  These are the unstable
spaces.  The sphere calculation of Burghelea and Lashof gives
\begin{equation}\label{eq:sphere-concordance-fiber}
 P_{\mathrm{PL}}(S^n)\simeq
 E(X_{n+1},X_n)
 =\operatorname{hofib}_{*}(X_n\longrightarrow X_{n+1});
\end{equation}
see \cite[Section~6, equation~(1), p.~43]{BurgheleaLashof1974II}.
Here a point of $E(X_{n+1},X_n)$ consists of $x\in X_n$ and a path
in $X_{n+1}$ from its basepoint to the image of $x$.

For clarity, this comparison uses the PL bundle fibration
\[
 \mathrm{PL}_n\longrightarrow\mathrm{PL}_{n+1}
       \longrightarrow S^n\times BP_{\mathrm{PL}}(S^n)
\]
and its comparison with the orthogonal sphere fibration; see
\cite[Proposition~5.3 and the preceding paragraph, p.~38]
{BurgheleaLashof1974II}.  The underlying exact sequence in
\cite[Theorem~1 and definitions, pp.~243--244]{KuiperLashof1966II}
has no high-dimensional stability hypothesis.  Its parameter simplices
are PL bundle isomorphisms on the entire parameter product, and its
sphere concordances fix the bottom pointwise while leaving the top
free.  On compact products these are precisely the finite, jointly PL
families used here.

The low-degree stabilization formula is
\[
 \pi_i(X_j)\xrightarrow{\ \cong\ }
       \pi_i(\mathrm{PL}/\mathrm O)\qquad(i<j),
\]
as recorded in \cite[Section~5, formula~(3), p.~37]
{BurgheleaLashof1974II}.  The stable space is connected and its
homotopy groups vanish through degree six
\cite[Annex~C, calculation following Theorem~5.4, p.~318]
{KirbySiebenmann1977}.  Consequently, for every $n\geq2$,
\[
                 \pi_0X_n=0,
                 \qquad \pi_1X_{n+1}=0.
\]
The component part of the homotopy-fiber sequence is
\[
 \pi_1X_{n+1}\longrightarrow
 \pi_0P_{\mathrm{PL}}(S^n)\longrightarrow\pi_0X_n.
\]
Exactness, including the fundamental-group action on fiber
components, now makes the middle pointed set a singleton.  This
proves that the component group is trivial; compatibility of
\eqref{eq:sphere-concordance-fiber} with a chosen multiplication is
unnecessary for this conclusion.

Finally, Theorems~\ref{result-40} and~\ref{result-4} transfer the
finite PL calculation and its fixed-bottom condition to the
stipulated definable models over every real closed field.
\end{proof}

The low-degree formula used here is separate from the stronger
high-dimensional range of Burghelea--Lashof Theorem~5.1.
The smooth-to-PL component comparison is not needed in base
dimensions three and four.  The theorem asserts nothing about
general closed manifolds in those dimensions.

\begin{corollary}\label{cor:simply-connected-three-components}
Let $M$ be a closed connected simply connected PL three-manifold.
Then $\pi_0P_{\mathrm{PL}}(M)=0$.  The corresponding component group
also vanishes for every closed definably simply connected definable
three-manifold, with its stipulated PL model, over every real closed
field and in every fixed o-minimal expansion.
\end{corollary}

\begin{proof}
A PL three-manifold has a compatible smoothing
\cite[Essay~I, p.~39]{KirbySiebenmann1977}.  The resulting smooth
manifold is orientable because it is simply connected.  Perelman's
elliptization conclusion identifies a closed three-manifold with
finite fundamental group as a spherical space form
\cite[Section~1.4, p.~3]{Perelman2003Extinction}.  Here the covering
group is trivial, so there is a homeomorphism $h:M\to S^3$.

To compare the given PL structures, apply Shalen's approximation
theorem \cite[Approximation Theorem, p.~34; proof in Section~9,
pp.~75--79]{Shalen1984}.  Since both boundaries are empty, it gives a
PL embedding $q:M\to S^3$ with no additional boundary condition.
Its image is open by invariance of domain and closed by compactness.
The sphere is connected, so $q$ is onto and is a PL homeomorphism.
Conjugation by $q\times\mathrm{id}_I$, also on every total parameter
product, identifies the two jointly PL concordance simplicial groups.
Theorem~\ref{thm:standard-sphere-components} therefore applies.

For the definable assertion, realize the finite abstract PL model
with rational vertices and choose one vertex as basepoint.  This is
a pointed semialgebraic set defined without parameters.
Baro--Otero's comparison identifies the o-minimal fundamental group
over the given field with the classical fundamental group of the
real realization
\cite[arXiv version~1, Corollary~4.3, p.~10]{BaroOtero2008v1}.
Thus that real PL three-manifold is simply connected.  Apply the
real calculation, then the relative component and field comparisons
of Theorems~\ref{result-4} and~\ref{result-40}.  They preserve the
fixed-bottom convention and composition.
\end{proof}


\subsection{A higher PL--TOP discrepancy on the three-sphere}

The component calculation for spheres does not extend to an unrestricted
higher PL--TOP comparison.  In the following statement
$\mathrm{TOP}(4)=\operatorname{Homeo}(\mathbb R^4;0)$ is the group of
global origin-fixing homeomorphisms with its compact-open topology, and
\[
 Q_4=\operatorname{hofib}
       \bigl(B\mathrm O(4)\longrightarrow B\mathrm{TOP}(4)\bigr).
\]
Thus $Q_4$, also denoted $\mathrm{TOP}(4)/\mathrm O(4)$ in smoothing
theory, is a homotopy fibre.  All loop spaces and positive homotopy
groups below use the identity basepoint.

\begin{theorem}\label{thm:three-sphere-higher-top-discrepancy}
For the standard PL three-sphere, with bottom fixed pointwise and
upper boundary free, there is a based weak equivalence
\[
                  P_{\mathrm{TOP}}(S^3)\simeq\Omega Q_4.
\]
Here the TOP concordance object is the singular complex of the
compact-open homeomorphism group, and the PL object uses jointly
finite PL families on the whole parameterized product.  Their
homotopy groups satisfy
\[
 \begin{gathered}
 \pi_0P_{\mathrm{PL}}(S^3)=\pi_1P_{\mathrm{PL}}(S^3)
          =\pi_2P_{\mathrm{PL}}(S^3)=0,\\
 \pi_0P_{\mathrm{TOP}}(S^3)=\pi_1P_{\mathrm{TOP}}(S^3)=0,
 \qquad \pi_2P_{\mathrm{TOP}}(S^3)\cong\mathbb Z/2.
 \end{gathered}
\]
Consequently the natural real PL-to-TOP map fails to be a weak
equivalence for this base.
\end{theorem}

\begin{proof}
Let $D\subset\mathbb R^4$ be the closed unit ball.  We use singular
complexes of topological groups throughout, and write
\[
 A=\operatorname{Sing}\operatorname{Homeo}(D;0),\qquad
 T=\operatorname{Sing}\operatorname{Homeo}(\mathbb R^4;0).
\]
Let $G$ be the simplicial group of parameterized germs at the zero
section.  A $q$-simplex is a germ of a homeomorphism between open
neighbourhoods of $\{0\}\times\Delta^q$ in
$\mathbb R^4\times\Delta^q$, preserving the parameter and fixing
the zero section.  Representatives are identified when they agree
on a neighbourhood of the entire zero section.  In particular,
compactness of the parameter simplex supplies common positive
source and image radii for each representative and its inverse.
There are restriction homomorphisms
\[
                       r:A\longrightarrow G,
                 \qquad g:T\longrightarrow G.
\]

\medskip\noindent\textit{A relative lifting statement.}
We first prove that these are Kan fibrations, with exact agreement
on the prescribed parameter faces.  The geometric input is a
continuous local ambient-extension operator.  For concentric compact
balls $C\subset\operatorname{int}U\subset U\subset\operatorname{int}D$,
there is a neighbourhood $\mathcal V$ of the inclusion in the
embedding space $\operatorname{Emb}(U,\operatorname{int}D)$ and a
continuous map
\[
 \sigma:\mathcal V\longrightarrow
       \operatorname{Homeo}(D\;\mathrm{rel}\;\partial D)
\]
such that $\sigma(v)|_C=v|_C$ and
$\sigma(\operatorname{id})=\operatorname{id}$.  Its support may be
kept in a fixed compact subset of $\operatorname{int}D$.
This is the canonical ambient extension obtained from
\cite[Theorem~5.1, p.~70; proof of Corollary~1.2, p.~79]
{EdwardsKirby1971}.  Their deformation depends continuously on the
input embedding, fixes the inclusion, ends at an embedding equal to
the identity on $C$, and is unchanged off a compact neighbourhood
of $C$ in $U$.  Composing the input with the inverse of that endpoint
and extending by the identity gives $\sigma$.  Shrinking the input
neighbourhood keeps all supports away from $\partial D$.

Let $H=\Lambda_i^q\subset\Delta^q$ be a horn.  Suppose that a
family $a:H\to\operatorname{Homeo}(D;0)$ lifts the restriction of
a given germ simplex represented by $f_t$.  There are finitely
many horn faces; equality of their parameterized germs therefore
gives, after one common shrinking,
\[
 f_t=a_t\quad\hbox{on a fixed ball, for all }t\in H.
\]
Choose a strong deformation $d_s$ of $\Delta^q$ from a retraction
onto $H$ to its identity, fixed on $H$.  A simplex collapses onto
each of its horns, so such a deformation exists.  Set
\[
 a^0_t=a_{d_0(t)},\qquad
 k_{t,s}=f_{d_s(t)}(a^0_t)^{-1}
\]
as germs.  Uniform shrinking, using joint continuity of the maps
and inverses on the compact parameter set, represents all these
germs and their inverses on common neighbourhoods.  On a smaller
common neighbourhood they satisfy the actual equalities
\[
       k_{t,0}=\operatorname{id},\qquad
       k_{t,s}=\operatorname{id}\quad(t\in H).
\]

Choose $U$ so small that every $k_{t,s}^{-1}$ is defined on $U$
and sends it into a common source neighbourhood for all the maps
$k_{t,s'}$.  Choose $C\subset\operatorname{int}U$, and then a
still smaller ball $V$ such that
\[
                  k_{t,s}(V)\subset\operatorname{int}C
                  \quad(t\in\Delta^q,\ s\in I).
\]
These choices follow from compactness and joint continuity at the
fixed zero section.  The increments
\[
          v_{t;s,s'}=k_{t,s}k_{t,s'}^{-1}|_U
\]
are continuous in the embedding topology and equal the inclusion
when $s=s'$.  Compactness of this diagonal supplies a single
$\eta>0$ such that $|s-s'|<\eta$ puts every increment in
$\mathcal V$.  In particular all those increments take values in
$\operatorname{int}D$.

Take a finite partition $0=s_0<\cdots<s_N=1$ with mesh less
than $\eta$.  Starting with $b_{t,0}=\operatorname{id}$, define
on each subinterval
\[
 b_{t,s}=\sigma(v_{t;s,s_j})b_{t,s_j},
                  \qquad s_j\leq s\leq s_{j+1}.
\]
The formulas fit at the endpoints because $\sigma(\operatorname{id})$
is the identity.  Inductively $b_{t,s}|_V=k_{t,s}|_V$: the input
to the next extension is $k_{t,s_j}(x)\in C$ for $x\in V$.
The extensions fix $0$.  On the whole horn every increment is
the identity, so $b_{t,s}$ is the identity there for every $s$.
Consequently
\[
                       \widetilde a_t=b_{t,1}a^0_t
\]
is a disk-homeomorphism family equal to $a$ on $H$, with germ
$f_t$ on a common ball.  This is the required horn lift.  For $g$,
extend the same supported corrections by the identity outside
$D$ and compose them with the original global family.  This proves
both fibration statements without assuming that the original
family itself is close to the identity.

\medskip\noindent\textit{The two fibres.}
Let $K=\ker r$.  For $0<\varepsilon<1$, write
\[
 K_\varepsilon=
 \operatorname{Sing}\operatorname{Homeo}
              (D\;\mathrm{rel}\;D_\varepsilon),
 \qquad D_\varepsilon=\varepsilon D.
\]
The kernel $K$ is the filtered union of these simplicial groups.
Restriction to the annulus, followed by a fixed product
parameterization, identifies $K_\varepsilon$ with
$P_{\mathrm{TOP}}(S^3)$.  Extension by the identity on the inner
ball is the continuous inverse.  The inner sphere is fixed and
the outer sphere is free, exactly as required.

If $0<\varepsilon<\varepsilon_0<1$, choose increasing radial
homeomorphisms $\rho_u$ with $\rho_0=\operatorname{id}$ and
\[
 \rho_u(D_{\varepsilon_0})=D_{v(u)},\qquad
 \varepsilon\leq v(u)\leq\varepsilon_0,
 \qquad v(1)=\varepsilon.
\]
For example, linearly interpolate the increasing piecewise affine
radius functions.  Conjugation by $\rho_u$ gives a homotopy within
$K_\varepsilon$ from the inclusion of $K_{\varepsilon_0}$ to an
isomorphism onto $K_\varepsilon$.  Every finite simplicial sphere
and homotopy in the filtered union uses a single common radius.
Hence
\[
          K_{\varepsilon_0}\xrightarrow{\ \simeq\ }K,
                 \qquad K\simeq P_{\mathrm{TOP}}(S^3).
\]

The global kernel $L=\ker g$ is weakly contractible.
Every finite simplicial family $h_t$ in $L$ is represented on finitely
many parameter simplices. Their identity-germ conditions therefore
give a common fixed ball $D_\varepsilon$. The contraction
\[
 c_s(h_t)(x)=
 \begin{cases}s^{-1}h_t(sx),&s>0,\\x,&s=0\end{cases}
\]
is continuous in the compact-open topology: on any fixed compact
set it is exactly the identity for all sufficiently small $s$,
uniformly in $t$, and the same assertion holds for its inverse.
It stays in the kernel, fixing $D_\varepsilon$ for $0\leq s\leq1$,
and keeps identity parameter values fixed.  This contracts all
finite homotopy representatives and proves the claim.

Write $T_0,G_0$ for the simplicial identity components.  The
fibration $g$ restricts to $T_0\to G_0$.  Every vertex of $G_0$
is joined to the identity by finitely many edges, which lift by
the fibration property.  Its fibre is $L$, so the fibre sequence
gives a weak equivalence
\[
                            T_0\xrightarrow{\ \simeq\ }G_0.
\]
No globalization of a germ in an arbitrary other component is
needed here.

\medskip\noindent\textit{Identifying the disk group.}
Restriction to the boundary gives
$\operatorname{Homeo}(D;0)\to\operatorname{Homeo}(S^3)$,
with the radial-coning section.  Its kernel is contractible by
the relative Alexander isotopy
\[
 a_u(x)=
 \begin{cases}
      u\,a(x/u),&\|x\|\leq u,\ u>0,\\
      x,&\|x\|\geq u\text{ or }u=0.
 \end{cases}
\]
Both formulas agree on the dividing sphere.  They fix the origin,
and the displacement of the map and its inverse is at most $2u$,
which gives uniform continuity at $u=0$.  Splitting by the coned
boundary map proves that boundary restriction is a homotopy
equivalence.

The actual topological three-sphere group has the homotopy type
given by the natural inclusion
$\mathrm O(4)\to\operatorname{Homeo}(S^3)$
\cite[Appendix, statements~(3) and~(16), pp.~605--606]{Hatcher1983}.
The cited appendix explicitly includes TOP, so this is not merely
the smooth Smale theorem.  Thus
\[
      \operatorname{Sing}\mathrm O(4)\xrightarrow{\ \simeq\ }A,
 \qquad \operatorname{Sing}\mathrm{SO}(4)\xrightarrow{\ \simeq\ }A_+,
\]
where $A_+$ is the orientation-preserving identity component.
Every element of $K$ fixes a neighbourhood of the origin and
therefore preserves orientation; it is exactly the fibre of
$A_+\to G_0$.

Extend a disk homeomorphism to all of $\mathbb R^4$ by its
boundary map outside the disk:
\[
                  e(a)(r\omega)=r\,a|_{S^3}(\omega)
                         \quad(r\geq1).
\]
Together with $e(a)=a$ on $D$, this defines a continuous group
homomorphism $e:A\to T$ with $ge=r$.  On $\mathrm O(4)$ it is
the usual linear inclusion.  The preceding weak equivalence and
fibration statements identify
\[
 \begin{aligned}
 K&\simeq\operatorname{hofib}(A_+\longrightarrow T_0)\\
  &\simeq\operatorname{hofib}(\mathrm{SO}(4)\longrightarrow T_0)
   \simeq\operatorname{hofib}_1(\mathrm O(4)\longrightarrow T).
 \end{aligned}
\]
In the last step an orientation-reversing orthogonal map cannot
be joined to the identity in $T$, while $\mathrm{SO}(4)$ lies
in $T_0$.  Thus precisely these source and target components
enter the based homotopy fibre.  Using the singular-group models
for classifying spaces, looping their defining homotopy-fibre
square identifies the last space with $\Omega Q_4$.  Together
with the description of $K$, this proves the asserted TOP model.

\medskip\noindent\textit{The homotopy groups.}
Only the $4$-connectedness of the stabilization
$Q_4\to\mathrm{TOP}/\mathrm O$ is needed here: it gives
isomorphisms on homotopy groups in degrees at most three.
This is established in the first part of the proof of
\cite[Theorem~8.7A, p.~128; proof on pp.~130--131]{FreedmanQuinn1990},
whose full statement is $5$-connectedness.
The corrected account
\cite[Section~1.1, Theorem~1.3 and the following paragraph,
p.~2]{GabaiGayHartmanKrushkalPowell2026}
specifies exactly the global origin-fixing group and
classifying-space homotopy-fibre convention used above and
supplies the correction needed for the final $5$-connected step.
The low stable groups are
\[
 \pi_i(\mathrm{TOP}/\mathrm O)=
 \begin{cases}
 \mathbb Z/2,&i=3,\\
 0,&0\leq i\leq6,\ i\ne3,
 \end{cases}
\]
by \cite[Essay~V, Theorem~5.5(II), pp.~251--252]{KirbySiebenmann1977}.
The TOP assertions follow by looping.

For PL, the jointly PL sphere fibre model
\eqref{eq:sphere-concordance-fiber} gives
\[
 P_{\mathrm{PL}}(S^3)\simeq
       \operatorname{hofib}(X_3\longrightarrow X_4),
       \qquad X_j=\mathrm{PL}_j/\mathrm O_j.
\]
The stabilization maps induce isomorphisms
$\pi_iX_j\cong\pi_i(\mathrm{PL}/\mathrm O)$ for $i<j$
\cite[Section~5, formula~(3), p.~37]{BurgheleaLashof1974II};
the stable groups vanish through degree six
\cite[Annex~C, p.~318]{KirbySiebenmann1977}.
The outer terms of
\[
 \pi_3X_4\longrightarrow\pi_2P_{\mathrm{PL}}(S^3)
                          \longrightarrow\pi_2X_3
\]
are therefore zero.  The same exact sequence gives vanishing
of the fundamental group.  For components, $X_3$ is connected
and $X_4$ is simply connected, so the based homotopy fibre is
connected as well.  This proves all the stated PL calculations
and the failure of the comparison on the second homotopy group.
\end{proof}


\subsection{Topological sphere concordances and finite-rank stabilization}
We express ordinary concordances of spheres by the failure of
finite-rank smoothing spaces to stabilize. The proof first uses
parametric germ extension to obtain a cone fibre, then identifies
its map with product by a line.


Let $\mathrm{TOP}(k)=\operatorname{Homeo}(\mathbb R^k;0)$ carry its
actual compact-open topology, and put
\begin{equation}\label{eq:finite-rank-top-smoothing-space}
 Y_k=\operatorname{hofib}_{*}
       \bigl(B\mathrm O(k)\longrightarrow B\mathrm{TOP}(k)\bigr).
\end{equation}
The classifying spaces may be formed from singular simplicial groups.
The orthogonal group has its usual topology, and its inclusion in
$\mathrm{TOP}(k)$ is the literal linear inclusion.  Product with a
line gives a based stabilization $Y_k\to Y_{k+1}$.

\begin{theorem}\label{thm:top-sphere-finite-rank-stabilization}
For every standard sphere $S^n$, $n\geq2$, the ordinary concordance
space with bottom fixed pointwise and top free satisfies
\begin{equation}\label{eq:top-sphere-concordance-fibre}
 P_{\mathrm{TOP}}(S^n)\simeq
       \operatorname{hofib}_{*}(Y_n\longrightarrow Y_{n+1}).
\end{equation}
These are the finite-rank spaces in
\eqref{eq:finite-rank-top-smoothing-space}.
\end{theorem}

\begin{proof}
We work with singular simplicial groups for actual homeomorphism
groups.  Write $\mathcal G_k$ for the simplicial group of germs at
the zero section: a simplex is represented on
$\Delta^q\times U$, for one open neighbourhood $U$ of $0$, by a
projection-preserving homeomorphism onto an open image that fixes
$\Delta^q\times\{0\}$.  Two representatives agree as germs when
they agree on a smaller common neighbourhood for the whole simplex.
This is the finite-parameter convention of
\cite[Section~3.4, equations~(3.5)--(3.7), pp.~14--15]
{SalvatoreTurchin2026}.

\medskip\noindent\textit{Global representatives and point stabilizers.}
Set $T_k=\operatorname{Sing}\mathrm{TOP}(k)$.  Restriction induces
a degreewise surjection
\[
                     T_k\longrightarrow\mathcal G_k.
\]
Indeed, a germ family has a common source ball and a common image
chart by compactness of its parameter simplex.  After fixed changes
of scale, Kuiper--Lashof's parametric extension lemma gives a
sphere homeomorphism family agreeing with it on a smaller ball and
fixing both distinguished points
\cite[Lemma~3, p.~11; TOP proof, pp.~12--13]{KuiperLashof1966I}.
Deleting the point at infinity gives the required global family.

A degreewise surjective map of simplicial groups is a Kan fibration.
Its kernel consists of families equal to the identity on a common
ball about $0$.  Sphere inversion identifies this kernel with
families of compactly supported homeomorphisms of $\mathbb R^k$
that fix $0$.  Every finite family has a common bounded support.
The formula
\[
 a_s(x)=s\,a(x/s)\quad(s>0),\qquad a_0(x)=x
\]
contracts such a family, together with its inverses, in the
compact-open topology.  For example, if the common support lies
in $D_R$, then the support at time $s$ lies in $D_{sR}$ and
$\sup\|a_s-\operatorname{id}\|\leq
s\sup\|a-\operatorname{id}\|$.  Identity parameter values stay
fixed.  Thus the kernel is weakly contractible and
$T_k\to\mathcal G_k$ is a weak equivalence on all components.

Fix $p\in S^n$ and an $\mathrm O(n)$-equivariant chart at $p$.
For $H_n=\operatorname{Sing}\operatorname{Homeo}(S^n;p)$ the
germ map $H_n\to\mathcal G_n$ is also degreewise surjective by
the same lemma.  Its kernel is the space of families equal to the
identity on a common neighbourhood of $p$.  On the complement of
$p$ these are compactly supported families on $\mathbb R^n$.
The same contraction works, without any additional fixed-point
condition.  Hence
\begin{equation}\label{eq:sphere-stabilizer-germ-equivalence}
 H_n\xrightarrow{\simeq}\mathcal G_n
                 \xleftarrow{\simeq}T_n.
\end{equation}
Both maps are the identity on the continuous orthogonal singular
subgroup specified by the chosen chart.

\medskip\noindent\textit{The disk germ fibre.}
Let $m=n+1$ and
$A_m=\operatorname{Sing}\operatorname{Homeo}(D^m;0)$.
The germ map $A_m\to\mathcal G_m$ is a Kan fibration.  We recall
the relative argument, which has no restriction on $m$.
The Edwards--Kirby local ambient-extension operator assigns to
every sufficiently small perturbation $v$ of the inclusion of a
compact ball $U\subset\operatorname{int}D^m$ a disk homeomorphism
$\sigma(v)$ equal to $v$ on a smaller ball $C$, supported away
from the outer boundary, with $\sigma(\operatorname{id})=
\operatorname{id}$
\cite[Theorem~5.1, p.~70; proof of Corollary~1.2, p.~79]
{EdwardsKirby1971}.  It depends continuously on $v$.

Suppose a germ simplex $f_t$ has a disk lift $a_t$ on a horn
$L\subset\Delta^q$.  Choose a deformation $d_s$ of the simplex
from a retraction onto $L$ to the identity, fixed on $L$.  On
common smaller neighbourhoods put
\[
 a^0_t=a_{d_0(t)},\qquad
 k_{t,s}=f_{d_s(t)}(a^0_t)^{-1}.
\]
The equalities $k_{t,0}=1$ and $k_{t,s}=1$ for $t\in L$ hold
on one common ball, after shrinking once.  Choose $U,C$ and a
still smaller ball $V$ so that all increments
$k_{t,s}k_{t,s'}^{-1}$ are defined on $U$ and
$k_{t,s}(V)\subset\operatorname{int}C$ for every $(t,s)$.
Joint continuity of maps and inverses and compactness give these
choices.  They also give a single finite partition of $I$ on
whose subintervals every increment belongs to the domain of
$\sigma$.  Successively multiplying their extensions produces
$b_{t,s}$ with germ $k_{t,s}$, equal to $1$ on the horn.  The
induction is exact on $V$, since the previous image of $V$ lies
in $C$.  Then $b_{t,1}a^0_t$ is a lift equal to $a_t$ on $L$.
Every extension fixes $0$ because the corresponding increment
does so near $0$.  This proves the fibration assertion with the
prescribed faces fixed exactly.

Its identity fibre is the filtered union of
\[
 K_\varepsilon=
 \operatorname{Sing}\operatorname{Homeo}
                  (D^m\;\mathrm{rel}\;D^m_\varepsilon),
 \qquad 0<\varepsilon<1.
\]
Restriction to the annulus identifies each $K_\varepsilon$ with
$P_{\mathrm{TOP}}(S^n)$: the inner boundary is fixed and the
outer boundary is free.  If $\varepsilon<\varepsilon_0$, radial
conjugations that move $D_{\varepsilon_0}$ to $D_\varepsilon$
give a homotopy in $K_\varepsilon$ from the inclusion
$K_{\varepsilon_0}\to K_\varepsilon$ to an isomorphism.
Each finite simplicial sphere or homotopy in the union has one
common fixed ball.  Therefore the fibre has the weak homotopy
type of $P_{\mathrm{TOP}}(S^n)$.

Write $G_n=\operatorname{Sing}\operatorname{Homeo}(S^n)$.
Boundary restriction $A_m\to G_n$ is a weak equivalence: radial
coning gives a section, and the boundary-fixed kernel, still
fixing $0$, contracts by the relative Alexander formula.
Extend a disk map outside the disk by coning its boundary map.
This gives a homomorphism $A_m\to T_m$ compatible with germs.
On the cone section it is the radial homomorphism
\[
 c:G_n\longrightarrow T_{n+1},\qquad
                  c(h)(r\omega)=r h(\omega).
\]
The preceding weak equivalences and the disk germ fibration give
\begin{equation}\label{eq:top-sphere-cone-fibre}
             P_{\mathrm{TOP}}(S^n)\simeq\operatorname{hofib}_1(c).
\end{equation}
The global germ equivalence here includes all components.

\medskip\noindent\textit{Removing orthogonal directions.}
The cone map is the linear inclusion on
$\operatorname{Sing}\mathrm O(n+1)$.  Taking right cosets by
this same simplicial subgroup preserves its homotopy fibre:
\[
 \operatorname{hofib}(c)\simeq
 \operatorname{hofib}\left(
   G_n/\operatorname{Sing}\mathrm O(n+1)
       \longrightarrow T_{n+1}/\operatorname{Sing}\mathrm O(n+1)
                         \right).
\]
This follows from the two simplicial subgroup fibrations, whose
map on fibres is the identity.  Such a quotient models the
homotopy fibre of the corresponding map of classifying spaces.
We use these simplicial subgroup quotients throughout.

The point stabilizer gives
\[
 H_n/\operatorname{Sing}\mathrm O(n)
       \xrightarrow{\simeq}
 G_n/\operatorname{Sing}\mathrm O(n+1).
\]
To see this directly, for a family $g_t$ on a simplex choose
rotations $R_t$ taking $p$ to $g_t^{-1}(p)$.  Such a continuous
choice exists because the pullback of
$\mathrm O(n+1)\to S^n$ to a simplex is a trivial principal
$\mathrm O(n)$-bundle.  Then $g_tR_t$ fixes $p$, and two such
representatives differ precisely by a family in $\mathrm O(n)$.
Together with \eqref{eq:sphere-stabilizer-germ-equivalence},
this identifies the source quotient with $Y_n$ and the target
quotient with $Y_{n+1}$.

\medskip\noindent\textit{Identifying the arrow.}
It remains to identify this map, not just its source and target.
Take $p=(0,\ldots,0,1)$.  If $h$ fixes $p$, its cone fixes the
whole ray through $p$.  Consequently
\[
 F_s(h)=\tau_{-sp}\,c(h)\,\tau_{sp},\qquad 0\leq s\leq1,
\]
is a homotopy of homomorphisms from $H_n$ to $T_{n+1}$, all
fixing $0$.  It is constant on each
$\operatorname{diag}(A,1)$ with $A\in\mathrm O(n)$.
Near the origin use
\[
 \theta(u)=(u,\sqrt{1-\|u\|^2}),\qquad
 \psi(u,t)=(1+t)\theta(u)-p.
\]
The cone preserves radius and applies $h$ to direction.  Thus
the coordinates centred at $p$ give the exact germ identity
\begin{equation}\label{eq:cone-germ-line-stabilization}
 \psi^{-1}F_1(h)\psi(u,t)
               =\bigl(\theta^{-1}h\theta(u),t\bigr).
\end{equation}
The right side is product with a line on the point germ of $h$.

The coordinate change does not alter the homotopy class of this
map.  For $r>0$, let $\psi_r(z)=r^{-1}\psi(rz)$ and set
$\psi_0=\operatorname{id}$.  These form a continuous family
of local diffeomorphisms with continuous inverses on common
neighbourhoods.  For example, writing $y=(v,w)$ and
$N_r=\|(rv,1+rw)\|$, the inverse has the nonsingular expression
\[
 \psi_r^{-1}(v,w)=
 \left(\frac{v}{N_r},\,
       \frac{2w+r(\|v\|^2+w^2)}{N_r+1}\right),
 \qquad N_0=1.
\]
This verifies the extension at $r=0$ directly.  The family
commutes with $\operatorname{diag}(A,1)$ for every
$A\in\mathrm O(n)$.  Conjugation by it therefore provides
the required homotopy of quotient maps, relative to the source
orthogonal subgroup.  A finite parameter family of $h$ uses
one smaller chart; identity parameter faces remain fixed.
Equation~\eqref{eq:cone-germ-line-stabilization} now identifies
the arrow in \eqref{eq:top-sphere-cone-fibre} with the actual
stabilization $Y_n\to Y_{n+1}$, proving the theorem.
\end{proof}

\begin{corollary}\label{cor:four-sphere-low-concordance-groups}
For the standard four-sphere,
\[
 \pi_iP_{\mathrm{TOP}}(S^4)=\pi_iP_{\mathrm{PL}}(S^4)=0
                       \qquad(0\leq i\leq3).
\]
For $i=0$, both component sets are singletons.
\end{corollary}

\begin{proof}
The stabilization $Y_4\to\mathrm{TOP}/\mathrm O$ is
$5$-connected by
\cite[Theorem~8.7A, p.~128]{FreedmanQuinn1990}, with the
corrected account in
\cite[Section~1.1, Theorem~1.3 and the following paragraph,
p.~2]{GabaiGayHartmanKrushkalPowell2026}.
The latter uses precisely the global origin-fixing group and
classifying-space fibre convention in
\eqref{eq:finite-rank-top-smoothing-space}.

For rank five, use the direct finite-rank statement
\cite[Annex~C, Theorem~5.3, p.~317]{KirbySiebenmann1977}:
\[
 \pi_i(\mathrm{TOP}_m/\mathrm{CAT}_m)
       \xrightarrow{\cong}\pi_i(\mathrm{TOP}/\mathrm{CAT})
       \quad(i<m,\ m\geq5),
       \qquad \mathrm{CAT}=\mathrm{PL}\text{ or }\mathrm{DIFF}.
\]
The source defines these as classifying-space fibres; its DIFF
version uses the orthogonal classifying space.  The orthogonally compatible germ
comparison above therefore identifies its rank-five map with
$Y_5\to\mathrm{TOP}/\mathrm O$.  It is an isomorphism on
$\pi_i$ for $i\leq4$.  Naturality of stabilization shows that
$Y_4\to Y_5$ is also an isomorphism in those degrees.  The
homotopy-fibre sequence of the theorem gives the TOP vanishing
through degree three, including the component assertion.

For PL use \eqref{eq:sphere-concordance-fiber} with $n=4$.
Here $X_j=\mathrm{PL}_j/\mathrm O_j$ denotes the classical
unstable smoothing homotopy type.  Its classifying-space
definition does not assert a literal inclusion of continuous
orthogonal singular chains into joint-PL germ chains; a precise
construction via piecewise differentiable germs is recalled in
\cite[Section~3.4, Remark~3.7 and footnote~16, p.~16]
{SalvatoreTurchin2026}.  The classical stabilization calculation
\cite[Section~5, formula~(3), p.~37]{BurgheleaLashof1974II}
and the stable vanishing
\cite[Annex~C, p.~318]{KirbySiebenmann1977} give
\[
              \pi_iX_4=0\ (i\leq3),\qquad
              \pi_iX_5=0\ (i\leq4).
\]
The exact sequence of $\operatorname{hofib}(X_4\to X_5)$
proves the PL assertion.  Its source is connected and its
target is simply connected, so the same conclusion includes
components.
\end{proof}

The argument does not determine $\pi_4P_{\mathrm{TOP}}(S^4)$:
the displayed stabilization information does not control
$\pi_5Y_4\to\pi_5Y_5$.  It also leaves the corresponding
ordinary concordance spaces for general four-manifold bases
unevaluated.


\subsection{Vanishing components for balls and specified surfaces}

\begin{theorem}\label{thm:low-concordance-components}
With the bottom and side boundary fixed and the upper face free,
\[
                         \pi_0P_{\mathrm{PL}}(B)=0
\]
in each of the following cases: a compact manifold of dimension zero
or one; a PL ball of any dimension; a compact connected orientable surface other
than $S^2$; or a closed connected nonorientable surface other than
$\mathbb{RP}^2$.  The assertion also holds for finite disjoint unions
of these manifolds and for their stipulated definable models over every
real closed field.
This is a statement about components.  The surface clauses do not
include $S^2$, $\mathbb{RP}^2$, or nonorientable surfaces with boundary.
\end{theorem}

\begin{proof}
We will combine PL isotopies by means of the piecewise linear function
\[
              r(s,t)=\max(0,s+t-1)\qquad(s,t\in I).
\]
It vanishes on the bottom and left edges, and restricts to $t$ and $s$
on the right and top edges.  Composition with this function keeps total
parameter maps PL; the product $st$ would not have that property.

\emph{Balls.}
Let $h\in P_{\mathrm{PL}}(D^m)$, and denote its upper-face restriction
by $a$.  This homeomorphism fixes $\partial D^m$.  There is a jointly
PL Alexander isotopy $a_u$ from the identity to $a^{-1}$ relative to
that boundary \cite[Theorem~5.2, pp.~95--96]{Epstein1966}.
One can see joint PL regularity directly on $D^m=[-1,1]^m$:
\[
 a_u(x)=
 \begin{cases}
  u\,a^{-1}(x/u),&\lVert x\rVert_\infty\leq u,\quad u>0,\\
  x,&\lVert x\rVert_\infty\geq u,
 \end{cases}
 \qquad a_0(x)=x.
\]
On a cell where $a^{-1}(y)=Ay+b$, the first expression is $Ax+ub$.
Its domain is the cone over that cell at height one.  These finitely
many cones and the outer region give a finite PL map; the formulas
match because $a$ fixes the boundary.  The analogous construction with
$a$ gives the inverse family.

The family
\[
              A_s(x,t)=(a_{r(s,t)}(x),t)
\]
starts at the identity and consists of concordances.  Its endpoint
makes $A_1h$ fix the whole boundary of the PL ball $D^m\times I$.
A second Alexander isotopy, now relative to that whole boundary,
connects $A_1h$ to the identity.  Concatenation proves the assertion for
balls, including dimension zero.

\emph{The circle.}
The only remaining connected one-manifold is $S^1$.  A concordance
of the annulus $S^1\times I$ has an orientation-preserving upper
restriction.  This circle homeomorphism admits a jointly PL isotopy
to the identity.  For an explicit construction, lift it to an
increasing PL map $a:\mathbb R\to\mathbb R$ with
$a(x+1)=a(x)+1$.  Subdivide one period at its breakpoints, triangulate
each resulting rectangle in the parameter strip, and extend the
bottom identity and the top map affinely on those triangles.  On each
horizontal piece, the slope is either $1$ or the positive slope of
$a$ on the corresponding interval.  Each horizontal slice is therefore
increasing, and periodicity makes the construction descend to the
circle.  Using $r(s,t)$ as above adjusts the upper restriction without
moving the bottom.

We may consequently assume that a given annulus concordance $h$ fixes
both boundary circles.  Let $\alpha=\{0\}\times I$.  Its image has an
integer winding $k$ relative to the endpoints, measured in
$\mathbb R\times I$.  The homeomorphism
\[
                   T_{-k}([x],t)=([x-kt],t)
\]
removes that winding and itself lies in the identity component of the
concordance group, by
\[
                   T_{-k,s}([x],t)=([x-kr(s,t)],t).
\]
The upper circle is allowed to move during this isotopy.  The adjusted
arc is homotopic to $\alpha$ relative to its endpoints, so Epstein's
relative arc theorem supplies a PL ambient isotopy fixed on both
boundary circles carrying it to $\alpha$
\cite[Theorem~3.1 and proof, pp.~89--90]{Epstein1966}.
Correcting the arc image does not yet fix its parameterization.
In a rectangular neighborhood $(t,z)\in I\times[-1,1]$ of the arc,
let $\lambda$ be the remaining interval homeomorphism and choose
a jointly PL interval isotopy $c_s$ from the identity to
$\lambda^{-1}$, fixed at both endpoints. The family
\[
 \Phi_s(t,z)=\bigl(c_{\max(0,s-|z|)}(t),z\bigr)
\]
is jointly PL, fixes the rectangle boundary, and extends by the
identity outside the rectangle. Its inverse uses the inverse
interval family. At its endpoint it corrects the central arc
pointwise; this is the regular-neighborhood step in the argument
on p.~91 of the same source.

The annular map is orientation preserving, so after this correction
it preserves the two sides of the arc. Cutting along the arc gives
a disk homeomorphism fixed on each copy of the cut arc and on the
remaining boundary. Its relative Alexander isotopy descends when
the two cut edges are glued. Thus a winding number does not survive
when the upper face is free.

\emph{Orientable surfaces.}
Let $F$ be compact, connected, orientable, and different from $S^2$.
Waldhausen works in the PL category and imposes exactly these surface
hypotheses in his product-line-bundle section
\cite[pp.~56 and~65]{Waldhausen1968}.
His Lemma~3.5, including its proof on p.~69, changes a
bottom-and-side-fixed homeomorphism of $F\times I$ by an isotopy
constant on the full boundary into a level-preserving map
\[
                         (x,t)\longmapsto(h_t(x),t).
\]
Here $h_0=\operatorname{id}$ and every $h_t$ fixes $\partial F$.
The jointly PL family
\[
                         (x,t,s)\longmapsto(h_{r(s,t)}(x),t,s)
\]
joins the identity to this level-preserving concordance.  Its inverse
uses the inverse level-preserving family.  This proves the orientable
surface clause.

\emph{Closed nonorientable surfaces.}
Now let $F$ be closed, nonorientable, and different from
$\mathbb{RP}^2$, and put $V=F\times I$.
We verify the hypotheses of the nonorientable three-manifold theorem.
The universal cover of $F$ is the plane and $\pi_1F$ is torsion free
\cite[Corollary~1.8, pp.~85--86; Lemma~4.3, p.~92]{Epstein1966}.
Thus the universal cover of $V$ is $\mathbb R^2\times I$.
A PL sphere in the interior lifts to a sphere bounding a PL ball in
this slab.  Distinct translates of its lifted sphere are disjoint,
so their bounded balls are either disjoint or nested.  Proper
discontinuity and infinite order of a nonidentity deck transformation
exclude nesting: after possibly taking its inverse, its iterates would
put an infinite orbit in one compact ball.  The ball therefore projects
injectively to $V$.  This proves irreducibility.

A two-sided embedded $\mathbb{RP}^2$ in $V$ would lift to the
orientable slab, since its fundamental group maps trivially to the
torsion-free group $\pi_1V$.  A two-sided surface in an orientable
three-manifold is orientable, a contradiction.
The middle surface $F\times\{1/2\}$ is two-sided and
$\pi_1$-injective.  It cannot bound a contractible submanifold because
its nontrivial fundamental group injects into that of $V$.
Consequently $V$ is $\mathbb{RP}^2$-irreducible and sufficiently large
in the precise sense of
\cite[Chapter~II, Definitions~1.2--1.5, pp.~15--17]{Laudenbach1974}.

Write a concordance as $h(x,t)=(u(x,t),v(x,t))$.  It is properly
homotopic to the identity in the boundary-pair sense: first change the
second coordinate by
\[
             (x,t)\longmapsto
             \bigl(u(x,t),(1-s)v(x,t)+st\bigr),
\]
and then change the first coordinate by
\[
             (x,t)\longmapsto\bigl(u(x,(1-s)t),t\bigr).
\]
Throughout these homotopies each boundary component maps to itself and
interior points map to the interior.  These are homotopies of maps;
there is no claim that their intermediate maps are homeomorphisms.

Laudenbach's Theorem~II.6.1 and proof, pp.~46--47, turn this proper
homotopy class into an isotopy class.  To use its smooth formulation
for the given PL endpoints, use the category comparisons in
Chapter~I, Section~1, pp.~7--9: a PL homeomorphism is TOP isotopic to
a smooth one, and the map from PL to TOP isotopy classes is
injective.  Composing the TOP isotopies and then applying this
injectivity gives a jointly PL isotopy $H_s$ from $h$ to the identity.
This argument does not require an auxiliary smooth map to be PL in a
fixed triangulation.

Finally enforce the fixed bottom.  Each $H_s$ preserves the two boundary
components individually.  Write $b_s$ for its restriction to the
bottom, identified with $F$, and replace it by
\[
                   \widehat H_s=(b_s^{-1}\times\operatorname{id}_I)H_s.
\]
The inverse of the total PL boundary family is PL, so this remains a
jointly PL isotopy.  It fixes the bottom, and its endpoints remain $h$
and the identity because $b_0=b_1=\operatorname{id}$.

The fixed bottom also prevents permutations of connected components.
Finite disjoint unions therefore give products.  The relative
PL--definable component comparison transfers the result, with the side
boundary included in the fixed subpolyhedron, to every real closed
field and the stipulated definable category.
\end{proof}


\begin{corollary}\label{cor:all-closed-surface-components}
For every closed PL surface, the bottom-fixed, top-free concordance
component group is zero. This includes disconnected surfaces and
their stipulated definable models over every real closed field.
\end{corollary}

\begin{proof}
Every connected closed surface is orientable or nonorientable.
The preceding theorem covers all such surfaces except $S^2$ and
$\mathbb{RP}^2$. Theorems~\ref{thm:standard-sphere-components}
and~\ref{thm:projective-plane-components} give those two cases.
The fixed bottom makes a finite disjoint union a product on both
groups and component relations.
\end{proof}


% Authored continuation for original A.130. Complete targeted proof.
\subsection{Infinite-rank subgroups in dimension three}

Lin, Xie and Zhang construct infinite-rank families of smooth
concordances detected by topological isotopy.  Their stronger detection
theorems also give the following PL and definable consequence.

\begin{theorem}\label{thm:infinite-rank-three-components}
Let $B$ be a closed connected orientable PL three-manifold with
infinite fundamental group.  Then there is a group monomorphism
\[
             \bigoplus_{j\geq1}\mathbb Z
                   \longrightarrow\pi_0P_{\mathrm{PL}}(B).
\]
Over the real field these classes remain independent in
$\pi_0\operatorname{Homeo}(B\times I)$, where no boundary is fixed.
The same subgroup conclusion holds for the stipulated definable
realizations of $B$ over every real closed field.  In particular it
holds for a closed definably orientable definable three-manifold
whose definable fundamental group is infinite, with its stipulated
finite PL model.
\end{theorem}

\begin{proof}
Choose a compatible smoothing of $B$
\cite[Essay~I, p.~39]{KirbySiebenmann1977}, and put
$Z=B\times I$ and $K=B\times\{0\}$.  Choose a differentiable
triangulation $t:|Z|\to Z$ compatible with $K$ and with the given
PL product structure.

We first specify the comparison on components.  Write
$P_{\mathrm{PD}}(Z,t,K)$ for the semisimplicial complex of
piecewise differentiable homeomorphisms from the triangulated
$|Z|$ to $Z$, prescribed to equal $t$ on $K$.  Its parameter maps
preserve the entire parameter simplex.  This complex has commuting
left actions by smooth maps and right actions by PL maps; no group
structure on it is required.  There are maps
\[
 \begin{aligned}
 j:P_{\mathrm{PL}}(B)&\longrightarrow P_{\mathrm{PD}}(Z,t,K),
       &p&\longmapsto t\circ p,\\
 i:P_{\mathrm{Diff}}(B)&\longrightarrow P_{\mathrm{PD}}(Z,t,K),
       &h&\longmapsto h\circ t.
 \end{aligned}
\]
The relative PL--PD triangulation comparison makes $j$ a homotopy
equivalence in every dimension
\cite[pp.~1--2; Theorem~1.1(c), p.~8]{BurgheleaLashof1974I}.
Consequently
\[
 c=(\pi_0j)^{-1}\pi_0i:
     \pi_0P_{\mathrm{Diff}}(B)\longrightarrow
     \pi_0P_{\mathrm{PL}}(B)
\]
is defined.  It is a group homomorphism: if the two PD paths join
$h_at$ to $tp_a$ for $a=1,2$, the two actions concatenate them to
\[
                  h_1h_2t\simeq h_1tp_2\simeq tp_1p_2.
\]
Injectivity of $\pi_0j$ gives the required multiplicativity.
Forgetting the PD structure turns these paths into topological
isotopies.  Thus the composite of $c$ with the PL-to-TOP map agrees
with the smooth-to-TOP map on components, both with and without the
bottom condition.

The convention in the cited comparison requires maps to be
homotopic to the identity as maps of the boundary pair, relative to
$K$.  It includes every concordance here.  Indeed, write
$h(x,u)=(a(x,u),b(x,u))$.  First linearly replace $b(x,u)$ by $u$;
then use
\[
                   (x,u)\longmapsto(a(x,(1-s)u),u),
                   \qquad 0\leq s\leq1.
\]
Both homotopies preserve the boundary pair and fix $K$.  They are
homotopies of maps, which is all that this source convention asks.

Theorems~C and~D of Lin--Xie--Zhang apply respectively to the
irreducible and reducible cases
\cite[arXiv version~1, Theorem~C, p.~13 and its deduction on p.~14;
Theorem~D, p.~28]{LinXieZhang2026}.
For a closed $B$, their target fixing only $I\times\partial B$
is the full homeomorphism group.  They show that the image of
\[
 \pi_0\operatorname{Diff}_{\partial}(B\times I)
       \longrightarrow\pi_0\operatorname{Homeo}(B\times I)
\]
has infinite rank.  Here $\operatorname{Diff}_{\partial}$ means
identity near the whole boundary.  Its component group is abelian:
two representatives can be isotoped into disjoint subintervals
of the cylinder and then commute; see also the explicit observation
on p.~1 of that paper.

Choose classes $d_j$ whose images in this abelian image group are
linearly independent over $\mathbb Z$.  Such a countable choice
exists because tensoring that image with $\mathbb Q$ has infinite
dimension.  Include the $d_j$ in the smooth free-top concordance
group and apply $c$.  Their PL images commute.  Every integral
relation among them would give the identical relation in the full
topological homeomorphism component group, and hence would have
all coefficients zero.  They therefore generate a copy of
$\bigoplus_{j\geq1}\mathbb Z$ in the asserted PL group.  This also
proves the stated persistence after forgetting all boundary
conditions; no injectivity assertion about $c$ on other classes
has been used.

Finally, Theorems~\ref{result-40} and~\ref{result-4} transport the
PL subgroup to the stipulated definable models and preserve the
group operations.  For the last formulation, a rational realization
of the finite abstract model has the same orientability, as read
from its finite simplex incidences.  Its real and definable
fundamental groups agree by
\cite[arXiv version~1, Corollary~4.3, p.~10]{BaroOtero2008v1},
so the real model satisfies the hypotheses just used.
\end{proof}

This theorem supplies a subgroup.  It does not describe the entire
three-dimensional concordance group, and its proof requires no
dimension-three smooth-to-PL isomorphism on components.


\subsection{Effective countable concordance models}

Classical decision procedures over real closed fields and the Moore
description of simplicial-group homotopy give an effective version of
the finite-model comparisons.  The output is a complete countable
presentation; it need not solve the word problem.

\begin{theorem}\label{thm:effective-countable-concordance-models}
Given an ordered finite simplex list for a finite abstract complex $K$
whose realization is a closed PL manifold, one can construct countable
simplicial groups $G^{\mathrm{PL}}(K)$ and $G^{\mathrm{SA}}(K)$ with
the following properties.  Their simplices have finite integer codes;
validity, equality and degeneracy are decidable.  All group operations,
faces and degeneracies are computable by terminating algorithms, as
is a simplicial homomorphism
\[
                   G^{\mathrm{PL}}(K)\longrightarrow G^{\mathrm{SA}}(K).
\]
For every real closed field $R$ the groups, interpreted over the real
algebraic numbers, give weak equivalences
\[
 \begin{aligned}
 G^{\mathrm{PL}}(K)&\longrightarrow P_{\mathrm{PL},R}(|K|_R),\\
 G^{\mathrm{SA}}(K)&\longrightarrow P_{\mathrm{SA},R}(|K|_R).
 \end{aligned}
\]
In these concordance groups the total map preserves its entire
parameter simplex, fixes the bottom pointwise, and preserves the
free top setwise. In the PL case the total map and its inverse
are finite piecewise linear; in the SA case they are semialgebraic.
In both cases they are jointly continuous on the whole product.

Uniformly in $K$ and $r\geq0$, an algorithm enumerates a complete
countable group presentation of each of these homotopy groups.
For $r\geq1$ it can instead enumerate a countable integer relation
matrix for the resulting abelian group, based at the identity.
Translation gives the homotopy groups at other components.
The PL presentation in degree zero also presents the stipulated
definable component group in every fixed o-minimal expansion.
If every component of $|K|$ has dimension at least five, the PL
construction is also a weak model for the real TOP concordance
space and presents all of its homotopy groups.

The empty complex gives the constant one-element simplicial group.
The manifold condition is a promise on the input; no recognition
algorithm for closed PL manifolds is asserted.  The presentations
need not be finite, and no terminating word problem, invariant-factor
calculation or complexity bound is claimed.
\end{theorem}

\begin{proof}
Let $F=\mathbb Q^{\mathrm{rc}}$.  Realize the ordered vertices of $K$
as distinct standard basis vectors, and put
\[
 B=|K|_F,\qquad Z=B\times[0,1],\qquad
 A=B\times\{0\},\qquad T=B\times\{1\}.
\]
These are finite rational polyhedra.  A $q$-simplex will be described
by a graph in the variables $(z,z',t)$, with $z,z'\in Z$ and
$t\in\Delta^q_F$.  The repeated parameter is implicit: its total
map is $(z,t)\mapsto(z',t)$.

\medskip\noindent\textit{Effective scalar and formula operations.}
Represent an element of $F$ by a nonzero integer polynomial and a
rational open interval containing exactly one of its real roots,
with neither endpoint a root.  Such codes are enumerable and their
validity is decidable by Sturm's theorem.  A finite list of algebraic
constants can alternatively be replaced in any formula by variables
subject to their root and isolating-interval conditions.  These
conditions have unique solutions.  Decision and quantifier
elimination over the rational coefficients therefore decide equality,
order and all first-order tests with the encoded constants.
To compute a scalar sum, product or inverse of a nonzero scalar,
enumerate valid scalar codes and test the appropriate equality; a
code exists and the search terminates.  These are classical effective
procedures; the exact locators used here are
\cite[online revision of 26 June 2016, Theorem~2.62, p.~62;
Algorithm~14.16, pp.~586--587; Algorithm~14.21 and Theorem~14.22,
pp.~591--593]{BasuPollackRoy2016Online}.

We also use the elementary linear version of elimination.  Put a
Boolean combination of affine equalities and inequalities in
disjunctive normal form.  Equalities with a nonzero eliminated
coefficient allow substitution.  In the other cases the inequalities
give finitely many upper and lower bounds for that variable.  Existence
means that every lower bound is at most every upper bound, strictly
so whenever either bound was strict.  With no bound on one side
there is no additional condition.  This eliminates one existential
variable; negation handles universal quantifiers.  Repetition gives
a terminating linear decision and elimination procedure over $F$.

\medskip\noindent\textit{PL graphs and their validity.}
A raw PL code is a finite list of closed simplices in the graph
coordinates, with vertices in $F$ and with affine independence
checked.  It describes the union $\Gamma$ of these simplices.
Barycentric coordinates express membership by a linear formula.
Require the following five decidable properties:
\begin{enumerate}
\item $\Gamma\subset Z\times Z\times\Delta^q_F$;
\item for each $(z,t)\in Z\times\Delta^q_F$ there is exactly one
      $z'$ with $(z,z',t)\in\Gamma$;
\item for each $(z',t)\in Z\times\Delta^q_F$ there is exactly one
      $z$ with $(z,z',t)\in\Gamma$;
\item $(z,z,t)\in\Gamma$ for every $z\in A$ and $t\in\Delta^q_F$;
\item on $\Gamma$, $z\in T$ if and only if $z'\in T$.
\end{enumerate}
All these statements are linear first-order sentences.  Equality
of two coded graphs is a linear sentence as well.

The two graph projections of a valid code are PL homeomorphisms.
Indeed, each is injective on each graph simplex, hence on its affine
hull: a nonzero kernel direction there would give two distinct points
in the relative interior with the same image.  Its inverse on the
projected simplex is therefore affine.  The projected simplices are
closed, form a finite cover, and their inverses agree on intersections.
Finite closed pasting gives continuity of both total maps over $F$.
The finite polyhedral family has a compatible triangulation over $F$:
intersect the finitely many face hyperplanes, then triangulate the
resulting polyhedral cells compatibly, for instance by their face
barycentres.  Thus the map and its inverse are jointly PL.  This
argument uses no ordinary compactness assertion for the ordered field.
Conversely, every jointly PL concordance on the total product has a
finite graph triangulation of this kind, so the raw codes represent
exactly $P_{\mathrm{PL},F}(B)_q$.

Encode the finite strings in a fixed finite alphabet and order them
first by length and then lexicographically.  Replace a valid raw code
by the first valid code in this order defining its graph.  It is
enough to examine its finitely many predecessors.  Validity and
graph equality make this a terminating canonicalization algorithm;
the canonical codes themselves form a decidable set.  They give
$G^{\mathrm{PL}}(K)_q$.

For multiplication, form the linear graph formula
\[
  \exists w\;\bigl(\Gamma_2(z,w,t)\ \wedge\ \Gamma_1(w,z',t)\bigr).
\]
For inversion interchange $z$ and $z'$. For a simplicial operator
$\alpha:\Delta^p_F\to\Delta^q_F$, define the pullback graph by
\[
 \Gamma_{\alpha^*g}(z,z',t)=
 [\,z\in Z,\ z'\in Z,\ t\in\Delta^p_F\,]
 \ \wedge\ \Gamma_g(z,z',\alpha(t)).
\]
The displayed domain conditions are part of the formula. All graph
equalities are tested in the full ambient coordinate space; in
particular, no points with $t\notin\Delta^p_F$ are admitted after an
affine parameter substitution. The identity is the diagonal graph
on $Z\times\Delta^q_F$. These operations are effective on formulas.
They describe jointly PL homeomorphisms: finite PL compositions,
inverses and these affine parameter pullbacks admit finite compatible
polyhedral subdivisions.  Enumerate valid raw codes and compare their
graphs with the resulting formula until equality holds, then
canonicalize.  Existence of such a graph code proves termination.
The identity has its diagonal graph code.  Consequently all operations
are total on valid inputs; an invalid input can be recognized and
rejected first.  Because the canonical graphs are actual maps,
the group and simplicial identities hold exactly on their codes.

\medskip\noindent\textit{Semialgebraic graphs.}
A raw SA code is a finite Boolean formula in polynomial equalities
and inequalities over the encoded field $F$.  Impose the same five
conditions and, in addition, continuity of the two total maps.
For example, the forward condition is the first-order sentence
\[
\begin{split}
 \forall (z,z',t)\in\Gamma\ \forall\varepsilon>0\ \exists\delta>0\
 \ \forall(w,w',u)\in\Gamma:\qquad\qquad\qquad\quad\\
 \|(w,u)-(z,t)\|^2<\delta^2
 \ \Longrightarrow\
 \|(w',u)-(z',t)\|^2<\varepsilon^2.
\end{split}
\]
Interchanging the source and target variables gives the inverse
condition.  This tests continuity in the parameter as well as in
the fibre.  Polynomial quantifier elimination decides all these
sentences and equality of two graphs.  Hence the same finite-predecessor
canonicalization gives a decidable set $G^{\mathrm{SA}}(K)_q$.
Every semialgebraic concordance gives such a graph and conversely.

Composition, inversion and parameter substitution are the same graph
operations as above.  Quantifier elimination puts their formulas in
the permitted raw form, after which canonicalization is effective.
These give the group and simplicial operations.  A PL graph has an
effective semialgebraic formula, so canonicalizing that formula gives
the required computable simplicial homomorphism.  In fact the two
constructed simplicial groups are coded copies of the respective
concordance groups over $F$.  Their asserted weak comparisons with
other real closed fields are exactly Theorem~\ref{result-40} and
Corollary~\ref{result-52}.

Degeneracy is also decidable in either model.  A $q$-simplex $a$,
with $q>0$, is degenerate precisely when
\[
                   a=s_i d_i a\quad\text{for some }0\leq i<q.
\]
If $a=s_i b$, the identity $d_i s_i b=b$ gives the displayed test;
the converse is immediate.  A degeneracy composite lies in the
image of some single $s_i$, and there are only finitely many indices
to test.  No zero-simplex is degenerate.

\medskip\noindent\textit{Complete homotopy presentations.}
Let $G$ be either coded simplicial group, and let $1$ denote the
identity in each degree.  Define decidable sets
\[
\begin{aligned}
 Z_0&=G_0,\\
 Z_r&=\{a\in G_r:d_i a=1\text{ for }0\leq i\leq r\}
                    \quad(r\geq1),\\
 N_{r+1}&=\{u\in G_{r+1}:d_i u=1\text{ for }1\leq i\leq r+1\}
                    \quad(r\geq0).
\end{aligned}
\]
The simplicial identities show $d_0N_{r+1}\subset Z_r$.
Conjugating $u\in N_{r+1}$ by $s_0a$, for $a\in Z_r$, stays in
$N_{r+1}$ and sends its first face to $a(d_0u)a^{-1}$.
Thus the image is a normal subgroup of $Z_r$.

The classical Moore description gives
\[
                 \pi_r(G,1)=Z_r/d_0N_{r+1},\qquad r\geq0,
\]
where degree zero means the component group.  For clarity about
conventions, May uses the last face: see
\cite[Theorem~17.1, pp.~67--68; both propositions numbered~17.3,
pp.~68--69; Proposition~17.4, p.~69]{May1967Simplicial}.
In the first-face convention an elementary homotopy from a cycle
$y$ to a cycle $x$ has a simplex $v$ with $d_0v=x$, $d_1v=y$,
and all remaining faces equal to $1$.  Then
$u=v(s_0y)^{-1}$ lies in $N_{r+1}$ and $d_0u=xy^{-1}$.
Conversely such a $u$ gives $v=u(s_0y)$.  Simplicial groups are
Kan, so this is precisely the usual cycle homotopy relation.
Multiplication of cycles gives the homotopy-group operation;
May's Proposition~17.2, p.~68, records this identification.
Theorem~16.6, p.~65, identifies these with the homotopy groups
of the realization.

Now enumerate the presentation
\[
 \left\langle x_a\ (a\in Z_r)\ \middle|\
 x_1=1,\quad x_ax_b=x_{ab}\ (a,b\in Z_r),\quad
 x_{d_0u}=1\ (u\in N_{r+1})\right\rangle.
\]
All generator and relator indices use the decidable tests and
terminating operations already constructed. At stage $m$, examine all
strings and pairs of strings of index at most $m$, printing the
valid cycle generators, multiplication relations, and boundaries.
Each stage terminates and every required item eventually occurs.
This does not wait for a next valid code when only finitely many exist.
The multiplication-table relations present $Z_r$ itself: the
assignments $x_a\mapsto a$ and $a\mapsto x_a$ are inverse
homomorphisms.  The boundary relations then quotient by exactly
the normal subgroup $d_0N_{r+1}$.  This proves completeness of the
presentation, not merely a surjection onto the desired group.

For $r\geq1$ this quotient is abelian, including $r=1$; this is
the first Proposition~17.3 on May's p.~68.  Therefore the equivalent
abelian presentation has free generators $e_a$ and relation rows
\[
                   e_1,\qquad e_a+e_b-e_{ab},\qquad e_{d_0u}.
\]
Every row has finite support and the rows can be enumerated uniformly.
This gives the asserted countable integer relation matrix.  The
boundary image is recursively enumerable; deciding membership in it
has not been used or asserted.

The degree-zero definable comparison and the real PL--TOP weak
comparison in the stated dimension range follow from
Theorem~\ref{result-4}.  They use a chosen real realization, not a
map of nonreal points into a TOP space.  The construction preserves
components of a disconnected base because the bottom is fixed.
The comparisons therefore apply componentwise.  For an empty base
use the one-element group in every degree.  This completes the
construction and all the assertions.
\end{proof}

The theorem constructs a computable PL-to-SA map but does not use that
map as an additional higher-comparison theorem.  Its full-space
identifications are the separately established coefficient comparisons.


\subsection{Real PL--TOP comparison in dimensions at most two}

The following comparison is a consequence of the classical joint-family
theorems for three-manifolds and the classical annulus calculation.
The component conventions in those inputs require attention because
the concordance groups here contain all components.

\begin{theorem}\label{thm:low-dimensional-real-concordance-comparison}
Let $B$ be a compact real PL manifold without boundary, every component
of which has dimension at most two.  The natural map
\[
 P_{\mathrm{PL}}(B)\longrightarrow
 \operatorname{Sing}P_{\mathrm{TOP}}(B)
\]
is a weak equivalence.  The source uses jointly PL homeomorphisms on
the whole product with a parameter simplex.  The target is the singular
complex of the compact-open topological concordance group.  In both
categories the bottom is fixed pointwise and the top is free.

If every component of $B$ has dimension zero or one, both concordance
objects are weakly contractible.  For every closed real PL surface $F$
there is also a weak equivalence on the full automorphism groups
\[
 \operatorname{Aut}_{\mathrm{PL}}(F)
 \longrightarrow \operatorname{Sing}\operatorname{Homeo}(F).
\]
Orientable and nonorientable surfaces are included.  The empty base
gives the constant one-element simplicial group.
\end{theorem}

\begin{proof}
For a compact PL manifold $N$, write $\mathcal H_{\mathrm{PL}}(N)$
for the simplicial group of all parameter-preserving joint PL
homeomorphisms of $N\times\Delta^q$, and write
$\mathcal H_{\mathrm{TOP}}(N)=\operatorname{Sing}\operatorname{Homeo}(N)$.
There is no restriction to orientation-preserving maps or to one
component in this notation.

\medskip\noindent\textit{The classical three-manifold comparison.}
Hamstrom allows compact manifolds with boundary and defines her
$\operatorname{Top}(N)$ to be the identity component
\cite[\S0.2, p.~208]{Hamstrom1974}.  Her PL convention is the
semisimplicial complex of joint PL maps on the total parameter
product, with homotopies preserving both the simplex and homotopy
parameters \cite[\S2.1, p.~222]{Hamstrom1974}.  Thus her finite
representatives and relative homotopies compute precisely the
positive homotopy groups of the simplicial group used here.  Her
Theorem~3.1.2, including its proof, gives the natural isomorphisms
\[
 \pi_j\mathcal H_{\mathrm{PL}}(N)
   \xrightarrow{\ \cong\ }
 \pi_j\mathcal H_{\mathrm{TOP}}(N),\qquad j\geq1,\quad\dim N=3,
\]
at the identity \cite[pp.~225--226]{Hamstrom1974}.

Here is the additional component argument.  The Moise--Bing
approximation theorem approximates every homeomorphism of a compact
PL three-manifold by PL homeomorphisms; its all-homeomorphism
formulation appears in
\cite[Chapter~I, \S1, Theorem~1.1, p.~7]{Laudenbach1974}.
Local contractibility of the TOP homeomorphism group
\cite[Corollary~1.1]{EdwardsKirby1971} implies that a sufficiently
close approximation is TOP-isotopic to the original homeomorphism.
Thus every TOP component contains a PL vertex.

For injectivity on components, suppose a TOP isotopy joins two PL
vertices.  Right multiplication by the inverse of its initial vertex
reduces to an isotopy $h_t$ starting at the identity.  Joint PL local
$0$-connectedness gives a $\delta>0$ such that every PL vertex in the
TOP identity component whose uniform distance from the identity is
less than $\delta$ admits a joint PL isotopy to the identity
\cite[Theorem~2.3.1, p.~224, with \S2.1]{Hamstrom1974}.
Subdivide the parameter interval so finely that consecutive values
of $h_t$ have uniform distance less than $\delta/3$.  Approximate
these values by PL maps $g_i$, preserving the two PL endpoints,
with errors less than $\delta/3$ and small enough that each $g_i$
is TOP-isotopic to its corresponding value of $h_t$.  The approximations
therefore belong to the TOP identity component and satisfy
\[
 d(g_{i+1}g_i^{-1},1)=d(g_{i+1},g_i)<\delta.
\]
The local joint PL isotopies of these differences, right-multiplied
by $g_i$, join consecutive approximations.  Their finite concatenation
is jointly PL.  Multiplying back by the initial vertex proves
injectivity on components.  Translation by a PL representative in
each component extends the positive-degree comparison there.  Hence
\begin{equation}\label{eq:three-manifold-full-automorphism-comparison}
 \mathcal H_{\mathrm{PL}}(N)\longrightarrow
 \mathcal H_{\mathrm{TOP}}(N)
 \quad\text{is a weak equivalence when }\dim N=3.
\end{equation}

\medskip\noindent\textit{A product decomposition for surface cylinders.}
First let $F$ be connected and closed, and put $N=F\times I$.
For either category let $\mathcal A(F)\subset\mathcal H(N)$ be the
subgroup preserving the bottom boundary component setwise.  A path
cannot change the permutation of the two boundary components.
Thus $\mathcal A(F)$ is a union of entire components, and
\eqref{eq:three-manifold-full-automorphism-comparison} restricts to
a weak equivalence on these subgroups.

Bottom restriction $r:\mathcal A(F)\to\mathcal H(F)$ has the
product section $s(a)(x,u)=(a(x),u)$.  There is a literal product
decomposition of the underlying simplicial sets,
\begin{equation}\label{eq:bottom-setwise-product-decomposition}
 \begin{aligned}
 \mathcal A(F)&\longrightarrow P(F)\times\mathcal H(F),&
 h&\longmapsto\bigl(hs(rh)^{-1},rh\bigr),\\
 P(F)\times\mathcal H(F)&\longrightarrow\mathcal A(F),&
 (p,a)&\longmapsto ps(a).
 \end{aligned}
\end{equation}
The formulas are inverse and preserve all parameters.  Joint PL
maps have joint PL inverses; in TOP, composition and inversion are
continuous on the compact-domain homeomorphism group.  Thus the TOP
formula is also a product decomposition of the underlying topological
spaces.  The multiplication law is a semidirect product, which does
not affect this underlying product decomposition.

Under \eqref{eq:bottom-setwise-product-decomposition}, the natural
comparison is the product of the concordance comparison and the
automorphism comparison for $F$.  A product of two maps between
nonempty spaces is a weak equivalence only if each factor is one:
bijectivity on the product of component sets implies bijectivity
on each set, and the product decomposition of the based homotopy
groups gives the positive-degree assertion.  Both factors are
nonempty because they contain the identity.  This proves the two
surface assertions for connected $F$.

\medskip\noindent\textit{The circle: positive PL homotopy groups.}
Scott's original definition uses joint PL homeomorphisms on
$\Delta^q\times N$ preserving its parameter projection, with the
usual faces and degeneracies
\cite[Chapter~III, Definition~0.1, p.~19]{Scott1968}.
For the annulus, his natural evaluation map at an interior point
induces isomorphisms
\[
 \pi_j\mathcal H_{\mathrm{PL}}(S^1\times I)
   \longrightarrow\pi_j(S^1\times I),\qquad j\geq1
\]
\cite[Chapter~IV, Theorem~3.1(b)(iii), p.~51,
and its annulus proof, p.~53]{Scott1968}.
Move the evaluation point along a path to the bottom circle.  Evaluating
along this path is a homotopy, with the corresponding basepoint
transport.  Bottom evaluation therefore induces the same
isomorphisms.  Evaluation for $\mathcal H_{\mathrm{PL}}(S^1)$
induces isomorphisms in positive degrees as well: the natural PL
rotation circle is a homotopy equivalence and evaluates with degree
one \cite[Theorem~2.2.4 and proof, pp.~223--224]{Hamstrom1974}.
Since the bottom circle includes into the annulus by a homotopy
equivalence, the evaluation square shows that bottom restriction
induces isomorphisms on every positive homotopy group.
Apply \eqref{eq:bottom-setwise-product-decomposition} to the annulus.
Its first factor $P_{\mathrm{PL}}(S^1)$ has zero positive homotopy
groups.

\medskip\noindent\textit{The circle: components.}
We give the component argument explicitly.  An annulus homeomorphism
fixed on the bottom is orientation preserving.  Indeed, naturality of
the boundary map from its relative fundamental class shows that its
orientation sign equals the degree on each boundary circle; the bottom
degree is one.  Its upper restriction therefore has degree one as well.  Every
orientation-preserving PL circle map has a joint PL isotopy to the
identity.  To construct one, lift it to an increasing PL map
$a:\mathbb R\to\mathbb R$ with $a(x+1)=a(x)+1$, subdivide one
period at all breakpoints, and triangulate each parameter rectangle.
Interpolate the identity on its lower edge and $a$ on its upper
edge affinely on the two triangles.  The spatial slope on either
triangle is respectively $1$ or the positive slope of $a$ on that
interval.  Consequently every horizontal slice is increasing.
The endpoint values respect translation by one, so the finite
joint PL construction descends to the circle.

Set $\rho(s,t)=\min\{s,t\}$.  If $b_s$ is a joint PL isotopy from
the identity to the inverse of the upper restriction, the family
\[
                    (x,t)\longmapsto(b_{\rho(s,t)}(x),t)
\]
is a joint PL concordance isotopy starting at the identity.  Composing
with its endpoint makes the given concordance fix both boundary
circles.  For $\alpha=\{0\}\times I$, the image of $\alpha$ then
has an integer winding $k$ relative to its endpoints in the universal
cover $\mathbb R\times I$.  Composing with
\[
 T_{-k}([x],t)=([x-kt],t)
\]
removes that winding.  This twist lies in the identity component
of the bottom-fixed concordance group, by the joint PL family
$([x],t)\mapsto([x-k\rho(s,t)],t)$.  The upper circle may move
during this family, as the definition permits.

The adjusted arc is homotopic to $\alpha$ relative to its endpoints.
Epstein's relative arc theorem gives a PL ambient isotopy fixed on
both boundary circles which carries the adjusted arc to $\alpha$
\cite[Theorem~3.1 and proof, pp.~89--90]{Epstein1966}.
To correct its parametrization, let $u:I\to I$ be the restriction
on the now setwise fixed arc, and choose a joint PL isotopy $b_s$
from the identity to $u^{-1}$ relative to the endpoints.  In a PL
strip $[-1,1]\times I$ about the arc, the family
\[
 (x,t)\longmapsto\bigl(x,b_{\min\{s,1-|x|\}}(t)\bigr)
\]
fixes the strip boundary and extends by the identity outside it.
At time one its restriction to the central arc is $u^{-1}$.
The formula and its inverse are jointly PL, so postcomposition
makes the arc pointwise fixed.
The adjusted annulus map is still orientation preserving, so it preserves
each of the two local sides of this pointwise fixed arc.  Consequently the
cut-open disk map fixes both copies of the arc pointwise, as well as the
rest of its boundary.
The PL Alexander isotopy of this disk map descends to the annulus
\cite[Theorem~5.2 and proof, pp.~95--96]{Epstein1966}.
This proves $\pi_0P_{\mathrm{PL}}(S^1)=0$ and hence weak
contractibility of the PL concordance group.

\medskip\noindent\textit{The TOP annulus.}
Yagasaki's compact-surface lemma implies that the identity component
of $\operatorname{Homeo}(S^1\times I\ \mathrm{rel}\ S^1\times\{0\})$
is contractible
\cite[Lemma~2.1(ii), p.~2, and proof, p.~3, author version]{Yagasaki2000}.
To check that this is the whole group, take a TOP concordance $h$.
Epstein's approximation isotopy changes it to a PL homeomorphism
$g$ \cite[Theorem~A4 and proof, p.~107]{Epstein1966}; one may choose
any fixed bottom point for the theorem's basepoint.  During that
isotopy $h_t$, the bottom boundary component remains the bottom
component.  If $a_t$ is its bottom restriction, the corrected isotopy
\[
                   (a_t^{-1}\times\operatorname{id}_I)h_t
\]
is fixed on the bottom, starts at $h$, and ends at a PL concordance.
The preceding PL component argument joins that endpoint to the
identity.  Thus the TOP concordance group is connected and equals
its contractible identity component.  The comparison for $S^1$
is now a weak equivalence between weakly contractible objects.

\medskip\noindent\textit{Points and finite disjoint unions.}
For a point, the TOP concordance group is the group of increasing
interval homeomorphisms fixing both endpoints.  Linear interpolation
with the identity contracts it continuously.  In the PL category,
use the joint parameter-preserving Alexander construction
\cite[Theorem~2.2.1 and proof, p.~223]{Hamstrom1974}.
More explicitly, represent a homotopy class by a joint PL map
$h:D\to D$, where $D=\Delta^q\times I$, which preserves the simplex
coordinate and is the identity on $\partial D$.
Let $a$ be the simplex barycenter, put $p=(a,\tfrac12,0)$ in $D\times I$,
and let $C=p*(D\times\{1\})$.  On this cone define
\[
 \widehat h\bigl((1-\lambda)p+\lambda(z,1)\bigr)
   =(1-\lambda)p+\lambda(h(z),1),
 \qquad 0\leq\lambda\leq1,
\]
and extend by the identity outside $C$.  The formula is the identity
on the lateral boundary because $h$ fixes $\partial D$.
Coning a finite triangulation for $h$ makes the extension jointly PL;
coning $h^{-1}$ gives its inverse.  The homotopy coordinate is $\lambda$.
For $z=(u,x)$, the simplex coordinate is
$(1-\lambda)a+\lambda u$ on both sides of the formula, so both
parameters are preserved.  The extension is the identity at time zero
and equals $h$ at time one.  This also applies when $q=0$ and proves
weak contractibility without bilinear interpolation.

A compact manifold has finitely many connected components, and a
bottom-fixed concordance preserves the cylinder of each one.
Both concordance objects are therefore finite products of their
component objects, even when the dimensions differ.  This proves
the stated comparison and weak contractibility for disconnected
bases and includes the empty product.  For a disconnected closed
surface, the same componentwise argument proves the automorphism
comparison on the component-preserving subgroup.  Other components
are indexed by permutations between homeomorphic surface components.
For a homeomorphism between two surface components, apply Moise--Bing
approximation to its product with $\operatorname{id}_I$.  A sufficiently
close PL approximation still carries the bottom boundary to the bottom
boundary; its restriction is a PL homeomorphism of the two surfaces.
Thus all realized permutations have PL representatives.  Translation
by those representatives gives the comparison on the full groups.
\end{proof}

\begin{corollary}\label{cor:effective-low-dimensional-top-models}
Given an ordered finite abstract simplex list promised to realize a
closed PL manifold whose components have dimension at most two,
there is an effective countable simplicial weak model for its real
TOP concordance space.  Its codes have decidable validity, equality
and degeneracy, and computable group and simplicial operations.
Uniformly in the input and the homotopy degree, a complete countable
group presentation can be enumerated; in positive degrees this may
be an integer relation matrix for the abelian homotopy group.
No terminating word problem or manifold-recognition algorithm is
asserted.
\end{corollary}

\begin{proof}
Compose the weak model in
Theorem~\ref{thm:effective-countable-concordance-models} with
Theorem~\ref{thm:low-dimensional-real-concordance-comparison}.
The codes and presentations remain the same; only their
identification with the real TOP homotopy groups is added.
\end{proof}


\subsection{A concordance obstruction over cyclic three-manifolds}

Let $B$ be a connected, closed, oriented real PL three-manifold.
Write $P_{\mathrm{TOP}}(B)$ for the compact-open space of
homeomorphisms of $B\times I$ that fix $B\times\{0\}$ pointwise;
the top face is free.  Write $P_{\mathrm{PL}}(B)$ for the corresponding
simplicial group of jointly PL families.  Thus a parameterized PL
family is required to be PL on its total space.

\begin{theorem}\label{thm:cyclic-three-concordance-obstruction}
Suppose that $\pi_1(B)$ is finite cyclic of even order.  Then there
is a homeomorphism
\[
 F:B\times I\longrightarrow B\times I,
 \qquad F|_{B\times\partial I}=\operatorname{id},
\]
whose component in $P_{\mathrm{TOP}}(B)$ contains no PL
concordance.  In particular,
\[
 \pi_0P_{\mathrm{PL}}(B)\longrightarrow
 \pi_0P_{\mathrm{TOP}}(B)
\]
is not surjective.  The same conclusion holds for the map from
definable concordance components, for every fixed o-minimal
expansion of $\mathbb R$.  The standard PL manifold
$\mathbb {RP}^{3}$ is an example.
\end{theorem}

We first describe the relative obstruction and the two classical
inputs used in the proof.  All cohomology groups without an
explicit coefficient group in this subsection have coefficients
$\mathbb Z/2$.

\paragraph{The formal PL difference class.}
If $X$ is a compact PL four-manifold and $f:X\to X$ fixes
$\partial X$ pointwise, its relative Casson--Sullivan class is
\begin{equation}\label{eq:relative-three-base-cs}
 \operatorname{cs}_{\mathrm{PL}}(f)
 =\sigma^{-1}\operatorname{ks}(M_f,\partial M_f)
 \in H^3(X,\partial X),\qquad
 M_f=(X\times I)\cup_f X.
\end{equation}
Here the two ends have the given PL tangent reductions, the
vertical boundary has their product reduction, and
$\operatorname{ks}$ is the obstruction to extending this
prescribed stable PL reduction across the five-dimensional
mapping cylinder.  The isomorphism
\[
 \sigma:H^3(X,\partial X)\xrightarrow{\cong}
             H^4(M_f,\partial M_f)
\]
is the relative suspension isomorphism.  Tangent microbundles
are identified using the canonical isomorphism induced by the
homeomorphism $f$.

This class measures the difference between the specified formal
stable tangent reductions.
Choose a smoothing compatible with the given PL structure,
whose existence is supplied by
\cite[Essay~I, p.~39]{KirbySiebenmann1977}.
The map $BO\to BPL$ identifies the degree-four
Kirby--Siebenmann obstruction, so \eqref{eq:relative-three-base-cs}
agrees with the smooth formal Casson--Sullivan class in
\cite[Remark~2.5 and Definition~2.11, pp.~2502--2504]{Galvin2026CassonSullivan}.
In particular it is invariant under an ordinary isotopy fixing
$\partial X$, and it vanishes for a PL homeomorphism fixing
$\partial X$.  The first assertion follows either by the
relative tangent-microbundle homotopy or by
\cite[Proposition~2.22, pp.~2507--2508]{Galvin2026CassonSullivan}.
For the second, the PL mapping cylinder extends the prescribed
boundary PL reduction over its entire interior, so its relative
obstruction vanishes.

\begin{lemma}\label{lem:compact-cyclic-relative-cs-realization}
Let $X$ be a compact, connected, oriented smooth four-manifold,
possibly with boundary, with finite cyclic fundamental group.
Every element of $H^3(X,\partial X)$ is
$\operatorname{cs}_{\mathrm{PL}}(f)$ for a homeomorphism
$f:X\to X$ which fixes $\partial X$ pointwise.
\end{lemma}

\begin{proof}
We spell out the boundary conditions in the normal-invariant
construction of
\cite[Section~4.2 and proof of Theorem~1.7, pp.~2524--2528]{Galvin2026CassonSullivan}.
Put $M=X\times I$ and retain the full boundary
\[
 \partial M=(X\times\partial I)\cup(\partial X\times I).
\]
Compatible collars, or a fixed rounding of corners, will be
understood.  The relative normal invariants have coordinates
\begin{equation}\label{eq:cyclic-relative-normal-coordinates}
 [M/\partial M,G/\mathrm{TOP}]
 \cong H^2(M,\partial M;\mathbb Z/2)
                  \oplus H^4(M,\partial M;\mathbb Z).
\end{equation}
Under the map to the relative stable smoothing obstruction,
the element $(0,y)$ has image $\operatorname{red}_2(y)$;
these are the coordinates and the characteristic-class formula
of Lemmas~4.8--4.9 in that source.  Only the zero first
coordinate is used here.

The coefficient exact sequence contains
\[
 H^4(M,\partial M;\mathbb Z)\xrightarrow{\operatorname{red}_2}
 H^4(M,\partial M;\mathbb Z/2)\longrightarrow
 H^5(M,\partial M;\mathbb Z)\xrightarrow{\,2\,}
 H^5(M,\partial M;\mathbb Z).
\]
The last integral group is $\mathbb Z$, by orientation and
Poincar\'e--Lefschetz duality.  Consequently reduction is
surjective.  Given $\eta\in H^3(X,\partial X)$, choose $y$
whose reduction is its suspension.

For a finite cyclic group $\pi$, Bak's calculation gives
$L_5^s(\mathbb Z[\pi])=0$, as recalled in
\cite[Section~4.4, p.~2528]{Galvin2026CassonSullivan}.
The surgery obstruction of $(0,y)$ therefore vanishes.
Exactness of the simple relative TOP surgery sequence
\cite[Theorem~11.3A and the simple-structure convention, p.~200]{FreedmanQuinn1990}
provides a simple homotopy equivalence of triads
\[
 \phi:(N;\partial_+N,\partial_0N,\partial_-N)
 \longrightarrow
 (X\times I;X\times\{0\},\partial X\times I,X\times\{1\}),
\]
whose restriction to the full boundary is a homeomorphism.
The smooth boundary markings are those induced by $\phi$.

Finite groups are good in the sense of Freedman--Quinn
\cite[p.~99]{FreedmanQuinn1990}.  The compact relative
five-dimensional $s$-cobordism theorem
\cite[Theorem~7.1A, p.~101]{FreedmanQuinn1990}, in the explicitly
marked form of
\cite[Proposition~4.10, pp.~2526--2527]{Galvin2026CassonSullivan},
gives a product homeomorphism
\[
 p:N\xrightarrow{\cong}\partial_+N\times I
\]
which is the identity on $\partial_+N$ and on its full vertical
boundary.  Write $p_-$ for its top restriction and $\phi_\pm$
for the end markings, regarded as maps to $X$.  Then
\[
 f=\phi_+\,p_-\,\phi_-^{-1}:X\longrightarrow X
\]
is the identity on $\partial X$, because the product structure
and the two end markings agree on the marked vertical boundary.
Moreover, $N$ is the mapping cylinder $M_f$ with exactly its
prescribed boundary reductions.  By the relative obstruction
identification in
\cite[Lemma~4.12, p.~2527]{Galvin2026CassonSullivan},
\[
 \operatorname{cs}_{\mathrm{PL}}(f)
   =\sigma^{-1}\operatorname{red}_2(y)=\eta.
\]
This proves the assertion with every boundary face fixed.
\end{proof}

\paragraph{Replacing an endpoint path by a jointly PL path.}
For a closed real PL three-manifold the automorphism comparison
\eqref{eq:three-manifold-full-automorphism-comparison} is a weak
equivalence.  Its parameter convention is the total-space PL
convention of
\cite[Section~2.1, p.~222; Theorem~3.1.2, pp.~225--226]{Hamstrom1974}.
We need only injectivity on components and surjectivity on
the fundamental group.

It follows that if $a_t$, $0\leq t\leq1$, is an ordinary path
of homeomorphisms with $a_0=\operatorname{id}$ and $a_1$ PL,
then there is a jointly PL path $b_t$ with these endpoints
which is ordinary homotopic to $a_t$ relative to its endpoints.
Indeed, component injectivity first supplies a jointly PL path
$b^0$ with the same endpoints.  The loop $a*(b^0)^{-1}$ is
ordinary homotopic to a jointly PL loop $\ell$, by surjectivity
on $\pi_1$.  Concatenating $\ell$ with $b^0$ gives the required
$b$.  Concatenations are made with piecewise affine
reparameterizations, which preserve the joint PL condition.

\begin{proof}[Proof of Theorem~\ref{thm:cyclic-three-concordance-obstruction}]
Give $B$ a smoothing compatible with its PL structure and put
$X=B\times I$.  Mod-two duality and the abelianization of its
fundamental group give
\[
 H^3(X,\partial X)\cong H^2(B)
 \cong H_1(B;\mathbb Z/2)\cong\mathbb Z/2.
\]
Apply Lemma~\ref{lem:compact-cyclic-relative-cs-realization} to
the nonzero class.  It gives a homeomorphism $F$ fixing both
ends with $\operatorname{cs}_{\mathrm{PL}}(F)\ne0$.

Suppose that an ordinary path $H_t$ in $P_{\mathrm{TOP}}(B)$
joins $F$ to a PL concordance $G$.  Its top restrictions
$a_t=H_t|_{B\times\{1\}}$ run from the identity to the PL
map $G|_{B\times\{1\}}$.  Define
\[
 T_t(x,s)=(a_{ts}(x),s),\qquad Q_t=T_t^{-1}H_t.
\]
These are ordinary TOP families; compactness of $B$ ensures
continuity of their inverse families.  Each $Q_t$ fixes both
ends.  Thus $Q_t$ is a relative-boundary isotopy from $F$ to
$T_1^{-1}G$.

Replace $a$ by the jointly PL path $b$ constructed above, and
write $T^b(x,s)=(b_s(x),s)$.  An endpoint-fixed homotopy from
$a$ to $b$ gives an isotopy of their traces with fixed boundary
restrictions.  After inversion and composition with $G$, it
gives an isotopy relative to the entire boundary from
$T_1^{-1}G$ to $(T^b)^{-1}G$.  The latter is a PL
homeomorphism fixing both ends.  Relative isotopy invariance
and PL vanishing in \eqref{eq:relative-three-base-cs} imply
$\operatorname{cs}_{\mathrm{PL}}(F)=0$, a contradiction.

For a fixed o-minimal expansion of $\mathbb R$, the relative
definable--PL component comparison of Theorem~\ref{result-4}
replaces every definable concordance by a PL one through a
definable isotopy with the bottom fixed pointwise.  Its
underlying ordinary isotopy shows that the definable and PL
images in $\pi_0P_{\mathrm{TOP}}(B)$ coincide.  The same
component obstruction therefore applies.
\end{proof}


\begin{corollary}\label{cor:cyclic-three-absolute-obstruction}
The homeomorphism $F$ in
Theorem~\ref{thm:cyclic-three-concordance-obstruction} can be chosen so
that its component in the absolute ordinary group
$\operatorname{Homeo}(B\times I)$ contains no PL homeomorphism and
no definable homeomorphism, for every fixed o-minimal expansion of
$\mathbb R$.
\end{corollary}

\begin{proof}
Set $X=B\times I$ and fix a compatible smoothing. The homeomorphism
constructed above has nonzero formal Casson--Sullivan class, so its
isotopy class relative to $\partial X$ has no diffeomorphism
representative. It preserves orientation: it fixes each boundary
point and hence has positive orientation sign in a boundary chart.
Galvin--Ladu's boundary-forgetting lemma
\cite[Lemma~5.3, pp.~17--18]{GalvinLadu2025Boundary}
implies that its absolute ordinary isotopy class has no
diffeomorphism representative either. The lemma applies to compact
oriented smooth four-manifolds with boundary, including the two
boundary components here. Its proof uses the boundary restriction
fibrations, injectivity of the three-dimensional boundary comparison
on components, and surjectivity on fundamental groups: the first
makes an absolute smooth representative relative to the boundary,
and the second removes the remaining smooth boundary twist.

Let $H$ be a PL self-homeomorphism of $X$, and write $\sigma$ for
the fixed compatible smoothing. Both $\sigma$ and $H^*\sigma$
are compatible smoothings of the same PL manifold. They are isotopic
as smoothings, by
\cite[Theorem~3.5(2), pp.~17--18; boundary convention p.~8]
{FriedlNagelOrsonPowell2025}.
Here isotopy of smoothings means that there is a diffeomorphism
\[
 u:(X,\sigma)\longrightarrow(X,H^*\sigma)
\]
whose underlying homeomorphism is isotopic to the identity
\cite[Preface, pp.~v--vi]{HirschMazur1974}.
By definition of pullback smoothing, $Hu$ is a diffeomorphism of
$(X,\sigma)$, and an isotopy from the identity to $u$ gives an
ordinary isotopy from $H$ to $Hu$. Thus a PL representative in the
absolute component of $F$ would contradict the preceding conclusion.

Finally, Shiota's given-map isotopy theorem
\cite[Complement to Theorem~2.1, pp.~17--18]{Shiota2010}
isotopes every definable homeomorphism of the compact polyhedron
$X$ to a PL homeomorphism. Over the real field this is also an
ordinary isotopy. Hence the same absolute component contains no
definable representative.
\end{proof}

\paragraph{The finite projective-space example.}
One explicit PL model of $\mathbb {RP}^{3}$ is the quotient
of the barycentric subdivision of the boundary of the
four-dimensional cross-polytope by the antipodal action.
Equivalently, take the order complex of the nonzero partial
sign vectors in $\{-1,0,1\}^4$, modulo simultaneous sign
reversal; order two orbits by compatible strict support
inclusion.  This is a PL quotient of $S^3$ by a free involution,
so it is the standard oriented PL projective three-space and
its fundamental group is $\mathbb Z/2$.

Its face vector is $(40,232,384,192)$.  More precisely, if
$S(r,k)$ denotes a Stirling number of the second kind, then
\[
 f_{k-1}=\sum_{r=k}^{4}\binom4r2^{r-1}k!S(r,k),
 \qquad 1\leq k\leq4.
\]
Choose the top support, its signs up to global reversal, and
the ordered nonempty blocks added along a chain to obtain
this formula.  The usual one-cell-per-dimension CW model has
integral cellular complex
\[
 0\longrightarrow\mathbb Z\xrightarrow{\,0\,}\mathbb Z
 \xrightarrow{\,2\,}\mathbb Z\xrightarrow{\,0\,}\mathbb Z
 \longrightarrow0.
\]
Thus $H_1(\mathbb {RP}^3;\mathbb Z)=\mathbb Z/2$ and
$H^2(\mathbb {RP}^3;\mathbb Z/2)=\mathbb Z/2$, as required.
The finite triangulation verifies the concrete base; the
existence of the obstructed concordance is the geometric
argument above.

\paragraph{Relation to the literature.}
For finite-cyclic elliptic three-manifolds, Galvin already
obtains a related failure of surjectivity on $\pi_1$ for
block diffeomorphism and block homeomorphism spaces
\cite[Theorem~9.0.1, Corollary~9.0.2, and Section~9.3,
pp.~109--117]{GalvinThesis2024}.
The finite cyclic case here uses his relative
Casson--Sullivan realization and makes explicit the passage
to ordinary free-top concordance components, using the
jointly PL replacement of the top path.


\subsection{A natural obstruction on three-dimensional concordance components}

The same relative class defines a homomorphism on all ordinary
concordance components, without an orientability assumption.
We use the real, joint-family, and free-top conventions of
Theorem~\ref{thm:cyclic-three-concordance-obstruction}.

\begin{theorem}\label{thm:natural-three-concordance-obstruction}
For every connected closed real PL three-manifold $B$ there
is a natural group homomorphism
\begin{equation}\label{eq:natural-three-concordance-class}
 c_B:\pi_0P_{\mathrm{TOP}}(B)\longrightarrow
                         H^2(B;\mathbb Z/2)
\end{equation}
which vanishes on the PL image and on the definable image
for every fixed o-minimal expansion of $\mathbb R$.
For a PL homeomorphism $f:B\to B'$ its naturality is
\begin{equation}\label{eq:three-concordance-typed-naturality}
 c_{B'}\bigl([(f\times1)H(f^{-1}\times1)]\bigr)
                    =(f^{-1})^*c_B([H]).
\end{equation}
If $B$ is oriented and $\pi_1(B)$ is finite cyclic, then
$c_B$ is surjective.  In particular, for even cyclic order,
\[
 \pi_0P_{\mathrm{TOP}}(B)/\ker c_B\cong\mathbb Z/2,
 \qquad
 \operatorname{im}\pi_0P_{\mathrm{def}}(B)\subseteq\ker c_B.
\]
\end{theorem}

\begin{proof}
Put $X=B\times I$, and abbreviate mod-two cohomology by
$H^*(-)$.  Let $R_B$ be the subgroup of TOP concordances
whose top map is PL.  For $h\in R_B$, both boundary maps of
$X$ are PL, so the mapping cylinder of $h$ has prescribed
PL tangent reductions on its entire boundary.  Define
\begin{equation}\label{eq:three-boundary-pl-class}
 k_B(h)=\sigma^{-1}\operatorname{ks}(M_h,\partial M_h)
                     \in H^3(X,\partial X).
\end{equation}
This extends \eqref{eq:relative-three-base-cs}; the side of
$M_h$ is now the mapping cylinder of the PL boundary map.

\emph{Common product collars.}
The boundary identification can be made literal on a product
collar.  If $a$ is the top map of $h$ and $0<r<1/2$, let
$e_r(x,s)=(x,r+(1-2r)s)$ and define
\[
 h_r(x,s)=
 \begin{cases}
  (x,s),&0\leq s\leq r,\\
  e_rhe_r^{-1}(x,s),&r\leq s\leq1-r,\\
  (a(x),s),&1-r\leq s\leq1.
 \end{cases}
\]
The pieces agree at both interfaces, and the inverse is the
same construction applied to $h^{-1}$.  Setting $h_0=h$
gives an ordinary isotopy with fixed boundary maps.  In a
product metric, uniform continuity bounds the displacement
from $h$ by $r+\omega_h(r)$; the same estimate holds for
the inverse.  A compact parameter family has a common
modulus of continuity.  Thus the construction works for
families as well.  At a fixed width it satisfies
$(uv)_r=u_rv_r$: inside it is conjugation by $e_r$, and
outside it is multiplication of the boundary maps.  A fixed
width also preserves joint PL families.  These collars let
us compare or compose any finite collection with identical
normal-line conventions.

\emph{The relative difference rule.}
The standard PL structure of $X$ supplies a section of the
stable tangent-reduction fibration.  Its fiber is
\[
 \mathrm{TOP}/\mathrm{PL}\simeq K(\mathbb Z/2,3)
\]
by \cite[Essay~V, Theorem~5.5(I), p.~251]{KirbySiebenmann1977}.
This coefficient system is untwisted, even when $B$ is
nonorientable.  Indeed the chosen section points each fiber;
the pointed self-equivalence space of $K(\mathbb Z/2,3)$
has one component, since $\operatorname{Aut}(\mathbb Z/2)=1$,
and its positive homotopy groups are
\[
 \widetilde H^{3-i}(K(\mathbb Z/2,3);\mathbb Z/2)=0
 \qquad(i>0).
\]
Consequently relative differences of its sections form the
ordinary group $H^3(X,\partial X)$.

Pull back the standard reduction by $h$ using the canonical
tangent-microbundle identification, with its prescribed PL
boundary identification.  Its difference from the original
reduction is \eqref{eq:three-boundary-pl-class}.  Inserting
an intermediate reduction in the comparison for a composite
gives
\begin{equation}\label{eq:three-relative-cocycle}
 k_B(uv)=k_B(v)+v^*k_B(u),\qquad u,v\in R_B.
\end{equation}
The boundary identifications compose because they are PL.
This is the relative PL difference-of-reductions argument
underlying the classical Casson--Sullivan cocycle; compare
\cite[Proposition~2.23 and its full proof, pp.~2508--2509]{Galvin2026CassonSullivan}.
It also proves invariance under an ordinary isotopy with
fixed PL boundary maps, and under one whose boundary is a
jointly PL family.  The common collars above provide the
literal relative boundary convention for such a family.
If $h$ is PL, its mapping cylinder is PL, so $k_B(h)=0$.

Every bottom-fixed concordance acts trivially on
$H^3(X,\partial X)$.  To see this without orientability,
its top map $a$ is homotopic to the identity by
$(x,t)\mapsto\operatorname{pr}_B h(x,t)$.  The restriction
maps $H^j(X)\to H^j(\partial X)$ are diagonal for $j=2,3$,
and the latter diagonal is injective.  The exact sequence
of the pair therefore gives
\begin{equation}\label{eq:three-relative-trivial-action}
 H^3(X,\partial X)
 \cong\operatorname{coker}\bigl(H^2(B)
                 \xrightarrow{\Delta}H^2(B)\oplus H^2(B)\bigr).
\end{equation}
The boundary action is the identity on the bottom and $a^*=1$
on the top.  It is therefore the identity on this cokernel.
Thus \eqref{eq:three-relative-cocycle} is additive on $R_B$.

\emph{Endpoint traces have zero class.}
If $\delta_s$ is an ordinary path from the identity to a PL
map, set $T_\delta(x,s)=(\delta_s(x),s)$.  The
three-manifold automorphism comparison
\eqref{eq:three-manifold-full-automorphism-comparison}
replaces $\delta$ by a jointly PL path with the same
endpoints and the same ordinary relative path class, as in
the preceding proof.  Its trace is PL, and the intervening
trace isotopy fixes both boundary maps.  Hence
\begin{equation}\label{eq:three-trace-class-zero}
 k_B(T_\delta)=0.
\end{equation}

\emph{Normalization and independence.}
Let $h\in P_{\mathrm{TOP}}(B)$ have top map $f$.  Surjectivity
on components of the automorphism comparison supplies a
PL map $a$ and an ordinary path $\gamma_s:f\to a$.  Put
\begin{equation}\label{eq:three-endpoint-normalization}
 E_\gamma(x,s)=(\gamma_s f^{-1}(x),s),
 \qquad h^\gamma=E_\gamma h\in R_B.
\end{equation}
Replacing $\gamma_s$ by $\gamma_{ts}$ gives a path from the
identity to $E_\gamma$ in the bottom-fixed group.  Thus
$h^\gamma$ belongs to the same ordinary free-top component
as $h$.  If $\gamma':f\to a'$ is another choice, then
\[
 h^{\gamma'}(h^\gamma)^{-1}=E_{\gamma'}E_\gamma^{-1}
            =T_\delta,\qquad \delta_s=\gamma'_s\gamma_s^{-1}.
\]
This path starts at the identity and ends at the PL map
$a'a^{-1}$.  Additivity and \eqref{eq:three-trace-class-zero}
give $k_B(h^{\gamma'})=k_B(h^\gamma)$.

\emph{Ordinary component invariance.}
Let $H_t$ be a path with $H_0=h'\in R_B$, $H_1=k$, and
top path $f_t$ starting at the PL map $a=f_0$.  Define
\[
 T_t(x,s)=(f_{ts}a^{-1}(x),s),\qquad Q_t=T_t^{-1}H_t.
\]
The family $Q_t$ fixes the bottom and has the constant PL
top map $a$.  Its relative class is therefore constant.
At $t=1$ it is the normalization of $k$ obtained from
\[
 \gamma_s=a f_s^{-1}f_1:f_1\longrightarrow a,
 \qquad E_\gamma(x,s)=(a f_s^{-1}(x),s)=T_1^{-1}(x,s).
\]
Hence the normalized value is constant along $H_t$.  For a
path whose initial map is not normalized, first prepend
the reverse of its normalization path.  This proves that
the value is well-defined on all ordinary components.

Compose it with the interval isomorphism
$H^3(X,\partial X)\cong H^2(B)$ to define $c_B$.
For arbitrary $h,k$, multiply their two normalization
paths.  The product of their normalized representatives
lies in $R_B$ and represents the component of $hk$.
Additivity on $R_B$ yields
\[
 c_B([hk])=c_B([h])+c_B([k]).
\]
For a PL concordance choose a constant normalization path,
giving zero.  Theorem~\ref{result-4} makes every definable
concordance ordinarily isotopic to a PL one, so its class
also vanishes.

Conjugation by $f\times1$, for a PL homeomorphism of bases
$f:B\to B'$, transports the product reductions, boundary
markings, and mapping cylinders.  Pullback of the relative
obstruction through the inverse change of coordinates is
$(f^{-1}\times1)^*$.  Naturality of interval suspension
then gives exactly \eqref{eq:three-concordance-typed-naturality}.

Finally, if $B$ is oriented with finite cyclic fundamental
group, Lemma~\ref{lem:compact-cyclic-relative-cs-realization}
realizes every class by a map fixing both ends.  Such a map
is already normalized, so $c_B$ is surjective.  For even
cyclic order, $H^2(B;\mathbb Z/2)=\mathbb Z/2$ by mod-two
duality and abelianization.  The first isomorphism theorem
gives the asserted quotient by $\ker c_B$.
\end{proof}

The cocycle and its vanishing on PL maps are classical
formal Casson--Sullivan properties.  The construction above
records their descent to the specified ordinary free-top
component group. Determining which zero-class components contain
a PL representative remains a separate question.


\subsection{Realization over three-manifolds with finite fundamental group}

The finite cyclic hypothesis can be removed from the surjectivity
assertion in Theorem~\ref{thm:natural-three-concordance-obstruction}.
The relevant surgery input is the relative obstruction formula of
Hambleton, Milgram, Taylor and Williams; the construction of the
homeomorphism is Galvin's relative realization construction.

\begin{theorem}\label{thm:finite-three-concordance-realization}
Let $B$ be a closed, connected, oriented real PL three-manifold
with finite fundamental group.  Then
\[
 c_B:\pi_0P_{\mathrm{TOP}}(B)\longrightarrow
                         H^2(B;\mathbb Z/2)
\]
is surjective.  Every class on the right has a representative
$F:B\times I\to B\times I$ fixing both ends pointwise.
Here $P_{\mathrm{TOP}}(B)$ has only the bottom face fixed,
and its components are ordinary topological path components.
The images of the jointly PL component map and of the definable
component map for every fixed o-minimal expansion of
$\mathbb R$ are contained in $\ker c_B$.
\end{theorem}

We use the formal relative PL difference class
$\operatorname{cs}_{\mathrm{PL}}(f)\in
H^3(X,\partial X;\mathbb Z/2)$ defined in
\eqref{eq:relative-three-base-cs}.  It records the prescribed
stable PL reductions on the entire boundary of the
five-dimensional mapping cylinder.

\begin{lemma}\label{lem:finite-sylow-relative-cs-realization}
Let $X$ be a compact, connected, oriented smooth four-manifold,
possibly with boundary.  Suppose that $\Gamma=\pi_1(X)$ is finite
and that a Sylow two-subgroup $\rho$ satisfies
\[
                    SK_1(\mathbb Z[\rho])=0.
\]
Then each $\eta\in H^3(X,\partial X;\mathbb Z/2)$ is
$\operatorname{cs}_{\mathrm{PL}}(f)$ for a homeomorphism
$f:X\to X$ fixing the entire boundary pointwise.
\end{lemma}

\begin{proof}
Set $M=X\times I$ and retain its full boundary, including
$\partial X\times I$.  Use fixed product collars, or round
the corners compatibly.  The relative normal-invariant
coordinates of
\cite[Lemmas~4.8--4.9, p.~2526]{Galvin2026CassonSullivan}
are
\begin{equation}\label{eq:finite-relative-normal-coordinates}
 [M/\partial M,G/\mathrm{TOP}]
 \cong H^2(M,\partial M;\mathbb Z/2)
                \oplus H^4(M,\partial M;\mathbb Z).
\end{equation}
The stable PL obstruction of $(0,y)$ is
$\operatorname{red}_2(y)$.  The coefficient exact sequence and
orientation give
\[
 H^4(M,\partial M;\mathbb Z)
 \xrightarrow{\operatorname{red}_2}
 H^4(M,\partial M;\mathbb Z/2)\longrightarrow
 H^5(M,\partial M;\mathbb Z)
 \xrightarrow{\,2\,}H^5(M,\partial M;\mathbb Z).
\]
Since $H^5(M,\partial M;\mathbb Z)\cong\mathbb Z$,
reduction is onto.  Choose $y$ reducing to the interval
suspension $\sigma\eta$.

We verify the vanishing of the surgery obstruction of this
particular normal invariant.  In the notation of
\cite[Theorem~A(b) and Definition~(0.6), p.~351]
{HambletonMilgramTaylorWilliams1988}, the permitted decoration
$U\subset\operatorname{Wh}(\Gamma)$ must contain
\[
 Y=\operatorname{im}\bigl(SK_1(\mathbb Z[\rho])
                       \longrightarrow\operatorname{Wh}(\Gamma)\bigr).
\]
Our hypothesis permits $U=0$, the simple decoration.
Apply that theorem to the oriented Poincar\'e pair
$(M,\partial M)$, taking the identity as the reference normal
map.  For a normal invariant $\nu$ it gives
\begin{equation}\label{eq:finite-five-dimensional-surgery-formula}
 \theta^s(\nu)=\kappa_3^s\bigl(c_*
                         \operatorname{ARF}_3(\nu)\bigr),
 \qquad c:M\longrightarrow B\Gamma.
\end{equation}
The class $\operatorname{ARF}_3$ is the degree-three homology
component of
\[
       \bigl(V_M^2\smile\nu^*k\bigr)\cap[M,\partial M],
       \qquad k=\{k_{4i+2}\}_{i\geq0}.
\]
Since $\dim M=5$, this component uses only $k_2$ and the
degree-zero Wu factor.  Its pullback vanishes for a normal
invariant with zero degree-two coordinate in
\eqref{eq:finite-relative-normal-coordinates}.  Therefore
\begin{equation}\label{eq:finite-zero-arf-realization}
                  \theta^s(0,y)=0.
\end{equation}
In particular, the integral degree-four coordinate does not
contribute to this obstruction.  This argument requires no
vanishing assertion for the whole group $L_5^s(\mathbb Z[\Gamma])$.

The simple relative TOP surgery sequence
\cite[Theorem~11.3A, p.~200]{FreedmanQuinn1990}
now realizes $(0,y)$ by a simple homotopy equivalence of triads
\[
 \phi:(N;\partial_+N,\partial_0N,\partial_-N)
   \longrightarrow
 (X\times I;X\times\{0\},\partial X\times I,X\times\{1\})
\]
whose restriction to the full boundary is the prescribed
homeomorphism.  The two end inclusions are simple homotopy
equivalences: this follows by comparing them under $\phi$
with the product inclusions, using additivity of Whitehead
torsion.  Consequently $N$ is a simple cobordism with its
labeled side product.

Finite groups are good
\cite[p.~99]{FreedmanQuinn1990}.  The compact relative
five-dimensional $s$-cobordism theorem
\cite[Theorem~7.1A, p.~101]{FreedmanQuinn1990}, in the marked
form stated in
\cite[Proposition~4.10, pp.~2526--2527]{Galvin2026CassonSullivan},
gives a product homeomorphism
\[
            p:N\xrightarrow{\cong}\partial_+N\times I
\]
which fixes the bottom and the entire labeled vertical
boundary.  Write $p_-$ for its restriction to the other end
and $\phi_\pm$ for the end markings, with the interval
coordinate suppressed.  Then
\[
                 f=\phi_+p_-\phi_-^{-1}:X\longrightarrow X
\]
All three boundary restrictions agree with the fixed side marking,
so $f|_{\partial X}=\mathrm{id}$.  To track the specified class,
put $\chi=(\phi_+\times\mathrm{id}_I)p:N\to X\times I$.
The maps $\phi$ and $\chi$ agree on the lower end and on the
entire vertical boundary.  Consequently
\[
 H=\operatorname{pr}_X\phi\chi^{-1}:X\times I\longrightarrow X
\]
is a homotopy relative to $\partial X$ from $\mathrm{id}_X$
to $f^{-1}$.  In particular, $f$ acts trivially on
$H^3(X,\partial X;\mathbb Z/2)$.

The two markings $\phi$ and $\chi$ induce the same map on
$H_1(N;\mathbb Z/2)$, because they agree on an end whose
inclusion is a homotopy equivalence.  Both carry the relative
mod-two fundamental class to that of $X\times I$.
Naturality of Poincar\'e--Lefschetz duality therefore identifies
their pullbacks on $H^4(X\times I,\partial(X\times I);\mathbb Z/2)$.
Thus the normal-invariant marking and the product marking assign
the same class $\operatorname{red}_2(y)$ to the relative
obstruction.

In the $\chi$ coordinates the lower reduction is the standard
one, the upper reduction is its pullback by $f^{-1}$, and the
side reduction is the prescribed product reduction.  The marked
mapping-cylinder identity
\cite[Lemma~4.12, p.~2527]{Galvin2026CassonSullivan} gives
\[
 \operatorname{cs}_{\mathrm{PL}}(f^{-1})
       =\sigma^{-1}\operatorname{red}_2(y)=\eta.
\]
The inverse rule for formal difference classes and the trivial
cohomology action give
\[
 \operatorname{cs}_{\mathrm{PL}}(f)
   =f^*\operatorname{cs}_{\mathrm{PL}}(f^{-1})=\eta,
\]
as required.
\end{proof}

\begin{proof}[Proof of Theorem~\ref{thm:finite-three-concordance-realization}]
We first verify the Sylow hypothesis of the lemma.  A finite
triangulation of $B$ lifts to a finite free cellular
$\mathbb Z[\Gamma]$-complex $C_*$ on its universal cover,
where $\Gamma=\pi_1(B)$.  The cover $\widetilde B$ is compact,
simply connected and oriented.  Simple connectivity and
three-dimensional Poincar\'e duality give
\[
 H_0(\widetilde B;\mathbb Z)
   =H_3(\widetilde B;\mathbb Z)=\mathbb Z,
 \qquad
 H_1(\widetilde B;\mathbb Z)
   =H_2(\widetilde B;\mathbb Z)=0.
\]
The deck group acts trivially on $H_3$, since the orientation
is pulled back from $B$.  Hence there is an exact sequence
of $\mathbb Z[\Gamma]$-modules
\begin{equation}\label{eq:finite-three-periodic-extension}
 0\longrightarrow\mathbb Z\xrightarrow{\iota}C_3
 \longrightarrow C_2\longrightarrow C_1
 \longrightarrow C_0\xrightarrow{\epsilon}\mathbb Z
 \longrightarrow0,
\end{equation}
with trivial modules at both ends.  Splice copies of this
extension, joining consecutive blocks by
$\iota\epsilon:C_0\to C_3$.  Exactness at a join follows
from
\[
 \ker(\iota\epsilon)=\ker\epsilon=\operatorname{im}d_1,
 \qquad
 \operatorname{im}(\iota\epsilon)
     =\operatorname{im}\iota=\ker d_3.
\]
This constructs a four-periodic free resolution and proves
that $\Gamma$ has periodic cohomology.

For a finite group with periodic cohomology, its Sylow
two-subgroups have trivial integral $SK_1$; this is the
specific arithmetic conclusion recorded in the last
introductory paragraph of
\cite[p.~352]{HambletonMilgramTaylorWilliams1988}.
The corresponding computations for cyclic and generalized
quaternion two-groups are also given in
\cite[Introduction, Example~2, p.~14; Example~14.4,
p.~333]{Oliver1988Whitehead}.
We use this assertion for the Sylow subgroup alone.

Choose a smoothing of $B$ compatible with its PL structure
\cite[Essay~I, p.~39]{KirbySiebenmann1977}
and set $X=B\times I$ with the product smoothing.
Given $a\in H^2(B;\mathbb Z/2)$, its interval suspension
lies in $H^3(X,\partial X;\mathbb Z/2)$.
Lemma~\ref{lem:finite-sylow-relative-cs-realization}
realizes this class by a homeomorphism $F:X\to X$
fixing both ends.  The normalization of $c_B$ in
Theorem~\ref{thm:natural-three-concordance-obstruction}
then gives $c_B([F])=a$.
The same theorem gives both kernel containments, with
ordinary top-free components and the stipulated joint PL
and real definable parameter conventions.
\end{proof}

The obstruction homomorphism is therefore onto. Determining
its kernel and the two comparison images remains a further question.


\subsection{Supported classes and connected-sum summands}

All homology and cohomology groups in this subsection have
coefficients $\mathbb Z/2$.  The ordinary concordance space
$P_{\mathrm{TOP}}(B)$ has its bottom face fixed and its top
face free, as in
Theorem~\ref{thm:natural-three-concordance-obstruction}.

\begin{theorem}\label{thm:supported-three-concordance-realization}
Let $B$ be a closed, connected real PL three-manifold, possibly
nonorientable.  Let $W_1,\ldots,W_s$ be pairwise disjoint compact
codimension-zero PL submanifolds of $B$.  Suppose that each
$W_i$ is PL homeomorphic to
$C_i\setminus\operatorname{int}D^3$, where $C_i$ is a closed,
connected, oriented PL three-manifold with finite fundamental
group.  For the inclusions $j_i:W_i\hookrightarrow B$, one has
\begin{equation}\label{eq:supported-three-image}
 \begin{split}
 \operatorname{im}c_B
   &\supset\sum_i(j_i)_!H^2(W_i,\partial W_i)\\
   &=\operatorname{PD}_B^{-1}
       \left(\sum_i(j_i)_*H_1(W_i)\right).
 \end{split}
\end{equation}
Here $(j_i)_!$ is extension by zero and
$\operatorname{PD}_B:H^2(B)\to H_1(B)$ is mod-two
Poincar\'e duality.  For each $a\in H^2(W_i,\partial W_i)$,
the class $(j_i)_!a$ is realized by a TOP concordance equal
to the identity outside $W_i\times I$ and on a neighborhood
of the entire boundary of $W_i\times I$.
\end{theorem}

\begin{proof}
First consider one punctured manifold
$W=C\setminus\operatorname{int}D^3$.  Van Kampen's theorem
gives $\pi_1(W)\cong\pi_1(C)$, since the gluing sphere is
simply connected.  The closed manifold $C$ supplies the
four-periodic free resolution
\eqref{eq:finite-three-periodic-extension}.  Hence the
Sylow two-subgroups of this finite group satisfy the
$SK_1$ hypothesis of
Lemma~\ref{lem:finite-sylow-relative-cs-realization}.
This use of the closed manifold is essential: the universal
cover of the punctured manifold need not be a homology sphere.

Regard $X=W\times I$ as a compact PL four-manifold with
boundary
\[
 \partial X=(\partial W\times I)\cup(W\times\partial I).
\]
Choose a smoothing compatible with this product PL structure
\cite[Essay~I, p.~39]{KirbySiebenmann1977}.  Product collars
specify its spatial and time faces.  The relative suspension
isomorphism identifies
\[
 H^2(W,\partial W)\cong H^3(X,\partial X).
\]
The relative realization lemma gives a homeomorphism
$F:X\to X$ fixing all of $\partial X$, with
$\operatorname{cs}_{\mathrm{PL}}(F)=\sigma a$ for the
chosen $a\in H^2(W,\partial W)$.  In particular both time
ends and the full spatial side are fixed.

\paragraph{Making the support interior.}
Choose a PL collar
$c:\partial X\times[0,\delta]\to X$ of the full boundary;
the existence theorem applies to this PL manifold with
compact boundary
\cite[Theorem~2.25 and Corollary~2.26, pp.~24--26]
{RourkeSanderson1972}.  The corners of the product cause no
exception: their union is the boundary of the underlying
PL manifold.  For $0\leq u\leq\delta/2$, set
\begin{equation}\label{eq:supported-full-collar-compression}
 \begin{split}
 \sigma_u(t)&=\min\{t+u,(t+\delta)/2\},\qquad 0\leq t\leq\delta,\\
 \sigma_u^{-1}(v)&=\max\{v-u,2v-\delta\},\qquad u\leq v\leq\delta.
 \end{split}
\end{equation}
The two affine branches meet at $t=\delta-2u$ and have
positive slopes.  Thus $\sigma_u$ maps $[0,\delta]$
increasingly onto $[u,\delta]$, fixes the outer endpoint,
and satisfies $\sigma_0=\operatorname{id}$.  Applying it
in the collar and using the identity elsewhere gives a
jointly PL family of embeddings $e_u:X\to X$.

Define
\[
 F_u(z)=
 \begin{cases}
   e_uF e_u^{-1}(z),&z\in e_u(X),\\
   z,&z\text{ in the removed collar}.
 \end{cases}
\]
The formulas agree on the moving interface because $F$
fixes $\partial X$.  Their inverse formulas are obtained
by replacing $F$ by $F^{-1}$.  To check joint continuity,
consider the closed regions of $[0,\delta/2]\times X$
given by $z\in e_u(X)$ and by collar depth at most $u$.
They cover the parameterized domain.  The inverse collar
formula in \eqref{eq:supported-full-collar-compression}
makes the first expression continuous on its region;
the two expressions agree on the intersection.  Closed
pasting proves continuity, also at $u=0$, and the same
argument applies to the inverse family.  Thus $F_u$ is
an ordinary isotopy relative to the whole boundary.
For any fixed $u>0$, its support is disjoint from a
neighborhood of that boundary.  Relative isotopy
invariance leaves its formal PL difference class unchanged.

Extend such an $F_u$ by the identity on
$(B\setminus\operatorname{int}W)\times I$ and denote the
result by $\widehat F$.  The collar identity gives a
homeomorphism and a continuous inverse, and
$\widehat F$ fixes both ends of $B\times I$.

\paragraph{Identifying the supported class.}
Put $K=B\setminus\operatorname{int}W$ and set
\[
 Y=B\times I^2,\qquad
 A=(K\times I^2)\cup(B\times\partial I^2).
\]
Write $\lambda$ for the stable PL reduction of the product.
On the two ends of the auxiliary mapping-cylinder interval, use
the prescribed reductions $\lambda$ and
$\widehat F^{*}\lambda$.  On the original time faces and on
$K\times I^2$, use their product reductions.  These data agree
on their intersections because $\widehat F$ is the identity
on a neighborhood of $K\times I$ and of both original time
faces.  They specify the boundary reduction over $A$, so the
relative mapping-cylinder obstruction determines a class
in $H^4(Y,A)$.
The quotient identification
\begin{equation}\label{eq:supported-three-relative-pair}
 Y/A\cong
 \frac{W\times I^2}
      {(\partial W\times I^2)\cup(W\times\partial I^2)}
 \cong (W/\partial W)\wedge(I^2/\partial I^2)
\end{equation}
is equivalently the relative excision isomorphism.
Naturality of the stable PL obstruction identifies the
restriction of the supported ambient class with the
local mapping-cylinder class of $F_u$.

Forgetting the extra relative spatial complement maps
$H^4(Y,A)$ to $H^4(B\times I^2,B\times\partial I^2)$.
Under double interval suspension this is exactly
\[
            j_!:H^2(W,\partial W)\longrightarrow H^2(B).
\]
The normalization of $c_B$ on maps fixing both ends
therefore gives
\begin{equation}\label{eq:supported-three-obstruction-value}
                      c_B([\widehat F])=j_!a.
\end{equation}
The relative pairs include the complete time-square boundary
and the complete spatial side.

The mod-two fundamental class of $B$ restricts by excision
to the relative fundamental class of $W$.  Naturality
of cap product consequently gives
\[
        \operatorname{PD}_B(j_!a)
          =j_*\operatorname{PD}_{W,\partial W}(a).
\]
These duality isomorphisms apply without an integral
orientation of $B$
\cite[Section~3.3, Theorem~3.43, p.~254]{Hatcher2002Algebraic}.
This proves the second equality in
\eqref{eq:supported-three-image} and the one-support
realization assertion.

For the pairwise disjoint $W_i$, the constructed extensions
have disjoint supports and commute.  The homomorphism
property of $c_B$ makes the values of their products add.
This proves the sum assertion.  The surgery construction
has been used only on the oriented local products
$W_i\times I$; there is no fundamental-group hypothesis
on the ambient $B$.
\end{proof}

\Needspace{12\baselineskip}
\begin{corollary}\label{cor:connected-sum-three-obstruction-quotient}
Suppose $B=C_1\#\cdots\#C_s\#D$, where the $C_i$ satisfy
the hypotheses of
Theorem~\ref{thm:supported-three-concordance-realization}
and $D$ is any closed, connected PL three-manifold.
Then the homomorphism
\begin{equation}\label{eq:connected-sum-three-obstruction-quotient}
 \pi_0P_{\mathrm{TOP}}(B)\xrightarrow{c_B}H^2(B)
 \xrightarrow{\operatorname{PD}_B}H_1(B)
 \longrightarrow\bigoplus_i H_1(C_i)
\end{equation}
is onto, and its kernel contains the PL and every
fixed-real-o-minimal definable comparison image.
\end{corollary}

\begin{proof}
Choose the disjoint punctured summands $W_i\subset B$.
Deleting a three-ball leaves first homology unchanged.
Mayer--Vietoris for each gluing sphere gives
\[
 H_1(B)\cong\bigoplus_i H_1(C_i)\oplus H_1(D):
\]
the sphere has zero first homology and its map on
zero-dimensional homology to the two connected pieces
is injective.  Iterating this argument proves the
displayed decomposition, with the inclusions of the
$W_i$ identifying the indicated summands.
The supported realizations prove surjectivity of
\eqref{eq:connected-sum-three-obstruction-quotient}.
The kernel containments follow from the PL and real
definable vanishing of $c_B$.
\end{proof}

Thus the quotient by the kernel of the displayed homomorphism
is $\bigoplus_iH_1(C_i)$.  Determining that kernel and the
two comparison images remains a further question.


\subsection{Vanishing of the three-torus obstruction}

For the standard real PL torus $T^3=\mathbb R^3/\mathbb Z^3$,
the obstruction of
Theorem~\ref{thm:natural-three-concordance-obstruction}
has the following value.

\begin{theorem}\label{thm:three-torus-concordance-obstruction}
With the ordinary bottom-fixed, top-free concordance convention,
\begin{equation}\label{eq:three-torus-obstruction-zero}
 c_{T^3}:\pi_0P_{\mathrm{TOP}}(T^3)
                   \longrightarrow H^2(T^3;\mathbb Z/2)
 \qquad\text{is zero.}
\end{equation}
\end{theorem}

\begin{proof}
Write $X=T^3\times I$.  We first make both end maps the
identity, and then construct a relative five-dimensional
structure whose normal invariant detects the given obstruction.

\emph{End maps and collars.}
For any $H\in P_{\mathrm{TOP}}(T^3)$, projecting $H(x,s)$
to $T^3$ gives a homotopy from the identity to its top map.
The endpoint normalization in
Theorem~\ref{thm:natural-three-concordance-obstruction}
permits us to assume that this top map $a$ is PL, without
changing $c_{T^3}([H])$.  The three-torus is irreducible and
contains an incompressible coordinate two-torus.  Hence
Waldhausen's homotopy--isotopy theorem applies to $a$:
$a$ is TOP isotopic to the identity
\cite[Theorem~7.1, pp.~80--81; proof pp.~81--85]{Waldhausen1968}.
The component injectivity in
\eqref{eq:three-manifold-full-automorphism-comparison}
therefore supplies a joint PL isotopy $\gamma_s$ from $a$
to the identity.  Its trace
\[
 E_\gamma(x,s)=(\gamma_s a^{-1}(x),s)
\]
is a PL concordance with bottom map the identity and top
map $a^{-1}$.  Replacing $H$ by $E_\gamma H$ makes both
end maps the identity and leaves $c_{T^3}$ unchanged.

We can moreover arrange that the resulting map $F$ is the
identity on product collars of both ends.  Explicitly, for
$0<r<1/2$ put $e_r(x,s)=(x,r+(1-2r)s)$, use
$e_rFe_r^{-1}$ on $T^3\times[r,1-r]$, and use the identity
on the two complementary collars.  At $r=0$ use $F$.
The pieces agree on their interfaces; applying the same
construction to $F^{-1}$ gives the inverse.  Uniform
continuity of $F$ and $F^{-1}$ on the compact cylinder
proves continuity at $r=0$.  This is an ordinary isotopy
relative to both ends, so it preserves the obstruction.

\emph{Removing the lift displacement.}
Since the bottom inclusion induces an isomorphism on
fundamental groups and $F$ fixes that bottom, $F_*$ is
the identity on $\mathbb Z^3$.  Choose the lift
\[
 \widetilde F(x,s)=(u(x,s),v(x,s)):
                    \mathbb R^3\times I\longrightarrow
                    \mathbb R^3\times I
\]
whose bottom restriction is $(x,0)$.  It is
$\mathbb Z^3$-equivariant.  On the upper collar its first
coordinate is $x+z$ for a constant $z\in\mathbb Z^3$,
whereas on the lower collar it is $x$.
Choose a PL function $\rho:I\to I$ that is zero on a
lower collar and one on an upper collar, and set
\begin{equation}\label{eq:three-torus-lift-correction}
 L_z([x],s)=([x-\rho(s)z],s).
\end{equation}
This is a PL homeomorphism: on the finitely many intervals
where $\rho$ is affine it is induced by an affine map,
and a finite subdivision accommodates the quotient
$\mathbb R^3\to T^3$.  It is the identity on both end
collars.  Its obstruction is zero.  The chosen lift of
$L_zF$ has first coordinate $u(x,s)-\rho(v(x,s))z$,
so it agrees with $x$ on both collars.  Replace $F$ by
$L_zF$ and retain the notation $F$.  We now have
\[
 \widetilde F(x,s)=(x+w(x,s),v(x,s)),
\]
where $w$ is $\mathbb Z^3$-periodic and $w=0$, $v=s$
on both end collars.  Consequently
\begin{equation}\label{eq:three-torus-relative-homotopy}
 k_t([x],s)=
 \bigl([x+t w(x,s)],(1-t)s+t v(x,s)\bigr)
\end{equation}
is a well-defined homotopy of maps from the identity to
$F$, relative to those collars.  The intermediate maps
need not be homeomorphisms.  Reparametrize $t$ so that
$k_t$ is constant near $t=0$ and near $t=1$.

\emph{The specified relative structure.}
Let $M_F$ be the mapping cylinder of $F:X\to X$, with
upper label $y=F(x)$, and put $Z=X\times I$.
Define
\begin{equation}\label{eq:three-torus-relative-structure}
 U:M_F\longrightarrow Z,
 \qquad U([x,t])=(k_t(x),t).
\end{equation}
On the upper boundary, $U$ is the identity in the
$y$-label; it is also the identity on the bottom and
side boundary.  The constant end portions of $k_t$ and
the collars on which $F$ is the identity make these
identifications products on neighborhoods of all faces,
including their intersections.  Thus $U$ preserves the
specified PL boundary reductions.  Both bottom inclusions
$X\to M_F$ and $X\to Z$ are homotopy equivalences, and
$U$ restricts to the identity on that bottom.  Hence
$U$ is a homotopy equivalence of compact five-manifolds
relative to their full boundary.

Transport the stable TOP normal bundle of $M_F$ across
$U$, together with the induced identification of its
Spivak normal fibration and the fixed boundary comparison.
This gives the relative normal invariant
\[
 \nu(U)\in[Z,\partial Z;G/\mathrm{TOP},*].
\]
Its surgery obstruction is zero because $U$ is a
homotopy equivalence.  The equality
$\operatorname{Wh}(\mathbb Z^3)=0$
\cite[Example~24.16, p.~278]{Ranicki1992TopologicalManifolds}
also identifies the simple and homotopy decorations here.

\emph{The exact assembly input.}
The relative normal-invariant sequence applies to the
five-dimensional manifold $Z$.  In its normal-invariant
term, Poincar\'e--Lefschetz duality identifies the surgery
obstruction map with quadratic assembly:
\begin{equation}\label{eq:three-torus-relative-assembly}
 \begin{split}
 [Z,\partial Z;G/\mathrm{TOP},*]
   &\cong H_5(Z;\mathbb L\langle1\rangle)
    \cong H_5(T^3;\mathbb L\langle1\rangle),\\
 \theta&=A:
 H_5(T^3;\mathbb L\langle1\rangle)
       \longrightarrow L_5(\mathbb Z[\mathbb Z^3]).
 \end{split}
\end{equation}
Here $\mathbb L\langle q\rangle$ is the $q$-connective
quadratic surgery spectrum.  The precise relative input
is \cite[Definition~18.2, pp.~197--198; the relative-version
paragraph on p.~198; Proposition~18.3 and Theorem~18.5,
pp.~198--201]{Ranicki1992TopologicalManifolds}.
The free-abelian assembly calculation in
\cite[Example~24.16, p.~278]{Ranicki1992TopologicalManifolds}
gives
\[
 H_5(T^3;\mathbb L\langle0\rangle)
                \xrightarrow{\ \cong\ }
 L_5(\mathbb Z[\mathbb Z^3]),
\]
since five is at least the rank three.
The cofiber of
$\mathbb L\langle1\rangle\to\mathbb L\langle0\rangle$
is $H\mathbb Z$.  Its exact sequence contains
\[
 H_6(T^3;\mathbb Z)\longrightarrow
 H_5(T^3;\mathbb L\langle1\rangle)\longrightarrow
 H_5(T^3;\mathbb L\langle0\rangle)\longrightarrow
 H_5(T^3;\mathbb Z).
\]
The outside terms vanish, so the middle comparison is
an isomorphism.  Thus $\theta$ in
\eqref{eq:three-torus-relative-assembly} is injective,
and $\nu(U)=0$.

\emph{Returning to the relative PL obstruction.}
Forgetting the spherical-fibration comparison in this
zero normal invariant gives a homotopy, relative to the
specified boundary comparison, from the transported
TOP normal reduction of $M_F$ to the standard normal
reduction of $Z$.  The latter has its standard PL lift.
Pulling that lift back along the homotopy and then along
$U$ gives a stable PL normal reduction on $M_F$ extending
its prescribed boundary reduction.  In particular its
relative normal Kirby--Siebenmann obstruction is zero.
The stable tangent and normal obstructions agree, since
their Whitney sum is trivial and the obstruction is
additive with mod-two coefficients.  Therefore
\[
 \operatorname{ks}(M_F,\partial M_F)=0.
\]
By the defining mapping-cylinder suspension of the
relative Casson--Sullivan class, followed by the interval
suspension $H^2(T^3;\mathbb Z/2)\cong
H^3(X,\partial X;\mathbb Z/2)$, this implies
$c_{T^3}([F])=0$.  Each preceding normalization preserved
$c_{T^3}$, which proves
\eqref{eq:three-torus-obstruction-zero}.
\end{proof}

The conclusion determines this particular component
obstruction.  A classification of the ordinary concordance
components would require further information beyond its value.

\paragraph{A proof by finite covers.}
There is also a direct covering argument for
Theorem~\ref{thm:three-torus-concordance-obstruction}.  Its geometric
input is the relative local deformation theorem of Edwards--Kirby
\cite[Theorem~5.1(1)--(2) and its boundary-pair Note, p.~70;
proof of Corollary~1.1, p.~77]{EdwardsKirby1971}.
We give the argument because it also describes the ordinary components
of sufficiently large lifts.

Put $B=T^3$, $M=B\times I$, and
$\mathcal P=P_{\mathrm{TOP}}(B)$.  For $N\geq1$ let
$p_N([x])=[Nx]$ and $q_N=p_N\times1_I$.
Every $H\in\mathcal P$ has a unique lift $H_N$ through $q_N$
which fixes the bottom pointwise.  Indeed $H_*=1$ on
$\pi_1M$, because the bottom inclusion is a homotopy equivalence.
The bottom-fixed universal lift and its inverse are therefore
$\mathbb Z^3$-equivariant homeomorphisms.  Write the lift as
\[
 \widetilde H(x,t)=(x+u(x,t),v(x,t)).
\]
The function $u$ is periodic and hence bounded on the compact
fundamental cylinder.  If its sup norm is at most $C$, then
\begin{equation}\label{eq:three-torus-cover-displacement}
 H_N([x],t)=
 \left([x+N^{-1}u(Nx,t)],v(Nx,t)\right),
 \qquad
 \sup_{x,t}d_B(\operatorname{pr}_BH_N(x,t),x)\leq C/N,
\end{equation}
where $d_B$ is the quotient sup metric.
The same bound holds for the top map $g_N$.

For a bottom-fixed homeomorphism $J$ with top map $j$, and $0<h\leq1$,
define
\[
 T_h(J)(x,t)=
 \begin{cases}
  (\operatorname{pr}_BJ(x,t/h),
                h\operatorname{pr}_IJ(x,t/h)),&0\leq t\leq h,\\
  (j(x),t),&h\leq t\leq1.
 \end{cases}
\]
The formulas agree at $t=h$.  The lower piece is conjugation
into the lower slab; the upper piece is a product map.
Replacing $J,j$ by $J^{-1},j^{-1}$ gives the inverse.
Thus $T_h(J)\in\mathcal P$ and $T_1(J)=J$.
For every fixed $h_0>0$ the formulas and their inverses are jointly
continuous for $h\in[h_0,1]$, by pasting on the two closed regions.
They give an ordinary path from $J$ to $T_{h_0}(J)$.
With the product max metric on $M$,
\[
 \sup_{z\in M}d_M(T_h(H_N)(z),z)\leq\max\{C/N,h\}.
\]
In the lower slab both time coordinates lie in $[0,h]$; in the
upper slab time is unchanged, and the horizontal estimate holds
on both pieces.

Edwards--Kirby's theorem gives a neighborhood $\mathcal U$ of the
identity whose bottom-fixed members contract to the identity through
bottom-fixed homeomorphisms.  To see the exact relative hypothesis,
take $C=U=M$ in that theorem and take
$D_1=V_1=B\times\{0\}$ in its boundary-pair Note.  The bottom
is both closed and open in $\partial M$, so $V_1$ is a neighborhood
of $D_1$ there.  For compact $M$, the proper self-embedding space
in the theorem is $\operatorname{Homeo}(M)$, as in its cited
Corollary~1.1 proof.  The deformation ends at the identity on all
of $M$ and preserves the prescribed bottom.
Choose $\varepsilon>0$ so that uniform displacement less than
$\varepsilon$ implies membership in $\mathcal U$, and fix
$0<h_0<\min\{1,\varepsilon\}$.  For every sufficiently large $N$,
the displayed estimate puts $T_{h_0}(H_N)$ in $\mathcal U$.
The compression path followed by the local deformation proves
\begin{equation}\label{eq:three-torus-large-cover-component}
 [H_N]=[1]\quad\text{in }\pi_0\mathcal P
 \quad\text{for every sufficiently large }N.
\end{equation}

It remains to verify the covering identity for the obstruction,
including its endpoint data.  Let $g$ be the top map of $H$.
The three-dimensional automorphism comparison, followed by
Waldhausen's theorem as above, gives an ordinary isotopy
$\beta_s$ from $g$ to the identity.  The maps
\[
 A_s(x,t)=(\beta_{st}g^{-1}(x),t),\qquad H_s=A_sH
\]
form a bottom-fixed isotopy from $H$ to a map $F=H_1$ fixing
both ends.  Their inverse formulas use $g\beta_{st}^{-1}$,
so the total inverse family is continuous as well.
Lift this entire isotopy through $q_N$, beginning at $H_N$.
Its endpoint $F_N$ fixes the bottom but can have upper map
translation by $a/N$, for some $a\in\mathbb Z^3$.

For any fixed $w\in\mathbb R^3$ put
$S_w([x],t)=([x-tw],t)$.  This is PL on the whole cylinder:
its universal lift is affine, and subdivision along the finitely
many quotient hyperplanes in a fundamental cube gives a finite
PL description.  Its inverse is $S_{-w}$, and it fixes the bottom.
Set
\[
 K_N=S_{a/N}F_N,\qquad K=S_aF.
\]
Both maps fix both ends, because $a$ is integral, and
\[
 q_NS_{a/N}=S_aq_N,\qquad q_NK_N=Kq_N.
\]
The obstruction is a component homomorphism and vanishes on PL
concordances.  Therefore
\[
 c_B[K_N]=c_B[H_N],\qquad c_B[K]=c_B[H].
\]

For the both-ends-fixed maps $K_N,K$, their marked auxiliary
mapping cylinders form a covering.  The tangent microbundle,
its PL boundary reductions, and the comparisons induced by these
maps all pull back: $q_N$ is locally PL and the displayed square
commutes.  Naturality of the relative obstruction in the definition
of $c_B$ gives
\[
 \operatorname{cs}_{\mathrm{PL}}(K_N)
       =q_N^*\operatorname{cs}_{\mathrm{PL}}(K)
       \quad\text{in }H^3(M,\partial M;\mathbb Z/2).
\]
The interval suspension $\sigma:H^2(B;\mathbb Z/2)\to
H^3(M,\partial M;\mathbb Z/2)$ satisfies
$q_N^*\sigma=\sigma p_N^*$, since $q_N$ is the identity
on the interval.  The defining both-ends-fixed formula for $c_B$
now proves
\begin{equation}\label{eq:three-torus-cover-pullback}
 c_B[H_N]=p_N^*c_B[H].
\end{equation}
This establishes the covering formula with its required markings.

Finally, if $x_1,x_2,x_3$ are the circle generators in mod-two
degree-one cohomology, then $p_N^*(x_ix_j)=N^2x_ix_j$.
Thus every odd $N$ acts identically on $H^2(T^3;\mathbb Z/2)$.
Choose such an $N$ large enough for
\eqref{eq:three-torus-large-cover-component}; then
\eqref{eq:three-torus-cover-pullback} gives $c_B[H]=0$.
The contraction of $H_N$ need not be equivariant under the deck
group, so this argument makes no assertion about an isotopy of
the original concordance or surjectivity of PL comparison.


\subsection{Projective three-space concordances and odd local homeomorphisms}

A concordance of real projective three-space lifts to an odd
homeomorphism of an annulus.  Filling its inner boundary by the
identity relates the concordance space to the group of odd local
homeomorphisms at the origin.  The distinction between extending
one local map and extending a whole parameter family is central
to this description.

Let $D=D^4$ be the closed Euclidean unit ball, let $S=\partial D$,
and let $\iota(x)=-x$.  Give
\[
 H=\operatorname{Homeo}_{\iota}(D),\qquad
 A=\operatorname{Homeo}_{\iota}(S)
\]
their compact-open topologies.  A member of $H$ fixes $0$, the
unique fixed point of $\iota$.  Define the simplicial group
$\mathcal G_{\iota}$ of \emph{joint odd germs} as follows.
A $q$-simplex is represented by a continuous, parameter-preserving
homeomorphism between neighborhoods of
$\Delta^q\times\{0\}$ in $\Delta^q\times\mathbb R^4$,
fixing that zero section and commuting with $\iota$ in each fibre.
Two representatives are equal when they agree on a neighborhood
of the entire zero section.  Compactness of $\Delta^q$ gives
common smaller balls for maps and inverses.  Faces and
degeneracies are parameter pullbacks.  Write $\mathcal G_4$
for the same construction with unrestricted local homeomorphisms.

An arrow from a topological group $X$ to a geometric realization
below denotes the derived map represented by the evaluation weak
equivalence $|\operatorname{Sing}X|\to X$ and the indicated realized
simplicial homomorphism.  Its homotopy fibre is computed from this
span.  In particular, the cone-germ map is represented by
\[
 A\ \xleftarrow{\simeq_w}\ |\operatorname{Sing}A|
 \longrightarrow |\operatorname{Sing}H|
 \xrightarrow{|q|}\ |\mathcal G_{\iota}|.
\]

\begin{theorem}\label{thm:projective-concordance-joint-germ}
Restriction to a joint germ defines a Kan fibration
\[
 q:\operatorname{Sing}H\longrightarrow\mathcal G_{\iota}
\]
which is onto vertices.  Its identity fibre is weakly equivalent
to $\operatorname{Sing}P_{\mathrm{TOP}}(\mathbb{RP}^3)$.
Radial coning and the linear inclusion give weak equivalences
\begin{equation}\label{eq:projective-concordance-germ-fibre}
 \begin{split}
 P_{\mathrm{TOP}}(\mathbb{RP}^3)
 &\simeq_w\operatorname{hofib}_1
       \bigl(A\longrightarrow|\mathcal G_{\iota}|\bigr)\\
 &\simeq_w\operatorname{hofib}_1
       \bigl(\mathrm O(4)\longrightarrow|\mathcal G_{\iota}|\bigr)\\
 &\simeq_w\Omega\bigl|\mathcal G_{\iota}/
                                   \operatorname{Sing}\mathrm O(4)\bigr|.
 \end{split}
\end{equation}
All orientation components are included.  Forgetting oddness
induces an isomorphism of groups
\begin{equation}\label{eq:projective-concordance-forgetful-kernel}
 \pi_0P_{\mathrm{TOP}}(\mathbb{RP}^3)
 \cong
 \ker\!\left(\pi_1|\mathcal G_{\iota}|
                 \longrightarrow\pi_1|\mathcal G_4|\right).
\end{equation}
In particular, the full concordance component group is abelian.
\end{theorem}

\begin{proof}
The geometric existence input is the equivariant annulus theorem
of Boyle--Chen--Conway: every free involution on $S^3\times I$
preserving its two boundary components is conjugate to the
antipodal product involution
\cite[Theorem~2.3, p.~6]{BoyleChenConway2025}.
Their proof identifies the quotient as a projective-space
$h$-cobordism and applies the topological product theorem of
Kwasik--Schultz
\cite[Theorem~2.1, p.~530; proof pp.~530--532]{KwasikSchultz1989}.
The ordinary four-dimensional annulus theorem supplies the
underlying annulus, as in
\cite[Theorem~1.4, p.~15]{BoyleChenConway2025}.
For the parameter families we construct a continuous local
extension operator from Edwards--Kirby's deformation theorem.

\paragraph{A normalized extension operator near the identity.}
Choose $0<a<b<c<d<1$.  In the punctured disk let
$C=\{b\leq\|x\|\leq c\}$, and choose a compact annulus
$U$ with $C\subset\operatorname{int}U\subset
U\subset\{a<\|x\|<d\}$.  Its free quotient lies in
\[
 M=(\operatorname{int}D\setminus\{0\})/\iota
           \cong\mathbb{RP}^3\times(0,1).
\]
An odd local homeomorphism $f$, defined on a ball containing
$U$, fixing $0$, and sufficiently close to the identity,
induces an embedding $\bar f:\bar U\to M$ near the inclusion.
We restrict the neighborhood of $f$ so that the relevant
closed source ball has image in $\operatorname{int}D$.

Edwards--Kirby's Theorem~5.1(1) deforms $\bar f$ through
embeddings $\varphi(\bar f,t)$, unchanged near $\partial\bar U$,
to an embedding equal to the identity on $\bar C$.  The
deformation fixes the inclusion and depends continuously on
$\bar f$ \cite[p.~70]{EdwardsKirby1971}.
The images of $\bar U$ are constant during this deformation:
the boundary is fixed, and invariance of domain makes the
membership of every point off that boundary constant along
the isotopy.  Thus the formula
\[
 \bar k_{f,t}=\bar f\,\varphi(\bar f,t)^{-1}
 \quad\hbox{on }\bar f(\bar U),
 \qquad \bar k_{f,t}=1\quad\hbox{elsewhere}
\]
defines an ambient isotopy.  Near the joining boundary the
two formulas agree.  This is the ambient correction in the
proof of \cite[Corollary~1.2, p.~79]{EdwardsKirby1971}.
After shrinking the embedding neighborhood, its support lies
in one compact annulus.  It starts at the identity, its
endpoint agrees with $\bar f$ on $\bar C$, and the entire
isotopy is the identity when $f$ is the identity.

Lift this isotopy through the double cover, starting at the
identity.  The lift $k_{f,t}$ is odd, is identical near $0$
and $\partial D$, and extends there by the identity.
Covering homotopy gives continuous dependence on $f$ and $t$.
Its endpoint $k_f$ agrees with $f$ on $C$: the two lifts
differ there by a constant deck transformation, and closeness
to the identity at one point of the connected collar selects
the identity deck transformation.  Fix $r\in(b,c)$ and set
\begin{equation}\label{eq:projective-normalized-local-extension}
 E(f)(x)=
 \begin{cases}
  f(x),&\|x\|\leq r,\\
  k_f(x),&\|x\|\geq r.
 \end{cases}
\end{equation}
The two maps agree on an annular neighborhood of the joining
sphere.  They map that sphere to the same locally flat sphere,
and both inner regions are the component containing $0$.
Consequently \eqref{eq:projective-normalized-local-extension}
is an odd homeomorphism of $D$.  Pasting gives joint evaluation
continuity, and inversion in the compact-domain homeomorphism
group gives joint inverse continuity.  This defines a single
continuous operator, equal to $f$ near $0$, with $E(1)=1$.

\paragraph{Extension of individual germs.}
Given an odd germ $f$, choose $\epsilon>0$ so that $f$ is
defined near $D_\epsilon$ and
$f(D_\epsilon)\subset\operatorname{int}D$.
The ordinary annulus theorem identifies
$D\setminus\operatorname{int}f(D_\epsilon)$ with $S^3\times I$.
The inherited involution is free and preserves both boundaries.
The equivariant annulus theorem therefore gives an odd product
parametrization.  Precompose its sphere coordinate with the
equivariant boundary homeomorphism needed to agree with $f$
on the inner sphere, using that same correction on every
product level.  Gluing this product to $f$ on $D_\epsilon$
gives an element of $H$ representing the germ.  This proves
surjectivity on vertices.

\paragraph{Relative parameter lifting.}
Let $T=\Delta^j$, let $Z\subset T$ be a horn, and suppose
a joint germ $g$ over $T$ has a prescribed continuous
$H$-lift $F_z$ on $Z$.  The finitely many faces paste to
one family on $Z$.  Choose a retraction $r:T\to Z$ and
a deformation $r_t:T\to T$ from $r$ to the identity,
fixed on $Z$.  Put $\widehat F_p=F_{r(p)}$ and, on a
common sufficiently small ball, put
\[
 \delta_p=\widehat F_p^{-1}g_p,\qquad
 \delta_{p,t}=\delta_{r_t(p)}.
\]
The joint-germ equality on each face, followed by the
finite intersection of the resulting neighborhoods, makes
$\delta_z=1$ on one common ball for all $z\in Z$.
Thus $\delta_{p,0}=1$ for all $p$, and
$\delta_{z,t}=1$ for all $(z,t)\in Z\times I$.

Compactness and joint inverse continuity give a common
ball on which all increments
\[
 \eta_{p,t}=\delta_{p,t}\delta_{p,t_i}^{-1},
                       \qquad t_i\leq t\leq t_{i+1},
\]
are defined.  Use a fixed smaller annular collar there
for a rescaled copy of the operator $E$.
Uniform continuity near the diagonal in the compact
parameter space of $(p,t_i,t)$ supplies a finite time
partition so that each increment belongs to its domain.
Choose a further common ball $V$ whose images under all
$\delta_{p,t}$ lie in the inner ball where $E$ agrees
with its argument.  Starting with $D_{p,0}=1$, define
successively
\[
             D_{p,t}=E(\eta_{p,t})D_{p,t_i}.
\]
The formulas agree at subdivision times because $E(1)=1$.
Induction shows that $D_{p,t}=\delta_{p,t}$ on $V$;
the image condition on $V$ makes each inductive composition
valid on this same ball.  On $Z$ every factor is the
identity.  Hence $\widehat F_pD_{p,1}$ lifts $g$ and
agrees exactly with the prescribed horn family.
This proves the Kan condition.

\paragraph{The identity fibre.}
Let $\mathcal F$ be the strict identity fibre.  Each of
its simplices is an actual $H$-family equal to the
identity on a common ball $D_\epsilon$.  Fix
$0<\epsilon_0<1$, and write $H_{\epsilon_0}$ for the
subgroup fixing $D_{\epsilon_0}$ pointwise.  Taking
the antipodal quotient of the outer annulus identifies
this group homeomorphically with
$P_{\mathrm{TOP}}(\mathbb{RP}^3)$.
Conversely a bottom-fixed projective concordance lifts
uniquely with identity lift on the inner sphere and
extends by identity across the ball.  Covering lifts
and inverse evaluations are continuous for compact
parameter families.  These identifications respect
composition.

The inclusion $\operatorname{Sing}H_{\epsilon_0}\to
\mathcal F$ is a weak equivalence.  A finite simplicial
representative or homotopy in $\mathcal F$ uses one
$\epsilon\leq\epsilon_0$.  Let $a_t$ be the odd radial
homeomorphism of $D$ whose radius function is affine
between the three points
\[
 (0,0),\quad
 (\epsilon,(1-t)\epsilon+t\epsilon_0),\quad(1,1).
\]
It fixes the outer sphere and expands all radii.
Conjugation by $a_t$ moves the given family into
$H_{\epsilon_0}$ at $t=1$.  It preserves that subgroup
throughout because $a_t(D_{\epsilon_0})\supset
D_{\epsilon_0}$.  This proves surjectivity on homotopy
groups.  Applied to a null-homotopy, the same conjugation
gives a null-homotopy in $H_{\epsilon_0}$ after attaching
its boundary conjugation track, which lies in that
subgroup.  This proves injectivity.  The argument also
treats components and keeps the identity basepoint fixed.

\paragraph{Replacing disk maps by orthogonal maps.}
Boundary restriction $H\to A$ has the homomorphic
section $c(a)(ru)=r a(u)$.  Its kernel contracts by
the odd Alexander formula
\[
 \mathcal A_t(h)(x)=
 \begin{cases}
  t h(x/t),&\|x\|\leq t,\quad t>0,\\
  x,&\|x\|\geq t,
 \end{cases}
 \qquad\mathcal A_0(h)=1.
\]
Maps and inverses are continuous in all compact
parameters.  The section gives a product decomposition
of the underlying spaces into this kernel and $A$,
so restriction is a homotopy equivalence.
On the section, $q$ is precisely the cone-germ map.
The Kan fibre sequence now proves the first equivalence
in \eqref{eq:projective-concordance-germ-fibre}.

Bamler--Kleiner prove that
$\operatorname{Isom}(N)\to\operatorname{Diff}^{\infty}(N)$
is a homotopy equivalence for every round spherical
three-manifold, including $N=\mathbb{RP}^3$
\cite[Theorem~1.3, p.~2; proof p.~108]{BamlerKleiner2019Metrics}.
Cerf's comparison of diffeomorphisms with homeomorphisms
applies to compact three-manifolds; its sphere Smale
hypothesis is supplied by that same theorem
\cite[III, \S3.2.1, Theorem~8, p.~365, and Corollary~1,
p.~366; proof pp.~367--370]{Cerf1961}.

We spell out the change of differentiability convention
in this application.  Cerf permits the $C^1$ topology,
whereas the preceding diffeomorphism group uses the
$C^{\infty}$ topology.  For a closed smooth manifold,
the inclusion of these two diffeomorphism groups is
a weak equivalence.  Indeed, embed the manifold in
Euclidean space and fix a smooth retraction of a
tubular neighborhood.  Given a compact parameter
family of $C^1$ diffeomorphisms, choose finitely many
nearby smooth maps by chartwise convolution and a
spatial partition of unity.  A parameter partition
of unity forms finite linear combinations of these
maps.  Composing with the tubular retraction produces
a family continuous in the $C^{\infty}$ topology.
It can be made uniformly $C^1$-close to the original
family.  The diffeomorphisms form an open subset of
the $C^1$ mapping space, so compactness gives a
uniform neighborhood in which both this family and
the retracted straight-line interpolation consist
of diffeomorphisms.

For a sphere or disk parameter pair, first make the
original family constant in the inward normal
parameter near its boundary.  If its boundary family
is already continuous in $C^{\infty}$, use that given
family unchanged on a smaller parameter collar and
interpolate to the finite approximation away from
the collar.  The interpolation weights vanish on
the smaller collar, so this approximation and its
$C^1$ homotopy are relative to the boundary.  The
given boundary family has bounded derivatives of
each order on its compact parameter set; hence the
pasted approximating family is continuous in every
smooth seminorm.  Relative approximation proves
surjectivity and injectivity on all homotopy groups,
including the component assertion.  This establishes
the stated change of convention for $\mathbb{RP}^3$.

Every homeomorphism of $\mathbb{RP}^3$ has two lifts
to $S^3$, both commuting with the deck involution.
The lift space is a double cover of its homeomorphism
group: close maps have a unique close lift after
one sheet is specified.  The same applies to
diffeomorphisms and isometries.  The compatible
diagram of double covers over
\[
 \mathrm{PO}(4)\longrightarrow
 \operatorname{Diff}^{\infty}(\mathbb{RP}^3)
 \longrightarrow\operatorname{Homeo}(\mathbb{RP}^3)
\]
therefore gives a weak equivalence $\mathrm O(4)\to A$.
Its cone-germ map is the literal linear-germ inclusion.
This proves the second equivalence in
\eqref{eq:projective-concordance-germ-fibre}.
Finally, the simplicial subgroup
$\operatorname{Sing}\mathrm O(4)\subset\mathcal G_{\iota}$
has a Kan coset fibration.  The homotopy fibre of its
fibre inclusion is the loop space of that coset space,
which proves the last equivalence.

\paragraph{Components and the forgetful kernel.}
Vertex surjectivity and $H\simeq A\simeq_w\mathrm O(4)$
show that $\pi_0\mathrm O(4)\to\pi_0\mathcal G_{\iota}$
is onto.  The local orientation character separates
the positive and negative linear germs, making it
bijective.  The exact sequence of the simplicial
group fibration gives an isomorphism of groups
\[
 \pi_0P_{\mathrm{TOP}}(\mathbb{RP}^3)
 \cong\operatorname{coker}\bigl(
     \pi_1\mathrm O(4)\longrightarrow
     \pi_1|\mathcal G_{\iota}|\bigr).
\]
The group law on the left is composition, by the
fixed-radius fibre identification.  Fundamental
groups of simplicial groups are abelian, so this
cokernel is an abelian group.

Put $V=\pi_1|\mathcal G_{\iota}|$,
$W=\pi_1|\mathcal G_4|$, and let
\[
 j:\pi_1\mathrm O(4)\longrightarrow V,
 \qquad f:V\longrightarrow W
\]
be the linear inclusion and forgetting map.
The global-to-joint-germ comparison in the proof
of Theorem~\ref{thm:top-sphere-finite-rank-stabilization}
identifies $W$ with $\pi_1\mathrm{TOP}(4)$ while
preserving the literal linear map.  The stabilization
range used in Corollary~\ref{cor:four-sphere-low-concordance-groups}
gives $\pi_1Y_4=\pi_2Y_4=0$ for the smoothing fibre
\eqref{eq:finite-rank-top-smoothing-space}.
The sequence $\mathrm O(4)\to\mathrm{TOP}(4)\to Y_4$
therefore makes $fj$ an isomorphism.  Since
$\pi_1\mathrm O(4)=\mathbb Z/2$, it follows that
\[
 V=j(\mathbb Z/2)\oplus\ker f,
 \qquad
 V/j(\mathbb Z/2)\xrightarrow{\ \cong\ }\ker f,
 \quad [v]\longmapsto v-j\bigl((fj)^{-1}f(v)\bigr).
\]
Combining these isomorphisms proves
\eqref{eq:projective-concordance-forgetful-kernel}.
\end{proof}

Let $K$ denote the kernel in
\eqref{eq:projective-concordance-forgetful-kernel}, and
let $I\subset K$ be the image of the PL concordance
component group.  Transporting the Casson--Sullivan
homomorphism gives an epimorphism
$\chi:K\to\mathbb F_2$ that vanishes on $I$, by
Theorems~\ref{thm:natural-three-concordance-obstruction}
and~\ref{thm:cyclic-three-concordance-obstruction}.
Since $K$ is abelian, this gives the exact sequence
\[
 0\longrightarrow\ker\chi/I\longrightarrow K/I
             \xrightarrow{\ \chi\ }\mathbb F_2\longrightarrow0.
\]
Determining $K$, its PL subgroup $I$, and the remaining
quotient $\ker\chi/I$ is the resulting comparison problem.


\section{PL and TOP concordances over four-manifolds}
We identify the relative comparison fibre and the precise obstruction
for a topological concordance over a four-manifold to have a PL
representative. The input is Kirby--Siebenmann's relative classification
of PL structures on open five-manifolds
\cite[Essay~V, (3.2), pp.~240--241, and \S5.6, pp.~252--253]{KirbySiebenmann1977}.
Extending a cylinder by its two exterior collars puts the relative
problem in that dimension. Keeping the top structure variable then
identifies the stabilization map, and tracing changes of the top
isotopy determines its ambiguity on components. The three steps are
proved in that order.

\subsection{Comparison fibres over a closed four-manifold}

Let $B$ be a closed compact real PL four-manifold.  Disconnected and
nonorientable manifolds are allowed.  All PL families below are PL on
their entire parameterized total spaces; the TOP objects are ordinary
singular complexes of compact-open homeomorphism groups.  Thus these
are ordinary parameterized automorphisms, not block automorphisms.
For $\mathcal C\in\{\mathrm{PL},\mathrm{TOP}\}$, put
\[
 A_{\mathcal C}=\operatorname{Aut}_{\mathcal C}(B),\quad
 P_{\mathcal C}=\operatorname{Aut}_{\mathcal C}
       (B\times I\ \mathrm{rel}\ B\times\{0\}),\quad
 H_{\mathcal C}=\operatorname{Aut}_{\mathcal C}
       (B\times I\ \mathrm{rel}\ B\times\partial I).
\]
Here ``rel'' means pointwise fixed, and the top of a member of
$P_{\mathcal C}$ is free.  Write $F_A,F_P,F_H$ for the homotopy fibres
over the identity of the corresponding PL-to-TOP maps.  A path-space
model uses pairs $(a,\gamma)$ with $a$ in the PL object and $\gamma$
a TOP path from the identity to its image.

\begin{theorem}\label{thm:four-dimensional-comparison-fibres}
There is a based homotopy fibre sequence
\begin{equation}\label{eq:four-dimensional-comparison-fibre-sequence}
 \operatorname{Map}(B,K(\mathbb Z/2,1))
       \simeq F_H\longrightarrow F_P\longrightarrow F_A.
\end{equation}
The mapping space is unbased and is pointed by the constant map.
Consequently $\pi_jF_P\to\pi_jF_A$ is an isomorphism for $j\geq3$,
and there is an exact sequence
\begin{equation}\label{eq:four-dimensional-comparison-low-sequence}
\begin{aligned}
0&\longrightarrow\pi_2F_P\longrightarrow\pi_2F_A
 \longrightarrow H^0(B;\mathbb Z/2)
 \longrightarrow\pi_1F_P\longrightarrow\pi_1F_A\\
 &\longrightarrow H^1(B;\mathbb Z/2)
 \longrightarrow\pi_0F_P\longrightarrow\pi_0F_A.
\end{aligned}
\end{equation}
At the component terms, exactness has its usual pointed-set meaning.
This statement does not evaluate the connecting maps.
\end{theorem}

\begin{proof}
We first verify the relative convention and the fibre sequence before
applying the classical PL-structure classification.

\emph{Top restriction.}
The map $e:P_{\mathcal C}\to A_{\mathcal C}$ is a Kan fibration
with literal fibre $H_{\mathcal C}$.  Here is a horn-lifting argument
that also checks the joint PL requirement.  For a horn
$J=\Lambda_i^q\subset\Delta^q$, $q\geq1$, use barycentric coordinates
$\lambda_0,\ldots,\lambda_q$ and set
\[
 m=\min_{j\ne i}\lambda_j,\qquad
 r_j=\lambda_j-m\ (j\ne i),\qquad r_i=\lambda_i+qm.
\]
This is a PL retraction $r:\Delta^q\to J$.  There is a joint PL
homotopy $d:\Delta^q\times I\to\Delta^q$ from $r$ to the identity,
fixed on $J$: triangulate the prescribed PL data on the bottom, top,
and $J\times I$, extend the triangulation over the prism, and assign
the remaining vertices values in the convex simplex before extending
affinely.  This construction avoids treating a merely slice-wise PL
linear interpolation as a joint PL map.

Suppose a horn family $p_J$ and a simplex $a$ in $A_{\mathcal C}$
have matching top restrictions.  Set
$p^0_t=p_{r(t)}$ and $\delta_t=a_t e(p^0_t)^{-1}$; then
$\delta_t=1$ on $J$.  The level-preserving family
\[
 E_t(x,u)=(\delta_{d(t,u)}(x),u),\qquad p_t=E_t p^0_t
\]
fixes the bottom, has top $a_t$, and agrees with $p_J$ on the horn.
The maps and inverses are joint PL in the PL case and continuous in
the TOP case.  This proves the lifting assertion.  Applying homotopy
fibres to the two compatible sequences
$H_{\mathcal C}\to P_{\mathcal C}\to A_{\mathcal C}$ gives
\begin{equation}\label{eq:relative-comparison-fibre-sequence}
 F_H\longrightarrow F_P\longrightarrow F_A.
\end{equation}
Equivalently, this follows by interchanging the two homotopy-fibre
constructions in their commutative square.

\emph{Replacing a fixed boundary by a fixed collar.}
Put $N=B\times I$, and choose $0<\delta<1/2$.  For
$0\leq s\leq\delta/2$, define on a bottom collar
\[
 \sigma_s(t)=\min\{t+s,(t+\delta)/2\},\qquad 0\leq t\leq\delta.
\]
It is an increasing PL homeomorphism from $[0,\delta]$ onto
$[s,\delta]$, with joint PL inverse
\[
 \sigma_s^{-1}(y)=\max\{y-s,2y-\delta\}.
\]
Use this map on the bottom collar, its reflected version on the top
collar, and the identity between them.  The resulting joint PL
embedding $c_s:N\to B\times[s,1-s]$ has $c_0=1$.
For $h\in H_{\mathcal C}$, conjugate by $c_s$ on its image and use
the identity on the two exterior collars.  Because $h$ fixes both
ends, the formulas agree along their common boundaries.  They define
a parameterized homotopy $R_s(h)$ of relative automorphisms.  The
same construction works for a whole simplex and preserves identity
faces.

For $s>0$, the result fixes collars of both ends uniformly over the
simplex.  It also preserves every previously fixed smaller collar:
on the bottom $\sigma_s(t)\geq t$, so that the inverse image of a
point in such a collar lies in that collar; the top is identical by
reflection.  Let $H_{\mathcal C}^{\mathrm{col}}$ be the simplicial
group in which each simplex fixes some common collar of both ends.
Every finite family of its simplices has a common smaller collar.
The preceding homotopy, applied to finite sphere and disk families,
therefore proves
\[
 H_{\mathcal C}^{\mathrm{col}}\longrightarrow H_{\mathcal C}
 \quad\text{is a weak equivalence}.
\]
The construction commutes with PL-to-TOP comparison.  In particular,
it also preserves its homotopy fibre.  This argument concerns the
specified simplicial groups; it does not require exchanging singular
complexes with an unrestricted topological direct limit.

\emph{The classical relative structure space.}
Adjoin exterior collars and put
\[
 M=B\times\mathbb R,\qquad C=M\setminus(B\times(0,1)).
\]
Now $M$ is a boundaryless PL five-manifold and $C$ is closed.
Extension by the identity identifies the collar-fixed automorphism
groups with the automorphism groups of $M$ fixed near $C$.
In the PL case, restriction to the compact strip converts locally
finite PL data into finite PL data, so these are the same families.
Let $\mathcal S$ denote the simplicial space of sliced PL structures
on $M$ equal near $C$ to the given structure $\alpha$.

The relative automorphism-to-structure fibration of
Kirby--Siebenmann, Essay~V, (1.5), p.~229, applies to this $M$ and $C$
\cite{KirbySiebenmann1977}.  More explicitly, if $X$ and $Y$ are the
PL and TOP automorphism simplicial groups fixed near $C$, the quotient
$Y/X$ is Kan and its natural inclusion in $\mathcal S$ has image a
union of components containing $\mathcal S_\alpha$.  Hence
\begin{equation}\label{eq:both-ends-fibre-structures}
 F_H\simeq\Omega\mathcal S_\alpha.
\end{equation}
Only the component of the standard structure is used here.  There
is no assertion that every relative PL structure is in its orbit.

The relative classification theorem, Essay~V, (3.2), pp.~240--241,
identifies $\mathcal S$ with the space of lifts of the tangent
classifying map to $B\mathrm{PL}(5)$, fixed at the given lift near
$C$.  Its dimension hypothesis is satisfied because $M$ is open of
dimension five.  By Theorem~5.5(I), p.~251, stabilization gives
\[
 \mathrm{TOP}(5)/\mathrm{PL}(5)\simeq
 \mathrm{TOP}/\mathrm{PL}\simeq K(\mathbb Z/2,3).
\]
The identification with an untwisted mapping space can be made
explicit.  Essay~V, \S5.6, pp.~252--253, identifies
$B\mathrm{PL}\to B\mathrm{TOP}$ with the homotopy-fibre construction
of the universal class
$k:B\mathrm{TOP}\to K(\mathbb Z/2,4)$.  The chosen PL structure
provides a reference nullhomotopy of the composite tangent map.
Comparing another nullhomotopy with this reference gives a map into
$\Omega K(\mathbb Z/2,4)=K(\mathbb Z/2,3)$, constant near $C$.
This is the relative version of the equivalence in \S5.6; it accounts
for the chosen lift rather than silently discarding a possible
twisting of the lift problem.

Maps constant near the two ends of $N$ have the same weak homotopy
type as maps constant on the ends.  Indeed, precompose with a collar
collapse.  For any finite parameter family this makes the map
constant on a common smaller collar; when a boundary parameter
family already has this property, choose the collapse inside its
fixed collar so that it is preserved.  Thus
\begin{equation}\label{eq:relative-pl-structure-mapping-space}
 \mathcal S\simeq
 \operatorname{Map}_*(N/\partial N,K(\mathbb Z/2,3)).
\end{equation}
Since $N/\partial N\cong\Sigma B_+$, looping its standard component
and using adjunction in \eqref{eq:both-ends-fibre-structures} yields
\[
\begin{aligned}
 F_H
 &\simeq\Omega\operatorname{Map}_*(\Sigma B_+,K(\mathbb Z/2,3))\\
 &\simeq\operatorname{Map}_*(B_+,\Omega^2K(\mathbb Z/2,3))
 \simeq\operatorname{Map}(B,K(\mathbb Z/2,1)).
\end{aligned}
\]
The first loop only uses the component of the constant map.
Its component group is $H^1(B;\mathbb Z/2)$, its fundamental group
at the constant map is $H^0(B;\mathbb Z/2)$, and its higher homotopy
groups vanish.  The long exact sequence of
\eqref{eq:relative-comparison-fibre-sequence} proves all assertions.
The empty manifold gives the same statements with all spaces
contractible and all displayed cohomology groups zero.
\end{proof}

\paragraph{The dimension restriction.}
The preceding argument is a consequence of the classical relative
PL-structure classification.  It does not classify PL structures on
$B$ itself and gives no surjectivity assertion for
$\pi_0P_{\mathrm{PL}}\to\pi_0P_{\mathrm{TOP}}$.  In particular,
one cannot apply the automorphism fibration directly to the compact
five-manifold $B\times I$: the printed hypothesis on p.~229 also
excludes a four-dimensional boundary.  The exterior collars remove
that issue.  Nor does the same argument prove the analogous formula
for a closed three-manifold base.  The open ambient manifold then
has dimension four, and the relative theorem on p.~240 requires
$\check H^4(M,C;\mathbb Z/2)=0$.  For this collared polyhedral pair
it is $H^3(B;\mathbb Z/2)$, which is nonzero on each nonempty closed
connected three-manifold.  That additional hypothesis fails.


\subsection{The stabilization map of the based comparison fibre}

Let $B$ be a closed compact real PL four-manifold, possibly
disconnected or nonorientable.  Use ordinary joint PL families and
ordinary singular families of compact-open TOP homeomorphisms.
For $\mathcal C\in\{\mathrm{PL},\mathrm{TOP}\}$, let
$A_{\mathcal C}=\operatorname{Aut}_{\mathcal C}(B)$ and let
$P_{\mathcal C}$ be the automorphisms of $B\times I$ fixing the
bottom pointwise, with free top.  Write $F_A,F_P$ for the identity
homotopy fibres of their PL-to-TOP comparison maps.  We specify a
map, as well as its domain and target.  Let $\mathcal S_B$ be the
simplicial space of sliced PL
structures on $B$, based at its given structure $\alpha$.  Thus a
simplex is a PL structure on $\Delta^q\times B$ whose projection is
a PL submersion.  Put
\[
 M=B\times\mathbb R,\qquad C_+=B\times[1,\infty),
\]
and let $\mathcal T$ be the space of germs of such structures near
$C_+$, based at the standard product structure.  Product with the
line defines $s:\mathcal S_B\to\mathcal T$.

Fix tangent classifying maps and their standard PL lifts.  The open
five-dimensional classification, with the standard lift as origin,
gives a based equivalence
\[
 \eta:\mathcal T\simeq
       \operatorname{Map}(B,K(\mathbb Z/2,3)).
\]
All classification maps may be interpreted as maps of based homotopy
types.  The orbit of $\alpha$ gives a map
$\rho:F_A\to\Omega\mathcal S_B$: a pair $(a,\gamma)$, where $a$ is
PL and $\gamma$ is a TOP path from the identity to $a$, is sent to
the loop $(\gamma_t)_*\alpha$.  Define
\begin{equation}\label{eq:four-dimensional-stabilization-map}
\begin{aligned}
 \beta_B:F_A&\xrightarrow{\rho}\Omega\mathcal S_B
   \xrightarrow{\Omega s}\Omega\mathcal T
   \xrightarrow{\Omega\eta}
      \Omega\operatorname{Map}(B,K(\mathbb Z/2,3))\\
 &\hspace{43mm}\simeq\operatorname{Map}(B,K(\mathbb Z/2,2)).
\end{aligned}
\end{equation}
The mapping spaces are unbased, pointed by their constant maps.
The loop spaces include every component of the based loop space.

\begin{theorem}\label{thm:four-dimensional-stabilization-fibre}
For the map \eqref{eq:four-dimensional-stabilization-map},
\begin{equation}\label{eq:four-dimensional-stabilization-hofibre}
 F_P\simeq\operatorname{hofib}_0(\beta_B).
\end{equation}
In particular there is an exact pointed-set sequence
\begin{equation}\label{eq:four-dimensional-stabilization-components}
 H^1(B;\mathbb Z/2)\longrightarrow\pi_0F_P
 \longrightarrow\pi_0F_A
 \xrightarrow{\pi_0\beta_B}H^2(B;\mathbb Z/2).
\end{equation}
For a representative $(a,\gamma)$, reparameterize $\gamma$ to be
constant near both ends.  The value of $\pi_0\beta_B$ is the
relative PL-structure difference class of
\[
 (x,u)\longmapsto(\gamma_u(x),u)
\]
on $B\times I$, with the standard cylinder structure as origin,
under the identification
\[
 H^3(B\times I,B\times\partial I;\mathbb Z/2)
       \cong H^2(B;\mathbb Z/2).
\]
Thus the image of $\pi_0F_P\to\pi_0F_A$ is precisely the
zero-obstruction locus of this specified map.
\end{theorem}

\begin{proof}
The empty case is immediate.  In all the structure spaces used below,
``near'' means a germ on an open neighborhood.  Every simplex has
one collar or neighborhood common to its entire parameter simplex;
finite families have a common smaller one.

\emph{Compact structure orbits.}
There is an injective orbit map
\[
 \operatorname{Aut}_{\mathrm{TOP}}(B)/
 \operatorname{Aut}_{\mathrm{PL}}(B)\longrightarrow\mathcal S_B.
\]
Its image is a union of components containing $\alpha$.  The
compactness hypothesis is the reason this assertion is valid in
dimension four.  The footnote to Kirby--Siebenmann, Essay~V,
Bundle Theorem~0.1, p.~217, states that a proper PL submersion is a
locally trivial bundle, so the bundle assertion for compact fibres
has no dimension restriction \cite{KirbySiebenmann1977}.
A simplex of $\mathcal S_B$ is such a proper submersion.  Trivialize
its PL bundle over the simplex.  If one vertex has the original PL
isomorphism type, choose the fibre identification there to obtain
a lift through the orbit map.  If a trivialization is prescribed on
a horn, its difference from the chosen trivialization is a horn in
the PL automorphism simplicial group.  Filling that horn corrects
the trivialization.  This proves the orbit and lifting assertions.
The quotient map of simplicial groups is a Kan fibration, and hence
\begin{equation}\label{eq:compact-four-orbit-loop}
 \rho:F_A\simeq\Omega(\mathcal S_B,\alpha).
\end{equation}
The loop space contains all based loops in the standard component;
no other PL isomorphism type has been placed in that orbit.

\emph{Product collars and cylinder orbits.}
Write $N=B\times I$.  Let $P_{\mathcal C}^{\mathrm{pc}}$ be the
concordance simplicial group whose maps are the identity near the
bottom and equal to $e(h)\times\operatorname{id}$ near the top.
There is a weak equivalence
$P_{\mathcal C}^{\mathrm{pc}}\to P_{\mathcal C}$ for both categories.
For completeness, choose $0<\delta<1/2$ and compress the two collars
using
\[
 \sigma_s(t)=\min\{t+s,(t+\delta)/2\},\quad
 \sigma_s^{-1}(y)=\max\{y-s,2y-\delta\},
 \qquad 0\leq s\leq\delta/2.
\]
Together with reflection at the top and the identity in the middle,
these give $c_s:N\to B\times[s,1-s]$.  Conjugate $h$ by $c_s$ on
the image, extend by the identity below it, and by
$e(h)\times\operatorname{id}$ above it.  The values agree at the
two seams.  The formula starts at $h$, has uniform product collars
for $s>0$, and is joint PL or continuous as appropriate.  Previously
existing product collars remain product because $\sigma_s(t)\geq t$
and the compression commutes with the boundary map times the
identity.  It preserves identity parameter faces.  Applying the
construction to finite sphere and relative disk families proves
the asserted weak equivalence.

Let $\mathcal S_N^{\mathrm{pc}}$ be the sliced PL structures on $N$
standard near the bottom and product near the top in the fixed
collar coordinates, with arbitrary induced top structure.  The
orbit map
\[
 P_{\mathrm{TOP}}^{\mathrm{pc}}/
 P_{\mathrm{PL}}^{\mathrm{pc}}\longrightarrow
 \mathcal S_N^{\mathrm{pc}}
\]
again has image a union of components containing the standard
cylinder structure $\alpha_N$.  Here are the relative details.
Trivialize a simplex as a proper PL bundle with compact fibre $N$,
using a chosen standard-cylinder identification at one vertex.
This also ensures that each boundary piece remains the corresponding
boundary piece throughout the connected parameter simplex.  The
trivialization may initially be nonproduct on the boundary collars.
Choose a width inside the common product collars of its source and
target, and apply the preceding compression, extending by the two
boundary maps times the identity.  The compression is PL in both
structures, since both are products on its support.  Thus the new
trivialization is still a PL isomorphism and is product at the ends.
Its bottom family $b_t$ is a joint PL automorphism of the standard
$B$.  Precomposition by $b_t^{-1}\times\operatorname{id}$ fixes the
bottom pointwise.  A prescribed horn is restored by filling the
difference horn in the simplicial group
$P_{\mathrm{PL}}^{\mathrm{pc}}$.  Consequently
\begin{equation}\label{eq:compact-cylinder-orbit-loop}
 F_P\simeq\Omega(\mathcal S_N^{\mathrm{pc}},\alpha_N).
\end{equation}
This uses the compact proper-submersion input, not a classification
theorem for a five-manifold with four-dimensional boundary.

\emph{The restriction square.}
Put $C_-=B\times(-\infty,0]$, and let $\mathcal S_-$ be the sliced
PL structures on $M$ equal to the standard structure near $C_-$.
Restriction to $C_+$ gives $r:\mathcal S_-\to\mathcal T$.
There is a strict pullback square
\begin{equation}\label{eq:four-dimensional-collar-pullback}
\begin{array}{ccc}
 \mathcal S_N^{\mathrm{pc}}&\longrightarrow&\mathcal S_-\\
 \big\downarrow&&\big\downarrow r\\
 \mathcal S_B&\xrightarrow{\ s\ }&\mathcal T.
\end{array}
\end{equation}
Indeed, a structure in the upper-left corner extends as its top
product on the entire positive exterior half-cylinder and as the
standard structure on the negative one.  Conversely, equality of
the top germs in the pullback means precisely that the structure
is this product on a neighborhood of all $C_+$.  Restriction to $N$
therefore has the required top collar.  Compactness of
$\Delta^q\times B$ supplies a common collar width.  These operations
are inverse, including on faces and degeneracies.  This is the
collar pullback construction in Essay~V, (4.1), pp.~242--243, with
the bottom held fixed; it leaves $\mathcal S_B$ unclassified.

The map $r$ is a Kan fibration.  Use restriction to the closed subset
$C_-\cup C_+$ in the boundaryless five-manifold $M$ and then take
the fibre over the standard germ at $C_-$.  The two subsets have
disjoint neighborhoods, so their germ spaces form a product.
The required restriction theorem is property~(4) in the proof of
Essay~V, (1.4), pp.~224--225, with the relative formulation (3.1)--(3.2),
p.~240.  Its dimension-five hypotheses hold.  This argument needs
no surjectivity on all components of $\mathcal T$.

The relative open classification gives
\begin{equation}\label{eq:one-ended-structure-classification}
 \mathcal S_-\simeq *,\qquad
 \eta:\mathcal T\simeq\operatorname{Map}(B,K(\mathbb Z/2,3)).
\end{equation}
For the first equivalence, $C_-\hookrightarrow M$ is a collared
homotopy equivalence, so its relative lift space is contractible.
For the second, the neighborhoods
$B\times(1-\varepsilon,\infty)$ are cofinal for every finite germ
family.  The entire positive half-cylinder is already in $C_+$,
and compactness at $B\times\{1\}$ supplies a uniform
$\varepsilon>0$.  Each such neighborhood is an open five-manifold
homotopy equivalent to $B$.  The natural classification of germs,
Essay~V, (2.2.1)--(2.3.1), pp.~234--235, identifies its restrictions
with the restrictions of tangent lifts.  These are weak equivalences
as the collar shrinks.  The finite-family germ colimit therefore has
the stated mapping-space type.  In both assertions use (3.2),
Theorem~5.5(I), and the universal-class trivialization in \S5.6,
pp.~240, 251--253, with the standard PL lift.  All manifolds being
classified are boundaryless and five-dimensional.

Since $r$ is a Kan fibration, \eqref{eq:four-dimensional-collar-pullback}
is a homotopy pullback.  Its upper-right corner is weakly
contractible, giving
\[
 \mathcal S_N^{\mathrm{pc}}\simeq
 \operatorname{hofib}_0\bigl(
     \eta s:\mathcal S_B\to
       \operatorname{Map}(B,K(\mathbb Z/2,3))\bigr).
\]
This includes exactly the components whose stabilized germs can be
joined to the standard germ.  Looping, using
\eqref{eq:compact-four-orbit-loop} and
\eqref{eq:compact-cylinder-orbit-loop}, proves
\eqref{eq:four-dimensional-stabilization-hofibre}, because based
loops preserve homotopy fibres.  Naturality of the classification
zigzags and the comparison choices in Essay~V, (2.3.1)--(2.3.2),
pp.~235--239, identify the corresponding models under compatible
changes of those choices.  A strict point-set representative of
$\beta_B$ is not part of the assertion.  The long exact sequence,
with $\pi_0\operatorname{Map}(B,K(\mathbb Z/2,2))=H^2(B;\mathbb Z/2)$,
now gives \eqref{eq:four-dimensional-stabilization-components}.

\emph{The geometric component class.}
For $(a,\gamma)$, make $\gamma$ constant near its two ends.
The structures $(\gamma_u)_*\alpha$ form a based loop in
$\mathcal S_B$.  Stabilizing and applying the relative classification
gives the class of the pushed-forward cylinder structure.  One can
check the tangent identification explicitly.  Extend $\gamma$
constantly outside $I$.  In the pair-coordinate model of the tangent
microbundle, fix the base image to be $(\gamma_u(x),u)$ and send the
nearby second point $(y,v)$ to
\[
 (\gamma_{u+t(v-u)}(y),v),\qquad 0\leq t\leq1.
\]
For fixed base and parameter this is a homeomorphism germ, with
inverse given by the inverse of the same $\gamma$ slice.  It
preserves the zero section.  At $t=0$ it is the fibrewise tangent
map with a trivial normal line; at $t=1$ it is the tangent map of
the full cylinder homeomorphism.  The homotopy is constant as a
germ near the ends, where $\gamma$ is constant.  Hence the relative
microbundle classes coincide.  Suspension identifies the resulting
$H^3(N,\partial N;\mathbb Z/2)$ class with the stated element of
$H^2(B;\mathbb Z/2)$, proving the component formula.
\end{proof}

The map $\beta_B$ is part of the fibre description and remains
unevaluated in general. The based comparison fibre describes the
component of the identity in the target together with its PL lifts;
the full image of PL concordance components is the next question.


\subsection{The complete component image with its endpoint condition}

Let $B$ be a closed compact real PL four-manifold, and put $N=B\times I$.
Use ordinary joint PL families and ordinary singular TOP families.
Write $A_{\mathcal C}=\operatorname{Aut}_{\mathcal C}(B)$ and
$P_{\mathcal C}=\operatorname{Aut}_{\mathcal C}
(N\ \mathrm{rel}\ B\times\{0\})$, with free top, for
$\mathcal C\in\{\mathrm{PL},\mathrm{TOP}\}$.  Denote the comparisons
by $c_A,c_P$ and top restriction by $e$.  Let $F_A$ be the identity
homotopy fibre of $c_A$ and use the specified map
$\beta_B:F_A\to\operatorname{Map}(B,K(\mathbb Z/2,2))$ of
\eqref{eq:four-dimensional-stabilization-map}.  Set
\[
 I_B=(\pi_0e)^{-1}\bigl(\operatorname{im}(\pi_0c_A)\bigr),
 \qquad J_B=\operatorname{im}(\pi_0\beta_B)
       \subseteq H^2(B;\mathbb Z/2).
\]
The first is a subgroup of $\pi_0P_{\mathrm{TOP}}$; subgroup status
of the second is part of the following assertion.

\begin{theorem}\label{thm:four-dimensional-component-image}
The subset $J_B$ is a subgroup, and there is a well-defined homomorphism
\begin{equation}\label{eq:four-dimensional-component-obstruction}
 \lambda_B:I_B\longrightarrow H^2(B;\mathbb Z/2)/J_B,
 \qquad \ker\lambda_B=\operatorname{im}(\pi_0c_P).
\end{equation}
For $h$ representing a component in $I_B$, choose a TOP isotopy
$\gamma:e(h)\leadsto a$ with $a$ PL, and lift it to a bottom-fixed
isotopy of $h$.  Let $h_\gamma$ be the endpoint.  After imposing
identity and product collars at the bottom and top, respectively,
define
\[
 \kappa(h_\gamma)=[(h_\gamma)_*\alpha_N]-[\alpha_N]
 \in H^3(N,\partial N;\mathbb Z/2)\cong H^2(B;\mathbb Z/2),
\]
where $\alpha_N$ is the standard cylinder structure and the class is
the relative PL-structure difference class.  Then
$\lambda_B([h])=\kappa(h_\gamma)\bmod J_B$, independent of all
these choices.

No surjectivity is asserted.  There is an injection
\begin{equation}\label{eq:four-dimensional-component-quotient}
 I_B/\operatorname{im}(\pi_0c_P)
       \lhook\joinrel\longrightarrow H^2(B;\mathbb Z/2)/J_B.
\end{equation}
The left-hand quotient is therefore a finite elementary abelian
$2$-group of order at most $2^{\dim_{\mathbb F_2}H^2(B;\mathbb Z/2)}$.
In particular, $h$ is TOP-isotopic through bottom-fixed, free-top
concordances to a PL concordance if and only if its top is
TOP-isotopic to a PL map and the displayed relative class vanishes
modulo $J_B$.  If $H^2(B;\mathbb Z/2)=0$, the endpoint condition
alone suffices.  The empty-base case has trivial groups.
\end{theorem}

\begin{proof}
All cohomology in the proof has coefficients $\mathbb Z/2$, and
$FG$ denotes $F\circ G$.  The product-collar normalization in the
preceding proof preserves both boundary restrictions and works for
finite parameter families.  It will be used whenever a relative
structure class is formed.

\emph{Boundary-PL maps and the zero test.}
Let $D_B$ be the group of TOP concordances whose top map is
literally PL.  For $F\in D_B$, product collars make $F_*\alpha_N$
standard near both ends.  Adjoining exterior collars places this
relative structure problem on the open five-manifold $B\times\mathbb R$
with both exterior half-cylinders fixed.  The relative classification
and orbit fibration used above define $\kappa(F)\in H^3(N,\partial N)$
and give the exact test
\begin{equation}\label{eq:relative-zero-test-boundary-pl}
 \kappa(F)=0
 \quad\Longleftrightarrow\quad
 F\text{ is TOP-isotopic with both boundary maps fixed to a PL map}.
\end{equation}
To see the nontrivial direction, a zero class gives a path of relative
structures from $\alpha_N$ to $F_*\alpha_N$.  Lift it through the
relative automorphism orbit map starting at the identity.  If $g$
is its endpoint, then $g_*\alpha_N=F_*\alpha_N$, so $g^{-1}F$ is
PL.  The lifted path gives the required isotopy, fixed near the
boundary.  Conversely, such an isotopy supplies the path of relative
structures.  The same argument, after common collar normalization,
shows that $\kappa$ is constant along paths with fixed PL boundary
maps, including paths joining two collar choices.  The exact
classical inputs are Kirby--Siebenmann, Essay~V, relative (1.5),
p.~229, (3.2), pp.~240--241, and \S5.6, pp.~252--253
\cite{KirbySiebenmann1977}.  The entire boundary is controlled before
they are used.

\emph{The coefficient action.}
Naturality and additivity of differences of PL lifts give the affine
cocycle rule
\[
 \kappa(FG)=\kappa(F)+F_*\kappa(G),\qquad F,G\in D_B.
\]
This is the addition of relative maps into $K(\mathbb Z/2,3)$ in
the standard-lift model of \S5.6.  The action is trivial for every
bottom-fixed concordance.  An explicit verification is useful.
Write $F(x,t)=(F_B(x,t),F_I(x,t))$.  The continuous maps
\[
 K_s(x,t)=\bigl(F_B(x,(1-s)t),\ (1-s)F_I(x,t)+st\bigr),
             \qquad 0\leq s\leq1,
\]
are a homotopy of pairs $(N,\partial N)$ from $F$ to the identity.
At the bottom they fix $(x,0)$; at the top their second coordinate
is $1$; at $s=1$ their first coordinate is $F_B(x,0)=x$.
No assertion that $K_s$ is a homeomorphism is needed.
It follows that $F_*$ is the identity on relative cohomology, and
therefore
\begin{equation}\label{eq:four-dimensional-kappa-additivity}
 \kappa(FG)=\kappa(F)+\kappa(G).
\end{equation}
This proof works without orientability and componentwise on
disconnected $B$.

\emph{Trace classes.}
For a TOP path $\delta_t$ from the identity to a PL map $b$, make
the path constant near its ends and put
\[
 T_\delta(x,t)=(\delta_t(x),t).
\]
It belongs to $D_B$ and is TOP-isotopic to the identity through
bottom-fixed, free-top concordances by
$(x,t)\mapsto(\delta_{st}(x),t)$.  These are continuous TOP
formulas, with no joint PL claim.  The geometric component formula
for $\beta_B$ gives
\begin{equation}\label{eq:four-dimensional-trace-obstruction}
 \kappa(T_\delta)=\pi_0\beta_B([b,\delta]).
\end{equation}
Every member of $J_B$ occurs this way by the path model for $F_A$.
Moreover $T_\delta T_\varepsilon=T_{\delta\varepsilon}$ and
$T_\delta^{-1}=T_{\delta^{-1}}$, where the paths on the right are
pointwise products and inverses.  Their endpoints are again PL.
Additivity proves that the set of their classes, namely $J_B$, is
a subgroup.  This also verifies subgroup status directly without
assuming a group structure compatibility for a chosen equivalence
model of $F_A$.

\emph{Endpoint correction and all choices.}
Let $f=e(h)$ and choose $\gamma:f\leadsto a$, with $a$ PL.
The endpoint lift may be taken explicitly as
\[
 E_{\gamma,s}(x,t)=(\gamma_{st}f^{-1}(x),t),\qquad
 h_\gamma=E_{\gamma,1}h.
\]
It starts at the identity correction, fixes the bottom, and changes
the top of $h$ to $\gamma_s$.  The inverse of the correction is
$(x,t)\mapsto(f\gamma_{st}^{-1}(x),t)$.  Another lift of the same
endpoint path differs by a path of both-ends-fixed maps starting
at the identity.  Its final corrected map is therefore isotopic
to $h_\gamma$ with both boundary maps fixed and has the same
$\kappa$.

For a second endpoint choice $\eta:f\leadsto b$, exact composition
before collar normalization gives
\[
 h_\eta h_\gamma^{-1}(x,t)=(\eta_t\gamma_t^{-1}(x),t).
\]
The path on the right starts at the identity and ends at the PL
map $ba^{-1}$.  By \eqref{eq:four-dimensional-kappa-additivity}
and \eqref{eq:four-dimensional-trace-obstruction}, the difference
$\kappa(h_\eta)-\kappa(h_\gamma)$ belongs to $J_B$.
Collar normalization preserves the classes, so it does not alter
this calculation.  Conversely, postcomposing $h_\gamma$ with any
trace adds the corresponding element of $J_B$ while staying in the
same ordinary free-top component and retaining a PL top.

The coset also depends only on the component of $h$.  It suffices
after endpoint correction to take a path $H_s$ from $F_0$ to $F_1$
in $P_{\mathrm{TOP}}$, with $F_0,F_1\in D_B$.  Put
$a=e(F_0)$, $f_s=e(H_s)$, and
\[
 T_s(x,t)=(f_{st}a^{-1}(x),t),\qquad Q_s=T_s^{-1}H_s.
\]
Every $Q_s$ has the same bottom and top maps, the identity and $a$.
After common collar normalization it is a family of relative
structures, so $\kappa(Q_s)$ is constant.  At $s=0$ it equals
$\kappa(F_0)$.  The map $T_1$ is a trace from the identity to
the PL map $e(F_1)a^{-1}$.  Therefore
\[
 \kappa(F_1)-\kappa(F_0)=\kappa(T_1)\in J_B.
\]
This proves independence of the representative and all the listed
choices.  Given two classes in $I_B$, choose corrected representatives
in $D_B$; their product is still in $D_B$ and represents the product
of the original classes.  Additivity then proves that $\lambda_B$
is a homomorphism.

\emph{The exact kernel.}
A PL concordance has PL top and zero relative class, so its component
lies in $\ker\lambda_B$.  Conversely, if $[h]\in I_B$ has zero
coset, choose $h_\gamma\in D_B$.  Its class lies in $J_B$.
Choose a trace with the opposite class, and postcompose it with
$h_\gamma$.  The result is in the same free-top component as $h$
and has zero $\kappa$.  The test
\eqref{eq:relative-zero-test-boundary-pl} makes it PL by a TOP
isotopy fixing both boundary maps.  This proves the asserted kernel.
A component outside $I_B$ cannot contain a PL map because top
restriction preserves paths.

The first isomorphism theorem gives
\eqref{eq:four-dimensional-component-quotient}.  A compact finite
PL manifold has finite-dimensional mod-$2$ cohomology, so the target
is a finite elementary abelian $2$-group and the stated order bound
follows.  The vanishing-cohomology and empty cases follow as well.
\end{proof}

The endpoint condition and $J_B$ remain unevaluated in general.
The argument uses relative structures in ambient dimension five
only after both boundary maps are PL.  It supplies neither a stable
criterion for a four-dimensional homeomorphism to be TOP-isotopic to a PL map
nor a general surjectivity result for TOP concordance components.


\subsection{Realizing the four-dimensional component obstruction}

We apply the realization theorem of Orson--Powell--Randal-Williams
to the component obstruction of
Theorem~\ref{thm:four-dimensional-component-image}.
Throughout this subsection $B$ is a connected, closed, oriented real
PL four-manifold, $N=B\times I$, and $\pi=\pi_1(B)$.
The spaces $A_{\mathrm{PL}}(B)$ and $P_{\mathrm{PL}}(B)$ use jointly
PL parameter families, while their TOP counterparts use the ordinary
compact-open topology. Concordances fix the bottom pointwise and have
free top. Write $c_A,c_P$ for the comparison maps and $e$ for top
restriction. Recall that
\[
 I_B=(\pi_0e)^{-1}\bigl(\operatorname{im}\pi_0c_A\bigr),
 \qquad
 \lambda_B:I_B\longrightarrow H^2(B;\mathbb Z/2)/J_B
\]
has kernel $\operatorname{im}\pi_0c_P$.
For a concordance whose two boundary maps are PL, its relative
structure difference is denoted by
\[
 \kappa(F)=[F_*\alpha_N]-[\alpha_N]
       \in H^3(N,\partial N;\mathbb Z/2)
       \cong H^2(B;\mathbb Z/2).
\]
Here one first imposes product collars and attaches exterior collars,
and uses the standard tangent lift as origin in the relative
open-five-dimensional PL-structure classification. The subgroup
$J_B$ consists of the classes $\kappa(T_\gamma)$ for TOP paths
$\gamma$ from the identity to a PL automorphism, where
$T_\gamma(x,t)=(\gamma_t(x),t)$.

Let $\delta_*$ be the Bockstein for
$0\to\mathbb Z_{(2)}\xrightarrow{2}\mathbb Z_{(2)}
\to\mathbb Z/2\to0$, and set
\begin{equation}\label{eq:four-dimensional-bockstein-subspace}
 L_B=\operatorname{PD}_B^{-1}
       \ker\bigl(\delta_*:H_2(B;\mathbb Z/2)
                      \longrightarrow H_1(B;\mathbb Z_{(2)})\bigr).
\end{equation}
We use precisely the quadratic assembly summand
\begin{equation}\label{eq:four-dimensional-assembly-summand}
 I_2:H_2(\pi;\mathbb Z_{(2)})
          \longrightarrow L^q_6(\mathbb Z[\pi])_{(2)}
\end{equation}
of \cite[Definition~5.8, p.~21]{OrsonPowellRandalWilliams2026}.
More explicitly, let $\mathcal L^q=\tau_{\geq0}L^q(\mathbb Z)$.
In the Taylor--Williams coordinates
\[
 H_6(B\pi;\mathcal L^q)_{(2)}\cong
 H_0(\pi;\mathbb Z/2)\oplus H_2(\pi;\mathbb Z_{(2)})
 \oplus H_4(\pi;\mathbb Z/2)\oplus H_6(\pi;\mathbb Z_{(2)}),
\]
the map is $I_2(b)=A_\pi\operatorname{TW}^{-1}(0,b,0,0)$,
where $A_\pi$ is quadratic assembly in the simple-surgery convention.

\begin{theorem}\label{thm:four-dimensional-obstruction-realization}
If $I_2=0$, then
\begin{equation}\label{eq:four-dimensional-obstruction-range}
       \operatorname{im}\lambda_B\supseteq(L_B+J_B)/J_B.
\end{equation}
If, in addition, $H_1(B;\mathbb Z_{(2)})$ has no $2$-torsion,
then $\lambda_B$ is surjective and induces an isomorphism
\begin{equation}\label{eq:four-dimensional-realized-quotient}
 I_B/\operatorname{im}(\pi_0c_P)
                 \cong H^2(B;\mathbb Z/2)/J_B.
\end{equation}
\end{theorem}

\begin{proof}
Choose a compatible smoothing of $B$, whose existence in dimension
four is recalled in \cite[Essay~I, p.~39]{KirbySiebenmann1977}.
Let $\sigma$ be its product smoothing on $N$, with underlying PL
structure $\alpha_N$.

\medskip\noindent\textit{Comparison of the relative classes.}
First let $F:N\to N$ fix both ends. Normalize it to be the identity
on common end collars. Put
\[
 W=N\times I,\qquad
 C=(N\times\{0\})\cup(\partial N\times I).
\]
The OPR boundary smoothing is $\sigma$ on $C$ and $F^*\sigma$ on
$N\times\{1\}$. Its relative Kirby--Siebenmann obstruction, followed
by Poincar\'e duality, is the invariant
$\operatorname{KS}_{\mathrm{OPR}}(F)\in H_2(B;\mathbb Z/2)$
\cite[pp.~2, 7--9]{OrsonPowellRandalWilliams2026}.

The universal square for $B\mathrm{PL}\to B\mathrm{TOP}$ identifies
a PL lift with a nullhomotopy of the universal map
$B\mathrm{TOP}\to K(\mathbb Z/2,4)$
\cite[Essay~V, Section~5.6, pp.~252--253]{KirbySiebenmann1977}.
With the standard lift fixed on $W$, changing the lift on the top
face by
\[
 d=[F^*\alpha_N]-[\alpha_N]
       \in H^3(N,\partial N;\mathbb Z/2)
       \cong H^3(\partial W,C;\mathbb Z/2)
\]
changes the relative obstruction by the connecting homomorphism
\begin{equation}\label{eq:relative-ks-interval-suspension}
 H^3(\partial W,C;\mathbb Z/2)
       \xrightarrow{\ \partial\ }
 H^4(W,\partial W;\mathbb Z/2).
\end{equation}
Indeed, the difference of the two boundary nullhomotopies is a
relative map to $K(\mathbb Z/2,3)$; its extension obstruction is
the connecting class in the cohomology sequence of the triple.
The inclusion $C\hookrightarrow W$ is a homotopy equivalence,
as simultaneous contraction in the last interval retracts both
spaces to $N\times\{0\}$. Thus
\eqref{eq:relative-ks-interval-suspension} is an isomorphism,
and is the interval suspension. Forgetting a smooth tangent lift
to a PL lift gives this same universal obstruction, so this
comparison identifies the classes themselves, not merely their
vanishing loci.

Since $F^*\alpha_N=(F^{-1})_*\alpha_N$, we have
$d=\kappa(F^{-1})$. A bottom-fixed concordance acts trivially on
the relative cohomology of $(N,\partial N)$. For completeness, if
$F(x,t)=(a(x,t),b(x,t))$, a homotopy of maps of pairs from $F$ to
the identity is
\[
 (x,t)\longmapsto
 \bigl(a(x,(1-s)t),(1-s)b(x,t)+st\bigr),\qquad 0\leq s\leq1.
\]
This is a homotopy of maps, and need not be an isotopy.
Additivity of the relative structure class therefore gives
$\kappa(F^{-1})=-\kappa(F)=\kappa(F)$.
Combining the two interval suspensions with mod-$2$ duality proves
\begin{equation}\label{eq:opr-relative-pl-identification}
 \operatorname{KS}_{\mathrm{OPR}}(F)
                    =\operatorname{PD}_B\kappa(F).
\end{equation}
The collar normalization is relative-boundary isotopic to the
original map and preserves both relative lift classes.

\medskip\noindent\textit{Realization with the ends fixed.}
Let $Q(B)$ be the OPR group of TOP pseudo-isotopies from the
identity to a diffeomorphism, with the equivalence relation in
\cite[Definition~2.1, p.~5]{OrsonPowellRandalWilliams2026}.
Theorem~2.3(D) of that paper says, under $I_2=0$, that the image of
\[
 \pi_1\widetilde{\operatorname{Homeo}}^+(B)
        \longrightarrow Q(B)
        \xrightarrow{\operatorname{KS}_{\mathrm{OPR}}}
                         H_2(B;\mathbb Z/2)
\]
contains $\ker\delta_*$.
The source is the block homeomorphism object. The assembly hypothesis
is the quadratic map in
\eqref{eq:four-dimensional-assembly-summand}, as specified in
\cite[Definition~5.8 and proof, pp.~21--22]{OrsonPowellRandalWilliams2026}.

Fix $z\in\ker\delta_*$. Choose a realizing element of $Q(B)$ and
represent it by a concordance $F$ from the identity to a
diffeomorphism $f$. Exactness of the block homotopy-fibre sequence
in \cite[equation~(1), p.~6]{OrsonPowellRandalWilliams2026}
implies that $f$ is smoothly pseudo-isotopic to the identity.
Choose such a smooth pseudo-isotopy $G$, with bottom $f$ and top
the identity. Then
\[
       G'=G\circ(f^{-1}\times\operatorname{id}_I),
       \qquad F_0=G'\circ F
\]
has bottom and top equal to the identity. The OPR invariant is a
homomorphism and vanishes on smooth concordances
\cite[Lemma~3.2 and the proof of Lemma~2.2, pp.~7--9]
{OrsonPowellRandalWilliams2026}. Consequently
\begin{equation}\label{eq:both-ended-obstruction-realization}
 \operatorname{KS}_{\mathrm{OPR}}(F_0)=z,
       \qquad \kappa(F_0)=\operatorname{PD}_B^{-1}(z).
\end{equation}
Smoothness of $G'$ is used to make its smoothing obstruction
zero. Its coordinate map need not be PL for the fixed
triangulation; the boundary maps of $F_0$ are literally the
identity, which is the condition needed in
\eqref{eq:opr-relative-pl-identification}.

The ordinary free-top component of $F_0$ lies in $I_B$, and its
image under $\lambda_B$ is
$\operatorname{PD}_B^{-1}(z)+J_B$. This proves
\eqref{eq:four-dimensional-obstruction-range}. If
$H_1(B;\mathbb Z_{(2)})$ has no $2$-torsion, exactness of the
coefficient sequence implies $\delta_*=0$. Hence
$L_B=H^2(B;\mathbb Z/2)$, and the previously established kernel
of $\lambda_B$ proves
\eqref{eq:four-dimensional-realized-quotient}.
\end{proof}

\begin{corollary}\label{cor:simply-connected-four-dimensional-trace}
If $B$ is a connected, closed, simply connected real PL
four-manifold, then
\[
              \pi_0P_{\mathrm{TOP}}(B)=1,
              \qquad J_B=H^2(B;\mathbb Z/2).
\]
For any such $B$, the ordinary trace homomorphism
\begin{equation}\label{eq:ordinary-four-dimensional-loop-obstruction}
 \eta_B:\pi_1A_{\mathrm{TOP}}(B)
       \longrightarrow H^2(B;\mathbb Z/2),
       \qquad [\gamma]\longmapsto\kappa(T_\gamma)
\end{equation}
is surjective and vanishes on the image of
$\pi_1A_{\mathrm{PL}}(B)$. Thus the ordinary automorphism
comparison is not surjective on $\pi_1$ if
$H^2(B;\mathbb Z/2)\ne0$.
\end{corollary}

\begin{proof}
Choose an orientation of $B$. Since $\pi=1$ and
$H_1(B;\mathbb Z_{(2)})=0$, both hypotheses of the theorem hold.
The corrected topological pseudo-isotopy theorem states that
every pseudo-isotopy of a compact simply connected four-manifold
is isotopic to the identity relative to the bottom and vertical
boundary \cite[Theorem~1.5, p.~3; conventions, p.~1]
{GabaiGayHartmanKrushkalPowell2026}. For closed $B$ this gives
$\pi_0P_{\mathrm{TOP}}(B)=1$.

Given $v\in H^2(B;\mathbb Z/2)$, choose the both-end-fixed
$F_0$ in \eqref{eq:both-ended-obstruction-realization} with
$\kappa(F_0)=v$. Take an ordinary path $H_s$ in
$P_{\mathrm{TOP}}(B)$ from the identity to $F_0$. The top maps
$\gamma_s=e(H_s)$ form a based ordinary loop, because both
ends of $F_0$ are the identity. Set
\[
 T_s(x,t)=(\gamma_{st}(x),t),\qquad Q_s=T_s^{-1}\circ H_s.
\]
The path $Q_s$ fixes both boundary components and joins the
identity to $T_\gamma^{-1}F_0$. Normalize the whole compact
parameter family on common end collars. Relative isotopy
invariance and additivity of $\kappa$ give
\[
            0=\kappa(T_\gamma^{-1}F_0)
                  =\kappa(T_\gamma)+\kappa(F_0).
\]
Hence $\eta_B[\gamma]=v$. This passage uses the ordinary
pseudo-isotopy theorem to obtain ordinary loops from the
realization statement.

Homotopies of based loops give relative-boundary isotopies of
their traces, so $\eta_B$ is well defined. Pointwise product of
loops induces product of traces; its operation on $\pi_1$ is
the usual loop operation. Additivity of $\kappa$ therefore
makes $\eta_B$ a homomorphism. A jointly PL loop has a jointly
PL trace, so its relative PL class is zero. Finally ordinary
loops are among the paths defining $J_B$, and surjectivity of
$\eta_B$ gives $J_B=H^2(B;\mathbb Z/2)$.
\end{proof}

\begin{example}\label{ex:four-dimensional-product-sphere-trace}
For $B=S^2\times S^2$, the target of
\eqref{eq:ordinary-four-dimensional-loop-obstruction} is
$(\mathbb Z/2)^2$. A finite PL model is the ordered product of
two tetrahedron boundaries. Its simplices are chains in the
coordinatewise order on $\{0,1,2,3\}^2$ whose two projections
have at most three vertices. These are the compatible ordered
subdivisions of the products of faces, so they triangulate the
product PL manifold itself.

One can count the $k$-simplices by the numbers $r,s$ of vertices
in their two projections. The successive steps increase the
first coordinate, the second coordinate, or both. Their numbers
are respectively $k+1-s$, $k+1-r$, and $r+s-k-2$. Thus
\[
 f_k=\sum_{r,s=1}^{3}\binom4r\binom4s
       \binom{k}{\,k+1-s,\ k+1-r,\ r+s-k-2\,},
 \qquad 0\leq k\leq4,
\]
where a term with a negative lower entry is zero. This gives
$f=(16,84,216,240,96)$. The product CW decomposition has one
$0$-cell, two $2$-cells and one $4$-cell, with zero cellular
differential; hence its integral homology is free of ranks
$(1,0,2,0,1)$. It is simply connected because each factor is.
The corollary therefore applies to this explicit finite base.
\end{example}

For general four-manifolds, the endpoint subgroup $I_B$, the trace
subgroup $J_B$, and the higher concordance groups remain to be determined.


\subsection{Ordinary trace classes with prescribed support}

Let $B$ be a connected, closed, oriented real PL four-manifold.
Keep the ordinary TOP and jointly PL family conventions, and the
subgroups $I_B,J_B$, of
Theorem~\ref{thm:four-dimensional-component-image}.
For a based ordinary loop $\gamma$ in $A_{\mathrm{TOP}}(B)$, set
\[
 T_\gamma(x,t)=(\gamma_t(x),t),\qquad
 \eta_B[\gamma]=\kappa_B(T_\gamma)
                   \in H^2(B;\mathbb Z/2).
\]
Homotopies of based loops give relative-boundary isotopies of
their traces. Additivity of the relative PL difference class,
and multiplication of loops, make $\eta_B$ a homomorphism.
It vanishes on the image of $\pi_1A_{\mathrm{PL}}(B)$, since
a jointly PL loop has a jointly PL trace. These assertions use
no assumption on the fundamental group.

\begin{theorem}\label{thm:supported-four-dimensional-traces}
Let $j:W\hookrightarrow B$ be a compact, connected, simply
connected, codimension-zero PL submanifold with collared boundary.
Allow $\partial W=\varnothing$, in which case $W=B$.
Let
\[
 j_!:H^2(W,\partial W;\mathbb Z/2)
                       \longrightarrow H^2(B;\mathbb Z/2)
\]
be extension by zero. Then
\begin{equation}\label{eq:supported-four-trace-image}
 \operatorname{im}\eta_B\supseteq\operatorname{im}j_!
 =\operatorname{PD}_B^{-1}
   \operatorname{im}\bigl(j_*:H_2(W;\mathbb Z/2)
                                  \to H_2(B;\mathbb Z/2)\bigr).
\end{equation}
Each class in this subgroup has a representing ordinary based
TOP loop which is the identity outside $\operatorname{int}W$
and on one common collar of $\partial W$ for the entire loop.
There is no assembly hypothesis on $\pi_1(B)$.

If $j_*$ on $H_2(-;\mathbb Z/2)$ is surjective, then
\begin{equation}\label{eq:supported-four-component-image}
 \operatorname{im}\eta_B=J_B=H^2(B;\mathbb Z/2),
 \qquad \operatorname{im}(\pi_0c_P)=I_B.
\end{equation}
If this cohomology group is nonzero, the automorphism comparison
is not surjective on $\pi_1$.
\end{theorem}

\begin{proof}
We retain all boundary faces in the realization argument.
Write $N=W\times I$, so that
\[
 \partial N=(W\times\partial I)\cup(\partial W\times I).
\]
Use compatible collars at the corners, or round the corners once.
The relative PL-structure convention is obtained by attaching
exterior collars and working on the resulting open
five-manifold, with the standard structure fixed near a closed
set. This is the relative formulation in
\cite[Essay~V, Section~3, equation~(3.2), p.~240]{KirbySiebenmann1977}.
For $F:N\to N$ fixing $\partial N$ pointwise, the relative
difference from the standard product PL structure $\alpha_N$ is
\begin{equation}\label{eq:relative-four-supported-class}
 \kappa_W(F)=[F_*\alpha_N]-[\alpha_N]
 \in H^3(N,\partial N;\mathbb Z/2)
       \cong H^2(W,\partial W;\mathbb Z/2).
\end{equation}
Common-collar normalization makes this construction invariant
under relative-boundary isotopy. A bottom-fixed homeomorphism
of $N$ acts trivially on its absolute homology: the inclusion
of the bottom is a homotopy equivalence and is fixed by the
homeomorphism. Poincar\'e--Lefschetz duality over $\mathbb Z/2$
then gives trivial action on the relative group in
\eqref{eq:relative-four-supported-class}. Consequently the
relative affine-lift cocycle is additive and
$\kappa_W(F^{-1})=\kappa_W(F)$.

\emph{Relative realization and ordinary loops.}
Choose one compatible smoothing of $W$, using
\cite[Essay~I, p.~39]{KirbySiebenmann1977}.
The Orson--Powell--Randal-Williams conventions allow compact
four-manifolds with boundary and fix the vertical boundary
pointwise throughout
\cite[Section~1.1, p.~1; Sections~2.1--2.2, pp.~4--5]
{OrsonPowellRandalWilliams2026}.
For a boundary-fixed $F$ their relative invariant satisfies
\begin{equation}\label{eq:supported-four-ks-comparison}
 \operatorname{KS}_{\mathrm{OPR}}(F)
             =\operatorname{PD}_W\kappa_W(F)
                         \in H_2(W;\mathbb Z/2).
\end{equation}
Here is the class comparison with the vertical faces included.
Put $Z=N\times I$ and
$C=(N\times\{0\})\cup(\partial N\times I)$.
The standard lift extends over $Z$. The changed top lift gives
a class in
$H^3(\partial Z,C;\mathbb Z/2)\cong H^3(N,\partial N;\mathbb Z/2)$.
The universal Kirby--Siebenmann square identifies its obstruction
with the connecting class in $H^4(Z,\partial Z;\mathbb Z/2)$
\cite[Essay~V, Section~5.6, pp.~252--253]{KirbySiebenmann1977}.
Contraction in the last interval shows $C\hookrightarrow Z$
is a homotopy equivalence, so this connecting map is the
interval-suspension isomorphism. Pulling back the standard
structure gives the class $\kappa_W(F^{-1})$, which equals
$\kappa_W(F)$ by the preceding trivial-action calculation.
The two suspensions followed by duality on $Z$ give precisely
the relative duality in \eqref{eq:supported-four-ks-comparison}.

Simple connectivity of $W$ makes both the assembly domain
$H_2(\pi_1W;\mathbb Z_{(2)})$ and the Bockstein target
$H_1(W;\mathbb Z_{(2)})$ zero. The realization theorem
\cite[Theorem~2.3(D), p.~6; Definition~5.8 and proof,
pp.~21--22]{OrsonPowellRandalWilliams2026}
therefore realizes every $z\in H_2(W;\mathbb Z/2)$ through
the block homotopy-fibre sequence. To obtain a literal
boundary-fixed representative, represent its element of $Q(W)$
by $F$ from the identity to a diffeomorphism $f$.
Exactness says that $f$ is smoothly pseudo-isotopic to the
identity relative to $\partial W$. Choose such a smooth
pseudo-isotopy $G$ with bottom $f$ and top identity.
Then $G'=G\circ(f^{-1}\times\mathrm{id}_I)$ has bottom
identity and top $f^{-1}$, and $F_0=G'F$ fixes every face
of $\partial N$. The OPR invariant is additive and is zero
on $G'$, so $\operatorname{KS}_{\mathrm{OPR}}(F_0)=z$.
By \eqref{eq:supported-four-ks-comparison}, every
$v\in H^2(W,\partial W;\mathbb Z/2)$ is thus
$\kappa_W(F_0)$ for some such $F_0$.

The ordinary topological pseudo-isotopy theorem for compact
simply connected four-manifolds, including those with boundary,
gives a path $H_s$ from the identity to $F_0$ which fixes
$(W\times\{0\})\cup(\partial W\times I)$ throughout
\cite[Theorem~1.5, p.~3; relative convention, p.~1]
{GabaiGayHartmanKrushkalPowell2026}.
Its top restrictions form a based ordinary loop $\gamma_s$
on $W$ fixing $\partial W$ pointwise. Set
\[
 T_s(x,t)=(\gamma_{st}(x),t),\qquad Q_s=T_s^{-1}H_s.
\]
Every $Q_s$ fixes the whole boundary of $N$;
$Q_0=\mathrm{id}$ and $Q_1=T_\gamma^{-1}F_0$.
Relative isotopy invariance, with common collars for this
compact family, gives $\kappa_W(Q_1)=0$. Additivity therefore
gives $\kappa_W(T_\gamma)=v$. This construction gives an ordinary TOP loop.

\emph{A common support collar.}
If the boundary is nonempty, choose a PL collar
$\partial W\times[0,1]\hookrightarrow W$.
For $0\le r\le1/2$ define an embedding $c_r:W\to W$
which is the identity off the collar and on it is
\[
 c_r(y,t)=(y,r+(1-r)t).
\]
Its image is $W$ with the open collar of width $r$ removed.
On this image conjugate $\gamma_s$ by $c_r$, and extend
by the identity over the removed collar; denote the result
by $h_{r,s}$. The two formulas agree on the interface because
$\gamma_s$ fixes $\partial W$. Each $h_{r,s}$ is a
homeomorphism, with inverse given by the same construction
from $\gamma_s^{-1}$, and
$h_{r,0}=h_{r,1}=\mathrm{id}$.

Continuity at $r=0$ can be checked uniformly. Give the compact
manifold a metric for which motion along a collar interval has
its stated length. Both $c_r$ and its inverse on $c_r(W)$
move points by at most $r$. If $\omega$ is a common modulus
of continuity for the compact families $\gamma_s$ and
$\gamma_s^{-1}$, then
\[
 \sup_{s,x}d(h_{r,s}(x),\gamma_s(x))\le r+\omega(r),
 \qquad
 \sup_{s,x}d(h_{r,s}^{-1}(x),\gamma_s^{-1}(x))
                                      \le r+\omega(r).
\]
On the removed collar use the fixed boundary point with the
same collar coordinate $y$ for this estimate. Thus $h$ is
a based ordinary loop homotopy. At $r=1/2$ the loop has
one common identity collar. This construction concerns
ordinary continuous TOP families.

Reparameterize the resulting loop by a PL function
$\rho:I\to I$ which is zero near $0$ and one near $1$.
The straight homotopy from $t$ to $\rho(t)$ is relative
to the endpoints. Hence the relative trace class remains
$v$, while its trace is now the identity near every face
of $\partial N$. When $\partial W$ is empty only this
time reparameterization is needed.

\emph{Extension and its class.}
Extend the loop by the identity to $B$ and call it $\Gamma$.
Gluing on the closed sets $W$ and $B\setminus\operatorname{int}W$
gives a jointly continuous family; the same construction
for inverses verifies that it is an ordinary TOP loop.
The fixed common collar makes the relative PL lifts agree
literally on the gluing region. Naturality of their difference
therefore gives
\begin{equation}\label{eq:supported-four-extension-zero}
 \eta_B[\Gamma]=\kappa_B(T_\Gamma)
                              =j_!\kappa_W(T_\gamma)=j_!v.
\end{equation}
At the cylinder level this is extension by zero from
$H^3(N,\partial N;\mathbb Z/2)$ to
$H^3(B\times I,B\times\partial I;\mathbb Z/2)$;
interval suspension identifies it with $j_!$ in degree two.

Finally,
$\operatorname{PD}_B j_!=j_*\operatorname{PD}_W$.
For example, choose a finite triangulation compatible with
$W$ and its boundary. Extend a relative cochain by zero
and cap with the fundamental chain of $B$. Simplices outside
$W$ contribute zero, and the remaining cap is the inclusion
of the cap with the relative fundamental chain of $W$.
This proves the equality in
\eqref{eq:supported-four-trace-image} and completes the
support assertion.

If $j_*$ is onto, then $\eta_B$ is onto. Every based loop
is among the paths defining $J_B$, so $J_B$ is all of
$H^2(B;\mathbb Z/2)$. The kernel statement in
Theorem~\ref{thm:four-dimensional-component-image} gives
\eqref{eq:supported-four-component-image}. Since $\eta_B$
vanishes on jointly PL loops, its nonzero target also proves
the stated failure of surjectivity on $\pi_1$.
\end{proof}

\begin{corollary}\label{cor:connected-sum-supported-traces}
Let $B=C\#D$, where $C$ is a connected, closed, simply
connected, oriented PL four-manifold and $D$ is a connected,
closed, oriented PL four-manifold with
$H_2(D;\mathbb Z/2)=0$. Then
\eqref{eq:supported-four-component-image} holds.
In particular it holds for $C\#k(S^1\times S^3)$ for every
integer $k\ge0$.
\end{corollary}

\begin{proof}
In the connected-sum decomposition take
$W=C\setminus\operatorname{int}D^4$, with boundary $S^3$.
Van Kampen applied after adding collar neighborhoods to
$C=W\cup D^4$ gives $\pi_1(W)=\pi_1(C)=0$.
Removing the open ball does not change $H_2$.
The Mayer--Vietoris sequence for the two punctured summands
has adjacent groups $H_2(S^3;\mathbb Z/2)=0$ and
$H_1(S^3;\mathbb Z/2)=0$. Hence
\[
 H_2(B;\mathbb Z/2)\cong H_2(W;\mathbb Z/2)
                                 \oplus H_2(D;\mathbb Z/2),
\]
and $j_*$ is onto. The same sequence and the product homology
of $S^1\times S^3$ prove the finite connected-sum case;
$k=0$ is the simply connected case with $W=B$.
\end{proof}

The endpoint subgroup $I_B$ consists of the components whose top
restriction is TOP-isotopic to a PL automorphism. Determining this
subgroup and the remaining concordance groups for the nonsimply
connected examples is a separate problem.


\subsection{The relative cylinder over an aspherical four-manifold}

Asphericity lets us turn a family of cylinder automorphisms into a
relative homotopy problem.  Surgery assembly then detects the normal
invariant of its mapping cylinder.  The resulting comparison separates
the assembly hypotheses needed for the component and fundamental-group
kernels from the additional hypothesis needed for component surjectivity.

Orson--Powell--Randal-Williams use realization of the centre by
automorphism loops for aspherical smooth four-manifolds in their
stabilized smooth pseudo-isotopy construction
\cite[Section~8.1, Remark~8.2 and Example~8.3, pp.~25--26;
Section~8.2, pp.~26--27]{OrsonPowellRandalWilliams2026}.
Here the same central evaluation condition enters the ordinary
PL-to-TOP comparison for the fully boundary-fixed cylinder, through
the relative structure classification and the assembly maps below.

Let $B$ be a closed connected oriented real PL four-manifold
which is aspherical, let $\pi=\pi_1B$, and put $N=B\times I$.
Use the ordinary TOP and joint PL family conventions of
Theorem~\ref{thm:four-dimensional-comparison-fibres}, and write
\[
 H_C=\operatorname{Aut}_C(N\ \mathrm{rel}\ \partial N),
 \qquad C\in\{\mathrm{PL},\mathrm{TOP}\}.
\]
Thus both ends are fixed pointwise.  Homotopy groups are
based at the identity; the TOP groups are replaced by
their ordinary singular complexes for comparison.
Suppose that $\operatorname{Wh}(\pi)=0$, and write
\begin{equation}\label{eq:aspherical-four-assembly}
 A_j:H_j(B;\mathbb L\langle1\rangle)
                   \longrightarrow L_j(\mathbb Z[\pi])
\end{equation}
for quadratic assembly with trivial orientation character
and the decoration of the geometric surgery sequence.
Whitehead-group vanishing identifies the simple and
homotopy decorations needed here.

\begin{theorem}\label{thm:aspherical-four-relative-comparison}
If $A_7$ is injective, the natural relative structure
boundary map identifies, as groups,
\begin{equation}\label{eq:aspherical-four-component-kernel}
 \ker\bigl(\pi_0H_{\mathrm{PL}}\longrightarrow
                  \pi_0H_{\mathrm{TOP}}\bigr)
                 \cong H^1(B;\mathbb Z/2).
\end{equation}
If $A_8$ is also injective, there is an exact sequence
\begin{equation}\label{eq:aspherical-four-fundamental-sequence}
 0\longrightarrow H^0(B;\mathbb Z/2)
 \longrightarrow\pi_1H_{\mathrm{PL}}
 \longrightarrow\pi_1H_{\mathrm{TOP}}
 \longrightarrow0,
\end{equation}
and the comparison is an isomorphism on $\pi_j$ for
every $j\geq2$.

If, in addition to $A_7$, the map $A_6$ is injective and
every element of $Z(\pi)$ is the evaluation class of a
based joint PL loop of automorphisms of $B$, then
\begin{equation}\label{eq:aspherical-four-component-sequence}
 0\longrightarrow H^1(B;\mathbb Z/2)
 \longrightarrow\pi_0H_{\mathrm{PL}}
 \longrightarrow\pi_0H_{\mathrm{TOP}}
 \longrightarrow0
\end{equation}
is exact as a sequence of groups.  Neither exact sequence
is asserted to split.

Let $\eta_B$ denote the ordinary-loop contribution to
\eqref{eq:four-dimensional-stabilization-components}:
\begin{equation}\label{eq:aspherical-four-ordinary-loop-map}
 \eta_B:\pi_1\operatorname{Aut}_{\mathrm{TOP}}(B)
 \longrightarrow\pi_0F_A
 \xrightarrow{\pi_0\beta_B}H^2(B;\mathbb Z/2).
\end{equation}
Under $A_6$-injectivity, its value on a loop is zero
whenever that loop's evaluation class is represented
by a based joint PL loop.  In particular the preceding
central evaluation hypothesis implies $\eta_B=0$.
That hypothesis is automatic if $Z(\pi)=0$.
\end{theorem}

\begin{proof}
All ordinary cohomology groups in this proof have
coefficients $\mathbb Z/2$.

\emph{A relative mapping-cylinder calculation.}
For $r\geq1$ put $V=B\times I^r$.  Suppose that a TOP
homeomorphism $F:V\to V$ fixes $\partial V$ and is
homotopic to the identity relative to that boundary.
Its stable PL reduction difference is
\[
 \kappa(F)\in H^3(V,\partial V)\cong H^{3-r}(B).
\]
We claim that
\begin{equation}\label{eq:aspherical-four-relative-zero-lemma}
 A_{5+r}\text{ injective}\quad\Longrightarrow\quad
 \kappa(F)=0.
\end{equation}

Choose a PL collar of the entire boundary of the
underlying PL manifold $V$, including the product corners
\cite[Theorem~2.25 and Corollary~2.26,
pp.~24--26]{RourkeSanderson1972}.
On its interval coordinate $0\leq t\leq\delta$, use
\[
 \sigma_s(t)=\min\{t+s,(t+\delta)/2\},
 \qquad 0\leq s\leq\delta/2,
\]
and use the identity off the collar.  This defines an
embedding $c_s:V\to V$ omitting the outer collar of
width $s$, with inverse collar coordinate
$\max\{t-s,2t-\delta\}$.  Conjugate $F$ by $c_s$ on
its image and extend by the identity on the omitted
collar.  The formulas agree at their common boundary;
the identical construction for $F^{-1}$ gives the inverse.
They define an isotopy relative to $\partial V$.
For a given relative homotopy of maps $k_u$, the formula
$c_sk_uc_s^{-1}$, likewise extended by the identity,
is a homotopy of maps with a common fixed collar.
Continuity follows by pasting on the two closed regions
in $(s,x,u)$; the displayed inverse collar coordinate
is continuous on the first region, including $s=0$.
This construction does not require $k_u$ to be invertible.

After this compression and an end-constant reparametrization,
take $k_0=1$ and $k_1=F$, both fixed on that collar.
The labeled mapping cylinder and its target are
\[
 M_F=(V\times I)\cup_FV,\qquad W=V\times I,
 \qquad U([x,u])=(k_u(x),u).
\]
At the upper face the label is $y=F(x)$, so $U$ is
the identity in that label.  It is also the identity
on the bottom and side faces.  The end-constant
homotopy and the fixed collar make these identifications
products near every face and corner intersection.
The bottom inclusions are homotopy equivalences and
$U$ restricts to their common identity.  Hence $U$ is
a homotopy equivalence of the full boundary pairs.

Its source stable TOP normal reduction, transported
with the Spivak normal-fibration comparison and the
specified boundary data, gives a relative normal
invariant satisfying
\begin{equation}\label{eq:aspherical-four-normal-invariant}
 \nu(U)\in[W,\partial W;G/\mathrm{TOP},*]
       \cong H_{5+r}(B;\mathbb L\langle1\rangle),
 \qquad \theta\nu(U)=0.
\end{equation}
The dimension here is $5+r\geq6$.  The compact relative
normal-invariant sequence and its identification with
assembly are precisely
\cite[Definition~18.2, pp.~197--198; relative-version
paragraph on p.~198; Proposition~18.3 and Theorem~18.5,
pp.~198--201]{Ranicki1992TopologicalManifolds}.
Thus injectivity of $A_{5+r}$ gives $\nu(U)=0$.

This zero class includes the fixed boundary comparison.
Forgetting the spherical-fibration comparison gives a
relative homotopy of the source TOP normal reduction
to the pullback of the target's standard reduction.
The latter has a PL lift.  Pulling this lift along the
homotopy gives a PL normal reduction of $M_F$ agreeing
with all its prescribed boundary reductions.  Its
relative Kirby--Siebenmann obstruction is consequently
zero.  The tangent and normal obstructions agree by
the Whitney-sum identity and stable triviality of their
sum.  The mapping-cylinder definition identifies this
class with the interval suspension of $\kappa(F)$.
Since
\[
 H^3(V,\partial V)\xrightarrow{\ \cong\ }
 H^4(V\times I,\partial(V\times I))
\]
is an isomorphism, the claim
\eqref{eq:aspherical-four-relative-zero-lemma} follows.

\emph{Relative homotopies for at least two interval factors.}
If $r\geq2$, the boundary $B\times\partial I^r$ is
connected and its fundamental group maps onto that
of $V$.  A homeomorphism fixed there induces the
identity on $\pi_1V$.  The full preimage of the boundary
in the universal cover is connected.  The lift fixing
one boundary point therefore fixes this entire preimage.
The universal cover of $V$ is contractible by asphericity.
A relative cellular construction, selecting one cell
in each deck orbit, gives a deck-equivariant homotopy
from that lift to the identity.  All extension
obstructions vanish in the contractible target.
It descends to a homotopy relative to the boundary.
The same argument fixes a closed collar whenever the
original map fixes that collar.

\emph{The family maps in the relative comparison.}
Let $\mathcal S_N$ be the sliced relative PL structure
space from the proof of
Theorem~\ref{thm:four-dimensional-comparison-fibres}.
Its construction uses the exterior-collared manifold
$B\times\mathbb R$, fixed outside the compact strip;
the classical orbit fibration therefore has an open
five-dimensional ambient manifold.  The relevant
identification is
\begin{equation}\label{eq:aspherical-four-structure-space}
 \mathcal S_N\simeq
 \operatorname{Map}_*(N/\partial N,K(\mathbb Z/2,3)),
 \qquad \pi_q\mathcal S_N=H^{2-q}(B).
\end{equation}
The orbit quotient contains the entire component of
the standard structure.  The orbit map sends a
homeomorphism to its transported standard structure.

Represent an element of $\pi_qH_{\mathrm{TOP}}$, for
$q=1,2$, by a based cube family $h_p$, $p\in I^q$.
The collar compression in the proof of the cited
comparison theorem fixes one common spatial collar
for the family.  Precomposing each parameter coordinate
by a PL function constant near its endpoints also
makes the family the identity on a parameter-boundary
collar.  Its trace
\begin{equation}\label{eq:aspherical-four-family-trace}
 \mathcal T_h:N\times I^q\longrightarrow N\times I^q,
 \qquad (x,p)\longmapsto(h_p(x),p)
\end{equation}
is an actual TOP homeomorphism fixed near the full
boundary.  Its inverse uses $h_p^{-1}$ in the same slice.

The image of the family under the orbit map is its
trace's stable PL difference in
\[
 H^3(N\times I^q,\partial(N\times I^q))
                 \cong H^{2-q}(B).
\]
Here the tangent data of an ordinary family can be
checked directly.  For a base point $(x,p)$ mapped
to $(h_p(x),p)$, send a nearby second point $(y,p')$ to
\[
 \bigl(h_{p+t(p'-p)}(y),p'\bigr),\qquad 0\leq t\leq1.
\]
Extend the family constantly outside the parameter
cube using its identity collar.  For every fixed base
point this is a homeomorphism germ, with inverse using
the inverse of the indicated slice.  It preserves
the zero section and the fixed boundary germs.
At $t=0$ it is the vertical tangent map with trivial
parameter directions; at $t=1$ it is the full tangent
map of the trace.  Their relative PL obstructions are
therefore equal.  This verifies the family-to-trace
identification with the specified ordinary TOP convention.

For the trace, $r=q+1\geq2$.  The preceding relative
homotopy argument and
\eqref{eq:aspherical-four-relative-zero-lemma} show that
\[
 \pi_1H_{\mathrm{TOP}}\longrightarrow H^1(B)
 \quad\text{vanishes under }A_7\text{-injectivity},
\]
and that
\[
 \pi_2H_{\mathrm{TOP}}\longrightarrow H^0(B)
 \quad\text{vanishes under }A_8\text{-injectivity}.
\]

\emph{Exactness and the group law.}
The comparison fibre is
$F_H\simeq\operatorname{Map}(B,K(\mathbb Z/2,1))$.
Consequently its exact sequence contains
\begin{equation}\label{eq:aspherical-four-relative-long-sequence}
\begin{aligned}
0&\longrightarrow\pi_2H_{\mathrm{PL}}
 \longrightarrow\pi_2H_{\mathrm{TOP}}
 \longrightarrow H^0(B)
 \longrightarrow\pi_1H_{\mathrm{PL}}\\
 &\longrightarrow\pi_1H_{\mathrm{TOP}}
 \longrightarrow H^1(B)
 \longrightarrow\pi_0H_{\mathrm{PL}}
 \longrightarrow\pi_0H_{\mathrm{TOP}}.
\end{aligned}
\end{equation}
The two zero maps give
\eqref{eq:aspherical-four-component-kernel},
\eqref{eq:aspherical-four-fundamental-sequence}, and
the degree-two isomorphism.  Since the fibre's higher
homotopy groups vanish, the higher isomorphisms follow
from the same sequence.

We verify the asserted group structure at the component
kernel.  The path-model fibre consists of a PL map $a$
and a TOP path from the identity to $a$.  Pointwise
composition of the paths makes its map to
$\pi_0H_{\mathrm{PL}}$ a group homomorphism.
The identification of its components with $H^1(B)$
uses the relative PL difference of those path traces,
whose boundary maps are PL.  The difference of a
composite is the sum of the differences with the usual
pullback action.

That action is trivial.  Every element of $H_C$ is
homotopic to the identity as an unrestricted map of
$N$, since it fixes the bottom.  Moreover the exact
sequence of the pair gives naturally
\[
 H^j(N,\partial N)\cong
 \operatorname{coker}\bigl(H^{j-1}(B)
       \xrightarrow{\Delta}H^{j-1}(B)^{\oplus2}\bigr).
\]
Both boundary restrictions are identities, so the
induced relative cohomology action is the identity.
For a path trace the same assertion follows after
the interval suspension: it preserves the parameter,
and acts identically on the relative cohomology of
every slice.  Equivalently, the relative product
filtration has only the single interval column.
The trace difference of a product therefore adds in
ordinary cohomology.  The identification with $H^1(B)$
and its injective boundary map in
\eqref{eq:aspherical-four-component-kernel} are group
homomorphisms, as claimed.

\emph{Components and central evaluation.}
For the single interval the boundary has two components.
If $F(x,t)=(a_t(x),v(x,t))$ fixes both ends, projecting
onto $B$ gives a loop $a_t$ in
$\operatorname{Map}(B,B)_{\mathrm{id}}$.  Each $a_t$
is homotopic to the identity and hence a homotopy
equivalence.  For an aspherical connected $B$,
the classical evaluation calculation gives
\cite[Lemma~4.3, p.~3336]{EnkelmannEtAl2018Atheory}
\begin{equation}\label{eq:aspherical-four-central-evaluation}
 \pi_1\operatorname{Map}(B,B)_{\mathrm{id}}\cong Z(\pi).
\end{equation}
Indeed a loop is a map $B\times S^1\to B$ restricting
to the identity on $B\times\{1\}$.  Its fundamental-group
homomorphism sends the extra circle generator to a
central element.  Conversely any such homomorphism
extends that identity and is realized uniquely up to
relative homotopy by the cellular $K(\pi,1)$ construction.
Loop concatenation corresponds to multiplication in
the centre.

Choose a joint PL automorphism loop $b_t$ representing
the same central element.  Its trace $T_b(x,t)=(b_t(x),t)$
is PL and fixes both ends.  The composition $T_b^{-1}F$
has first coordinate $b_{v(x,t)}^{-1}(a_t(x))$.
Interpolating $v(x,t)$ to $t$, both in this subscript
and in the output interval coordinate, is a homotopy
of maps relative to the two ends.  Its result is the
trace of the loop of maps $b_t^{-1}a_t$.  This loop
has zero central evaluation class by
\eqref{eq:aspherical-four-central-evaluation}, so is
nullhomotopic as a based loop of maps.  Thus
$T_b^{-1}F$ is homotopic to the identity relative to
both ends.  The relative zero lemma with $r=1$ and
$A_6$-injectivity gives $\kappa(T_b^{-1}F)=0$.
Since $T_b$ is PL and the relative cohomology action
is trivial, $\kappa(F)=0$.

In the five-dimensional relative structure space
\eqref{eq:aspherical-four-structure-space}, a zero
difference means membership in the standard component.
That entire component lies in the orbit quotient.
Its Kan fibration therefore lifts a structure path
to a TOP isotopy, with the prescribed fixed collar,
ending at a PL homeomorphism.  The preliminary collar
compression preserves the original relative component.
This proves component surjectivity and
\eqref{eq:aspherical-four-component-sequence}.

Finally, the trace identification in
Theorem~\ref{thm:four-dimensional-stabilization-fibre}
says that $\eta_B$ on an ordinary loop is the PL
difference of its trace.  Removing a joint PL loop
with the same evaluation class and applying the same
$r=1$ argument makes this difference zero.  Only
$A_6$-injectivity is needed for this assertion about
the specified loop.  In particular it vanishes on
the kernel of evaluation.  This proves the theorem.
\end{proof}

\begin{corollary}\label{cor:aspherical-four-explicit-groups}
All the hypotheses of
Theorem~\ref{thm:aspherical-four-relative-comparison}
hold for the standard $T^4$.  Its relative component
comparison is surjective with kernel $(\mathbb Z/2)^4$;
its fundamental-group comparison is surjective with
kernel $\mathbb Z/2$; its comparisons in degrees at
least two are isomorphisms; and $\eta_{T^4}=0$.

More generally, the Whitehead and assembly hypotheses
hold for every $B$ as above whose fundamental group
is word-hyperbolic or acts properly, cocompactly and
isometrically on a finite-dimensional CAT(0) space.
Thus the relative kernel, fundamental-group sequence
and higher isomorphisms hold throughout this class.
The component-surjectivity conclusion additionally
requires the stated realization of central elements
by joint PL loops, and holds in particular for
trivial centre.
\end{corollary}

\begin{proof}
For $T^4$, the free-abelian calculation
\cite[Example~24.16, p.~278]{Ranicki1992TopologicalManifolds}
gives $\operatorname{Wh}(\mathbb Z^4)=0$ and the
zero-connective quadratic assembly isomorphism in
degrees at least four.  The cofiber of
$\mathbb L\langle1\rangle\to\mathbb L\langle0\rangle$
is $H\mathbb Z$.  For $j=6,7,8$ the groups
$H_{j+1}(T^4;\mathbb Z)$ and $H_j(T^4;\mathbb Z)$
vanish, so the one-connective comparison in those
degrees is an isomorphism.  The translation loop
$b_t([x])=[x+tz]$ for $z\in\mathbb Z^4$ is jointly
PL and evaluates to $z$.  To see the joint assertion, work on
$[0,1]^4\times I$ and cut it by the finitely many affine hyperplanes
$x_i+t z_i=m$, with $m\in\mathbb Z$ meeting that compact domain.
On each resulting polytope the map modulo the lattice is affine.
A common triangulation, compatible with the quotient faces, makes
its trace jointly PL; the inverse uses $-z$ and the same construction.
The standard product
calculation gives $H^1(T^4;\mathbb Z/2)=(\mathbb Z/2)^4$
and $H^0(T^4;\mathbb Z/2)=\mathbb Z/2$.

For the general group class, first observe that $\pi$
is torsion-free.  The contractible universal cover
of the finite aspherical four-complex $B$ provides a
finite-length free $\mathbb Z\pi$ resolution of
$\mathbb Z$.  Restriction to a cyclic subgroup of
prime order would give a finite-length free resolution
over its group ring, contradicting its nonzero
cohomology in arbitrarily high degrees.  Every
nontrivial torsion element has a power of prime order.

Bartels--L\"uck's class contains both types of groups
required here
\cite[definition on p.~632]{BartelsLuck2012Borel}.
Their Theorem~B supplies coefficient K-theory
assembly isomorphisms in degrees at most zero,
surjectivity in degree one, and ultimate L-theory
assembly in every degree.  Proposition~0.3(i) and
its proof explicitly give, for torsion-free groups,
\begin{equation}\label{eq:aspherical-four-simple-assembly}
 \operatorname{Wh}(\pi)=0,\qquad
 H_j(B\pi;\mathbb L^s_{\mathbb Z})
       \xrightarrow{\ \cong\ }L_j^s(\mathbb Z\pi)
 \quad(j\in\mathbb Z)
\end{equation}
\cite[Theorem~B, pp.~633--634; Proposition~0.3 and
proof, pp.~634--635]{BartelsLuck2012Borel}.
The proof also gives vanishing of the reduced
projective class group and the negative K-groups.
Those are the specified inputs for the decoration
comparison and passage from virtually cyclic to
ordinary $B\pi$ assembly.  We use the trivial
orientation character and the algebraic conclusion
(i) of that proposition.

Finally, the cofiber $D$ of
$\mathbb L\langle1\rangle\to\mathbb L^s_{\mathbb Z}$
has no positive homotopy groups.  Filter the finite
four-complex $B$ by its cells.  A cell of dimension
$p\leq4$ contributes a suspension $\Sigma^pD$;
induction on the cells gives $H_j(B;D)=0$ for $j>4$.
The cofiber exact sequence, with outside terms
$H_{j+1}(B;D)$ and $H_j(B;D)$, therefore makes the
one-connective comparison an isomorphism for
$j=6,7,8$.  Equation~\eqref{eq:aspherical-four-simple-assembly}
proves all three required assembly isomorphisms.
\end{proof}

For a finite model of the torus, take the staircase
triangulation of the product of four three-edge
polygonal circles.  It has
\[
 (f_0,f_1,f_2,f_3,f_4)=(81,1215,4050,4860,1944).
\]
Indeed a $k$-simplex with minimal cubical carrier of
dimension $p$ corresponds to an ordered partition
of the $p$ varying coordinates into $k$ nonempty
blocks.  Thus, for $k\geq1$,
\[
 f_k=3^4\sum_{p=k}^4\binom4p\,k!\,S(p,k),
\]
where $S(p,k)$ is a Stirling number of the second kind.
The integral cellular chain complex is the tensor
product of four circle complexes, each chain homotopy
equivalent to one copy of $\mathbb Z$ in degrees zero
and one with zero differential.  Hence its integral
homology is free of ranks $(1,4,6,4,1)$.
This finite calculation describes the product model;
the assembly and automorphism comparisons follow from
the preceding geometric and algebraic arguments.

The conclusions evaluate the fully boundary-fixed
cylinder comparison and the ordinary-loop part of
$\beta_B$.  The remaining values of $\beta_B$, the
free-top comparison fibre, and the endpoint-image
condition require separate information.


\section{Mapping cylinders and attaching maps}
Mapping cylinders describe neighbourhoods of a closed core while
keeping its outer face fixed.

\subsection{Mapping cylinders around a definable zero locus}
\label{subsec:zero-locus-mapping-cylinders}
Mapping-cylinder neighborhoods of closed semialgebraic sets over real
closed fields are classical.  Delfs and Knebusch prove their existence
in locally complete ambient spaces
\cite[Theorem~2, p.~48; \S3, pp.~63--65]{DelfsKnebusch1984Retractions};
their Proposition~4.1, p.~66, also gives relative homotopy uniqueness
of extensions.  The construction below records the compact definable
version from scalar triangulation, with the specified outer face,
the continuation of a prescribed strict tube, and the empty cases
made explicit.

Let $K$ be a real closed field with a fixed o-minimal expansion.
Compactness in this subsection means definable compactness, equivalently
closedness and boundedness for the affine definable sets used here.
For a continuous definable map $\varphi:B\to Y$, put
\[
 M_\varphi=(B\times[0,1]\amalg Y)/((b,0)\sim\varphi(b)).
\]
An \emph{admissible mapping-cylinder description} of a pair $(V,B)$
with core $Y$ identifies $M_\varphi$ with $V$, is the identity on $Y$,
and identifies $(b,1)$ with the specified point $b$ of the outer face.
The last convention fixes the domain of the attaching map.

\begin{theorem}[Mapping-cylinder descriptions of strict tubes]
\label{thm:zero-locus-mapping-cylinder}
Let $X\subset K^n$ be closed and bounded, and let
$f:X\to K_{\ge0}$ be continuous and definable.  Put $Y=f^{-1}(0)$.
For all sufficiently small $\epsilon>0$, the sublevel set
$V_\epsilon=f^{-1}([0,\epsilon])$ admits an admissible
mapping-cylinder description with outer face
$B_\epsilon=f^{-1}(\epsilon)$.  It is a neighborhood of $Y$ in
$X$, and the displayed product structure off $Y$ extends across
$B_\epsilon$ to an exterior collar.  Its pointed quotient by $Y$
is a rigid cone.

More generally, let $V\subset X$ be a closed strict tube about $Y$ whose
pointed quotient by $Y$ has a specified definable rigid-cone
description with base $B$.  Then $(V,B)$ admits an admissible
mapping-cylinder description, possibly different from the initial
cone description.  For a fixed $(V,B,Y)$, all attaching maps
$B\to Y$ obtained from admissible descriptions are definably homotopic.
\end{theorem}
In the second assertion, the quotient-cone data give a continuous
definable radius $\rho:V\to[0,1]$, with $\rho^{-1}(0)=Y$, and a
specified definable homeomorphism
\[
 e:B\times(0,1]\longrightarrow V\setminus Y,
 \qquad \rho(e(b,t))=t,\qquad e(b,1)=b.
\]
The strict tube also includes an exterior collar of its specified
outer face. Since $V$ is closed in $X$, it is closed and bounded.
These are the data used in the construction below.
\begin{proof}
The triangulation theorem for a continuous definable scalar
function on a compact definable set
\cite[Theorem~2.2(1), p.~6]{Shiota2010} gives a finite polyhedron
$P$, a definable homeomorphism $h:P\to X$, and a nonnegative
piecewise affine function $g=f\circ h$.  Choose a triangulation
on which $g$ is affine on every simplex.  Its zero locus $Z$ is
the subcomplex consisting of faces spanned by zero-valued vertices.
If any positive vertex values occur, choose $\epsilon>0$ with
$2\epsilon$ smaller than all of them.

On a simplex having both zero and positive vertices, write the
barycentric coordinates of $x$ as $(\lambda_v)_v$, and put
\[
 a(x)=\sum_{g(v)=0}\lambda_v,
 \qquad
 y(x)=\frac{1}{a(x)}\sum_{g(v)=0}\lambda_vv.
\]
Whenever $0<g(x)\le2\epsilon$, both the zero and positive
coordinate sums are nonzero.  Indeed
\[
 g(x)=\sum_{g(v)>0}\lambda_vg(v)
   >2\epsilon\sum_{g(v)>0}\lambda_v,
\]
so the latter sum is less than one.  Simplices with only positive
vertices do not meet this sublevel set.  Thus $y(x)$ is defined
there and is continuous across simplex faces: coordinates of
vertices outside a common face are zero, leaving the same formula.

Let $L=g^{-1}(\epsilon)$ and set $\varphi(b)=y(b)$ for $b\in L$.
For $0\le t\le1$ define
\[
 \Psi(b,t)=\varphi(b)+t\bigl(b-\varphi(b)\bigr).
\]
This lies in the simplex containing $b$ and satisfies
$g(\Psi(b,t))=t\epsilon$.  Together with the identity on $Z$,
it defines a continuous definable map $M_\varphi\to g^{-1}([0,\epsilon])$.
For a point $x$ with $r=g(x)>0$, the inverse follows the same
segment to height $\epsilon$:
\[
 b=y(x)+\frac{\epsilon}{r}\bigl(x-y(x)\bigr),
 \qquad t=r/\epsilon.
\]
The total positive barycentric weight of $b$ is
$\epsilon/g(z)<1$, where $z$ is the normalized positive-face
component of $x$; hence $b$ remains in that simplex.  The normalized
zero-face component stays equal to $y(x)$, so these formulas are
mutual inverses off $Z$.  On $Z$ the mapping-cylinder equivalence
relation gives exactly the identity.  A continuous definable
bijection from a closed bounded definable model of this mapping
cylinder to the Hausdorff set $g^{-1}([0,\epsilon])$ is a
homeomorphism.  For example, a closed bounded model is supplied
by $(tb,\varphi(b),t)$ together with $(0,y,0)$ for $y\in Z$.

The same segments continue to height $2\epsilon$: the positive
vertex values are larger than that height, so the same inverse
formula applies.  Consequently
\[
 L\times(0,2\epsilon]\longrightarrow g^{-1}((0,2\epsilon]),
 \qquad (b,r)\longmapsto
 \varphi(b)+(r/\epsilon)(b-\varphi(b))
\]
is a definable product description.  It provides both the inner
collar and the exterior collar at height $\epsilon$.
The set $g^{-1}([0,\epsilon))$ is open in $P$ and contains $Z$,
so this is a neighborhood of the core.  Transport by $h$ proves
the first assertions.  Collapsing the core of a mapping cylinder
gives the cone on its outer face.

The exceptional cases are included.  Components contained in $Z$
need no outer-face preimage.  If $Y=X$, take $B_\epsilon=\varnothing$
and the mapping cylinder $Y$.  If $Y=\varnothing$, then either
$X=\varnothing$ or $f$ has a positive minimum; choose $\epsilon$
below that minimum and the neighborhood and outer face are empty.
For quotients use the pointed convention $X/\varnothing=X\amalg\{o\}$,
so the empty-neighborhood quotient is the cone on the empty set.

For the prescribed strict tube, its cone radius pulls back to a
continuous definable function $\rho:V\to[0,1]$ with zero set $Y$.
The given rigid cone identifies $V\setminus Y$ with
$B\times(0,1]$, preserving $\rho$.  Apply the construction just
proved to obtain an admissible mapping cylinder below a small
height $\eta>0$.  Let $j:B\to\rho^{-1}(\eta)$ be the restriction
of the given positive-radius product.  Transport the new attaching
map along $j$.  Below height $\eta$, use the new mapping cylinder
with radial coordinate rescaled by $\eta$; above height $\eta$,
use the original product.  Both formulas equal $j(b)$ on the
common level.  They glue to a continuous definable bijection
$M_\varphi\to V$ that is the identity on the core and outer face,
and definable compactness again makes it a homeomorphism.
This proves the assertion without requiring the original cone
rays to extend continuously to individual points of $Y$.

Finally, two admissible descriptions give strong definable
deformation retractions $H_i:V\times[0,1]\to V$ onto $Y$,
with final retractions $r_i:V\to Y$.  The homotopy
\[
 (x,s)\longmapsto r_0(H_1(x,s))
\]
runs from $r_0$ to $r_1$ and fixes $Y$ pointwise, since
$r_0|_Y$ is the identity.  Its restriction to the fixed outer
face $B$ joins the two attaching maps.  This proves independence
of their homotopy class; it makes no assertion that the space of
descriptions is contractible.
\end{proof}


\subsection{Boundary bundles and mapping-cylinder parametrizations}
\label{subsec:boundary-bundle-rulings}
All spaces and maps in this subsection are definable in the fixed
o-minimal expansion of a real closed field $K$, and the objects are
closed and bounded.  A bordant pair $(E,B)$ has a collared boundary
subset $B$.  A boundary-bundle object consists of $(E,B)$, a core
space $Y$, and a specified locally trivial definable bundle
$\varphi:B\to Y$.  Its morphisms are homeomorphisms of the pairs
and of the cores whose boundary restrictions satisfy the bundle
square.  The unrigidified collapse identifies each boundary fibre
with its point in $Y$ and remembers only the resulting pair $(X,Y)$.

Collapsing a fibre can lose the limiting directions needed to lift a
homeomorphism back to the boundary. Specifying the product
parametrization makes the extension of a boundary-compatible map
unique; cutting off the open inner tube then gives the inverse
construction.

\begin{proposition}[Boundary collapse and product parametrizations]
\label{prop:boundary-bundle-retained-equivalence}
The unrigidified collapse functor is not full as an ordinary functor,
already for a one-point core.  In the stipulated definable simplicial
categories it also fails to be a Dwyer--Kan equivalence.

For the corrected target, an object consists of closed bounded definable sets
$(X,Y,V,B)$ with $Y\subset V\subset X$ and
$B\subset V\setminus Y$, a specified locally trivial definable bundle
$\varphi:B\to Y$, and a specified homeomorphism
$\eta:M_\varphi\to V$ that is the identity on the core $Y$ and
identifies $B\times\{1\}$ with the outer face $B$.
Require that $V\setminus B$ be open in $X$ and that $B$ have an
exterior collar in $X\setminus(V\setminus B)$.
A target morphism preserves these marked sets and the full
presentation: on the tube it has the form
\[
 [b,t]\longmapsto[h_B(b),t],
 \qquad h_Y\varphi=\varphi'h_B,
 \qquad 0<t\leq1.
\]
Here $h_B$ is its restriction to the outer face and $h_Y$ its
restriction to the core.
With these object data and morphisms, attachment
\[
 (E,B,Y,\varphi)\longmapsto E\cup_B M_\varphi
\]
is an equivalence and induces isomorphisms of the definable
simplicial mapping spaces.
\end{proposition}
\begin{proof}
For ordinary fullness take $Q=[-1,1]^2$, $B=\partial Q$, and
the cylinder $B\times[0,1]$ with its inner boundary collapsed to
one point.  The quotient is $Q$, with quotient map
\[
 q(b,t)=tb.
\]
The semialgebraic homeomorphism
$h:Q\to Q$, $h(x,y)=(x,y^3)$, fixes the distinguished point.
Any lift to a cylinder homeomorphism is forced away from the
inner boundary, because $q$ is a homeomorphism there.
On the right edge, for $b=(1,s)$ and $t>0$, the forced lift is
\[
 ((1,s),t)\longmapsto((1,t^2s^3),t).
\]
Indeed the maximum norm of $(t,t^3s^3)$ is $t$.
Continuity at $t=0$ would send every point of that inner edge to
the same point $((1,0),0)$, which excludes a homeomorphism.
The argument uses polynomial identities and order inequalities,
so it works over every $K$ as stated.

The square homeomorphism is itself joined to the identity by
\[
 h_u(x,y)=\bigl(x,(1-u)y+uy^3\bigr),\qquad0\le u\le1.
\]
Each map in this definable family is a homeomorphism: its second
coordinate is strictly increasing from $-1$ to $1$.
For the enriched conclusion we use the six-torus obstruction of
Theorem~\ref{result-5}.  Here the boundary bundle over a point is
automatically locally trivial, and equisingularity along that core
places no additional restriction.  To recall the precise mapping
space comparison, put $T=(S^1)^6$ and
\[
 \begin{split}
 G&=\operatorname{Homeo}_{\rm def}
       (T\times I;T\times\{0\}\text{ setwise}),\\
 H&=\operatorname{Homeo}_{\rm def}(CT;o).
 \end{split}
\]
Restriction to the outer boundary maps both groups to
$\operatorname{Homeo}_{\rm def}(T)$ and has sections given by
constant extension and by coning.  These sections identify the
underlying simplicial spaces with the product of the restriction
kernel and $\operatorname{Homeo}_{\rm def}(T)$.  The kernel for
$G$ is the group fixing the upper boundary, which becomes
$P_{\rm def}(T)$ after interval reversal.  The kernel for $H$
fixes the base and vertex and contracts by
Theorem~\ref{result-1}.  Boundary collapse commutes with these
restriction maps.  Its homotopy fibre is therefore
$P_{\rm def}(T)$, whose component group is nonzero by
Theorem~\ref{result-3}.  Thus the induced mapping-space map is
not a weak equivalence.  This is the required Dwyer--Kan obstruction
inside the one-point-core subcategory.

For the positive assertion attach the full mapping cylinder along
its outer face.  The product collar on the exterior side comes
from the bordant pair; the mapping-cylinder side supplies the
inner product and its limiting core.  Thus the specified tube is
strict.  Local triviality of $\varphi$ identifies its neighborhood
of the core, locally over $Y$, with the product of an open subset
of $Y$ and the cone on the bundle fibre.

Conversely, from a target object $(X,Y,V,B)$ with the stipulated
strict mapping-cylinder description recover the exterior
\[
 E=X\setminus(V\setminus B).
\]
The complement is closed, because $V\setminus B$ is the open
inner part of the strict tube.  Its outer collar makes $(E,B)$
a bordant pair.  Reattaching $M_\varphi$ reconstructs the original
target with its given markings and parametrization.

A resolved morphism $(h_E,h_Y)$ has boundary restriction $h_B$.
It extends to the tube by the displayed product formula, and to
the core by $h_Y$.  At $t=0$ the compatibility equation
$h_Y\varphi=\varphi'h_B$ makes this well-defined and continuous
on the quotient.  At $t=1$ it agrees with $h_E$ on the glued
outer face.  Applying the same construction to the inverse data
proves that the extension is a homeomorphism.  A target morphism
with the stipulated product formula restricts to exactly those
resolved maps, so restriction and extension are inverse on
morphism sets and commute with composition.

The same formulas apply to a joint definable family over every
standard simplex $\Delta^p_K$.  Continuity at the core can also
be checked in the closed bounded model consisting of the points
$(tb,\varphi(b),t)$ together with the entire core copy
$\{(0,y,0):y\in Y\}$. The output first coordinates are $t h_B(b)$,
with $h_B$ uniformly bounded, and the core coordinate is
$h_Y(\varphi(b))$, with the parameter included in both maps.
These formulas are jointly continuous at $t=0$.  Gluing along
the closed outer face gives the required family, and the same
argument applies to the inverse.  Parameter pullback leaves
the formulas unchanged.
Thus extension and restriction commute with faces and
degeneracies and are inverse on the simplicial mapping spaces.
If $B$ is empty, the construction is the disjoint union $E\amalg Y$,
and the same restriction and extension statements hold.

The full product condition is necessary for this particular
inverse construction.  For instance, with $B=[-1,1]$ and
one-point core, put $a(t)=t(1-t)$.  The map
\[
 [b,t]\longmapsto
 [(1-a(t))b+a(t)b^3,t]
\]
is a definable radius-preserving tube automorphism, equal to the
identity on its outer face and core, and nonidentity inside.
Strict monotonicity in $b$ and the cone topology give its
continuous inverse.  It extends by the identity on the exterior.
Consequently retaining only the radius would leave distinct
target morphisms with the same resolved restriction.
\end{proof}


\subsection{A limitation of the fibre-dimension condition}

Let $Y$ be a closed definable subset of a closed bounded definable
space $X$.  An admissible mapping-cylinder description of a strict
tube $(V,Y,B)$ is a definable homeomorphism
\[
 M_\varphi=
 \bigl(B\times[0,1]\amalg Y\bigr)/((b,0)\sim\varphi(b))
 \longrightarrow V
\]
which is the identity on $Y$ and uses the prescribed identification
of $B\times\{1\}$ with the outer face $B$.  Here $V$ is a closed
subset of $X$, $V\setminus B$ is an open neighborhood of $Y$ in $X$,
and $B$ has the exterior collar required for a strict tube.
Write $q(b,t)$ for the associated map on $B\times[0,1]$.
In particular, its restriction at positive $t$ is a homeomorphism
onto $V\setminus Y$.

The obstruction comes from the junction of a disk and a transverse
segment. A finite attachment fibre would separate a connected
surface by finitely many points. We first justify the PL charts used
in that argument; a second proof below detects the same obstruction
by a winding circle and requires no manifold assumption on the
outer face.

\begin{lemma}[Compact marked charts and boundary preservation]
\label{lem:tube-marked-compact-charts}
Let $P$ be a compact polyhedron over an o-minimal real closed field.
A definable manifold chart at a point $p\in P$ gives a PL manifold
chart at $p$ in the existing PL structure on $P$.
For a boundary chart the conclusion respects its marked boundary.
Consequently a definable homeomorphism between compact definable
manifolds with boundary preserves their boundaries.
\end{lemma}
\begin{proof}
Subdivide so that $p$ is a vertex.  In the given chart choose the
inverse image $C$ of a small closed coordinate cube, with $p$ mapping
to its centre and with $C$ containing a relative neighborhood of $p$
in $P$.  Apply \cite[Theorem~2.2(1)--(2), author p.~6]{Shiota2010}
to the constant scalar function on $P$, marking $C$.
It gives a definable self-homeomorphism $\tau$ of $P$ preserving the
chosen triangulation and such that $C'=\tau^{-1}(C)$ is a
subpolyhedron.  Since $p$ is a vertex, $\tau(p)=p$, and $C'$ still
contains a neighborhood of $p$ in the original polyhedron.
The chart composed with $\tau$ identifies the compact marked pair
$(C',\{p\})$ definably with the closed cube and its centre.
Marked Hauptvermutung
\cite[Theorem~2.1, author p.~5]{Shiota2010}
provides a PL homeomorphism of these pairs.  It is a PL Euclidean
chart at the specified point $p$.

At a boundary point use a closed half-cube with $p$ in the relative
interior of its flat face, marking that face as well as $p$.
The same construction respects the boundary.  Thus a triangulation
of a compact definable manifold together with its designated boundary
is a PL manifold with that PL boundary.

Now let $h:M\to M'$ be a definable homeomorphism of two such compact
manifolds, and let $x\in M$.  In their compact polyhedral models,
subdivide to make $x$ and $h(x)$ vertices.  Apply marked
Hauptvermutung to the pairs $(M,\{x\})$ and $(M',\{h(x)\})$.
The resulting PL homeomorphism takes $x$ to $h(x)$ and preserves
intrinsic PL boundary.  Hence $x\in\partial M$ if and only if
$h(x)\in\partial M'$.  The argument does not assume beforehand
that $h$ preserves the boundary.
\end{proof}


\begin{theorem}\label{thm:tube-fibre-dimension-obstruction}
There is a compact nowhere dense semialgebraic subset $Y\subset S^3$,
with $Y\subset\overline{S^3\setminus Y}$, and a point $a\in Y$ such
that every admissible mapping-cylinder description of a strict tube
about $Y$ has
\[
 \dim\varphi^{-1}(a)>0,
 \qquad
 \dim_a S^3-\dim_aY-1=0.
\]
Thus the equality
\begin{equation}\label{eq:tube-fibre-dimension}
 \dim\varphi^{-1}(y)=\dim_yX-\dim_yY-1
\end{equation}
cannot always be imposed.  It holds for the link projection of a
locally trivial conical bundle over a pure-dimensional manifold.
Even when it holds, it does not imply contractibility of the derived
space of strict conical choices used in Theorem~\ref{result-12}.
\end{theorem}

\begin{proof}
In a stereographic chart of $S^3$, take
\[
 D=\{(x_1,x_2,0):x_1^2+x_2^2\leq1\},\qquad
 A=\{(0,0,z):0\leq z\leq1\},\qquad Y=D\cup A,
\]
and let $a=(0,0,0)$.  The stereographic inclusion gives a compact
semialgebraic subset of $S^3$.  Its interior is empty, and its local
dimension at $a$ is two.

Suppose $(V,Y,B)$ has an admissible description.  The positive product $B\times(0,1)$ is homeomorphic to an open
subset of $S^3$. Triangulate $B$ and take a compact product slab
$|B|\times[a,b]$ with $0<a<b<1$. At an interior point of each
vertical edge, the ambient definable three-manifold chart and
Lemma~\ref{lem:tube-marked-compact-charts} give a PL chart in
the product's existing PL structure. The link of the vertical
edge above a vertex $v$ of $B$ is, up to subdivision, the vertex
link of $v$ in $B$. It is a PL circle. Hence every vertex star
of $B$ is a disk, and $B$ is a closed definable surface without
boundary.

We next show that $B$ is connected.  Put
$r=(x_1^2+x_2^2)^{1/2}$.  For sufficiently small $\delta>0$, the
closure of
\[
 N=\{r<1+\delta,\ |z|<\delta\}
   \ \cup\ \{r<\delta,\ -\delta<z<1+\delta\}
\]
lies in $V\setminus B$ and in the chosen chart.  Such a $\delta$
exists because these neighborhoods shrink to the compact set $Y$.
The complement $N\setminus Y$ is path connected.  Here is an
explicit verification.  Below the disk, points can move radially
to the circle $r=1+\delta/2$ at height $-\delta/2$.  Above the disk,
points in the thin tube can first move off the axis, then descend
at a fixed positive radius smaller than $\delta$ to height
$\delta/2$.  Points in the enlarged disk neighborhood can move
within its upper half to that height.  They then move radially to
$r=1+\delta/2$.  This outer circle connects the upper and lower
parts through $z=0$ without meeting $D$.  Points initially on
$z=0$ outside $Y$ already have $r>1$.  These paths also cover the
axis above the tip of $A$, by first moving off that axis.

Every positive-radius ray of the mapping cylinder approaches a
point of $Y$, so it meets $N\setminus Y$.  Every point of
$V\setminus Y$ can therefore be joined to this connected set.
It follows that $V\setminus Y$, and hence $B$, is connected.

The map $q$ has compact, hence closed, image in $V$ and its image
contains $V\setminus Y$.  Since $Y$ is nowhere dense in $V$, this
image is all of $V$.  Positive-radius points map outside $Y$, so
$\varphi$ is onto.  The punctured set $Y\setminus\{a\}$ is the
disjoint union of the two nonempty relatively open and closed sets
$D\setminus\{a\}$ and $A\setminus\{a\}$.  Their inverse images
disconnect $B\setminus\varphi^{-1}(a)$.  A zero-dimensional
definable fibre is finite, whereas deleting finitely many points
from a connected surface does not disconnect it: paths can be
altered in finitely many disk charts to avoid those points.  This
contradiction proves $\dim\varphi^{-1}(a)>0$.

For the positive assertion, a local conical-bundle chart has the
form $U\times CL$, with $U$ open in a pure $m$-dimensional manifold
and the base given by $U$ times the cone vertex.  For nonempty $L$,
its local dimension along the base is $m+\dim L+1$, and the
attachment projection has fibre $L$.  This proves
\eqref{eq:tube-fibre-dimension}.  With the conventions
$C\varnothing=\{\mathrm{vertex}\}$ and
$\dim\varnothing=-1$, the same equality holds for the empty link.

Finally, take the point core of the cone on $T^6=(S^1)^6$.
For every strict neighborhood of its vertex, the punctured
neighborhood is an open seven-dimensional definable manifold and
is homeomorphic to $B\times(0,1)$. Thus $\dim B=6$.
Its attachment has the one vertex as target, so the sole fibre
is all of $B$ and satisfies \eqref{eq:tube-fibre-dimension}.
The numerical condition therefore includes the entire derived
component under consideration.  The derived choice component
specified in Theorem~\ref{result-12} has homotopy type
$BP_{\mathrm{def}}(T^6)$, whose fundamental group is nonzero by
Theorem~\ref{result-3}.  The numerical equality therefore does not
give categorical uniqueness in that precise derived model.
\end{proof}

\paragraph{Alternative proof of the dimension obstruction by winding.}
In a Euclidean chart of $S^3$, take
\[
 Y=\{(x,y,0):x^2+y^2\leq1\}
   \cup\{(0,0,z):0\leq z\leq1\},
 \qquad p=(0,0,0).
\]
This core is compact, semialgebraic, and nowhere dense.
At $p$ the proposed formula requires fibre dimension $3-2-1=0$.
Suppose an admissible tube
$q:B\times I\to V$ had a zero-dimensional fibre
$F=\varphi^{-1}(p)$.  The closed image of $q$ contains the dense
complement of $Y$ in $V$, so $q$ and $\varphi$ are onto and $F$ is
nonempty and finite.

Triangulate $(B,F)$ with the points of $F$ as vertices.
Small neighborhoods of the form $\{\lambda_f>c_f\}$, with
$c_f>1/2$, are contractible and pairwise disjoint; their straight-line
contractions to $f$ stay in their simplexes.
Choose them sufficiently small and choose $\delta>0$ so that, for
each $f\in F$,
\[
 q(N_f\times[0,\delta])\subset W,
\]
where $W$ is a small ambient ball about $p$, contained in $V$ and
of radius less than one in the displayed chart.
Continuity at the finitely many points $(f,0)$ supplies these choices.
The complement of
$\bigcup_f N_f\times[0,\delta)$ in $B\times I$ is closed and
bounded and its image under $q$ misses $p$.
Consequently there is a still smaller ball $W'$ about $p$ with
\[
 q^{-1}(W')\subset\bigcup_{f\in F}N_f\times[0,\delta).
\]
These are definable closed-image arguments over a general field.

Choose $a,h>0$ small enough that the horizontal circle
\[
 \gamma:S^1\longrightarrow W'\setminus Y,
 \qquad \gamma(u,v)=(au,av,h)
\]
lies in this smaller ball.
Its positive-radius inverse under $q$ is contained in one
$N_f\times(0,\delta)$, since $S^1$ is definably connected and the
sets in that finite union are disjoint and open.
This product is definably contractible, so applying $q$ gives a
nullhomotopy of $\gamma$ in $W\setminus Y$.

The plane section $W\cap\{z=0\}$ lies in the disk part of $Y$.
The connected image of the nullhomotopy must therefore remain in
$\{z>0\}$, where the normal segment removes the entire axis inside
$W$.  Angular projection
\[
 (x,y,z)\longmapsto
 \frac{(x,y)}{\sqrt{x^2+y^2}}
\]
is defined there and takes $\gamma$ to the identity of $S^1$.
This contradicts its nonzero degree, or equivalently its nontrivial
fundamental-group class.  In the arbitrary-field definable category
the same circle obstruction follows from
\cite[Corollary~3.4, author version p.~8]{BaroOtero2008v1}.
Thus every such fibre has positive dimension.
This argument imposes no manifold condition on $B$.

\paragraph{An explicit strict tube for this witness.}
There is also a direct finite construction showing that the example
has admissible strict tubes.
Choose a finite triangulation of $(S^3,Y)$, and barycentrically
subdivide so that the core is a full subcomplex $A$ of a finite
complex $L$.  For the semialgebraic witness the triangulation can
be chosen semialgebraic.  Work in this polyhedral model and write
\[
 w(x)=\sum_{v\in A}\lambda_v(x),\qquad
 V=\{w\geq1/2\},\qquad B=\{w=1/2\}.
\]
Fullness gives $w^{-1}(1)=|A|$, and $V$ contains the open
neighborhood $\{w>1/2\}$ of $|A|$.
For $b\in B$ put
\[
 a(b)=2\sum_{v\in A}\lambda_v(b)v,
 \qquad
 c(b)=2\sum_{v\notin A}\lambda_v(b)v.
\]
These are the normalized core and noncore parts of $b$, and
$b=(a(b)+c(b))/2$.
Define
\[
 q(b,t)=(1-t/2)a(b)+(t/2)c(b),\qquad 0\leq t\leq1,
 \qquad \varphi(b)=a(b).
\]
Both terms lie in faces of a common simplex, so this defines a
continuous semialgebraic map into $V$, with $q(b,1)=b$ and
$q(b,0)=\varphi(b)$.
For $x\in V\setminus|A|$, recover the two normalized parts by
\[
 a=\frac{\sum_{v\in A}\lambda_v(x)v}{w(x)},\qquad
 c=\frac{\sum_{v\notin A}\lambda_v(x)v}{1-w(x)}.
\]
The inverse of the positive-radius map is
\[
 x\longmapsto\left(\frac{a+c}{2},\,2(1-w(x))\right).
\]
It is continuous and semialgebraic.
Every point of the core lies in a simplex having a noncore coface:
the ambient complex is a triangulated three-manifold and the core
is nowhere dense.  Choosing a noncore vertex in such a coface
shows that $\varphi$ is onto.
The quotient identifies exactly those bottom points with the same
$\varphi$-value.  The finite polyhedral mapping-cylinder model and
closed-bounded definable compactness therefore identify its induced
map with a homeomorphism onto $V$, acting as the identity on the
specified copy of the core.
Finally the same displayed formula for $0<t<2$ has range
$\{0<w<1\}$ and the same inverse.
It gives a collar on both sides of $B$, including an exterior
collar for $t>1$.  The function $\rho=2(1-w)$ is zero exactly on
the core, and $V=\{\rho\leq1\}$, $B=\{\rho=1\}$.
Transporting this finite model back to $S^3$ proves strictness.


\subsection{Attachment fibres at a boundary core}

The real cell-like conclusion below is a boundary-core instance
of Lacher's relation between local connectivity of a mapping
cylinder and the neighborhood homotopy properties of its fibres
\cite[Theorem~2.2, pp.~652--653, and Corollary~2.4, p.~654]
{Lacher1970CellLikeII}.  We give the neighborhood argument explicitly
and then use definable triangulation to contract inside each fibre.
For an arbitrary o-minimal real closed field, all compactness,
homotopies, and manifold charts in the definable assertion have
their definable meanings.

\begin{theorem}\label{thm:boundary-attachment-fibres}
Let $V$ be a closed bounded definable manifold with boundary, let
$Y$ be a union of boundary components, and let $B$ be a closed
bounded definable set.  Suppose that a continuous definable map
\[
 q:B\times[0,1]\longrightarrow V
\]
satisfies $q(b,0)=\varphi(b)\in Y$ and restricts to a homeomorphism
$B\times(0,1]\longrightarrow V\setminus Y$.  Then $\varphi$ is
onto and each fibre $\varphi^{-1}(y)$ is definably contractible.

Over $\mathbb R$, suppose instead that $V$ is a compact topological
manifold with boundary and $B$ is a compact metric ANR, with the
same continuous-map hypotheses.  Then every fibre of $\varphi$ is
cell-like.  In the definable case over $\mathbb R$, these fibres
are contractible finite polyhedra up to definable homeomorphism.
\end{theorem}

\begin{proof}
The image of $q$ is closed and contains the dense set
$V\setminus Y$, so $q$ and $\varphi$ are onto.  

Fix $y\in Y$, put $F=\varphi^{-1}(y)$, and let $O$ be an open
neighborhood of $F$ in $B$, definable in the definable case.
The closed image $\varphi(B\setminus O)$ avoids $y$.  Choose an
open neighborhood $U$ of $y$ in $Y$ such that
$\varphi^{-1}(U)\subset O$.  For $0<\epsilon<1$, the set
\[
 N=U\ \cup\ q\bigl(\varphi^{-1}(U)\times[0,\epsilon)\bigr)
\]
is open in $V$. Indeed, its complement is the image under $q$ of
the closed bounded set
\[
 \bigl((B\setminus\varphi^{-1}(U))\times[0,1]\bigr)
 \ \cup\ \bigl(B\times[\epsilon,1]\bigr).
\]
That image is closed, also in the definable setting over an
arbitrary real closed field.  Choose a boundary half-ball $W$ about $y$ contained in $N$.
Because $Y$ is a union of boundary components, we may choose this
chart so that $W\cap Y$ is its whole flat boundary.  The set
$W\setminus Y$ is then homeomorphic to a convex open half-ball.
In particular, it has a contraction within itself, definable when
the data are definable.

There exists $0<\delta<\epsilon$ with
$q(F\times\{\delta\})\subset W\setminus Y$.  Indeed, the set
\[
 C=\{(b,t)\in F\times[0,1]:q(b,t)\notin W\}
\]
is compact and its projection to $[0,1]$ is closed and avoids zero.
Thus it is disjoint from an initial interval $[0,\delta]$ after
decreasing $\delta$.  This proof also applies to definably compact
sets and definable projections.

The open set
\[
 O_0=\{b\in B:q(b,\delta)\in W\}
\]
contains $F$.  Since $W\setminus Y\subset N\setminus Y$ and $q$
is injective at positive radius, $O_0\subset\varphi^{-1}(U)\subset O$.
Apply a contraction of $W\setminus Y$ to $q(b,\delta)$, and then
compose with
\[
 W\setminus Y\xrightarrow{\ q^{-1}\ }
 \varphi^{-1}(U)\times(0,\epsilon)
 \xrightarrow{\operatorname{pr}_B}\varphi^{-1}(U)\subset O.
\]
This is a nullhomotopy of the inclusion $O_0\hookrightarrow O$.
Its restriction is a contraction of $F$ inside $O$.

For the definable assertion, triangulate the compact pair $(B,F)$
using \cite[Theorem~2.2(1), p.~6]{Shiota2010}, and take a barycentric
subdivision.  Write the resulting pair as $(|K|,|L|)$, with $L$ a
full subcomplex of $K$.  If $a_v(x)$ denotes the global barycentric
coordinate of a vertex $v$ of $K$, put
\[
 O_L=\left\{x\in|K|:\sum_{v\in L}a_v(x)>0\right\},\qquad
 r(x)=\frac{\sum_{v\in L}a_v(x)v}{\sum_{v\in L}a_v(x)}.
\]
The vertices from $L$ having positive coordinates at $x$ span a
face in $L$, by fullness.  Hence $r$ is a continuous semialgebraic
retraction $O_L\longrightarrow|L|$.  The segment from $x$ to $r(x)$
stays in a simplex containing $x$.  Transport the retraction to
a definable neighborhood of $F$ in $B$, use the contraction of $F$
inside that neighborhood just constructed, and compose with the
retraction.  This contracts $F$ inside itself.

In the real topological case, the neighborhood construction proves
that $F\subset B$ has property $UV^\infty$.
Also, $B$ embeds as the compact slice $q(B\times\{1/2\})$ in $V$,
so it is finite dimensional.  A locally compact finite-dimensional
metric ANR is an ENR, and Lacher's criterion now makes $F$ cell-like
\cite[preliminary lemma, p.~718, and Theorem~1.1, p.~719]
{Lacher1969CellLike}.  This last conclusion requires no ANR
assumption on the fibre itself.
\end{proof}

\subsection{The numerical class of fixed-tube parametrizations}

Let $Y$ be a closed bounded definable manifold without boundary,
put $X=Y\times[0,2]$, and fix
\[
 V=Y\times[0,1],\qquad Y_0=Y\times\{0\},\qquad B=Y\times\{1\}.
\]
An admissible parametrization is a continuous definable map
$q:Y\times[0,1]\to V$ such that
\[
 q(y,1)=(y,1),\qquad q(y,0)=(\varphi(y),0),\qquad
 q|_{Y\times(0,1]}:Y\times(0,1]\xrightarrow{\sim}V\setminus Y_0.
\]
It induces the mapping-cylinder identification fixing the copy of
$Y_0$.  For the simplicial space of such maps, a simplex means a
jointly continuous definable family over the closed parameter
simplex, preserving that parameter.  Faces and degeneracies are
the usual affine pullbacks.  Write $P_{\mathrm{def}}(Y)$ for the
corresponding simplicial group of homeomorphisms of $Y\times[0,1]$
fixed pointwise on $Y\times\{0\}$.

\begin{theorem}\label{thm:numerical-tube-concordances}
The simplicial space of admissible parametrizations satisfying
\[
 \dim\varphi^{-1}(y)=\dim_yX-\dim_yY-1=0\qquad(y\in Y)
\]
is naturally isomorphic, by interval reversal, to
$P_{\mathrm{def}}(Y)$.  In particular, for $Y=T^6$ it has component
group
\[
 \bigoplus_{v\in\mathbb Z^6\setminus\{0\}}\mathbb F_2,
\]
and the numerical condition gives neither a connected nor a
contractible space of choices, even with $V$ and its outer-face
identification fixed.
\end{theorem}

\begin{proof}
By Theorem~\ref{thm:boundary-attachment-fibres}, every fibre of
$\varphi$ is nonempty and definably contractible.  Under the
dimension hypothesis each fibre is finite.  A nonempty finite
definably connected set is a point, so $\varphi$ is a continuous
definable bijection of compact sets, hence a homeomorphism.
It follows that $q$ is itself a continuous definable bijection
of compact sets, hence a homeomorphism of the cylinder fixing
its top pointwise.

Conversely, Lemma~\ref{lem:tube-marked-compact-charts} shows that
a definable homeomorphism of the cylinder preserves its whole
boundary over the given real closed field. If it fixes the top
pointwise, bijectivity prevents a bottom point from mapping to
the top. It therefore preserves the bottom setwise, including
when $Y$ is disconnected.  Its bottom restriction
is therefore a homeomorphism $\varphi$, and it supplies an
admissible parametrization with singleton fibres.  For a family
over a closed parameter simplex $\Delta$, the total map
\[
 \Delta\times Y\times[0,1]\longrightarrow
 \Delta\times Y\times[0,1],\qquad (s,y,t)\longmapsto(s,q_s(y,t))
\]
is a continuous definable bijection of compact sets.  Its inverse
is consequently jointly continuous and definable.  Thus the
identification holds for the stipulated families, as well as at
vertices.  Conjugation by $(y,t)\mapsto(y,1-t)$ changes the fixed
top to a fixed bottom and commutes with all affine parameter
pullbacks.  This is the required simplicial isomorphism.  The
torus calculation is Theorem~\ref{result-3}.  If $Y$ is empty,
both simplicial spaces are the one-point simplicial space.
\end{proof}

\paragraph{A noninjective attachment with a manifold mapping cylinder.}
The numerical hypothesis above has content.  Let the source circle
have cyclic vertices $b_0,b_1,b_2,b_3$ and the target circle have
cyclic vertices $c_0,c_1,c_2$.  Define a simplicial map by
\[
 b_0,b_1\longmapsto c_0,\qquad b_2\longmapsto c_1,\qquad
 b_3\longmapsto c_2.
\]
Its mapping cylinder is the annulus with the following seven
triangles:
\[
\begin{gathered}
 (b_0,b_1,c_0),\ (b_1,b_2,c_0),\ (b_2,c_1,c_0),\
 (b_2,b_3,c_1),\\
 (b_3,c_2,c_1),\ (b_3,b_0,c_2),\ (b_0,c_0,c_2).
\end{gathered}
\]
For each noncollapsed source edge the two triangles are the usual
triangulated square.  The collapsed edge has quotient parametrization
\[
 (u,t)\longmapsto (1-t)c_0+t\bigl((1-u)b_0+ub_1\bigr).
\]
This collapses exactly its bottom edge and is injective for $t>0$.
These parametrizations agree on the vertical edges, proving the
mapping-cylinder assertion.

For a direct planar verification of the annulus, assign coordinates
\[
\begin{aligned}
 b_0&=(-2,-2),&b_1&=(2,-2),&b_2&=(2,2),&b_3&=(-2,2),\\
 c_0&=(1,-1),&c_1&=(0,1),&c_2&=(-1,-1).
\end{aligned}
\]
The seven triangles have disjoint interiors, meet pairwise in
their indicated common faces, and meet the inner triangle only
in the indicated inner faces.  These assertions follow by the
linear half-plane inequalities of their edges.  Their areas, in
the order displayed, are
\[
 2,\quad2,\quad\tfrac52,\quad2,\quad\tfrac52,\quad2,\quad1.
\]
They sum to $14$, the area of the outer square minus the inner
triangle.  Their finite closed union is therefore exactly that
closed planar annulus: a missing interior point would leave a
nonempty open region of positive area, and both polygonal
boundaries are covered.  The outer four-cycle and inner
three-cycle are its two boundary components.  In particular,
the target circle is collared although the fibre over $c_0$ is
the nontrivial interval $[b_0,b_1]$.  This example is semialgebraic
over the rational numbers.  Its finite incidence inequalities
and quotient formulas remain valid over every real closed field.
Products give the same phenomenon with higher-dimensional
manifold boundary cores.

\subsection{Distinct components with unrestricted attachment fibres}

\begin{theorem}\label{thm:unrestricted-tube-components}
Let $Y=T^6=(S^1)^6$, put $V=Y\times[0,1]$, and fix its core
$Y_0=Y\times\{0\}$ and outer face $B=Y\times\{1\}$.
Let $\mathcal R$ be the space of all continuous maps $q:V\to V$
satisfying
\[
 q(y,1)=(y,1),\qquad q(Y_0)\subset Y_0,\qquad
 q|_{V\setminus Y_0}:V\setminus Y_0\xrightarrow{\sim}V\setminus Y_0,
\]
with the compact-open topology.  Its subspace of homeomorphisms
is $G=\operatorname{Homeo}(V\;\mathrm{rel}\;B)$.
There is a locally constant map
\[
 \lambda:\mathcal R\longrightarrow\pi_0G,
 \qquad\lambda(g)=[g]\quad(g\in G),
\]
where the target has the discrete topology.  Consequently
$\pi_0G\longrightarrow\pi_0\mathcal R$ is injective, and
$\mathcal R$ is not path connected.

Two fixed finite PL parametrizations witnessing this failure
also cannot be joined by a jointly continuous definable family
in any o-minimal expansion of $\mathbb R$, or by a semialgebraic
family over any real closed field.  Thus allowing arbitrary
attachment fibres does not give essential uniqueness for this
fixed strict tube in $Y\times[0,2]$.
\end{theorem}

\begin{proof}
We first record a metric observation.  Let $Z$ be a compact metric
space and $H$ a subgroup of $\operatorname{Homeo}(Z)$ for which
the identity path component contains an open neighborhood of the
identity.  The uniform metric
\[
 d_\infty(f,h)=\sup_{z\in Z}d(f(z),h(z))
\]
induces the compact-open topology on continuous self-maps and
on the homeomorphism group.  On homeomorphisms it is right
invariant.  In particular,
\begin{equation}\label{eq:tube-uniform-right-invariance}
 d_\infty(hg^{-1},\operatorname{id})=d_\infty(h,g).
\end{equation}
Choose $\epsilon>0$ such that every $h\in H$ with
$d_\infty(h,\operatorname{id})<\epsilon$ is in the identity
path component.  Equation~\eqref{eq:tube-uniform-right-invariance}
shows that any two elements of $H$ at distance less than
$\epsilon$ are in the same path component.

Let $\overline H$ denote its closure in $C(Z,Z)$.  Given
$f\in\overline H$, choose $h\in H$ with
$d_\infty(f,h)<\epsilon/4$ and set $\lambda_H(f)=[h]$.
Two such choices have distance less than $\epsilon/2$, so this
definition is independent of the choice.  If
$d_\infty(f,f')<\epsilon/4$, approximating homeomorphisms for
$f$ and $f'$ have distance less than $3\epsilon/4$, showing
$\lambda_H(f)=\lambda_H(f')$.  Thus $\lambda_H$ is locally
constant and extends the component map of $H$.

We apply this observation to $G$.  Edwards--Kirby local
contractibility applies to the homeomorphism group of the compact
manifold with boundary $V$
\cite[Corollary~1.1, p.~63; proof, p.~77]{EdwardsKirby1971}.
Here is the passage to the fixed-top subgroup.  The subgroup
$H\subset\operatorname{Homeo}(V)$ preserving both boundary
components is open.  If $h\in H$, write
$h(y,1)=(a_h(y),1)$ and define
\[
 R(h)=(a_h^{-1}\times\operatorname{id}_{[0,1]})\circ h.
\]
Restriction, inversion, and composition are continuous in the
compact-open group topology.  Hence $R:H\to G$ is continuous
and is the identity on $G$.  Take a sufficiently small local
contraction near the identity whose image lies in $H$ and
compose it with $R$.  It follows that a neighborhood of the
identity in $G$ contracts inside $G$.  In particular, $G$ has
the radius $\epsilon$ required above.

It remains to prove $\mathcal R\subset\overline G$.  Fix
$q\in\mathcal R$.  The real assertion of
Theorem~\ref{thm:boundary-attachment-fibres} shows that the
fibres of its bottom restriction are cell-like.  All other
fibres of $q$ are singletons.  The map $q$ is onto by compactness
and density, and is proper.  Moreover its restriction to the
whole boundary is a proper cell-like map onto the whole
boundary: it is the bottom attachment on $Y_0$ and the identity
on $B$.

Siebenmann's approximation theorem now applies to the compact
seven-manifolds with boundary
\cite[Theorem~A, p.~271]{Siebenmann1972Approximation}.
The theorem requires a proper cell-like map and a cell-like
boundary restriction.  Its condition~(a), which excludes
dimensions $3,4,5$, holds here.  Thus, for every $\eta>0$, there
is a homeomorphism $h:V\to V$ with
$d_\infty(h,q)<\eta$.

Use a product metric on $Y\times[0,1]$.  Choose $\eta$ smaller
than half the distance between the boundary components.
A homeomorphism preserves the manifold boundary and maps each
of its connected components onto a boundary component.  Since
$q|_B=\operatorname{id}_B$ and $d_\infty(h,q)<\eta$, no point
of $B$ can be sent to $Y_0$.  Hence $h(B)=B$ and $h(Y_0)=Y_0$.
The product metric and the same estimate on $B$ give
$d_\infty(a_h,\operatorname{id}_Y)<\eta$.  In addition,
\[
 d_\infty(a_h^{-1},\operatorname{id}_Y)
   =d_\infty(a_h,\operatorname{id}_Y)<\eta.
\]
Consequently $g=R(h)$ is in $G$ and
$d_\infty(g,q)<2\eta$.  This proves the required density
without making a continuous selection of approximations.
Restricting $\lambda_G$ to $\mathcal R$ gives $\lambda$.
Any path in $\mathcal R$ has constant $\lambda$, proving
injectivity on path components.

After interval reversal, $G$ is the ordinary topological
concordance group of $T^6$.  Theorem~\ref{result-3} gives
\[
 \pi_0G\cong
 \bigoplus_{v\in\mathbb Z^6\setminus\{0\}}\mathbb F_2.
\]
Choose a nonzero class represented by a finite PL concordance,
using the PL--TOP comparison in Theorem~\ref{result-4}.
By the finite coefficient model of
Theorem~\ref{thm:effective-countable-concordance-models}, we
may choose the representative $p$ over the field $F$ of real
algebraic numbers, on a fixed finite PL model of the torus.
Reversing the interval makes $p$ top-fixed.  The identity and
$p$ have different $\lambda$-values.  Both give admissible
mapping-cylinder identifications fixing the specified copy
of the core and the outer face.

A jointly continuous definable path over $\mathbb R$ is an
ordinary continuous path into $C(V,V)$, because $V$ is compact.
It therefore cannot join these two points of $\mathcal R$.
The same applies to a finite chain of definable one-simplices,
which can be concatenated using finitely many linear changes
of the parameter.

For the semialgebraic assertion, let $K$ be any real closed
field containing its canonical ordered copy of $F$.  Suppose
there were a semialgebraic family joining the two fixed
endpoints over $K$.  Its graph has a finite formula in the
language of ordered fields with a finite coefficient tuple
from $K$.  The conditions that it be a function, be jointly
continuous, have the specified endpoints, fix the top, send
the bottom into the bottom, and give a homeomorphism on the
complement of the bottom in every fibre are first-order
conditions on that tuple.  For the last condition one may
express bijectivity and continuity of the inverse graph by
the usual quantified $\varepsilon$--$\delta$ formulas.
All coefficients specifying the endpoints and the torus
belong to $F$.  Tarski's quantifier-elimination theorem makes
embeddings of real closed fields elementary
\cite[Theorem~3.2 and Corollary~3.3, pp.~2--3]{SwanTarski}.
Existential quantification of the finite tuple
and the elementary embedding $F\subset K$ give such a family
over $F$; the elementary embedding $F\subset\mathbb R$ then
gives a semialgebraic family over $\mathbb R$.  This contradicts
the preceding paragraph.  This transfer concerns semialgebraic
families over arbitrary real closed fields; the assertion for
general o-minimal families here is over $\mathbb R$.
\end{proof}

\begin{theorem}[Ordinary components of unrestricted rigidifications]
\label{thm:unrestricted-tube-component-bijection}
For the ordinary real six-torus, let
$V=T^6\times[0,1]$, let $B=T^6\times\{1\}$, and let
$G=\operatorname{Homeo}(V\;\mathrm{rel}\;B)$.
Let $\mathcal R$ have the compact-open topology and consist of the
continuous maps $q:V\to V$ that fix $B$ pointwise, map the bottom
into itself, and restrict to a homeomorphism off the bottom.
Every $q\in\mathcal R$ admits a continuous path
$\gamma:[0,1]\to\mathcal R$ with
\[
 \gamma(0)=q,\qquad \gamma(t)\in G\quad(0<t\leq1).
\]
Consequently inclusion induces a bijection
\[
 \pi_0G\xrightarrow{\ \cong\ }\pi_0\mathcal R.
\]
Under composition this is an isomorphism of component monoids;
in particular $\pi_0\mathcal R$ is the group identified with
$\pi_0G$.  The inverse bijection is induced by the locally constant
label of Theorem~\ref{thm:unrestricted-tube-components}.
\end{theorem}
\begin{proof}
The proof of Theorem~\ref{thm:unrestricted-tube-components} gives
$G\subset\mathcal R\subset\overline G$ for uniform convergence on
$V$, and gives a locally constant label
$\lambda:\mathcal R\to\pi_0G$ restricting to $g\mapsto[g]$.
Thus inclusion is already injective on path components.
We prove surjectivity with a controlled limiting isotopy.
This is the mechanism explicitly stated in the published Complement
to Theorem~A of
\cite[p.~272]{Siebenmann1972Approximation}.
We spell out the fixed-top construction.

For every $\epsilon>0$ there is $\delta>0$ such that any
$u\in G$ with $d_\infty(u,\mathrm{id})<\delta$ can be joined to
$\mathrm{id}$ inside the $\epsilon$-ball of $G$.
Indeed, Edwards--Kirby local contractibility
\cite[Corollary~1.1, pp.~63, 77]{EdwardsKirby1971}
provides contractions in arbitrarily small neighborhoods of the
identity in the manifold homeomorphism group.
Restrict to the open subgroup preserving the two boundary components
and compose its contraction with the continuous retraction
\[
 R(h)=(a_h^{-1}\times\mathrm{id}_{[0,1]})\circ h,
 \qquad a_h=h|_{T^6\times\{1\}}.
\]
Since $R$ fixes $G$ and the identity, continuity allows the
contraction to take values in any prescribed uniform neighborhood
in $G$.  This proves the claimed controlled path property.

For $n\geq1$ choose $\delta_n>0$ for
$\epsilon_n=2^{-n}$ and set
\[
 \eta_n=\min\{2^{-n},\delta_1/4,\ldots,\delta_n/4\}.
\]
For the specified $q\in\mathcal R$, density gives $g_n\in G$ with
$d_\infty(g_n,q)<\eta_n$.
Surjectivity of $g_n$ gives the exact identity
\[
 d_\infty(g_{n+1}g_n^{-1},\mathrm{id})
   =d_\infty(g_{n+1},g_n)<\delta_n/2.
\]
Choose a path $u_n$ from $g_{n+1}g_n^{-1}$ to the identity in the
$\epsilon_n$-ball, and put $\alpha_n(s)=u_n(s)g_n$.
It joins $g_{n+1}$ to $g_n$ and satisfies
\[
 d_\infty(\alpha_n(s),g_n)<\epsilon_n
 \qquad(0\leq s\leq1).
\]
On $[2^{-n},2^{-(n-1)}]$ define $\gamma$ by the affine
reparametrization of $\alpha_n$, with lower endpoint $g_{n+1}$
and upper endpoint $g_n$.  Define $\gamma(0)=q$.
The endpoints agree at every joining time.  On that interval,
\[
 d_\infty(\gamma(t),q)<\epsilon_n+\eta_n\longrightarrow0.
\]
Hence $\gamma$ is continuous at zero and everywhere else, with
$\gamma(1)=g_1$.  All positive-time maps are in $G$.
Compactness of $V$ makes this compact-open path equivalent to a
jointly continuous family $[0,1]\times V\to V$.
This proves surjectivity on path components and, with the label,
the stated bijection.

Composition preserves all defining conditions of $\mathcal R$
and is continuous for compact-open maps on compact $V$.
Inclusion is a monoid homomorphism, so the resulting bijection of
component monoids transports the group structure of $\pi_0G$.
\end{proof}

\paragraph{Scope.}
The paths just constructed are ordinary real paths obtained by
countably many successive isotopies.  No definable or semialgebraic
path-surjectivity, arbitrary-field analogue, higher weak equivalence,
or continuous choice of the path as $q$ varies is asserted.
Siebenmann's published Complement gives this limiting-isotopy
conclusion. Its subsequent announcement of nearly canonical
dependence on the map is not needed here.


\section{The collar and cone categories}
This section computes the maps produced by attachment, collapse, and forgetting a choice.

\begin{theorem}[Rigidified blow-up and forgetting cone rigidification]\label{result-6}
Use strict collars and strict conical neighborhoods, and require morphisms
to preserve the chosen product or cone parametrizations. Then
\[
 PBCR\simeq EPCR\simeq EPC.
\]
The first is an equivalence induced by collapsing the lower collar
boundary and blowing up the cone vertex. The second is a weak equivalence
on every mapping space and is essentially surjective. These assertions
hold with any finite number of disjoint labeled boundary pieces or points.
They do not imply equivalence with $EP$.
\end{theorem}
\begin{proof}
For $(X,B,c)$ collapse $B$ and identify the image of $c$ with
$C\dot C$ by its radial parameter. Conversely for $(Z,o,C,\psi)$ put
$E=Z\setminus(C\setminus\dot C)$ and form
\[
 \widetilde Z=E\amalg_{\dot C}(\dot C\times[0,1]),
\]
where $\dot C\times\{1\}$ is attached to $\dot C\subset E$.
The new lower boundary is $\dot C\times\{0\}$ and its displayed
product is the rigidified collar. Strictness is exactly the condition
that the outer face is collared in $E$. Each of these two constructions
followed by the other has the canonical indicated homeomorphism to the
starting object. A morphism respecting radial parameters extends
uniquely across the blown-up lower boundary; a product-preserving
morphism descends uniquely. These statements hold on total spaces of
families and respect simplicial operators. They prove the first
equivalence, including its inverse on all maps.

For the second functor fix any two rigidified cone objects. Given a map
$f$ preserving the unparametrized cone, replace its restriction inside
that cone by
\[
 \psi'\,C(f|_{\dot C})\,\psi^{-1}
\]
and leave it equal to $f$ outside. The formulas agree on $\dot C$.
The discrepancy is a homeomorphism of the source cone fixing base and
vertex. Contract it by Theorem~\ref{result-1} and extend the contraction as the identity
outside the cone. This gives a deformation retraction of mapping spaces
onto the maps preserving the rigidification. Every unrigidified cone has
a rigidification by its definition, proving essential surjectivity.
For finitely many disjoint cones use the formulas separately on each
closed cone; they agree on the complementary region. Permutations of
labels commute with the constructions. Empty boundary pieces correspond
to isolated cone vertices under $C\varnothing=\{o\}$.
\end{proof}

\begin{proof}[Proof of Theorem~\ref{result-7}]
Precomposition makes $E=\operatorname{Col}^{\mathrm{rig}}(X,B)$ a free
simplicial $P_{\mathrm{def}}(B)$-set. Two parametrizations have exactly
the same collar and outer face if and only if their transition map is
in this group. Thus the stated unparametrized space is the degreewise
quotient $E/P_{\mathrm{def}}(B)$.

This is a principal Kan fibration. Here are the needed details for the
family convention. Both the group and the embedding space satisfy horn
filling: glue the compatible maps on a finite simplicial horn and pull
the resulting family back by a PL retraction of the simplex onto the
horn. Such a retraction is obtained by an elementary simplicial collapse.
For a horn lifting problem for the quotient, choose a representative of
the given quotient simplex. Its transition maps to the prescribed horn
lifts form a horn in the group, since the action is free. Fill that
group horn and adjust the representative. This proves the lifting
property. Theorem~\ref{result-2} contracts the total space. The long exact homotopy
sequence and the universal principal-bundle construction therefore
identify the base with the classifying space of the group. The
nonvanishing follows from Theorem~\ref{result-3}.
\end{proof}

\begin{theorem}[Local collars glue over a real closed field]\label{result-8}
Let $B$ be closed in a definable subset $X$ of an affine space over a
real closed field. If $B$ has a finite definable local collar atlas, then
there is a definable open collar $B\times[0,1)\to X$. If $X$ and $B$
are definably compact, it restricts after rescaling to a strict closed
collar. Thus I a) holds for the compact bordant pairs used here.
\end{theorem}
\begin{proof}
Brown's finite gluing lemma is \cite[Lemmas 1--4, pp.~333--337]{Brown1962}:
the union of two relatively open collared subsets is collared. Its proof
uses distance functions, minima, finite closed shrinkings of an open
cover, and piecewise affine changes of the interval coordinate. In a
definable subset over a real closed field these operations are definable.
Distances to definable sets exist by definable completeness; a distance
to the complement is positive at every point of an open set. Closed
shrinkings and separated neighborhoods are given by the same distance
inequalities. Brown's spindle modification is the identity below half
a positive height and conjugates the transition map by the linear
collapse above that height. Its inverse has the displayed same finite
piecewise formula, so its continuity proof uses no sequential
completeness or infinite process. Hence Lemma 4 holds verbatim for a
finite definable collar atlas. Induction on its finite size yields an
open collar. In the compact case its restriction to $[0,1/2]$ is a
closed embedded collar, and restriction to $[0,3/4]$ supplies an outer
extension, proving strictness.
\end{proof}

\begin{theorem}[Existence of strict pointed cones]\label{result-9}
Every compact definable pointed space $(X,o)$ admits a strict closed
conical neighborhood of $o$, including when $o$ is isolated.
\end{theorem}
\begin{proof}
Take a finite definable triangulation compatible with $\{o\}$, so that
$o$ is a vertex. Its barycentric coordinate $\lambda_o$ is a continuous
PL function, extended by zero off the closed star. For $0<\epsilon<1$,
the set $C=\{\lambda_o\ge\epsilon\}$ is a cone with vertex $o$ on the
level set $\{\lambda_o=\epsilon\}$. Varying the coordinate in
$[\epsilon/2,\epsilon]$ gives an exterior collar of its outer face.
The complement of that outer face in $C$ is open around $o$. If the
vertex is isolated the level set is empty and $C=\{o\}$. The finite
triangulation input is the definable triangulation theorem used in
\cite[\S1]{Shiota2010}; the cone and strictness are given explicitly here.
\end{proof}

\begin{theorem}[Attaching a collar and attaching a cone are enriched equivalences]\label{result-10}
Let $F:PB\to PBCR$ attach $B\times[-1,0]$ to $X$ along
$B\times\{0\}=B$, with new boundary $B\times\{-1\}$ and its
displayed rigid collar. Let $K:PB\to EPCR$ attach $CB$ to $X$
along its base. Both functors are equivalences of the specified
definable simplicial categories, with actual isomorphisms on mapping
spaces. In particular the valid corrected comparison is
\[
 PB\simeq PBCR\simeq EPCR\simeq EPC.
\]
This conclusion does not include $EP$.
\end{theorem}
\begin{proof}
The complement of the attached collar minus its outer face is exactly
the original $X$. A map respecting the displayed rigid collar restricts
to a homeomorphism of $(X,B)$; conversely any such homeomorphism extends
uniquely by the product of its boundary restriction and the interval
identity. These operations are inverse on all parameter families.
The inverse construction on an arbitrary rigid collar object is
\[
 (Z,B,C,c)\longmapsto
 \bigl(Z\setminus(C\setminus\dot C),\dot C\bigr).
\]
Strictness says precisely that this is a bordant pair. The chosen
parametrization identifies the collar reattached to it with the
original collar, giving a natural isomorphism to the starting object.
Thus $F$ is fully faithful and essentially surjective, with its stated
inverse.

For $K$ use the same exterior and attach a cone. A rigid cone map is
uniquely the cone of its base map, and its exterior restriction is a
bordant-pair map. Hence the identical argument proves the assertion
for $K$. The other equivalences in the display are Theorem~\ref{result-6}. Our pointed
quotient convention is $X/B=(X\amalg\{o\})/(B\sim o)$; when
$B=\varnothing$ this is $X\amalg\{o\}$. Thus empty boundary pieces
correspond to isolated marked points and all inverse constructions
remain valid. Finitely many disjoint labeled pieces are treated by
the same formula on each piece; the formulas commute with permutations
of the labels.
\end{proof}

\begin{theorem}[The naturality qualification for the two displayed $2$-cells]\label{result-11}
On $PB$ with all homeomorphisms, neither the collar-absorption
isomorphism $UF\Rightarrow\id$ nor the cone-absorption isomorphism
$VK\Rightarrow Q$ is a strict natural isomorphism. Here $U$ forgets
the attached collar, $V$ forgets the attached cone, and $Q$ collapses
the original boundary. They become strict natural isomorphisms after
passing to pairs with a chosen extended collar and maps preserving its
product coordinates. This is the role of the manuscript's marginal
instruction to precompose with $PBR$ at cote~119, sheet~149 (PDF p.~150).
\end{theorem}
\begin{proof}
Take $X=[0,1]$ and $B=\{0\}$. Every increasing definable
homeomorphism of this interval fixes the newly attached collar or cone
pointwise, because it fixes $B$. The common fixed set of these
homeomorphisms on the original interval is exactly $\{0,1\}$:
each interior point is moved by an increasing piecewise affine
homeomorphism. Naturality would force the image of the entire attached
interval under either proposed isomorphism to lie in this two-point
set. No homeomorphism has this property. This proves both failures.

Now choose $c:B\times[0,2]\hookrightarrow X$. In the union of its
image with the attached collar use the coordinate $t\in[-1,2]$ and
send it to $2(t+1)/3\in[0,2]$, keeping $b$ fixed; use the identity
outside this extended collar. The formulas agree at $t=2$, give a
homeomorphism, and send the new boundary to the old one. They commute
with every map preserving the chosen product coordinates. This is
$\alpha$ on that category. Quotienting its source and target by their
respective boundary pieces gives $\beta$, because the source quotient
is exactly the attached-cone construction. This also proves the
compatibility of the two $2$-cells. The piecewise construction applies
independently to each of finitely many disjoint chosen collars.
\end{proof}

\begin{proof}[Proof of Theorem~\ref{result-12}]
Give $CB$ radial interval $[0,2]$ and choose the smaller cone
$C_1B$ with radii $[0,1]$. Its exterior is $B\times[1,2]$.
Let $H$ be the simplicial group of homeomorphisms of $CB$ fixing
the vertex, and let $J\subset H$ preserve the smaller cone.
For this torus link the outer base is the intrinsic manifold boundary
of the cone's manifold locus, so every member of $H$ preserves it.
Restriction to the outer base gives split homomorphisms
\[
 H\longrightarrow\Homeo(B),\qquad
 J\longrightarrow\Homeo(B).
\]
The first kernel is Alexander-contractible by Theorem~\ref{result-1}. Write $J_0$
for the second kernel. Restriction of $J_0$ to the shell gives a
split homomorphism
\[
 J_0\longrightarrow
 \Homeo(B\times[1,2];B\times\{2\}),
\]
whose section cones the map induced on the inner face. Its kernel
is the group of maps of $C_1B$ fixed on its base and vertex, again
Alexander-contractible. The target is $P_{\mathrm{def}}(B)$ after
reversing and rescaling the interval. These split decompositions are
decompositions of simplicial spaces; hence $J_0\simeq P_{\mathrm{def}}(B)$.
It follows, by applying classifying spaces to the split group maps,
that the homotopy fibre of $BJ\to BH$ is $BJ_0\simeq BP_{\mathrm{def}}(B)$.
The component of an object in the classifying space of a simplicial
groupoid is the classifying space of its automorphism group. Therefore
this is exactly the indicated component of the derived choice space.
Surjectivity of $\pi_0J\to\pi_0H$ makes this fibre connected. Other
components of the cone-object category cannot remove its nonzero
fundamental group.

Theorem~\ref{result-3} gives $\pi_1BP_{\mathrm{def}}(B)=\pi_0P_{\mathrm{def}}(B)\ne0$.
Finally Theorem~\ref{result-6} identifies the rigidified and unrigidified cone categories
over the pointed-space category by weak equivalences on mapping spaces,
so their derived choice spaces agree. This argument does not invoke an
unproved parametrized isotopy-extension theorem to identify a different,
pointwise-defined topology on a set of cone neighborhoods with this
derived space.
\end{proof}

\begin{proof}[Proof of Theorem~\ref{result-13}]
Only the three remaining collar arrows require verification. Put
$X=B\times[0,2]$, with boundary at $0$ and chosen collar $[0,1]$.
Let $J_c$ be its collar-preserving automorphism group. Restriction to
the upper half gives a split map
\[
 J_c\longrightarrow G_+=\Homeo(B\times[1,2];B\times\{1\}
 \text{ setwise}).
\]
Its section is constant extension of the inner boundary map across
the lower half. Its kernel is the lower-half group fixed on its top,
which is $P_{\mathrm{def}}(B)$ after reversal. The automorphism group
with rigid collar is exactly this section $G_+$. Since the underlying
space of $J_c$ is a product of $G_+$ with this kernel, the section
fails to induce a surjection on components. This disproves
$PBCR\to PBC$ at the displayed object.

The inclusion of the rigid-collar automorphisms into all automorphisms
of $(X,B\times\{0\})$ is a homotopy equivalence in this example.
To see it explicitly, write an automorphism as the constant extension
of its lower boundary map times a bottom-fixed $p$. For $0\le s\le1$
compress $p$ into $[s,2]$: use the identity below $s$ and
$D_s pD_s^{-1}$ above $s$, where
$D_s(b,t)=(b,s+(2-s)t/2)$. The formulas agree at $s$ because $p$
fixes the bottom, and vary continuously and definably, including
$s=0$. At $s=1$ the result preserves the rigid collar. If it already
preserved that collar, this entire homotopy stays within the same
subgroup. It therefore proves the claimed homotopy equivalence.
If forgetting the unrigidified collar were also a weak equivalence
on this mapping space, the two-out-of-three property would contradict
the failure of the preceding section. Thus $PBC\to PB$ fails.

After collapse, restriction to the unchanged upper shell identifies
the cone-preserving group with $G_+$ up to the Alexander-contractible
kernel of maps of the small cone fixed on its base and vertex.
The map from $J_c$ commutes with this restriction. Its lower-half
kernel $P_{\mathrm{def}}(B)$ maps into that contractible cone kernel.
The same split product decompositions show that its homotopy fibre
is $P_{\mathrm{def}}(B)$. Hence $PBC\to EPC$ fails too. The other
entries are the stated earlier results. Taking one nontrivial boundary
piece and leaving all other labels fixed gives the same failures in
every finite multipair variant in which that piece is permitted.
\end{proof}

\section{An intrinsic partition with saturated singular loci}
Weak regular filtrations can have incompatible layers. We first give
two explicit examples, then prove the canonical theorem with its
frontier and singular-saturation conditions.

\subsection{Two obstructions to a coarsest weak filtration}
\label{subsec:weak-filtration-counterexamples}
A \emph{weak regular equidimensional filtration} is a finite increasing
chain of closed subsets whose nonempty successive differences are
manifolds without boundary, each of a single dimension.  Its layers
need not be connected.  Refinement is refinement of the associated
layer partitions.  These requirements alone do not force a frontier
condition or singular-locus saturation for closed unions.

\begin{proposition}[A compact counterexample in three-dimensional space]
\label{prop:weak-filtration-r3}
There is a compact semialgebraic set $X\subset\mathbb R^3$ with two
closed filtrations having layers of dimensions $0,1,2$, all embedded
smooth manifolds, such that the only common coarsening of their layer
partitions is $\{X\}$.  The set $X$ is not a manifold of a single
dimension.  Consequently the two filtrations have no common weak
regular coarsening, for either embedded smooth or intrinsic PL
regularity.
\end{proposition}
\begin{proof}
Let
\[
 S=\{(x,y,z):x^2+y^2=z^2(1-z),\ 0\le z\le1\},
 \qquad p=(0,0,0),\qquad \delta=\tfrac14,
\]
and define
\[
 c(t)=(t\sqrt{1-t},0,t)\quad(-\delta\le t\le\delta).
\]
Put $J=c([-\delta,0])$, $C=c([0,\delta])$, $q=c(-\delta)$,
$r=c(\delta)$, and $X=S\cup J$.  The set $S$ is closed and bounded:
the equation bounds $x^2+y^2$ when $z\in[0,1]$.  The arcs are
continuous semialgebraic images of compact intervals.  Thus $X$ is
compact and semialgebraic.  The last coordinate of $c$ is $t$,
so $c$ is an injective smooth parametrization with nonzero tangent.
The defining equation shows that $C\subset S$, whereas the negative
part of $J$ is disjoint from $S$.  Hence $J\cap S=\{p\}$.

There is a global smooth chart for the surface away from $p$.
On the open unit disk $D=\{(u,v):u^2+v^2<1\}$ put $a=1-u^2-v^2$.
The polynomial map
\[
 \Phi:D\longrightarrow S\setminus\{p\},\qquad
 \Phi(u,v)=(au,av,a)
\]
is bijective, with smooth inverse $(x,y,z)\mapsto(x/z,y/z)$.
Indeed every point of $S\setminus\{p\}$ has $z>0$, and its
coordinates satisfy $(x/z)^2+(y/z)^2=1-z$.  This also proves that
$\Phi$ is an embedding: differentiating the inverse identity gives
rank two for its derivative.  Thus $S\setminus\{p\}$ is an embedded
smooth surface without boundary, including the cap $(0,0,1)$.

Consider the closed chains
\[
 \{p,q\}=F_0\subset F_1=J\subset F_2=X,
 \qquad
 \{q,r\}=G_0\subset G_1=J\cup C\subset G_2=X.
\]
Their one-dimensional layers are respectively
$c((-\delta,0))$ and $c((-\delta,\delta))$.  Their
two-dimensional layers are $S\setminus\{p\}$ and $S\setminus C$.
The latter is open in $S\setminus\{p\}$, since $C$ is closed in
$S$ and contains $p$.  All the displayed layers are therefore
embedded smooth manifolds without boundary, with dimensions $0,1,2$.
They are semialgebraically locally Euclidean, so the local
triangulation and pointed Hauptvermutung argument in the proof of
Theorem~\ref{result-14} also makes them intrinsically PL regular;
the classical input is the marked compact Hauptvermutung and
triangulation theorem of Shiota~\cite[Theorems~2.1--2.2, pp.~5--6]{Shiota2010}.

Any partition coarser than both layer partitions must put $p,q$ in
one block, by $F_0$, and $q,r$ in one block, by $G_0$.  Since $r$
belongs to $F_2\setminus F_1$, this same block contains all of
$S\setminus\{p\}$.  Since $p$ belongs to $G_1\setminus G_0$, it
also contains $c((-\delta,\delta))$.  With $q$ these sets exhaust
$X$, so the common coarsening has only one block.

At a point in the interior of $J$ a sufficiently small neighborhood
in $X$ is one-dimensional.  At a point in $S$ away from $C$ it is
two-dimensional.  Invariance of dimension
\cite[Theorem~2.26, p.~126]{Hatcher2002Algebraic} therefore prevents
$X$ from being a manifold of a single dimension, and its indiscrete
partition is not weakly regular.
\end{proof}

The same formulas give a closed and bounded semialgebraic example
over every real closed field, with definable smoothness and the
intrinsic finite PL convention.  The chart and its inverse are
defined over that field, and the final contradiction uses local
definable dimension.  Ordinary compactness is asserted here only
over $\mathbb R$.  This three-dimensional construction has a
one-dimensional part; the next example retains pure dimension two.

\begin{proposition}[A pure two-dimensional compact counterexample]
\label{prop:weak-filtration-pure-r4}
There is a compact semialgebraic set $X\subset\mathbb R^4$, of
local dimension two at every point, with two closed filtrations
whose layers are embedded smooth manifolds of dimensions $0,1,2$
and whose only common coarsening is $\{X\}$.  The set $X$ is not
an intrinsic PL manifold or a topological manifold.  The two
filtrations therefore have no common weak regular coarsening.
\end{proposition}
\begin{proof}
Let $N$ be the north pole of $S^2$ and use stereographic coordinates
$(x,y)$ on $S^2\setminus\{N\}$.  Define
\[
 h(x,y)=
 \frac{\sqrt{(x-1)^2+y^2}+\sqrt{(x+1)^2+y^2}-2}
 {1+x^2+y^2},\qquad h(N)=0.
\]
The triangle inequality gives $h\ge0$, with equality in this chart
precisely on $A=[-1,1]\times\{0\}$.  If
$\rho=\sqrt{x^2+y^2}$, the numerator is at most $2\rho$,
because each distance to $(\pm1,0)$ is at most $\rho+1$.
Hence $0\le h\le2\rho/(1+\rho^2)$, which tends to zero at
$N$.  Thus $h$ is continuous and semialgebraic on $S^2$ and is
smooth away from $N$ and the endpoints $a_-,a_+$ of $A$.

In $\mathbb R^3\times\mathbb R$ set
\[
 S_0=\{(z,0):z\in S^2\},\qquad
 S_h=\{(z,h(z)):z\in S^2\},\qquad X=S_0\cup S_h.
\]
Both sheets are compact semialgebraic copies of the sphere, and
their intersection is $A\cup\{N\}$.  Each has local dimension two
at every point, so $X$ has that same local dimension everywhere.
Let $C\subset S_0$ be the great circle given by $y=0$ in the
finite chart, together with $N$.  Consider
\[
 \begin{array}{lll}
 F_0=\{a_-,a_+,N\},& F_1=A\cup\{N\},& F_2=X,\\
 G_0=\{N\},&G_1=C,&G_2=X.
 \end{array}
\]
Write $P_i=F_i\setminus F_{i-1}$ and
$Q_i=G_i\setminus G_{i-1}$, with $F_{-1}=G_{-1}=\varnothing$.
The zero-dimensional layers are finite sets; $P_1$ is the open
arc between the endpoints, and $Q_1=C\setminus\{N\}$ is a smooth
one-manifold.  The layer $P_2$ is the union of the two disjoint
smooth sheets off $A\cup\{N\}$.  The layer $Q_2$ consists of the
two open hemispheres of $S_0\setminus C$ and the smooth graph
$S_h\setminus(A\cup\{N\})$.  At each point of either two-layer,
the sheets are separated by a positive height locally.  Each layer
is consequently an embedded smooth manifold without boundary;
the same local definable-chart argument gives intrinsic PL
regularity.

The bipartite graph of nonempty intersections of the two partitions
contains the edges
\[
 P_0Q_0,\quad P_0Q_1,\quad P_1Q_1,
 \quad P_2Q_1,\quad P_2Q_2.
\]
They connect all six vertices.  Representatives are, in order, the
north pole, an endpoint of $A$, an interior point of $A$, a point
of the base great circle outside $A\cup\{N\}$, and a point of the
base sphere outside $C$.  In a common coarsening, two intersecting
layers must lie in the same block.  Connectivity therefore forces
the sole block $X$.

At an interior point of $A$ the intrinsic germ of $X$ is two disks
glued along a diameter: projection to the sphere identifies each
sheet with a disk and identifies the sheets exactly on that
diameter.  This local topological gluing description is
independent of their ambient tangency.  A link is the graph
$\Theta_4$ with two vertices joined by four edges.  Orienting the
edges in the same direction gives cellular boundary columns
$(-1,1)$.  The three differences $e_2-e_1,e_3-e_1,e_4-e_1$
form an integral basis for its kernel, so
\[
 H_2(X,X\setminus\{x\};\mathbb Z)
 \cong\widetilde H_1(\Theta_4;\mathbb Z)
 \cong\mathbb Z^3.
\]
A two-manifold has local degree-two homology $\mathbb Z$; a
manifold of another dimension has a different local homology
pattern~\cite[p.~126]{Hatcher2002Algebraic}.  Thus $X$ is not a manifold, and its sole common
coarsening is not regular.
\end{proof}

These examples concern the weak filtration convention just stated.
The canonical construction in Theorem~\ref{result-14} requires a
frontier condition and saturation of singular loci of closed unions.
The second example fails the frontier condition for its literal
layer partition: $\overline{P_1}=A$ contains $a_-$ and $a_+$ but
omits $N$ from $P_0=\{a_-,a_+,N\}$.  Thus the stronger
hypotheses have substantive content.  Requiring a coarser chain
to be a subsequence only further restricts the possible coarsenings,
so the negative conclusions remain valid under that convention
as well.


\subsection{The canonical intrinsic partition}
Intrinsic CS stratifications provide the classical background for the
coarsening question.  Friedman describes their construction from pointed
neighbourhood germs, records its earlier appearances in Handel, and credits
King's use of CS spaces in intersection-homology invariance
\cite[p.~144]{FriedmanIntrinsic2022}.
Marked coarsenings also have an antecedent: for paracompact, dimensionally
homogeneous CS spaces fully adapted to a maximal torsion-sensitive
coefficient system, Friedman's relative
construction preserves a specified closed union of positive-codimension
strata and coarsens every CS stratification satisfying the same data
\cite[Definition~6.1 and Proposition~6.4, arXiv v3 pp.~43--46]
{Friedman2023Torsion}.

Theorem~\ref{result-14} concerns the stated finite definable PL class over
an arbitrary real closed field.  It requires saturation of the
non-PL-manifold locus of every closed saturated union and allows a finite
family of closed marks, without requiring conical charts along strata.
Proposition~\ref{result-intrinsic-not-cs} shows that the coarsest partition
in this class can fail the CS condition, even for connected strata on a
compact real polyhedron.  This establishes a difference between the
admissibility classes and their outputs.  The all-field naturality below
uses compatible definable triangulation and the compact marked
o-minimal Hauptvermutung.

A different intrinsic construction is Beke's finite stratification of
real o-minimal sets by locally compact pieces
\cite[Lemma~1.8 and Proposition~1.12, author manuscript pp.~4--6]
{Beke2011Tame}. Its regularity operator tests local compactness.
For the compact interval $[0,1]$ that construction has one stratum:
every point already has a compact neighbourhood. Here the two
endpoints must instead form a saturated subset, because they are
precisely the points without a PL-manifold neighbourhood without
boundary. With no marks, the partition into $\{0,1\}$ and $(0,1)$ is admissible
and coarser than every partition satisfying our conditions.
Thus even this example distinguishes the two regularity requirements.

The coarsest frontier partition under the weaker axioms for a
closure-stable Boolean family is treated in Hua's account of cote~121
\cite[Theorems~71.1--71.2, pp.~341--342]{HuaMontpellier2026}.
\begin{proof}[Proof of Theorem~\ref{result-14}]
The empty case follows from the stated convention, so assume
$X\ne\varnothing$.
We first justify the intrinsic operator, including for noncompact
sets. A finite family of definable sets admits a model by unions
of open faces of a finite compact polyhedron. To reduce to the
compact triangulation theorem, first make all sets bounded by the
coordinate homeomorphism $t\mapsto t/\sqrt{1+t^2}$.
Take the finitely many cells comprising their bounded union in
a cell decomposition adapted to the sets. Each cell is locally
closed, and hence is the difference of its compact closure and
its compact frontier. Triangulate these finitely many closed
sets simultaneously inside their compact union. This gives the
required model from \cite[Theorem 2.2(1), p.~6]{Shiota2010},
applied to the constant function $f=0$.
All closures used subsequently are relative to $X$.

Call a point intrinsically regular if it has a definable
neighborhood definably homeomorphic to an open Euclidean ball,
without boundary. The dimension is allowed to vary between local
components of a closed set; a regular isolated point of a
mixed-dimensional set is not declared singular. In any of the above polyhedral models this
is equivalent to PL regularity. Here is the local argument that
also proves independence of the model. Suppose a point $p$ is
intrinsically regular in a union of open faces of a finite
complex $K$. A chart supplies the image $D$ of a closed Euclidean
ball, with $p$ the image of its centre and the image of its interior
a neighbourhood of $p$ in the modeled set. The set $D$ is closed and
bounded by definable compactness. Subdivide $K$ so $p$ is a vertex.
Apply \cite[Theorem 2.2(2), p.~6]{Shiota2010} to the constant
function on $|K|$, with $D$ as a marked compact definable subset.
It gives a definable self-homeomorphism $\tau$ of $|K|$
preserving each simplex, and therefore preserving
our union of open faces and fixing $p$, such that $\tau^{-1}(D)$
is polyhedral. The relative Hauptvermutung
\cite[Theorem 2.1, p.~5]{Shiota2010}, applied with the centers
marked, gives a PL homeomorphism from a standard ball to this
polyhedron carrying its center to $p$. Since $\tau^{-1}(D)$
contains a relative neighborhood of $p$, the ambient germ is
PL Euclidean. The converse follows because a PL chart is definable.
Consequently this regularity is preserved by every definable
homeomorphism, even when the whole space is not locally compact.
Here compact over a non-Archimedean field means closed and bounded.

For a definable $Y$, the PL germ is constant on the relative
interior of each face in a finite model: it is the product of
that open face with its transverse conical germ. Its singular
locus is therefore a finite union of open faces and is definable.
The regular locus is open by definition, so the singular locus
is closed relative to $Y$. For nonempty $Y$ its dimension is smaller
than that of $Y$, since interiors of the maximal included faces are
regular and dense; $\operatorname{Sing}_{\rm PL}(\varnothing)
=\varnothing$. This verifies every regularity input needed below.

For completeness we give the finite existence and coarsening
argument with this intrinsic operator. Start with a compatible
finite triangulation of $(X;\Phi)$ and its open-face partition.
Freeze strata in decreasing dimension. Suppose strata of dimension
greater than $k$ are already frozen and the unfrozen remainder
is closed in $X$, partitioned with the frontier condition and
compatibly with the closures of all frozen strata. For every union
$A_J$ of closures of frozen strata and every current $k$-stratum
$S$, partition $S$ according to the finitely many sets
$S\cap\operatorname{Sing}_{\rm PL}(A_J)$. Remove their relative
boundaries and freeze the remaining open pieces, taking components
if the connected version is wanted. The removed boundaries have
dimension less than $k$. Every current $k$-stratum is open in the
unfrozen remainder, and hence so are the pieces just frozen. The
new remainder is therefore closed. Triangulate this remainder
compatibly with all frozen closures and the previous pieces. It
cannot accumulate onto a frozen stratum, because it is closed and
disjoint from the frozen strata. This maintains the frontier
condition and the invariant needed for the next dimension. Finiteness
of each stage and the dimension decrease make the process finite.
In the connected version, components are taken before freezing,
and the closures of these individual components are included
among the marks for the next lower-dimensional refinement.
Thus no assertion that taking components preserves an already
finished frontier partition is required.

For a final $k$-stratum $S$ contained in a closed saturated $Y$,
the germ of $Y$ along $S$ is either the regular germ of $S$ alone
or the germ of a union $A_J$ of closures of higher-dimensional
strata of $Y$. Indeed the frontier condition makes each such
closure either contain all of $S$ or miss it, and finitely many
closures missing a given point can be avoided locally. In the
second case these higher strata had already been frozen when $S$
was selected, so the construction made
$\operatorname{Sing}_{\rm PL}(A_J)\cap S$ either empty or all of
$S$. Thus the singular locus of every closed saturated union is
saturated, not merely the singular locus of each stratum closure.

For uniqueness, write $\mathcal B_P$ for the finite Boolean
algebra of unions of strata of one such partition. This algebra is
closed under closure and under $Y\mapsto\operatorname{Sing}_{\rm PL}(Y)$
on closed members. For two admissible partitions $P,Q$, the
intersection $\mathcal B_P\cap\mathcal B_Q$ has those properties
and contains $\Phi$. Its atoms are locally closed. In fact, for an
atom $A$ its frontier $\overline A\setminus A$ and the closure of
that frontier belong to the algebra. If this closed frontier met
$A$, atomicity would imply that it contained $A$, contrary to
the strict frontier dimension inequality. Thus $A$ is open in
$\overline A$. The singular locus of $\overline A$ is saturated
and cannot contain the dense atom $A$, so $A$ is a PL manifold.
The atoms and their closures satisfy the frontier condition.

Purity is also preserved in this intersection. For every closed
saturated $Y$, the closed definable set
\[
 D_k(Y)=\{y\in Y:\dim_yY\ge k\}
\]
is saturated: its value on a stratum $S$ is determined by the
maximum dimension of a stratum whose closure contains $S$.
Thus both algebras, and their intersection, are closed under the
operators $D_k$. For an atom $A$, each $A\cap D_k(\overline A)$
is empty or all of $A$, so its manifold dimension is constant.
In the connected version the algebras are also closed under
individual connected components. A connected stratum lies in one
component of every saturated set containing it; components of such
sets are therefore saturated. This property passes to the algebra
intersection and forces its atoms to be connected.

Choose an admissible partition with the smallest possible number
of strata, which exists by the first construction and the
well-ordering of the positive integers. Its common coarsening with
any other admissible partition, constructed by the intersection
of algebras, cannot have fewer strata. It is therefore this same
partition, establishing coarseness and uniqueness. This proves the assertion for the intrinsic PL operator.

A definable homeomorphism of marked data preserves closure, local
dimension, the PL regularity operator and connected components.
It therefore transports the unique coarsest partition to the
unique coarsest partition on the target, proving naturality.
The frontier of a stratum has strictly smaller dimension, so
$X_{\le k}$ contains the closure of every stratum it contains
and is closed. Each $k$-stratum is closed in
$L_k=X_{\le k}\setminus X_{\le k-1}$, because its frontier
has smaller dimension. There are finitely many such strata,
so each is also open in $L_k$. Thus $L_k$, and every union
of selected $k$-strata, is a pure $k$-manifold or empty.
Finally, each $Y_j$ is saturated, so the
differences in the displayed refinement are exactly
\[
 (Y_j\setminus Y_{j-1})\cap
 (X_{\le k}\setminus X_{\le k-1}).
\]
They are unions of selected $k$-strata, hence pure $k$-manifolds
or empty. All sets and their order are canonically specified.
\end{proof}

\begin{proposition}[Singularity saturation does not force conical strata]
\label{result-intrinsic-not-cs}
There is a compact finite polyhedron $X\subset\mathbb R^5$ whose
non-PL-manifold locus is a PL circle $S$, such that $X\setminus S$ is a
connected PL surface and both versions of
Theorem~\ref{result-14}, with no marks, give the partition
\[
                 \{X\setminus S,S\}.
\]
The pointed ambient neighbourhood germ is not locally constant along
$S$.  In particular this partition is not a CS stratification, where
CS has the locally conelike meaning of
\cite[Definition 1.10, pp.~5--6]{ChataurSaralegiTanre2019},
with codimension-one strata allowed.  Thus the theorem's saturation condition does not imply
local conical triviality, even with connected strata and compact
real polyhedra.
\end{proposition}
\begin{proof}
Begin with a PL circle $S$ and three PL closed discs $D_1,D_2,D_3$.
Identify each boundary $\partial D_i$ with $S$ by a PL homeomorphism.
Choose a proper closed arc $I\subset S$ with two distinct endpoints.
Take a PL two-sphere $P$ containing a PL embedded arc $J$, and identify
$J$ with $I$.  Choose finite triangulations making all these
identifications simplicial.  The quotient is a compact finite
two-dimensional polyhedron.  Outside $S$ it is the disjoint union of
the three open discs and the slit sphere $P\setminus J$, each a PL
surface without boundary.

We connect these four regular pieces without changing any germ along
$S$.  Choose three pairwise disjoint small closed PL discs
$B_{0i}\subset P\setminus J$ and one small closed PL disc
$B_i\subset\operatorname{int}D_i$ for each $i=1,2,3$.
Remove their interiors and attach an annulus between
$\partial B_{0i}$ and $\partial B_i$, using PL boundary
homeomorphisms.  These attachments take place away from $S$.
The resulting space, still denoted by $X$, is a compact finite
two-dimensional polyhedron.  Its complement
$A=X\setminus S$ is a connected PL surface without boundary: each
punctured disc is joined to the punctured slit sphere, and the
annuli fill all newly created boundary circles.  Removing these
finitely many discs does not disconnect the slit sphere, as paths
can be rerouted along their boundary circles.

For completeness, $X$ can be realized as a finite polyhedron in
$\mathbb R^5$.  Subdivide the finite gluing complex to an abstract
simplicial complex and place its vertices at distinct rational
points of the moment curve
$t\mapsto(t,t^2,t^3,t^4,t^5)$.  Any six of these vertices are
affinely independent.  Since two simplices of dimension at most two
have at most six vertices in their union, their convex hulls meet
exactly in their common face.  The affine realization is therefore
an embedding.  In particular $X$ is compact and semialgebraic.

Write $\Theta_r$ for the graph with two vertices and $r$ edges
joining them.  The links at points of $S$, unaffected by the
annulus attachments, are
\[
\begin{array}{c|c|c}
\text{position of }x&\text{link at }x&
\operatorname{rank}H_2(X,X\setminus\{x\};\mathbb Z)\\ \hline
x\in S\setminus I&\Theta_3&2\\
x\in\partial I&\Theta_3\vee S^1&3\\
x\in\operatorname{int}_S I&\Theta_5&4.
\end{array}
\]
Indeed, the three discs contribute three half-disc pages along
$S$.  At an interior point of $I$, the sphere contributes two
more pages.  At an endpoint of $I$, its link circle meets the
three-disc link at only one of the two spine directions.  The
displayed ranks follow from
$H_2(X,X\setminus\{x\};\mathbb Z)
\cong\widetilde H_1(\operatorname{lk}_X x;\mathbb Z)$.
They exclude a manifold germ of any dimension at every point of
$S$.  All points outside $S$ are regular by construction, so
$\operatorname{Sing}_{\rm PL}(X)=S$.

Both $A$ and $S$ are connected manifolds of a single dimension,
$\overline A=X$, and $S$ is closed.  The only closed unions in the
two-stratum partition are $\varnothing,S,X$, whose singular loci
are respectively $\varnothing,\varnothing,S$.
Thus this partition is admissible in both classes of
Theorem~\ref{result-14}.  Conversely, every admissible partition
must make $\operatorname{Sing}_{\rm PL}(X)=S$ saturated, so each
of its strata lies in $S$ or in $A$.  The displayed partition is
therefore coarser than every admissible partition in either class.

In a conical chart $U\times cL$ along a stratum, sufficiently
nearby points of $U\times\{\text{vertex}\}$ have homeomorphic
pointed ambient neighbourhoods, by translating inside a smaller
coordinate ball in $U$.  Their local homology groups consequently
agree.  The ranks in the table change in every neighbourhood in
$S$ of either endpoint of $I$.  Hence $S$ cannot be a CS stratum.
The finer filtration
$X_0=\partial I$, $X_1=S$, $X_2=X$ does give a CS structure:
the open arcs have the displayed product-page charts and the
two vertices have their polyhedral cone charts.  Thus the
difference is between partition conditions on the same compact
polyhedron, not a failure of the space to admit a CS structure.
\end{proof}


\subsection{Coarsest partitions for embedded smooth regularity}
\label{subsec:coarsest-embedded-smooth-partitions}

In this subsection $X\subset\mathbb R^N$ is an arbitrary semialgebraic
set, and $\Phi$ is a finite family of semialgebraic subsets of $X$.
Neither closedness nor local closedness of $X$ is assumed.  All closures
used to define partitions are relative to $X$.
For a semialgebraic set $A\subset\mathbb R^N$, let $\operatorname{Sm}(A)$
be the set of points at which $A$ agrees locally with an embedded
$C^\infty$ manifold without boundary, and put
\[
 \Sigma(A)=A\setminus\operatorname{Sm}(A).
\]
The local manifold dimension may vary on different components.
Thus $\operatorname{Sm}(A)$ uses embedded smoothness at every local
dimension.  In the semialgebraic setting these local manifolds are Nash
manifolds; see \cite[pp.~629--630]{Fernando2018}.

A finite semialgebraic partition is \emph{cellular} if every stratum is
locally closed in $X$ and its closure is a union of strata.  It is
\emph{$F$-regular} if $\Sigma(Y)$ is a union of strata for every union
$Y$ of strata that is closed in $X$.  This requirement includes all
closed unions, not just individual stratum closures.  Adaptation to
$\Phi$ means that every member of $\Phi$ is a union of strata.  We order
partitions by coarsening: $P$ is coarser than $Q$ when every stratum of
$Q$ is contained in a stratum of $P$.

The construction must control the singularities of every closed
union of strata at once. We first choose smooth pieces in decreasing
dimension, testing the finitely many relevant unions before choosing
the next pieces. To obtain coarseness, we intersect the Boolean
algebras of saturated sets. Stability under closure and the singular
locus makes this common coarsening admissible; minimizing the number
of strata then gives the coarsest partition.

We first record the exact classical semialgebraic input.

\begin{lemma}[The smooth locus in every local dimension]
\label{lem:all-dimensional-semialgebraic-smooth-locus}
For every semialgebraic $A\subset\mathbb R^N$, the set
$\operatorname{Sm}(A)$ is semialgebraic and relatively open in $A$.
Its complement $\Sigma(A)$ is relatively closed and nowhere dense.
If $A\ne\varnothing$, then $\dim\Sigma(A)<\dim A$.
These assertions also hold when $A$ is nonpure or nonclosed.
\end{lemma}

\begin{proof}
Stasica's theorem states separately for each integer $k$ that the set
$A^{(k)}$ of smooth points of dimension $k$ is semialgebraic
\cite[Theorem~4, p.~152]{Stasica2003Smooth}.  Consequently
\[
 \operatorname{Sm}(A)=\bigcup_{k=0}^N A^{(k)}
\]
is semialgebraic.  This is also the smooth-locus convention
$\operatorname{Sth}(A)$ and $\operatorname{Sth}_k(A)$ in
\cite[\S2.A, pp.~635--636]{Fernando2018}.  It is relatively open by
the local definition.  No global maximal-dimension restriction is
imposed on this union.

The same semialgebraicity conclusion follows from Koike--Shiota's
theorems on the nonsmooth points of fibres, applied to the constant
map $A\to\{0\}$
\cite[Theorems~2.1--2.2 and Corollary~2.4, p.~956]{KoikeShiota2007}.
Their Remark~0.1 on p.~953 identifies the semialgebraic
$C^\infty$ and analytic manifold conventions, and their argument
on pp.~961--962 includes sets that are not locally closed.

For the dimension assertion choose a finite pure-dimensional Nash
frontier partition of $A$.  Such partitions, adapted to any prescribed
finite family, are classical
\cite[Theorem~1.3 and the frontier refinement following it,
pp.~3--4]{Coste2007}.  Every stratum of maximal dimension is open in
$A$: the closure of any other stratum cannot meet it, by the frontier
condition and the strict drop of dimension on proper frontiers.
These maximal-dimensional strata lie in $\operatorname{Sm}(A)$.
Their complement has smaller dimension, so the same is true of
$\Sigma(A)$.

This global dimension drop also gives the required local density even
for a nonpure set.  If $\Sigma(A)$ had nonempty relative interior,
it would contain a nonempty relatively open semialgebraic subset $U$
of $A$, obtained by intersecting with a sufficiently small open ball.
Smoothness is local, so
$\Sigma(U)=U\cap\Sigma(A)=U$, contradicting the dimension drop
applied to $U$.  Thus $\Sigma(A)$ is nowhere dense.  The argument
applies to arbitrary semialgebraic subsets, so a set closed only in
$X$ is within its scope.
\end{proof}

\begin{theorem}[Coarsest smooth regular partition]
\label{thm:coarsest-embedded-smooth-partition}
There is a unique coarsest finite semialgebraic cellular $F$-regular
partition of $X$ adapted to $\Phi$.  Every one of its strata is
smooth.  The assertion makes no connectedness, local triviality, or
equisingularity requirement.
\end{theorem}

\begin{proof}
If $X$ is empty, its empty partition has all the asserted properties.
Suppose henceforth that $X$ is nonempty.  We first construct one
admissible partition and then prove the coarsest property.

Choose a finite pure-dimensional smooth frontier partition adapted
to $\Phi$.  We freeze pieces in decreasing order of dimension.  At
each stage keep a closed semialgebraic remainder $R\subset X$,
disjoint from the already frozen pieces, with a smooth frontier
partition adapted to $\Phi$, to the earlier pieces, and to every
closure in $X$ of a frozen piece.  Initially $R=X$.

Let $k=\dim R$, and let $S$ be one of the current $k$-dimensional
strata.  The frontier condition makes $S$ open in $R$.  All previously
frozen strata have dimension greater than $k$.  For every subset $J$
of the finite set of frozen strata, form the closed set
\[
 A_J=\bigcup_{T\in J}\overline T^{\,X}.
\]
The finitely many sets $\Sigma(A_J)\cap S$ are semialgebraic.
Remove all their relative boundaries in $S$.  Each of these
boundaries has dimension less than $k$, by the semialgebraic
dimension theorem in a $k$-manifold.  On the complement, each
membership pattern in the sets $\Sigma(A_J)\cap S$ is an open
semialgebraic subset of $S$.  Freeze every nonempty such piece.
These pieces are smooth of dimension $k$, and they are open in $R$.

Perform this operation simultaneously on all current $k$-strata.
The new remainder is closed in $R$, hence closed in $X$, and has
dimension less than $k$.  Choose on it a smooth frontier partition
adapted to the preceding partition and to all frozen closures.
The new remainder, being closed and disjoint from the frozen pieces,
cannot accumulate at a point of a frozen piece.  Within a given
$k$-stratum the newly frozen open pieces do not lie in each other's
closures.  Adaptation to their closures supplies all future frontier
relations.  The same statements hold between distinct $k$-strata by
the preceding frontier condition.  These observations preserve the
stated induction invariant.  The dimension decreases at each stage,
so after finitely many stages the frozen pieces give a finite smooth
frontier partition $P$ adapted to $\Phi$.

We verify the condition on every closed saturated union.  Let $S$ be
a final stratum of dimension $k$ and let $Y$ be a union of final
strata that is closed in $X$.  If $S$ is disjoint from $Y$, it is also
disjoint from $\Sigma(Y)$.  Suppose $S\subset Y$.  Collect the
strata $T\subset Y$ of dimension greater than $k$ for which
$S\subset\overline T^{\,X}$.  The frontier condition implies that
these are precisely the strata other than $S$ whose closures can
meet $S$.  At each point of $S$, a neighborhood avoids the finitely
many other stratum closures.  If the collection is empty, the germ
of $Y$ there is the smooth germ of $S$.  Otherwise the germ of $Y$
is the germ of the union $A_J$ of the collected closures: closedness
of $Y$ ensures that each of those closures lies in $Y$, and their
union already contains $S$.  All the collected strata had been
frozen before $S$ was constructed.  The membership test at that
stage therefore makes $\Sigma(A_J)\cap S$ either all of $S$ or
empty.  Locality of smoothness gives the same conclusion for
$\Sigma(Y)\cap S$.  Thus $P$ is $F$-regular.

For coarseness, associate to each admissible partition $P$ the finite
Boolean algebra $\mathcal B_P$ of all its saturated subsets.
Cellularity implies that this algebra is closed under closure in
$X$, because closure commutes with finite unions.  Its closed
members are also stable under $Y\mapsto\Sigma(Y)$, by
$F$-regularity.  If $P$ and $Q$ are adapted admissible partitions,
their intersection
\[
 \mathcal B=\mathcal B_P\cap\mathcal B_Q
\]
has the same properties and contains $\Phi$.  Its nonempty atoms
form a finite semialgebraic common coarsening $R$.

We check local closedness of those atoms.  For an atom $A$, both
$\overline A^{\,X}$ and
$F=\overline A^{\,X}\setminus A$ belong to $\mathcal B$, as does
$\overline F^{\,X}$.  If $A$ met $\overline F^{\,X}$, atomicity
would imply $A\subset\overline F^{\,X}$.  But the semialgebraic
frontier inequality gives
\[
 \dim\overline F^{\,X}=\dim F<\dim A,
\]
a contradiction; see \cite[Proposition~1.9, p.~5]{Coste2007}.
Hence $A=\overline A^{\,X}\setminus\overline F^{\,X}$ is open
in its closure.  Each atom closure is a union of atoms, so $R$ is
cellular.  The stability of $\mathcal B$ under $\Sigma$ on closed
members makes $R$ $F$-regular.

Moreover $\Sigma(\overline A^{\,X})\cap A$, being in
$\mathcal B$, is empty or all of $A$.  The latter is impossible:
$A$ is nonempty and open in its closure, while the nonsmooth locus
is nowhere dense there by
Lemma~\ref{lem:all-dimensional-semialgebraic-smooth-locus}.
Thus $A$ is open in the smooth locus of its closure, and is itself
an embedded smooth manifold, with local dimension allowed to vary.

Choose an adapted admissible partition $P$ with the smallest number
of strata; existence of such a minimum follows from the construction
and the well-ordering of the nonnegative integers.  The common
coarsening of $P$ and any other adapted admissible $Q$ just constructed
cannot have fewer strata than $P$.  Since it is coarser than $P$,
it must equal $P$.  Consequently $P$ is coarser than every $Q$.
Two partitions with this property coarsen each other and are equal,
which proves uniqueness.
\end{proof}

\begin{theorem}[Connected and pure-dimensional versions]
\label{thm:coarsest-smooth-partition-variants}
The preceding existence and uniqueness assertion holds separately
when each stratum is required to be connected, when each stratum
is required to have a single local dimension, and when both
requirements are imposed.  A connected smooth stratum automatically
has a single local dimension.
\end{theorem}

\begin{proof}
The construction in
Theorem~\ref{thm:coarsest-embedded-smooth-partition} already produces
pure-dimensional pieces.  For the connected version, freeze the
connected components of each open piece instead.  There are finitely
many such components and they are semialgebraic
\cite[Proposition~1.6, p.~5]{Coste2007}; they are open in the smooth
piece.  Refining the lower-dimensional remainder by their closures
preserves the same frontier and closed-remainder arguments.
Thus each of the stated classes is nonempty.

For connected strata, each connected component of a saturated set
is saturated: a connected stratum lies wholly in one component.
Consequently $\mathcal B_P$ is closed under taking connected
components, and so is $\mathcal B_P\cap\mathcal B_Q$ for two
connected partitions.  A disconnected atom would have a nonempty
proper connected component in that algebra, which contradicts
atomicity.  The common coarsening constructed above therefore has
connected strata.  Its smooth strata have locally constant manifold
dimension, so connectedness also makes that dimension constant.

For the class with pure-dimensional strata, use the local dimension sets
\[
 D_k(Y)=\{y\in Y:\dim_yY\ge k\}
\]
for each closed saturated $Y$.  For a pure-dimensional frontier
partition $P$, one has the exact formula
\[
 D_k(Y)=
 \bigcup_{\substack{T\in P,\ T\subset Y\\\dim T\ge k}}
       \overline T^{\,X}.
\]
Indeed the local dimension at a point of a stratum $S\subset Y$
is the maximum of the dimensions of the strata $T\subset Y$ whose
closures contain $S$.  A neighborhood avoids every other stratum,
and every neighborhood meeting such a pure-dimensional $T$ meets
an open subset of $T$ of its full dimension.  The displayed finite
union is closed in $X$ and saturated.  It shows directly that each
$\mathcal B_P$, and hence the intersection algebra for two pure
partitions, is stable under $D_k$ on closed members.

For an atom $A$ of that intersection, each set
$A\cap D_k(\overline A^{\,X})$ is empty or all of $A$.
The earlier proof makes $A$ open in the smooth locus of its closure,
so the local dimension of $\overline A^{\,X}$ on $A$ is the
manifold dimension of $A$.  The preceding all-or-none tests for all
$k$ force that dimension to be constant.  Thus the common coarsening
also belongs to the pure-dimensional class.  Applying the same
minimum-number-of-atoms argument in each class proves the assertions.
\end{proof}


\subsection{The naturality class of embedded smooth partitions}
\label{subsec:embedded-smooth-naturality}

\begin{proposition}[Smooth-locus preservation and the cusp]
\label{prop:embedded-smooth-partition-naturality}
Let $X\subset\mathbb R^N$ and $X'\subset\mathbb R^{N'}$ be
semialgebraic sets, with finite marked families $\Phi$ and $\Phi'$.
A semialgebraic homeomorphism $h:X\to X'$ that carries the marks to
the corresponding marks and satisfies
\[
 h(\Sigma(Y))=\Sigma(h(Y))
 \quad\text{for every semialgebraic }Y\subset X
 \text{ closed in }X
\]
carries each of the coarsest partitions in
Theorems~\ref{thm:coarsest-embedded-smooth-partition}
and~\ref{thm:coarsest-smooth-partition-variants} to its counterpart
on $X'$.  In particular this holds for the restrictions of ambient
Nash diffeomorphisms.  It fails for arbitrary semialgebraic
homeomorphisms, already for the interval and the compact cusp.
\end{proposition}

\begin{proof}
A homeomorphism transports relative closures, local closedness,
connectedness, and the frontier condition.  A semialgebraic
homeomorphism also preserves local dimension, by applying
semialgebraic invariance of dimension to corresponding neighborhoods
\cite[Proposition~1.8, p.~5]{Coste2007}.  The displayed assumption
therefore transports the $F$-regularity condition for every closed
saturated union.  Its inverse has the same property, because every
closed semialgebraic subset of $X'$ is the image of one in $X$.
Thus transport gives mutually inverse, order-preserving bijections
between the relevant classes of adapted partitions.  It carries
their unique coarsest elements to each other.  Preservation of the
smooth-locus operator by Nash diffeomorphisms is precisely
\cite[Lemma~2.3, p.~636]{Fernando2018}.

For failure of general homeomorphism invariance, let
\[
 I=[-1,1],\qquad
 C=\{(t^2,t^3):-1\le t\le1\},\qquad
 \gamma(t)=(t^2,t^3).
\]
The map $\gamma:I\to C$ is a semialgebraic homeomorphism.
It is bijective, and its inverse is $y/x$ away from the origin,
extended by zero at the origin.  Continuity there follows from
$|\gamma^{-1}(x,y)|=\sqrt{x}$ on $C$.
The origin is an embedded smooth point of $I$, but not of $C$.
Indeed $C$ has local dimension one at the origin, and its tangent
cone there is the nonnegative horizontal half-line.  For example
\[
 \frac{\gamma(t)}{\|\gamma(t)\|}
   =\frac{(1,t)}{\sqrt{1+t^2}}
   \longrightarrow(1,0)
 \quad (t\longrightarrow0,\ t\ne0).
\]
Every nonzero tangent-cone direction is therefore positive horizontal,
and every positive horizontal vector occurs by rescaling.
An embedded $C^1$ one-manifold without boundary has a full tangent
line, so this half-line excludes such a manifold germ.

The nonsmooth loci of the two compact sets are exactly
\[
 \Sigma(I)=\{-1,1\},\qquad
 \Sigma(C)=\{\gamma(-1),(0,0),\gamma(1)\}.
\]
The endpoints have one-sided manifold germs and hence are nonsmooth
under the convention without boundary.  Away from those endpoints
and the cusp origin, $\gamma$ is a regular Nash parametrization of
an embedded open arc.

In either case the partition into the displayed finite nonsmooth
set and its complement is pure-dimensional, cellular, and
$F$-regular.  Its only closed saturated unions are the empty set,
the finite zero-dimensional stratum, and the whole space; their
nonsmooth loci are respectively empty, empty, and that finite
stratum.  Every admissible partition must saturate the nonsmooth
locus of the whole space, so it refines this two-piece partition.
These are therefore the coarsest pure-dimensional partitions.
The open interval is one stratum, while the cusp origin belongs
to a zero-dimensional stratum containing both endpoints as well.
Since the finite strata have two and three
points respectively, no homeomorphism can identify the two
partitions.  Separating the finite points and the connected open arcs
gives admissible partitions with respectively three and five strata:
every closed saturated union is a union of closed subarcs and isolated
vertices, and its nonsmooth locus is a union of vertices.  A connected
partition refining either two-piece partition also refines its
component partition.  Hence these are the two coarsest connected
versions.
\end{proof}



\section{Finite variation and spatial regularization}
The next construction turns a homeomorphism into the quotient of two uniformly Lipschitz maps, while preserving every prescribed fixed point.

\begin{theorem}[Finite spatial regularization with fixed marks]\label{result-15}
Let $M$ be an o-minimal real closed field, let $X$ be a compact
finite PL polyhedron, and fix finitely many marked subpolyhedra.
Let $P=[0,1]^k$ and let $U_p:X\to X$ be a jointly continuous
definable family of homeomorphisms preserving the marks and
equal to the identity for $p\in\partial P$. There are jointly
continuous definable families of homeomorphisms $r_p,s_p$,
preserving the marks and equal to the identity on $\partial P$,
such that
\begin{equation}
 U_p=s_pr_p^{-1},
 \label{equation-1}
\end{equation}
and both families have one common uniform spatial Lipschitz
bound. If $U$ fixes a marked subpolyhedron pointwise, both factors
fix it pointwise. The construction uses any fixed finite
target dimension.

\end{theorem}

\begin{proof}
Use maximum norms.

\textbf{Finite variation and its stability.}\quad
For a definable family of continuous scalar functions on
$[0,1]$, o-minimal monotonicity gives an ordinary finite
uniform bound $N$ for the number of monotonicity intervals.
For each member $f$ and $0\le t\le1$ its variation is exactly
\begin{equation}
 V_f(t)=
 \max_{0=t_0\le t_1\le\cdots\le t_N=t}
       \sum_{\nu=1}^N|f(t_\nu)-f(t_{\nu-1})|.
 \label{equation-2}
\end{equation}
Repeated partition points are allowed. A partition by the
monotonicity intervals realizes the ordinary variation, so
the displayed finite maximum is exact. It is definable,
and its maximum exists by definable compactness. Comparison
on the same partition gives
\begin{equation}
 |V_f(t)-V_g(t)|\le2N\|f-g\|_\infty.
 \label{equation-3}
\end{equation}
For a finite vector of $d$ scalar functions, sum their
variations; the comparison constant is then $2dN$.
Variation is additive on adjacent intervals, and hence
\begin{equation}
 \|f(v)-f(u)\|_\infty\le V_f(v)-V_f(u)
 \quad(0\le u\le v\le1).
 \label{equation-4}
\end{equation}
The variation depends jointly continuously on the parameters
and on $t$. To check the endpoint variable, apply \eqref{equation-3}
to the rescaled lines $z\mapsto f(tz)$ on the fixed unit
interval. Definable uniform continuity on a compact product
gives the required uniform convergence of these lines.
These arguments work over every o-minimal real closed field.

\textbf{The mixed coordinate-clock induction.}\quad
Let $C=[0,1]^n$, and let
$E:C\times P\to[0,1]^a$ be any continuous definable vector
map. Begin with
\[
 \chi^0_p=\mathrm{id}_C,\qquad
 F^0(x,p)=(x,E(x,p)).
\]
At stage $i$, suppose the vector
$F^{i-1}=(\chi^{i-1},E\chi^{i-1})$ is uniformly Lipschitz
in every already processed coordinate $x_j$, $j<i$.
Fix the other coordinates and $p$ as a profile $b$, and
write $f_b(t)$ for its $i$th coordinate line. Define
\begin{equation}
 A_b(t)=t+V_{f_b}(t),\qquad
 D_b=A_b(1),\qquad
 \xi_b(u)=A_b^{-1}(D_bu).
 \label{equation-5}
\end{equation}
The function $A_b$ is continuous and strictly increasing;
$A_b(v)-A_b(u)\ge v-u$. Its inverse exists and depends
jointly continuously on the profile. Replace $x_i$ by
$\xi_b(x_i)$, leaving all other coordinates and $p$ fixed,
to obtain a parameter-preserving homeomorphism $\Psi_i$.
Its inverse changes that coordinate to $A_b(x_i)/D_b$.
Both maps preserve each coordinate face. Set
\[
 \chi^i=\chi^{i-1}\Psi_i,\qquad
 F^i=F^{i-1}\Psi_i.
\]

Here is the needed mixed estimate. Let two profiles $b,c$
differ in a previously processed coordinate by $\delta$,
with old Lipschitz bound $K$. If the vector has $d$
components and its line monotonicity bound is $N$, put
$C_0=2dN$. Formula \eqref{equation-3} gives
\[
 \sup_t|A_b(t)-A_c(t)|\le C_0K\delta,\qquad
 |D_b-D_c|\le C_0K\delta.
\]
For $h=\xi_b(u)$ and $h'=\xi_c(u)$,
\[
 |A_c(h)-A_c(h')|\le2C_0K\delta.
\]
Using \eqref{equation-4} to control the change along the old line,
rather than assuming that line was Lipschitz, yields
\begin{equation}
 \|f_b(h)-f_c(h')\|_\infty
       \le(1+2C_0)K\delta.
 \label{equation-6}
\end{equation}
For a change in the new coordinate $u$,
\begin{equation}
 \|f_b(\xi_b(u))-f_b(\xi_b(u'))\|_\infty
       \le D_b|u-u'|.
 \label{equation-7}
\end{equation}
The heights $D_b$ have a uniform bound by finite variation
and bounded coordinate ranges. Thus the induction preserves
all earlier coordinate estimates and adds the new one.
After $n$ stages, changing coordinates one at a time proves
uniform global spatial Lipschitz bounds for both $\chi$
and $E\chi$. The evolving domain map is included in the
vector throughout; omitting it would not prove a bound
for the first factor.

If all vector components have standard bounded ranges,
these final constants have standard integer bounds.
Each line variation is bounded by its range times the
ordinary finite monotonicity count. Equations \eqref{equation-6}--\eqref{equation-7}
then involve only finitely many ordinary numerical operations.
In particular this conclusion holds on a fixed standard
unit-cube model. No integer is chosen larger than an
arbitrary infinite field element.

If $E_p=\mathrm{id}_C$ for $p\in\partial P$ and $a=n$,
then on this boundary the evolving vector is $(x,x)$.
Each processed line has variation $2t$, so $A_b(t)=3t$,
$D_b=3$, and $\xi_b(u)=u$. Hence every clock fixes the
parameter boundary pointwise.

\textbf{A PL cubical model with all marks.}\quad
Triangulate $X$ compatibly with the finite marks and number
the vertices by $1,\ldots,v$. For the standard realization
of the abstract complex $K$, set
\begin{equation}
 Q=\{z\in[0,1]^v:
       \operatorname{supp}(z)\in K,\ \max_i z_i=1\}.
 \label{equation-8}
\end{equation}
It is the finite union of the coordinate faces
$z_i=1$, $z_j=0$ for $j\notin\sigma$, where
$\sigma\in K$ and $i\in\sigma$. A marked subcomplex is
the corresponding union of such faces.

There is a PL homeomorphism $b:X\to Q$ preserving these
marks. On the barycentric subdivision of the standard
realization, send the barycentre of $\sigma$ to its
characteristic vector $\mathbf1_\sigma$ and extend
affinely over chains. These chains give the usual
monotone-coordinate triangulations of the displayed
coordinate faces, and agree on their intersections.
They triangulate all of $Q$. Thus this is a finite
PL homeomorphism. One may check the underlying topology
also by $z\mapsto z/\sum z_i$ and
$\lambda\mapsto\lambda/\max\lambda_i$, but those rational
maps are not the conjugacy used here. The map $b$
and its inverse are Lipschitz on their finite polyhedra.

\textbf{Extending a vector map, not homeomorphisms.}\quad
Conjugate $U$ to $U^Q$ on $Q$. On the closed polyhedral
subcomplex
\[
 Z=(Q\times P)\cup(C\times\partial P),
 \qquad C=[0,1]^v,
\]
prescribe $U^Q$ on the first piece and the identity on the
second. The prescriptions agree and define a continuous
definable map $E_0:Z\to C$.

Take a finite compatible triangulation of $C\times P$ and
barycentrically subdivide so that $Z$ is full. On a simplex,
let $\alpha(w)$ be the sum of its barycentric weights at
vertices lying in $Z$. If $\alpha>0$, let $q(w)$ be the
normalized barycentric sum of those vertices. Fullness
ensures $q(w)\in Z$. For a fixed $c\in C$, put
\begin{equation}
 E(w)=\alpha(w)E_0(q(w))+(1-\alpha(w))c
      \quad(\alpha(w)>0),\qquad E(w)=c\quad(\alpha(w)=0).
 \label{equation-9}
\end{equation}
The formulas agree on shared faces and extend $E_0$.
They are continuous at $\alpha=0$ because $E_0$ is bounded;
their values lie in the convex cube $C$.

Apply the coordinate clocks to this vector map. They
preserve every coordinate face, hence $Q$ and every mark.
Their restrictions to $Q$ are homeomorphisms. Define
\begin{equation}
 r_p=b^{-1}\chi_pb,\qquad s_p=U_pr_p.
 \label{equation-10}
\end{equation}
Since $E=U^Q$ on $Q$, the spatial bounds for $\chi$ and
$E\chi$ give the required bounds for $r,s$ after fixed
PL conjugation. They satisfy \eqref{equation-1}, preserve all marks,
and are identity on the parameter boundary. If $U$ fixes a mark pointwise, then on every coordinate face of
that mark the extension $E$ is the identity.  Inductively, the evolving
vector on the face is $(\mathrm{id},\mathrm{id})$.  Along any coordinate
free on that face its variation is $2t$, hence its normalized clock is
$(t+2t)/3=t$.  A coordinate fixed by the face stays at its endpoint.
Every clock therefore fixes the face pointwise.  Thus $r=s=\mathrm{id}$
on every pointwise fixed mark.
The total maps and their inverses are jointly continuous;
this also follows from the continuous bijection of compact
definable products. This proves the factorization.

\end{proof}


\section{Why a homeomorphic parameter change is insufficient}
The following explicit family distinguishes spatial regularity from parameter regularity.

\begin{theorem}[A based semialgebraic family with no homeomorphic Hausdorff-Lipschitz reparameterization]\label{result-16}
Work over the real field. Put
\[
 C=\partial[-1,1]^2,\quad B=C^d\ (d\ge1),\quad
 X=B\times[-1,1],\quad P=[-1,1]^2.
\]
There is an explicit jointly continuous semialgebraic family
$U_p:X\to X$ with the following properties.
\begin{enumerate}
\item Every $U_p$ and its inverse are spatially Lipschitz with one
uniform Euclidean constant strictly less than $2$. They are the
identity on both outer half-interval collars in the last coordinate.
\item The family is the identity at $p=0$ and on $\partial P$, and
has an explicit semialgebraic nullhomotopy relative to those parameters
and the spatial boundary.
\item No homeomorphism $h:P\to P$, even a nondefinable TOP
homeomorphism, makes $p\mapsto\Gamma U_{h(p)}$ Lipschitz in the
Euclidean Hausdorff metric.
\item For no such $h$ can one write $U_{h(p)}=s_pr_p^{-1}$ with
two families of homeomorphisms of $X$ that are both uniformly
forward parameter-Lipschitz. This impossibility does not require
the factors to be definable or to preserve any marks.
\item There is a noninjective Lipschitz semialgebraic
$g:P\to P$, equal to the identity on $\partial P$ and at $0$,
such that $U\circ g$ is uniformly forward parameter-Lipschitz.
The straight homotopy from the identity to $g$ gives the allowed
based homotopy from $U$ to this representative.
\end{enumerate}
In particular the example is available on the literal PL six-torus
$B=C^6$. It obstructs a homeomorphic reparameterization or an exact
factorization of the given representative; it is not an obstruction
to the homotopy regularization established later.
\end{theorem}

\begin{proof}
\textbf{The tent and the homeomorphisms.}
For $p\ne0$ define
\begin{equation}
 r=\|p\|_\infty,\qquad \theta=p/r\in C,\qquad
 a=r(1-r),\qquad
 f_p(z)=\max\{0,a-\|z-\theta\|_\infty\}\quad(z\in C),
 \label{equation-11}
\end{equation}
and put $f_0=0$. These are rational semialgebraic formulas, using
division only on $r>0$. The identity
\begin{equation}
           \frac14-r(1-r)=(r-\tfrac12)^2
 \label{equation-12}
\end{equation}
shows $0\le f_p\le a\le1/4$. The bound by $a$ proves continuity
at $p=0$ and at $\partial P$, where $a=0$. Elsewhere continuity
follows directly from the formula. Every spatial slice $f_p$ is
$1$-Lipschitz in the maximum norm and hence also in the Euclidean
norm of the ambient plane.

Write $b=(z_1,\ldots,z_d)\in B$, and set
\begin{equation}
 \eta(y)=\max\{0,1-2|y|\},\qquad \varepsilon=\tfrac14,\qquad
 U_p(b,y)=\bigl(b,y+\varepsilon f_p(z_1)\eta(y)\bigr).
 \label{equation-13}
\end{equation}
For fixed $b,p$, put $c=\varepsilon f_p(z_1)$, so
$0\le c\le1/16$. The last-coordinate map is
\begin{equation}
 T_c(y)=
 \begin{cases}
 y,&-1\le y\le-\tfrac12,\\
 (1+2c)y+c,&-\tfrac12\le y\le0,\\
 (1-2c)y+c,&0\le y\le\tfrac12,\\
 y,&\tfrac12\le y\le1.
 \end{cases}
 \label{equation-14}
\end{equation}
The common join values are $-1/2,c,1/2$, and all slopes belong to
$[7/8,9/8]$. Therefore $T_c$ is an increasing homeomorphism.
Its inverse is the identity outside $[-1/2,1/2]$, and on the
two remaining intervals it is
\begin{equation}
 T_c^{-1}(Y)=
 \begin{cases}
 (Y-c)/(1+2c),&-\tfrac12\le Y\le c,\\
 (Y-c)/(1-2c),&c\le Y\le\tfrac12.
 \end{cases}
 \label{equation-15}
\end{equation}
Positive denominators and the matching join values show that
this inverse is jointly continuous and semialgebraic. Thus
\eqref{equation-13} is a genuine joint family of homeomorphisms.

We give uniform spatial estimates, independent of $d$.
For two input points let $D=\|b-b'\|_2$ and $E=|y-y'|$.
The last-coordinate difference is at most
$D/4+9E/8$. Consequently
\begin{equation}
 \|U_p(b,y)-U_p(b',y')\|_2
 \le \tfrac54D+\tfrac98E
 \le\tfrac{5\sqrt2}{4}\sqrt{D^2+E^2}<2\sqrt{D^2+E^2}.
 \label{equation-16}
\end{equation}
For the inverse, monotonicity with slope at least $7/8$ gives
\[
 |y-y'|\le\tfrac87|Y-Y'|+\tfrac27\|b-b'\|_2.
\]
The inverse Euclidean Lipschitz constant is therefore at most
$9\sqrt2/7<2$. The strict inequalities follow from
$25/8<4$ and $162/49<4$.

\textbf{The local angular lower bound.}
Let $p,q$ be nonzero interior parameters. Write their amplitudes as
$a,a'$ and their directions as $\theta,\theta'$. If
\[
             \delta=\|\theta-\theta'\|_\infty
                            \le\min\{a,a'\},
\]
evaluation of the two tents at $\theta$ and $\theta'$ gives
\begin{equation}
 \|f_p-f_q\|_\infty
       \ge\max\{|a-a'+\delta|,\ |a-a'-\delta|\}
       =|a-a'|+\delta\ge\delta.
 \label{equation-17}
\end{equation}
The condition on $\delta$ ensures that both subtracted tent values
are on their nonzero affine branches. The equality is the elementary
two-sign identity for real numbers and $\delta\ge0$.

For maps of $X$ define
$d_\infty(V,W)=\sup_{x\in X}\|V(x)-W(x)\|_2$.
Since $\eta(0)=1$ and every $z_1\in C$ occurs in $B$, we have
the exact equality
\begin{equation}
           d_\infty(U_p,U_q)
                     =\varepsilon\|f_p-f_q\|_\infty.
 \label{equation-18}
\end{equation}
Write $d_H$ for the Euclidean Hausdorff distance between their
compact graphs in $X\times X$. Matching the same input point gives
$d_H(\Gamma U_p,\Gamma U_q)\le d_\infty(U_p,U_q)$.
In the other direction, for any $x$ choose a closest point
$(x',U_q(x'))$ in the other graph. The spatial Lipschitz bound $2$
gives
\begin{equation}
 \begin{split}
 \|U_p(x)-U_q(x)\|_2
 &\le\|U_p(x)-U_q(x')\|_2+
                         2\|x'-x\|_2\\
 &\le3d_H(\Gamma U_p,\Gamma U_q).
 \end{split}
 \label{equation-19}
\end{equation}
Thus \eqref{equation-17}--\eqref{equation-19} imply
\begin{equation}
 \|\theta-\theta'\|_\infty
       \le12d_H(\Gamma U_p,\Gamma U_q)
 \quad\text{if }\|\theta-\theta'\|_\infty\le\min\{a,a'\}.
 \label{equation-20}
\end{equation}
The graph metric and the forward supremum metric have been kept
distinct; \eqref{equation-19} uses the uniform spatial bound.

\textbf{An arbitrary parameter homeomorphism cannot repair the metric.}
Suppose, towards a contradiction, that
$p\mapsto\Gamma U_{h(p)}$ is $L$-Lipschitz for a homeomorphism
$h:P\to P$. Put $q=h^{-1}(0)$. A homeomorphism of the closed square
preserves its boundary: invariance of domain applied to the
restrictions of $h$ and $h^{-1}$ to interiors proves that both
send interior points to interior points. Thus $q\in\operatorname{int}P$.

Choose $\rho>0$ so that
\begin{equation}
 Q_\rho=q+[-\rho,\rho]^2\subset\operatorname{int}P,\qquad
                    \rho<\frac{1}{100(1+L)}.
 \label{equation-21}
\end{equation}
On the boundary loop of this small square define
$\alpha(v)=\theta(h(v))$. The image $h(\partial Q_\rho)$ avoids both $0$ and
$\partial P$. Therefore the continuous amplitude on that loop
has a positive minimum $\nu$. Uniform continuity of $\alpha$
allows every boundary arc to be partitioned into finitely many
successive pieces whose angular endpoint distance is at most
$\nu$. Retain the four corners in such a partition.
For every successive pair \eqref{equation-20} applies, and the hypothesized
Hausdorff estimate bounds its angular distance by
$12L$ times the Euclidean distance between the domain points.
Summing along an arc, whose length is at most the square perimeter
$8\rho$, gives
\begin{equation}
                   \operatorname{diam}_\infty\alpha(\partial Q_\rho)
                                      \le96L\rho.
 \label{equation-22}
\end{equation}
No regularity of $h$ and no length theorem for its image is used:
only finite subdivision of the original square loop is involved.

On the other hand this angular image is all of $C$.
Indeed $K=h(Q_\rho)$ is compact and contains $0$ in its interior,
and $\partial K=h(\partial Q_\rho)$. For each $u\in C$, the ray
$\{tu:t\ge0\}$ has a last point $t_*u$ in $K$; here $t_*>0$.
The last point belongs to $\partial K$, and its maximum-norm
direction is $u$. This proves $\alpha(\partial Q_\rho)=C$.
Its maximum-norm diameter is $2$, contradicting
\begin{equation}
      2\le96L\rho<
             \frac{96L}{100(1+L)}<\frac{24}{25}<2.
 \label{equation-23}
\end{equation}
For $L=0$ the contradiction is already immediate from \eqref{equation-22}.
This proves the assertion for every TOP parameter homeomorphism.

\textbf{Exact factorization would give the forbidden graph bound.}
Suppose $U_{h(p)}=s_pr_p^{-1}$, with forward parameter bounds
$L_r,L_s$ for two families of homeomorphisms. Its graph is
\[
      \Gamma U_{h(p)}
            =\{(r_p(x),s_p(x)):x\in X\}.
\]
Match the same $x$ at two parameter values, in both directions.
This gives
\begin{equation}
 d_H(\Gamma U_{h(p)},\Gamma U_{h(q)})
                 \le\sqrt{L_r^2+L_s^2}\,\|p-q\|_2,
 \label{equation-24}
\end{equation}
which has just been excluded. The reasoning does not use
definability or any marking conditions on the factors.

\textbf{The family nevertheless has an explicit relative nullhomotopy.}
For $0\le\lambda\le1$ replace $\varepsilon$ in \eqref{equation-13} by
$(1-\lambda)\varepsilon$. The same four slopes still lie in
$[7/8,9/8]$. This is a semialgebraic homotopy through
homeomorphisms from $U$ to the identity, fixed at $p=0$,
on all of $\partial P$, and on the spatial outer collars.
Rescaling $[-1,1]$ affinely to $I$ makes it a based
bottom-fixed concordance cube in the stipulated convention.

\textbf{A noninjective radial clock gives a Lipschitz representative.}
Fix $0<r_0<1$, for example $r_0=1/2$. Define
\begin{equation}
 R(r)=\max\{0,(r-r_0)/(1-r_0)\},\qquad
 g(p)=
 \begin{cases}
 0,&\|p\|_\infty\le r_0,\\
 R(\|p\|_\infty)\,p/\|p\|_\infty,
                         &\|p\|_\infty\ge r_0 .
 \end{cases}
 \label{equation-25}
\end{equation}
The two formulas agree at the interface. This is a continuous
semialgebraic map, collapses the central square to $0$, and fixes
$\partial P$ and $0$.

Here is a uniform estimate that also verifies its Lipschitz property.
For parameters of radii at least $r_0$, put
$\Delta=\|p-q\|_\infty$. Then
\begin{equation}
 |R(r)-R(r')|\le\Delta/(1-r_0),\qquad
       \|\theta(p)-\theta(q)\|_\infty\le2\Delta/r_0.
 \label{equation-26}
\end{equation}
The second inequality follows by subtracting $p/r-q/r'$ and
using $|r-r'|\le\Delta$. Since $0\le R\le1$, these inequalities
make $g$ Lipschitz with constant
$K_0=(1-r_0)^{-1}+2/r_0$ in the maximum norm.
If one radius is below $r_0$, its image is zero and
$R(r)\le|r-r'|/(1-r_0)$ gives the same bound.

The amplitude function $a(u)=u(1-u)$ is $1$-Lipschitz on $I$,
since $|a(u)-a(v)|=|u-v|\,|1-u-v|\le|u-v|$.
Directions in the noncollapsed region are unchanged by $g$.
The tent formula therefore gives
\begin{equation}
             \|f_{g(p)}-f_{g(q)}\|_\infty
                                     \le K_0\|p-q\|_\infty.
 \label{equation-27}
\end{equation}
For a pair straddling the collapsed region use
$0\le f_{g(p)}\le R(r)\le
 \|p-q\|_\infty/(1-r_0)$; inside it both functions vanish.
Combining with \eqref{equation-18}, the forward parameter bound of
$U\circ g$ is $\varepsilon K_0$.
For $r_0=1/2$ this is $3/2$, also a valid Euclidean
parameter bound. The homotopy
$g_t(p)=(1-t)p+tg(p)$ stays in $P$, fixes $\partial P$ and $0$,
and yields the required based homotopy $U\circ g_t$.
For every $t<1$ this parameter map is itself a homeomorphism:
its radial function $(1-t)r+tR(r)$ is strictly increasing from
$0$ to $1$. Only the final map collapses the central square.
The factors $\operatorname{id}$ and $U\circ g$ now satisfy
the forward parameter condition. Thus this example separates
homeomorphic exact reparameterization from the permitted homotopy
regularization.

\textbf{Relation to round circles.}
The main construction already has a literal polyhedral torus base.
If one wants a round-circle model, the maps
\begin{equation}
 R_1:S^1\to C,\quad R_1(z)=z/\|z\|_\infty,\qquad
 R_1^{-1}(u)=u/\|u\|_2
 \label{equation-28}
\end{equation}
are mutually inverse semialgebraic bi-Lipschitz homeomorphisms.
Their constants can be bounded by $\sqrt2+2<4$ and $2$,
respectively, using denominators $\|z\|_\infty\ge1/\sqrt2$
and $\|u\|_2\ge1$. Their finite products, times the interval
identity, remain bi-Lipschitz. For any such spatial conjugacy $J$,
\[
        \Gamma(JVJ^{-1})=(J\times J)(\Gamma V).
\]
Both Lipschitz constants compare the two graph Hausdorff metrics,
so the impossibility of a parameter homeomorphism is preserved.
No equality between a Euclidean-distance tent and the conjugate
of the maximum-norm tent is asserted.
\end{proof}

\Needspace{12\baselineskip}
\section{Finite joins and exact relative parameter control}
We first construct a finite resolution of a proper surjection; the metric conclusion then follows from the coordinate clocks.

\begin{theorem}[A universal finite-join resolution and a relative metric consequence]\label{result-17}
Let \(M\) be any o-minimal expansion of a real closed field. All sets, maps, homotopies, compactness assertions, and products below are definable over \(M\). Here compact means definably compact. Let \(f:K\to Y\) be a continuous surjection between compact definable sets. If \(Y\ne\varnothing\) and \(d=\dim Y\), there are a compact definable set \(R\), a continuous surjection \(q:R\to Y\), an integer \(1\le m\le d+1\), and a continuous map
\begin{equation}
 \ell:R\longrightarrow J_m(f),\qquad \pi\ell=q,
 \label{equation-29}
\end{equation}
where \(J_m(f)\) is the labelled \(m\)-fold fibrewise join of \(f\). The map \(q\) is universally a definable homotopy equivalence: for every continuous map \(T\to Y\) from a compact definable set, \(R\times_YT\to T\) is a definable homotopy equivalence.

If \(Z\subset Y\) is closed and \(b:Z\to J_a(f)\) is a continuous section of its join projection, the resolution can be enlarged to a compact \(R'\) containing \(Z\) as a closed subspace, with \(q'|_Z=\operatorname{id}_Z\), a universal homotopy equivalence \(q':R'\to Y\), and a map
\begin{equation}
 \ell':R'\longrightarrow J_{m+a}(f),\qquad
 \pi\ell'=q',\qquad \ell'|_Z=b
 \label{equation-30}
\end{equation}
(the last equality uses the last \(a\) slots). In particular, for \(Y=P=[0,1]^k\), \(Z=\partial P\), there is a continuous definable \(w:P\to J_{m+a}(f)\) with \(w|_{\partial P}=b\), and \(\pi w\) is definably homotopic to \(\operatorname{id}_P\) relative to \(\partial P\).

For the following metric consequence impose an additional hypothesis, which is not supplied by the topological assertion. Suppose \(C_p\), \(p\in P\), is a bounded Hausdorff-continuous definable family of nonempty compact subsets of a fixed \(M^n\). Suppose there is a compact definable \(K\) and a Lipschitz surjection \(\phi:K\to P\) such that \(s\mapsto C_{\phi(s)}\) is Hausdorff-Lipschitz. If the boundary family \(C|_{\partial P}\) is Hausdorff-Lipschitz, there is a definable Lipschitz map \(g:P\to P\), equal to the identity on \(\partial P\), such that \(p\mapsto C_{g(p)}\) is Hausdorff-Lipschitz. The convex homotopy from \(\operatorname{id}_P\) to \(g\) fixes \(\partial P\). This consequence holds over every \(M\) as above, conditional on the displayed Lipschitz-resolution hypothesis; the topological resolution alone does not supply that metric hypothesis.
\end{theorem}

\begin{proof}
We use only finite constructions. The required all-field topology is provided by Barton--Commelin \cite{BartonCommelin2021}: Proposition 6.17, printed p.~29, gives existence and pullback stability of closed-by-proper pushouts; Propositions 8.5--8.6, pp.~47--48, give the homotopy extension property for closed embeddings of compact definable sets and the resulting cofibration category, whose weak equivalences are definable homotopy equivalences; Proposition 9.28, p.~60, gives a finite closed splitting filtration for every surjective definable map; Proposition 9.31, p.~63, identifies the proper pushout square used below. No ordinary finite subcover or subdivision by an infinitesimal mesh is invoked.

\textbf{Finite mapping-cylinder facts.}
For maps \(A\xrightarrow{a}E\) and \(A\xrightarrow{b}F\) between compact definable sets, write
\begin{equation}
 D(E\leftarrow A\to F)=
 E\ \cup_a\ (A\times[0,1])\ \cup_b\ F
 \label{equation-31}
\end{equation}
for their double mapping cylinder. It is a compact definable space and can be realized as a closed bounded definable set. For completeness, after embedding the three sets in finite coordinate spaces, map the cylinder point \((x,t)\) to
\begin{equation}
 ((1-t)a(x),\,t b(x),\,t(1-t)x,\,t).
 \label{equation-32}
\end{equation}
Include the endpoint copies \((e,0,0,0)\), \(e\in E\), and \((0,z,0,1)\), \(z\in F\). For \(0<t<1\), the third coordinate recovers \(x\); at the two ends exactly the identifications in \eqref{equation-31} occur. The image is compact and definable, and the induced proper quotient is the required realization. Equivalently, \eqref{equation-31} is formed by two closed-by-proper pushouts.

We use two finite homotopy-pushout facts. An objectwise homotopy equivalence between two spans induces a homotopy equivalence of their double mapping cylinders. Also, if \(A\to E\) is a closed embedding, the canonical map from \eqref{equation-31} to the ordinary pushout \(E\cup_AF\) is a homotopy equivalence. These are the gluing consequences of the cofibration category in Proposition 8.6: replace a leg by its mapping cylinder, whose free endpoint is a closed cofibration, and apply the gluing lemma for cofibrations. Thus these facts concern finite definable mapping cylinders, not ordinary topological realizations. They also apply after any compact definable base change. Indeed, each of the two endpoint pushouts in \eqref{equation-31} is pullback-stable by Proposition 6.17, and a closed embedding remains a closed embedding after base change.

\textbf{The labelled join.}
Embed \(K\) in some \(M^N\). Define \(J_m(f)\) as the image of
\begin{equation}
 \{(s_1,\ldots,s_m)\in K^m:f(s_1)=\cdots=f(s_m)\}
          \times\Delta^{m-1}
 \longrightarrow M^{m(N+1)},\qquad
 (s,t)\longmapsto(t_i,t_i s_i)_{i=1}^m.
 \label{equation-33}
\end{equation}
It is compact and definable. A zero-weight slot forgets its \(s_i\). Conversely, any such inactive slots can be completed using one active slot, so this model is exactly the labelled fibrewise join. Its projection \(\pi\) takes the common value \(f(s_i)\) of an active slot. The quotient construction makes \(\pi\) continuous. Adding zero-weight slots gives the usual inclusions among the joins.

\textbf{Induction constructing a universal resolution.}
If \(Y=\varnothing\), then \(K=\varnothing\); use \(R=\varnothing\) and the unique maps. If \(Y\) is nonempty of dimension zero, it is finite; choose one point of each fibre and take \(R=Y\), \(q=\operatorname{id}\), \(m=1\). This case is universal under every base change.

For the induction step, a definable choice section of \(f\) is continuous on a relatively open set \(U\subset Y\) whose closed complement \(C\) has dimension less than \(\dim Y\). This is precisely the dimension-dropping step in Proposition 9.28. Set \(K_C=f^{-1}(C)\) and
\begin{equation}
 E=\overline{s(U)}^{\,K}\cup K_C,\qquad
 A=E\cap K_C=K_C.
 \label{equation-34}
\end{equation}
The addition of \(K_C\) avoids needing a separate density assertion. The restriction \(E\to Y\) is proper and onto, and is a homeomorphism over \(U\): the graph of the continuous section is closed over the relatively open set \(U\). Thus the square
\begin{equation}
 \begin{matrix}
 A&\longrightarrow&E\\
 \downarrow&&\downarrow\\
 C&\longrightarrow&Y
 \end{matrix}
 \label{equation-35}
\end{equation}
is a pullback and a pushout by Proposition 9.31. The upper arrow is a closed embedding.

If \(C=\varnothing\), simply take \(R=E\), \(q=f|_E\), \(m=1\). Otherwise apply the induction hypothesis to \(K_C\to C\), obtaining
\[
 q_C:R_C\to C,\qquad
 \ell_C:R_C\to J_{m_C}(K_C\to C),\qquad
 m_C\le\dim C+1.
\]
Put \(A'=A\times_CR_C\), and form
\begin{equation}
 R=D(E\leftarrow A'\to R_C).
 \label{equation-36}
\end{equation}
All three pieces map to \(Y\), with a common value on each cylinder line; hence they define \(q:R\to Y\). The projection \(A'\to A\) is a homotopy equivalence by the \emph{universal} induction hypothesis. Therefore the map of spans from \eqref{equation-36} to \(E\leftarrow A\to C\) is an objectwise homotopy equivalence. The first mapping-cylinder fact compares their double mapping cylinders, and the second compares the latter with \eqref{equation-35}. It follows that \(q\) is a homotopy equivalence.

Here is the full universality check. Given compact \(T\to Y\), append a subscript \(T\) to each base change. Pullback stability of the two endpoint pushouts identifies \(R\times_YT\), as a definable space, with
\[
 D(E_T\leftarrow A'_T\to (R_C)_T).
\]
The maps \(A'_T\to A_T\) and \((R_C)_T\to C_T\) are base changes of \(q_C\), so are homotopy equivalences. The map \(A_T\to E_T\) is still a closed embedding, and \eqref{equation-35} remains a pushout with lower-right corner \(T\). The same two finite mapping-cylinder facts therefore show \(R\times_YT\to T\) is a homotopy equivalence. No density or section assertion about \(U_T\) is required.

To define \(\ell\), use the point \(e\in E\subset K\) in the first slot on the \(E\)-piece, and use \(\ell_C(r)\) in the remaining \(m_C\) slots on the \(R_C\)-piece. On a cylinder point \((a,r,t)\), with \(f(a)=q_C(r)\), give the first slot weight \(1-t\), and multiply every weight of \(\ell_C(r)\) by \(t\). These points lie in one common fibre. The formula agrees with the endpoint values, so descends to a continuous definable map \(R\to J_{m_C+1}(f)\) over \(Y\). Since \(m_C+1\le d+1\), this completes the induction.

\textbf{A prescribed section on a closed boundary.}
Given \(Z\subset Y\) and \(b\) as in \eqref{equation-30}, set \(R_Z=R\times_YZ\). It is closed in \(R\), and \(R_Z\to Z\) is a homotopy equivalence by universality. Form
\begin{equation}
 R'=D(R\leftarrow R_Z\to Z).
 \label{equation-37}
\end{equation}
The inclusion \(R\to R'\) is a homotopy equivalence: compare \eqref{equation-37} to the double mapping cylinder \(D(R\leftarrow R_Z\to R_Z)\), which retracts onto \(R\), using the homotopy equivalence \(R_Z\to Z\). Thus \(q':R'\to Y\) is a homotopy equivalence by two-out-of-three. The identical argument after compact base change shows it is universal. The endpoint copy \(Z\subset R'\) is a closed embedding and \(q'|_Z=\operatorname{id}\).

Map \(R\) to the first \(m\) join slots by \(\ell\), and \(Z\) to the last \(a\) slots by \(b\). On \(R_Z\times I\), multiply the first weights by \(1-t\) and the last weights by \(t\). Again all active slots have the same base value, so this defines \(\ell'\) over \(Y\).

When \(Y=P\) and \(Z=\partial P\), choose a definable homotopy inverse \(v:P\to R'\), together with a homotopy \(v q'\simeq\operatorname{id}_{R'}\). Its restriction to the endpoint copy of \(\partial P\) is a homotopy from \(v|_{\partial P}\) to that inclusion. Apply the definable homotopy extension property of \(\partial P\subset P\) to change \(v\) to \(u:P\to R'\) with \(u|_{\partial P}\) exactly the endpoint inclusion. Set \(w=\ell'u\). Then \(w|_{\partial P}=b\), while \(\pi w=q'u\) fixes \(\partial P\). Because \(P\) is convex, the straight homotopy from \(\operatorname{id}_P\) to \(q'u\) remains in \(P\) and fixes \(\partial P\). We do not assert that \(w\) is a strict interior lift of \(\operatorname{id}_P\).

\textbf{The conditional metric assertion.}
Use maximum norms. Add to \(K\) a positively separated tagged copy of \(\partial P\); map that copy to \(\partial P\) by the identity and assign its original family \(C_p\). Boundedness, positive separation, and the assumed boundary Lipschitz estimate make both the resulting projection \(\phi^+:K^+\to P\) and its pulled-back compact-set family globally Lipschitz. The boundary section into one join slot is affine in \(p\) in the coordinates \eqref{equation-33}.

For any finite \(m\), both the projection and the pulled-back family on \(J_m(\phi^+)\) are Lipschitz in its actual Euclidean coordinates. Here is the estimate. Let \(M_0=\max_{s\in K^+}\|s\|_\infty\). For two join points \(z=(t_i,v_i)\), \(z'=(t_i',v_i')\), put \(\delta=\|z-z'\|_\infty\). If \(\delta<1/(2m)\), choose \(t_i\ge1/m\). Then \(t_i'>0\), and
\begin{equation}
 \left\|\frac{v_i}{t_i}-\frac{v_i'}{t_i'}\right\|_\infty
 \le m(1+M_0)\delta.
 \label{equation-38}
\end{equation}
Evaluate the projection or the compact-set family on this active slot. For \(\delta\ge1/(2m)\), use the bounded diameter of \(P\), respectively a bound for all Hausdorff distances of the family. This proves a global Lipschitz estimate without selecting slots continuously.

The preceding relative construction supplies a continuous definable \(w:P\to J_m(\phi^+)\) with the literal affine boundary value \(b\), for some finite \(m\le k+2\). We now make \(w\) Lipschitz without any subdivision argument. Apply the total-variation coordinate-clock construction of Theorem~\ref{result-15} to the finite-dimensional vector map \((\operatorname{id}_P,w)\). That construction applies to any finite target dimension: its proof uses only a finite maximum formula for variation, uniform o-minimal monotonicity, and the estimate
\[
 |\operatorname{Var}(h)-\operatorname{Var}(h')|
       \le 2N\|h-h'\|_\infty
\]
for scalar line functions with at most \(N\) monotonicity intervals. It gives a definable homeomorphism \(\chi:P\to P\) preserving every coordinate face such that \(w\chi\) is globally Lipschitz.

In this application \(\chi\) fixes the boundary pointwise. At each stage, a processed coordinate line contained in a boundary face has \(w=b\), which is affine. Its summed coordinate variation is therefore \(ct\) for a constant \(c\) on that line, and the normalized clock \((t+ct)/(1+c)\) is exactly \(t\). A boundary face which fixes the processed coordinate is fixed because the interval clock fixes its two endpoints. Induction preserves the affine boundary values at every stage. This argument is valid over every real closed field.

Finally put \(g=\pi w\chi\). The join estimates make \(g\) Lipschitz and make \(C_g\) Hausdorff-Lipschitz. Its boundary value is the identity. The convex relative homotopy completes the proof. If the parameter cube has dimension zero the assertion is immediate, and empty spatial families, when separately allowed, have their unique trivial regularization.
\end{proof}



\section{Compact Lipschitz resolutions over real closed fields}
A finite collection of selected parameter gradients controls an entire uniformly Lipschitz spatial family.

\begin{theorem}[A uniform spatial family has a finite compact Lipschitz parameter resolution]\label{result-18}
Let \(M\) be an arbitrary o-minimal expansion of a real closed field, let
\(P\subset[0,1]^k\) be a nonempty compact definable set, and let
\[
 F:P\times[0,1]^n\longrightarrow[0,1]^r
\]
be continuous and definable. Use maximum norms. Suppose
\begin{equation}
 \|F(p,x)-F(p,x')\|_\infty\le K_0\|x-x'\|_\infty
 \quad(p\in P,\ x,x'\in[0,1]^n)
 \label{equation-39}
\end{equation}
for some \(K_0\in M_{\ge0}\). There are standard integers
\(J,D,C_\phi,C\), a compact set
\[
 K=\{1,\ldots,J\}\times[0,1]^D,
\]
and a continuous definable surjection \(\phi:K\to P\), such that
\(\phi\) is \(C_\phi\)-Lipschitz and
\begin{equation}
 (z,x)\longmapsto F(\phi(z),x)
 \quad\hbox{is }C(1+K_0)\hbox{-Lipschitz}.
 \label{equation-40}
\end{equation}
The standard integers may depend on the given definable family; no
uniformity across all formulas or all parameter values of an additional
unfixed family is claimed.

Consequently every Hausdorff-continuous definable family of nonempty
compact subsets \(C_p\subset[0,1]^n\), with compact definable parameter
set \(P\subset[0,1]^k\), has a compact Lipschitz resolution with standard
integer projection and Hausdorff constants. Combined with Theorem~\ref{result-17}, this
gives relative Hausdorff-Lipschitz regularization on a parameter cube
over every o-minimal real closed field: if the boundary family is
Hausdorff-Lipschitz, there is a definable Lipschitz \(g:[0,1]^k\to[0,1]^k\),
equal to the identity on the boundary, such that \(C_{g(p)}\) is
Hausdorff-Lipschitz. The convex homotopy from the identity to \(g\) fixes
the boundary.
\end{theorem}

\begin{proof}
The only external differentiable input is the generic
\(\mathcal C^1\) boundary-extension claim in the proof of Lemma 2.4 of
Kocel-Cynk--Pawlucki--Valette \cite{KocelCynkPawluckiValette2019},
printed pp.~102--103, DOI \(10.4153/CMB-2018-030-x\).
Its introduction, printed p.~99, explicitly permits any o-minimal
expansion of any real closed field. The claim says that a definable
\(\mathcal C^1\) scalar function on \(W\times(0,1)\), where \(W\) is
open, with bounded derivative in the last variable, extends
\(\mathcal C^1\) to \(W\times[0,1)\) after deleting a closed nowhere
dense subset of the boundary parameter set \(W\). We use this claim,
not an unproved all-field extension of the 2014 Hausdorff-resolution
theorem. We supply the needed version with a nonsmooth exceptional
set explicitly.

\textbf{Generic bounded parameter gradients.}
Let \(U\subset M^d\) be open and definable, and let
\(H:U\times[0,1]^n\to[0,1]^r\) be continuous and definable with the
uniform spatial bound \eqref{equation-39}. Choose an open dense definable subset
\(\mathcal D\subset U\times(0,1)^n\) on which all components of \(H\)
are \(\mathcal C^1\). By dimension of fibres, after removing a closed
nowhere dense subset of \(U\), each fibre
\(\mathcal D_t\) is dense in \((0,1)^n\).
We claim that, after deleting a further closed nowhere dense subset,
all the sets
\begin{equation}
 \left\{\frac{\partial H_j}{\partial t_i}(t,x):
                   x\in\mathcal D_t\right\}
 \quad(1\le i\le d,\ 1\le j\le r)
 \label{equation-41}
\end{equation}
are bounded.

It suffices to rule out an open parameter set \(W\) where one of
these sets is unbounded. Definable choice, followed by a finite
\(\mathcal C^1\) cell decomposition with the parameter coordinates
first, gives, after shrinking \(W\), a definable \(\mathcal C^1\) map
\(\alpha(t,s)\) on \(0<s<\epsilon(t)\), with
\((t,\alpha(t,s))\in\mathcal D\), whose selected parameter derivative
has absolute value tending to infinity as \(s\to0\).
For example, choose its absolute value greater than \(1/s\) before
the cell refinement. The bounded curve \(\alpha(t,s)\) has a
definable limit \(a(t)\in[0,1]^n\); shrink \(W\) so \(a\) is
\(\mathcal C^1\).

After one more finite refinement, one spatial coordinate
\(\alpha_l\) has nonzero \(s\)-derivative of fixed sign and dominates
all the other absolute \(s\)-derivatives. Such a nonzero coordinate
exists, since an eventually constant curve could not carry a
derivative tending to infinity. Write
\[
 v=\varepsilon_l(\alpha_l(t,s)-a_l(t))>0,
 \qquad \varepsilon_l\in\{1,-1\}.
\]
Inverting this coordinate gives a \(\mathcal C^1\) map
\(\widetilde\alpha(t,v)\) for \(0<v<\beta(t)\), with \(\beta>0\)
definable and \(\mathcal C^1\), and
\[
 \left|\frac{\partial\widetilde\alpha_i}{\partial v}\right|
 \le1\quad(1\le i\le n).
\]
Set \(b(t)=\beta(t)/(2(1+\beta(t)))\) and
\(\gamma(t,u)=\widetilde\alpha(t,b(t)u)\), \(0<u<1\).
Then \(|\partial\gamma_i/\partial u|\le1\), the image still lies in
\(\mathcal D\), and the selected parameter derivative still tends
to infinity at \(u=0\). Apply the cited generic boundary-extension
claim to every component of \(\gamma\). They extend
\(\mathcal C^1\) at generic boundary parameters. The function
\(h(t,u)=H_j(t,\gamma(t,u))\) is \(\mathcal C^1\), and
\[
 \left|\frac{\partial h}{\partial u}\right|\le nK_0.
\]
Apply the same claim to \(h\). At a parameter where all these
extensions are \(\mathcal C^1\), the chain rule gives
\begin{equation}
 \frac{\partial H_j}{\partial t_i}(t,\gamma(t,u))
 =
 \frac{\partial h}{\partial t_i}(t,u)
 -\sum_{\nu=1}^n
   \frac{\partial H_j}{\partial x_\nu}(t,\gamma(t,u))
   \frac{\partial\gamma_\nu}{\partial t_i}(t,u).
 \label{equation-42}
\end{equation}
Every term on the right is bounded as \(u\to0\): the spatial
partials have absolute value at most \(K_0\), and the other
partials extend continuously. This contradicts the choice of
\(\gamma\), proving the claim. All bad-parameter sets are definable;
a definable set with empty interior has nowhere dense closure,
which gives the asserted closed exceptional set. When \(n=0\),
there is no spatial supremum, and the assertion follows directly
from generic \(\mathcal C^1\) regularity.

\textbf{A finite norming collection of parameter gradients.}
On a good open parameter set, define
\[
 \mathcal G_t=
 \{e_1,\ldots,e_d\}
 \ \cup
 \{\nabla_tH_j(t,x):x\in\mathcal D_t,\ 1\le j\le r\}.
\]
For each \(t\), this is a bounded definable subset of \(M^d\).
The number
\[
 A(t)=\sup\{|\det(g_1,\ldots,g_d)|:g_i\in\mathcal G_t\}
\]
is finite and at least \(1\). Definable choice supplies an ordered
tuple \(g_1(t),\ldots,g_d(t)\in\mathcal G_t\) with
\begin{equation}
 |\det(g_1(t),\ldots,g_d(t))|>\tfrac12 A(t).
 \label{equation-43}
\end{equation}
For every \(g\in\mathcal G_t\), Cramer's rule then gives
\begin{equation}
 g=\sum_{i=1}^d c_i g_i(t),\qquad |c_i|\le2.
 \label{equation-44}
\end{equation}
Indeed, each numerator determinant replaces one selected vector
by another member of \(\mathcal G_t\), so is at most \(A(t)\).

Choose witnesses for the selected gradients. After a finite
parameter refinement, each selected vector is either a fixed
basis vector \(e_l\), or
\(\nabla_tH_{j_i}(t,\lambda_i(t))\), where
\((t,\lambda_i(t))\in\mathcal D\). Off a lower-dimensional
parameter set the \(\lambda_i\) are \(\mathcal C^1\), and the
finite labels \(j_i,l\) are fixed. In the second case put
\(h_i(t)=H_{j_i}(t,\lambda_i(t))\). In the basis-vector case put
\(\lambda_i(t)=0\) and \(h_i(t)=t_l\).
For a differentiable parameter curve \(t=t(u)\), these choices
give the exact formulas
\begin{equation}
 \begin{aligned}
 g_i(t)\cdot\partial_{u_\nu}t
 &=\partial_{u_\nu}h_i
   -\nabla_xH_{j_i}(t,\lambda_i(t))
                  \cdot\partial_{u_\nu}\lambda_i
      &&\text{for a function witness},\\
 g_i(t)\cdot\partial_{u_\nu}t
 &=\partial_{u_\nu}t_l
      &&\text{for a basis witness}.
 \end{aligned}
 \label{equation-45}
\end{equation}
Thus it suffices to control the finitely many bounded coordinate
functions \(t,\lambda_i,h_i\); no bound on the size of the original
parameter gradients is required.

\textbf{Standard bounds from the finite coordinate clocks.}
We record the quantitative part of Theorem~\ref{result-15} that is used here.
If \(q_0:[0,1]^e\to[0,1]^a\) is continuous and definable, there is
a definable homeomorphism \(\chi\) of the source cube, preserving
every coordinate face, for which \(q=q_0\chi\) is Lipschitz with
a standard integer constant \(L_q\). Apply the coordinate clocks
to \((\operatorname{id},q_0)\). Every component stays in \([0,1]\).
At a stage with \(N\) monotonicity pieces and \(a+e\) components,
the total variation is at most \((a+e)N\), so the new coordinate
bound is at most \(1+(a+e)N\). Previously obtained bounds are
multiplied by at most \(1+4(a+e)N\), by the variation comparison
in Theorem~\ref{result-15}. The number of stages and all monotonicity counts are
ordinary finite integers. Starting from the first coordinate
bound therefore gives a standard integer bound. This assertion
does not rely on the field being Archimedean and does not claim
that a field-valued Lipschitz constant can always be rounded to
a standard integer.

\textbf{Resolving one parameter stratum.}
We prove by induction on dimension a slightly stronger statement:
for every definable \(S\subset P\), finitely many compact cube
charts have images in \(\overline S\), cover \(S\), have standard
Lipschitz projection constants, and have pulled \(F\) satisfying
\eqref{equation-40}. Zero-dimensional sets are finite; constant charts work,
with spatial bound \(K_0\). The empty set needs no charts.

Decompose a positive-dimensional \(S\) into finitely many
\(\mathcal C^1\) cells. Lower-dimensional cells are covered by
induction. A cell of dimension \(d\) has a definable
\(\mathcal C^1\) parametrization
\[
 p:(0,1)^d\longrightarrow S_0
\]
which is a diffeomorphism onto that cell. Its image is bounded
in \([0,1]^k\), although its derivatives need not be bounded.
Apply the preceding arguments to \(H(t,x)=F(p(t),x)\).
Remove the lower-dimensional exceptional parameter set and
refine so all selected witnesses and their finite labels have
the stated regularity. The images under \(p\) of the removed
sets have dimension less than \(d\), so induction will cover
them.

On each remaining open parameter piece \(W\), form the bounded
graph
\begin{equation}
 \Gamma_\Phi=
 \{(t,p(t),\lambda_1(t),\ldots,\lambda_d(t),
                    h_1(t),\ldots,h_d(t)):t\in W\}.
 \label{equation-46}
\end{equation}
All its coordinates belong to unit cubes. Triangulate the compact
closure \(\overline{\Gamma_\Phi}\) compatibly with
\(\Gamma_\Phi\). Take the closures of all simplices whose relative
interiors are mapped into \(\Gamma_\Phi\). These finitely many
closed simplices cover \(\overline{\Gamma_\Phi}\): a compatible
finite triangulation represents the graph by a union of open
simplices, and its closure is the union of their closed
simplices. Parametrize each such closed \(e\)-simplex by a fixed
PL homeomorphism from \([0,1]^e\), and compose with the definable
triangulation map. Apply the preceding coordinate-clock
construction to this finite vector map. The resulting chart
\[
 q:[0,1]^e\longrightarrow\overline{\Gamma_\Phi}
\]
is Lipschitz with a standard integer constant \(L_q\), has the
same image as before, and maps its cube interior into
\(\Gamma_\Phi\). Let \(\phi_q\) be its \(p\)-coordinate.
Its image lies in \(\overline{S_0}\subset P\), and it is
\(L_q\)-Lipschitz.

\textbf{Controlling the entire pulled spatial family.}
Set \(\widetilde F(u,x)=F(\phi_q(u),x)\). This is continuous on
the compact product cube. On the cube interior the coordinate
relations \eqref{equation-46} hold. The chart \(q\) is \(\mathcal C^1\) on
a dense open subset. Moreover, for every good \(t\), the spatial
exceptional fibre \((0,1)^n\setminus\mathcal D_t\) has dimension
less than \(n\). Thus its pullback under \(u\mapsto t(u)\) has
dimension less than \(e+n\), by dimension of fibres. On a dense
open subset of the product cube the chart and the composed
function are therefore \(\mathcal C^1\), and all formulas
\eqref{equation-44}--\eqref{equation-45} apply.

Each coordinate derivative of \(q\) is bounded by \(L_q\).
For a selected function witness, \eqref{equation-45} gives
\[
 |g_i(t(u))\cdot\partial_{u_\nu}t|
       \le L_q(1+nK_0);
\]
for a basis witness the bound \(L_q\) suffices. Hence, for every
scalar component of the pulled family,
\begin{equation}
 \left|\frac{\partial\widetilde F_j}{\partial u_\nu}\right|
 \le 2d L_q(1+nK_0),\qquad
 \left|\frac{\partial\widetilde F_j}{\partial x_\mu}\right|
 \le K_0
 \label{equation-47}
\end{equation}
on that dense open subset. This remains valid for a chart
simplex of dimension \(e<d\); the spatial exceptional-fibre
bound was imposed for every parameter of \(W\), not merely
for generic parameters of the ambient \(d\)-space.

For clarity, bounded derivatives on such a dense open set imply
the required global Lipschitz bound here. A continuous definable
scalar function on a convex cube, with all coordinate derivatives
bounded on a dense open subset, is Lipschitz with the sum of
those bounds. Given an interior segment direction, project the
lower-dimensional bad set along that direction. Outside a
lower-dimensional set of parallel offsets, a parallel segment
meets the bad set in finitely many points, by dimension of
fibres. Split at those finitely many points, apply the Mean
Value Theorem on the good open intervals, and use continuity
at the splitting points. Arbitrarily small such parallel
translations exist. Pass to the given segment by continuity,
and then to boundary endpoints by interior approximation.
Every partition in this argument is finite and definable;
neither integration nor an ordinary finite mesh is used.
Applying this argument componentwise to \eqref{equation-47} proves a
Lipschitz bound
\[
 C_q(1+K_0)
\]
with a standard integer \(C_q\), for example any standard
integer at least \(2edL_q(n+1)+n+1\).

This produces the required charts over \(W\), and the
lower-dimensional induction covers all removed parameter
sets. All resulting images remain in \(\overline S\).
For \(S=P\), their union is therefore exactly \(P\).
Pad lower-dimensional source cubes by unused coordinates to
one common dimension \(D\), and put the finitely many charts
on the distinct integer tags \(1,\ldots,J\). Between different
tags the source distance is at least \(1\), whereas all
projection and function values lie in unit cubes. The
within-chart constants thus give global standard constants
\(C_\phi,C\) as in \eqref{equation-40}. This proves the first assertion.

\textbf{Compact-set families and the relative consequence.}
For \(C_p\subset[0,1]^n\), define
\[
 F(p,x)=\operatorname{dist}_\infty(x,C_p),
 \qquad x\in[0,1]^n.
\]
The distance is a definable minimum, belongs to \([0,1]\), is
\(1\)-Lipschitz in \(x\), and is jointly continuous by Hausdorff
continuity of the original family. For nonempty compact
\(A,B\subset[0,1]^n\),
\begin{equation}
 d_H(A,B)=
 \sup_{x\in[0,1]^n}
 |\operatorname{dist}_\infty(x,A)
             -\operatorname{dist}_\infty(x,B)|.
 \label{equation-48}
\end{equation}
The upper inequality is the triangle inequality; the reverse
inequality is obtained by evaluating at a point of \(A\) or
\(B\) attaining the directed Hausdorff distance. These
maxima exist by definable compactness. Applying \eqref{equation-40} with
\(K_0=1\) and then \eqref{equation-48} gives a compact resolution with
standard Hausdorff and projection constants.

The initial metric hypothesis of Theorem~\ref{result-17} is now fulfilled over
every o-minimal real closed field. Apply its relative
construction, including the affine tagged boundary section
and the normalized finite coordinate clocks. This supplies
the stated \(g\) and its relative convex homotopy. The empty
parameter set has the unique empty resolution; zero-dimensional
parameter cubes and separately allowed empty spatial families
have their immediate trivial versions.
\end{proof}



\subsection{Increasing clocks and based sphere parameters}
The compact resolution also gives increasing changes of a scalar
parameter and relative regularizations on a sphere.  The first
construction keeps a prescribed finite set of times fixed; the second
keeps the basepoint fixed.


Over the reals, the Hausdorff-resolution theorem of
Kocel-Cynk--Paw\l ucki--Valette provides a Lipschitz parameter
resolution for a bounded family of compact sets
\cite[Main Theorem, author manuscript p.~1]{KocelCynkPawluckiValette2014}.
Valette's generic Lipschitz proposition instead gives joint bounds
on compact subsets of finitely many parameter pieces
\cite[Proposition~5.7, pp.~25--26; scope on p.~22]{Valette2022RegularVectors}.
The scalar consequence below gives one increasing clock fixing
any prescribed finite set. It follows over every real closed field
from the finite cube resolution proved above. The sphere and based
homeomorphism-family consequences below also hold over every such field.

\begin{corollary}[An increasing definable clock]
\label{cor:increasing-definable-clock}
Let $M$ be an o-minimal expansion of a real closed field, let
$Z\subset M^n$ be closed, bounded and definable, and let
$F:[0,1]\times Z\to M$ be continuous and definable.  Suppose
\[
 |F(t,z)-F(t,z')|\le L\|z-z'\|_2
 \qquad(t\in[0,1],\ z,z'\in Z)
\]
for some $L\in M_{\ge0}$.  For every finite set $E\subset[0,1]$
there is an increasing definable Lipschitz homeomorphism
$h:[0,1]\to[0,1]$, fixing $E$, and $K\in M_{\ge0}$ such that
\[
 |F(h(s),z)-F(h(s'),z')|
 \le K|s-s'|+L\|z-z'\|_2.
\]
If $(A_t)_{t\in[0,1]}$ is a Hausdorff-continuous definable family
of nonempty closed bounded sets contained in one bounded box, such
an $h$ also satisfies
\[
 d_H(A_{h(s)},A_{h(s')})\le K|s-s'|.
\]
The constants $L,K$ may be arbitrary elements of $M$.
\end{corollary}

\begin{proof}
We first extract an increasing clock from the finite cube resolution.
Suppose $\phi:K\to[0,1]$ is a continuous definable surjection from
a finite disjoint union of closed cubes, and that a metric-valued
family $G_t$ has a Lipschitz pullback $G_\phi$.  In each cube on
which $\phi$ is nonconstant, choose a point of minimum and a point
of maximum, and join them by an affine segment $\alpha:[0,1]\to K$.
The continuous definable function $\phi\alpha$ has image the full
interval between that cube's minimum and maximum.  Partition its
domain into finitely many closed intervals on each of which it is
constant or strictly monotone.

The nondegenerate image intervals of these monotone pieces cover
$[0,1]$.  Indeed, their union is closed.  Every value omitted from
the union belongs to the finite set consisting of constant-piece
values and images of constant cubes.  A nonempty relatively open
subset of $[0,1]$ cannot be a finite set, so nothing is omitted.
This argument includes zero-dimensional cubes.

Partition $[0,1]$ at all endpoints of these image intervals and all
points of $E\cup\{0,1\}$.  Each successive interval $[a_i,a_{i+1}]$
lies in the image of one strict monotonicity piece: take an image
interval containing its midpoint, and use the fact that every
endpoint of every image interval is a partition point.  Let
$u_i,v_i$ be the preimages of $a_i,a_{i+1}$ in that piece, and put
\[
 h(s)=\phi\alpha_i\left(
 u_i+\frac{v_i-u_i}{a_{i+1}-a_i}(s-a_i)\right)
 \qquad(a_i\le s\le a_{i+1}).
\]
If the piece is decreasing, then $u_i>v_i$, so the displayed
composition still increases.  Its endpoint values are $a_i$ and
$a_{i+1}$.  The formulas therefore join to a continuous strictly
increasing definable bijection of $[0,1]$, hence a homeomorphism,
fixing every partition point.

The argument of $\phi$ on each such interval is an affine path in
one cube with a finite field-valued Lipschitz constant.  The
Lipschitz bound for $G_\phi$ gives a bound for $G_h$ on that interval.
Take the maximum of the finitely many bounds.  For parameters in
different intervals, split the interval between them at the
partition points and apply the triangle inequality.  The lengths
of the pieces add to the original parameter distance, so this
maximum is a global bound.  The same reasoning shows that $h$ is
Lipschitz whenever $\phi$ is Lipschitz.  Only finite monotonicity,
definable extrema and the definable intermediate value theorem
have been used; all are valid over $M$.

If $Z$ is empty, the function assertion is vacuous and $h$ may be
the identity.  Otherwise place $Z$ in a closed bounded coordinate
box $Q$, and extend the function to this box by
\[
 \widetilde F(t,x)=
 \min_{z\in Z}\bigl(F(t,z)+L\|x-z\|_2\bigr).
\]
Definable compactness gives the minimum.  The given spatial
inequality proves equality with $F$ on $Z$, and the triangle
inequality gives the same spatial constant $L$ on $Q$.  Moreover,
\[
 |\widetilde F(t,x)-\widetilde F(t',x')|
 \le \sup_{z\in Z}|F(t,z)-F(t',z)|+L\|x-x'\|_2.
\]
Definable uniform continuity on the closed bounded domain proves
joint continuity.  Normalize the box and the bounded range by
affine changes of coordinates and apply Theorem~\ref{result-18}.
Norm equivalence changes only the bound used for that theorem.
It gives a finite cube resolution for which the pullback is
jointly Lipschitz.

Apply the preceding clock construction to $G_t=F(t,-)$ with the
supremum metric.  The resulting parameter estimate is uniform
in $z$.  Combining it with the original spatial estimate gives
the asserted inequality with the original coefficient $L$.
For the compact-set assertion, apply the Hausdorff consequence of
Theorem~\ref{result-18} and the same construction with $G_t=A_t$.
No point matching between different compact sets is required.
\end{proof}

For an explicit example with a fixed intermediate time, let
$F(t,z)=\sqrt t+Lz$ on $[0,1]^2$ and require $h(1/2)=1/2$.
Writing $a=1/\sqrt2$, the clock
\[
 h(s)=
 \begin{cases}
  2s^2,&0\le s\le1/2,\\
  \bigl(a+2(1-a)(s-1/2)\bigr)^2,&1/2\le s\le1
 \end{cases}
\]
is increasing and fixes $0,1/2,1$.  The function $\sqrt{h(s)}$
is affine on each half, with slopes $\sqrt2$ and $2-\sqrt2$.
Thus the parameter constant may be taken to be $\sqrt2$, while
the spatial constant stays $L$.

\begin{corollary}[Hausdorff regularization on a based sphere]
\label{cor:based-sphere-hausdorff}
Let $R$ be a real closed field with a fixed o-minimal expansion.
In this statement compact means closed, bounded and definable, and
all norms are maximum norms.  Let
$\Sigma^k=\partial[-1,1]^{k+1}$, $k\ge1$, with a chosen basepoint $b$,
and let $(C_p)_{p\in\Sigma^k}$ be a Hausdorff-continuous definable
family of nonempty compact sets with bounded incidence.
There is a definable Lipschitz map $g:\Sigma^k\to\Sigma^k$ which
fixes $b$, is definably homotopic to the identity through maps fixing
$b$, and makes $p\mapsto C_{g(p)}$ Hausdorff-Lipschitz.
The Lipschitz constants may be arbitrary elements of $R$.

On a parameter cube, a Hausdorff-Lipschitz family on the entire
boundary admits the corresponding regularization fixing that
whole boundary.  In particular, let $F$ be a closed bounded
definable set and let $(U_p)$ be a jointly continuous definable
family of homeomorphisms of $F$ on a parameter cube, with
$U_p=\operatorname{id}_F$ on its entire boundary.
There is a definable homotopy relative to that boundary through
homeomorphism families to one with Hausdorff-Lipschitz graphs.
Every spatial condition common to the original family persists.
The full map and its inverse family are jointly continuous in
the field topology throughout the homotopy.  The empty $F$ case
is allowed separately.
\end{corollary}

\begin{proof}
The cube assertion is the relative conclusion of
Theorem~\ref{result-18}.  We derive the sphere assertion using
finite cells and estimates over $R$.

First, for $d\ge2$, any $x,y\in\partial[-1,1]^d$ can be joined
in that boundary by at most three segments, each lying in one
facet, with total maximum-norm length at most
$3\|x-y\|_\infty$.  Choose coordinates $x_i=\epsilon$ and
$y_j=\delta$ with $\epsilon,\delta\in\{-1,1\}$.
If $i\ne j$, change the $j$-coordinate of $x$ to $y_j$ and
then join to $y$.  Each of the two segments has length at most
$\|x-y\|_\infty$ and lies in one of the specified facets.
If $i=j$ with the same sign, use the common-facet segment.
If the signs are opposite, the distance is $2$; choose a
different coordinate, change it to $1$ at $x$ and at $y$,
and use the resulting three segments, each of length at most $2$.

Consequently a map into any $R$-valued metric space that is
$L$-Lipschitz on each facet and agrees on their intersections
is $3L$-Lipschitz on the whole boundary.  Affine normalization
gives the same conclusion for a rectangular cell, with a
constant depending on its positive side lengths.
For a one-dimensional cell boundary, the two values are
Lipschitz directly: divide their distance by the positive
distance between the endpoints.

We also need a Lipschitz extension of a boundary map.
For an $L$-Lipschitz map
$a:\partial[-1,1]^d\to[-1,1]^d$, put
\[
 E(0)=0,\qquad
 E(x)=\|x\|_\infty\,a(x/\|x\|_\infty)\quad(x\ne0).
\]
It extends $a$ and takes the cube into itself.
For $r=\|x\|_\infty\ge s=\|y\|_\infty>0$ and
$D=\|x-y\|_\infty$,
\[
 s\|x/r-y/s\|_\infty
 \le (s/r)D+(s/r)(r-s)\le2D.
\]
Since $r-s\le D$ and $\|a\|_\infty\le1$, this proves
$\|E(x)-E(y)\|_\infty\le(1+2L)D$.
The case $s=0$ is immediate.  Affine conjugation gives
the corresponding extension on every nondegenerate
rectangular cell.

In ambient coordinate $i$, subdivide $[-1,1]$ at the distinct
elements of $\{-1,b_i,1\}$.  The cells lying in $\Sigma^k$
form a finite rectangular complex, and $b$ is a vertex.
Construct $g$ in increasing cell dimension, fixing every
vertex.  Require that $g$ take every cell into itself, agree
on common faces, and that both $g$ and the pulled family be
Lipschitz on each constructed cell.

Suppose the construction is complete on the proper faces
of a cell $Q$.  The pasted map on $\partial Q$ and its
pulled family are Lipschitz by the boundary estimate above,
with the separate two-point argument when $\dim Q=1$.
Cone the boundary map to a definable Lipschitz
$E_Q:Q\to Q$.  The family
\[
 x\longmapsto C_{E_Q(x)}
\]
has nonempty compact fibres, bounded incidence, and
Hausdorff continuity, and its entire boundary family is
Hausdorff-Lipschitz.  Apply the relative cube theorem in
affine coordinates on $Q$.  It supplies a definable
Lipschitz self-map $h_Q$ fixing $\partial Q$ such that
$x\mapsto C_{E_Q(h_Q(x))}$ is Hausdorff-Lipschitz.
Set $g|_Q=E_Qh_Q$.  This agrees exactly with all earlier
faces and preserves $Q$.  Finite induction completes $g$.

These cellwise bounds are global.  On an original facet
of $[-1,1]^{k+1}$, split a straight segment at its finitely
many grid crossings.  Each piece lies in one closed
cell, and its lengths sum to the original distance.
The maximum of the finitely many cell constants therefore
bounds the restriction to the entire facet.  The preceding
boundary-path estimate in dimension $k+1\ge2$ gives the
global bounds on $\Sigma^k$, for both $g$ and the
Hausdorff metric of its pulled family.
All these constants lie in $R$; infinitesimal grid widths
only change the finite affine-normalization constants.

The formulas agree on closed cell intersections, so
finite closed pasting gives continuity.  For $p$ in
any cell $Q$, both $p$ and $g(p)$ belong to $Q$.  Hence
\[
 H(s,p)=(1-s)p+s\,g(p)
\]
stays in the same convex cell and thus in $\Sigma^k$.
It is a jointly continuous definable homotopy from the
identity to $g$, fixing all vertices and in particular $b$.

For the homeomorphism-family assertion, assume first
$F\ne\varnothing$ and write $P$ for its parameter cube.
Definable uniform continuity on the closed bounded
domain $P\times F$ gives Hausdorff continuity of
$\Gamma_p=\{(x,U_p(x)):x\in F\}$, since
\[
 d_H(\Gamma_p,\Gamma_q)
 \le \max_{x\in F}\|U_p(x)-U_q(x)\|_\infty.
\]
These graphs are nonempty closed bounded definable sets
with bounded incidence.  On $\partial P$ they are the
constant diagonal.  Apply the relative cube theorem to
obtain a boundary-fixed Lipschitz $h:P\to P$ for which
$\Gamma_{h(p)}$ is Hausdorff-Lipschitz, and put
\[
 U^s_p(x)=U_{(1-s)p+s h(p)}(x).
\]
Every member is one of the original homeomorphisms,
so every common spatial condition persists, and the
homotopy is relative to the whole parameter boundary.

For completeness, inverse continuity over $R$ follows
from definable extrema, without ordinary compactness.
If $\Psi:X\to Y$ is a continuous definable bijection and
$X$ is closed and bounded, then, for each $\epsilon>0$,
the nonempty set of pairs $(x,x')\in X^2$ with
$\|x-x'\|_\infty\ge\epsilon$ is closed bounded definable.
On it the continuous function
$\|\Psi(x)-\Psi(x')\|_\infty$ has a positive minimum:
the minimum is attained and injectivity excludes zero.
Thus $\Psi^{-1}$ is uniformly continuous; when the
set of pairs is empty the assertion is immediate.
The general-field extrema statement and uniform-continuity
fact are respectively Theorem~3.4 and Exercise~3.5 of
\cite[pp.~26--27; field convention pp.~9--10]{Coste1999}.

Apply this inverse argument to
\[
 (s,p,x)\longmapsto(s,p,U^s_p(x)),
\]
a continuous definable bijection of the closed bounded
definable set $[0,1]\times P\times F$ onto itself.
It gives the jointly continuous definable inverse
family, as required.  A zero-dimensional parameter
cube already gives a constant graph family.
For $F=\varnothing$ there is a unique family and
constant homotopy; its empty graphs have self-distance
zero when this convention is required.
\end{proof}

The parameter map in the sphere and cube assertions need not be
injective.  The conclusions control the graphs as compact sets;
they do not choose homeomorphisms matching their points.


\part{Tannakian reconstruction, freeness, and polynomial tensors}
We first distinguish the tensor and unit conditions in categorical
reconstruction, then classify the polynomial tensor maps. The
Noetherian classification permits zero divisors in the base and
torsion in its coefficient extension; subsequent sections describe
positive cases and obstructions beyond Noetherian rings. For an
inclusion $R\subset L$ we write
\[
 \Omega_n(R,L)=\{f\in L[X_1,\ldots,X_n]:f(R^n)\subset R\},
 \qquad \operatorname{Int}(R)=\Omega_1(R,\operatorname{Frac}(R)).
\]
The map $\theta_n$ is multiplication in disjoint variables from
$\Omega_1(R,L)^{\otimes_R n}$ to $\Omega_n(R,L)$. The term
\emph{unique polynomial operations} means that $\theta_2$ is injective
and its image contains $f(X+Y)$ and $f(XY)$ for every $f\in\Omega_1(R,L)$.
A ring epimorphism is an epimorphism in the category of commutative
unital rings; equivalently, its multiplication map on the tensor square
is an isomorphism. A plethory is a biring together with associative
composition and identity $X$, compatible with its represented ring
operations. We establish those compatibilities by all the maps
$\theta_n$, rather than by the binary existence condition alone.

The positive local calculation proves
Theorem~\ref{thm:int-analytic-free} for every one-dimensional
analytically irreducible Noetherian local domain. Elliott's original
nonflatness conjectures are in
\cite[Conjectures~2.10, 2.15 and 2.17, author version pp.~11--12]{Elliott2011IntegerTClosure};
\cite[p.~2]{ElliottPresentations2011} recalls them, and
\cite[p.~2]{Elliott2016RingsModules} later discusses nonfree examples.
The two-branch example of Ma \cite[\S\S3--4, pp.~2--5]{Ma2026}
provides such a nonfree module, since it is not even flat.
The branch criterion below identifies why that example does not
belong to the positive class.



\section{The unit constraint in the pro-algebra reconstruction}

Throughout this section $k$ is a field and $L/k$ is finite. Tensor
categories and tensor functors are symmetric and unital. Completed
vector spaces and tensor products are the continuous duals of
discrete vector spaces and their coalgebra constructions; ideals
used in a completed quotient are closed.

\begin{lemma}\label{lem:idealizer-takeuchi}
Let $s:L\to A$ be a unital inclusion into a $k$-algebra. In
$D=A\otimes_kA$ let $J$ be the right ideal generated by
$s(\lambda)\otimes1-1\otimes s(\lambda)$, $\lambda\in L$.
Then $D/J$ is the tensor product of the two left $L$-modules $A$.
The quotient of the left idealizer of $J$ is the algebra
\[
 A\times_LA=\left\{\sum a_i\otimes_Lb_i:
 \sum a_is(\lambda)\otimes_Lb_i
 =\sum a_i\otimes_Lb_is(\lambda)\text{ for all }\lambda\right\}.
\]
The statement also holds for the completed products.
\end{lemma}
\begin{proof}
The right ideal imposes exactly the relations
$s(\lambda)a\otimes b=a\otimes s(\lambda)b$. An element $d$
idealizes it precisely when its product on the left with each
displayed generator belongs to $J$, which is the stated invariance
condition. The idealizer is a subalgebra and $J$ is two-sided
inside it. Thus factorwise multiplication descends there, without
assigning multiplication to the entire balanced tensor product.
The same argument with closed relations proves the completed
version, since multiplication is continuous.
This is the equal-source-and-target specialization of Takeuchi's
product~\cite[Section~1, Proposition~3.1 and Definition~4.5]{Takeuchi1977}.
\end{proof}

\begin{example}\label{ex:missing-unit-character}
The algebra and coproduct conditions alone do not force the
character constraint on the counit, even in finite dimension over
$\mathbb F_2$. Let $E=\{a,b,z\}$ be the commutative semigroup
with $a^2=a$, $b^2=b$, $ab=z$, and $z$ absorbing. Its semigroup
algebra has unit $e_a+e_b-e_z$ and the unital augmentation
$e_s\mapsto1$. On its dual $A=k^E$, use pointwise multiplication
and $\Delta f(s,t)=f(st)$. This coproduct is unital,
multiplicative, coassociative and cocommutative. Its coalgebra
counit is $e(f)=f(a)+f(b)-f(z)$, which is not multiplicative:
$e(p_a)e(p_b)=1$ but $e(p_ap_b)=0$.
The proposed tensor category has no allowed unit. Its three
one-dimensional simples satisfy $V_s\otimes V_t=V_{st}$, while
no element of $E$ is a semigroup identity.
\end{example}

\begin{theorem}\label{thm:corrected-pro-bialgebroid-reconstruction}
Let $A$ be a pseudocompact $k$-algebra with continuous multiplication,
source map $s:L\to A$, left $L$-linear continuous counit $e:A\to L$,
and continuous, unital, multiplicative, cocommutative and
coassociative coproduct into the completed Takeuchi product.
Assume the scalar and ordinary counit compatibilities, and impose
\[
                   e(1)=1,\qquad e(ab)=e(a\,s(e(b))).
\]
These are the data of a continuous cocommutative left
$L$-bialgebroid with equal source and target. Its finite-dimensional
continuous modules form a symmetric abelian tensor category with
exact faithful fibre functor to finite-dimensional $L$-spaces.
Conversely an essentially small abelian $k$-linear symmetric
tensor category with such a fibre functor reconstructs these data,
and the two constructions are inverse up to equivalence respecting
the fibre functor and the induced continuous pro-algebra isomorphism.
\end{theorem}
\begin{proof}
Define the action on $M\otimes_LN$ by the coproduct. The Takeuchi
condition makes it balanced, and all sums act through finite
quotients. Multiplicativity gives the action law; coassociativity
and cocommutativity give the usual associativity and symmetry.
The unit object is $L$ with action
$a\cdot\lambda=e(a s(\lambda))$. The character equation gives
the module law, while $e(1)=1$ and left $L$-linearity give the
unit law and the prescribed scalar action. Apply the counit
identity to $\Delta(a s(\lambda))$ to check that both ordinary
unit constraints are $A$-linear. Coherence is the ordinary linear
coherence. Kernels and cokernels remain finite and continuous,
and tensoring over the field $L$ is exact.

For the converse let $F$ be the fibre functor, set
$W=\operatorname{Res}_{L/k}F$, and take
$A=\operatorname{Nat}_k(W,W)$ with its finite topology. Hom spaces
are finite-dimensional over $k$, since $L/k$ is finite and $F$ is
faithful. Exact faithfulness bounds the length of a strict chain
of subobjects by $\dim_LF(M)$. Thus the category is locally finite.
The coalgebra reconstruction theorem identifies it with finite
continuous $A$-modules; its hypotheses are precisely those just
checked~\cite[6.1--6.2]{Deligne1990}\cite[Theorem~2.1.2]{Hai2008}.
Scalar multiplication gives $s:L\to A$.

A natural endomorphism of $(M,N)\mapsto M\otimes_LN$ is determined
on pairs of finite regular quotients $A/I,A/J$ by its value at
$1\otimes1$. The maps $a\mapsto am$ then determine its value on
every pair of modules. Naturality for right multiplication by
$s(\lambda)$ is exactly the Takeuchi condition. Compatible values
on all such quotients are consequently the completed Takeuchi
product. The action on tensor-product objects reconstructs
$\Delta$; the action on the unit reconstructs $e$. The action law,
associativity and symmetry give all the asserted equations,
including the character identity. Continuous dualization gives
the commutative unital $L$-algebra dual to the coproduct; the
transpose of $1_A$ is its augmentation. The same regular-quotient
argument recovers $A$ as the natural endomorphism algebra, proving
the inverse comparison. The coend formulation is
Hai's~\cite[Corollary~2.2.5]{Hai2008}; no finite-dimensional dual
is substituted for the continuous dual without a finiteness proof.
\end{proof}

\begin{proposition}\label{prop:restricted-diagonal-fibre-functor}
If the preceding category is rigid and $k=\operatorname{End}(1)$,
its scalar extension to $L$ has the exact faithful fibre functor
induced by $F$. In the pro-algebra description this is the diagonal
functor on the corresponding finite continuous
$(A\widehat\otimes_kL)$-modules. The conclusion includes inseparable
$L/k$. The diagonal functor on unrestricted
$(L\otimes_kL)$-modules need not be exact and need not be faithful.
\end{proposition}
\begin{proof}
Deligne identifies the category with representations of the affine
groupoid over $\operatorname{Spec}L$ whose arrow scheme is faithfully
flat over $\operatorname{Spec}(L\otimes_kL)$
\cite[Theorem~1.12]{Deligne1990}. After scalar extension, restrict
an equivariant vector bundle to the diagonal $L$-point. Local
freeness makes this exact; transitivity makes it faithful, since
after the faithfully flat base change of arrows from that point,
every equivariant map is transported from its value there.
Deligne's scalar-extension compatibility identifies this with
$M\mapsto M\otimes_{L\otimes_kL}L$
\cite[5.18--5.21]{Deligne1990}.

For a purely inseparable extension of degree $p$,
$L\otimes_kL=L[\delta]/(\delta^p)$. Tensoring the injection
$(\delta^{p-1})\hookrightarrow L[\delta]/(\delta^p)$ with the
diagonal quotient $L$ gives a nonzero source and zero map.
For a nontrivial separable extension the diagonal component is an
idempotent factor; the other factor is nonzero and is killed.
These are unrestricted-module counterexamples, not objects with
the required equivariance.
\end{proof}

\begin{proposition}\label{prop:continuous-dual-and-rigidity}
Choose $0\ne\ell:L\to k$ linear, and put
$U=\operatorname{Hom}^{\rm cont}_L(A,L)$, with
$(au)(b)=u(ba)$. Then
\[
 \operatorname{Hom}_A(M,U)\simeq\operatorname{Hom}_L(M,L),\qquad
 \operatorname{End}_A(U)\simeq A^{\rm op}.
\]
A rigid tensor structure induces a continuous anti-isomorphism
fixing $s(L)$ and hence a second commuting action on $U$.
Conversely such an anti-isomorphism gives rigidity exactly when,
for every finite continuous module, the ordinary evaluation and
coevaluation maps are $A$-linear. Their triangle identities then
hold automatically. Testing only the infinite pro-module $A$
does not replace the finite-module tests.
\end{proposition}
\begin{proof}
Evaluation at $1$ is the first map; its inverse sends $f$ to
$\phi_f(m)(b)=f(bm)$. The pairing $\ell(xy)$ on $L$ is
nondegenerate: for $x\ne0$ choose $z$ with $\ell(z)\ne0$ and
take $y=x^{-1}z$. Thus $u\mapsto\ell\circ u$ identifies $U$
with the continuous $k$-dual of $A$, compatibly with the action.
Continuous duality identifies its endomorphism ring with the
opposite of the continuous right-regular endomorphism ring of $A$,
which is $A^{\rm op}$ by evaluation at $1$.

For a categorical dual $DM$, identify $F(DM)$ with the $L$-dual
of $F(M)$ and use the nondegenerate $k$-pairing
$\ell(\phi(m))$. Take the adjoint of the action on $DM$.
Naturality gives an element of $A$, transposition reverses
products, and the bidual isomorphism gives the inverse. Scalars
are self-adjoint, so the map fixes $s(L)$; the construction is
continuous in the finite topology. It realizes the dual module
through $\operatorname{Hom}_A(M,U)$.

Conversely use the anti-isomorphism and the commuting action to
give that Hom space its dual action. On underlying spaces this
is ordinary exact $L$-duality. With dual bases $m_i,\phi_i$,
the two maps are evaluation and $1\mapsto\sum_i m_i\otimes\phi_i$.
The triangle identities are
$\sum_i m_i\phi_i(m)=m$ and
$\sum_i\phi(m_i)\phi_i=\phi$. They give categorical duality
precisely when both maps are $A$-linear. Explicitly these conditions
are
\[
 \sum\operatorname{ev}(a_{(1)}m\otimes a_{(2)}\phi)
       =e(a s(\operatorname{ev}(m,\phi))),
 \quad
 \sum_{i,(a)}a_{(1)}m_i\otimes a_{(2)}\phi_i
       =e(a)\sum_i m_i\otimes\phi_i.
\]
For separable $L/k$ the trace is a permissible canonical choice of
$\ell$. In a purely inseparable extension the trace is zero; the
argument uses an arbitrary nonzero functional instead.
\end{proof}

\paragraph{Comparison with established structures.}\label{par:qualified-tannaka-comparison}
Takeuchi's Definition~4.5 requires both the coproduct and the map
to $\operatorname{End}_k(L)$ to be algebra maps. The latter map is
$a\mapsto(\lambda\mapsto e(a s(\lambda)))$, so its multiplicativity
is exactly the missing character equation. Schauenburg's
Theorem~5.1 and Remark~5.3 identify bialgebroid structures with
monoidal module structures and permit transport from a non-strict
underlying functor~\cite{Schauenburg1998}; the regular module in
that proof becomes the family of finite regular quotients here.
Hai's coend theorems provide the abelian reconstruction and
rigid refinements~\cite[Theorems~2.1.2, 2.2.4, 2.2.6 and
Corollary~2.2.5]{Hai2008}. The weak-bialgebra reformulation requires
a Frobenius-separable base~\cite[Section~3 and Theorems~5.1, 5.5]{Schauenburg2003};
a finite separable field extension has that property, while a
purely inseparable extension does not. This restriction is distinct
from the unrestricted reconstruction and restricted diagonal
statements. These comparisons make no priority or independence claim.

\section{Analysers and the extra plethory data}

\begin{proposition}\label{prop:divided-power-analyser}
The positive divided-power lattice
$\Gamma_+=\bigoplus_{n\ge1}\mathbb Z\gamma_n\subset\mathbb Q[T]$,
$\gamma_n=T^n/n!$, is closed under multiplication and composition,
with composition identity $T$. It acts on every divided-power
ideal, including ideals in rings with torsion. Its constant ring
is zero. Its additive decomposition does not take values in
$\Gamma_+\otimes\Gamma_+$; it takes values in the augmentation
ideal of the tensor product of its unitizations.
\end{proposition}
\begin{proof}
Multiplication has coefficients $\binom{m+n}m$. If
$G=\sum a_i\gamma_i$, the coefficient of $\gamma_N$ in
$\gamma_m(G)$ indexed by $\sum r_i=m$, $\sum ir_i=N$ is
\[
                 \frac{N!}{\prod_i r_i!(i!)^{r_i}}
                              \prod_i a_i^{r_i}.
\]
The integer factor counts set partitions with the prescribed block
sizes. These identities follow from the divided-power axioms
without division in the coefficient ring, so they give the
asserted action even with torsion. The ordinary polynomial
identities in $\mathbb Q[T]$ give composition associativity.
The additive formula is
\[
 \gamma_n(X+Y)=\gamma_n(X)+\gamma_n(Y)
                    +\sum_{i=1}^{n-1}\gamma_i(X)\gamma_{n-i}(Y),
\]
while $\gamma_n(XY)=n!\gamma_n(X)\gamma_n(Y)$.
Every factor in $\Gamma_+$ vanishes at zero, so a sum of products
from $\Gamma_+\otimes\Gamma_+$ would vanish when $Y=0$; the
left side does not. This proves the additive qualification and
also the failure of faithful evaluation on the zero constant ring.
\end{proof}

\begin{theorem}\label{thm:analyser-to-plethory}
For a unital analyser $P$ with specified constant subring $k$, a
compatible $k$-plethory consists exactly of a $k$--$k$ biring
structure realizing its additive and multiplicative decompositions,
with scalar evaluation $\beta(c)(F)=F\circ c$, such that composition
$P\odot_kP\to P$ and the unit $k[T]\to P$ are biring morphisms.
Existence of unspecified finite decompositions does not impose the
biring coherence identities. If the full constant ring is a field
$k$, evaluation on it is faithful, and the unital analyser axioms
and both decomposition axioms hold, there is a unique compatible
$k$-plethory.
\end{theorem}
\begin{proof}
The first assertion unpacks the monoid definition of a plethory
and the relations of the composition product
\cite[1.1, 1.3, 1.8]{BorgerWieland2005}. Linearity and
multiplicativity in the outer argument, the two decomposition
relations in the inner argument, and the scalar relation are
exactly what allows composition to descend. The remaining
condition is that composition and the unit respect the biring maps.

For the sufficient criterion regard $P$ as a subalgebra of $k^k$.
Evaluation embeds $P^{\otimes_k n}$ into $k^{k^n}$ for every $n$.
For two factors, a finite independent family of functions in the
first factor has evaluation vectors spanning its dual; sufficiently
many evaluations recover each coefficient function. Induction
gives all $n$. The two decompositions therefore define unique
coproducts. Every biring identity can be checked on constants,
where it is an ordinary field identity, and injectivity makes
the identity hold in the tensor product. The counits are
evaluation at $0,1$, the coinverse is composition with $-T$,
and the scalar map is evaluation at constants. Composition
respects each structure by the same evaluation test; its
associativity and unit are the analyser identities.
\end{proof}

\begin{proposition}\label{prop:constant-evaluation-criterion}
For a polynomial analyser $\Omega(R,S)$ the constants are exactly
$R$, and faithful evaluation on them is equivalent to
\[
 \bigcap_{r\in R}\ker\bigl(S^{d+1}\to S,
                   (a_i)\mapsto\sum_i a_ir^i\bigr)=0
                            \quad(d\ge0).
\]
If $S$ is a domain this holds exactly when $R$ is infinite.
More generally it holds whenever, for each $d$, there is a
$(d+1)$-point Vandermonde determinant regular in $S$.
\end{proposition}
\begin{proof}
A constant under composition equals its composition with zero,
so is an ordinary constant in $R$. The kernel of evaluation is
the null-polynomial ideal, giving the coefficient criterion.
A nonzero polynomial over a domain has at most its degree many
roots, while for finite $R$ the monic product
$\prod_{r\in R}(T-r)$ is a nonzero null polynomial. The adjugate
of a Vandermonde matrix annihilates each null coefficient by its
determinant, proving the regular-determinant implication.
\end{proof}

\begin{example}\label{ex:finite-analyser-without-biring}
The ring $P=\mathbb F_2[T]/((T^2-T)^3)$ has $64$ elements and
is a unital analyser satisfying both finite decomposition axioms,
but admits no compatible plethory structure.
\end{example}
\begin{proof}
For every polynomial $G$, $G^2-G$ vanishes at $0,1$ and is
divisible by $T^2-T$. Consequently substitution preserves the
outer defining relation. Changing the inner representative also
preserves it, by the polynomial difference identity. Ordinary
substitution therefore descends, with its unit, associativity and
outer ring laws. Binomial and diagonal monomial expansions give
the two finite decomposition formulas. The constants are
$\mathbb F_2$, by composing with zero.
Put $q=T^2-T$. A compatible unit forces $q$ to be primitive for
coaddition, so
\[
                  \Delta_+(q^3)=q^2\otimes q+q\otimes q^2.
\]
This is nonzero because $q,q^2$ are linearly independent, whereas
$q^3=0$ in $P$. Hence a compatible coaddition cannot be an algebra
map. This also explains why the faithful-evaluation hypothesis
in the preceding theorem is necessary.
\end{proof}


\section{Integer-valued polynomial operations}
This section gives a criterion for the two polynomial operations in
Grothendieck's cote 153, sheet 4.  Let $R\subset S$ be a unital inclusion
of commutative rings and set
\[
 A=\Omega(R,S)=\{f\in S[X]: f(r)\in R\text{ for all }r\in R\}.
\]
Define $\theta_2:A\otimes_R A\to S[X,Y]$ by
$f\otimes g\mapsto f(X)g(Y)$.  The question is whether, for every
$f\in A$, each of $f(X+Y)$ and $f(XY)$ has exactly one preimage under
$\theta_2$.  Uniqueness includes $f=0$, and consequently requires
$\theta_2$ to be injective.  This binary question is distinct from all
the associativity and composition identities of a plethory.

Elliott's Theorem 12 in \cite{Elliott2010}, pp.~56--57, gives sufficient
tensor hypotheses and a necessary weak polynomial compositeness condition
for compatible biring structures over domains.  Jiang, Li, Sun, Wang,
Xiao and Yu give a negative domain example in \cite[\S3.4]{Rethlas2026}.
Hua gives the free-module coordinate argument and the tensor cooperation
theorem for localization coefficients in
\cite[\S34, pp.~155--157]{HuaMontpellier2026}.  His broader assertion
using only a ring epimorphism and freeness is explicitly listed there
as unverified \cite[Theorem~34.6, p.~160]{HuaMontpellier2026}.
The construction below concerns arbitrary coefficient inclusions and
measures the binary tensor obstruction using a length-three double
complex.  We make no assertion of bibliographic priority.
\begin{theorem}[A universal coefficient-ring obstruction without Noetherian hypotheses]\label{result-19}
For any unital inclusion $R\subset S$ and $d\ge1$,
let $D_d(R)$ be the ideal generated by all
Vandermonde determinants on $d+1$ elements of $R$,
and put
\begin{equation}
 C_d=(R:_S D_d(R)),\qquad C=\bigcup_{d\ge1}C_d.           \label{equation-49}
\end{equation}
Then $C$ is a subring of $S$. It is exactly the
set of coefficients that occur in integer-valued
polynomials, and also exactly the set of their
possible leading coefficients, with zero included. Every multivariable
integer-valued polynomial has its coefficients in
$C$. If the second polynomial tensor map is
injective, the inclusion $R\to C$ is a ring
epimorphism. If $C=R$, all tensor maps are the
ordinary polynomial isomorphisms.
\end{theorem}
\begin{proof}
The ideals $D_d$ decrease with $d$, since a
larger Vandermonde product is a multiple of
one on a subset of its nodes. Thus the $C_d$
increase, and $C_1=R$ because the nodes $0,1$
have determinant one.
Cramer's rule says that every coefficient of
an integer-valued polynomial of degree at most
$d$ lies in $C_d$.

There is an explicit monic polynomial $P_d$
whose every value belongs to $D_d$. For $d=1$
take $P_1=1$. For $d\ge2$, take the Vandermonde
product on the polynomial nodes
\[
                     0,\ 1,\ X,\ X^2,\ldots,X^{d-1}.
\]
Every factor with positive degree is monic,
and the only degree-zero factor is $1$; hence
$P_d$ is monic over every base ring.
For any $c\in C_d$, the polynomial $cP_d$
takes $R$ into $R$ and has leading coefficient
$c$. This proves both coefficient characterizations.
Sums of leading coefficients are obtained after
multiplying the corresponding polynomials by
powers of $X$ to equalize their degrees; products
come from multiplication of polynomials.
Therefore their set $C$ is a subring.

For a multivariable polynomial, substitute
$X_i=T^{N_i}$ with sufficiently separated
positive exponents so that its finitely many
monomials remain distinct. The resulting
one-variable polynomial is integer-valued on
$R$, and has exactly the same coefficients.
Thus all coefficients lie in $C$.

For $c\in C_d$, both $P_d$ and $cP_d$ belong
to $\Omega(R,S)$. The tensor
$cP_d\otimes P_d-P_d\otimes cP_d$ has zero
polynomial image. If the tensor map is
injective it vanishes. Extract the top
coefficient in each factor, now as maps
$\Omega(R,S)\to C$, to get
$c\otimes1=1\otimes c$ in $C\otimes_R C$.
This for every $c\in C$ makes multiplication
$C\otimes_R C\to C$ an isomorphism.
Finally $C=R$ gives $\Omega_n=R[X_1,\ldots,X_n]$
by the coefficient statement, proving the
last assertion.
\end{proof}\begin{proof}[Proof of Theorem~\ref{result-20}]
We first check the matrix test and then construct the invariants and
prove their meaning by explicit chain calculations.
If \eqref{equation-51} holds, balancing its entries shows
$c\otimes1=1\otimes c$ in $C\otimes_R C$: use $YM$ to move the first
sum to the second factor, and then $MZ$ to move it back.
Conversely suppose this tensor equality holds. A finite presentation
of that equality gives finitely many $s_i\in C$, with labeled
$s_0=1,s_1=c$, and a finite matrix $(m_{ij})$ over $R$ such that
\begin{equation}
 \sum_i m_{i0}s_i=c,\qquad \sum_i m_{i1}s_i=-1,
 \qquad \sum_i m_{ij}s_i=0\ (j\ne0,1),\qquad
                       \sum_jm_{ij}s_j=0\ (\text{all }i).
                                                        \label{equation-57}
\end{equation}
Indeed present the second copy of $C$ by the free module on its
elements. The vector having entries $c$ at $s_0$, $-1$ at $s_1$,
and zero elsewhere lies in the image of its relation module tensored
with the first copy of $C$. Only finitely many relations and elements
are needed, giving \eqref{equation-57}, after enlarging the common index set.
Remove the row and column indexed by zero, and call the remaining
matrix $M'$, with row $Y'=(s_i)_{i\ne0}$ and column
$Z'=(s_j)_{j\ne0}$. Then $Y'M'$ and $M'Z'$ have entries in $R$,
and $Y'M'Z'=m_{00}-c$. Appending a one-by-one block $m_{00}$ to
$-M'$, and a coordinate $1$ to both $Y'$ and $Z'$, gives \eqref{equation-51}.
This proves the equivalence, including the finite certificate.
It is the Stacks Lemma 10.107.11 and Remark 10.107.12
\cite{StacksEpicMatrices};
the argument also supplies the final elementary step omitted there.
The coefficient-ring lemma proves that failure of this test precludes uniqueness.

Assume henceforth \eqref{equation-51}. Since $C\otimes_R C\simeq C$, one has
$C\otimes_R(C/R)=0$. The module $E$ is a quotient of $(C/R)[X]$,
and $\mathcal F\to E$ is surjective. Right exactness therefore gives
\begin{equation}
 \mathcal F\otimes_R E=E\otimes_R\mathcal F=0,
                         \qquad E\otimes_R E=0.          \label{equation-58}
\end{equation}
In particular $V_0=W_0=0$, while
$U_0=\mathcal F\otimes_R\mathcal F=C[X,Y]$ by coefficient
multiplication. Hence $Z_0=U_0$ as claimed.

Here is a fixed, functorial free construction, so that the
transgressions in \eqref{equation-55} are unambiguous. For any $R$-module $M$,
let $P_0(M)$ be the free $R$-module on its underlying set, augmented
by $[m]\mapsto m$. Take $P_1(M)$ to be free on the elements of the
kernel, with its evident map to $P_0(M)$, and repeat on successive
kernels. A module map sends the indicated generators to their
images at each stage, making $P_\bullet$ functorial. Every generator
of a kernel is a finite linear relation in the preceding free
module. Thus this construction is explicit even without finite
generation or any choice of a tensor presentation of $\Omega$.

Write $P=P_\bullet(\mathcal F)$, $Q=P_\bullet(E)$, and
$u=P_\bullet(\varepsilon):P\to Q$. For nonnegative vertical degree
$q$, form three complexes of free $R$-modules
\begin{equation}
\begin{aligned}
 \mathcal T_{0,\bullet}&=\operatorname{Tot}(P\otimes_RP),\\
 \mathcal T_{1,\bullet}&=\operatorname{Tot}(P\otimes_RQ)
                         \oplus\operatorname{Tot}(Q\otimes_RP),\\
 \mathcal T_{2,\bullet}&=\operatorname{Tot}(Q\otimes_RQ).
\end{aligned}                                           \label{equation-59}
\end{equation}
Here $\operatorname{Tot}(P\otimes Q)_q=
\bigoplus_{a+b=q}P_a\otimes_RQ_b$, with vertical differential
$d=d_P\otimes1+(-1)^a1\otimes d_Q$ on that summand.
Define the horizontal chain maps
\begin{equation}
 h_0=(1\otimes u,u\otimes1),\qquad
 h_1=(u\otimes1,-1\otimes u).                            \label{equation-60}
\end{equation}
They commute with $d$ and satisfy $h_1h_0=0$.
Their induced maps on vertical homology are exactly $a_q,b_q$.
The vertical homology identifications with Tor follow by the free
double-resolution argument, as in Stacks Lemma 10.75.5
\cite{StacksTor}.

For $q=0,1$ and $[z]\in Z_q$, choose a vertical cycle
$z\in\mathcal T_{0,q}$ representing it. Since $a_q[z]=0$,
there exists $w\in\mathcal T_{1,q+1}$ such that
\[
                dw=(-1)^{q+1}h_0z.
\]
The element $(-1)^{q+1}h_1w$ is a vertical cycle in
$\mathcal T_{2,q+1}$. Define
\begin{equation}
        \delta_q[z]=[(-1)^{q+1}h_1w]
                  \pmod{\operatorname{im}b_{q+1}}\in T_{q+1}.
                                                        \label{equation-61}
\end{equation}
Changing $w$ by a vertical cycle changes this class by an image of
$b_{q+1}$. Changing $z$ to $z+dv$, with
$v\in\mathcal T_{0,q+1}$, permits replacing $w$ by
$w+(-1)^{q+1}h_0v$, which leaves $h_1w$ unchanged.
This proves independence of choices and $R$-linearity. Standard
comparison maps, or the same calculation with homotopic chain maps,
identify these definitions for any other free resolutions; the fixed
functorial resolution already defines them without that ambiguity.
For $q=0$ an arbitrary coefficient polynomial $H\in C[X,Y]$ is
represented by any preimage in $\mathcal T_{0,0}$, and \eqref{equation-61}
specifies its obstruction using one finite chain lift $w$.

We now identify what these invariants measure. Give
$\mathcal T_{p,q}$ total homological degree $q-p$, and set
\begin{equation}
                      D=d+(-1)^q h.                     \label{equation-62}
\end{equation}
The commutation $dh=hd$ and $h^2=0$ give $D^2=0$.
This total complex has homology in degree $n$ equal to
$\operatorname{Tor}_n^R(A,A)$, where $A=\ker\varepsilon=\Omega(R,S)$,
and zero homology in negative degrees. We verify this assertion
without a convergence assumption. The two-column total complex
$L$ formed from $P\xrightarrow{u}Q$, in degrees $q-p$ with
$p=0,1$, has $L_n=P_n\oplus Q_{n+1}$ and differential \eqref{equation-62}.
The exact sequence $0\to A\to\mathcal F\to E\to0$ shows
$H_0(L)=A$ and all other homology zero. Also $L_n=0$ for $n<-1$,
and $L_{-1}=Q_0$ is free. The surjection $L_0\to L_{-1}$ therefore
splits, separating a contractible two-term complex. The remainder
is a nonnegative projective resolution of $A$: its degree-zero
module is a direct summand of the free module $L_0$, and its
higher modules are free. Tensoring two copies of this decomposition
shows that $L\otimes_RL$ computes $\operatorname{Tor}_\bullet^R(A,A)$.

There is an explicit sign isomorphism from $L\otimes_RL$ to the
three-column complex \eqref{equation-59}--\eqref{equation-62}. On the component whose
first factor has column $p_1\in\{0,1\}$ and whose second factor
has vertical resolution degree $b$, multiply by $(-1)^{p_1b}$.
For the second vertical differential this changes its sign
$(-1)^{a-p_1}$ to $(-1)^a$; for the first horizontal differential
it changes $(-1)^a$ to $(-1)^{a+b}$; for the second horizontal
differential it leaves $(-1)^{a+b-p_1}$, which is the sign in
\eqref{equation-60}--\eqref{equation-62}. The first vertical sign is unchanged.
This verifies the chain isomorphism term by term and proves the
homology assertion. In particular
\begin{equation}
                 H_0(\operatorname{Tot}\mathcal T)=A\otimes_RA,
                 \qquad H_{-1}(\operatorname{Tot}\mathcal T)=0.
                                                        \label{equation-63}
\end{equation}
The edge map obtained by taking the column-zero component and
then its vertical homology is the natural map
$A\otimes_RA\to\mathcal F\otimes_R\mathcal F=C[X,Y]$.
Indeed $H_0(L)=A$ is identified with the kernel of $\varepsilon$
by the augmentation to $\mathcal F$. Thus this edge map is
precisely $\theta_2$.

For completeness the low-degree exactness can now be checked
directly. A total degree-zero cycle consists of
\begin{equation}
 x_0\in\mathcal T_{0,0},\quad
 x_1\in\mathcal T_{1,1},\quad
 x_2\in\mathcal T_{2,2},\qquad
             dx_1=-h_0x_0,\quad dx_2=h_1x_1.             \label{equation-64}
\end{equation}
It follows from \eqref{equation-61} that its edge value lies in
$\ker\delta_0$. Conversely, if $\delta_0[x_0]=0$, choose $x_1$
solving the first equation. The zero class of $-h_1x_1$ modulo
$\operatorname{im}b_1$ allows adding a vertical cycle to $x_1$
so that $h_1x_1$ becomes a vertical boundary; then choose $x_2$
solving the second equation. Therefore
$\operatorname{im}\theta_2=\ker\delta_0$.
If $y\in\mathcal T_{2,1}$ is a vertical cycle, it is a total
degree-minus-one cycle. By \eqref{equation-63}, it is the boundary of a
triple $(x_0,x_1,x_2)$ of total degree zero. The column-one
equation is $dx_1=-h_0x_0$, and the column-two equation is
$-h_1x_1+dx_2=y$. Thus $\delta_0[x_0]=[y]$ in $T_1$.
Every class in $T_1$ has such a representative, proving surjectivity
and the canonical cokernel identification
\begin{equation}
                    \operatorname{coker}\theta_2\simeq T_1.
                                                        \label{equation-65}
\end{equation}

If a total degree-zero cycle has zero edge value, its $x_0$ is a
vertical boundary. Subtracting a total boundary makes $x_0=0$.
Then $x_1$ is a vertical cycle, and \eqref{equation-64} says that its class
lies in $\ker b_1$. Different ways to remove $x_0$ change that
class by an image of $a_1$. This defines a canonical surjection
\begin{equation}
                       \ker\theta_2\longrightarrow K_1. \label{equation-66}
\end{equation}
Surjectivity follows by starting with a representative of a class
in $\ker b_1$ and lifting $h_1x_1$ to $dx_2$.
If the class in $K_1$ is zero, subtracting further total boundaries
makes $x_1=0$ as well. What remains is a vertical cycle
$x_2\in\mathcal T_{2,2}$. Such a pure column-two cycle is a total
boundary exactly when, modulo a vertical boundary, it is of the form
$h_1w$, where
\begin{equation}
 dz=0,\quad z\in\mathcal T_{0,1},\qquad
 w\in\mathcal T_{1,2},\quad dw=h_0z.                     \label{equation-67}
\end{equation}
Indeed the column-zero and column-one components of a boundary
of total degree one must vanish, giving precisely these equations;
the column-two component is $h_1w$ plus a vertical boundary.
Changing $w$ by a cycle contributes $\operatorname{im}b_2$,
and \eqref{equation-67} is exactly the defining lift for $\delta_1$.
Consequently \eqref{equation-66} has kernel $K_2$, with a canonical exact
sequence
\begin{equation}
             0\longrightarrow K_2\longrightarrow\ker\theta_2
                     \longrightarrow K_1\longrightarrow0.   \label{equation-68}
\end{equation}
This proves that uniqueness is equivalent to $K_1=K_2=0$.

All coefficient vectors of integer-valued polynomials belong to
$C$ by Theorem~\ref{result-19}. The two explicit arrays in \eqref{equation-56} are exactly the
coefficient polynomials of $F(X+Y)$ and $F(XY)$.
Equations \eqref{equation-64}--\eqref{equation-65} say that they lift exactly when their
$\delta_0$ classes vanish. If \eqref{equation-68} is zero, these lifts are
unique and recover the requested tensors. This proves both
directions of \eqref{equation-56} for every original inclusion $R\subset S$.

Every complex used to define the criterion is a free complex built
from $\mathcal F$, $E$ and $\varepsilon$, which themselves come
directly from \eqref{equation-50}. Relations in the free resolutions and the
lifts in \eqref{equation-61}, \eqref{equation-64}, and \eqref{equation-67} have finite support.
Only $U_1,V_1,W_1$, $V_2,W_2$ and boundaries in vertical degree
three are needed. Thus the test is an explicit arithmetic syzygy
criterion, though infinite rings may give infinitely many such
tests. When $C$ is flat, positive Tor groups involving $\mathcal F$ vanish.
The criterion then has $K_1=0$, $K_2=\operatorname{Tor}_2^R(E,E)$ and
$T_1=\operatorname{Tor}_1^R(E,E)$.  This follows directly from the
definitions; no flatness is used in the general argument.


\end{proof}

\subsection{Faithful evaluation and bases that work for every extension}
The coefficient ideals also distinguish polynomial identities from
identities of polynomial functions. This distinction is necessary even
when a coefficient module is free.

\begin{theorem}[The exact ideal test for faithful evaluation]\label{result-47}
Let $R\subset S$ be any unital inclusion of commutative rings, and
let $D_d(R)$ be the Vandermonde ideals of Theorem~\ref{result-19}.
Evaluation $\Omega(R,S)\to\operatorname{App}(R,R)$ is injective if and
only if
\[
              \operatorname{Ann}_S D_d(R)=0\qquad(d\ge1).
\]
There is no finiteness, domain, reducedness or torsion-free hypothesis.
\end{theorem}
\begin{proof}
Suppose that $F=\sum_{i=0}^d a_iX^i$ vanishes at every point of $R$.
Evaluation on any $d+1$ nodes, followed by the adjugate of their
Vandermonde matrix, shows that the determinant annihilates each $a_i$.
Thus $D_d(R)a_i=0$ for every coefficient. Under the asserted
annihilator condition all coefficients vanish; a constant null
polynomial vanishes already by evaluation at zero.

Conversely, if $0\ne a\in\operatorname{Ann}_S D_d(R)$, take the
explicit monic polynomial $P_d$ constructed in the proof of
Theorem~\ref{result-19}. Every value $P_d(r)$ belongs to $D_d(R)$, so
$aP_d$ is a null polynomial. It belongs to $\Omega(R,S)$ and is
nonzero because its leading coefficient is $a$. Evaluation therefore
fails to be injective. This proves both implications, including
arbitrary zero divisors in $S$.
\end{proof}

\begin{theorem}[The bases that work for every coefficient extension]\label{result-48}
For a nonzero commutative ring $R$, the following are equivalent:
\begin{enumerate}[label=(\roman*)]
\item every residue field $R/\mathfrak m$ at a maximal ideal is infinite;
\item $\Omega(R,S)=R[X]$ for every unital inclusion $R\subset S$;
\item for every such inclusion the two polynomial substitutions have
unique tensor representatives.
\end{enumerate}
In the positive case every multivariable integer-valued polynomial has
its coefficients in $R$, and all tensor maps are the ordinary polynomial
isomorphisms. In particular every ring containing an infinite field has
this property. A single finite residue field produces a square-zero
extension for which uniqueness fails.
\end{theorem}
\begin{proof}
Assume (i), and let $F\in S[X]$ have degree at most $d$ and take $R$
into $R$. At a maximal ideal $\mathfrak m$, choose $d+1$ elements of
$R$ with distinct residues. Their Vandermonde determinant becomes a
unit in $R_{\mathfrak m}$. Solving the evaluation equations over
$R_{\mathfrak m}$ shows that every coefficient of $F$ has zero class
in $(S/R)_{\mathfrak m}$. An element of a module that vanishes at
every maximal localization is zero: otherwise a maximal ideal
containing its proper annihilator would contradict that localization.
Hence the coefficients lie in $R$. Repeating the same argument on a
rectangular grid proves the assertion in any finite number of
variables. Thus (i) implies (ii), and the ordinary polynomial tensor
isomorphisms give (iii).

For both converses, suppose $k=R/\mathfrak m=\mathbb F_q$. Form
$S=R\oplus k\epsilon$ with
\[
 (r,a\epsilon)(s,b\epsilon)=(rs,(\bar r b+\bar s a)\epsilon),
 \qquad \epsilon^2=0.
\]
The natural map $R\hookrightarrow S$ is unital. Monic division by
$X^q-X$ identifies the null polynomials on $k$, and gives the direct
sum of $R$-modules
\[
 \Omega(R,S)=R[X]\oplus\epsilon(X^q-X)k[X].
\]
In particular (ii) fails. Put $e=\epsilon(X^q-X)$. Dividing the
second summand by $X^q-X$ and taking its constant coefficient
exhibits $ke$ as a direct summand. Consequently $e\otimes e$ is
nonzero: its projection to $ke\otimes_Rke\simeq k$ is $1$.
Its polynomial image is zero because $\epsilon^2=0$, so uniqueness
for the zero polynomial fails and (iii) is false. Finally a unital
map from a field into any residue field is injective, proving the
assertion about rings containing an infinite field.
\end{proof}

\begin{corollary}[The explicit flat-coefficient Tor sequence]\label{result-49}
For an arbitrary unital inclusion $R\subset S$, let
$C=\bigcup_d(R:_S D_d(R))$, put $Q=C/R$, and define
\[
 E=\left\{\left(r\longmapsto\sum_i c_i r^i+R\right):
                    (c_i)\in\bigoplus_{i\ge0}C\right\}
       \subset\operatorname{App}(R,Q).
\]
Assume that $R\to C$ is flat and epic, and write $A=\Omega(R,S)$.
There is a canonical exact sequence
\[
 0\longrightarrow\operatorname{Tor}_2^R(E,E)
 \longrightarrow A\otimes_R A\xrightarrow{\theta_2}C[X,Y]
 \xrightarrow{\partial}\operatorname{Tor}_1^R(E,E)
 \longrightarrow0.
\]
Here $\partial$ is the explicit map $\delta_0$ of
\eqref{equation-61}: represent the coefficient polynomial by
$z\in\mathcal T_{0,0}$, choose $w\in\mathcal T_{1,1}$ with
$dw=-h_0z$, and take the class of $-h_1w$.
Thus the two operations exist uniquely precisely when
$\operatorname{Tor}_2^R(E,E)=0$ and
\[
 \partial\!\left(\sum_{i,j\ge0}\binom{i+j}{i}a_{i+j}X^iY^j\right)=0,
 \qquad
 \partial\!\left(\sum_{i\ge0}a_iX^iY^i\right)=0
\]
for every finite coefficient vector $(a_i)$ in $C$ with
$\sum_i a_i r^i\in R$ for every $r\in R$; missing coefficients are
zero. If only flatness is initially assumed, epicness is an additional
necessary condition, and epicness together with these displayed tests
is sufficient. This criterion concerns the two binary operations.
\end{corollary}
\begin{proof}
The evaluation sequence is
$0\to A\to\mathcal F=C[X]\to E\to0$. Since $C$ is epic,
$C\otimes_R Q=0$. The module $E$ is a quotient of $Q[X]$, whence
$\mathcal F\otimes_R E=E\otimes_R\mathcal F=E\otimes_RE=0$.
Flatness of $C$ makes $\mathcal F$ flat, so all positive Tor groups
with a factor $\mathcal F$ vanish. In the explicitly proved
filtration of Theorem~\ref{result-20}, this gives
\[
 K_1=0,\qquad K_2=\operatorname{Tor}_2^R(E,E),\qquad
 T_1=\operatorname{Tor}_1^R(E,E).
\]
Equations~\eqref{equation-65} and \eqref{equation-68} are therefore
exactly the asserted five-term sequence, with the stated chain formula
and its already proved independence of choices. The two displayed
polynomials are the coefficients of $F(X+Y)$ and $F(XY)$.
Exactness proves their existence tests, and the initial term proves
the uniqueness test. Finally Theorem~\ref{result-19} makes epicness
necessary if uniqueness is assumed. No flatness of $A$ has been used.
\end{proof}

\begin{example}[A free polynomial module need not have unique tensors]
Let $R=\mathbb F_2$, $S=\mathbb F_4=\mathbb F_2[\alpha]$
with $\alpha^2+\alpha+1=0$, and $q=X^2-X$.  Monic division gives
\[
 A=R[X]_{<2}\oplus qS[X].
\]
Indeed the remainder of an integer-valued polynomial has its two
coefficients in $R$, as follows by evaluating at zero and one; conversely
every polynomial in the displayed module is integer-valued.  The two
elements $q$ and $\alpha q$ are linearly independent over $R$.  Extending
them to an $R$-basis of $A$ shows that
\[
 (\alpha q)\otimes q-q\otimes(\alpha q)
\]
is nonzero in $A\otimes_R A$.  Its polynomial image is zero, so uniqueness
fails even though $A$, as a vector space over $R$, is free.
\end{example}


\section{Rational coefficients and tensor reconstruction}
\begin{theorem}[The precise tensor criterion and the valid free-module theorem]\label{result-21}
For an inclusion $R\subset S$, let
$\Omega=\{F\in S[T]:F(R)\subset R\}$ and let
\[
 \theta_2:\Omega\otimes_R\Omega\longrightarrow S[X,Y],
 \qquad G\otimes H\longmapsto G(X)H(Y).
\]
The polynomial identities of cote 153, sheet 4 determine one and only one
tensor for every $F$ if and only if $\theta_2$ is injective and both
$F(X+Y)$ and $F(XY)$ belong to its image for every $F\in\Omega$.
For a domain $R$ with fraction field $S$, freeness of $\Omega$ over
$R$ implies these conditions and, more strongly, gives isomorphisms
\[
 \theta_n:\Omega^{\otimes_R n}\xrightarrow{\sim}
 \{P\in S[T_1,\ldots,T_n]:P(R^n)\subset R\}
 \quad(n\ge1).
\]
In this case the coaddition, comultiplication, and ordinary
composition give a unique compatible $R$-plethory structure.
\end{theorem}
\begin{proof}
Existence of the requested tensors is exactly the image condition.
Uniqueness for $F=0$ already forces $\ker\theta_2=0$, and injectivity
then gives uniqueness for every $F$. This proves the exact criterion
for an arbitrary inclusion, including rings with zero divisors.

Now let $S=\operatorname{Frac}(R)$. If $R$ is finite it is a field,
$S=R$, and $\Omega=R[T]$, so the conclusion is immediate. Otherwise
$R$ is infinite. Let $\mathcal B$ be an $R$-basis of $\Omega$.
Since $R[T]\subset\Omega\subset S[T]$, localization gives
$S\otimes_R\Omega=S[T]$, and $\mathcal B$ is an $S$-basis of
$S[T]$. A polynomial $P(X,Y)$ integer-valued on $R^2$ has a unique
finite expansion $\sum_i b_i(X)h_i(Y)$ with $b_i\in\mathcal B$ and
$h_i\in S[Y]$. For every $y\in R$, the polynomial $P(X,y)$ belongs
to $\Omega$. Uniqueness of the $S$-basis expansion then implies
$h_i(y)\in R$. Therefore $h_i\in\Omega$, which proves surjectivity
of $\theta_2$. Injectivity follows because the tensor products of
the basis vectors remain linearly independent after localization.
The same argument in successive variables gives every $\theta_n$.

Transport the polynomial substitutions $T\mapsto X+Y$ and
$T\mapsto XY$ across $\theta_2$. The biring identities hold in
$S[X_1,\ldots,X_n]$, so injectivity of $\theta_n$ proves them in
the tensor powers. Evaluation at constants, negation, and the
inclusion $R[T]\subset\Omega$ provide the remaining biring maps.
Composition closes because an integer-valued polynomial sends
$R$ into $R$, and its compatibility identities are ordinary
substitution identities, again transported across the injective
$\theta_n$. Hence it defines the plethory structure. This proves
the manuscript's free-module argument under the fraction-field
hypothesis it originally had, but not under the later, weaker
hypothesis. Elliott's Theorem 12, pp.~56--57, discusses the related
weak polynomial compositeness condition \cite{Elliott2010}.
The branch obstruction below gives failures of multiplicative existence.
\end{proof}

\begin{theorem}[Complete answer for a finite base ring]\label{result-22}
Let $R\ne0$ be a finite commutative ring and let $R\subset S$ be a
unital inclusion. For the module $\Omega$ and map $\theta_2$ of Theorem~\ref{result-21},
the following are equivalent: $\theta_2$ is injective; the identities
on 153/004 determine unique coaddition and comultiplication tensors;
and $S=R$. In the last case $\Omega=R[T]$ with its usual plethory
structure. This classifies arbitrary extensions of a finite base,
including extensions that are infinite or have zero divisors.
\end{theorem}
\begin{proof}
Set $P(T)=\prod_{r\in R}(T-r)$ and $d=|R|$. This is a monic
polynomial of degree $d$ vanishing on $R$. Division by $P$, valid over
any commutative coefficient ring because $P$ is monic, gives
\[
 \Omega=E\oplus P S[T],\qquad
 E=\{H\in S[T]:\deg H<d,\ H(R)\subset R\}.
\]
Indeed $F=H+PQ$ is integer-valued on $R$ exactly when $H$ is.
The summand $SP$ of $PS[T]$ is itself split: project $Q$ to its
constant coefficient. Thus $S\otimes_R S$ embeds as a direct summand
of $\Omega\otimes_R\Omega$ by $s\otimes t\mapsto(sP)\otimes(tP)$.
On this summand $\theta_2$ is
\[
 s\otimes t\longmapsto stP(X)P(Y).
\]
Injectivity of $\theta_2$ forces multiplication
$\mu:S\otimes_R S\to S$ to be injective. It is always surjective,
so it is an isomorphism.

Write $J$ for the Jacobson radical of $R$. It is nilpotent: every
$x\in J$ has repeating powers, say $x^a=x^{a+b}$, and $1-x^b$ is a
unit, whence $x^a=0$; a finite generating set of this nil ideal gives
$J^N=0$ by the pigeonhole principle. The Chinese remainder theorem
identifies $R/J$ with a finite product of finite fields. After base
change, multiplication for $T=S/JS$ over $R/J$ is still an
isomorphism. On a field factor $k$, either $T_k=0$, or choose a
$k$-linear functional $\ell:T_k\to k$ with $\ell(1)=1$.
For every $a\in T_k$, injectivity of multiplication gives
$a\otimes1=1\otimes a$. Applying $\ell\otimes1$ yields
$a=\ell(a)1$. Thus every nonzero factor equals $k$ and
$R/J\to S/JS$ is surjective. Consequently $S=R+JS$.
Successive substitution gives $S=R+J^nS$ for all $n$, and $J^N=0$
gives $S=R$. Conversely, if $S=R$ then $\Omega=R[T]$ and every
$\theta_n$ is the standard polynomial-ring isomorphism. Theorem~\ref{result-21} now
proves the equivalence with the requested uniqueness statement.
This assertion concerns uniqueness of the polynomial tensor
expansions; it does not infer nonexistence of every conceivable
abstract biring structure merely from a nonzero tensor kernel.
\end{proof}

\begin{theorem}[All field bases and the rational-coefficient reduction]\label{result-23}
Let $R$ be an infinite domain, $K=\operatorname{Frac}(R)$, and let
$R\subset S$ with $S$ torsion-free as an $R$-module. Inside
$K\otimes_R S$ define the overring
\[
 S_{\mathrm{rat}}=S\cap(K\cdot1).
\]
Then
\[
 \Omega(R,S)=\{F\in S_{\mathrm{rat}}[T]:F(R)\subset R\}.
\]
Moreover $\ker\theta_2$ is exactly the $R$-torsion submodule of
$\Omega\otimes_R\Omega$. In particular, for an infinite field base
$R$, every unital extension $R\subset S$ gives $\Omega=R[T]$ and
unique standard co-operations. Together with Theorem~\ref{result-22}, this is a
complete classification for every field base: an infinite field
allows every extension; a finite field allows unique polynomial
tensor expansions exactly for the identity extension.
\end{theorem}
\begin{proof}
Torsion-freeness embeds $S$ in $K\otimes_R S$. If $F$ has degree
at most $d$, choose distinct $r_0,\ldots,r_d\in R$. Its coefficient
vector is obtained from $(F(r_0),\ldots,F(r_d))\in R^{d+1}$ by the
inverse Vandermonde matrix over $K$. Therefore each coefficient
belongs to $K\cdot1$ as well as to $S$, proving the first equality.
The converse inclusion is immediate. Since
$R[T]\subset\Omega\subset K[T]$, localization identifies
$K\otimes_R\Omega$ with $K[T]$; after localization $\theta_2$ is the
injective standard map $K[X]\otimes_KK[Y]\to(K\otimes_RS)[X,Y]$.
Its kernel is therefore torsion. Conversely a torsion tensor maps
to zero because the target $S[X,Y]$ is $R$-torsion-free. This proves
the asserted equality. If $R=K$ is an infinite field, then
$S_{\mathrm{rat}}=R$ and every polynomial in $R[T]$ preserves $R$,
which proves the field assertion.
\end{proof}

\begin{theorem}[The free-module repair for every torsion-free extension of an infinite domain]\label{result-24}
Let $R$ be an infinite domain and $R\subset S$ a unital inclusion
with $S$ torsion-free over $R$. If
$\Omega=\{F\in S[T]:F(R)\subset R\}$ is free over $R$, then
for every $n\ge1$ the natural map is an isomorphism
\[
 \Omega^{\otimes_R n}\xrightarrow{\sim}
 \{P\in S[T_1,\ldots,T_n]:P(R^n)\subset R\}.
\]
Consequently the two requested tensors and the compatible plethory
structure exist uniquely. This extends Theorem~\ref{result-21} from a fraction field
to every torsion-free overring of an infinite domain; the finite-base theorem shows why the word infinite is essential.
\end{theorem}
\begin{proof}
By Theorem~\ref{result-23} replace the coefficient ring by
$L=S_{\mathrm{rat}}\subset K=\operatorname{Frac}(R)$.
The same reduction holds in any finite number of variables:
use a rectangular grid of distinct elements of $R$ and invert
the Vandermonde matrix in each variable. Every coefficient of an
integer-valued multivariable polynomial then lies in $K\cdot1$
and hence in $L$.

Let $\mathcal B$ be an $R$-basis of $\Omega$. Localization shows
that it is a $K$-basis of $K[T]$. It is in fact an $L$-basis of
$L[T]$: independence follows after embedding in $K[T]$, and
spanning follows because every monomial $T^n$ already lies in
$\Omega$ and thus has a finite expansion with coefficients in $R$.
Now expand $P(X,Y)\in L[X,Y]$, integer-valued on $R^2$, uniquely
as $\sum_i b_i(X)h_i(Y)$ using this $L$-basis. The coefficients
$h_i$ lie in $L[Y]$. At every $y\in R$, the polynomial $P(X,y)$
belongs to $\Omega$, so uniqueness in the $K$-basis forces
$h_i(y)\in R$. Thus $h_i\in\Omega$. This gives surjectivity.
The tensors of the $R$-basis remain linearly independent over $K$,
which proves injectivity. Iteration gives the assertion for every
$n$. The same transport of polynomial identities as in Theorem~\ref{result-21}
then gives the biring and plethory axioms, proving the conclusion.
\end{proof}

\begin{theorem}[A finite dual-basis criterion, including every projective case]\label{result-25}
In the infinite-domain, torsion-free-extension setting of Theorem~\ref{result-23},
assume the following condition on the $R$-module $\Omega$: for every
finite subset $E\subset\Omega$ there are finitely many
$b_i\in\Omega$ and $R$-linear maps $\lambda_i:\Omega\to R$ such that
\[
 x=\sum_i b_i\lambda_i(x)\qquad(x\in E).
\]
Then every multivariable tensor map in Theorem~\ref{result-24} is an isomorphism,
and $\Omega$ has the unique compatible plethory structure.
In particular this holds if $\Omega$ is projective over $R$;
freeness is unnecessary.
\end{theorem}
\begin{proof}
Reduce to the coefficient overring $L\subset K$ by Theorem~\ref{result-23} and
its grid argument. Every $\lambda_i$ extends uniquely to a
$K$-linear functional on $K\otimes_R\Omega=K[T]$.
It maps $L[T]$ into $L$: its value on each monomial is in $R$,
and linearity gives the assertion for arbitrary coefficients in $L$.

Let $P(X,Y)\in L[X,Y]$ be integer-valued on $R^2$. Choose a nonzero
$d\in R$ clearing the denominators of its finitely many coefficient
polynomials in $X$, so those coefficients of $dP$ belong to $R[X]$
inside $\Omega$. Choose the finite dual-basis identity on this finite
set. After extension to $K$, it holds on their $K$-linear span,
so
\[
 P(X,Y)=\sum_i b_i(X)h_i(Y),\qquad
 h_i(Y)=\lambda_{i,K}\bigl(P(X,Y)\bigr).
\]
The preceding observation gives $h_i\in L[Y]$. At $y\in R$,
$P(X,y)\in\Omega$, so $h_i(y)=\lambda_i(P(X,y))\in R$.
Thus $h_i\in\Omega$, proving surjectivity.

For injectivity, let $z=\sum_j x_j\otimes y_j$ have zero polynomial
image and choose a finite dual-basis identity on the $x_j$. Then
\[
 z=\sum_i b_i\otimes\left(\sum_j\lambda_i(x_j)y_j\right).
\]
Applying $\lambda_{i,K}$ in the first variable to the zero polynomial
image gives $\sum_j\lambda_i(x_j)y_j=0$ in $K[Y]$, hence in $\Omega$.
Thus $z=0$. The same argument one variable at a time proves both
assertions for all tensor powers. Polynomial identities then give
the plethory as in Theorem~\ref{result-21}.

Finally a projective module is a summand of a free module. Embed
$\Omega$ in a free module and project back. The finitely many free
coordinates occurring in a finite subset $E$ give coordinate
functionals $\lambda_i$ on $\Omega$ and projected free basis vectors
$b_i$ with exactly the required identity. This verifies the
hypothesis rather than assuming that an arbitrary torsion-free
module has a dual basis.
\end{proof}

\section{Localization and the punctured spectrum}
\begin{theorem}[Polynomial localization in every Noetherian dimension]\label{result-26}
Let $R$ be a Noetherian domain, $K=\operatorname{Frac}(R)$,
and $R\subset L\subset K$ an overring. Localization holds
without a dimension restriction, and for localization at any
multiplicatively closed subset $U\subset R\setminus\{0\}$:
\[
 U^{-1}\Omega_n(R,L)=\Omega_n(U^{-1}R,U^{-1}L).
\]
The inclusion from left to right holds for every domain, without
Noetherianity.
\end{theorem}
\begin{proof}
We prove the inclusion by induction on the total degree $d$ of
$F\in\Omega_n(R,L)$. A constant such polynomial belongs to $R$.
For $s\in U$ set
\[
 G_s(X_1,\ldots,X_n)
   =s^dF(X_1,\ldots,X_n)-F(sX_1,\ldots,sX_n).
\]
All terms of total degree $d$ cancel; hence $G_s$ has smaller
degree and still takes $R^n$ to $R$. Every tuple in $(U^{-1}R)^n$
can be written $a/s$ with a common $s\in U$ and $a\in R^n$.
The induction hypothesis and
\[
 s^dF(a/s)=G_s(a/s)+F(a)
\]
show that $F(a/s)\in U^{-1}R$. Dividing $F$ by an additional
element of $U$ preserves this property. This proves the inclusion
without any approximation or dimension assumption.

For the converse, take a polynomial with coefficients in
$U^{-1}L$ and values in $U^{-1}R$. Clear its finitely many
coefficient denominators by $t\in U$ to obtain $F\in L[X]$,
where $X$ denotes all $n$ variables. Choose $d\ne0$ with
$dF\in R[X]$. The ideal $(dF(r):r\in R^n)$ is finitely generated
and is contained in $dU^{-1}R$ after localization.
Multiplying by one element $u\in U$ puts every generator in $dR$,
so $uF\in\Omega_n(R,L)$ and the original polynomial is
$(uF)/(ut)$. This proves the converse with one denominator for all values.
\end{proof}

\begin{theorem}[The punctured-spectrum coefficient ring and the depth test]\label{result-27}
Let $(R,\mathfrak m)$ be a nonfield Noetherian local domain with
fraction field $K$. Define the explicitly specified overring
\[
 R^\sharp=\{x\in K:\mathfrak m^Nx\subset R
                    \text{ for some }N\ge0\}.
\]
Then
\[
 R^\sharp=\bigcap_{\mathfrak p\ne\mathfrak m}R_{\mathfrak p}
 \quad\hbox{inside }K,
 \qquad
 R^\sharp=R\ \Longleftrightarrow\ \operatorname{depth}R\ge2.
\]
If the residue field is finite, and $R\subset L\subset K$ is an
overring, put $E=L\cap R^\sharp$. Then
\[
 \Omega_n(R,L)=\Omega_n(R,E)\quad(n\ge1),
 \qquad
 \Omega_1(R,L)=R[X]\ \Longleftrightarrow\ E=R.
\]
Thus the ring controlling whether nonpolynomial operations can
occur is determined without referring to tensor maps.
\end{theorem}
\begin{proof}
The displayed set is closed under addition and multiplication,
using a common sufficiently high power of $\mathfrak m$.
For $x\in K$ its nonzero denominator ideal
$I_x=\{r\in R:rx\in R\}$ detects membership in every
$R_{\mathfrak p}$. If $x$ belongs to all localizations with
$\mathfrak p\ne\mathfrak m$, either $I_x=R$ or
$\sqrt{I_x}=\mathfrak m$. Noetherianity then gives
$\mathfrak m^N\subset I_x$. Conversely if
$\mathfrak m^Nx\subset R$ and $\mathfrak p\ne\mathfrak m$,
choose $a\in\mathfrak m\setminus\mathfrak p$; the equality
$x=(a^Nx)/a^N$ places $x$ in $R_{\mathfrak p}$.
This proves the intersection formula.

Suppose $a,b\in\mathfrak m$ is a regular sequence. If
$R^\sharp\ne R$, choose $x\notin R$ with
$\mathfrak m x\subset R$: starting with any element of
$R^\sharp\setminus R$, use a least annihilating power and a
suitable element of the preceding power. Then $ax,bx\in R$ and
$b(ax)=a(bx)$. The nonzerodivisor $b$ on $R/aR$ forces
$ax\in aR$; cancellation by nonzero $a$ gives $x\in R$,
a contradiction. Hence depth at least two implies $R^\sharp=R$.
Conversely, if no regular sequence of length two exists, choose
$0\ne a\in\mathfrak m$. Every element of $\mathfrak m$ is a
zerodivisor on $R/aR$. The zerodivisors of this finite Noetherian
module form the union of its finitely many associated primes.
Prime avoidance forces $\mathfrak m$ to be one of those primes.
Thus there is $z\notin aR$ with $\mathfrak m z\subset aR$;
$x=z/a$ belongs to $R^\sharp\setminus R$.
The associated-prime facts used in this short argument are the
Stacks Project, Section 10.63, Lemmas 10.63.5 and 10.63.9,
and the depth-zero formulation in Lemma 10.72.5.

For the coefficient assertion take $F\in\Omega_n(R,L)$ and a
prime $\mathfrak p\ne\mathfrak m$. The domain $R/\mathfrak p$
is not a field and is therefore infinite. Choose as many elements
of $R$ with distinct residues modulo $\mathfrak p$ as the degree
in each variable requires. The corresponding Vandermonde
matrices are invertible over $R_{\mathfrak p}$. Applying them
to the values of $F$ proves that every coefficient lies in
$R_{\mathfrak p}$. The intersection formula puts every coefficient
in $E$, proving the module equality. If $E=R$ this is the usual
polynomial ring.

Conversely let $x\in E\setminus R$. Choose $N$ with
$\mathfrak m^Nx\subset R$. The quotient $R/\mathfrak m^N$ is
finite: its finite filtration has finite-dimensional vector-space
quotients over the finite residue field. The monic polynomial
$P$ formed from all of its residue representatives satisfies
$P(R)\subset\mathfrak m^N$. Hence $xP\in\Omega_1(R,L)$, but
its leading coefficient $x$ is not in $R$. This proves the last
equivalence.
\end{proof}

\section{Integer-valued polynomials over analytically irreducible local domains}
\label{sec:integer-free}
This section proves freeness by combining two free modules across the
conductor square. The module over the normalization has the familiar
valuation-ordering basis, and the module over the conductor quotient
is a union of finite free function modules with split refinement maps.
The gluing step is the free-module case of Milnor's patching theorem
\cite[\S2, Theorems~2.1--2.3]{Milnor1971KTheory}; for the statement
including arbitrary projective modules, see
\cite[Theorem~2.3, p.~4]{ChenLe2020Milnor}. We include its elementary
block-matrix proof to make the infinite-rank construction explicit.

\begin{lemma}[Milnor patching for free modules]\label{result-28}
Let $V\to D$ be a surjective ring homomorphism, let $k\to D$ be a
second homomorphism, and put $R=V\times_D k$. Fix a set $I$ and an
isomorphism $\alpha:D^{(I)}\to D^{(I)}$. Define
\[
 M_\alpha=\{(b,c)\in V^{(I)}\oplus k^{(I)}:
                  \bar b=\alpha(\bar c)\}.
\]
Then $M_\alpha$ is projective over $R$, even for infinite $I$.
More precisely $M_\alpha\oplus M_{\alpha^{-1}}$ is free over $R$.
\end{lemma}
\begin{proof}
All matrices below act on direct sums, so each column has finite
support. Matrix products are well defined: applying a matrix to a
finite-support vector again gives a finite-support vector.
In the endomorphism ring of $D^{(I)}$ one has
\[
 \begin{pmatrix}\alpha&0\\0&\alpha^{-1}\end{pmatrix}
 =
 \begin{pmatrix}1&\alpha\\0&1\end{pmatrix}
 \begin{pmatrix}1&0\\-\alpha^{-1}&1\end{pmatrix}
 \begin{pmatrix}1&\alpha\\0&1\end{pmatrix}
 \begin{pmatrix}0&-1\\1&0\end{pmatrix}.
\]
Direct block multiplication proves this equality. Lift the finitely
many entries in each column of $\alpha$ and $\alpha^{-1}$ arbitrarily
through $V\to D$, choosing zero lifts of zero entries. Replace the
two off-diagonal matrices in the displayed product by these lifts.
Each elementary block matrix is invertible over $V$, regardless of
whether the lifted off-diagonal matrices are inverses. Their product
is an automorphism $\Gamma$ of $V^{(I\sqcup I)}$ reducing to
$\operatorname{diag}(\alpha,\alpha^{-1})$.

The direct sum $M_\alpha\oplus M_{\alpha^{-1}}$ is precisely the
module patched with that block diagonal gluing map. Applying
$\Gamma^{-1}$ to its $V$-component and the identity to its
$k$-component identifies it with
\[
 \{(b,c)\in V^{(I\sqcup I)}\oplus k^{(I\sqcup I)}:
                       \bar b=\bar c\}=R^{(I\sqcup I)}.
\]
The last equality is coordinatewise; all supports are finite.
Thus $M_\alpha$ is a direct summand of a free module, as claimed.
\end{proof}

\begin{theorem}[Every finite-conductor subring of one discrete valuation ring]\label{result-29}
Let $V$ be a discrete valuation ring with fraction field $K$,
uniformizer $\pi$, and finite residue field of cardinality $Q$.
Let $R\subset V$ be a unital subring containing $\pi^mV$ for some
$m\ge1$. Then $R$ is a one-dimensional Noetherian local domain,
$V$ is its finite normalization, and
$A=\operatorname{Int}(R)$ is free over $R$, of countably infinite
rank. Consequently
all multivariable tensor maps are isomorphisms and $A$ has the
unique compatible plethory. This includes incomplete and mixed
characteristic rings, and arbitrary conductor gaps.
\end{theorem}
\begin{proof}
Put $J=\pi^mV$ and $R_0=R/J\subset D=V/J$. The ring $D$ is finite.
Its finitely many representatives generate $V$ over $R$, since
$J\subset R$. If $x\in R$ has valuation zero, multiplication by
its residue is injective on the finite ring $R_0$, hence surjective.
There is $y\in R$ with $xy=1+j$, $j\in J$. The inverse of $1+j$
in $V$ lies in $1+J\subset R$, so $x$ is a unit of $R$.
Thus $R$ is local with maximal ideal $R\cap\pi V$.
Also $V\subset\operatorname{Frac}(R)$ because $\pi^mV\subset R$.

For any nonzero ideal $H\subset R$, choose $h\in H$ of least
valuation. Then $H\subset hV$ and $JhV\subset H$.
The quotient $hV/JhV$ is finite, while $JhV$ is a finite $R$-module
because $V$ is. Finitely many representatives therefore generate
$H$ over $R$. This proves that $R$ is Noetherian. Every nonzero
prime has finite domain quotient, since it contains some $JhV$;
that quotient is a field. The nonzero maximal ideal gives dimension
one. Finally $V$ is finite, hence integral, over $R$ and is
integrally closed in $K$, so it is exactly the normalization.

Set $B=\operatorname{Int}(R,V)$. It is free over $V$ by the
valuation-ordering construction credited to Bhargava in
\cite[Definition~4 and Proposition~11, author version pp.~7--8]{Chabert2014RegularBases}.
We give the construction in the form needed here.
Choose $a_0\in R$. Given distinct
$a_0,\ldots,a_{n-1}$, choose $a_n\in R$ minimizing the finite value
\[
 v_\pi\left(\prod_{i<n}(a_n-a_i)\right).
\]
Such a minimum exists: $R$ is infinite, so some value is finite,
and finite valuations are nonnegative integers. Put
\[
 b_n(X)=\frac{\prod_{i<n}(X-a_i)}{\prod_{i<n}(a_n-a_i)}.
\]
Minimality gives $b_n(R)\subset V$, and $b_n(a_n)=1$.
Successive evaluation at the nodes proves
$B=\bigoplus_{n\ge0}Vb_n$, as follows: after subtracting the previous terms, evaluation at
$a_n$ makes the next coefficient a member of $V$ because $b_n(a_n)=1$.
After the degree of the polynomial has been reached, the remainder
vanishes at more points than its degree and is zero.

We next compute $B/JB$ by congruence-class approximation. This
use of polynomial approximations to locally constant functions
belongs to the Stone--Weierstrass theory of integer-valued
polynomials; see \cite[\S4, pp.~52--53]{CahenGrazziniHaouat1982}
and \cite[Proposition~4.3, author version pp.~18--19]{CahenChabert1999Skolem}.
The following explicit indicators suffice for our quotient.
Let $E_N$ be the finite set of
congruence classes modulo $\pi^NV$ met by $R$, and let
$\mathcal F(R,D)$ be the union of the function rings $D^{E_N}$
under pullback along refinement. Evaluation induces
\[
 B/JB\xrightarrow{\sim}\mathcal F(R,D).
\]
Injectivity is immediate: a polynomial with values in $J$ is
$\pi^m$ times one in $B$. Every polynomial in $B$ is constant
modulo $J$ on a sufficiently fine congruence partition. Indeed,
if its coefficients have valuations at least $-h$, then
$v(F(x)-F(y))\ge v(x-y)-h$ for $x,y\in V$.

For surjectivity we construct all congruence-class indicators.
The polynomial
\[
 \delta(X)=\frac{X^Q-X}{\pi}
\]
is integer-valued on $V$. If $v(x-y)=r\ge1$, factor its numerator
difference to obtain
\[
 v(\delta(x)-\delta(y))=r-1:
\]
the other factor reduces to $Q\bar x^{Q-1}-1=-1$ in the residue
field. Hence the residues of
$X,\delta(X),\ldots,\delta^{\circ(N-1)}(X)$ distinguish precisely
the classes modulo $\pi^N$. For a chosen class, let $c_j\in V$
represent its $j$th residue digit and put
\[
 H(X)=\prod_{j=0}^{N-1}
       \left(1-(\delta^{\circ j}(X)-c_j)^{Q-1}\right).
\]
Modulo $\pi$, this is its characteristic function on $V$.
A sufficiently high $Q$-power makes it the characteristic function
modulo $\pi^m$. Explicitly $H^{Q^m}$ suffices: when $H\in\pi V$
its power has valuation at least $Q^m\ge m$; when $H=1+z$,
$v(z)\ge1$, each $Q$-power increases $v((1+z)^Q-1)$ by at least
one. For the latter assertion, factor $Q$ as a power of the
residue characteristic $p$ and use the binomial formula for
$(1+z)^p$: every intermediate coefficient is divisible by $p$,
and the final term has valuation $p\,v(z)\ge v(z)+1$.
The argument also covers characteristic $p$, where these
intermediate coefficients vanish. These indicators, multiplied
by arbitrary representatives in $V$ and added finitely, prove
surjectivity.

Put $C=A/JB$. The ideal $JB=\pi^mB$ is contained in $A$; no
identification of $JA$ with $JB$ is assumed.
Its image is exactly $\mathcal F(R,R_0)\subset\mathcal F(R,D)$.
One implication follows from values in $R$. Conversely a
polynomial whose residue function has values in $R_0$ has all
its actual values in $R$, since $J\subset R$; and the preceding
indicator construction lifts every such function. Thus
\[
 C=A/JB\simeq\mathop{\mathrm{colim}}_N R_0^{E_N}.
\]
This is a free $R_0$-module. To see it directly, select one child
in $E_{N+1}$ over each class in $E_N$. Restriction to the selected
children retracts the pullback
$R_0^{E_N}\to R_0^{E_{N+1}}$, with free kernel on the unselected
children. The union is the direct sum of the first finite free
module and these finite free kernels. Its rank is countably
infinite, since the numbers $|E_N|$ are unbounded for infinite $R$.

Tensor product commutes with these unions and with their finite
free terms. Consequently multiplication gives an isomorphism
\[
 C\otimes_{R_0}D\xrightarrow{\sim}B/JB.
\]
The ring square is $R=V\times_D R_0$, and $A$ is the module
pulled back from $B$ and $C$ along this isomorphism. Both modules
are free of countable rank, so Lemma~\ref{result-28} makes $A$
a direct summand of a countably generated free $R$-module.
Kaplansky's theorem says that every projective module over a local
ring is free \cite[Theorem~2]{Kaplansky1958}; its arbitrary-rank
form is also \cite[Theorem~10.85.4, tag 0593]{StacksProjectiveLocal}.
Thus $A$ is free and countably generated. Localizing at the
nonzero elements of $R$ gives $K\otimes_R A=K[X]$, so its
rank is countably infinite. Theorem~\ref{result-25} proves the
tensor and plethory conclusions.
\end{proof}

\subsection{From analytic irreducibility to a single finite branch}
The following classical equivalence identifies the local domains to which
the conductor calculation applies.  We give the passage explicitly because
finite normalization by itself does not exclude several branches.
The forward locality implication is also the proposition of
\cite[\S2, p.~49, proof pp.~49--50]{CahenGrazziniHaouat1982}.

\begin{proposition}[Analytic irreducibility and finite local normalization]
\label{prop:analytic-one-branch}
Let $(R,\mathfrak m)$ be a one-dimensional Noetherian local domain,
let $K$ be its fraction field, and let $V$ be its integral closure in $K$.
Then $\widehat R$ is a domain if and only if $V$ is a finite $R$-module
and is local.  In that case $V$ is a discrete valuation ring and
$\pi^mV\subset R$ for some uniformizer $\pi$ and integer $m\geq1$.
\end{proposition}
\begin{proof}
Suppose first that $\widehat R$ is a domain.  It is reduced, so $R$
is analytically unramified.  The normalization-finiteness theorem
\cite[Lemma 10.162.10(4), tag 032Y]{StacksCompletionAnalytic} makes $V$
a finite $R$-module.  Completion is flat, and hence the inclusion
$V\subset K$ gives an injection
\[
 V\otimes_R\widehat R\ \hookrightarrow\
 K\otimes_R\widehat R
 =(R\setminus\{0\})^{-1}\widehat R.
\]
The localization on the right is a domain: faithful flatness keeps
every nonzero element of $R$ nonzero in $\widehat R$.  Thus
$B=V\otimes_R\widehat R$ is a domain.  It is finite over the complete
Noetherian local ring $\widehat R$, so it is a finite product of
complete local rings \cite[Lemma 10.160.2, tag 0325]{StacksCompletionAnalytic}.
A domain has no nontrivial idempotents, and therefore this product
has exactly one factor.  It follows that $B$ is local.

Reduction modulo $\mathfrak m\widehat R$ identifies
$B/\mathfrak mB$ with $V/\mathfrak mV$.  Every maximal ideal of the
finite integral $R$-algebra $V$ lies above $\mathfrak m$, so its
maximal ideals correspond to those of $V/\mathfrak mV$.  This quotient
is local because $B$ is local.  Hence $V$ is local.  It is Noetherian
by finiteness, normal by definition, and one-dimensional by integrality.
The discrete-valuation criterion
\cite[Lemma 10.119.7, tag 00PD]{StacksNormalDomains} now shows that $V$
is a discrete valuation ring.

Conversely, suppose that $V$ is finite and local.  The same criterion
makes it a discrete valuation ring.  Its maximal ideal $\mathfrak n$
is the radical of $\mathfrak mV$, so the $\mathfrak mV$-adic and
$\mathfrak n$-adic topologies are equivalent.  Exactness of completion
on finite modules sends $R\hookrightarrow V$ to
$\widehat R\hookrightarrow\widehat V$.  Here $\widehat V$ is the
completion of a discrete valuation ring in its maximal ideal, and is
again a discrete valuation ring.  In particular it is a domain, and
its subring $\widehat R$ is a domain.

Finally, choose finitely many $R$-module generators for $V$ in $K$.
Clearing their denominators gives $0\ne d\in R$ with $dV\subset R$.
In the discrete valuation ring, $dV=\pi^aV$ for some $a\geq0$.
Taking $m\geq\max(1,a)$ gives the asserted conductor inclusion.
\end{proof}

\begin{proof}[Proof of Theorem~\ref{thm:int-analytic-free}]
If $R/\mathfrak m$ is infinite, let $f\in\operatorname{Int}(R)$ have
degree at most $n$.  Choose $n+1$ elements of $R$ with distinct
residues.  Their Vandermonde determinant is a unit.  Solving for the
coefficients of $f$ from its values at these elements puts every
coefficient in $R$.  Therefore $\operatorname{Int}(R)=R[X]$ in this case.

Suppose that $R/\mathfrak m$ is finite.  By
Proposition~\ref{prop:analytic-one-branch}, its normalization $V$ is
a finite local discrete valuation ring and contains a power of its
maximal ideal inside $R$.  The residue field of $V$ is finite:
it is a quotient of the finite-dimensional $R/\mathfrak m$-algebra
$V/\mathfrak mV$.  Theorem~\ref{result-29} therefore proves freeness
of $\operatorname{Int}(R)$, by projective patching followed by
Kaplansky's theorem.

In both cases the rank is countably infinite.  In the finite-residue
case the patching construction makes the module a direct summand of
a countably generated free module, and hence countably generated.
Moreover
\[
 K\otimes_R\operatorname{Int}(R)=K[X],
\]
because $R[X]\subset\operatorname{Int}(R)\subset K[X]$.
The dimension of this vector space over $K$ is countably infinite.
The tensor and compatible plethory conclusions follow from
Theorem~\ref{result-24}, applied to the fraction-field extension.
\end{proof}


\Needspace{9\baselineskip}
\subsection{Explicit refutations of Elliott's conjectures}
The next identities directly negate the proposed antecedents in
Elliott's Propositions~2.14(1) and 2.16(1), respectively for the
residue-extension ring and the cusp ring
\cite[author version p.~12]{Elliott2011IntegerTClosure}.
Those propositions are implications and remain valid; it is the
nonflatness conjectures that Theorem~\ref{thm:int-analytic-free}
disproves. The identities give a small direct calculation behind
the two positive examples. Elliott's later announcement of cusp
freeness is credited in Section~1. An answer by user2035 to his
original 2011 question already gave an integer-valued polynomial
whose derivative at zero lies outside $t^{-1}\mathbb F_2[[t]]$
\cite[answer of 12 September 2011]{Elliott2013CuspAnnouncement}.
The paired identities below prove the required ideal membership
directly.
\begin{proposition}[Explicit witnesses for the two singular examples]
\label{prop:elliott-explicit-witnesses}
Put $f(X)=X^2+X$.
\begin{enumerate}
\item Let $D=\mathbb F_2[[t^2,t^3]]$. The polynomials
\begin{align*}
 h_2&=\frac{f^4+t^4f^2}{t^{10}},&
 h_1&=\frac{f}{t^2}+t h_2
\end{align*}
belong to $\operatorname{Int}(D)$ and satisfy
\begin{equation}\label{eq:cusp-explicit-witness}
 X^2+X=t^2h_1+t^3h_2.
\end{equation}
\item Let $D=\mathbb F_2+t\mathbb F_4[[t]]$, where
$\mathbb F_4=\mathbb F_2(\alpha)$ and $\alpha^2+\alpha+1=0$.
The polynomials
\begin{align*}
 k_2&=\frac{f^2+tf}{t^2},&
 k_1&=\frac{\alpha f^2+\alpha^2tf}{t^2}
\end{align*}
belong to $\operatorname{Int}(D)$ and satisfy
\begin{equation}\label{eq:residue-extension-explicit-witness}
 X^2+X=t k_1+\alpha t k_2.
\end{equation}
\end{enumerate}
In particular $X^2+X\in\mathfrak m\operatorname{Int}(D)$ in both
examples, where $\mathfrak m$ is the maximal ideal of $D$.
\end{proposition}
\begin{proof}
For the first ring, put $V=\mathbb F_2[[t]]$. Separating even and
odd exponents gives $D=\mathbb F_2+t^2V$, with maximal ideal
$(t^2,t^3)$. For $x=a+t^2y\in D$, where $a\in\mathbb F_2$ and
$y\in V$, put $g=f(x)/t^2=y+t^2y^2\in V$. Write
$g=b+ct+t^2z$ with $b,c\in\mathbb F_2$ and $z\in V$. Then
$(g^2+g)/t\in V$ has constant coefficient $c$. Its square has
only even exponents and is congruent to $c$ modulo $t^2$.
Consequently
\[
 h_2(x)=\left(\frac{g^2+g}{t}\right)^2\in D,
 \qquad h_1(x)=g+t h_2(x)\equiv b\pmod{t^2V}.
\]
Thus $h_1(x)\in D$ as well. Multiplication by $t^2$ proves
\eqref{eq:cusp-explicit-witness}, since the two terms $t^3h_2$
cancel. For example, at $x=t^3$ the values are
\[
 h_2(t^3)=1+t^2+t^6+t^{14},\qquad
 h_1(t^3)=t^3+t^4+t^7+t^{15}.
\]

For the second ring, put $V=\mathbb F_4[[t]]$. Its maximal ideal
is $tV=(t,\alpha t)D$: write every element of $V$ as
$A+\alpha B$ with $A,B\in\mathbb F_2[[t]]\subset D$.
If $x=a+ty\in D$, then $g=f(x)/t=y+ty^2\in V$.
Write $g\equiv b+c\alpha\pmod{tV}$ with $b,c\in\mathbb F_2$.
The relation $\alpha^2+\alpha=1$ gives
\[
 k_2(x)=g^2+g\equiv c\pmod{tV},\qquad
 k_1(x)=g+\alpha k_2(x)\equiv b\pmod{tV}.
\]
Both values therefore belong to $D$, and the last expression for
$k_1$ proves \eqref{eq:residue-extension-explicit-witness}.
All four polynomials have coefficients in the appropriate fraction
field: $t=t^3/t^2$ in the first case and
$\alpha=(\alpha t)/t$ in the second.
\end{proof}

\begin{lemma}[Freeness need not respect polynomial degree]
\label{lem:cusp-no-regular-basis}
Let $D=\mathbb F_2[[t^2,t^3]]$, $V=\mathbb F_2[[t]]$, and
$A=\operatorname{Int}(D)$. Then
\[
 A_{\le2}=D\oplus DX\oplus V(X^2+X),
\]
where $A_{\le2}$ denotes the polynomials of degree at most two.
In particular the ideal of coefficients of $X^2$ in $A_{\le2}$
is $V$, which is not a principal fractional $D$-ideal. Thus $A$
has no $D$-module basis consisting of one polynomial of each
degree, even though $A$ is free as a $D$-module.
\end{lemma}
\begin{proof}
The map $z\mapsto z^2+z$ sends $D$ onto its maximal ideal
$\mathfrak m=t^2V$. Indeed its values lie in $\mathfrak m$, and
for $w\in\mathfrak m$ the convergent series
$z=\sum_{j\ge0}w^{2^j}$ belongs to $\mathfrak m$ and satisfies
$z^2+z=w$.

Every polynomial of degree at most two has a unique expression
$a(X^2+X)+bX+c$ over the fraction field. If it is integer-valued,
evaluation at $0$ and $1$ gives $c,b\in D$. The remaining
condition is $a\mathfrak m\subset D$. For $a\ne0$, write
$a=t^r u$ with $u\in V^\times$. Since $a\mathfrak m=t^{r+2}V$,
this inclusion holds exactly when $r\ge0$: for $r=-1$ or $-2$
the left side contains $t$, and for $r<-2$ it contains an element
of negative valuation. This proves the displayed decomposition.

If $V=aD$, comparison of the least valuations gives $a\in V^\times$.
Multiplying by $a^{-1}$ would then give $D=V$, a contradiction.
Finally, if $(f_n)_{n\ge0}$ were a basis with $\deg(f_n)=n$, the
coefficient of $X^2$ in every polynomial of degree at most two
would be a $D$-multiple of the leading coefficient of $f_2$.
There can be no cancellation with higher basis elements, because
the highest degree in a finite expansion occurs only once.
The ideal $V$ would be principal, contradicting what was proved.
\end{proof}


\Needspace{8\baselineskip}
\subsection{What changes when the normalization has several branches}
\begin{theorem}[The branch criterion for finite normalization in dimension one]\label{result-30}
Let $(R,\mathfrak m)$ be a one-dimensional Noetherian local domain
with finite residue field, fraction field $K$, and finite
normalization $T\subset K$. Then the following are equivalent:
\begin{enumerate}
\item $T$ is local;
\item every map $\operatorname{Int}(R)^{\otimes_R n}
\to\operatorname{Int}(R^n)$ is an isomorphism;
\item the polynomial identities of cote 153, sheet 4 give unique coaddition
and comultiplication tensors;
\item $F(XY)$ is a sum of terms $G(X)H(Y)$ with
$G,H\in\operatorname{Int}(R)$ for every $F\in\operatorname{Int}(R)$.
\end{enumerate}
If $T$ is not local, the proof gives an explicit higher-variable
polynomial obstruction, which rules out condition 4 by finite
iteration. Thus the criterion does not merely assert failure of
projectivity or of a sufficient condition.
\end{theorem}
\begin{proof}
The finite normalization is a one-dimensional Noetherian normal
semilocal domain. Its maximal ideals
$\mathfrak p_0,\ldots,\mathfrak p_s$ all contract to
$\mathfrak m$, and all their residue fields are finite.
Its localizations are discrete valuation rings, by the discrete valuation criterion for one-dimensional normal
Noetherian local domains \cite[Lemma 10.119.7, tag 00PD]{StacksNormalDomains}. If $s=0$, it is a discrete
valuation ring. A nonzero conductor ideal exists: clear the
denominators of finitely many $R$-module generators of $T$.
In a discrete valuation ring it contains $\pi^mT$ for some $m$.
Theorem~\ref{result-29} therefore proves condition 2. The implications
$2\Rightarrow3\Rightarrow4$ are immediate from Theorem~\ref{result-21}.

Suppose now $s\ge1$. We verify the conductor and approximation
steps explicitly. For each $i$, choose $\pi_i\in T$ congruent
to a valuation-one element modulo $\mathfrak p_i^2$, and
congruent to $1$ modulo every other maximal ideal. The Chinese
remainder theorem supplies it. Thus
$v_i(\pi_i)=1$ and $v_j(\pi_i)=0$ for $j\ne i$.
Localizing at all maximal ideals shows
$\mathfrak p_i=\pi_iT$; no general claim that a nonprincipal
ideal has a chosen generator is being used.
Write the nonzero conductor as
\[
 I=(R:T)=\prod_i\mathfrak p_i^{c_i}=aT,
 \qquad a=\prod_i\pi_i^{c_i}.
\]
The equality follows by localizing its two sides at the finitely
many discrete valuation rings. Every $c_i$ is positive. Otherwise
some conductor element lies outside $\mathfrak p_i$, hence outside
$\mathfrak m$ in $R$, so it is a unit of $R$. The conductor would
then contain $1$ and force $T=R$, contrary to $s\ge1$.
In particular $a\in I\subset R$ and $v_0(a)=c_0$.

Set $r=c_0+1$ and choose $N>r c_0$.
The unit $\pi_1$ has finite order $L$ in
$(T/\mathfrak p_0^N)^*$, since this quotient is finite.
For $n\ge1$ put
\[
 u_n=a\pi_1^{Ln}\in I\subset R.
\]
Then $u_n\equiv a\pmod{\mathfrak p_0^N}$ and
$v_1(u_n)\longrightarrow\infty$. For every finite family
$\mathcal H\subset\operatorname{Int}(R)$, sufficiently large $n$
gives
\[
 h(u_n)-h(0)\in\mathfrak p_1\cap R=\mathfrak m
 \qquad(h\in\mathcal H).
\]
Indeed every nonconstant term of each of the finitely many
polynomials tends to zero for $v_1$, and both
values lie in $R$.

Use the digit-indicator construction in Theorem~\ref{result-29} at the discrete
valuation ring $T_{\mathfrak p_0}$ to produce a polynomial $H$
which is $1$ modulo $\mathfrak p_0$ on the class of $a^r$ modulo
$\mathfrak p_0^N$, and $0$ on the zero class. These classes are
distinct because $v_0(a^r)=r c_0<N$. In that construction the
residue representatives can be chosen in $T$, and its only
denominators are powers of $\pi_0$. Thus
$H\in T[\pi_0^{-1}][X]$, it is integer-valued in
$T_{\mathfrak p_0}$, and its coefficients are integral at every
other maximal ideal of $T$. Consequently $H(R)\subset T$.
Put
\[
 F=H-H(0),\qquad g=aF.
\]
Then $F(R)\subset T$, $F(0)=0$, and
$g(R)\subset aT=I\subset R$. Also
\[
 v_0(g(u_n^r))=c_0\quad\hbox{for every }n,
\]
since $u_n^r\equiv a^r\pmod{\mathfrak p_0^N}$ and $F$ has
unit residue on this class.

The polynomial $g(X_1\cdots X_r)$ is not in the image of the
$r$-fold tensor map. Otherwise write it as a finite sum of
products $\prod_{j=1}^r h_{ij}(X_j)$ and choose $n$ large for
all these factors in the displayed difference condition.
Its iterated difference on the $2^r$ points with coordinates
$0,u_n$ would be
\[
 \sum_i\prod_{j=1}^r(h_{ij}(u_n)-h_{ij}(0))\in\mathfrak m^r.
\]
The other expression gives exactly $g(u_n^r)$, because $g(0)=0$.
Every element of $\mathfrak m^r$ has $v_0\ge r>c_0$, a
contradiction. Finally, if condition 4 held, repeated binary
expansion, first separating $X_1$ from $X_2\cdots X_r$ and
then expanding every remaining factor, would put
$g(X_1\cdots X_r)$ in that image. The expansions are finite at
each of the $r-1$ steps, so no coherence or injectivity assumption
is needed for this implication. Thus condition 4 fails.
This proves all the equivalences and the negative direction.
\end{proof}

\Needspace{15\baselineskip}
\section{Characteristic ideals over Krull domains}

The characteristic-ideal product formula for Krull domains is
classical~\cite[Proposition~2.1, author pp.~2--3]{Chabert2003Factorial}.
The criterion for a regular basis in terms of principal characteristic
ideals is classical as well~\cite[Proposition~2.4 and Corollary~2.5,
author p.~4]{Chabert2003Factorial}; see also
\cite[Proposition~5]{Chabert2014RegularBases}.
The proof below also computes the intersections with an arbitrary
coefficient overring. In the Noetherian normal case, the height-one
intersection and local projectivity inputs can be found in
\cite[Tags~031T and 00NX]{StacksNormalDomains}.

\begin{theorem}\label{thm:krull-characteristic-decomposition}
Let $R$ be a Krull domain with fraction field $K$, and let
$R\subset L\subset K$. For each height-one prime $\mathfrak p$ with
finite residue field, put $q_{\mathfrak p}=|R/\mathfrak p|$ and
\[
 w_q(n)=\sum_{j\ge1}\left\lfloor\frac n{q^j}\right\rfloor,
 \qquad I_n=\prod_{\substack{\operatorname{ht}\mathfrak p=1\\
                         |R/\mathfrak p|<\infty}}
                      \mathfrak p^{-w_{q_{\mathfrak p}}(n)}.
\]
Each product has finite support. There are monic polynomials
$P_n\in R[X]$, $\deg P_n=n$, such that
\[
 \operatorname{Int}(R)=\bigoplus_{n\ge0}I_nP_n,
 \qquad
 \Omega(R,L)=\bigoplus_{n\ge0}(I_n\cap L)P_n.
\]
Every $I_n\cap L$ is an invertible fractional ideal. Consequently
$\Omega(R,L)$ is projective and all its multivariable tensor maps
are isomorphisms. The same conclusion applies to every torsion-free
coefficient extension of an infinite Krull domain, after rational
coefficient reduction. In particular it applies to every Noetherian
normal domain.
\end{theorem}

\begin{proof}
The finite-field case is immediate, so assume that $R$ is infinite.
We use the defining height-one intersection and finite-character
properties of a Krull domain.

A height-one prime $\mathfrak p$ with finite quotient is maximal and
invertible. Choose $b\in R$ with $v_{\mathfrak p}(b)=1$. There are
only finitely many other height-one primes $\mathfrak q_i$ containing
$b$. Choose $a_i\in\mathfrak q_i\setminus\mathfrak p$ and put
$a=\prod_i a_i^{v_{\mathfrak q_i}(b)}$. Then $x=a/b$ has value $-1$
at $\mathfrak p$ and nonnegative value at every other height-one
prime. The intersection formula gives $x\mathfrak p\subset R$.
The ideal $\mathfrak p\mathfrak p^{-1}$ contains $\mathfrak p$ and
$xb=a\notin\mathfrak p$, so it is $R$. This proves invertibility
without a finite-generation assumption on arbitrary prime ideals.

For each fixed finite cardinality $q$, choose $a\in R$ with
$a^q-a\ne0$. Every prime with residue field of size $q$ contains
this nonzero element, so finite character gives only finitely many
height-one such primes. Since $w_q(n)=0$ for $q>n$, the product
defining $I_n$ is finite.

We recall the local coefficient calculation. In a DVR $V$ with
residue cardinality $q$, choose representatives $r_0=0,r_1,\ldots,r_{q-1}$
and a uniformizer $\pi$. For $n=\sum n_jq^j$ put
$a_n=\sum r_{n_j}\pi^j$ and $N_n(X)=\prod_{i<n}(X-a_i)$.
For every $x\in V$, counting the nodes in its congruence class at
each level gives
\[
 v(N_n(x))\ge\sum_{j\ge1}\lfloor n/q^j\rfloor=w_q(n),
 \qquad v(N_n(a_n))=w_q(n).
\]
Thus $N_n/\pi^{w_q(n)}$ is integer-valued, and its value at $a_n$
is a unit. Successive evaluation at $a_0,a_1,\ldots$ shows that
these normalized Newton polynomials form a regular $V$-basis.
In particular the degree-$n$ leading coefficients form
$\pi^{-w_q(n)}V$. The same argument works for polynomials evaluated
only on a subring containing the chosen nodes.

Apply this calculation at a finite-residue height-one prime of $R$.
Representatives and a valuation-one element can be chosen in $R$,
so all nodes belong to $R$. At every remaining height-one prime,
$R/\mathfrak p$ is infinite; a unit Vandermonde matrix on sufficiently
many $R$-points forces every coefficient to lie in $R_{\mathfrak p}$.
Hence the leading coefficient of an integer-valued polynomial of
degree at most $n$ belongs to $I_n$.

For the finitely many primes contributing to $I_n$, use the Chinese
remainder theorem coefficient by coefficient to choose one monic
$P_n\in R[X]$ of degree $n$ congruent to the corresponding $N_n$
modulo $\mathfrak p^{w_{q_{\mathfrak p}}(n)}R_{\mathfrak p}$.
The reduction map to this local quotient is surjective because
$\mathfrak p$ is maximal. It follows that $I_nP_n$ is integer-valued
at every height-one localization, and hence is integer-valued on $R$.
Subtracting the leading coefficient times $P_n$ and descending in
degree proves the first direct-sum decomposition. Since each $P_n$
is monic with coefficients in $R$, the same subtraction preserves
$L[X]$ and gives the second decomposition.

It remains to verify projectivity for $L$ without assuming that it is
finitely generated. At a contributing prime, the overring
$L_{\mathfrak p}$ of the DVR $R_{\mathfrak p}$ is either that DVR
or $K$. Localization commutes with the intersection of two submodules.
At every other maximal ideal, $(I_n)_{\mathfrak m}=R_{\mathfrak m}$.
Testing equality at all maximal ideals therefore gives
\[
 I_n\cap L=
 \prod_{\substack{\operatorname{ht}\mathfrak p=1,\ |R/\mathfrak p|<\infty\\
                              L_{\mathfrak p}=K}}
                \mathfrak p^{-w_{q_{\mathfrak p}}(n)}.
\]
This is a finite product of invertible fractional ideals. Their direct
sum is projective, and the finite dual-basis theorem gives all tensor
isomorphisms and plethory compatibility.
\end{proof}

\begin{corollary}\label{cor:regular-basis-principal-characteristic}
In the preceding theorem, $\Omega(R,L)$ has a regular basis, with
one basis polynomial in each degree, if and only if every ideal
$I_n\cap L$ is principal. In particular this holds when these ideals
are principal, for example over a local base in the stated class.
\end{corollary}
\begin{proof}
A generator $c_n$ gives the basis element $c_nP_n$. Conversely the
leading coefficient of the degree-$n$ member of a regular basis
generates the degree-$n$ characteristic ideal. This is a statement
about regular bases; it does not assert that an arbitrary free-module
basis must respect the degree filtration.
\end{proof}


\section{Completion and finite-rank modules}
\begin{theorem}[Completion in every Noetherian dimension]\label{result-31}
Let $(R,\mathfrak m)$ be a Noetherian local domain with
fraction field $K$, let $A=\widehat R$, $Q=K\otimes_R A$,
and let $R\subset L\subset K$ be any overring.
Put $B=L\otimes_R A\subset Q$. For every $n$,
\begin{equation}
                    \Omega_n(R,L)\otimes_R A
                         \simeq\Omega_n(A,B).              \label{equation-69}
\end{equation}
The tensor maps and existence of the two polynomial
co-operations pass through this identification. There is
no dimension or finite-normalization hypothesis.
\end{theorem}
\begin{proof}
If the residue field is infinite, a grid of points with
distinct residues gives an invertible Vandermonde
matrix over both local rings, so the integer-valued
modules are the ordinary polynomial rings. Assume
the residue field is finite.

For $j\ge0$ define
\[
 C_j=\{x\in K:\mathfrak m^jx\subset R\},\qquad
 C'_j=\{x\in Q:(\mathfrak m A)^jx\subset A\}.
\]
Flatness gives $C_j\otimes_R A=C'_j$.
Indeed, choose finitely many generators $a_1,\ldots,a_s$
of $\mathfrak m^j$ and tensor the kernel description
\[
 C_j=\ker\bigl(K\longrightarrow(K/R)^s,
                       x\longmapsto(a_ix\bmod R)_i\bigr).
\]

Every coefficient of an integer-valued polynomial in
$K[X_1,\ldots,X_n]$ lies in some $C_j$, by the
punctured-spectrum Vandermonde argument of Theorem~\ref{result-27}.
We verify the corresponding assertion over $A$ even
though $A$ need not be a domain and $Q$ need not be
Artinian. Each coefficient $x\in Q$ has a denominator
$d\in R\setminus\{0\}$, which is a nonzerodivisor
of $A$ by flatness. If $\mathfrak p$ is a nonmaximal
prime of $A$, its contraction to $R$ is not
$\mathfrak m$, since otherwise it contains
$\mathfrak m A$, the maximal ideal of $A$.
The image of $R$ modulo this contraction is infinite.
Choose a rectangular Vandermonde grid of $R$-points
whose determinant $v$ is outside $\mathfrak p$.
The values of the polynomial are in $A$, so Cramer's
rule gives $vx\in A$. Consequently the ideal
\[
                         I_x=\{a\in A:ax\in A\}
\]
is not contained in any nonmaximal prime. If it is
proper, its radical is the maximal ideal, and finite
generation gives $(\mathfrak m A)^j\subset I_x$
for some $j$. Thus $x\in C'_j$.
Taking a common $j$ handles all coefficients of a
given polynomial.

Fix this $j$ and a finite rectangular degree bound
with $N$ allowed monomials. Let $E_j\subset R$
represent the finite ring $R/\mathfrak m^j$.
For coefficients in $C_j$, the value modulo $R$
depends only on the inputs modulo $\mathfrak m^j$:
a monomial difference has a factor in that ideal,
which sends each coefficient into $R$.
The desired bounded module is therefore
\[
 \ker\bigl(C_j^N\longrightarrow(C_j/R)^{E_j^n}\bigr).
\]
Since $R/\mathfrak m^j\simeq A/(\mathfrak m A)^j$,
the same finite tests apply over $A$. Tensoring the
kernel gives exactly the bounded module there.
Take the filtered union over $j$ and the degree
bounds to prove \eqref{equation-69} for $L=K$.
For arbitrary $L$, intersect each bounded coefficient module with
$L^N$ inside $K^N$. This intersection is the kernel of the difference
map from the direct sum of the two modules to $K^N$, so flat base
change preserves it. The resulting coefficient module is the same
intersection with $B^N$ inside $Q^N$. Taking the filtered union proves
the assertion. Multiplication in disjoint variables commutes with these
identifications. Every element after base change is a finite
$A$-linear combination of original polynomials, so existence of the
two expressions is preserved as well.

In dimension one, $Q$ is the total quotient ring of $A$. Choose
$0\ne t\in\mathfrak m$. For every $0\ne r\in R$, some power of $t$
belongs to $rR$, hence $Q=A[1/t]$. The latter is zero-dimensional
because $A/tA$ is Artinian and $\dim A=1$, and is therefore Artinian.
Every nonzerodivisor of $A$ remains a nonzerodivisor there and is thus
a unit. These arguments use exactness and faithful flatness of local
completion \cite[Lemmas 10.97.2--3, tags 00MB and 00MC]{StacksCompletionAnalytic}.

\end{proof}

\begin{theorem}[Finite-rank modules over a complete discrete valuation ring]\label{result-32}
Let $V$ be a complete discrete valuation ring with
uniformizer $t$ and fraction field $F$. Let $W$ be a
finite-dimensional $F$-vector space, and let $M\subset W$
be a $V$-submodule containing a full $V$-lattice.
There is an $F$-vector subspace
$D=\bigcap_{j\ge0}t^jM$ such that $M/D$ is a finite
free $V$-module. More precisely, after a choice of
$F$-basis, $M$ is a direct sum of finitely many
fractional ideals $t^aV$ and finitely many copies of $F$.
\end{theorem}
\begin{proof}
Choose a full lattice $\Lambda\subset M$ and identify
$(W,\Lambda)$ with $(F^d,V^d)$. Put $E=F/V$.
We first record the elementary duality facts used below.
The divisible module $E$ is injective over the principal
ideal domain $V$. For a direct proof of the extension
property, extend a homomorphism maximally by Zorn's
lemma. To adjoin a new element $x$, its annihilator
relative to the old submodule is an ideal $(a)$.
If $a\ne0$, divisibility solves $ay=f(ax)$; if the
ideal is zero choose any $y$. Sending $x$ to $y$
extends the map, a contradiction unless the domain was
already the whole module.
Also $E$ detects every nonzero torsion module: a nonzero
element has annihilator $(t^s)$, its cyclic submodule
maps nontrivially to $t^{-s}V/V$, and injectivity extends
this map to the whole module.

An endomorphism of $E$ is multiplication by a unique
element of $V$. Indeed its value on $t^{-s}+V$ is
$a_s t^{-s}+V$ for $a_s\in V/t^sV$, and compatibility
is exactly $a_{s+1}\equiv a_s\pmod {t^s}$. Completeness
identifies the compatible sequence with an element of
$V$. Consequently every homomorphism $E^d\to E$
is induced by a vector in $V^d$.

Let
\[
 M^\#=\{\phi\in W^*: \phi(M)\subset V\}
                                      \subset\Lambda^*.
\]
Then
\begin{equation}
 M=\{x\in W: \phi(x)\in V\text{ for every }
                                      \phi\in M^\#\}.       \label{equation-70}
\end{equation}
Only the reverse inclusion needs proof. If $x\notin M$,
the nonzero torsion element $x+M$ in $W/M$ is detected
by a homomorphism to $E$. Composing with
$W/\Lambda\simeq E^d$ and using the preceding
endomorphism computation gives an $F$-linear form
$\phi$ integral on $M$ but not on $x$.

The submodule $M^\#\subset\Lambda^*$ is finite free,
since $V$ is Noetherian. Smith reduction of its matrix
inside the finite free module $\Lambda^*$ puts it in
the form
$t^{a_1}Ve_1^*\oplus\cdots\oplus t^{a_s}Ve_s^*$
with zero remaining coordinates. Formula \eqref{equation-70}
then gives
\[
 M=t^{-a_1}Ve_1\oplus\cdots\oplus t^{-a_s}Ve_s
           \oplus Fe_{s+1}\oplus\cdots\oplus Fe_d.
\]
Its intersection of all $t^j$ multiples is exactly the
last vector-space summand, proving every assertion.
\end{proof}

\begin{theorem}[Epimorphic overrings of a complete one-dimensional order]\label{result-33}
Let $V$ be a complete discrete valuation ring with finite
residue field, let $A$ be a local commutative $V$-algebra
finite and torsion-free over $V$, and put
$Q=A\otimes_V\operatorname{Frac}(V)$.
For any overring $A\subset B\subset Q$, the multiplication
map $B\otimes_A B\to B$ is an isomorphism if and only if
\[
                            B=A+eQ
\]
for an idempotent $e\in Q$. Moreover injectivity of the
second polynomial tensor map for $\Omega(A,B)$ forces
this ring-epimorphism condition.
\end{theorem}
\begin{proof}
For the last assertion choose any $b\in B$. There is
a nonzerodivisor $d\in A$ with $db\in A$; a power
of the uniformizer of $V$ suffices. The finite ring
$A/dA$ has a monic null polynomial, for instance the
product of $X-a$ over a complete set of representatives.
Call it $P$. Then $P,bP\in\Omega(A,B)$. The tensor
\[
                       bP\otimes P-P\otimes bP
\]
has zero polynomial image, hence is zero if the tensor
map is injective. Taking the coefficient of degree
$\deg P$ in both factors gives
$b\otimes1=1\otimes b$ in $B\otimes_A B$.
This for every $b$ makes multiplication an isomorphism,
with inverse $b\mapsto b\otimes1$.

Now assume multiplication is an isomorphism.
Apply Theorem~\ref{result-32} to $B\subset Q$, using the full finite
free lattice $A$. The divisible part
$D=\bigcap_j t^jB$ is an ideal of $B$ and a vector
space over $\operatorname{Frac}(V)$. It is in fact
an ideal of $Q$, since $Q=A[1/t]$. The quotient
$C=B/D$ is finite over $V$, hence finite over $A$.
The composite $A\to C$ is a ring epimorphism because
$A\to B$ and $B\to C$ are. If $C\ne0$, its fiber
$C/\mathfrak m_A C$ is nonzero by Nakayama and is a
finite-dimensional algebra over $k=A/\mathfrak m_A$.
The multiplication isomorphism on this fiber implies
$(\dim_k C/\mathfrak m_A C)^2=\dim_k C/\mathfrak m_A C$,
so the dimension is one and the unit generates it.
Nakayama applied to $C/\operatorname{im}(A)$ shows
that $A\to C$ is surjective. The same conclusion is
automatic when $C=0$. Therefore $B=A+D$.

Put $I=A\cap D$. It is a finite $V$-module. From
$B/A\simeq D/I$ and the epimorphism condition,
\[
              B\otimes_A(D/I)=B\otimes_A(B/A)=0.
\]
As $D$ is already a $Q$-module,
$B\otimes_A D\simeq Q\otimes_QD=D$.
Right exactness now gives $D=BI=I+DI$.
The quotient $D/DI$ is a vector space over
$\operatorname{Frac}(V)$ generated as a $V$-module
by the image of the finite module $I$. Such a vector
space must be zero: multiplication by $t$ is
surjective, and Nakayama applies to a finite
$V$-module. Hence $D=DI\subset D^2\subset D$.
An idempotent ideal in the Artinian ring $Q$ is
$eQ$: in each local Artinian factor a proper ideal
is nilpotent, so an idempotent ideal is either zero
or the whole factor. This proves necessity.

Conversely, for $B=A+eQ$, write $f=1-e$. Then
$B=eQ\times fA$. The map $A\to fA$ is a quotient,
so $fA\otimes_A fA\simeq fA$. Localization gives
$eQ\otimes_A eQ\simeq eQ$, and the cross terms
$eQ\otimes_A fA$ and $fA\otimes_A eQ$ vanish,
because after inverting $t$ they involve the
orthogonal factors $eQ$ and $fQ$. Multiplication
is therefore an isomorphism on the whole product.
\end{proof}

\begin{theorem}[Epimorphic enlargements supported at the closed point]\label{result-34}
Let $(A,\mathfrak m)$ be a complete Noetherian local
ring, containing a complete discrete valuation ring
$V$ with uniformizer $t$, with $t\in\mathfrak m$
a nonzerodivisor on $A$. Assume the residue field
of $V$ maps isomorphically to $A/\mathfrak m$.
Let $W\subset A$ be a multiplicative set of
nonzerodivisors containing $t$, and put $Q=W^{-1}A$.
Suppose
\[
 A\subset B\subset Q,\qquad
 \text{every element of }B/A\text{ is killed by a power
 of }\mathfrak m.
\]
Then there is a finite-dimensional
$F=\operatorname{Frac}(V)$ vector space $D\subset B$
which is an ideal of $Q$, such that $B/D$ is a finite
$A$-module and $A\cap D$ is a finite $V$-module.
If $A\to B$ is a ring epimorphism, then
\begin{equation}
                         B=A+eQ,\qquad D=eQ                 \label{equation-71}
\end{equation}
for an idempotent $e\in Q$, and $eQ$ is finite-dimensional
over $F$. If the common residue field is finite,
injectivity of the polynomial tensor map for
$\Omega(A,B)$ implies this epimorphism condition.
\end{theorem}
\begin{proof}
Put $T=B/A$. Since every class is killed by a power
of $\mathfrak m$, it is killed by a power of $t$,
and $B\subset A[1/t]$.
The $t$-torsion subgroup $T[t]$ embeds in $A/tA$:
send the class of $a/t$ to $a\bmod tA$.
Its image is contained in the submodule of $A/tA$
killed elementwise by powers of $\mathfrak m$.
The ascending chain of annihilator submodules
$(0:_{A/tA}\mathfrak m^j)$ stabilizes, since
$A/tA$ is Noetherian. The stable module is killed
by one power of $\mathfrak m$, and hence has finite
length. In particular $T[t]$ is a finite-dimensional
vector space over $V/tV$.

Let $E=F/V$. By the injectivity of $E$ proved in
Theorem~\ref{result-32}, an embedding $T[t]\hookrightarrow E[t]^s$
extends to a map $T\to E^s$. This extension is
injective: any nonzero kernel element, being killed
by a power of $t$, has a nonzero multiple killed
by $t$, contradicting injectivity on $T[t]$.
The same double-annihilator argument as in Theorem~\ref{result-32}
classifies submodules of $E^s$. Explicitly, annihilate
$T$ by a submodule of
$\operatorname{Hom}_V(E^s,E)=V^s$; cogeneration by
$E$ shows that taking the annihilator again gives
$T$. Smith reduction of that finite submodule of
$V^s$ shows
\[
                            T\simeq E^r\oplus T_0,
\]
where $T_0$ has finite length over $V$.
Its maximal divisible submodule $T_{\mathrm{div}}$
is $E^r$ and is canonical, hence stable under $A$.

The ring $A$ is also $t$-adically complete and
separated. To see completeness, a $t$-adic Cauchy
sequence is $\mathfrak m$-adically Cauchy; its limit
has the required congruences because each ideal
$t^jA$ is $\mathfrak m$-adically closed.
Closedness and separatedness are the Krull
intersection theorem applied to finite modules.
There is consequently a natural isomorphism
\begin{equation}
 \operatorname{Hom}_V(E,A[1/t]/A)\simeq A.                  \label{equation-72}
\end{equation}
Indeed the image of $t^{-j}+V$ is
$a_j/t^j+A$ with $a_j\in A/t^jA$; the homomorphism
condition is the compatibility of these residues,
and completeness gives a unique $a\in A$.

Choose the $r$ coordinate injections of
$T_{\mathrm{div}}\simeq E^r$ into $A[1/t]/A$.
By \eqref{equation-72} they are multiplication by elements
$a_1,\ldots,a_r\in A$. Put
$D=Fa_1+\cdots+Fa_r$.
These vectors are $F$-linearly independent:
a relation, after scaling to coefficients in
$V$ with at least one unit, would give a nonzero
kernel on $E^r$.
Their fractional multiples lie in $B$ because
their classes lie in $T$. Moreover
\[
              A\cap D=Va_1+\cdots+Va_r,\qquad
              D/(A\cap D)=T_{\mathrm{div}}.
\]
The first equality follows from injectivity of
the coordinate map $E^r\to T$.
Since $B/(A+D)\simeq T_0$ is finite over $V$,
$B/D$ is a finite $A$-module.

Krull intersection on this finite module shows
$\bigcap_j t^j B=D$. Thus $D$ is an ideal of
$B$, and in particular stable under $A$.
Every $w\in W$ acts injectively on the finite
$F$-vector space $D$, because it is a unit of $Q$.
It therefore acts bijectively. Hence $D$ is stable
under $W^{-1}A=Q$, so it is an ideal of $Q$.

Suppose now $A\to B$ is a ring epimorphism.
The finite $A$-algebra $B/D$ is an epimorphic image
of $A$. If it is zero the map from $A$ is surjective;
otherwise the residue-field dimension argument
and Nakayama proof of Theorem~\ref{result-33} show that $A\to B/D$
is surjective. Thus $B=A+D$. With $I=A\cap D$,
the same tensor computation gives
\[
 B\otimes_A(B/A)=0,\qquad B\otimes_A D=D,\qquad
                         D=BI=I+DI.
\]
Here $B\otimes_A D=D$ follows by localizing at
$W$ on the already $Q$-module $D$.
The quotient $D/DI$ is an $F$-vector space generated
over $V$ by the finite module $I$, so it vanishes.
Therefore $D=D^2$.
Since $D$ is finite-dimensional over $F\subset Q$,
it is finitely generated as an ideal of $Q$.
A finitely generated idempotent ideal is generated
by an idempotent: write generators $d_i$ as
$d_i=\sum a_{ij}d_j$ with $a_{ij}\in D$;
the adjugate identity says $\det(1-(a_{ij}))D=0$.
Write this determinant as $1-e$, with $e\in D$.
Then $(1-e)e=0$, so $e^2=e$, and $D=eQ$.
This proves \eqref{equation-71}.

Finally suppose the residue field is finite and
the polynomial tensor map is injective. For any
$b\in B$ choose $j$ with $\mathfrak m^j b\subset A$.
A monic null polynomial of the finite ring
$A/\mathfrak m^j$ gives $P,bP\in\Omega(A,B)$.
Top-coefficient extraction from
$bP\otimes P-P\otimes bP$ gives
$b\otimes1=1\otimes b$ exactly as in Theorem~\ref{result-33}.
Thus $B\otimes_A B\simeq B$.
\end{proof}

\subsection{Polynomial lattices and orthogonal components}
The next two lemmas strengthen the negative direction: they detect a
kernel already in the binary tensor map, without assuming that either
polynomial cooperation exists.

\begin{lemma}[Polynomial lattices have infinitely many residue classes]\label{result-45}
Let $V$ be a discrete valuation ring with uniformizer $t$ and fraction
field $F$. Let $C$ be a nonzero local $V$-algebra, finite and
torsion-free over $V$, with maximal ideal $\mathfrak n$ and residue
field $\kappa$. If a $C$-submodule
\[
 C[X]\subset M\subset
 \{f\in C[1/t][X]:f(C)\subset C\}
\]
contains all ordinary polynomials, then $M/\mathfrak nM$ is
infinite-dimensional over $\kappa$.
\end{lemma}
\begin{proof}
For each $d$, the submodule $M_{\le d}$ is finite over $V$.
Evaluate its polynomials at $d+1$ distinct elements of the infinite
domain $V$. The Vandermonde matrix is invertible over $F$, while
all evaluations lie in the finite $V$-module $C^{d+1}$.
Its inverse bounds the coefficients in a fixed finite $V$-lattice
inside $C[1/t]^{d+1}$. Every submodule of that lattice is finite.

Suppose $M/\mathfrak nM$ were finite-dimensional. Lift a finite
spanning set and let $G$ be its $C$-span, so $M=G+\mathfrak nM$.
The local ring $C/tC$ is Artinian, hence
$\mathfrak n^s\subset tC$ for some $s$. Iterating the equality
gives $M=G+tM$. All elements of $G$ have degree at most some $d$.
Starting with $f_0=X^{d+1}$, choose recursively
$f_i=g_i+tf_{i+1}$ with $g_i\in G$ and $f_{i+1}\in M$.
Every $f_i$ has degree $d+1$ and leading coefficient $t^{-i}$.
These coefficients cannot lie in the fixed finite lattice bounding
$M_{\le d+1}$. This contradiction proves the assertion.
No finite-generation assumption on $M$ was used.
\end{proof}

\begin{lemma}[Orthogonal generic components force a binary kernel]\label{result-46}
Let $A,V,t,F,W,Q$ satisfy the coefficient-ring and localization
hypotheses of Theorem~\ref{result-34}, with finite common residue
field $k$. For $M=\Omega(A,B)$, the map
$M\otimes_A M\to B[X,Y]$ is not injective in either of the following
cases.
\begin{enumerate}[label=(\alph*)]
\item $B=A+eQ$, where $e\ne0,1$ is idempotent,
$I=A\cap eQ$ is finite over $V$, and $I[1/t]=eQ$.
In particular this applies to every proper nonzero summand
furnished by Theorem~\ref{result-34}.
\item $A$ is finite over $V$, $B=Q=A[1/t]$, and $Q$ has a
nontrivial idempotent.
\end{enumerate}
\end{lemma}
\begin{proof}
First consider (a). Put
\[
 C=eA,\qquad C'=(1-e)A,\qquad J=A\cap(1-e)Q.
\]
The action of $C$ on $I$ is faithful: an element annihilating
$I$ annihilates $I[1/t]=eQ$, hence is zero. Thus
$C\subset\operatorname{End}_V(I)$ is finite over $V$.
Both $C$ and $C'$ are nonzero local quotients of $A$, with
residue field $k$, and $t$ is regular on them.
There is an exact pullback description
\[
 A=C\times_T C',\qquad T=A/(I+J)\simeq C/I\simeq C'/J.
\]
Indeed two chosen lifts of matching residues differ by an element
of $I+J$, which permits adjustment to a common lift in $A$.
Since $C[1/t]=I[1/t]=eQ$, the finite $V$-module $C/I$ is
$t$-primary torsion. Therefore $T$ is a finite ring. It is
nonzero, because $I$ and $J$ are proper ideals of the local ring
$A$. Also $B=eQ\times C'$.

Let $M_e$ be the first projection of $M$. It is a $C$-submodule
of $\Omega(C,eQ)$ containing $C[X]$. Put
\[
 N_e=\{f\in eQ[X]:f(C)\subset I\}.
\]
These are exactly the polynomials whose supported pairs $(f,0)$
belong to $M$. Every $f\in M_e$ induces a function $T\to T$.
To verify this, lift $f$ to a pair $(f,g)\in M$ and, for two
elements $r,r'\in C$ with the same $T$-class, choose a common
$s\in C'$ with that class. Both $(r,s)$ and $(r',s)$ belong
to $A$, so evaluation shows $f(r)\equiv g(s)\equiv f(r')$
in $T$. The kernel of this function map is exactly $N_e$.
Hence $M_e/N_e$ is finite over $V$.

If every supported pair $(f,0)$ belonged to $\mathfrak mM$,
then $N_e\subset\mathfrak m_C M_e$. The preceding finite
quotient would make $M_e/\mathfrak m_C M_e$ finite-dimensional,
contrary to Lemma~\ref{result-45}. Choose therefore $(f,0)$
with nonzero class in $M/\mathfrak mM$.

A monic null polynomial $P\in C'[X]$ for the finite ring
$T=C'/J$ gives $(0,P)\in M$. Its class modulo $\mathfrak mM$
is nonzero: the second projection of $\mathfrak mM$ has every
coefficient in $\mathfrak m_{C'}$, while $P$ is monic.
The tensor
\[
                    (f,0)\otimes(0,P)
\]
is nonzero, as its image in
$(M/\mathfrak mM)\otimes_k(M/\mathfrak mM)$ is the pure
tensor of two nonzero vectors. Its polynomial image is zero,
because its two factors have orthogonal component supports.
This proves (a).

For (b), choose a nontrivial idempotent $e$. Both $eA$ and
$(1-e)A$ are finite over $V$, and the two support ideals are
full lattices in their respective components: clear powers
of $t$ in $Q=A[1/t]$. The same finite pullback ring $T$ is
available. Both polynomial components may now have denominators,
but the common-input argument above still defines the function
$T\to T$ in each projection. Applying Lemma~\ref{result-45}
to both projections gives supported elements $(f,0)$ and
$(0,g)$ with nonzero classes modulo $\mathfrak mM$.
Their pure tensor is nonzero by the same residue-field
detector and has zero polynomial image. This proves (b).
\end{proof}

\Needspace{9\baselineskip}
\section{A tensor obstruction from generic nilpotents}
\begin{theorem}[A generic nilpotent forces a tensor kernel after completion]\label{result-35}
Let $(A,\mathfrak m)$ be a one-dimensional complete
Noetherian local ring with finite residue field. Assume that
$A$ is Cohen--Macaulay and that either $A$ has characteristic
$p>0$, or $A$ has characteristic zero and $p$ is a
nonzerodivisor, where $p$ is its residue characteristic.
If $A$ is not reduced, then
\[
 \operatorname{Int}(A)\otimes_A\operatorname{Int}(A)
       \longrightarrow\operatorname{Int}(A^2)
\]
is not injective. Here coefficients belong to the total
quotient ring of $A$. The proof constructs a nonzero tensor
and an explicit balanced linear functional detecting it.
\end{theorem}
\begin{proof}
We first construct the finite lattice in which the argument
will take place, including the needed compatible coefficient
field. There is a complete discrete valuation subring
$V_0\subset A$ such that $A$ is finite and torsion-free over
$V_0$, and such that some nonreduced local factor $Q_0$ of
$Q(A)$ has residue field $F$ separable over
$K_0=\operatorname{Frac}(V_0)$.

In mixed characteristic take the Cohen coefficient ring.
It is a complete discrete valuation ring with uniformizer
$p$, and $p$ is a parameter. The finite-module assertion
follows by lifting finitely many generators modulo $p$ and
successively approximating in the $p$-adic topology;
$A$ is $p$-adically complete because $pA$ is
$\mathfrak m$-primary. Separability holds in characteristic
zero.

Here is the extra parameter choice needed in characteristic
$p$. The roots of $X^{|A/\mathfrak m|}-X$ give a coefficient
field $k\subset A$ (or use the Cohen structure theorem).
Choose a minimal prime $\mathfrak p$ corresponding to a
nonreduced generic factor. Such a factor exists because a
parameter is a nonzerodivisor, so a nonzero nilpotent remains
nonzero in $Q(A)$. The complete one-dimensional domain
$C=A/\mathfrak p$ has finite normalization, a complete
discrete valuation ring $V=l[[u]]$ with $l/k$ finite.
Here finiteness follows from the Nagata property of complete
Noetherian local rings; a finite domain over a complete local
ring is local, because a finite complete algebra decomposes
as a product of complete local algebras. A nonzero conductor
then gives $u^N\in C$ for all sufficiently large $N$.
Choose such an $N$ with $p\nmid N$ and lift $u^N$ to
$t_0\in\mathfrak m$.

We can adjust this lift to a parameter without changing its
image in $C$. Choose
$a\in\mathfrak p$ outside every other minimal prime, by
prime avoidance. At each other minimal prime, the nonzero
images of $a,a^2,a^3,\ldots$ are distinct: $a$ lies in
the maximal ideal of the corresponding local domain.
Thus $t_0+a^j$ belongs to each such prime for at most one
$j$. Choosing $j$ outside this finite exceptional set gives
a parameter $t$ mapping to $u^N$. If there is no other
minimal prime, use $t=t_0$.
Set $V_0=k[[t]]\subset A$. This map is injective, since
a nonzero power series is $t^r$ times a unit and $t$ is a
nonzerodivisor. Successive approximation modulo $t^m$ shows
that $A$ is finite over $V_0$. In the chosen factor the
residue field is $F=l((u))$, a finite separable extension
of $k((u^N))$ since $p\nmid N$.

Since $F/K_0$ is finite separable, it embeds in $Q_0$ as
a coefficient field compatibly with $K_0$. To see this
directly, choose a primitive element of $F/K_0$ and lift
it to $Q_0$. Its separable minimal polynomial has invertible
derivative. Newton correction modulo successive powers of
the nilradical, which is nilpotent, gives an actual root
lifting the primitive element. Its $K_0$-algebra is the
required copy of $F$. Write
\[
 \rho:Q_0\longrightarrow F,\qquad N=\ker\rho,
 \qquad N^h=0\quad(h\ge2).
\]
Let $V\subset F$ be the integral closure of $V_0$ in $F$.
It is a finite complete discrete valuation ring, with
uniformizer $u$ and finite residue field. The same notation
$V=l[[u]]$ applies in the equal-characteristic construction.
The ring
\[
 O=V[\operatorname{pr}_0(A)]\subset Q_0
\]
is finite and torsion-free over $V$, contains $V$, satisfies
$\rho(O)=V$, and spans $Q_0$ over $F$.
Finiteness follows because $\operatorname{pr}_0(A)$ is
finite over $V_0$ and hence integral over $V$.
Let $e$ denote the idempotent of the selected factor in
$Q(A)$. Since $eO$ is a finite $V_0$-module in $Q(A)$,
clearing finitely many denominators gives an integer $c\ge0$
such that
\begin{equation}
                         u^c eO\subset A.                 \label{equation-73}
\end{equation}
Indeed a sufficiently high power of the uniformizer of
$V_0$ works, and in $V$ that power is a unit times a
power of $u$; the unit preserves $O$.

Put $J=N\cap O$. There is a primitive element
$z_0\in J\cap\operatorname{Ann}_{Q_0}(N)$: start with a
nonzero element in the last nonzero power of $N$, multiply
it into $O$, and divide by the greatest possible power of
$u$ in the finite free $V$-module $O$. Choose a
$V$-linear map $\ell:O\to V$ with $\ell(z_0)=1$ and
extend it $F$-linearly to $Q_0$. Such a map exists because
a primitive vector in a finite free module over a discrete
valuation ring extends to a basis.

There is an integer $b\ge0$, independent of
$G\in\operatorname{Int}(A)$, such that, with
$x_G=\operatorname{pr}_0(G'(0))$,
\begin{equation}
                  Jx_G\subset u^{-b}O,\qquad
                  \rho(x_G)\in u^{-b}V.                   \label{equation-74}
\end{equation}
Here is a finite interpolation proof. Choose finitely many
$V$-generators $n_1,\ldots,n_s$ of $J$. By \eqref{equation-73},
$u^c e n_j\in A$. For any $h$ fixed distinct scalars
$s_1,\ldots,s_h\in V_0$, the values
$G(s_i u^c e n_j)$ lie in $A$ and their selected components
lie in $O$. As a polynomial in the scalar $s$, this
evaluation has degree at most $h-1$, because $n_j^h=0$.
Inverting this one fixed Vandermonde matrix bounds its
linear coefficient $u^c n_jx_G$ in a fixed fractional
lattice $u^{-d}O$. Increasing $b$ gives the first inclusion
for every generator, hence all of $J$. Since $z_0N=0$,
$z_0x_G=z_0\rho(x_G)$. Apply $\ell$ to obtain the second
inclusion from the first.

Write $\lambda(G)=\rho(x_G)$. Define
\begin{equation}
 \beta(G,H)=\ell(x_G)\lambda(H)\pmod {u^{-2b}V}
                         \quad\hbox{in }F/u^{-2b}V.
                                                               \label{equation-75}
\end{equation}
This additive pairing is $A$-balanced. If $a\in A$ and
$a_0=\operatorname{pr}_0(a)$, the difference between
$\beta(aG,H)$ and $\beta(G,aH)$ equals
\[
 \ell\bigl((a_0-\rho(a_0))x_G\bigr)\lambda(H).
\]
The first factor lies in $u^{-b}V$ by \eqref{equation-74}, since
$a_0-\rho(a_0)\in J$, and the second does also.
Their product therefore vanishes in the indicated quotient.
The quotient is an $A$-module through $\rho\operatorname{pr}_0$,
whose image is in $V$.

Choose an integer $r>2c+2b$. We construct a polynomial
$g_r\in F[X]$ with
\begin{equation}
 g_r(0)=0,\quad g_r'(0)=(-u^{-1})^r,\quad
 g_r(V)\subset V,\quad g_r(O)\subset V+u^{-r}J.             \label{equation-76}
\end{equation}
In characteristic $p$, choose a power $D$ of the residue
cardinality with $D\ge h$ and set
$\delta(X)=(X^D-X)/u$, $g_r=\delta^{\circ r}$.
The Frobenius identity and $N^h=0$ give
$\delta(a+n)=\delta(a)-n/u$ for $a\in V,n\in N$;
all assertions follow.

In mixed characteristic choose a power $D$ of the residue
cardinality so large that $D\ge h$ and
\[
 v(D)\ge (r-1)(h-2)+\max_{1\le j<h}v(j).
\]
Again use $g_r=\delta^{\circ r}$. The elementary identity
$\binom Dj=(D/j)\binom{D-1}{j-1}$ gives
$v(\binom Dj)\ge v(D)-v(j)$.
For $0\le i<r$, $a\in V$ and $n\in u^{-i}J$, expand
\[
 \delta(a+n)=\delta(a)+\frac{Da^{D-1}-1}{u}n
          +\sum_{2\le j<h}\frac{\binom Dj a^{D-j}}u n^j.
\]
The linear term belongs to $u^{-(i+1)}J$, and so does
each higher term: its valuation bound is
$v(\binom Dj)-ij-1\ge-i-1$.
Thus $\delta(V+u^{-i}J)\subset V+u^{-(i+1)}J$.
Induction proves \eqref{equation-76}; its derivative follows from
$\delta(0)=0$ and $\delta'(0)=-u^{-1}$.

Extend polynomials from $Q_0$ to $Q(A)$ by zero on the
other factors, and put
\[
       F_r=u^c e z_0g_r,\qquad H_r=u^{r+c}eg_r.
\]
Both are in $\operatorname{Int}(A)$. For $H_r$ this is
\eqref{equation-73} and \eqref{equation-76}. For $F_r$, multiplication by $z_0$
kills the nilpotent part of $g_r(a)$, leaving
$u^cz_0g_r(\rho(a))\in u^cO$.
The polynomial image of
\[
                   F_r\otimes H_r-H_r\otimes F_r
\]
is zero, because $F_r=u^{-r}z_0H_r$ in $Q(A)[X]$.
Its value under \eqref{equation-75} is $u^{2c-r}$: the second term
vanishes since $\rho(z_0)=0$, and the two signs cancel
in the first. This is nonzero modulo $u^{-2b}V$ by the
choice of $r$. The required tensor kernel is proved.

The structural inputs are the Cohen structure theorem,
complete finite-algebra decomposition, and the Nagata property of
complete local rings \cite[Theorem 10.160.8, Lemmas 10.160.2 and
10.162.8; tags 032A, 0325, 032W]{StacksCompletionAnalytic}. The lattice and obstruction arguments above are
given in full and do not invoke flatness of integer-valued
polynomial rings.
\end{proof}

\Needspace{12\baselineskip}
\section{The complete local and global criteria}
\begin{theorem}[Complete one-dimensional fraction-field classification without finite normalization]\label{result-36}
Let $R$ be an infinite one-dimensional Noetherian domain,
and put $K=\operatorname{Frac}(R)$.
The compatible coaddition and comultiplication tensors for
$\operatorname{Int}(R)$ are uniquely determined if and only
if, at each maximal ideal with finite residue field,
the completion $\widehat{R_{\mathfrak m}}$ is a domain.
Equivalently, at each such ideal the normalization of
$R_{\mathfrak m}$ is finite and local. In every positive
case all multivariable tensor maps are isomorphisms and
the compatible plethory is unique.
\end{theorem}
\begin{proof}
By Theorem~\ref{result-26} the tensor question localizes, so fix a
finite-residue maximal ideal and write its local ring
as $R$. If its completion is nonreduced, Theorem~\ref{result-35} applies:
the completion is one-dimensional Cohen--Macaulay, and
in mixed characteristic the nonzero element $p\in R$
remains a nonzerodivisor by flatness. Theorem~\ref{result-31} then proves
failure of tensor injectivity over $R$.

If the completion is reduced, $R$ is analytically
unramified and its normalization $T$ is finite
\cite[Lemma 10.162.10(4), tag 032Y]{StacksCompletionAnalytic}. Theorem~\ref{result-30} proves that the tensors work exactly
when $T$ is local. This is equivalent to the completion
of $R$ being a domain. In one direction, when $T$ is
local it is a discrete valuation ring, and exact
completion embeds $\widehat R$ in the complete discrete
valuation ring $\widehat T$. Conversely, if
$\widehat R$ is a domain, flatness embeds
$T\otimes_R\widehat R$ into
$K\otimes_R\widehat R$, which is a field. But the
completion of the finite semilocal algebra $T$ is the
product of its nonzero completions at its maximal
ideals. A subring of a field cannot have the nontrivial
idempotents of a product with more than one factor.
Thus $T$ is local.

At maximal ideals of infinite residue field the ordinary
Vandermonde argument gives the polynomial ring, as
already proved. In the positive finite-residue case,
Theorem~\ref{result-29} gives every multivariable tensor isomorphism.
Testing kernels and cokernels at every maximal ideal
proves the global assertion, and Theorem~\ref{result-21} supplies the
unique plethory compatibility.
\end{proof}

\begin{theorem}[All one-dimensional Noetherian domains and every torsion-free coefficient extension]\label{result-37}
Let $R$ be an infinite one-dimensional Noetherian domain,
let $K=\operatorname{Frac}(R)$, and let $R\subset S$
be an $R$-torsion-free unital extension. Put
$L=S_{\mathrm{rat}}\subset K$ as in Theorem~\ref{result-23}.
The binary map
$\Omega(R,S)\otimes_R\Omega(R,S)\to\Omega_2(R,S)$ is
injective if and only if at each maximal ideal $\mathfrak m$
with finite residue field either
\begin{equation}
 \begin{gathered}
 L_{\mathfrak m}=R_{\mathfrak m},\qquad\hbox{or}\\
 L_{\mathfrak m}=K\quad\hbox{and}\quad
 \widehat{R_{\mathfrak m}}\hbox{ is a domain}.
 \end{gathered}                                            \label{equation-77}
\end{equation}
Equivalently, in the second alternative the local
normalization is finite and local. These conditions are also
equivalent to unique polynomial coaddition and comultiplication.
In every positive case all tensor maps are isomorphisms and the
compatible plethory is unique. There is no
finite-normalization assumption on $R$.
\end{theorem}
\begin{proof}
Theorems~\ref{result-23} and \ref{result-26} reduce to the rational
coefficient overring and localize the tensor map. Infinite-residue
local rings give ordinary polynomials. Fix a finite-residue local
ring $R$, and put
\[
 A=\widehat R,\qquad Q=K\otimes_R A=Q(A),\qquad
 B=L\otimes_R A\subset Q.
\]
Theorem~\ref{result-31} identifies the completed polynomial modules
and their tensor maps. The ring $A$ is one-dimensional
Cohen--Macaulay and is finite and torsion-free over a complete
coefficient discrete valuation ring with the same finite residue
field, as in the proof of Theorem~\ref{result-35}.

Assume only injectivity of the binary map. Theorem~\ref{result-33}
gives $B=A+eQ$. A proper nonzero $e$ is impossible by
Lemma~\ref{result-46}(a): $A\cap eQ$ is a finite lattice and
$(A\cap eQ)[1/t]=eQ$. If $e=0$, then $B=A$, and faithful
flatness gives $L=R$.

It remains to consider $e=1$, so $B=Q$. If $Q$ has more than one
local Artinian factor, Lemma~\ref{result-46}(b) contradicts
injectivity. Thus $Q$ is local. If $A$ is nonreduced,
Theorem~\ref{result-35} supplies a nonzero kernel. Its characteristic
hypotheses hold: in mixed characteristic the nonzero element
$p\in R$ remains regular after flat completion. Hence $A$ is
reduced. The local Artinian ring $Q$ is then a field, and the
inclusion $A\hookrightarrow Q$ proves that $A$ is a domain.
Faithful flatness applied to $K/L$ gives $L=K$.
This proves necessity of \eqref{equation-77} from injectivity alone.

Conversely $L=R$ gives ordinary polynomials. If $L=K$ and the
completion is a domain, Theorem~\ref{result-36} gives every tensor
isomorphism and identifies the hypothesis with finite local
normalization having one branch. Testing kernels and cokernels
at maximal ideals proves the global positive assertion.
All tensor isomorphisms give the two unique polynomial operations;
conversely uniqueness for the zero polynomial forces injectivity.
The compatible plethory is provided by Theorem~\ref{result-21}.
\end{proof}

\begin{theorem}[The higher-dimensional local classification without finite normalization]\label{result-38}
Let $(R,\mathfrak m)$ be a Noetherian local domain
of dimension at least two with finite residue field,
let $K=\operatorname{Frac}(R)$ and $R\subset L\subset K$.
Put
\[
 R^\sharp=\{x\in K:\mathfrak m^jx\subset R
                       \text{ for some }j\},\qquad
 E=L\cap R^\sharp.
\]
The binary tensor map for $\Omega(R,L)$ is injective if and
only if $E=R$, equivalently $\Omega(R,L)=R[X]$. These conditions
are also equivalent to unique polynomial coaddition and
comultiplication. No assumption on normalization is needed.
\end{theorem}
\begin{proof}
Theorem~\ref{result-27} gives
$\Omega(R,L)=\Omega(R,E)$ and characterizes $E=R$ as the
ordinary polynomial case, where all conclusions hold.
For necessity assume only binary injectivity and put
\[
 A=\widehat R,\qquad Q=K\otimes_R A,
 \qquad B=E\otimes_R A\subset Q.
\]
Theorem~\ref{result-31} preserves the polynomial modules and
injectivity. Every element of $B/A$ is killed by a power of
$\mathfrak m_A$.

Choose the coefficient DVR required by Theorem~\ref{result-34}.
In characteristic $p$, use the coefficient field $k\subset A$
and $V=k[[t]]$ for $0\ne t\in\mathfrak m_R$. In characteristic
zero, use the Cohen coefficient ring with uniformizer $p$.
Flatness of completion makes the chosen uniformizer regular,
and it belongs to $W=R\setminus\{0\}$ defining $Q=W^{-1}A$.
The residue field agrees with that of $A$. No assertion that
$A$ is finite over $V$ is needed.

Theorem~\ref{result-34} gives $B=A+eQ$, where
$I=A\cap eQ$ is finite over $V$. Since $eQ\subset A[1/t]$,
clearing powers of $t$ gives $I[1/t]=eQ$.
A proper nonzero $e$ contradicts Lemma~\ref{result-46}(a).
If $e=1$, then $I=A$ is finite over the DVR $V$, so
$\dim A=1$, contradicting $\dim A=\dim R\ge2$.
Thus $e=0$ and $B=A$. Faithful flatness applied to $E/R$
gives $E=R$. Finally unique polynomial coaddition and
comultiplication imply injectivity by the zero polynomial,
while the ordinary polynomial ring supplies their existence
and uniqueness. This proves all the stated equivalences.
\end{proof}

\Needspace{19\baselineskip}
\begin{proof}[Proof of Theorem~\ref{result-39}]
The implication (ii)$\Rightarrow$(iii) follows because the two
polynomial substitutions preserve integer-valuedness. The implication
(iii)$\Rightarrow$(i) follows from uniqueness for the zero polynomial.

Assume (i). The rational-coefficient reduction of
Theorem~\ref{result-23} replaces $S$ by $L$, and
Theorem~\ref{result-26} carries injectivity to every maximal
localization. Theorems~\ref{result-37} and \ref{result-38} then
prove exactly the finite-residue conditions in (iv), from
injectivity alone. Infinite-residue local rings give the ordinary
polynomial ring.

Conversely, under (iv), the positive cases in those same theorems
give all tensor isomorphisms locally. Testing kernels and cokernels
at every maximal ideal gives (ii). The polynomial identities and
composition argument of Theorem~\ref{result-21} gives the unique
compatible plethory. Finally Theorem~\ref{result-36} identifies the
dimension-one domain-completion condition with finite local
normalization having one branch, and Theorem~\ref{result-27}
identifies $R_{\mathfrak m}^{\sharp}=R_{\mathfrak m}$ with depth
at least two. These arguments are local and require no uniform
bound on the dimensions of maximal localizations.
\end{proof}

\subsection{Finite normalization and conductor specializations}
The completed criterion retains the earlier finite-normalization
statements, including a distinction between failure of uniqueness and
failure of the existence of multiplicative tensors.

\begin{corollary}[The finite-normalization criteria in full]\label{result-50}
Let $R$ be a nonfield Noetherian domain, let $K$ be its fraction field,
and suppose its normalization $T$ is finite over $R$. For an arbitrary
unital $R$-torsion-free extension $R\subset S$, put $L=S_{\rm rat}$.
At a maximal ideal $\mathfrak m$ write $R_m=R_{\mathfrak m}$,
$L_m=L_{\mathfrak m}$ and $T_m=T_{R\setminus\mathfrak m}$.
The following are equivalent: binary tensor injectivity; all tensor
maps being isomorphisms; unique coaddition and comultiplication; and
these conditions at every maximal ideal with finite residue field:
\begin{enumerate}[label=(\roman*)]
\item if $\dim R_m=1$, then $L_m=R_m$, or $L_m=K$ and $T_m$ is local;
\item if $\dim R_m\ge2$, then
\[
 L_m\cap R_m^\sharp=R_m,\qquad
 R_m^\sharp=\{x\in K:(\mathfrak mR_m)^j x\subset R_m
                              \text{ for some }j\ge0\}.
\]
\end{enumerate}
There is no condition at an infinite-residue maximal ideal. Every
positive case has the unique compatible plethory. For $S=K$, condition
(ii) is exactly $\operatorname{depth}R_m\ge2$.

If $R$ is one-dimensional and $S=K$, these conditions are further
equivalent to the existence, for every $F\in\operatorname{Int}(R)$,
of a finite expression
\[
                    F(XY)=\sum_i G_i(X)H_i(Y),
       \qquad G_i,H_i\in\operatorname{Int}(R).
\]
Thus in this fraction-field case mere multiplicative existence is
equivalent to one normalization branch above each finite-residue maximal
ideal. If that branch condition holds, the criterion for arbitrary
$R$-torsion-free $S$ reduces to $L_m=R_m$ or $L_m=K$ at each such
ideal; the choice may differ from ideal to ideal.
\end{corollary}
\begin{proof}
Theorem~\ref{result-39} gives the four equivalent conditions without
finite normalization. Its one-dimensional domain-completion condition
is equivalent, by Theorem~\ref{result-36}, to finite local normalization.
The finite normalization of $R_m$ is exactly $T_m$, so this gives (i).
Theorem~\ref{result-27} gives the displayed higher-dimensional criterion
and its fraction-field depth formulation. This proves every assertion
except the additional mere-existence equivalence in dimension one.

If the branch condition holds, every tensor map is an isomorphism,
so the displayed multiplicative expressions exist. Conversely assume
the expressions exist globally, and localize at a maximal ideal
$\mathfrak m$. Theorem~\ref{result-26} writes every local integer-valued
polynomial as $G/s$ for a global one and $s\notin\mathfrak m$.
An expression for $G(XY)$, divided by $s$ in one factor, is therefore
an expression for the given local polynomial. The fourth condition of
Theorem~\ref{result-30} then forces $T_m$ to be local whenever the
residue field is finite. That theorem proves this stronger existence
obstruction by finite iteration of binary expressions; no inference
from noninjectivity alone is being made. The final endpoint description
is now (i), with its branch condition already satisfied.
\end{proof}

\begin{corollary}[Single-branch conductor rings and their coefficients]\label{result-51}
Let $V$ be a DVR with fraction field $K$, uniformizer $t$ and finite
residue field, and let $R\subset V$ be a unital subring containing
$t^mV$ for some $m\ge1$. For every unital $R$-torsion-free extension
$R\subset S$, with $L=S_{\rm rat}$, unique polynomial operations are
equivalent to
\[
                         L=R\quad\hbox{or}\quad L=K.
\]
Both cases have all tensor isomorphisms and the unique compatible
plethory. Every proper intermediate overring $R\subsetneq L\subsetneq K$
has a noninjective binary tensor map.

In particular, if $k\subset l$ are any finite fields and
$R=k+t^m l[[t]]$, then $R$ is a one-dimensional Noetherian local domain
with fraction field $l((t))$, and $\operatorname{Int}(R)$ is free
over $R$. All its tensor maps are isomorphisms and its compatible
plethory is unique. This includes, for $m\ge1$, both the case $k=l$
and finite residue extensions. The rings
$\mathbb F_2[[t^2,t^3]]$ and $\mathbb F_2+t\mathbb F_4[[t]]$
are therefore free, and hence flat, examples.
\end{corollary}
\begin{proof}
Theorem~\ref{result-29} proves that $R$ is a one-dimensional Noetherian
local domain with finite local normalization $V$, and proves freeness
of $\operatorname{Int}(R)$. Theorem~\ref{result-37} now gives precisely
the two coefficient endpoints, including the failure of binary
injectivity at every strict intermediate overring. The all-arity and
plethory conclusions are part of the positive direction there.

For the displayed family take $V=l[[t]]$. The finite-conductor
hypothesis is literal. Its fraction field is $l((t))$: every element
$v\in V$ is the quotient of $t^m v\in R$ by $t^m\in R$.
Theorem~\ref{result-29} supplies freeness. Finally
$\mathbb F_2[[t^2,t^3]]=\mathbb F_2+t^2\mathbb F_2[[t]]$,
since every integer at least two is generated by two and three.
The second named ring uses $k=\mathbb F_2$, $l=\mathbb F_4$, $m=1$.
Both examples have one normalization branch.
\end{proof}

\subsection{Existence failures with nonintegral coefficients}
The tensor criteria above also admit a stronger conclusion in the
finite-normalization cases. Failure of injectivity alone would not
give the following failure of existence.

\begin{theorem}[Nonintegral coefficients obstruct multiplicative existence]\label{result-55}
Let $(R,\mathfrak m)$ be a Noetherian local domain with finite residue
field, fraction field $K$, and finite normalization $T\subset K$.
\begin{enumerate}
\item If $\dim R=1$ and $R\subset S$ is any unital $R$-torsion-free
extension, put $L=S_{\rm rat}$. Unique polynomial coaddition and
comultiplication exist exactly when
$L=R$, or when $T$ is local and $L=K$. In these cases all arity
tensor maps are isomorphisms and the compatible plethory is unique.
If $T$ has at least two maximal ideals and $L\not\subset T$, there
is an $F\in\Omega(R,S)$ for which $F(XY)$ has no expression
$\sum_iG_i(X)H_i(Y)$ with $G_i,H_i\in\Omega(R,S)$.
\item If $\dim R\ge2$, $R\subset L\subset K$, and
$E=L\cap R^\sharp$, unique polynomial operations exist exactly when
$E=R$, equivalently $\Omega(R,L)=R[X]$. If $R\ne E\subset T$,
the binary tensor map is not injective. If $E\not\subset T$, there
is an $F\in\Omega(R,L)$ with no multiplicative expression as above.
For $L=K$ the positive criterion is $\operatorname{depth}R\ge2$.
\end{enumerate}
\end{theorem}
\begin{proof}
The uniqueness and positive all-arity assertions follow from
Theorems~\ref{result-37}--\ref{result-38} and
Corollary~\ref{result-50}; Theorem~\ref{result-27} identifies the depth
condition. Theorem~\ref{result-23} replaces $S$ by $L$ without changing
any polynomial module. Theorem~\ref{result-38} already gives the stated
noninjectivity. We prove the two stronger existence failures together.

In the first case choose $b\in L\setminus T$; in the second choose
$b\in E\setminus T$. The finite normalization $T$ is a Noetherian
normal semilocal domain, with all maximal ideals lying over
$\mathfrak m$. In dimension one all its height-one primes are maximal.
In higher dimension a negative valuation of $b$ can occur only at a
height-one maximal ideal: if a height-one prime $\mathfrak q$ is not
maximal, its contraction $\mathfrak q\cap R$ is not $\mathfrak m$
by integrality, and
$b\in R^\sharp\subset R_{\mathfrak q\cap R}\subset T_{\mathfrak q}$.
Here we use the normal-domain intersection formula
$T=\bigcap_{\operatorname{ht}\mathfrak q=1}T_{\mathfrak q}$,
with each factor a discrete valuation ring
\cite[Section 10.157, Lemma 10.157.6, and Section 10.119]{StacksNormalDomains}.
Since $b\notin T$, at least one negative valuation occurs.

For every height-one maximal ideal $\mathfrak p$, the Chinese remainder
theorem supplies $\pi_{\mathfrak p}\in T$ of valuation one at
$\mathfrak p$ and congruent to $1$ at every other maximal ideal.
Localizing at all maximal ideals gives
$\mathfrak p=\pi_{\mathfrak p}T$. In particular this element has no
other prime divisor. Let
\[
 c=\prod_{v_{\mathfrak p}(b)<0}
             \pi_{\mathfrak p}^{-v_{\mathfrak p}(b)},\qquad d=bc.
\]
The intersection formula puts $d$ in $T$. At each maximal ideal either
$c$ or $d$ is a unit; hence $(c,d)=T$. From $pc+qd=1$ we get
$1/c=p+qb\in T[b]$. Fix a negative branch $\mathfrak p_0$ and write
$\pi=\pi_{\mathfrak p_0}$, $v_0=v_{\mathfrak p_0}$. Then
$\pi^{-1}=(c/\pi)/c\in T[b]$.
There is another maximal ideal $\mathfrak p_1$: this is a hypothesis
in the first case; in the second, otherwise $T$ would be a local
one-dimensional ring, contrary to $\dim T=\dim R\ge2$.

Choose $0\ne a\in (R:T)$. This conductor exists by clearing a finite
set of module denominators. It is contained in $\mathfrak m$, since
otherwise it contains a unit of $R$ and forces $T=R$. Put
$c_0=v_0(a)>0$, $r=c_0+1$, and choose $N>r c_0$.
The digit construction in Theorem~\ref{result-29}, with residue
representatives chosen in $T$, gives
$H\in T[\pi^{-1}][X]$ which is integer-valued at
$T_{\mathfrak p_0}$, has residue one on the class of $a^r$ modulo
$\mathfrak p_0^N$, and has residue zero on the zero class.
The two classes differ because $v_0(a^r)<N$. Its coefficients are
integral at every other maximal ideal, so $H(R)\subset T$. Thus
\[
 g=a(H-H(0))\in\operatorname{Int}(R),\qquad g(0)=0.
\]
Moreover every coefficient of $g$ lies in
$aT[\pi^{-1}]\subset aT[b]\subset R[b]$. Consequently
$g\in\Omega(R,L)$ in both cases.

Suppose $g(X_1\cdots X_r)$ had a finite separated-variable expression
with factors in $\Omega(R,L)\subset\operatorname{Int}(R)$, and call
the finite set of its factors $\mathcal H$. Choose a nonzero common
coefficient denominator $e\in R$ for them and put
$e'=e/\pi^{v_0(e)}$. The intersection formula gives $e'\in T$, and
$e'\notin\mathfrak p_0$. The ideals $e'\mathfrak p_1$ and
$\mathfrak p_0^N$ are comaximal; choose
\[
 x\in e'\mathfrak p_1,\qquad x\equiv1\pmod{\mathfrak p_0^N},
 \qquad u=ax\in(R:T)\subset R.
\]
Because $\pi$ is a unit at $\mathfrak p_1$, we have
$u/e\in\mathfrak p_1T_{\mathfrak p_1}$. Every nonconstant term
of $h(u)-h(0)$, for $h\in\mathcal H$, therefore lies in that ideal.
Both values lie in $R$, whence
\[
 h(u)-h(0)\in\mathfrak p_1T_{\mathfrak p_1}\cap R=\mathfrak m.
\]
Apply the difference operator in every variable between zero and $u$
to the proposed expression. The left side is $g(u^r)$ because
$g(0)=0$; the right side belongs to $\mathfrak m^r$, so has valuation
at least $r$ at $\mathfrak p_0$. But $u^r\equiv a^r$ modulo
$\mathfrak p_0^N$, and the construction of $H$ gives
$v_0(g(u^r))=c_0<r$, a contradiction.

Finally, if every element had a binary multiplicative expression,
repeated finite substitution into its factors would give such an
expression in any finite number of variables. The contradiction for
$g(X_1\cdots X_r)$ therefore proves that some element has no binary
expression. It does not assert that every element lacks one.
\end{proof}


\section{Binary injectivity over an arbitrary Noetherian base}
\label{sec:binary-injective-classification}

This section proves that binary tensor injectivity determines every
polynomial tensor map over a Noetherian base. Let $R\subset S$ be a
unital inclusion of commutative rings, with
$R\ne0$.  Put
\[
 \Omega_n(R,S)=\{F\in S[X_1,\ldots,X_n]:F(R^n)\subseteq R\},
 \qquad \Omega(R,S)=\Omega_1(R,S),
\]
and let
\[
 \theta_n:\Omega(R,S)^{\otimes_R n}\longrightarrow\Omega_n(R,S)
\]
multiply one-variable polynomials in distinct variables.  All tensor
products in this section are algebraic tensor products.  Coefficients
in $S$ may have $R$-torsion.

\begin{theorem}[Binary injectivity and the local classification]
\label{thm:binary-injective-classification}
Suppose that $R$ is Noetherian.  For every maximal ideal
$\mathfrak m$ with finite residue field, define the subring
\[
 C_{\mathfrak m}
 =\bigcup_{j\ge0}
   (R_{\mathfrak m}:_{S_{\mathfrak m}}
                    (\mathfrak mR_{\mathfrak m})^j).
\]
The following conditions are equivalent.
\begin{enumerate}[label=\textup{(\roman*)}]
\item $\theta_2$ is injective.
\item $\theta_n$ is an isomorphism for every $n\ge1$.
\item For every $F\in\Omega(R,S)$, each of $F(X+Y)$ and $F(XY)$
      has exactly one preimage under $\theta_2$.
\item At each maximal ideal $\mathfrak m$ with finite residue field,
      either $C_{\mathfrak m}=R_{\mathfrak m}$, or
      $R_{\mathfrak m}$ is a one-dimensional domain whose completion
      is a domain and
      \[
       C_{\mathfrak m}=\operatorname{Frac}(R_{\mathfrak m})
       \subseteq S_{\mathfrak m}.
      \]
      The last inclusion is the canonical one extending
      $R_{\mathfrak m}\subseteq S_{\mathfrak m}$.
\end{enumerate}
There is no additional condition at maximal ideals with infinite
residue field.  When these conditions hold, $\Omega(R,S)$ has a unique
compatible $R$-plethory structure.
\end{theorem}

The implication \textup{(iii)}$\Rightarrow$\textup{(i)} already follows
from uniqueness for the zero polynomial.  The substantive necessity
statement is that binary injectivity alone forces \textup{(iv)}.
We first establish the base-change statements and the tensor
obstructions needed for that implication.

\Needspace{8\baselineskip}
\subsection{Base change and the effective coefficients}

\begin{lemma}[Localization and completion with arbitrary coefficients]
\label{lem:arbitrary-coefficient-basechange}
Let $R$ be Noetherian and $U\subset R$ multiplicatively closed.  Then
\[
 U^{-1}\Omega_n(R,S)
 =\Omega_n(U^{-1}R,U^{-1}S).
\]
If $(R,\mathfrak m)$ is local and $A=\widehat R$, then
\[
 A\otimes_R\Omega_n(R,S)
 =\Omega_n(A,A\otimes_R S).
\]
These identifications respect the maps $\theta_n$ and polynomial
substitution.  For a local base with finite residue field, every
coefficient occurring in any $\Omega_n(R,S)$ lies in
\[
 C=\bigcup_{j\ge0}(R:_S\mathfrak m^j).
\]
Moreover,
\[
 A\otimes_R C
 =\bigcup_{j\ge0}
   (A:_{A\otimes_R S}(\mathfrak mA)^j).
\]
For a local base with infinite residue field,
$\Omega_n(R,S)=R[X_1,\ldots,X_n]$.
\end{lemma}
\begin{proof}
For the forward localization inclusion, induct on total degree.
If $F$ has total degree $d$ and $s\in U$, then
$s^dF(\mathbf X)-F(s\mathbf X)$ has smaller total degree and still
takes $R^n$ into $R$.  Evaluating at $\mathbf a/s$, with
$\mathbf a\in R^n$, proves the assertion by induction.  Every tuple
in $(U^{-1}R)^n$ admits a common denominator.

Conversely, clear the finitely many coefficient denominators of a
polynomial over $U^{-1}S$, obtaining $F\in S[\mathbf X]$ whose values
have integral images after localization.  The coefficient classes
generate a finite $R$-submodule $M\subseteq S/R$.  The classes of
the values of $F$ generate a submodule of $M$, hence a finite
submodule.  It vanishes after localization, so one element of $U$
annihilates all its generators.  Multiplication by that element
makes $F$ integer-valued on $R^n$.

For completion, fix a finite rectangular set of monomials, of
cardinality $N$.  The evaluation rows
$(\mathbf a^\alpha)_\alpha$, for $\mathbf a\in R^n$, span a
submodule of $R^N$.  Noetherianity implies that finitely many actual
evaluation rows generate this submodule.  Integer-valuedness within
the chosen degree bound is therefore the kernel of a finite map
\[
 S^N\longrightarrow(S/R)^r.
\]
Tensoring with the flat $R$-algebra $A$ identifies this kernel with
the polynomials over $A\otimes_R S$ having integral values on $R^n$.
If $F$ is such a polynomial, let $M_F$ be the finite $A$-module
generated by its coefficient classes in $(A\otimes_R S)/A$.
For $\mathbf x\in A^n$, choose
$\mathbf a_k\in R^n$ with
$\mathbf x-\mathbf a_k\in(\mathfrak m^kA)^n$.  Then
$[F(\mathbf x)]\in\mathfrak m^kM_F$ for every $k$.
Krull intersection gives $[F(\mathbf x)]=0$.  Taking the union
over the degree bounds proves the completion formula.

Suppose next that $(R,\mathfrak m)$ has finite residue field, and
let $c$ be a coefficient of an integer-valued polynomial.
For each nonmaximal prime $\mathfrak p$, the domain
$R/\mathfrak p$ is infinite.  Rectangular Vandermonde interpolation
on a grid of distinct residues modulo $\mathfrak p$ gives
$v\notin\mathfrak p$ with $vc\in R$; Cramer's rule suffices,
so no coefficient torsion is discarded.  Thus
$\{a\in R:ac\in R\}$ is either $R$ or an ideal with radical
$\mathfrak m$.  It contains a power of $\mathfrak m$, proving
$c\in C$.  If the residue field is infinite, the grid may instead
be chosen with unit differences, giving $c\in R$ directly.

Write $C_j=(R:_S\mathfrak m^j)$.  These submodules contain $R$
and satisfy $C_jC_k\subseteq C_{j+k}$.  Thus $C$ is a subring.
For generators $a_1,\ldots,a_r$ of $\mathfrak m^j$,
\[
 C_j=\ker\bigl(S\longrightarrow(S/R)^r,\ 
                    c\longmapsto(a_ic\bmod R)_i\bigr).
\]
Flat completion preserves this kernel.  Taking the filtered union
proves the asserted formula for $C$.  All identifications arise
from the same polynomial evaluations and multiplications, giving
the stated compatibility.
\end{proof}
\begin{lemma}[Binary injectivity forces a ring epimorphism]
\label{lem:binary-effective-epimorphism}
Let $(A,\mathfrak m)$ be Noetherian local with finite residue field,
let $A\subset C$, and suppose that $C/A$ is killed elementwise by
powers of $\mathfrak m$.  If the binary tensor map for
$\Omega(A,C)$ is injective, then
$A\to C$ is a ring epimorphism.  Equivalently, multiplication
$C\otimes_A C\to C$ is an isomorphism.
\end{lemma}
\begin{proof}
For $c\in C$, choose $j\ge1$ with $\mathfrak m^jc\subseteq A$.
Choose one lift in $A$ of each element of the finite ring
$A/\mathfrak m^j$, and let $P$ be the product of the corresponding
linear factors.  This monic polynomial satisfies
$P(A)\subseteq\mathfrak m^j$, so $P,cP\in\Omega(A,C)$.
The tensor $cP\otimes P-P\otimes cP$ has zero polynomial image.
Injectivity and extraction of the coefficient of $X^{\deg P}$
in each factor give
$c\otimes1=1\otimes c$ in $C\otimes_A C$.
Consequently every $c\otimes d$ equals $cd\otimes1$.
This is the ring-epimorphism criterion
\cite[Lemmas 10.107.1--10.107.3, Tag 04VM]{StacksRingEpimorphisms}.
\end{proof}

\begin{lemma}[The positive local cases]
\label{lem:binary-positive-local}
If $C=A$, every polynomial tensor map is an isomorphism.
The same holds for $C=\operatorname{Frac}(A)$ when $A$ is a
one-dimensional analytically irreducible Noetherian local domain.
\end{lemma}
\begin{proof}
The first case is the ordinary polynomial ring.  In the second case,
Theorem~\ref{thm:int-analytic-free} makes
$M=\operatorname{Int}(A)$ free over $A$.  Let $(f_i)$ be a basis.
After extending scalars to $K=\operatorname{Frac}(A)$ it is a
$K$-basis of $K[X]$, since $A[X]\subseteq M\subseteq K[X]$.
Expand $F\in\operatorname{Int}(A^n)$ in the first variable as
\[
 F(X_1,\ldots,X_n)=\sum_i f_i(X_1)G_i(X_2,\ldots,X_n),
\]
with finitely many terms.  Specializing the other variables to
$A$ gives an element of $M$, so uniqueness of its basis
coefficients shows that each $G_i$ is integer-valued.
Induction gives surjectivity of $\theta_n$.
The products of the $f_i$ in distinct variables are linearly
independent over $K$, and hence over $A$, proving injectivity.
\end{proof}

\subsection{Finite-precision approximation through nilpotents}

For a one-dimensional Cohen--Macaulay local ring, a
\emph{regular parameter} means a nonzerodivisor in the maximal ideal
whose principal ideal is primary to that maximal ideal.
If $A$ is complete and $t$ is such a parameter, then
$Q(A)=A[1/t]$, and the $t$-adic and maximal-ideal topologies agree.

We use the following classical approximation statement:
for a complete one-dimensional Noetherian local domain $D$ with
finite residue field, $\operatorname{Int}(D)$ is uniformly dense in
the continuous functions $D\to D$.  In particular it approximates
every clopen indicator to any prescribed parameter-adic precision.
This is the complete-domain case of
\cite[Proposition 4.3, author version pp.\ 18--19]{CahenChabert1999Skolem};
see also \cite[Section 4, pp.\ 52--53]{CahenGrazziniHaouat1982}.

\begin{lemma}[Indicators with small derivatives]
\label{lem:binary-primary-indicators}
Let $(A,\mathfrak m)$ be a complete one-dimensional Noetherian
Cohen--Macaulay local ring with finite residue field and exactly
one minimal prime $N$.  Put $D=A/N$ and $Q=Q(A)$.
Fix a regular parameter $t$, and write $\bar t$ for its image in $D$.
For every clopen subset $U\subseteq D$ and every $b\ge1$, there is
$H\in Q[X]$ such that, for all $a\in A$,
\[
 H(a)-\mathbf1_U(\bar a)\in t^bA,
 \qquad H'(a)\in t^bA.
\]
In particular $H\in\operatorname{Int}(A)$.
\end{lemma}
\begin{proof}
The ring $Q$ is Artinian local, with nilradical $NQ$ and residue
field $K=\operatorname{Frac}(D)$.  Choose $h\ge1$ with $(NQ)^h=0$;
when $N=0$, take $h=1$.  Choose $e\ge1$ with
$\mathfrak m^e\subseteq tA$.  If $p$ is the residue characteristic,
then
\[
 p^{eM}\in t^MA\qquad(M\ge0),
\]
whether $p$ is regular, nilpotent, or a nonzero zero divisor in $A$.

Approximate $\mathbf1_U$ modulo $\bar tD$ by a polynomial
$f\in K[X]$.  Lift its finitely many coefficients arbitrarily along
$Q\to K$, obtaining $F\in Q[X]$.  For each $a\in A$ write
\[
 F(a)=\chi+tb(a)+n(a),\qquad
 \chi=\mathbf1_U(\bar a)\in\{0,1\},\quad
 b(a)\in A,\quad n(a)\in NQ.
\]
Choose $L,L'\ge0$ such that the coefficients of $F$ lie in
$t^{-L}A$ and $F'(A)\subseteq t^{-L'}A$.  The displayed
decomposition then has $n(a)\in t^{-L}A$ uniformly in $a$.
No ring section $K\to Q$ is used.

Set $q=p^r$.  Choose $r$ so large that all of the following finite
conditions hold:
\begin{align*}
 q&\ge\max\{h,\ b+(L+1)(h-1)\},\\
 r&\ge v_p(j)+e(b+Lj) &&(1\le j<h),\\
 r&\ge v_p(j)+eb &&(1\le j<b),\\
 r&\ge e\bigl(b+L(h-1)+L'\bigr).
\end{align*}
For $1\le j<q$, the integer identity
$j\binom qj=q\binom{q-1}{j-1}$ gives
$v_p(\binom qj)\ge r-v_p(j)$.
If $\chi=0$, nilpotence yields
\[
 F(a)^q=\sum_{0\le j<h}\binom qj
                  (tb(a))^{q-j}n(a)^j,
\]
whose terms lie in $t^{q-j-Lj}A\subseteq t^bA$.
If $\chi=1$, put $v=1+tb(a)$.  Then
\[
 F(a)^q=v^q+\sum_{1\le j<h}\binom qj v^{q-j}n(a)^j.
\]
The second line of bounds makes each nilpotent term lie in $t^bA$.
In the binomial expansion of $v^q-1$, the terms with $j\ge b$
already lie in $t^bA$, while the third line handles the remaining
finitely many terms.  Thus $F(a)^q\equiv\chi\pmod{t^bA}$ uniformly.

For every integer $m\ge0$, the same nilpotent expansion gives
$F(a)^m\in t^{-L(h-1)}A$ uniformly.  Consequently
\[
 (F^q)'(a)=p^rF(a)^{q-1}F'(a)
       \in p^rt^{-L(h-1)-L'}A\subseteq t^bA.
\]
Taking $H=F^q$ proves both conclusions.
\end{proof}

\begin{lemma}[Small reduced values and a unit derivative]
\label{lem:binary-primary-small-values}
Under the hypotheses of Lemma~\ref{lem:binary-primary-indicators},
for every $n\ge1$ there is $g_n\in\operatorname{Int}(A)$ such that
\[
 g_n(A)\subseteq t^nA+N,\qquad
 g_n'(0)\in1+t^nA\subseteq A^\times .
\]
\end{lemma}
\begin{proof}
Choose representatives $d_1,\ldots,d_s$ of the finite ring
$D/\bar t^nD$ and lifts $a_i\in A$.  Apply the preceding lemma
at precision $n$ to the clopen partition
$U_i=d_i+\bar t^nD$, obtaining $H_i$.
Then
\[
 g_n(X)=X-\sum_i a_iH_i(X)
\]
has the required reduced values and derivative.
\end{proof}

\Needspace{8\baselineskip}
\subsection{Nilpotent and multibranch tensor kernels}

\begin{theorem}[Nilpotents and unbounded parameter torsion]
\label{lem:nilpotent-without-section}
Let $(A,\mathfrak m)$ be a complete one-dimensional Noetherian
Cohen--Macaulay local ring with finite residue field.  Put
$M=\operatorname{Int}(A)\subseteq Q(A)[X]$ and let $t$ be any
regular parameter.  If $A$ is nonreduced, then
\[
 \ker\bigl(M\otimes_A M\longrightarrow\operatorname{Int}(A^2)\bigr)
\]
is nonzero.  Every element of this kernel is annihilated by some
power of $t$, but no one power annihilates the whole kernel.
This holds in every characteristic, including when the residue
characteristic is a nonzero zero divisor.
\end{theorem}
\begin{proof}
First suppose that $A$ has a unique minimal prime $N$.
Write $Q=Q(A)$, $D=A/N$, $K=\operatorname{Frac}(D)$, and let
$\rho:Q\to K$ be reduction.  Choose $h\ge2$ with $(NQ)^h=0$.
There is one fixed regular element
\[
 \delta=\prod_{0\le i<j<h}(t^j-t^i)\in A
\]
such that
\begin{equation}\label{eq:binary-uniform-nil-derivative}
 \delta zF'(0)\in A\qquad(z\in N,\ F\in M).                 
\end{equation}
Indeed, $F(zT)$ has degree at most $h-1$.  Its values at
$1,t,\ldots,t^{h-1}$ belong to $A$, and its linear coefficient
is $zF'(0)$.  The adjugate of this fixed Vandermonde matrix gives
\eqref{eq:binary-uniform-nil-derivative}.  Every factor of $\delta$ is a power of $t$ times a unit,
so $\delta$ is regular.  In particular,
\begin{equation}\label{eq:binary-nil-submodule-derivative}
 \delta G'(0)\in A\qquad(G\in NM).                         
\end{equation}

Choose a nonzero socle element of the Artinian local ring $Q$.
Its socle lies in $NQ$, since $Q$ is nonreduced.
After multiplication by a power of $t$, it gives $u\in A$ with
\[
 u\ne0,\qquad uNQ=0,\qquad u^2=0,\qquad
 \operatorname{Ann}_A(u)=N.
\]
For $\lambda\in K$, multiplication of $u$ by any lift of
$\lambda$ to $Q$ is independent of the lift.  Thus
\[
 L_u=\{\lambda\in K:u\lambda\in A\}
\]
is a $D$-module.  Multiplication by $u$ identifies it with
$A\cap uQ$, a finite $A$-module annihilated by $N$.
Hence $L_u$ is a finite $D$-submodule of $K$.
Choose $0\ne b\in D$ with $bL_u\subseteq D$, and choose a lift
$\widetilde b\in A$.  This lift is regular.

For $F\in M$, evaluation at the square-zero element $u$ gives
$uF'(0)=F(u)-F(0)\in A$, so $\rho(F'(0))\in L_u$.
Therefore
\begin{equation}\label{eq:binary-nil-detector}
 \Phi:M\otimes_A M\longrightarrow M/NM,\qquad
 \Phi(F\otimes G)=b\rho(F'(0))[G]                          
\end{equation}
is well-defined.  The target is a $D$-module, and the scalar in
\eqref{eq:binary-nil-detector} belongs to $D$.  Balancing is explicit: for $a\in A$,
\[
 b\rho((aF)'(0))[G]
 =b\rho(a)\rho(F'(0))[G]
 =b\rho(F'(0))[aG].
\]
Thus the detector is balanced over $A$ itself.

Take $g_n$ from Lemma~\ref{lem:binary-primary-small-values} and put
\[
 h_n=\frac{u}{t^n}g_n,\qquad
 w_n=h_n\otimes g_n-g_n\otimes h_n.
\]
The reduced-value condition gives $g_n(A)\subseteq t^nA+N$.
Since $uN=0$, one has $h_n(A)\subseteq uA\subseteq A$.
Both tensor factors therefore lie in $M$.  Their antisymmetric
tensor has zero polynomial image because $h_n$ is a scalar
multiple of $g_n$ in $Q[X]$.

Fix $c\ge0$.  If $t^cw_n=0$, apply \eqref{eq:binary-nil-detector}.  Since
$\rho(h_n'(0))=0$, this gives
\[
 \widetilde b\,t^c g_n'(0)h_n\in NM.
\]
Differentiation at zero and \eqref{eq:binary-nil-submodule-derivative} yield
\[
 \delta\widetilde b\,u\,t^{c-n}(g_n'(0))^2\in A.
\]
The derivative $g_n'(0)$ is a unit of $A$, so, for $n>c$,
\begin{equation}\label{eq:binary-nil-krull}
 \delta\widetilde b\,u\in t^{n-c}A.                       
\end{equation}
The left side is a fixed nonzero element, because $\delta$ and
$\widetilde b$ are regular.  Krull intersection shows that \eqref{eq:binary-nil-krull}
fails for all sufficiently large $n$.  Thus $t^cw_n\ne0$ for
all sufficiently large $n$.

Now allow several minimal primes.  Since $A$ is one-dimensional
Cohen--Macaulay, its associated primes are its minimal primes.
Choose the isolated primary components $\mathfrak q_i$ of zero
and put $A_i=A/\mathfrak q_i$.  Then
\[
 A\hookrightarrow\prod_iA_i,\qquad
 Q(A)=\prod_iQ(A_i).
\]
Each $A_i$ is complete, primary, one-dimensional Cohen--Macaulay,
and has the same finite residue field.  The image $t_i$ of $t$
is a regular parameter.  The cokernel of the displayed inclusion
is finite and supported only at $\mathfrak m$: localization at
each minimal prime leaves only the corresponding factor and
gives an isomorphism.  Consequently, for one fixed $d\ge0$,
\begin{equation}\label{eq:binary-primary-conductor}
 t^d\prod_iA_i\subseteq A.                               
\end{equation}
If $A$ is nonreduced, at least one $A_i$ is nonreduced.
Let $e_i$ be its generic idempotent and let $w_n$ be the tensors
constructed for $A_i$ above.  The polynomials
\[
 \widetilde h_n=t^de_i h_n,\qquad
 \widetilde g_n=t^de_i g_n
\]
preserve $A$ by \eqref{eq:binary-primary-conductor}.  Their antisymmetric tensor has zero image.
Projection of coefficients gives a map
$\operatorname{Int}(A)\to\operatorname{Int}(A_i)$, since every
input in $A_i$ lifts to $A$.  The induced map of tensor products
sends the supported tensor to $t_i^{2d}w_n$.  More generally its
$t^c$-multiple maps to $t_i^{c+2d}w_n$, which is nonzero for
large $n$.  This proves the required unboundedness for $A$.

Finally $M[1/t]=Q(A)[X]$, so the localized polynomial tensor map
is an isomorphism.  Every individual kernel element is therefore
annihilated by some power of $t$, as asserted.
\end{proof}

\begin{lemma}[Transfer through a regular conductor]
\label{lem:binary-selected-factor-transfer}
Let $A\subset C$ be a commutative inclusion, and let $e\in C$
be an idempotent such that $eC=Q(A_0)$, where $A_0=eA$ is a
complete one-dimensional Cohen--Macaulay local ring with finite
residue field.  Suppose that $A_0$ is nonreduced and, for a
regular parameter $t_0$ of $A_0$ and some $d\ge0$,
\begin{equation}\label{eq:binary-selected-conductor}
 t_0^deA_0\subseteq A.                                   
\end{equation}
Then the binary tensor map for $\Omega(A,C)$ is not injective.
The same conclusion holds if \eqref{eq:binary-selected-conductor} is obtained from a finite full
order $A_0\subseteq O\subseteq Q(A_0)$ and a regular conductor
$a eO\subseteq A$, with $a$ a nonzerodivisor of $O$.
If this conductor lies in an ideal of $A$, the supported
polynomials in the proof have all their values in that ideal.
\end{lemma}
\begin{proof}
Projection gives an $A$-linear map
$\Omega(A,C)\to\operatorname{Int}(A_0)$ and hence a map from its
binary tensor product to
$\operatorname{Int}(A_0)\otimes_{A_0}\operatorname{Int}(A_0)$.
By Theorem~\ref{lem:nilpotent-without-section}, choose an
antisymmetric kernel tensor
$w=h\otimes g-g\otimes h$ whose $t_0^{2d}$-multiple is nonzero.
The supported polynomials $t_0^deh$ and $t_0^deg$ preserve $A$
by \eqref{eq:binary-selected-conductor}.  Their antisymmetric tensor has zero polynomial image,
while projection sends it to $t_0^{2d}w\ne0$.

For the order formulation, $O/aO$ has finite length as an
$A_0$-module.  Thus $t_0^dO\subseteq aO$ for some $d$,
which implies \eqref{eq:binary-selected-conductor} and preserves any specified ideal containing
the conductor.  No tensor-preservation of injections is used.
\end{proof}

\begin{lemma}[Several reduced branches give a binary kernel]
\label{lem:binary-reduced-branches}
Let $A$ be a complete reduced one-dimensional Noetherian local
ring with finite residue field.  If $A$ has more than one minimal
prime, the binary tensor map for $\operatorname{Int}(A)$ is not
injective.
\end{lemma}
\begin{proof}
Such a ring is Cohen--Macaulay.  Put $D_i=A/\mathfrak p_i$
for its minimal primes, and identify $A$ with its image in
$\prod_iD_i$.  Let $e_i$ be the corresponding generic idempotents.
The conductor argument \eqref{eq:binary-primary-conductor} shows that
\[
 I_i=\{a\in D_i:e_i a\in A\}
\]
is a nonzero ideal of $D_i$.  Choose $0\ne v_i\in I_i$ and
$u_i\in I_i\setminus\mathfrak m_iI_i$.
Polynomial approximation on the complete domain $D_i$ gives
$f_i\in\operatorname{Int}(D_i)$ such that
\[
 f_i(v_i)-f_i(0)\equiv1\pmod{\mathfrak m_i}.
\]
Indeed a sufficiently small clopen neighborhood of $v_i$ misses
zero, and its indicator can be approximated modulo a parameter.
Then $H_i=e_i u_i f_i$ belongs to $\operatorname{Int}(A)$.

For $F\in\operatorname{Int}(A)$ write $F_i$ for its $i$th
component.  The difference
\[
 \Delta_i(F)=F_i(v_i)-F_i(0)
\]
belongs to $I_i$, because $F(e_i v_i)-F(0)\in A$ is supported
only on component $i$.  This defines an $A$-linear map to $I_i$.
Compose with a residue-field linear functional
$I_i/\mathfrak m_iI_i\to\kappa$, where $\kappa=A/\mathfrak m$,
taking $[u_i]$ to $1$.  The resulting map
$r_i:\operatorname{Int}(A)\to\kappa$ satisfies $r_i(H_i)=1$.
For distinct $i,j$, the tensor $H_i\otimes H_j$ has zero
polynomial image, but $r_i\otimes r_j$ sends it to
$1\in\kappa\otimes_A\kappa=\kappa$.  It is therefore nonzero.
\end{proof}

\begin{corollary}[Complete local curves]
\label{cor:binary-complete-cm}
For a complete one-dimensional Noetherian Cohen--Macaulay local
ring $A$ with finite residue field, the binary tensor map for
$\operatorname{Int}(A)$ is injective if and only if $A$ is a
domain.  In that case all its polynomial tensor maps are
isomorphisms.
\end{corollary}
\begin{proof}
The nonreduced and reduced multibranch cases are excluded by the
two preceding obstructions.  A complete local domain is
analytically irreducible, so the converse follows from
Lemma~\ref{lem:binary-positive-local}.
\end{proof}

\subsection{The complete local coefficient ring}

\begin{lemma}[A complete local epimorphic extension]
\label{lem:complete-epic-decomposition}
Let $(A,\mathfrak m)$ be complete Noetherian local with finite
residue field.  Let $A\subset C$ be a ring epimorphism such that
$C/A$ is killed elementwise by powers of $\mathfrak m$.  Define
\[
 D=\bigcap_{j\ge0}\mathfrak m^jC,\qquad I=A\cap D.
\]
Then
\[
 C=A+D,\qquad D=\mathfrak mD=D^2=eC
\]
for an idempotent $e\in C$.  In particular $D$ is a finitely
generated $C$-ideal, and
\[
 C\simeq D\times A/I,\qquad a\longmapsto(ea,a+I)
\]
is an isomorphism of $A$-algebras.
\end{lemma}
\begin{proof}
A nonzero epimorphic algebra over a field is that field.
Indeed, $b\otimes1=1\otimes b$, and applying a linear functional
which takes $1$ to $1$ in one factor forces $b$ to be a scalar.
Since ring epimorphisms are preserved by base change, either
$C/\mathfrak mC=0$ or $C/\mathfrak mC=A/\mathfrak m$.
In the first case $\mathfrak mC=C$, so $D=C$, $I=A$, and $e=1$.

In the other case $\mathfrak n=\mathfrak mC$ is maximal and
$A\to C_{\mathfrak n}$ is a local ring epimorphism.
Ferrand's complete-local epimorphism lemma
\cite[Lemma 2.1, p.\ 5]{Ferrand1968NoetherianMonomorphisms}
implies that it is surjective.  The lemma requires the source to
be complete Noetherian local; it imposes no Noetherian hypothesis
on the target.  Write $C_{\mathfrak n}=A/J$, with $J\subseteq
\mathfrak m$.

The epimorphism $A/J\to C/JC$ has the retraction induced by
localization to $C_{\mathfrak n}$.  A split monomorphism which is
also a ring epimorphism is an isomorphism: if $i$ and $r$ are the
section and retraction, then $ir$ and the identity agree after
precomposition with the epimorphism $i$.
Consequently
\[
 C/JC=A/J,\qquad C=A+JC,\qquad A\cap JC=J.
\]
The ideal $J$ has finitely many $A$-generators, all zero in
$C_{\mathfrak n}$.  Multiplying finitely many localization
denominators gives $v\in C\setminus\mathfrak n$ with $vJC=0$.
Its image in $C/JC=A/J$ is a unit.  Choose $w\in C$ so that
$e=1-vw\in JC$.  For $z\in JC$, one has $z=ez$, whence
$e^2=e$ and $JC=eC$.  Thus
\[
 C=eC\times(1-e)C,\qquad (1-e)C=A/J.
\]
Because $e\in JC$ and $J\subseteq\mathfrak m$, multiplication
by $e$ in an expression for $e$ shows
$\mathfrak m(eC)=eC$.  Krull intersection on the finite
$A$-module $A/J$ now gives
\[
 \bigcap_{j\ge0}\mathfrak m^jC
   =eC\times\bigcap_{j\ge0}\mathfrak m^j(A/J)
   =eC\times0.
\]
Therefore $D=eC$ and $I=J$, proving every assertion, including
the named intersection and the finite generation of $D$.
\end{proof}

\begin{lemma}[An empty closed fiber is a punctured curve]
\label{lem:binary-empty-closed-fiber}
Let $(B,\mathfrak m_B)$ be a nonzero Noetherian local ring and
let $B\subset T$ satisfy
$T/B$ elementwise $\mathfrak m_B$-power torsion and
$\mathfrak m_BT=T$.  Then $B$ is one-dimensional
Cohen--Macaulay and $T=Q(B)$ canonically as a $B$-algebra.
\end{lemma}
\begin{proof}
For every $x\in\mathfrak m_B$, localization gives $B_x=T_x$.
If $x_1,\ldots,x_r$ generate $\mathfrak m_B$, their principal
opens cover $\operatorname{Spec}(T)$, because
$(x_1,\ldots,x_r)T=T$.  These canonical identifications glue to
\[
 \operatorname{Spec}(T)
 \simeq\operatorname{Spec}(B)\setminus\{\mathfrak m_B\}.
\]
The punctured spectrum is affine.  The affine-puncture criterion
\cite[Lemma 31.17.1, Tag 0BCR]{StacksAffinePuncture}
gives $\dim B\le1$.  Dimension zero is impossible, since then
$\mathfrak m_B$ is nilpotent while $\mathfrak m_BT=T\ne0$.

If $\mathfrak m_B^j b=0$, use
$1\in\mathfrak m_B^jT$ to obtain $b=0$ in $T$, and hence in $B$.
Thus the maximal ideal is not an associated prime of $B$.
It follows that $B$ has depth one and is Cohen--Macaulay.
For a regular parameter $t$, the punctured spectrum equals
$D(t)$, so the canonical isomorphism above gives
$T=B[1/t]=Q(B)$.  The last equality follows because every
regular nonunit has a maximal-ideal-primary principal ideal.
\end{proof}

\begin{lemma}[The selected component and its conductor ideal]
\label{lem:selected-punctured-curve}
In Lemma~\ref{lem:complete-epic-decomposition}, suppose that
$D\ne0$ and put $A_0=eA\subset D$.
Then $A_0$ is a complete one-dimensional Cohen--Macaulay local
ring and $D=Q(A_0)$.  The ideal $I$ embeds as an ideal of $A_0$
containing a nonzerodivisor.  If $I$ is proper in $A$, its image
is proper in $A_0$.
\end{lemma}
\begin{proof}
The image $A_0$ is a complete Noetherian local quotient of $A$.
The quotient $D/A_0$ is maximal-ideal-power torsion and
$\mathfrak m_0D=D$.  Lemma~\ref{lem:binary-empty-closed-fiber}
therefore gives the dimension, depth, and total quotient ring.
For $i\in I$, the identity $ei=i$ proves injectivity of
$I\to A_0$ and identifies its image as an ideal.
Since $D/I=C/A$, the class of $e$ is killed by a power of
$\mathfrak m$.  A corresponding power of any regular parameter
of $A_0$ therefore lies in the image of $I$.
If this image were all of $A_0$, then $e\in I\subseteq A$.
A nonzero idempotent in the local ring $A$ equals $1$,
which would give $I=A$.
\end{proof}

\begin{lemma}[A proper closed-fiber factor gives a binary kernel]
\label{lem:binary-proper-closed-fiber}
Under the hypotheses of Lemma~\ref{lem:complete-epic-decomposition},
if $D\ne0$ and $I\ne A$, the binary tensor map for $\Omega(A,C)$
is not injective.
\end{lemma}
\begin{proof}
Put $B=eA$, $T=A/I$, and $J=\ker(A\to B)$.  Then
$I\cap J=0$.  The image $L$ of $I$ in $B$ is proper and
contains a regular element by
Lemma~\ref{lem:selected-punctured-curve}.  Thus $L$ is primary
to the maximal ideal of the one-dimensional ring $B$, and
\[
 E=B/L=A/(I+J)
\]
is a nonzero finite Artinian local ring.  The usual two-ideal
fiber-product identity gives the original inclusion in the form
\begin{equation}\label{eq:binary-closed-fiber-product}
 A=B\times_E T\ \subseteq\ Q(B)\times T=C.                 
\end{equation}
All these local rings have residue field $\kappa=A/\mathfrak m$.

Let $B_i$ be the isolated primary factors of $B$, with generic
idempotents $\varepsilon_i$ in $Q(B)$, and put
\[
 L_i=\{z\in B_i:\varepsilon_i z\in L\}.
\]
Choose a regular parameter $a\in L$.  The finite
primary-component conductor gives
$a^d\prod_iB_i\subseteq B$ for some $d$; multiplying by $a$
shows that
$a_i^{d+1}B_i\subseteq L_i$.  Hence each $L_i$ is a nonzero
finite ideal.  Fix one $i$, put $v_i=a_i^{d+1}$, and choose
$u_i\in L_i\setminus\mathfrak m_iL_i$.
The reduced image of $v_i$ is nonzero.  A clopen neighborhood
separating it from zero, together with
Lemma~\ref{lem:binary-primary-indicators}, gives
$f_i\in\operatorname{Int}(B_i)$ with
\[
 f_i(v_i)-f_i(0)\equiv1\pmod{\mathfrak m_i}.
\]
The polynomial
$H=(\varepsilon_i u_i f_i,0)$ preserves $A$, since its first
component has values in $L$ and its second component is zero.

The element $z=(\varepsilon_i v_i,0)$ belongs to $A$.
For $F\in\Omega(A,C)$, the difference $F(z)-F(0)$ lies in $A$
and is supported only on component $i$ of the first factor
in \eqref{eq:binary-closed-fiber-product}.  Therefore
$\Delta_i(F)=F_i(v_i)-F_i(0)\in L_i$.
This is $A$-linear.  Composing with a $\kappa$-linear functional
$L_i/\mathfrak m_iL_i\to\kappa$ taking $[u_i]$ to $1$ gives
\[
 r:\Omega(A,C)\longrightarrow\kappa,\qquad r(H)=1.
\]

Choose in $T$ one lift $\widetilde e$ of each $e\in E$, and set
\[
 P(X)=\prod_{e\in E}(X-\widetilde e)\in T[X].
\]
It is monic of degree $|E|$ and has values in $\ker(T\to E)$.
Thus $K=(0,P)\in\Omega(A,C)$.
Taking the coefficient of $X^{|E|}$ in the $T$-component and
reducing to $\kappa$ gives an $A$-linear map $s$ with $s(K)=1$.
The tensor $H\otimes K$ has zero polynomial image because
its coefficient components are orthogonal.  However,
$r\otimes s$ sends it to
$1\in\kappa\otimes_A\kappa=\kappa$, so it is nonzero.
This construction also applies when the selected primary factor
$B_i$ is nonreduced.
\end{proof}

\begin{proposition}[Complete local classification]
\label{prop:complete-arbitrary-coefficients}
Let $(A,\mathfrak m)$ be complete Noetherian local with finite
residue field, and let $A\subset C$ have $C/A$ killed elementwise
by powers of $\mathfrak m$.  The binary tensor map for
$\Omega(A,C)$ is injective if and only if
\[
 C=A\quad\text{or}\quad
 \bigl(A\text{ is a one-dimensional domain and }
                   C=\operatorname{Frac}(A)\bigr).
\]
In either case all polynomial tensor maps are isomorphisms and
the coaddition and comultiplication preimages are unique.
\end{proposition}
\begin{proof}
Binary injectivity makes $A\to C$ a ring epimorphism by
Lemma~\ref{lem:binary-effective-epimorphism}.
Use the notation of Lemma~\ref{lem:complete-epic-decomposition}.
If $D=0$, then $C=A$.
If $D\ne0$, Lemma~\ref{lem:binary-proper-closed-fiber}
rules out $I\ne A$.  Hence $I=A$, $e=1$ and $C=D$.
Lemma~\ref{lem:selected-punctured-curve} now gives
$\dim A=\operatorname{depth}A=1$ and $C=Q(A)$.
Corollary~\ref{cor:binary-complete-cm} forces $A$ to be a domain.
The converse and the all-arity assertion follow from
Lemma~\ref{lem:binary-positive-local}.
\end{proof}

\begin{corollary}[A completed domain with arbitrary coefficient torsion]
\label{cor:binary-completed-domain}
Let $(R,\mathfrak m_R)$ be a nonfield Noetherian local domain
with finite residue field, let $K=\operatorname{Frac}(R)$, and
put $A=\widehat R$ and $Q=K\otimes_R A$.
If $A\subset C$ and $C/A$ is elementwise killed by powers of
the maximal ideal of $A$, then binary injectivity, all-arity
tensor isomorphisms, and unique polynomial operations are
equivalent.  They hold exactly when
\[
 C=A\quad\text{or}\quad
 \bigl(\dim R=1,\ A\text{ is a domain},\ C=Q\bigr),
\]
with the second equality understood as an equality of
canonical $A$-algebras.
\end{corollary}
\begin{proof}
Apply Proposition~\ref{prop:complete-arbitrary-coefficients}.
In its nontrivial case $\dim R=\dim A=1$.
Choose $0\ne a\in\mathfrak m_R$.
Its image is nonzero in the one-dimensional local domain $A$,
so $A[1/a]=\operatorname{Frac}(A)$.
Therefore $K\otimes_R A=\operatorname{Frac}(A)$ canonically.
This proves the exact stated form.  The other implications
follow from the proposition and uniqueness for the zero
polynomial.
\end{proof}

\subsection{Canonical descent and globalization}

\begin{proof}[Proof of Theorem~\ref{thm:binary-injective-classification}]
Suppose first that $\theta_2$ is injective.
Lemma~\ref{lem:arbitrary-coefficient-basechange} localizes this
condition at every maximal ideal.  At an infinite-residue
maximal ideal the integer-valued ring is the ordinary polynomial
ring, so there is no further restriction.

Fix a finite-residue maximal ideal, write
$A=R_{\mathfrak m}$ and $S'=S_{\mathfrak m}$, and let $C$ be
the stated effective coefficient subring of $S'$.
The coefficient assertion of the base-change lemma gives
$\Omega_n(A,S')=\Omega_n(A,C)$ for every $n$.
Set $\widehat C=\widehat A\otimes_A C$ and
$\widehat S'=\widehat A\otimes_A S'$.
Flatness embeds $\widehat A\subseteq\widehat C\subseteq\widehat S'$,
and the same lemma identifies the completed polynomial tensor
maps with those over $\widehat A$.
The quotient $\widehat C/\widehat A$ is elementwise killed by
powers of the maximal ideal.  Proposition
\ref{prop:complete-arbitrary-coefficients} applies.

If $\widehat C=\widehat A$, faithful flatness gives $C=A$.
Otherwise $\widehat A$ is a one-dimensional domain and
$\widehat C=\operatorname{Frac}(\widehat A)$.
The original $A$ embeds in its completion, so it is a domain;
dimension invariance under completion gives $\dim A=1$.
It remains to descend the coefficient subring.
For every $0\ne a\in A$, multiplication by $a$ on
$\widehat C$ is bijective.  Its kernel and cokernel are the
faithfully flat base changes of those of multiplication by
$a$ on $C$, so $a:C\to C$ is bijective.
In particular there is an element $c_a\in C$ with $ac_a=1$.
Thus all nonzero elements of $A$ have their unique inverses
inside $C\subseteq S'$, giving the canonical embedding
\[
 \operatorname{Frac}(A)\subseteq C\subseteq S'.
\]
Conversely, for $c\in C$, choose $j$ with
$\mathfrak m_A^jc\subseteq A$ and a nonzero
$a\in\mathfrak m_A^j$; for $j=0$ use $a=1$.
Then $c=(ac)/a$ lies in that canonical fraction field.
Hence $C=\operatorname{Frac}(A)$, proving condition
\textup{(iv)} without any initial torsion-free hypothesis on $S'$.

Conversely, assume \textup{(iv)}.
At every finite-residue maximal ideal, the localized polynomial
maps are isomorphisms by Lemma~\ref{lem:binary-positive-local};
at every infinite-residue maximal ideal they are ordinary
polynomial tensor maps.  Localization detects the kernels and
cokernels of arbitrary module homomorphisms.  Thus every global
$\theta_n$ is an isomorphism, proving \textup{(ii)}.
Both substitutions $F(X+Y)$ and $F(XY)$ belong to
$\Omega_2(R,S)$, so \textup{(ii)} implies \textup{(iii)}.
As observed above, uniqueness for $F=0$ gives
\textup{(iii)}$\Rightarrow$\textup{(i)}.

The tensor isomorphisms transport polynomial coaddition,
comultiplication, scalar operations and the evaluations at
$0$ and $1$ to $\Omega(R,S)$.
Their identities hold after applying the appropriate $\theta_n$,
and injectivity gives the same identities before applying it.
Composition also preserves $\Omega(R,S)$, because both
polynomials take $R$ into $R$.
Transporting these composition identities gives the compatible
plethory structure.  Every compatible operation has the stated
polynomial image, and the tensor maps are injective, so this
structure is unique.
\end{proof}

\begin{proof}[Proof of Theorem~\ref{thm:arbitrary-noetherian-coefficients}]
The assertion is the full four-way equivalence in
Theorem~\ref{thm:binary-injective-classification}, together with
its plethory conclusion.
\end{proof}

\begin{corollary}[Domain bases with coefficient torsion]
\label{cor:domain-torsion-coefficients}
Let $R$ be a nonfield Noetherian domain with fraction field $K$,
and let $R\subset S$ be any unital commutative inclusion.
Binary injectivity, all-arity tensor isomorphisms, and unique
polynomial operations are equivalent.  They hold exactly when
every finite-residue maximal ideal $\mathfrak m$ satisfies
\[
 C_{\mathfrak m}=R_{\mathfrak m}
 \quad\text{or}\quad
 \bigl(\dim R_{\mathfrak m}=1,\
       \widehat{R_{\mathfrak m}}\text{ is a domain},\
       C_{\mathfrak m}=K\bigr).
\]
The equality with $K$ specifies its canonical copy in
$S_{\mathfrak m}$.  No torsion-free hypothesis on $S$ is needed.
\end{corollary}
\begin{proof}
The fraction field of any localization of a domain is $K$,
so this is precisely the domain specialization of
Theorem~\ref{thm:binary-injective-classification}.
\end{proof}

\begin{corollary}[A selected factor with a nonzero conductor]
\label{cor:binary-selected-nonzero-conductor}
Let $(A,\mathfrak m)$ be complete Noetherian local with finite
residue field, and let $A\subset C$ be a unital commutative
inclusion.  Suppose that an idempotent $e\in C$ selects a
nonreduced local Artinian ring $Q_0=eC$, that $A_0=eA$ is
one-dimensional Cohen--Macaulay, and that a finite full order
\[
 A_0\subseteq O\subseteq Q_0,\qquad Q(O)=Q_0,
\]
has a nonzero conductor into $A$:
\[
 \mathfrak c=\{x\in O:e xO\subseteq A\}\ne0.
\]
Then the binary tensor map for $\Omega(A,C)$ is not injective.
Here the conductor is only required to be nonzero; it may
consist entirely of zero divisors.
\end{corollary}
\begin{proof}
The image $A_0$ is a complete local quotient of $A$ with the
same finite residue field.  A regular parameter $t_0$ of $A_0$
is not nilpotent.  Its image in the Artinian local ring $Q_0$
is therefore a unit.  Every regular element of $A_0$ is likewise
a unit in $Q_0$, so $Q(A_0)$ embeds in $Q_0$.
The Artinian ring $Q(A_0)$ cannot have a nontrivial idempotent,
because $Q_0$ has none.  Hence $A_0$ has exactly one minimal
prime.

The finite $A_0$-algebra $O$ is $t_0$-torsion-free, and
$O[1/t_0]$ is Artinian.  Every regular element of $O$ remains
regular there and hence becomes a unit.  Thus
$O[1/t_0]=Q(O)=Q_0$.
The conductor $\mathfrak c$ is an ideal of $O$, is contained
in $A_0$, and remains nonzero after inverting $t_0$.
A nonzero ideal of a local Artinian ring meets its socle.
Choose such a socle element in $\mathfrak c[1/t_0]$ and clear
a power of $t_0$, obtaining
\[
 0\ne u\in\mathfrak c\cap\operatorname{Soc}(Q_0).
\]
The socle of the nonreduced $Q_0$ is contained in its nilradical,
so $u^2=0$.  It also annihilates the nilradical of the subring
$Q(A_0)$.

Apply Lemma~\ref{lem:binary-primary-small-values} to $A_0$ and
$t_0$, and define
\[
 h_n=\frac{u}{t_0^n}g_n\in Q(A_0)[X].
\]
As before, $h_n(A_0)\subseteq uA_0\subseteq\mathfrak c$.
The supported polynomial $e h_n$ therefore belongs to
$\Omega(A,C)$.  For arbitrarily large $n$, it does not belong
to $A[X]$: otherwise projection of its derivative at zero
would give
$u g_n'(0)/t_0^n\in A_0$, hence
$u\in t_0^nA_0$, contradicting separatedness since $g_n'(0)$
is a unit.

Thus $\Omega(A,C)\ne A[X]$, while $eu\in A$ is a nonzero
square-zero element.  If the binary tensor map were injective,
Theorem~\ref{thm:binary-injective-classification}, applied to
this complete local base, would force either the ordinary
polynomial case or a domain base.  Both conclusions contradict
what was just proved.
\end{proof}


\subsection{Conductors over more general local bases}

The residue field and the regularity of the conductor play different
roles in the tensor obstruction.  The next two examples separate them.

\begin{proposition}\label{prop:infinite-residue-nilpotent-conductor}
There is a complete one-dimensional Noetherian Cohen--Macaulay local
ring $A$ and an inclusion $A\subset C$ with a selected nonreduced
Artinian factor and a nonzero conductor, consisting entirely of
zero divisors, for which every polynomial tensor map is an
isomorphism.  Its residue field is infinite.
\end{proposition}
\begin{proof}
Let $k$ be an infinite field, $D=k[[t]]$, and $K=k((t))$.  Set
\[
 A=D[\varepsilon,\eta]/(\varepsilon^2,\varepsilon\eta,\eta^2),
 \qquad
 C=(K[\varepsilon]/\varepsilon^2)
       \times(D[\eta]/\eta^2).
\]
Embed $A$ in $C$ by
\[
 a+b\varepsilon+c\eta\longmapsto
             (a+b\varepsilon,a+c\eta).
\]
This is an injective unital ring homomorphism.  The ring $A$ is
finite free over $D$, complete local with maximal ideal
$(t,\varepsilon,\eta)$, and has reduced quotient $D$.
Thus it is Noetherian of dimension one.  The parameter $t$ is
regular, proving the Cohen--Macaulay assertion.

For $e=(1,0)$, the selected image is
$A_0=eA=D[\varepsilon]/\varepsilon^2$ and the selected factor is
$Q_0=eC=K[\varepsilon]/\varepsilon^2=Q(A_0)$.
Take $O=A_0$.  Its conductor into $A$ is
\[
 \{x\in O:exO\subset A\}=D\varepsilon.
\]
Indeed, if $(a+b\varepsilon,0)$ belongs to $A$, its common scalar
coefficient is zero, so $a=0$; conversely every supported
$\varepsilon$-multiple lies in $A$.  The conductor is nonzero,
square zero, and contains no regular parameter multiple of $O$.

For any polynomial of degree at most $d$ in an arbitrary coefficient
extension of $A$, choose $d+1$ distinct elements of $k$.  Their
Vandermonde determinant is a unit of $A$.  If all polynomial values
belong to $A$, the inverse Vandermonde matrix therefore puts every
coefficient in $A$.  Tensor products of these matrices prove the
same assertion for a finite grid in several variables.  Hence
\[
 \Omega_n(A,C)=A[X_1,\ldots,X_n]
\]
for every $n$, and every tensor map is an isomorphism.
\end{proof}

\begin{proposition}\label{prop:nonnoetherian-nilpotent-conductor}
Let $k=\mathbb F_q$, $D=k[[t]]$, $K=k((t))$, and let $N$ be a
$D$-module.  Give
\[
 A=D\oplus D\varepsilon\oplus N,
 \qquad A_1=D\oplus N
\]
the indicated square-zero extensions of $D$, and embed $A$ in
\[
 C=(K[\varepsilon]/\varepsilon^2)\times A_1,
 \qquad
 d+a\varepsilon+n\longmapsto(d+a\varepsilon,d+n).
\]
The binary tensor map for $\Omega(A,C)$ is not injective.
For $N=\prod_{j\ge1}D$, the ring $A$ is complete separated local
and non-Noetherian, with finite residue field and a nonzero
square-zero selected conductor.
\end{proposition}
\begin{proof}
The map is a unital injection, and the selected order for $e=(1,0)$
is $A_0=D[\varepsilon]/\varepsilon^2$.  Its conductor is
$D\varepsilon$, by the same common-scalar calculation as above.

Write the first component of $F\in\Omega(A,C)$ as
$f_0+\varepsilon f_1$, with $f_i\in K[X]$.
The scalar component of its second polynomial lies in $D[X]$.
Equality of the two scalar values on the infinite domain $D$
forces $f_0\in D[X]$.  Evaluation at $d\in D$ then gives
$f_1(d)\in D$ for every $d$, so $f_1\in\operatorname{Int}(D)$.
Define
\[
 \lambda:\Omega(A,C)\longrightarrow k,
 \qquad \lambda(F)=f_1(t)-f_1(0)\pmod {tD}.
\]
This is $A$-linear through $A\to k$.  In fact, multiplication by
$d+a\varepsilon+n$ changes $f_1$ to $df_1+af_0$, and
$f_0(t)-f_0(0)\in tD$.

The polynomial
\[
 H=\left(\frac{\varepsilon(X^q-X)}t,0\right)
\]
belongs to $\Omega(A,C)$ because $d^q-d\in tD$ for $d\in D$
and the nilpotent part of the argument is annihilated by
$\varepsilon$.  Its value under $\lambda$ is
$t^{q-1}-1\pmod t=-1$.
Since $H(X)H(Y)=0$, the tensor $H\otimes H$ is in the binary
kernel.  It is nonzero: the $A$-balanced map
\[
 \Omega(A,C)\otimes_A\Omega(A,C)\longrightarrow k,
 \qquad F\otimes G\longmapsto\lambda(F)\lambda(G)
\]
sends it to $1$.

For the specified product module, $A$ is local and
\[
 \mathfrak m^r=t^rD\oplus t^{r-1}(D\varepsilon\oplus N)
 \qquad(r\ge1).
\]
Both $D$ and $\prod_{j\ge1}D$ are complete and separated
$t$-adically; the displayed equality proves the same assertion
for the maximal-ideal topology of $A$.
The vector space $N/tN=\prod_{j\ge1}k$ is infinite-dimensional.
Thus $N$ is not finitely generated as a $D$-module.  The
square-zero ideal $N\subset A$ has its $A$-action through $D$,
so it is not finitely generated as an $A$-ideal either.
Consequently $A$ is not Noetherian.
\end{proof}


\section{Further criteria without Noetherian hypotheses}

\subsection{Bases with no proper injective epimorphism}

\begin{proposition}\label{prop:epimorphism-rigid-bases}
Every injective ring epimorphism $R\to C$ is an isomorphism in each
of the following situations:
\begin{enumerate}
\item $R$ is absolutely flat.
\item $R/N$ is absolutely flat and $N$ is a T-nilpotent ideal.
\item $R$ is a directed union of unital self-injective subrings.
\item $R$ is local with nil maximal ideal $\mathfrak m$, and there is
an ideal $I$ with $\mathfrak m^NI=0$ for some $N$, such that every
injective epimorphism from $R/I$ is an isomorphism.
\end{enumerate}
For any such $R$ and any inclusion $R\subset S$, unique polynomial
co-operations are equivalent to
\[
              (R:_S D_d(R))=R\quad\text{for every }d,
 \qquad\text{equivalently}\qquad\Omega(R,S)=R[X].
\]
Every positive case has all multivariable tensor isomorphisms.
\end{proposition}
\begin{proof}
In the absolutely flat case $C/R$ is flat, so tensoring
$0\to R\to C\to C/R\to0$ with any module preserves the initial
injection. Taking that module to be $C/R$, epicness gives
$C\otimes_R(C/R)=0$, hence $C/R=0$.

For (2), the nil ideal $NC$ contracts to $N$: if its contraction
were larger, its image in the reduced ring $R/N$ would contain a
nonzero nilpotent. The induced injective epimorphism
$R/N\to C/NC$ is therefore an isomorphism, giving $C=R+NC$.
For a fixed element of $C$, repeatedly expand each coefficient by
this equality. This produces a finitely branching tree; label its
edges by the elements of $N$ occurring in the expansions and prune
every branch whose product is zero. An infinite branch would
contradict T-nilpotence. K\"onig's lemma makes the remaining tree
finite, so the original element is a finite sum of elements of $R$.

For (3), write $R=\bigcup R_\alpha$. The equality
$c\otimes1=1\otimes c$ in $C\otimes_R C$ uses only finitely many
balancing relations, so it already holds over one $R_\alpha$.
Self-injectivity extends the identity of $R_\alpha$ to an
$R_\alpha$-linear retraction $\rho:C\to R_\alpha$. Applying
$\rho\otimes1$ gives $c=\rho(c)\in R$.

For (4), $\mathfrak mC$ is nil and therefore proper. Reduction gives
an epimorphism from the residue field to the nonzero algebra
$C/\mathfrak mC$, hence an isomorphism. Thus $C=R+\mathfrak mC$.
Multiplication by $I$ and finite iteration give
$IC=I+\mathfrak mIC=I+\mathfrak m^NIC=I$. Consequently
$R/I\subset C/I$ is an injective epimorphism, so $C/I=R/I$ and
$C=R$.

In every case apply the coefficient-ring obstruction to
$C=\bigcup_d(R:_S D_d(R))$. Tensor injectivity forces epicness,
and the preceding conclusions force $C=R$. The converse is the
ordinary polynomial case. The equivalence between absolute
flatness and reduced zero-dimensionality is standard
\cite[Tag~092F]{StacksAbsolutelyFlat}.
\end{proof}

\begin{corollary}\label{cor:boolean-and-unbounded-nilpotent-bases}
For a Boolean base ring $R$, unique polynomial co-operations for
$R\subset S$ occur exactly when $S=R$. The preceding proposition
also determines every coefficient extension of
\[
 R=\mathbb F_2[z,x_1,x_2,\ldots]/
       (z^2,zx_i,x_i^{2^{i+1}}:i\ge1).
\]
This ring is local, zero-dimensional and non-Noetherian. Its maximal
ideal is not T-nilpotent, and its elementwise nilpotence has no
uniform bound.
\end{corollary}
\begin{proof}
For a Boolean ring $D_2(R)=0$: every prime quotient is $\mathbb F_2$,
so every three-node determinant vanishes in all prime quotients,
and the ring is reduced. Thus the effective coefficient ring is
$S$, and the absolutely flat case applies.

For the displayed ring put
$A_n=\mathbb F_2[x_1,\ldots,x_n]/(x_i^{2^{i+1}})$ and $A=\bigcup A_n$.
The coefficient of the top monomial defines a perfect associative
pairing on $A_n$, pairing each monomial with its complement. It
identifies $A_n$ with its vector-space dual as an $A_n$-module.
That dual is injective because
$\operatorname{Hom}_{A_n}(M,\operatorname{Hom}_{\mathbb F_2}(A_n,\mathbb F_2))
\simeq\operatorname{Hom}_{\mathbb F_2}(M,\mathbb F_2)$ is exact.
Apply (3) to $A$, then (4) to $R=A\oplus\mathbb F_2z$ with
$I=\mathbb F_2z$ and $\mathfrak mI=0$.
Every element of $\mathfrak m$ uses finitely many variables and is
nilpotent. The independent classes of the $x_i$ in
$\mathfrak m/\mathfrak m^2$ show non-Noetherianity; the products
$x_1\cdots x_n$ never vanish, and $x_i$ has exact nilpotence order
$2^{i+1}$. These prove all remaining assertions.
\end{proof}

\Needspace{9\baselineskip}
\subsection{Regular interpolation and depth}

\begin{theorem}\label{thm:regular-vandermonde-criterion}
Let $Q(R)$ be the total quotient ring of a nonzero commutative ring
$R$. Suppose every $D_d(R)$ contains a regular element and
$(R:_{Q(R)}D_d(R))=R$. For every inclusion $R\subset S$, the
following are equivalent: unique polynomial co-operations; binary
tensor injectivity; faithful evaluation on $R$; the conditions
$\operatorname{Ann}_S D_d(R)=0$ for all $d$; and
$\Omega(R,S)=R[X]$. All positive cases have every tensor
isomorphism. The hypotheses hold for every infinite domain with
$\operatorname{Int}(R)=R[X]$, for every ring $A[Z,W]$, and for
$A[Z]$ whenever each $D_d(A)$ contains a regular element.
\end{theorem}
\begin{proof}
Localize $S$ at the regular elements of $R$. Cramer's rule says
$D_d(R)a_i\subset R$ for every coefficient of an integer-valued
polynomial. A regular element of this ideal makes its rationalized
coefficient lie in $Q(R)$, and the assumed colon identity puts it
in $R$. Rationalization therefore splits
\[
              \Omega(R,S)=R[X]\oplus\mathcal N,
\]
where $\mathcal N$ is the module of null polynomials. To identify
the kernel, clear the regular annihilators of the finitely many
coefficients; its values lie in $R$, so regularity forces them
to vanish. Conversely Cramer's rule kills each null coefficient
by $D_d(R)$ and therefore by a regular element.

If $\mathcal N\ne0$, choose a nonzero null coefficient $a$ with
$D_d(R)a=0$. The monic polynomial $P_d$ from the coefficient-ring
theorem has values in $D_d(R)$, so $aP_d$ is nonzero and null.
The tensor $aP_d\otimes P_d-P_d\otimes aP_d$ lies in the two
different mixed summands and is nonzero; extraction of the monic
coefficient detects its first component. Its polynomial image is
zero. This proves the injectivity implication. The other
equivalences follow from the faithful-evaluation criterion and
ordinary polynomials.

For an infinite domain, $D_d(R)\ne0$, and the coefficient theorem
identifies the colon condition with $\operatorname{Int}(R)=R[X]$.
For $R=A[Z]$, evaluate a rational integer-valued polynomial on
$A$-points. A regular $v\in D_d(A)$ bounds every coefficient in
$A[1/v][Z]$. Evaluate next at $Z^N$, choosing $N$ larger than all
coefficient degrees in $Z$. The degree intervals are disjoint;
the fact that the value belongs to $A[Z]$ forces all coefficients
to belong to $A[Z]$. Finally $D_d(A[W])$ contains a monic
polynomial: use the determinant on $0,1,W,\ldots,W^{d-1}$
(and $1$ for $d=1$). A monic polynomial is regular, proving
the assertion for arbitrary $A[Z,W]$.
\end{proof}

\begin{proposition}\label{prop:power-generated-and-regular-sequence-bases}
Let $R$ be local with finite residue field and maximal ideal
$\mathfrak m$.
If $\mathfrak m$ is generated by its $N$th powers for every $N$,
then unique polynomial co-operations are equivalent to
$(R:_S\mathfrak m)=R$, equivalently $\Omega(R,S)=R[X]$.
This applies in particular to local rings of characteristic $p$
with surjective Frobenius. If instead $\mathfrak m$ contains a
regular sequence of length two, the same equivalence holds with
$\Omega(R,S)=R[X]$, for every coefficient extension $S$.
\end{proposition}
\begin{proof}
In the power-generated case the determinant calculation in the
proof of Theorem~\ref{thm:valuation-conductor-arbitrary-coefficients}
uses only the local finite residue field and generation by powers:
$D_d(R)=R$ below the residue cardinality and $D_d(R)=\mathfrak m$
thereafter. The same proof with $C=(R:_S\mathfrak m)$ gives
$\mathfrak mC=\mathfrak m$ under epicness and then $C=R$.
Surjective Frobenius gives generation by every sufficiently large
$p$-power and consequently by every fixed $N$th power.

For a regular sequence $a,t$, both $a^e$ and $t^e$ belong to
$D_d(R)$ for a suitable $e$, by evaluating the same determinant
polynomial at $a$ and $t$. For any effective coefficient $b$,
write $x=a^eb\in R$ and $y=t^eb\in R$. Then $t^ex=a^ey$.
Regularity of $t^e$ modulo $a^e$ gives $x=a^ez$ for $z\in R$.
Rationalization at regular elements therefore sends $b$ to $z$.
This defines a ring retraction of the effective coefficient ring
to $R$. Under tensor injectivity that inclusion is epic, so the
retraction forces equality. The converse is ordinary polynomials.
\end{proof}

\subsection{A principal maximal ideal, valuation bases and binomial bases}

\begin{theorem}\label{thm:principal-maximal-coefficients}
Let $(R,\mathfrak m)$ be local with finite residue field and
$\mathfrak m=\pi R$. No regularity of $\pi$ is assumed. Put
$C=\{s\in S:\pi^ns\in R\text{ for some }n\}$.
If $\pi$ is nilpotent, unique polynomial co-operations occur
exactly when $S=R$. Otherwise they occur exactly when
\[
 C=R\quad\text{or}\quad
 \pi\text{ is regular on }R\text{ and }C=R[1/\pi]
 \text{ canonically inside }S.
\]
In every positive case all tensor maps are isomorphisms. If $\pi$
is regular on $R$, the criterion is equivalently
$\operatorname{Ann}_S(\pi)=0$.
\end{theorem}
\begin{proof}
If $\pi$ is nilpotent, the finite residue field and the filtration
by its powers make $R$ finite, so the finite-base theorem applies.
Assume otherwise and set $P=\bigcap_{n\ge0}\pi^nR$.
An element outside $P$ is a finite power of $\pi$ times a unit.
Consequently $\operatorname{Ann}_R(\pi)\subset P$, and
$P=\pi P$: if $x=\pi y\in P$, comparison with expressions
$x=\pi^{n+1}z$ gives $y\in\pi^nR+\operatorname{Ann}_R(\pi)$
for every $n$. It follows that $R/P$ is a DVR with parameter
$\bar\pi$. Indeed every nonzero element has a finite valuation,
products have the sum of valuations, and an ideal is generated
by an element of least valuation.

Finite grids of distinct powers of $\pi$ show that every coefficient
of an integer-valued polynomial lies in $C$. Conversely a power
of a monic polynomial null modulo $\pi R$ places every element of
$C$ in the effective coefficient ring. The ideals satisfy $PC=P$:
if $\pi^nc=a$ and $x\in P$, divide $x$ by arbitrarily high powers
of $\pi$ inside $P$ to obtain $xc\in P$. Reduction gives, in all
arities,
\[
 \Omega_n(R,C)/P\Omega_n(R,C)
       \simeq\Omega_n(R/P,C/P),\qquad
 P\Omega_n(R,C)=P[X_1,\ldots,X_n].
\]
Surjectivity follows by lifting coefficients; the preimage of
$R/P$ in $C$ is $R$. The kernel formula follows from $PC=P$.

For completeness the tensor kernel has the same elementary
description:
\[
 P\Omega(R,C)^{\otimes_Rn}\simeq P[X_1,\ldots,X_n].
\]
Given $h\in P$ times a pure tensor, write $h=\pi^Nh'$ with
$h'\in P$ and take $N$ large enough to clear the finitely many
coefficients of every factor, distributing this power among the
factors. Balancing reduces the expression to ordinary monomials
with coefficients in $P$. Conversely the polynomial coefficient
map on those ordinary monomials is injective. This proves the
displayed identification and shows that tensor injectivity and
surjectivity hold upstairs exactly when they hold after reduction.
Existence of either substituted tensor also lifts: correct any
lift by its ordinary polynomial error with coefficients in $P$.

The arbitrary-coefficient DVR case of the Noetherian theorem now
says $C/P=R/P$ or $C/P=\operatorname{Frac}(R/P)$ canonically.
The first case gives $C=R$. In the second, lift an inverse of
$\bar\pi$ to get $\pi c=1+p$, $p\in P$. Since $1+p$ is a unit
of $R$, $\pi$ is invertible in $C$, hence regular on the embedded
ring $R$, and $C=R[1/\pi]$. These prove necessity and sufficiency.
When $\pi$ is regular on $S$, $C$ embeds in $R[1/\pi]$; any
element outside $R$ has a numerator equal to a unit times a
finite power of $\pi$, so it generates $1/\pi$. Thus exactly
the two positive alternatives occur. Conversely either alternative
has no $\pi$-torsion, and every $\pi$-torsion element of $S$ lies
in $C$. This proves the last equivalence.
\end{proof}

\begin{corollary}\label{cor:valuation-arbitrary-coefficients}
Let $V$ be a valuation domain that is not a field. For any inclusion
$V\subset S$, unique polynomial co-operations occur exactly when
the residue field is infinite or
$\operatorname{Ann}_S(\mathfrak m_V)=0$. All positive cases have
every tensor isomorphism.
\end{corollary}
\begin{proof}
Infinite residue fields give ordinary polynomials. A principal
maximal ideal is covered by the preceding theorem, since its
generator is regular on $V$. If the maximal ideal is nonprincipal,
the ordered value group has no least positive element. Repeatedly
choosing a smaller positive value gives, for every $\epsilon>0$
and $N$, a $\delta>0$ with $N\delta<\epsilon$.
Thus $\mathfrak m_V$ is generated by its $N$th powers for every
$N$, and the power-generated criterion applies.
Here $(V:_S\mathfrak m_V)=V\oplus\operatorname{Ann}_S(\mathfrak m_V)$.
Indeed rationalization sends this colon into $K$: multiply by any
nonzero element of the maximal ideal. Its rational image has
nonnegative valuation, since arbitrarily small positive values
exclude a negative one. Subtract that element of $V$; the remainder
is killed by the whole maximal ideal. The sum is direct because
$V$ is a domain. This proves the equivalence.
\end{proof}

\begin{corollary}\label{cor:valuation-free-newton-basis}
Let $V$ be an infinite valuation domain with fraction field $K$,
and let $V\subset S$ be a torsion-free extension. Identify
$L=S\cap K$ inside $K\otimes_VS$. If the residue field is infinite,
or the maximal ideal is nonprincipal, then $\Omega(V,S)=V[X]$.
If the nonzero maximal ideal is $\pi V$ and the residue field has
$q$ elements, then
\[
 \Omega(V,S)=
 \begin{cases}
 V[X],&L=V,\\
 \operatorname{Int}(V),&L\ne V,
 \end{cases}
\]
and $\operatorname{Int}(V)$ has the free Newton basis described
below. Thus all these modules are free and every multivariable
polynomial tensor map is an isomorphism, in arbitrary valuation
rank.
\end{corollary}
\begin{proof}
Rationalization reduces the coefficient extension to $L$.
If $V$ is a field, or its residue field is infinite, Vandermonde
interpolation with pairwise unit differences gives $V[X]$.
Suppose its maximal ideal is nonprincipal. For every positive
value $\gamma$ and integer $N>0$ one can choose a positive value
$\delta$ with $N\delta<\gamma$. Indeed any positive value admits
a smaller one, and replacing $0<\xi<\epsilon$ by the smaller of
$\xi$ and $\epsilon-\xi$ gives a positive value at most
$\epsilon/2$ in the ordered-group sense. Finite iteration proves
the assertion without an Archimedean hypothesis.

If a coefficient of a degree-$d$ integer-valued polynomial had
value $-\gamma<0$, choose $a\in V$ with positive value $\delta$
such that $d(d+1)(d+2)\delta/6<\gamma$. The Vandermonde determinant
on $0,a,a^2,\ldots,a^d$ has exactly this smaller value: every
difference of positive powers has the smaller exponent's value.
Its adjugate bounds the negative of each coefficient value by
that determinant value, a contradiction. Hence
$\operatorname{Int}(V)=V[X]$.

In the remaining case choose representatives
$c_0=0,c_1,\ldots,c_{q-1}$ modulo $\pi V$ and, for
$n=\sum_jn_jq^j$, set $a_n=\sum_jc_{n_j}\pi^j$. Define
\[
 w_q(n)=\sum_{j\ge1}\lfloor n/q^j\rfloor,
 \qquad
 b_n(X)=\frac{\prod_{i<n}(X-a_i)}{\prod_{i<n}(a_n-a_i)}.
\]
First differing digits show that the denominator has value
$w_q(n)v(\pi)$. Every element of $V$ has an arbitrarily long
finite string of residue digits. At least $\lfloor n/q^j\rfloor$
of the first $n$ nodes share its first $j$ digits. Counting these
finitely many contributions shows that the numerator has at
least the denominator's value. This remains valid if a difference
has value greater than every integer multiple of $v(\pi)$.
Thus $b_n$ is integer-valued, vanishes at earlier nodes, and has
value $1$ at $a_n$. Successive evaluation gives
$\operatorname{Int}(V)=\bigoplus_{n\ge0}Vb_n$.

Each $b_n$ has coefficients in $V[1/\pi]$. If $b\in L\setminus V$,
the least positive value $v(\pi)$ satisfies $-v(b)\ge v(\pi)$,
so $(\pi b)^{-1}\in V$ and $\pi^{-1}=(\pi b)^{-1}b\in L$.
Consequently the entire basis is available over every proper
overring $L$. Iterating the triangular evaluation argument in
each variable gives the tensor product basis, proving all the
asserted tensor isomorphisms.
\end{proof}

\begin{corollary}\label{cor:binomial-arbitrary-coefficients}
Let $R$ be a nonzero binomial ring. For every inclusion $R\subset S$,
unique polynomial co-operations occur if and only if $S$ is
$\mathbb Z$-torsion-free. All positive cases have all multivariable
tensor isomorphisms. In particular this applies to
$\operatorname{Int}(\mathbb Z)$ and to
$\mathbb Z\oplus\mathbb Q\epsilon$, $\epsilon^2=0$.
\end{corollary}
\begin{proof}
Put $(X)_n=X(X-1)\cdots(X-n+1)$. The integral Stirling identities
give, without canceling integers in $S$,
\[
 F(r)=\sum_{n=0}^{\deg F}\Delta^nF(0)\binom rn,
 \qquad
 \Omega(R,S)=\bigoplus_{n\ge0}(R:_S n!)(X)_n.
\]
In several variables the coefficient ideal is
$(R:_S\prod_i n_i!)$. Preservation in a fixed degree is a finite
test on the integer grid, so these modules commute with arbitrary
localization. Localizations of $R$ remain binomial: degree induction
with $s^dF(X)-F(sX)$ proves preservation on localized inputs,
applied to $F=\binom Xn$ in $\mathbb Q\otimes R$.

At a local base of residue characteristic zero every integer is
a unit, giving ordinary polynomials. At residue characteristic
$p$, the identity $a^p-a\in pR$ makes the local ring $R/pR$
equal to $\mathbb F_p$: every element of its maximal ideal
satisfies $a(1-a^{p-1})=0$ with the second factor a unit.
Thus the maximal ideal is $pR$, and $p$ is regular on $R$.
The principal-maximal theorem gives precisely the condition
$\operatorname{Ann}_{S_{\mathfrak m}}(p)=0$.
These local conditions are equivalent to integer torsion-freeness
of $S$. A nonzero integer-torsion element supplies a nonzero
$p$-torsion element for some prime $p$; localize at a maximal ideal
containing its annihilator to detect it. Conversely regularity of
each integer persists under localization. The local tensor
isomorphisms glue by testing kernels and cokernels.

The ring $\operatorname{Int}(\mathbb Z)$ is binomial because
$\binom{f(n)}j$ is an integer for every integer $n$. For
$r=a+b\epsilon\in\mathbb Z\oplus\mathbb Q\epsilon$,
$\binom rj=\binom aj+b(\binom Xj)'(a)\epsilon$ lies in the same
ring. Its square-zero ideal is not finitely generated over the
base, so this example is non-Noetherian as well as nonreduced.
\end{proof}

\subsection{A bounded nilpotent maximal ideal}

\begin{theorem}\label{thm:bounded-p-nilpotence}
Let $R$ be a local $\mathbb F_p$-algebra with finite residue field
$k=\mathbb F_q$, and suppose $x^p=0$ for every $x\in\mathfrak m$.
For any inclusion $R\subset S$, unique polynomial co-operations
are equivalent to the two conditions
\[
                 \mu:S\otimes_R S\longrightarrow S
                       \text{ is an isomorphism},
                 \qquad(R:_S\mathfrak m)=R.
\]
Here $\mu(s\otimes t)=st$ is the multiplication map.
All positive cases have every tensor isomorphism. If
$\mathfrak m$ is nilpotent, these conditions force $S=R$.
\end{theorem}
\begin{proof}
Frobenius identifies its image $R^p$ with the residue field,
giving a coefficient field. Put $Q=(X^q-X)^p$, a null polynomial
on $R$, and choose the Hermite basis $h_{a,j}$, $a\in k$, $0\le j<p$,
dual to the Hasse derivatives of order below $p$ at these points.
Write $C_j=(R:_S\mathfrak m^j)$. Monic division and Hermite
interpolation give
\[
                     \Omega(R,S)=\bigoplus_{a,j}C_jh_{a,j}
                                      \oplus Q S[X].
\]
To check the value condition, evaluate at $a+\lambda n$ for
$\lambda\in\mathbb F_p$ and $n\in\mathfrak m$. The unit
Vandermonde matrix separates the coefficients of $n^j$.
Polarization, using that $j!$ is a unit for $j<p$, shows that
these powers generate $\mathfrak m^j$. This proves the formula.

Every element of $S$ is a leading coefficient of $cQ$, so tensor
injectivity makes $R\to S$ epic. For $c\in C_1$ and $x\in\mathfrak m$,
$(xc)^p=0$, whence $xc\in\mathfrak m$. Thus $\mathfrak mC_1=\mathfrak m$
and $C_1$ is a ring. The split Hermite summand $C_1h_{a,1}$ makes
$C_1\otimes_R C_1\to C_1$ injective under tensor injectivity.
Reduction of this epimorphism modulo $\mathfrak m$ gives the
identity residue field, forcing $C_1=R$.

Conversely $C_1=R$ implies all $C_j=R$ by iterated colons.
In several variables, mixed differences, the same scalar
interpolation and polarization put every low Hermite coefficient
in $(R:_S\mathfrak m^{\sum j_i})=R$. A term divisible by $Q$ in
at least one variable has arbitrary coefficient in $S$. The tensor
power of the displayed splitting has exactly these coefficient
modules because $S\otimes_R S=S$, proving every tensor
isomorphism. If $\mathfrak m^N=0$, the iterated colon identity
gives $S=(R:_S\mathfrak m^N)=R$.
\end{proof}



\section{Conductor ideals in valuation domains}

The following statement includes valuation groups of arbitrary rank and
does not require a coefficient field in the base ring.

\begin{theorem}\label{thm:valuation-conductor-arbitrary-coefficients}
Let $V$ be a valuation domain with fraction field $K$ and infinite
residue field $l$. Suppose that $R\subset V$ contains a nonzero proper
$V$-ideal $J$ and that $R/J$ is a finite field. For every inclusion of
commutative rings $R\subset S$, the following conditions are equivalent:
\[
\begin{gathered}
 \text{the standard coaddition and comultiplication tensors exist uniquely},\\
 \theta_2\text{ is injective},\qquad (R:_S J)=R,\\
 \operatorname{Ann}_{S/R}(J)=0,\qquad \Omega(R,S)=R[X].
\end{gathered}
\]
In these cases all multivariable tensor maps are isomorphisms, and the
compatible plethory is the ordinary polynomial plethory.
\end{theorem}

We first establish the valuation-ideal calculation used in the proof.
It is a form of the classical Clifford regularity of ideals in a
valuation domain; see Halter-Koch~\cite[Section~5, Lemmas~5.1--5.3 and
Theorem~5.4(1)]{HalterKoch2009}.

\begin{lemma}\label{lem:valuation-ideal-division}
For every nonzero ideal $J$ of a valuation domain and every $N\ge1$,
\[
                         J^N(J:_KJ^N)=J.
\]
\end{lemma}
\begin{proof}
Put $U=(J:J)$, a valuation overring, and let $\mathfrak p$ be its
maximal ideal. If $J=aU$, then $J=J^2a^{-1}U$. Otherwise
$J\mathfrak p=J$: an element $a\in J\setminus J\mathfrak p$ would
generate $J$, because $b\notin aU$ implies $a/b\in\mathfrak p$.
In the nonprincipal case, $J(U:J)=\mathfrak p$. This product is proper,
since equality with $U$ would make $J$ invertible and therefore
principal over the local valuation ring. Conversely, for
$z\in\mathfrak p$, equality $zJ=J$ would imply $z^{-1}\in U$.
Choose $a\in J\setminus zJ$. Comparability gives $zJ\subset aU$,
so $z/a\in(U:J)$ and $z\in J(U:J)$.
Thus in either case $J=J^2A$ for a fractional ideal $A$. Iterating
gives $J=J^NA^{N-1}$, and $A^{N-1}\subset(J:J^N)$. This proves one
inclusion; the other is the definition of the ideal quotient.
\end{proof}

\begin{proof}[Proof of Theorem~\ref{thm:valuation-conductor-arbitrary-coefficients}]
The ring $R$ is local with maximal ideal $J$. Indeed $R/J$ injects
into $V/\mathfrak m_V$, and the inverse of $1+j$, for $j\in J$, lies
in $1+J$. Write $q=|R/J|$.

We first dispose of the idempotent case. If $J^2=J$, then for every
$N$ the ideal $J$ is generated as an $R$-ideal by the $N$th powers
of its elements. To see this, write $0\ne r\in J=J^{N+1}$ as a
finite sum of products and choose a product of value at most $v(r)$.
For its factor $y$ of smallest value, $r/y^N$ belongs to $J$.
For such a local ring the Vandermonde ideal $D_d(R)$ is $R$ when
$d<q$ and is $J$ when $d\ge q$. The first assertion uses distinct
residue classes. For the second, all determinants belong to $J$;
evaluating the determinant polynomial on $0,1,a,a^2,\ldots$ gives
$a^{d(d-1)(d+1)/6}$ times a unit for $a\in J$, and the generation
by powers gives the reverse inclusion. Consequently the effective
coefficient ring is $C=(R:_S J)$. If $R\to C$ is epic, then $JC$ is
an ideal of $R$ containing $J$. It cannot equal $R$: an equation
$1=\sum a_ic_i$ would make $J$ finitely generated, and $J=J^2$ would
contradict Nakayama's lemma. Hence $JC=J$. Reducing the epimorphism
modulo $J$ gives the injective field epimorphism $R/J\to C/J$, so
$C=R$. The coefficient-ring obstruction already proved then gives
the theorem in this case.

Assume now $J^2\ne J$ and fix $0\ne c\in J\setminus J^2$. Set
$W=V[1/c]=R[1/c]$. For every $0\ne y\in J$ one has
$2v(y)>v(c)$; comparison of the factors in a product therefore gives
\begin{equation}\label{eq:conductor-power-control}
                        J^{2h+1}\subset c^hJ\quad(h\ge0).
\end{equation}
The ideals $D_d(R)$ are cofinal with the powers of $J$. Many nodes
in the $q$ residue classes force arbitrarily many differences into
$J$, proving one direction. For the other, choose $d+1$ elements
$u_i\in V$ with distinct residues. The determinant on $cu_i$ is
$\delta=c^hu$ with $h=d(d+1)/2$ and $u\in V^\times$; then
$J^{2h+1}\subset c^hJ=\delta J\subset\delta R$.
Thus the effective coefficient ring is
\[
                            C=\bigcup_{n\ge1}(R:_S J^n).
\]
Iterating colons shows that $(R:_S J)=R$ implies $C=R$, which proves
all positive conclusions. Suppose for the rest of the proof that
$\theta_2$ is injective. The coefficient-ring obstruction makes
$R\to C$ an epimorphism.

Rationalization has image $L\subset W$ and kernel $T$ annihilated
elementwise by powers of $J$. If every $b\in L$ has
$v(b)\ge-v(c)$, then $J^3L\subset J$ is a common ideal containing
$J^3$. The injective epimorphism
$R/J^3L\to L/J^3L$ is an isomorphism: reduction by the nilpotent
maximal ideal gives an epimorphism from a field, and finite iteration
of that ideal proves surjectivity. Hence $L=R$. Otherwise choose
$b\in L$ with $v(b)<-v(c)$. Since $b^{-1}V\subset J$, one obtains
$V\subset L$ and $c^{-1}=b(cb)^{-1}\in L$, whence $L=W$.
The case $L=R$ retracts the epimorphism $R\to C$ and gives $C=R$.

In the remaining case lift $c^{-1}$ to $a\in C$ and put $z=ca$.
Some $c^n$ annihilates $z-1$, so $z^n(z-1)=0$. Therefore $e=z^n$
is idempotent and $e=z^N$ for all $N\ge n$. It annihilates $T$,
and rationalization identifies $eC$ with $W$. The element $1-e$ is
killed by some $J^M$. The same nilpotent-ideal argument on the other
factor gives
\begin{equation}\label{eq:valuation-coefficient-splitting}
                 C=W\times R/I,\qquad J^M\subset I\subset R,
\end{equation}
with the diagonal inclusion of $R$; $I=R$ is allowed.

Choose representatives $B\subset R$ of $R/J$ and put
$P(X)=\prod_{b\in B}(X-b)$. Differences of representatives are
units, $P(R)\subset J$, and $P'(b)$ is a unit for every $b\in B$.
The monic basis $X^iP^j$, $0\le i<q$, gives the exact decomposition
\begin{equation}\label{eq:valuation-conductor-coefficients}
 \Omega(R,K)=\bigoplus_{i<q}RX^i\ \oplus\!
 \bigoplus_{\substack{i<q\\j\ge1}}(J:_KJ^j)X^iP^j.
\end{equation}
Here is the necessity argument, including the boundary of any
valuation cut. Subtract interpolation on $B$, leaving only $j\ge1$.
For fixed $0\ne y\in J$, divide the substitution $X=b+yZ$ by a
coefficient of minimum value among the $a_{ij}y^j$. Reduction gives
\[
 \sum_{j\ge1}\left(\sum_{i<q}\overline{a_{ij}y^j/\mu}\,\bar b^i\right)
                         (\overline{P'(b)}Z)^j.
\]
For some $b$ this is a nonzero polynomial, by interpolation on the
finite field, and it has zero constant term. A negative minimum
contradicts $V$-valuedness. A zero minimum would give a nonconstant
polynomial from the infinite field $l$ into a finite set. A positive
minimum outside $v(J)$ gives a value in $R\cap\mathfrak m_V$ outside
$J$. Thus $a_{ij}y^j\in J$. Products of $j$ elements of $J$ are
divisible by the $j$th power of their least-valued factor, proving
$a_{ij}J^j\subset J$. Sufficiency is immediate. All the coefficient
ideals displayed above lie in $W$, since $c^ja_{ij}\in J$.

For $F$ vanishing on $B$, define the ideal
\[
              \nu(F)=\sum_{i<q,j\ge1}a_{ij}J^j.
\]
If $F(b+Y)=\sum_{h\ge1}g_{bh}Y^h$, the same interpolation calculation
at each fixed $y\in J$ proves
\begin{equation}\label{eq:valuation-ideal-norm}
                         \nu(F)=\sum_{b\in B,h\ge1}g_{bh}J^h.
\end{equation}
Indeed the two sides have equal least coefficient values after
replacing every $J^h$ by $y^hV$; the union over $y$ computes the
powers of a valuation ideal. No numerical infimum is used.

If $I=R$ in \eqref{eq:valuation-coefficient-splitting}, the coefficient
module of $P$ in \eqref{eq:valuation-conductor-coefficients} is
$U=(J:J)$, with $JU=J$. Choose $\beta\in V$ whose residue is not in
$R/J$. The tensor $\beta P\otimes P-P\otimes\beta P$ has zero
polynomial image. Projection to $U\otimes_RU$ and reduction modulo
$J$ give $\bar\beta\otimes1-1\otimes\bar\beta\ne0$ over the field
$R/J$. Hence this case contradicts injectivity.

Suppose $I\subsetneq R$. Put
\[
\begin{aligned}
 H&=\{F\in W[X]:F(R)\subset I\},& E_0&=R[X]+H,\\
 H_0&=\{F\in H:F|_B=0\},& K_0&=\sum_{F\in H_0}\nu(F).
\end{aligned}
\]
The first projection of $\Omega(R,C)$ is $E_0$, by lifting its
second component to $R[X]$. Interpolation on $B$ maps $H$ to $H_0$,
and the polynomials $aP$, $a\in J^{M-1}$, show
$J^M\subset K_0\subset J$.

We claim $H\not\subset JE_0$. Otherwise, for $F\in H_0$ and
$N\ge1$, substitute
$\psi=b+\lambda P^N$, where $b\in B$ and $\lambda\in(J:J^N)$.
This preserves $R$, and $F\circ\psi\in H_0$ has only $P$-degrees
at least $N$. Lemma~\ref{lem:valuation-ideal-division}, pure-power
generation of powers of valuation ideals, and
\eqref{eq:valuation-ideal-norm} give
\[
 \sum_{b,\lambda}\nu(F\circ\psi)
   =\sum_{b,h}g_{bh}(J:J^N)^hJ^{Nh}
   =\sum_{b,h}g_{bh}J^h=\nu(F).
\]
Write $F\circ\psi=A+\sum r_tF_t$, with $A\in JR[X]$, $r_t\in J$,
and $F_t\in H$. Subtract their interpolants to arrange $F_t\in H_0$.
Projection to $P$-degrees at least $N$ gives
$K_0\subset J^{N+1}+JK_0$. Taking $N=M$ yields $K_0=JK_0$.
Iteration would imply $J^M\subset K_0\subset J^{M+1}$.
But $J^M=J^{M+1}$ forces $J=J^2$: express $r^M\in J^{2M}$ as a
finite sum of products and choose a least-valued factor $y$ of a
product of value at most $Mv(r)$; then $2v(y)\le v(r)$ and
$r\in y^2V$. This contradicts our choice of $c$, proving the claim.

Choose $F\in H\setminus JE_0$. The two integer-valued polynomials
$f=(F,0)$ and $g=(0,P^M\bmod I)$ survive modulo $J\Omega(R,C)$:
use the first projection for $f$ and the monic coefficient of the
second projection for $g$. Their tensor remains nonzero over the
residue field, while its polynomial image is zero because their
coefficient supports are disjoint. This is the final contradiction.
Hence injectivity forces $C=R$. Iterated colons and the universal
coefficient description finish every asserted equivalence.
\end{proof}

\begin{corollary}\label{cor:valuation-conductor-specializations}
The theorem applies, with arbitrary coefficient torsion, to
$R=\mathbb F_q+t^al[[t]]$ for any infinite extension
$l/\mathbb F_q$ and $a\ge1$; to $R=\mathbb F_q+xV$ for an arbitrary
valuation domain of characteristic $p$, infinite residue field, and
$0\ne x\in\mathfrak m_V$; and to $R=\mathbb F_q+J$ for every
nonzero proper ideal in such a valuation domain. In particular no
attained or rational cut is required for a real-valued valuation,
and no rank restriction is required in the general case.
\end{corollary}



\section{Explicit existence and uniqueness obstructions}
The following examples separate failure of existence from failure of
uniqueness, and show why the coefficient-ring condition alone is insufficient.

\begin{lemma}\label{lem:finite-conductor-coefficient-obstruction}
Let $R\subset S$ be a finite injective extension, with $R$ nonzero
and local. If the conductor $J=(R:S)$ has finite quotient $R/J$,
then the binary polynomial tensor map is injective exactly when
$S=R$.
\end{lemma}
\begin{proof}
The monic product $P$ over representatives of $R/J$ has
$P(R)\subset J$, and monic division gives
$\Omega(R,S)=E\oplus PS[X]$, where $E$ has degree below $\deg P$.
In particular $SP$ is a direct summand. Restriction to its tensor
square forces $S\otimes_RS\to S$ to be injective, hence an
isomorphism. A finite ring epimorphism is surjective
\cite[Lemma~10.107.6]{StacksRingEpimorphisms}, so injectivity of
$R\to S$ gives $R=S$. The converse is immediate.
\end{proof}

\begin{lemma}\label{lem:binary-mixed-difference-obstruction}
Let $(R,\mathfrak m)$ be a local domain. Suppose that $u_n\in R$
has the following property: for every finite family
$\mathcal H\subset\operatorname{Int}(R)$, all sufficiently large
$n$ satisfy $h(u_n)-h(0)\in\mathfrak m$ for every $h\in\mathcal H$.
If $g\in\operatorname{Int}(R)$ has $g(0)=0$ and
$g(u_n^2)\notin\mathfrak m^2$ for all $n$, then $g(XY)$ is not a
finite sum of products $h_i(X)k_i(Y)$ with integer-valued factors.
\end{lemma}
\begin{proof}
The mixed difference at $(0,0),(u_n,0),(0,u_n),(u_n,u_n)$ of such
an expansion would be
$\sum_i(h_i(u_n)-h_i(0))(k_i(u_n)-k_i(0))\in\mathfrak m^2$.
For $g(XY)$ the same difference is $g(u_n^2)$, a contradiction.
\end{proof}

\begin{example}\label{ex:all-finite-fields-two-branches}
For every finite field $k=\mathbb F_q$, let
\[
 T=k[t]_{k[t]\setminus((t)\cup(t-1))},\quad
 m=t(t-1),\quad R=k+mT,\quad K=k(t).
\]
Then $R$ is a one-dimensional Noetherian local domain with finite
normalization $T$. Put
\[
 a(X)=\frac{X^q-X}{m},\qquad g(X)=a(X)^q-a(X),\qquad
 u_n=t(t-1)^{n+1}.
\]
The polynomial $g$ is integer-valued on $R$, has degree $q^2$,
and $g(XY)$ has no multiplicative tensor expansion.
\end{example}
\begin{proof}
Let $A=k[m]_{(m)}$. The equation $t^2=t+m$ gives
$T=A\oplus At$: every inverted polynomial $f(t)$ has
$f(t)f(1-t)\in k[m]$ with nonzero constant term, and its
inverse is therefore already in $A+At$.
Also $R=A\oplus A(mt)$ with
$(mt)^2=m(mt)+m^3$. Hence $R$ is finite over the DVR $A$ and
is Noetherian of dimension one. Its maximal ideal is $mT$ and
its residue field is $k$. Its fraction field is $K$, and the
normal finite overring $T$ is its normalization.

For $r\in R$, $r^q-r\in mT$, so $a(r)\in T$; reduction at
both maximal ideals of $T$ gives $g(r)\in mT\subset R$.
Write $v_0,v_1$ for the two normalized valuations. One has
\[
 v_0(u_n)=1,\quad v_1(u_n)=n+1,\quad
 v_0(g(u_n^2))=1,\quad v_1(g(u_n^2))=2n+1.
\]
For the last two equalities, $u_n^2$ has positive value at both
branches, so the leading term of $x^q-x$ is $-x$; division by
$m$ subtracts one from the value, and the same observation applies
to $a^q-a$. For finitely many polynomials, every nonconstant term
of $h(u_n)-h(0)$ tends to zero at $v_1$, while the difference lies
in $R$. It therefore belongs to $(t-1)T\cap R=mT$ eventually.
Since $(mT)^2$ has $v_0$ at least two, the preceding lemma applies.
\end{proof}

\begin{example}\label{ex:two-dimensional-split-tensor-kernel}
Let $T=\mathbb F_2[x,y]_{(x,y)}$, $\mathfrak n=(x,y)T$, and
$R=\mathbb F_2+\mathfrak n^2$. This is a two-dimensional
Noetherian local domain with finite normalization $T$. For
$q(X)=X^2-X$,
\[
 \operatorname{Int}(R)=R\oplus RX\oplus q(X)T[X],\qquad
              (xq)\otimes q-q\otimes(xq)\ne0,
\]
although the displayed tensor has zero polynomial image.
\end{example}
\begin{proof}
Over $A=\mathbb F_2[x^2,y^2]_{(x^2,y^2)}$, the ring $T$ is finite
with generators $1,x,y,xy$, and $R$ is generated by
$1,xy,x^3,xy^2,x^2y,y^3$. Thus both are Noetherian, and integrality
gives dimension two. The maximal ideal of $R$ is $\mathfrak n^2$;
its fraction field equals that of $T$, since $x=x^3/x^2$ and
$y=y^3/y^2$. The normal ring $T$ is therefore the finite
normalization.

At each height-one prime of $T$, the image of $R$ is infinite,
so Vandermonde interpolation puts every integer-valued coefficient
in that localization. Normality gives coefficients in $T$.
Since $q(R)\subset\mathfrak n^2$ and this is a $T$-ideal, monic
division by $q$ reduces the value condition to a linear remainder;
evaluation at $0,1$ puts its coefficients in $R$. This proves the
decomposition. In its split summand $Tq\otimes_RTq$, reduction
modulo $\mathfrak n^2$ gives $D\otimes_{\mathbb F_2}D$, where
$D=\mathbb F_2[x,y]/(x,y)^2$. The tensor $x\otimes1-1\otimes x$
is nonzero in its nine-dimensional monomial basis and maps to
zero under multiplication. This detects the displayed tensor.
\end{proof}

\begin{lemma}\label{lem:derivation-detects-polynomial-torsion}
Let $R$ be a domain with fraction field $K$, let $V\subset K$ be
a DVR containing $R$, and let $D:K\to K$ be a derivation with
$D(R)\subset V$. Suppose $u_n\in R$ satisfies $v(u_n)\to\infty$
and $D(u_n)=1$. Then $F'(0)\in V$ for every
$F\in\operatorname{Int}(R)$, and
\[
       \beta(F,G)=D(F'(0))G'(0)+V\quad\text{in }K/V
\]
is balanced over $R$. If $P\in R[X]$, $cP\in\operatorname{Int}(R)$,
and $D(c)P'(0)^2\notin V$, then
$(cP)\otimes P-P\otimes(cP)$ is nonzero and has zero polynomial
image.
\end{lemma}
\begin{proof}
The chain rule gives
$D(F(u_n))=(DF)(u_n)+F'(u_n)$. Its left side lies in $V$;
the constant coefficient of $DF$ is $D(F(0))\in V$, and its other
terms eventually have nonnegative valuation. Thus $F'(u_n)\in V$,
and the same estimate for $F'(u_n)-F'(0)$ gives $F'(0)\in V$.
For $r\in R$, the balancing error is
$D(r)F'(0)G'(0)\in V$. Finally evaluation on the displayed tensor
cancels the two $D(P'(0))$ terms and leaves
$D(c)P'(0)^2+V\ne0$. Its polynomial image cancels identically.
\end{proof}

\begin{example}\label{ex:akizuki-reid-polynomial-kernel}
Let $A_0$ be a DVR with parameter $t$ and finite residue field
$\mathbb F_q$. Choose units $a_i\in A_0$ and integers
$n_0=0$, $n_{i+1}\ge2n_i+2$, such that
$z=\sum_i a_it^{n_i}\in\widehat A_0$ is transcendental over
$\operatorname{Frac}(A_0)$. Put
\[
 z_i=\frac{z-\sum_{j<i}a_jt^{n_j}}{t^{n_i}},\quad
 C=A_0[t(z-a_0),(z_i-a_i)^2:i\ge0],\quad
 M=(t,t(z-a_0))C,\quad R=C_M.
\]
Then $R$ is a one-dimensional Noetherian local domain whose
normalization is a DVR with parameter $t$ and the same residue
field, but is not finite over $R$. For $c=z-a_0$ and $P=X^q-X$,
\[
 cP\in\operatorname{Int}(R),\qquad
              (cP)\otimes P-P\otimes(cP)\ne0
\]
with zero polynomial image. This holds in mixed characteristic
as well as positive characteristic.
\end{example}
\begin{proof}
The ring assertions are Reid's two theorems on page~2, under
precisely these hypotheses; their proofs occupy pages~2--4
\cite{Reid1995Akizuki}. We give the obstruction and a separate
nonfiniteness argument. On $K=\operatorname{Frac}(A_0)(z)$ set
$D=t^{-1}\partial/\partial z$. It kills $A_0$, sends
$t(z-a_0)$ to $1$, and
\[
              D((z_i-a_i)^2)=2t^{-n_i-1}(z_i-a_i)
\]
has value at least $n_{i+1}-2n_i-1\ge1$ (or is zero).
The quotient rule gives $D(R)\subset V$. The elements
$u_i=t(z-\sum_{j\le i}a_jt^{n_j})\in R$ have derivative $1$
and value $n_{i+1}+1\to\infty$.

Writing $u=t(z-a_0)$, one has $ct=u$ and $cu=t(z-a_0)^2\in C$,
so $cM R\subset R$. Since $P(R)\subset MR$, $cP$ is
integer-valued. The preceding lemma detects its tensor by
$D(c)P'(0)^2=1/t\notin V$.
If $V$ were finite over $R$, the derivatives of a finite module
generating set would bound the valuations of all $D(V)$ below,
because $D(R)\subset V$. But $z_i\in V$ and
$D(z_i)=t^{-n_i-1}$, a contradiction.

An entirely explicit characteristic-two instance is
$A_0=\mathbb F_2[t]_{(t)}$, $a_i=1$, $n_i=(i+1)!$ for $i\ge1$.
The gap condition holds. To verify transcendence, a hypothetical
relation $P(t,z)=0$ of $z$-degree $d$ and coefficient degrees at
most $h$ would make its value on the $N$th truncation have
valuation at least $n_{N+1}$ and degree at most $h+dn_N$.
For large $N$ it must vanish identically. Infinitely many distinct
truncations cannot be roots of the same nonzero polynomial over
$\mathbb F_2(t)$, giving the required contradiction.
\end{proof}

\begin{proposition}\label{prop:coefficient-basechange-not-sufficient}
For any $R\subset S$, let $C$ be its effective coefficient ring,
$Q=C/R$, and
\[
 N(R,Q)=\{(q_i)\in\bigoplus_{i\ge0}Q:
                         \sum_i r^iq_i=0\text{ for every }r\in R\}.
\]
Binary tensor injectivity implies
\[
 C\otimes_R\Omega(R,S)\simeq C[X],\qquad
 C\otimes_RN(R,Q)=0.
\]
These two conditions are equivalent even without tensor injectivity.
They are not sufficient, even when $C$ is a flat epic localization
of a domain. In that scope a nonzero tensor can be killed by a
monic polynomial in one variable.
\end{proposition}
\begin{proof}
The coefficient reduction sequence is
$0\to R[X]\to\Omega\to N(R,Q)\to0$: lift a finite vector of
coefficients and use its defining equations for surjectivity.
After tensoring by $C$, the natural map $C[X]\to C\otimes_R\Omega$
is split by multiplication. Its cokernel is $C\otimes_RN(R,Q)$,
proving the equivalence.
If $z=\sum c_i\otimes F_i$ maps to zero in $C[X]$, choose a common
monic $H\in R[X]$ with every $c_iH\in\Omega$, using the monic
coefficient witnesses and their product. The tensor
$\sum c_iH\otimes F_i$ has zero polynomial image. Under tensor
injectivity it is zero, and extracting its first leading
coefficient gives $z=0$.

For the counterexample let $R=k+tV$, $V=l[[t]]$, $K=l((t))$,
where $k=\mathbb F_q\subset l$ and $l$ is infinite. Put $P=X^q-X$.
The valuation-conductor decomposition gives
\[
 \operatorname{Int}(R)=\bigoplus_{i<q}RX^i\ \oplus\!
 \bigoplus_{i<q,j\ge1}t^{1-j}VX^iP^j.
\]
Every element of $K$ occurs as a leading coefficient, so $C=K$
and $K\otimes_R\operatorname{Int}(R)=K[X]$.
Choose $\beta\in l\setminus k$. The tensor
$z=\beta P\otimes P-P\otimes\beta P$ has zero polynomial
image but projects to the nonzero element
$\beta\otimes1-1\otimes\beta\in l\otimes_kl$.
Nevertheless $P^2/t$ is integer-valued and $t,t\beta\in R$, so
\[
 \beta P^2\otimes P=(P^2/t)\otimes(t\beta P)
                         =P^2\otimes\beta P.
\]
Thus $P(X)z=0$.
More generally, for $R=k+t^aV$, every proper rational coefficient
overring fails uniqueness by the valuation-conductor theorem:
an overring inside $V$ lies in $(R:_S t^aV)$, and an overring
outside $V$ is $K$. Indeed, a sufficiently large power of an
element of negative value has inverse in $t^{a+1}V$; multiplying
the conductor by that power puts both $V$ and $t^{-1}$ in the
overring. For $F=P^2/t^a$, the split coefficient module
$t^{-a}V$ detects
$(\beta F)\otimes F-F\otimes(\beta F)\ne0$ after reduction to
$l\otimes_kl$. These rings are non-Noetherian, since
$t^aV/t^{2a}V$ has infinite dimension over $k$.
\end{proof}

\begin{corollary}\label{cor:flat-cusp-nonzero-tor-obstruction}
For the cusp with $a=1$ in the preceding proof, the value module is
\[
 E=\bigoplus_{i<q}(K/R)X^i\ \oplus\!
   \bigoplus_{i<q,j\ge1}(K/t^{1-j}V)X^iP^j.
\]
The universal obstruction $K_2=\operatorname{Tor}_2^R(E,E)$ is
nonzero. Its nonzero class is the double connecting lift of
$\beta\otimes1-1\otimes\beta\in V\otimes_RV$ from the exact
sequence $0\to V\to K\to K/V\to0$.
\end{corollary}
\begin{proof}
The displayed description follows by quotienting the split monic
coefficient modules. The $P$-summand is $K/V$. The indicated tensor
is nonzero by residue reduction and dies in $K\otimes_RV=K$.
Since $K$ is flat, two long exact Tor sequences identify that
kernel with $\operatorname{Tor}_2^R(K/V,K/V)$. The split inclusion
of $K/V$ into $E$ preserves this nonzero class. The sign is the
connecting convention already fixed in the universal theorem;
nonvanishing is independent of that sign.
\end{proof}

\begin{corollary}\label{cor:full-bivariate-flat-obstruction}
Under the flat epic coefficient hypothesis, evaluation modulo
$C/R$ induces a map
$\lambda:\operatorname{Tor}_1^R(E,E)\to\operatorname{App}(R^2,C/R)$
through the chosen connecting map. There is an exact sequence
\[
 0\to\operatorname{Tor}_2^R(E,E)\to\Omega^{\otimes_R2}
                \to\Omega_2\to\ker\lambda\to0.
\]
Thus the full bivariate tensor map is an isomorphism exactly when
$\operatorname{Tor}_2^R(E,E)=0$ and $\lambda$ is injective.
\end{corollary}
\begin{proof}
Evaluation modulo $C/R$ vanishes on the tensor image, so it factors
uniquely through the cokernel and defines $\lambda$. A class of
$\ker\lambda$ lifts through the surjective connecting map to a
polynomial whose values are in $R$, hence to $\Omega_2$.
The stated sequence and criterion follow. This retains the
distinction between the two specified substitution tests and
surjectivity for every bivariate integer-valued polynomial.
\end{proof}


\section{A nonflat positive example}
The next construction shows that all polynomial tensor maps can be
isomorphisms even when the polynomial module is not flat. Its finite
tensor-balancing identity is analogous to the one in Lazard's
epimorphism example \cite[Section 110.36, Tag 06RH]{StacksLazardExample};
the all-arity polynomial calculation is given here.

\begin{theorem}\label{thm:nonflat-forest-plethory}
There is a proper inclusion $R\subset S$ of zero-dimensional local
$\mathbb F_2$-algebras for which every nilpotent has square zero and
$R$ is not Noetherian, such that
\[
          \Omega(R,S)^{\otimes_R n}\xrightarrow{\sim}\Omega_n(R,S)
                         \qquad(n\ge1).
\]
Both $S$ and $\Omega(R,S)$ are nonflat over $R$. The effective
coefficient ring is $S$. Thus flatness is not a necessary condition
for the existence of the unique compatible polynomial plethory.
\end{theorem}
\begin{proof}
Let the vertices $v=(j,w)$ range over countably many full binary
trees, where $j\ge1$ and $w$ is a finite binary word. Denote the
children by $v0,v1$, and define
\[
 S=\mathbb F_2[x_v,y_v:v]/(x_v^2,y_v^2,y_v-x_vy_{v0}y_{v1}:v).
\]
This presentation has an injective direct-limit model. At level
$n$ use the squarefree polynomial algebra on $x_v$ of depth below
$n$ and $y_v$ of depth $n$. Replace a boundary $y_v$ at the next
level by $x_vy_{v0}y_{v1}$. The new triples are pairwise disjoint
and disjoint from the retained variables. Distinct squarefree
monomials therefore have distinct nonzero images. Every finite
expression and every displayed relation occurs at a finite level,
so the direct limit is exactly $S$.

Give $x_v$ degree $-1$ and $y_v$ degree $1$, and take
$R=\bigoplus_{d\le0}S_d$. This is a unital subring, and it is proper
because each nonzero $y_v$ has degree $1$. The augmentation ideals
of both rings consist of square-zero elements: this follows from
the squarefree monomial model and the Frobenius identity. Every
element with constant term $1$ is its own inverse. Hence both
rings are local, with one prime and residue field $\mathbb F_2$.
Write $N$ for the maximal ideal of $R$.

Products of degree-zero monomials $x_vy_{v0}$ in distinct tree
components are nonzero. Consequently $N^k\ne0$ for every $k$.
A finitely generated ideal generated by square-zero elements is
nilpotent, so $N$ is not finitely generated and $R$ is not Noetherian.

The relations also make $R\to S$ epic. In $S\otimes_R S$,
\[
 y_v\otimes1
 =y_{v0}\otimes x_vy_{v1}
 =x_vy_{v0}\otimes y_{v1}
 =1\otimes y_v,
\]
because both $x_vy_{v0}$ and $x_vy_{v1}$ belong to $R$.
The analogous identity for $x_v$ is automatic. These generators
give $S\otimes_R S\simeq S$.

For every $k\ge1$ one has the exact colon identity
\[
                              (R:_S N^k)=R.
\]
Indeed, if $c\notin R$, its nonzero positive-degree part involves
only finitely many components. Multiply it by a product of $k$
degree-zero monomials from fresh components. This product belongs
to $N^k$ and cannot kill that positive-degree part: at a finite
level multiplication across disjoint components is the tensor
product of nonzero vectors over $\mathbb F_2$, and the transition
maps are injective. The product remains outside $R$.

Set $Q(X)=X^4+X^2=X^2(X-1)^2$. For $r=a+n\in R$, with
$a\in\mathbb F_2$ and $n\in N$, one has $Q(r)=0$. Monic division
by $Q$ reduces the value condition to a polynomial $H$ of degree
below $4$. Since
$H(a+n)=H(a)+H'(a)n$, the colon identity says precisely that
\[
                        H(0),H(1),H'(0),H'(1)\in R.
\]
The Chinese remainder theorem for $(X^2)$ and $((X-1)^2)$ identifies
these four coordinates with the four coefficients, over
$\mathbb F_2$ and hence over both coefficient rings. Thus
\begin{equation}\label{eq:forest-module-splitting}
                  \Omega(R,S)=R[X]_{<4}\oplus Q(X)S[X].
\end{equation}

In several variables use the monic basis
$b_i=X^i$ for $i<4$ and $b_i=QX^{i-4}$ for $i\ge4$ in each
coordinate. Terms with a high index vanish on $R^n$ and have
arbitrary $S$-coefficients. For the remaining polynomial, take
finite differences about $a\in\{0,1\}^n$ in a set $J$ of variables
with arbitrary increments $n_j\in N$. The difference equals
$(\partial_JH)(a)\prod_{j\in J}n_j$. Its values lie in $R$;
the products generate $N^{|J|}$, so the colon identity places each
mixed derivative in $R$. The case $J=\varnothing$ is evaluation.
The tensor product of the four-coordinate Hermite isomorphisms
therefore puts all coefficients of the low-index remainder in $R$.

Hence $\Omega_n$ has coefficient module $R$ at every all-low index
and $S$ at every other index. The tensor power of
\eqref{eq:forest-module-splitting} has precisely the same modules,
since $S\otimes_R S=S$. The polynomial tensor map is the identity
on these coefficient modules, proving the asserted isomorphisms.
For reference, the substitutions on $Q$ are
\[
 Q(X+Y)=Q(X)+Q(Y),\qquad
 Q(XY)=Q(X)Q(Y)+Q(X)Y^2+X^2Q(Y).
\]
The all-variable isomorphisms transport every polynomial identity
and composition identity, yielding the unique compatible plethory.

If $S$ were flat, the local inclusion with nonzero residue fibre
would be faithfully flat. A faithfully flat epimorphism is an
isomorphism, contradicting $y_v\notin R$; see
\cite[Lemmas~10.39.16 and 10.107.7]{StacksFaithfulEpimorphisms}.
The split copy of $S$ in \eqref{eq:forest-module-splitting} then
shows that $\Omega(R,S)$ is not flat. Finally $cQ$ is integer-valued
with leading coefficient $c$ for every $c\in S$, so the effective
coefficient ring is all of $S$.
\end{proof}

\begin{proposition}\label{prop:forest-syzygy-invariants}
For the preceding inclusion, put $D=S/R$ and use the modules and
maps of Theorem~\ref{result-20}. Then
\[
                   K_1=K_2=0,\qquad T_1\simeq D^{16}.
\]
The obstruction map records, up to its stipulated sign, the sixteen
low-low coefficients modulo $R$. Both polynomial substitutions have
zero obstruction. No vanishing of $\operatorname{Tor}_q^R(S,S)$
is required for this calculation.
\end{proposition}
\begin{proof}
In the monic basis just used,
\[
 \mathcal F=\bigoplus_{i\ge0}Sb_i,\qquad
 E=\bigoplus_{i<4}Db_i,
\]
and $\varepsilon$ is the quotient $S\to D$ on the four low
summands and zero on the others. Let
$M_q=\operatorname{Tor}_q^R(S,S)$ for $q\ge1$.
The exact sequence $0\to R\to S\to D\to0$ and epicness give
\[
 M_q\simeq\operatorname{Tor}_q^R(S,D)
       \simeq\operatorname{Tor}_q^R(D,S),
\]
as well as
\[
 0\to M_1\to\operatorname{Tor}_1^R(D,D)\to D\to0,
 \qquad \operatorname{Tor}_q^R(D,D)\simeq M_q\quad(q\ge2).
\]
The degree-one row of the universal construction splits by pairs
of basis indices. At a low-low pair it is the diagonal map of
$M_1$, followed by the difference and the displayed injection into
$\operatorname{Tor}_1(D,D)$. Its middle homology is zero and its
cokernel is $D$. At a mixed pair the first map is an isomorphism;
at a high-high pair the other two modules vanish. Hence $K_1=0$
and $T_1=D^{16}$. In degree two the final row map is surjective
at every pair, giving $T_2=K_2=0$.
The cokernel identifies the obstruction with the low-low
coefficients modulo $R$. In the two displayed identities for $Q$,
every term has a high index in at least one coordinate. Multiplying
by the substituted coefficient polynomial preserves this property.
Ordinary $R$-polynomials have low-low coefficients in $R$.
Therefore both substitutions have zero obstruction, as claimed.
\end{proof}


\part{Comparison over real closed fields}
The following arguments combine compact parameter resolution,
finite joins and exact relative PL replacement.  The replacement uses
literal product prisms simultaneously respecting both ends of the
cylinder.  Its underlying finite fragmentation construction is due
to Mn\"ev \cite{Mnev2007}; the marked adaptation and the field-transfer
steps are made explicit below.  Together they prove that finite jointly
PL and finite jointly definable concordance families have the same weak
homotopy type over every o-minimal real closed field, with exact control
of the prescribed parameter boundary.

For context, the ordinary definable homotopy comparison of Baro--Otero
\cite[Theorem~3.1 and Corollary~3.3, pp.~5--7]{BaroOtero2008v1}
concerns homotopies of maps between definable sets.  Here every parameter
value must remain a homeomorphism and the parameter projection must stay
fixed.  Shiota's comparison of PL and definable homotopies and isotopies
\cite[Lemma~3.3, p.~37, and relative isotopy argument, pp.~44--46]{Shiota2010}
supplies related finite-polyhedron results.  The additional relative
family requirements are addressed by the explicit product-prism argument
below.  These comparisons do not constitute a novelty claim for the
result assembled here.
\section{Automorphisms under elementary extension}

We prove that elementary extension preserves the entire weak homotopy
type of marked definable automorphism families.

Let $\mathcal M\preccurlyeq\mathcal N$ be an elementary extension of
o-minimal structures in the same language, both expanding real closed
fields.  Definability over either structure permits all of its
parameters.  For a closed bounded definable set $D\subset M^n$, let
$\mathcal S$ be a finite list of definable subsets of $D$, with each
subset designated either setwise preserved or pointwise fixed.  The
marked subsets need not be closed.  Write
\[
 G_{\mathcal M}(D;\mathcal S)
   =\operatorname{Aut}_{\mathrm{def},\mathcal M}(D;\mathcal S)
\]
for the simplicial group whose $q$-simplices are definable
homeomorphisms of $D\times\Delta^q_M$ preserving the entire parameter
projection and satisfying the designated conditions on
$S\times\Delta^q_M$ for every $S\in\mathcal S$.  Both the total map
and its inverse are jointly continuous.  Faces and degeneracies are
affine parameter pullbacks, and multiplication is composition.

\begin{theorem}\label{thm:elementary-extension-marked-automorphisms}
Extension of defining formulas induces a weak equivalence
\[
 G_{\mathcal M}(D;\mathcal S)
 \longrightarrow
 G_{\mathcal N}(D_{\mathcal N};\mathcal S_{\mathcal N}).
\]
This includes the empty set, for which both groups are trivial.
In particular, for any closed bounded definable manifold $B$ over
$\mathcal M$, it gives a weak equivalence
\[
 P_{\mathrm{def},\mathcal M}(B)
 \longrightarrow P_{\mathrm{def},\mathcal N}(B_{\mathcal N}),
\]
where the bottom of $B\times[0,1]$ is fixed pointwise and its free
top is preserved setwise.

If $B$ is the rational realization of a finite abstract closed PL
manifold, the square with the PL
coefficient-extension map commutes.  In each homotopy degree,
injectivity, surjectivity, or bijectivity of the PL-to-definable
comparison over $\mathcal M$ is equivalent to the corresponding
property over $\mathcal N$.  Thus the comparison is a weak
equivalence over one structure exactly when it is so over the other.
\end{theorem}

We first isolate the continuity assertion required to move
coefficients.  Elementary transfer alone would not provide this
movement.

\begin{lemma}\label{lem:coefficient-cells-joint-evaluation}
Fix a nonempty closed bounded definable set $E$ over $\mathcal M$
and a formula $\varphi(a;x,y)$ with fixed parameters in
$\mathcal M$.  Let $T$ be the definable set of tuples $a$ for which
the formula defines a homeomorphism $h_a:E\to E$ satisfying any
specified finite list of first-order conditions.  There is a finite
partition of $T$ into cells defined over $\mathcal M$ such that
both maps
\[
 (a,x)\longmapsto h_a(x),\qquad
 (a,x)\longmapsto h_a^{-1}(x)
\]
are jointly continuous on every cell times $E$.  The same partition
and continuity assertions hold after elementary extension.
\end{lemma}

\begin{proof}
Homeomorphism validity is first order: functionality and bijectivity
are first-order graph conditions, and continuity of the map and its
inverse are expressed by epsilon--delta formulas in all coordinates.
Thus the stated set $T$ is definable; quantifier elimination for the
o-minimal theory is not required.

Consider the graph family
\[
 \Gamma=\{(a,x,y):a\in T,\ x,y\in E,\ \varphi(a;x,y)\}.
\]
Definable Hardt triviality, followed by cell decomposition of the
base, gives a finite cell partition with a trivialization over each
cell.  The precise classical inputs are
\cite[Theorem~5.22, pp.~62--63; Theorem~2.10, p.~19]{Coste1999}.
These statements apply to a fixed o-minimal expansion of an arbitrary
real closed field; the graph family need not be proper over its base.
Choose $a_0\in C(M)$ for a nonempty cell $C$.  By identifying the
trivializing fibre with the fibre at $a_0$, write the trivialization as
\[
 \theta:C\times\Gamma_{a_0}\xrightarrow{\ \cong\ }\Gamma|_C
\]
over $C$.  The set $\Gamma_{a_0}$ is closed and bounded, since it is
the graph of a continuous map from the closed bounded set $E$ to itself.

Compose $\theta$ with the projection to the input coordinate.  The
result is a continuous definable bijection
\[
 J:C\times\Gamma_{a_0}\longrightarrow C\times E.
\]
We verify inverse continuity without using ordinary compactness of
a non-Archimedean field.  If $J^{-1}$ were discontinuous at
$z=(a,x)$, there would be $\epsilon>0$ such that $z$ lies in the
closure of the definable set of points $z'$ with
$\|J^{-1}(z')-J^{-1}(z)\|\geq\epsilon$.
Definable curve selection provides a curve $z(t)$ in this set tending
to $z$ as $t\downarrow0$.  Write
$J^{-1}(z(t))=(a(t),\gamma(t))$.  Its first coordinate tends to $a$,
and its second lies in the closed bounded set $\Gamma_{a_0}$.
Coordinatewise monotonicity gives a limit
$\gamma\in\Gamma_{a_0}$: bounded definable one-variable functions
have finite one-sided limits, and closedness retains the limit in
the set.  Continuity of $J$ gives $J(a,\gamma)=z$.  Bijectivity
then identifies $(a,\gamma)$ with $J^{-1}(z)$, contradicting the
chosen distance bound.  The curve-selection and definable-compactness
inputs are \cite[Theorems~3.2 and~3.4, pp.~26--27]{Coste1999}.

Consequently the output projection of $\theta J^{-1}$ is joint
continuous evaluation on $C\times E$.  Repeating the argument with
the input and output interchanged proves joint inverse evaluation.
All the objects just constructed and their continuity assertions are
defined over $\mathcal M$.  Elementarity transfers them to
$\mathcal N$.
\end{proof}

\begin{proof}[Proof of Theorem~\ref{thm:elementary-extension-marked-automorphisms}]
The empty case is immediate, so suppose $D$ is nonempty.
If two graph formulas define the same map over $\mathcal M$, the
sentence asserting their equality remains true over $\mathcal N$.
Functionality, bijectivity, joint continuity in both directions,
parameter preservation, and the marked conditions are first order.
For a setwise mark the condition is equality of its image with that
mark, equivalently membership in the mark if and only if the image
belongs to it.  Thus no closedness assumption on the marks enters.
Composition, graph reversal for inversion, and affine pullback
commute with interpretation.  This proves that the extension map
is a well-defined homomorphism of simplicial groups.

We use the classical finite-simplex description of their homotopy
groups.  For a simplicial group $G$, put
\[
 N_qG=\bigcap_{i=1}^q\ker d_i,\qquad
 Z_qG=\bigcap_{i=0}^q\ker d_i\quad(q\geq1),\qquad Z_0G=G_0.
\]
Then $\pi_q(|G|,1)=Z_qG/d_0N_{q+1}G$, including the component
group for $q=0$.  Simplicial groups are Kan; the multiplication and
normal-boundary assertions needed for this formula are the classical
results in \cite[Theorem~16.6, p.~65; Theorem~17.1 and
Propositions~17.2--17.4, pp.~67--69]{May1967Simplicial}.
Reversing the simplex vertices converts that source's last-face
Moore convention into the displayed first-face convention.  In particular, a positive-degree
class has a single simplex representative equal to the identity
on every face.

\medskip\noindent\textit{Surjectivity.}
Take such a representative $h$ over $\mathcal N$ in degree $q$,
where $q=0$ is allowed.  Regard it as a homeomorphism of the closed
bounded set $E=D\times\Delta^q$.  Choose one graph formula for $h$
and replace its finitely many additional parameters by variables
$a$.  In the coefficient set $T$ impose parameter preservation, all
marked conditions, and, when $q\geq1$, the identity on every face
of $D\times\Delta^q$.  These are finite first-order requirements.
The tuple defining $h$ belongs to $C(N)$ for one cell of the
partition in Lemma~\ref{lem:coefficient-cells-joint-evaluation}.

Every such cell is definably homeomorphic over $\mathcal M$ to
$M^d$ for its dimension $d$
\cite[Proposition~2.5, pp.~18--19]{Coste1999}.
Choose $a_0\in C(M)$ and, in the extended coordinate chart, join
the tuple defining $h$ to $a_0$ by a straight segment.  Pulling back
the joint evaluation and inverse evaluation along this definable
path gives a total homeomorphism on
$D\times\Delta^q\times[0,1]$, preserving all parameters and the
required marks and fixing $D\times\partial\Delta^q\times[0,1]$.
At the last endpoint it is the extension of a simplex over
$\mathcal M$.  Triangulating the parameter prism by finitely many
rational simplices turns this family into a simplicial homotopy;
the inverse restricts on the same simplices.  For $q=0$ it is simply
a path between the two vertices.  Thus the two representatives have
the same homotopy class, proving surjectivity in every degree.

\medskip\noindent\textit{Injectivity.}
Let $h\in Z_qG_{\mathcal M}$ have trivial image in the homotopy
group over $\mathcal N$.  The displayed Moore description gives
$w\in N_{q+1}G_{\mathcal N}$ with $d_0w=h_{\mathcal N}$.
Choose one formula for the graph of $w$ and replace its extra
parameters by a finite variable tuple.  Being a valid marked total
homeomorphism, satisfying $d_iw=1$ for $1\leq i\leq q+1$, and
satisfying $d_0w=h$ are one first-order condition with parameters
in $\mathcal M$.  The condition is satisfiable in $\mathcal N$,
hence in $\mathcal M$.  Its solution is a simplex
$w_0\in N_{q+1}G_{\mathcal M}$ with $d_0w_0=h$.  This proves
injectivity, also for $q=0$.  Translation deals with every other
base component.  The extension map is therefore a weak equivalence.

Apply this result with $D=B\times[0,1]$ and the two marked ends to
obtain the concordance assertion.  For a finite rational PL model,
extending a PL map and extending its definable graph give exactly
the same map.  The resulting square commutes, and its PL vertical
map is a weak equivalence by Theorem~\ref{result-40}.  Its other
vertical map has just been proved to be a weak equivalence.
The assertions about each comparison homomorphism follow.
\end{proof}

\begin{corollary}\label{cor:complete-theory-concordance-comparison}
Fix a complete o-minimal theory in the given language having a real
model, and a rational realization $B$ of a finite abstract closed PL
manifold.  In each homotopy degree, the
injectivity, surjectivity, and bijectivity properties of its
PL-to-definable concordance comparison are the same over all models
of that theory, including the real model.
\end{corollary}

\begin{proof}
Two models of a complete theory have elementary embeddings into a
common model.  Indeed, place their elementary diagrams in disjoint
sets of added constants.  A finite part of either diagram is the
instance of an existential sentence consistent with the complete
theory, hence implied by it.  The two finite parts can therefore be
realized simultaneously.  Compactness gives the common model and
the elementary embeddings; these classical inputs are
\cite[Theorem~2.1.4, pp.~34--35; Lemma~2.3.3, p.~44]{Marker2002ModelTheory}.  The rational PL model is preserved by
both embeddings.  Applying the theorem to the resulting zigzag
proves the assertion.
\end{proof}

The real-model hypothesis is substantive.  In the language of
ordered fields with one named constant, the complete theory of a
real closed field containing a named positive infinite element is
o-minimal but has no real model.  The corollary also does not supply
a comparison property which has not been proved over its real model.

\begin{corollary}\label{cor:countable-definable-automorphism-model}
If the language is countable, the marked automorphism group in
Theorem~\ref{thm:elementary-extension-marked-automorphisms} has a
countable simplicial group as a weak model.
\end{corollary}

\begin{proof}
Start with the finitely many parameters defining $D$ and its marks.
Repeatedly close a countable set under the language's functions and
adjoin a witness to every existential formula with parameters in
the current set which is true in $\mathcal M$.  The union of this
countable sequence is a countable elementary substructure by the
Tarski--Vaught criterion
\cite[Proposition~2.3.5, p.~45]{Marker2002ModelTheory}.  The set and its marks are defined over it,
and the assertion that $D$ is closed and bounded transfers.  There are only
countably many formulas and finite parameter tuples in each
simplicial degree.  The theorem applied to this elementary
substructure proves the claim.
\end{proof}

This last conclusion concerns cardinality.  An arbitrary elementary
diagram need not be decidable, so it supplies neither an effective
test for homeomorphism codes nor an oracle-free recursive
presentation of the homotopy groups.

Countability alone also gives no PL-to-definable weak equivalence.


\subsection{Finite formula families and infinitesimal approximation}
We describe the local topology of a finite definable family of
homeomorphisms and the separate limitation of finite PL approximation.

Let $R$ be a real closed field with a fixed o-minimal expansion,
and let $K\subseteq R^n$ be closed, bounded, and definable.
Fix finitely many formulas for graphs of maps $K\to K$, with
finite parameter tuples.  Retain the parameters whose graphs are
homeomorphisms satisfying a prescribed finite list of definable
marked-set or boundary conditions.  Identify parameters that define
the same whole map, also across different formulas, and write
$\mathcal H$ for the resulting definable quotient.  No closure under
composition or inversion within the chosen formats is assumed.
For nonempty $K$, give it the metric
\begin{equation}\label{eq:finite-format-uniform-metric}
 d(h,g)=\max_{x\in K}\|h(x)-g(x)\|
          +\max_{x\in K}\|h^{-1}(x)-g^{-1}(x)\|.
\end{equation}
All topologies and inequalities here are over $R$, including when
$R$ is non-Archimedean.  The following application uses Walsberg's
affineness theorem for definable metric spaces
\cite[Theorem~7.1, p.~16]{Walsberg2015}.

\begin{theorem}\label{thm:finite-format-homeomorphism-contractions}
The metric space $(\mathcal H,d)$ is definably homeomorphic to a
definable subset of some $R^N$ with its Euclidean topology.
For every $h_0\in\mathcal H$ and every $\epsilon>0$, there are a
definable $d$-open neighborhood $U$ of $h_0$, a $\delta>0$, and a
definable continuous contraction $C:U\times I\to U$ satisfying
\[
 B_d(h_0,\delta)\subseteq U\subseteq B_d(h_0,\epsilon),\qquad
 C(h,0)=h,\quad C(h,1)=h_0,\quad C(h_0,s)=h_0.
\]
Both $(h,s,x)\mapsto C(h,s)(x)$ and
$(h,s,x)\mapsto C(h,s)^{-1}(x)$ are jointly definable and continuous.
Thus the contraction consists of actual homeomorphisms in the same
finite family of formats satisfying the prescribed marks.  If $K$
is empty, the quotient is empty or a singleton and the assertion is
interpreted accordingly.
\end{theorem}

\begin{proof}
Admissibility of a parameter tuple is definable: graph uniqueness,
bijectivity, continuity, inverse continuity, and the finite marked
requirements are first-order conditions.  A finite label combines
the different formats into a single definable parameter space.
Equality of their whole graphs is a definable equivalence relation.
Elimination of imaginaries for o-minimal expansions of fields
realizes its quotient as a definable set; this is also the quotient
convention in \cite[p.~1]{Walsberg2015}.  Evaluation and inverse
evaluation are definable on that quotient.  Definable compactness
of $K$ gives the two maxima in
\eqref{eq:finite-format-uniform-metric}; their sum is a definable
metric.

We verify the hypothesis of Walsberg's theorem over the given
field: every definable $d$-discrete subset is finite.  Suppose
instead that $D\subseteq\mathcal H$ is an infinite definable
$d$-discrete subset.  Apply definable Hardt triviality to the
family of its compact graphs and refine the parameter partition
into cells \cite[Theorem~5.22, pp.~62--63]{Coste1999}.  One cell
$A\subseteq D$ has positive dimension, and over it there is a
definable graph trivialization
\[
 \theta:A\times G\longrightarrow
 \{(a,x,h_a(x)):a\in A,\ x\in K\},
\]
where $G$ is one closed bounded graph fibre.  Projection after
$\theta$ to $(a,x)$ is a continuous definable bijection
$A\times G\to A\times K$.

This bijection is locally a homeomorphism on the parameter cell.
Indeed, a cell is locally closed, and each of its points has a
closed bounded definable neighborhood $L$ within the cell
\cite[Proposition~3.7 and Exercise~3.8, pp.~27--28]{Coste1999}.
On $L\times G$ the projection is a continuous definable bijection
between closed bounded definable sets, so its inverse is continuous.
One can check this without ordinary compactness: the image of
every closed definable subset is closed by definable compactness,
which proves inverse continuity on definable basic neighborhoods.
Projection to $(a,y)$ gives the same conclusion for the inverse
graphs.  Consequently both $(a,x)\mapsto h_a(x)$ and
$(a,y)\mapsto h_a^{-1}(y)$ are jointly continuous on $A\times K$.

On each $L\times K$ these maps are uniformly continuous, by
definable compactness.  For example, for a continuous definable
map on a closed bounded set, the pairs of points whose images
differ by at least $\epsilon$ form a closed bounded set disjoint
from the diagonal.  If nonempty, the distance of its points from
the diagonal has a positive minimum.  This supplies a uniform
modulus at $\epsilon$; the empty case is immediate.  Applied to
the two evaluation maps, it shows that
$a\mapsto h_a$ is continuous from the Euclidean cell $A$ to
$(\mathcal H,d)$.  A sufficiently small $d$-ball about any $a\in A$
meets $D$ only in $a$.  Its inverse image would make $\{a\}$ open
in the positive-dimensional cell, which is impossible.  This
contradiction proves definable separability in the exact sense
used by Walsberg.  Theorem~7.1 therefore gives a definable
homeomorphism $\varphi:(\mathcal H,d)\to X\subseteq R^N$.
No ordinary separability or Archimedean hypothesis has been used.

We spell out the strong local contraction, including the possibility
that $X$ is not locally closed.  Replace the ambient Euclidean
space by a bounded definably homeomorphic copy if necessary.
Put $x_0=\varphi(h_0)$ and
$V=\varphi(B_d(h_0,\epsilon))$.  Triangulate the closed bounded
set $\overline X$ compatibly with $X$, $V$, and $\{x_0\}$
\cite[Theorem~4.4, pp.~42--44]{Coste1999}.  The inverse image of
$x_0$ is a vertex $v$.  In the triangulation, intersect the open
star of $v$ with the inverse image of $X$.  This is an open
neighborhood of $v$ in that inverse image.  It is contained in the
inverse image of $V$: otherwise a compatible open simplex lying
in $X\setminus V$ would have $v$ as a face, contrary to the
relative openness of $V$ at $x_0$.

Straight contraction to $v$ preserves this neighborhood.  For a
point in an open simplex containing $v$ as a face, its segment to
$v$ remains in that open simplex until the endpoint.  Thus no
missing faces of $X$ are introduced.  Transport the contraction
back through the triangulation and $\varphi$.  This gives $C$ and
$U$ with the stated endpoint identities and with $U$ invariant
throughout the contraction.  As $U$ is $d$-open at $h_0$, it contains
a ball of some positive radius $\delta$.

Finally, evaluation is continuous for the metric topology directly:
at $(g,y)$,
\[
 \|h(x)-g(y)\|
 \leq\max_{z\in K}\|h(z)-g(z)\|+\|g(x)-g(y)\|.
\]
The first term is controlled by $d(h,g)$ and the second by
continuity of $g$.  Applying the same argument to inverses proves
joint inverse evaluation.  Composition with $C$ proves both
family assertions.  No continuous choice of coefficient tuples
is needed.
\end{proof}

\begin{proposition}\label{prop:infinitesimal-pl-approximation-obstruction}
If $p:[0,1]\to R$ is continuous and affine on at most $m$ intervals,
where $m\geq1$ is an ordinary finite integer, then
\begin{equation}\label{eq:finite-piece-quadratic-lower-bound}
 \max_{x\in[0,1]}|x^2-p(x)|\geq\frac{1}{8m^2}.
\end{equation}
The coefficients and breakpoints may be arbitrary elements of $R$.
Hence no finite PL map approximates $x\mapsto x^2$ to a positive
infinitesimal uniform tolerance over a non-Archimedean real closed
field.  Nevertheless $x\mapsto x^2$ is semialgebraically isotopic
to the identity relative to the endpoints.
\end{proposition}

\begin{proof}
One affine interval $[a,b]$ has length at least $1/m$.  Put
$c=(a+b)/2$ and $e(x)=x^2-p(x)$.  Affineness gives
\[
 e(a)+e(b)-2e(c)=\frac{(b-a)^2}{2}.
\]
If $E$ is the uniform error, the left-hand side has absolute value
at most $4E$.  Thus $E\geq(b-a)^2/8\geq1/(8m^2)$.
The argument is ordered-field arithmetic, so allowing nonstandard
coefficients or breakpoints does not change it.  A positive
infinitesimal is smaller than each of these positive rational bounds
as $m$ ranges over ordinary finite integers.

The semialgebraic family
\[
 H_s(x)=(1-s)x^2+sx,\qquad 0\leq s\leq1,
\]
fixes both endpoints.  If $x<y$ in $[0,1]$, then
\[
 H_s(y)-H_s(x)=(y-x)\bigl((1-s)(x+y)+s\bigr)>0.
\]
Each slice is therefore an increasing semialgebraic homeomorphism
of the interval.  The total map $(x,s)\mapsto(H_s(x),s)$ is a
continuous definable bijection of the closed bounded square, so
definable compactness makes its inverse continuous.  It is the
required genuine semialgebraic isotopy.
\end{proof}

The contraction neighborhood depends on the finite formula family.
No contraction on the union of all formats is asserted.  The
infinitesimal approximation obstruction concerns finite PL
approximation and does not contradict the displayed isotopy or
establish a failure of a higher concordance comparison.


\Needspace{12\baselineskip}
\section{Piecewise linear coefficient extension}
\begin{theorem}[All PL concordance homotopy groups are independent of the real closed field]\label{result-40}
Let $F\subset K$ be an extension of real closed fields and let $B$
be a closed compact PL manifold given by a finite polyhedral model
over $F$. Extension of coefficients induces a weak equivalence
of simplicial groups
\[
 P_{{\rm PL},F}(B)\longrightarrow P_{{\rm PL},K}(B_K).
\]
Consequently a fixed finite abstract PL model has the same PL
concordance homotopy type over any two real closed fields, by the
zigzag through a realization over the field of real algebraic
numbers. This is a zigzag of coefficient-extension maps; it is
not a literal map from a non-Archimedean field's families to
ordinary topological families.
\end{theorem}
\begin{proof}
We give the parameter-preserving argument. A statement only about
unparametrized homeomorphisms would not suffice.

\emph{Extension on product cells.}
Let $Z$ be the underlying space of a finite complex of convex
cells, let $Q$ be a finite parameter polyhedron and $Q_0$ a
subcomplex. Suppose a PL homeomorphism of $Z\times Q_0$ is over
$Q_0$ and preserves each cell of $Z$. It extends to a PL
homeomorphism of $Z\times Q$ over $Q$, still preserving every
cell. More generally it extends relative to any already
specified union of product cells on which the same conditions hold.
To see this, triangulate $Q$ and proceed over the cells
$\sigma\times\tau$ in increasing total dimension. On the boundary
of a cell not already specified the homeomorphism has been
defined, maps that boundary to itself, and preserves the parameter
coordinate. Choose an interior point
$c=(c_\sigma,c_\tau)$ of this convex product cell and extend
by coning:
\[
 c+t(z-c)\longmapsto c+t(H(z)-c),
 \qquad z\in\partial(\sigma\times\tau),\quad 0\le t\le1.
\]
This is PL after starring a triangulation of the boundary at
$c$, and is a homeomorphism by the same construction for the
inverse boundary map. Its parameter coordinate is unchanged,
since
\[
 \operatorname{pr}_Q(c+t(H(z)-c))
 =(1-t)c_\tau+t\operatorname{pr}_Q(z)
 =\operatorname{pr}_Q(c+t(z-c)).
\]
Thus it is coning on the \emph{total product cell}, not separate
radial coning in each fibre. The maps agree on faces and the
finite induction proves the claim. It works over every ordered
field. In particular cell-preserving PL families can be contracted
through such families. Two families agreeing on a union of
product cells can be contracted with that agreement retained:
first perform this same extension on the common part, then on
the remaining cells separately.

Put $Z=B\times[0,1]$, with a fixed finite convex-cell
decomposition $\mathcal C_Z$ over $F$ that contains its bottom
and top as subcomplexes. Represent a homotopy group element by
a PL homeomorphism
\[
 h:Z_K\times P_K\longrightarrow Z_K\times P_K,
 \qquad P=[0,1]^q,
\]
over $P_K$, fixing the bottom and equal to the identity on
$Z_K\times\partial P_K$. For $q=0$ the last condition is empty.
Fix a cell decomposition $\mathcal C_P$ of $P$ over $F$.

\emph{Finite transferred data.}
Apply simultaneous subdivision to the one-way chain consisting
of two distinct copies of the total space,
\[
 Z_K\times P_K\xrightarrow{h}Z_K\times P_K
          \xrightarrow{\operatorname{pr}_P}P_K.
\]
The source and target subdivisions $M,N$ can be required to
refine the same fixed product cells
$\mathcal C_Z\times\mathcal C_P$, and the base subdivision $L$
refines $\mathcal C_P$. Then $h:M\to N$ is a simplicial
isomorphism and both projections to $L$ are simplicial, the
source projection being the composite of the displayed maps.
This is exactly the permitted one-way case of
\cite[Theorems 2.14--2.15, pp.~17--19]{RourkeSanderson1972};
it does not ask for a simplicial subdivision of a cycle.
On the bottom and on $Z_K\times\partial P_K$, where $h$ is
the identity, the induced source and target subdivisions coincide.

Keep the finite abstract complexes and all their carrier labels
fixed, but regard their vertex coordinates as unknowns. The
requirements that they realize subdivisions of the designated
fixed product cells, that the abstract map $M\to N$ is realized
by a PL homeomorphism, and that each source vertex and its
corresponding target vertex project to their designated vertex
of $L$, are a finite first-order condition over $F$. To spell
out the finiteness, affine independence is expressed by nonzero
minors; simplex membership uses finitely many barycentric
coordinates; intersection along the prescribed common faces
and coverage of each fixed cell are finite quantified field
formulas. Add equality of the corresponding source and target
vertices on the bottom and on the parameter boundary. The
given configuration satisfies this condition over $K$.
Quantifier elimination for real closed fields, hence
$F\preccurlyeq K$, gives a configuration $M',N',L'$ over $F$
satisfying the same condition. It determines a PL family
$h_F$ of concordances, based on $\partial P$, over $F$.

There are matching-vertex PL homeomorphisms over $K$
\[
 u:M\longrightarrow M',\qquad
 v:N\longrightarrow N',\qquad
 \beta:L\longrightarrow L',
 \qquad h_F=vhu^{-1}.
\]
Each of $u,v$ preserves every fixed product cell, and $\beta$
preserves every fixed base cell. The projection identities
hold on entire simplices, not just their vertices:
\[
 \operatorname{pr}_P u=\beta\operatorname{pr}_P,\qquad
 \operatorname{pr}_P v=\beta\operatorname{pr}_P.
\]
Indeed a simplex of $M$ projects into one simplex of $L$,
on which $\beta$ is affine. For
$x=\sum_i\lambda_i x_i$ in this simplex, the vertex-incidence
equalities give
$\operatorname{pr}_P u(x)=\sum_i\lambda_i\beta(\operatorname{pr}_P x_i)
=\beta(\operatorname{pr}_P x)$; the proof for $N$ is identical.

Set
\[
 u_0=(\operatorname{id}_Z\times\beta^{-1})u,\qquad
 v_0=(\operatorname{id}_Z\times\beta^{-1})v.
\]
They preserve the parameter pointwise and every product cell.
They agree on the bottom and on $Z_K\times\partial P_K$,
because the matching-vertex maps agree there. Apply the
product-cell extension construction over the enlarged parameter
$P_K\times[0,1]$ to obtain PL isotopies $u_s,v_s$ from the
identity to $u_0,v_0$ which retain this agreement for every $s$.
Concretely extend their common boundary data first on the union
of these product subcomplexes, then extend over the other cells.
The families
\[
 v_s h u_s^{-1}
\]
are therefore concordances fixed on the bottom and based on
the parameter boundary throughout. Their last member is
$(\operatorname{id}_Z\times\beta^{-1})h_F
(\operatorname{id}_Z\times\beta)$.
The same product-cell construction gives an isotopy of the
cell-preserving map $\beta$ to the identity, preserving every
base cell. Conjugating by that isotopy connects this last member
to $h_F$ while keeping the parameter boundary based. This proves
surjectivity on every homotopy group, including components.

\emph{Injectivity.}
Suppose a based PL family over $F$ becomes nullhomotopic over
$K$. Take one finite PL triangulation of its $K$-nullhomotopy.
Its vertex coordinates again give a finite first-order
configuration problem over $F$, with the extra requirement that
on the parameter boundary its PL map equal the prescribed
$F$-family or the identity, as appropriate. That equality is
itself a finite field formula, because the prescribed family
has finitely many affine pieces over $F$. Transfer therefore
gives an $F$-PL nullhomotopy with exactly the same boundary data.
The same argument applies to a path between two zero-simplices.
Thus the maps on all homotopy groups are injective.

Finally simplicial groups are Kan, and their homotopy groups
are computed by these finite parameter families. The established
isomorphisms prove the weak equivalence. Real algebraic numbers
embed as an ordered real closed subfield in every real closed
field; realizing a finite abstract complex by standard rational
vertices gives the asserted comparison zigzag.
\end{proof}\section{Parameter regularity and finite fragmentation}
\begin{theorem}[A precise regularization criterion for higher PL comparison]\label{result-41}
Fix an o-minimal expansion of the real field. Let \(B\) be a compact
closed PL manifold, put \(N=B\times I\), and let \(k\geq1\).
A class in \(\pi_kP_{\mathrm{def}}(B)\) belongs to the image of
\begin{equation}
 \pi_kP_{\mathrm{PL}}(B)\longrightarrow\pi_kP_{\mathrm{def}}(B)
 \label{equation-80}
\end{equation}
if and only if it has a based definable representative
\(U:N\times P\to N\), \(P=[0,1]^k\), for which a constant \(L<\infty\)
satisfies
\begin{equation}
 \|U_p(x)-U_q(x)\|_\infty\leq L\|p-q\|_\infty
 \quad(x\in N,\ p,q\in P).
 \label{equation-81}
\end{equation}
Here \(U_p\) fixes \(B\times\{0\}\) pointwise, preserves the top end,
and \(U_p=\operatorname{id}\) for \(p\in\partial P\).
Only the forward parameter estimate is assumed.
In particular, real semialgebraic
\(\pi_2P_{\mathrm{PL}}(T^6)\to\pi_2P_{\mathrm{SA}}(T^6)\)
is surjective exactly when each based semialgebraic two-cube is based
semialgebraically homotopic to one satisfying \eqref{equation-81}.
This criterion does not itself assert existence of such representatives.
\end{theorem}

\begin{proof}
All metrics in this proof are maximum norms in fixed finite PL models.
Changes of such models multiply real Lipschitz constants by finite
constants. We give the finite construction, including its relative
boundary conventions; we do not assume a parametric Hauptvermutung.

\textbf{A definable fixed-point fact.}\quad
If \(T:[0,1]^a\to[0,1]^a\) is continuous and definable and has
contraction constant \(c<1\), it has one fixed point. Indeed,
\(\|T(t)-t\|_\infty\) has a minimum \(m\). At a minimizing \(t_0\),
the assumption \(m>0\) would give
\[
 \|T(T(t_0))-T(t_0)\|_\infty\leq cm<m.
\]
Uniqueness follows from the contraction estimate. If \(t,t'\) are
the fixed points for \(T,T'\) with a common contraction bound, then
\begin{equation}
 \|t-t'\|_\infty\leq
 (1-c)^{-1}\sup_z\|T(z)-T'(z)\|_\infty.
 \label{equation-82}
\end{equation}
Consequently a jointly continuous definable family of such maps has
a jointly continuous definable fixed point. Definability follows from
uniqueness; continuity follows from \eqref{equation-82} and uniform continuity on
the compact variable cube, locally in the other variables. This proof
uses a minimum, not an infinite limiting iteration.

\textbf{The double and the finite Hudson functions.}\quad
For a connected positive-dimensional \(B\), double \(N\) along its two
ends, obtaining the closed PL manifold
\[
 X=N^+\cup N^-\cong B\times S^1.
\]
Write \(A_0,A_1\) for the disjoint seams. Reflect \(U\) to the second
half. The resulting family \(V\) preserves \(N^\pm,A_0,A_1\), fixes
\(A_0\) pointwise, and is the identity on \(\partial P\).
The fixed PL reflection changes \eqref{equation-81} by only a finite factor.

Choose a finite PL triangulation \(Q\) containing all four marks as
subcomplexes. Choose the preliminary subdivisions needed for the
small-ball separation construction with at most \(k\) support stars;
this is the finite subdivision in Mn\"ev, Proposition~57,
printed pp.~89--90. It depends on \(Q,k\), before a number of time
levels is chosen. Let \(K\) be a sufficiently fine such subdivision,
also refining \(\operatorname{sd}Q\).
Enumerate its vertices \(v_1,\ldots,v_\nu\), and let \(\lambda_i\)
be their continuous PL barycentric functions.
For an integer \(M>0\), use the sequence
\begin{equation}
 f^{a\nu+i}(y)=\frac aM+
        \frac1M\sum_{j=1}^{i}\lambda_j(y),
 \quad 0\leq a<M,\quad0\leq i\leq\nu .
 \label{equation-83}
\end{equation}
The coincident values at successive block ends are identified.
There are \(R=M\nu+1\) values, with
\[
 0=f^0\leq f^1\leq\cdots\leq f^{R-1}=1.
\]
Each consecutive difference is supported in one closed vertex star.
The spatial Lipschitz constants of all \(f^j\) are bounded by
\(M^{-1}\sum_i\operatorname{Lip}(\lambda_i)\).
Choose \(M\) so that this common bound \(c\) satisfies \(cL<1\).
This is legitimate over the real field. The later all-field version
must separately justify a standard bound before choosing this integer.

\textbf{An explicit graph deformation.}\quad
For \(0\leq v\leq1\) and \(0\leq q\leq T_v\), put
\begin{equation}
 \rho_{y,v}(q)=
   \min_{0\leq j<R}\max\{f^j(y),q-2jv\},
 \qquad T_v=1+2(R-1)v.
 \label{equation-84}
\end{equation}
These are finite PL formulas. The values lie in \([0,1]\);
\(\rho_{y,v}(0)=0\), \(\rho_{y,v}(T_v)=1\),
and \(\rho_{y,0}(q)=q\) for \(0\leq q\leq1\).
They have spatial Lipschitz bound \(c\).
For \((x,p,v)\in X\times P\times I\), solve
\begin{equation}
 t_i=\rho_{V_t(x),v}(T_vp_i),\qquad 1\leq i\leq k,
 \qquad H_{p,v}(x)=V_t(x).
 \label{equation-85}
\end{equation}
The right side is a contraction in \(t\) with constant \(cL<1\).
The fixed-point fact supplies unique joint continuous definable
solutions. Moreover
\begin{equation}
 H_{p,v}^{-1}(y)=
 V^{-1}_{(\rho_{y,v}(T_vp_i))_{i=1}^k}(y).
 \label{equation-86}
\end{equation}
For this formula the proposed parameter vector is exactly the vector
required in \eqref{equation-85}; uniqueness verifies both inverse identities.
The total family inverse is jointly continuous, since its
total map is a continuous bijection of compact spaces. Thus \eqref{equation-86}
also proves joint continuity of the inverse family.

We have \(H_{p,0}=V_p\). The deformed maps preserve the halves and
seams. On \(A_0\), every \(V_t\) is the identity, so \(H\) is also
the identity. If a coordinate of \(p\) is an endpoint, the same
coordinate of \(t\) is that endpoint. Thus on every parameter boundary
face \(V_t=\operatorname{id}\), and the deformation is based.
No estimate for \(V^{-1}\) in the parameter variable has entered.

\textbf{Finite support labels at the terminal time.}\quad
Set \(D=2R-1\), and at \(v=1\) write \(q_i=Dp_i\).
For each moving band
\(2j\leq q_i\leq2j+1\), \(0\leq j<R-1\), formula \eqref{equation-84} is
\begin{equation}
 \rho_{y,1}(q_i)=
 \min\{f^{j+1}(y),\max\{f^j(y),q_i-2j\}\}.
 \label{equation-87}
\end{equation}
On the next plateau it equals \(f^{j+1}(y)\).
The last index contributes only the final constant plateau.
Within an interval meeting at most one moving band, the value in
\eqref{equation-87} is independent of \(q_i\) outside the star supporting
\(f^{j+1}-f^j\).

Cover the finite interval \([0,D]\) by finitely many relatively open
PL intervals, each meeting at most one moving band. Take their
coordinate products and all nonempty intersections. Each is a box.
For such a box \(E_a\), let \(S_a\) be its actual list of moving-band
support stars, one for each active coordinate; repetitions are harmless.
It has at most \(k\) members. Crucially, the label of an intersection
is computed from the bands met by that intersection, rather than
by taking the union of the original box labels. Hence
\begin{equation}
 E_a\subset E_b\quad\Longrightarrow\quad S_a\subset S_b .
 \label{equation-88}
\end{equation}
Formula \eqref{equation-86} is pointwise independent of \(p\in E_a/D\) outside
the union \(J_a\) of the stars in \(S_a\). These are finite support
data depending only on \(K,M,k\), not on the individual values of \(V\).
They are the open PL parameter topology and monotone support
presheaf used in the finite construction.

Mn\"ev's Proposition~57 replaces configurations of at most \(k\) small
stars by compatible configurations of pairwise disjoint conic PL balls
containing those stars. Its natural transformation is in the covering
order. We retain the entire finite indexing poset, including indices
with repeated labels, and make the inclusions strict by the following
finite version of the conic enlargement used in Proposition~35,
pp.~61--62. Start with maximal indices and enlarge their balls slightly,
keeping the balls in each configuration disjoint and consistent with
the fixed triangulation. Proceed in decreasing rank. At a lower
index each old ball is compactly contained in the already enlarged
higher balls that contain it. A sufficiently small conic enlargement
stays in all of these interiors and retains disjointness within the
lower configuration. The finitely many requirements can be imposed
simultaneously. This gives strict containment for proper indexing
relations. We use this explicit extension of the conic construction;
we do not apply Proposition~35 to a map lacking its strictness
hypothesis.

Write \(\xi_a\) for the resulting configurations. Outside their union,
the inverse family is pointwise constant on the corresponding box.
Continuity gives pointwise constancy on the boundaries as well.
Relative to one fixed parameter in the connected box, the resulting
homeomorphisms fix the complement and the boundaries of the disjoint
balls. They preserve each ball: the possible permutation of finitely
many ball interiors is locally constant in the parameter and starts
at the identity. Thus the inverse image of every ball is constant.
This is the configuration-prismaticity needed below.

\textbf{A split marked coarsening.}\quad
The preceding separation must respect the boundary of \(N\); applying
an unmarked closed-manifold statement and then restricting would not
suffice. In the present application the input coloring is constant:
\(Q(\lambda)=Q\) for every index to which Proposition~36 is applied.
Thus the marked \(Q\) and the requirement \(K\preceq\operatorname{sd}Q\)
were fixed before \(M\) and \(\xi\) were chosen. The rank construction
below produces new subdivision complexes, but it does not require
choosing \(K\) again or making it refine those new complexes.
All input cells already respect \(N^\pm,A_0,A_1\).
Mn\"ev's Proposition~36, pp.~62--63, supplies its conic-star
construction and the incidence property that a consistent ball
centered at \(b\) meets a \(Q(\lambda)\)-cell only when that cell
contains \(b\). A ball meeting a seam must therefore have its center
in that seam and cannot meet both disjoint seams.
Conic-star invariance, together with local flatness of the seam as
a PL subcomplex, identifies the marked pair with a ball and an
equatorial disk. This local-pair conclusion is an inference from
these facts, not an additional literal assertion in Proposition~36.
The two intersections with the halves are its closed half-balls.

In the coarsening construction of Lemma~6, pp.~63--65, replace a
crossing ball's single interior cell by the two half-ball interiors
and the middle-disk interior, retaining their refined boundary cells.
For a ball avoiding the marks use the original single interior cell.
Every new closed cell is a PL ball; its boundary is a union of the
retained cells. The fine complex refines this split complex because
it already respects the marks. This is a split marked coarsening,
not the unchanged inscribed-ball conclusion of Lemma~6.

The decreasing-rank proof is still finite. At a given rank triangulate
each boundary compatibly with its equatorial disk; extend that boundary
triangulation to the complement, refining the original target cells,
the previously constructed higher-rank subdivisions, and the fixed
marks. The interiors of the present balls are left unchanged.
For nested balls, intersection with the same global half or seam
preserves every inclusion, so their split cells have the required
refinement incidence. This is precisely the comparison needed in
the final step of the proof of Lemma~6.

There is also a direct check of prismaticity after splitting. If \(A\)
is a seam or a half, mark preservation gives
\begin{equation}
 H_p^{-1}(D_0\cap A)=H_p^{-1}(D_0)\cap A
 \label{equation-89}
\end{equation}
for every configuration ball \(D_0\). Thus the split interiors have
constant inverse images. Arbitrary chosen boundary subcells have
constant inverse images because the inverse is pointwise constant
on the ball boundary. That last assertion follows from constancy on
the exterior and continuity; the other balls have positive separation.
No assertion about connected components of a split preimage is needed.

\textbf{From the configurations to literal product prisms.}\quad
Apply the finite procedure in the proof of Mn\"ev's Lemma~8,
pp.~67--68, with no auxiliary \(G_i\), with \(Z=\varnothing\), and
with a constant marked initial target coloring.
Here is the assembly, so that its use does not conceal a new
approximation step. Replace the finite open PL parameter topology
by a strict closed PL approximation (Proposition~26). Strengthen it
with the closed-cell topology of a finite triangulation of \(P\).
Proposition~25 realizes the resulting finite order data on an order
complex, simultaneously recording the target coloring and \(\xi\).
On that order complex perform the decreasing-rank split construction
above, obtaining the two natural transformations consisting of a
subdivision and a coarsening. Paste their two mapping-cylinder
colorings as in Step~4 of Lemma~8. Restrict the resulting coloring
to its terminal copy of \(P\).

All choices just listed use only the finite target coloring,
configuration presheaf, and parameter topology. Their validity for
the present definable input follows from the inverse-image
constancies \eqref{equation-86}--\eqref{equation-89}: the construction never approximates
the input maps or differentiates them. Thus the same proof gives a
finite terminal coloring for this definable homeomorphism family.
If several families have the same data and constancies, the choices
may literally be made once for all of them.

The resulting target prisms are the products \(C\times\sigma\) of
Proposition~4, p.~21. Equation~(11) and Proposition~5, pp.~22--23,
identify their full inverse images with
\begin{equation}
 H^{-1}(C\times\sigma)=D_C^\sigma\times\sigma,
 \label{equation-90}
\end{equation}
where \(D_C^\sigma\) is independent of the parameter on the entire
closed simplex \(\sigma\). These are finite ball complexes.
The spatial cells in the source are at this stage definable balls.
The whole boundary of each product is a union of proper lower
dimensional prisms; on a lower parameter face its spatial
decomposition may be refined. This causes no omission of boundary
data in the induction below. The marks are unions of these prisms.

\textbf{One ambient change makes the spatial cells PL.}\quad
Take the finite family of all source spatial cell closures and all
cells of the fixed target complex \(Q\), including the marks.
A compatible definable triangulation
\(\tau:|K'|\to X\) makes every preimage in this list polyhedral.
Apply marked Hauptvermutung, Shiota, arXiv:1002.1508,
Theorem~2.1, printed p.~5, to the pair of compact polyhedra with
the finite marks \(\tau^{-1}Q_i\) and \(Q_i\).
It gives a PL homeomorphism \(a:|K'|\to X\) carrying those marks
to the corresponding \(Q_i\). Then
\begin{equation}
 \phi=\tau a^{-1}:X\to X
 \label{equation-91}
\end{equation}
preserves every \(Q_i\) setwise, and
\(\phi^{-1}(D)=a\tau^{-1}(D)\) is a subpolyhedron for every source
spatial cell. It is a PL ball: apply the same marked
Hauptvermutung to its definable ball chart and boundary pair.
Thus common conjugation by \(\phi\) makes the source and target
spatial cells PL without losing any fixed mark.

The cell-preserving map \(\phi\) is definably isotopic to the identity,
preserving every \(Q_i\). This follows from the finite cellular
extension construction in the next paragraph, with the parameter
space a point and PL endpoint the identity. It is not an appeal to
an unstated ambient isotopy complement of Hauptvermutung.
Conjugating \(H\) along this isotopy preserves the parameter boundary
identity and the pointwise bottom identity. Call the resulting
terminal family \(W\).

\textbf{Finite endpoint extension by total product coning.}\quad
Construct a PL map \(G\) with the same product-prism correspondence
as \(W\), by increasing total dimension. Prescribe the identity
first on
\[
 E=(A_0\times P)\cup(X\times\partial P).
\]
This is a subcomplex, and identity is the required correspondence
there. For a remaining source prism \(D_0\times\sigma\), all of
\[
 (\partial D_0\times\sigma)\cup(D_0\times\partial\sigma)
       =\partial(D_0\times\sigma)
\]
is already assigned. Make compatible finite boundary subdivisions
so that this prescribed boundary map is PL on the displayed products;
previous maps themselves are unchanged. Choose fixed PL charts of
\(D_0\) and the target spatial ball \(C\) to convex PL balls.
Cone the entire product boundary from centers
\((d_0,b_\sigma)\) and \((c_0,b_\sigma)\), using a common base
barycentre. The extension is a PL homeomorphism. It preserves the
base projection exactly, since corresponding cone segments have
parameter \((1-u)b_\sigma+up\).
This is the full-boundary extension in Mn\"ev's Lemma~3,
proof Steps~2--3, pp.~24--25. Coning only the spatial boundary
with a parameter-dependent formula is not substituted for this
total-product construction.
Finite pasting yields the desired jointly PL \(G\).

\textbf{Relative definable straightening of the remaining cellular map.}\quad
Set \(K=G^{-1}W\). It preserves each source prism and equals the
identity on \(E\). Construct a homotopy \(K\to\operatorname{id}\)
by increasing total prism dimension, relative to \(E\).
For one spatial ball \(C\) and base simplex \(\sigma\), put
\(D_1=\sigma\times I\), where the extra coordinate is homotopy time.
The already assigned face homotopies, initial \(K\), and final identity
give compatible data on
\begin{equation}
 (\partial C\times D_1)\cup(C\times\partial D_1).
 \label{equation-92}
\end{equation}
Use fixed PL ball charts to centered cubes. The prescribed spatial
boundary homeomorphisms extend radially to homeomorphisms \(R_p\)
of the whole spatial cube for every \(p\in D_1\).
For \(p\in\partial D_1\), divide the prescribed whole-ball map by
\(R_p\), obtaining a homeomorphism \(Q_p\) fixing the spatial boundary.
Identify the parameter ball \(D_1\) with a centered cube. Write
\(p=ru\), \(r=\|p\|_\infty\), \(u\in\partial D_1\), when \(r>0\).
The formula
\begin{equation}
 A_r(Q_u)(x)=
 \begin{cases}
   rQ_u(x/r),&\|x\|_\infty\leq r,\quad r>0,\\
   x,&\|x\|_\infty\geq r,
 \end{cases}
 \qquad A_0=\operatorname{id}
 \label{equation-93}
\end{equation}
is a homeomorphism fixing the spatial boundary; the two formulas
agree where they overlap. Its displacement is at most \(2r\),
so it is jointly continuous at the parameter center even though
\(u\) is undefined there. The family \(R_pA_r(Q_u)\), with value
\(R_0\) at the center, extends all of \eqref{equation-92}. It is definable,
and all inverse families are continuous by the same formulas
or the compact total-map argument.
Finite pasting proves the relative homotopy.

This extension argument also establishes the earlier isotopy of
\(\phi\) to the identity, by applying it successively to the cells
of \(Q\). No spatial or parameter smoothness is used.
Multiplying the constructed homotopy by \(G\) joins \(W\) to \(G\)
while preserving the marks and \(E\). Concatenate the graph
deformation and the conjugation homotopy, and restrict from \(X\)
to \(N^+\). This proves the sufficient direction of \eqref{equation-80}.

\textbf{Converse and degenerate cases.}\quad
A jointly PL map on a compact real polyhedron is globally Lipschitz
in fixed PL coordinates, hence satisfies \eqref{equation-81}. A class in the
comparison image therefore has such a representative.
For disconnected \(B\), basedness and connectedness of \(P\) force
each component to remain fixed; apply the construction on finitely
many positive-dimensional components and paste.
For a point component, the concordance is an increasing homeomorphism
of the interval fixing its ends. The convex homotopy
\((1-t)U_p(x)+tx\) is again such a homeomorphism, is based, and ends
at the PL identity. Empty bases have only their empty maps.
These observations complete all cases and the stated named
\(T^6,\pi_2\) criterion.
\end{proof}\Needspace{12\baselineskip}
\section{Exact relative piecewise linear filling}
\begin{theorem}[Relative PL filling in every real base dimension]\label{result-42}
Fix an o-minimal expansion of $\mathbb R$, a closed compact PL
manifold $B$, and $N=B\times I$. Let $P=[0,1]^m$, $m\ge1$,
and let $U_p:N\to N$ be a jointly continuous definable family
of homeomorphisms fixing $A_0=B\times\{0\}$ pointwise and
preserving $A_1=B\times\{1\}$ setwise. If its restriction to
$N\times\partial P$ is jointly PL, there is a jointly PL
family $W$ with the same spatial constraints and
\begin{equation}
 W_p=U_p\quad(p\in\partial P).
 \label{equation-94}
\end{equation}
Consequently $P_{\rm PL}(B)\to P_{\rm def}(B)$ is a weak
equivalence for every such real $B$, including dimensions
three and four. This does not identify its comparison with
the whole ordinary TOP concordance space.
\end{theorem}

\begin{proof}
The first, spatial-clock part works over every o-minimal real
closed field. The later use of the relative regularization consequence of the compact resolution theorem and the finite mesh in
the fragmentation part is over the real field.

\textbf{PL normalization of variation clocks.}\quad
For a fixed $D>3$, define, on $1\le d\le D$, $0\le v\le d$,
\begin{equation}
 T_D(d,v)=
 \begin{cases}
 \min\{v,1+(v-d)/3\},&1\le d\le3,\\
 \min\{v/3,1+(v-d)/D\},&3\le d\le D.
 \end{cases}
 \label{equation-95}
\end{equation}
Both formulas equal $v/3$ at $d=3$. For fixed $d$ this is a
strictly increasing PL homeomorphism $[0,d]\to[0,1]$.
Its vertical slopes lie in $[1/D,1]$. It is globally
$1$-Lipschitz in the maximum norm: every affine gradient
has sum of absolute coefficients at most1, and its domain
is a convex trapezoid. In particular,
\begin{equation}
 |v-v'|\le D|T_D(d,v)-T_D(d,v')|.
 \label{equation-96}
\end{equation}
At one coordinate stage of Theorem~\ref{result-15} choose
$D>\max\{3,\sup_b A_b(1)\}$ and use the inverse of
\begin{equation}
 \tau_b(t)=T_D(A_b(1),A_b(t))
 \label{equation-97}
\end{equation}
in place of the normalized raw clock. If $A_b(t)=3t$,
this clock is exactly the identity.

Here is the full mixed estimate. If two profiles $b,c$
differ by $\delta$ in a previously processed coordinate,
with old bound $K$, variation comparison gives
\[
 \sup_t|A_b(t)-A_c(t)|,\quad |A_b(1)-A_c(1)|
       \le C_0K\delta,\qquad C_0=2d_0N,
\]
where $d_0$ is the number of vector components and $N$ the
line monotonicity bound. Thus
$\sup_t|\tau_b(t)-\tau_c(t)|\le C_0K\delta$.
For $h=\tau_b^{-1}(u)$, $h'=\tau_c^{-1}(u)$,
\begin{equation}
 |A_c(h)-A_c(h')|
 \le D|\tau_c(h)-\tau_c(h')|
 \le DC_0K\delta.
 \label{equation-98}
\end{equation}
The old-coordinate output bound is therefore
$(1+DC_0)K$, and the new-coordinate bound is $D$.
Including the evolving domain map among the vector
components, exactly as in Theorem~\ref{result-15}, proves uniform spatial
Lipschitz bounds for both final factors.

If a vector input is jointly PL on a parameter subpolyhedron,
its line variation is jointly PL there. Its exact finite
maximum formula uses linear constraints on the partition
points, sums, absolute values, and finite piecewise-affine
evaluation. Linear quantifier elimination gives a semilinear
graph for the maximum. Its continuity, proved in Theorem~\ref{result-15},
therefore makes it a finite PL function. Hence the raw
clock, its endpoint height, \eqref{equation-97}, and the inverse total
clock are jointly PL on this parameter part. The inverse
is PL because it is the inverse of a PL homeomorphism of
the compact polyhedral product. Induction preserves this
property for every stage and both final factors.

The same argument respects coefficient fields. If boundary
data are PL over a subfield and $D$ belongs to that field,
all these boundary functions and inverses are PL over it.
On a unit-cube vector model one may choose $D$ to be a
standard integer larger than $3$ and $1+d_0N$.
This quantitative refinement is used later in Theorem~\ref{result-44}.

\textbf{Boundary factors and the fixed bottom.}\quad
Use the PL cubical map $b:N\to Q\subset C=[0,1]^n$
of Theorem~\ref{result-15}, compatible with both spatial ends. Thus
$U^Q=bUb^{-1}$ is jointly PL on $Q\times\partial P$.
Extend that restriction to a jointly PL vector map
$E_\partial:C\times\partial P\to C$: triangulate compatibly,
retain the prescribed vertex values, assign a fixed point
of $C$ at the other vertices, and extend affinely.
This is a vector extension, not a homeomorphic extension.

On $(Q\times P)\cup(C\times\partial P)$ prescribe $U^Q$
and $E_\partial$. The prescriptions agree. The explicit
full-subcomplex formula of Theorem~\ref{result-15} gives a continuous definable
extension $E:C\times P\to C$. Apply \eqref{equation-97} coordinatewise
to $(\mathrm{id},E)$. Restrict the resulting map to $Q$
and conjugate back. We obtain
\begin{equation}
 U=s r^{-1},
 \label{equation-99}
\end{equation}
with uniform spatial Lipschitz bounds, with both factors
jointly PL on the parameter boundary, and preserving both
end marks.

The factors individually fix the bottom. It is a union
of coordinate faces on which $E=\mathrm{id}$. Inductively,
on a free coordinate line within such a face the evolving
vector is $(\mathrm{id},\mathrm{id})$, so its variation is
$2t$ and its raw clock is $3t$. Formula \eqref{equation-97} is then
identity. A coordinate fixed by the face is an endpoint
and remains fixed. Thus every clock is pointwise identity
on every bottom face, and $r=s=\mathrm{id}$ on $A_0$.

The tagged graph family of these two factors has
Hausdorff-Lipschitz boundary, since its boundary maps are
jointly PL on a compact polyhedron. The exact relative
regularization consequence of Theorems~\ref{result-17}--\ref{result-18} gives $g:P\to P$, equal to identity on
the parameter boundary, with both pulled graph families
Hausdorff-Lipschitz. Their spatial bounds and the graph
matching estimate of Theorem~\ref{result-15} give uniform forward parameter
Lipschitz bounds for
\begin{equation}
 \bar r=r\circ g,\qquad \bar s=s\circ g.
 \label{equation-100}
\end{equation}
They are individually bottom-fixed and their parameter
boundary values remain exactly those of $r,s$.

\textbf{The extension-only relative fragmentation assertion.}\quad
We prove the assertion needed for each factor: a definable
concordance family $V$ which is uniformly forward
parameter-Lipschitz and jointly PL on $\partial P$ has
a jointly PL filling of that prescribed boundary.
No PL filling is assumed, and the intervening deformation
need not be fixed on the parameter boundary.

Double $N$ to $X=B\times S^1$, mark both halves and both
seams, and reflect $V$. Choose the Hudson functions and
finite target data of Theorem~\ref{result-41} with $cL<1$. The fixed-point
and terminal support arguments use this inequality and
the spatial marks, not the identity value on parameter
faces. For the boundary track there is a necessary
normalization refinement. Put
\[
 \rho_{y,v}(q)=\min_{0\le j<R}\max\{f^j(y),q-2jv\},
 \quad T_v=1+2(R-1)v,\quad D=2R-1,
\]
and use
\begin{equation}
 \Lambda(v,p)=\max\{p,D(p-1)+T_v\}.
 \label{equation-101}
\end{equation}
This is the jointly PL inverse, in the vertical coordinate,
of $(d,q)\mapsto(d,\min\{q,1+(q-d)/D\})$.
For fixed $v$ it takes $[0,1]$ homeomorphically to
$[0,T_v]$, and
\[
 \Lambda(0,p)=p,\quad\Lambda(1,p)=Dp,\quad
 \Lambda(v,0)=0,\quad\Lambda(v,1)=T_v.
\]
Use the fixed-point construction
\begin{equation}
 t_i=\rho_{V_t(x),v}(\Lambda(v,p_i)),\qquad
 H_{p,v}(x)=V_t(x).
 \label{equation-102}
\end{equation}
Its contraction constant is unchanged. Its inverse is
\begin{equation}
 H_{p,v}^{-1}(y)=
 V^{-1}_{(\rho_{y,v}(\Lambda(v,p_i)))_i}(y).
 \label{equation-103}
\end{equation}
Initial and terminal maps agree with those of Theorem~\ref{result-41}, since
the normalizations agree at $v=0,1$. In particular all
terminal support labels, strictification, and marked ball
coarsenings are unchanged.

The WHOLE parameter-boundary track is jointly PL.
If $p_i$ is0 or1, the corresponding coordinate in \eqref{equation-103}
is that same endpoint. Thus the formula evaluates the
given jointly PL boundary inverse of $V$. Both $\rho$
and $\Lambda$ are jointly PL, so their composition with
that inverse is PL. The total boundary inverse map is a
PL homeomorphism, hence its inverse $H$ is PL as well.
This proves the assertion without treating the original
bilinear normalization $T_vp_i$ as PL.

\textbf{Prescribing the boundary in the PL endpoint.}\quad
Write $H=H_{-,1}$. The finite construction of Theorem~\ref{result-41} gives
source and target prism decompositions over a triangulation
of $P$ compatible with $\partial P$, with cells
\begin{equation}
 D_\sigma\times\sigma\longrightarrow C_\sigma\times\sigma.
 \label{equation-104}
\end{equation}
These are literal products over the entire closed simplex:
Mn\"ev's Proposition4, p.21, describes the target cells;
equation(11) and Proposition5, pp.22--23, give
\[
 H^{-1}(C\times\sigma)=H_{v_\sigma}^{-1}(C)\times\sigma.
\]
Thus each source spatial cell is fixed over the simplex.
All product boundaries are unions of proper lower-total-
dimensional prisms. Over a proper parameter face, spatial
cells may refine $D_\sigma$; the boundary conclusion is
unchanged. All spatial marks are subcomplexes.

On the parameter boundary, source spatial cells are
already PL: $D_\sigma=H_p^{-1}(C_\sigma)$ for any
$p\in\sigma$, and $H_p$ is PL there. Choose the fixed
target PL spatial complex to refine all target cells,
all these boundary source cells, and all marks.
Compatible triangulation and marked Hauptvermutung give
a definable homeomorphism $\phi$ preserving every one of
these PL cells setwise, with all other source cells PL
after taking inverse images.

The conjugate $\phi^{-1}H\phi$ has PL prism decompositions.
On a parameter-boundary prism, its source and target sets
are still the ORIGINAL $D_\sigma,C_\sigma$. Therefore
the ORIGINAL map $H|_{\partial P}$ is a jointly PL
cellular map with exactly the required correspondence.
Prescribe it on $X\times\partial P$ and prescribe identity
on $A_0\times P$. They agree on their intersection.
For the bottom correspondence, any target cell
$C\subset A_0$ has source $H_p^{-1}(C)=C$ because $H$
fixes the bottom pointwise, and $\phi$ preserves this
target cell. Thus the converted source is still $C$.

Rename each converted source cell $\phi^{-1}(D_\sigma)$
as $D_\sigma$. It is a PL ball: it is polyhedral and
definably homeomorphic, as a pair with boundary, to its
target ball; marked Hauptvermutung supplies a PL ball
chart. Products and incidence are retained by conjugation.

Order the unassigned prisms by total dimension. All of
\[
 \partial(D_\sigma\times\sigma)
 =(\partial D_\sigma\times\sigma)\cup
                      (D_\sigma\times\partial\sigma)
\]
has already been assigned a PL boundary homeomorphism
preserving the parameter. Choose compatible finite
subdivisions making it simplicial; refinements of earlier
triangulations change no assigned map or cell set.
Use fixed spatial PL charts to convex balls and cone
the TOTAL product boundary from $(d_0,b_\sigma)$ to
$(c_0,b_\sigma)$. The same base barycentre is used in
both balls. This is a PL homeomorphism and preserves
the parameter exactly: the base of corresponding points
is $(1-u)b_\sigma+up$. This is Mn\"ev's Lemma3,
proof Steps2--3, pp.24--25, not separate radial scaling
of a parameterized spatial map.
Finite induction gives a jointly PL family $G$, with
$G|_{\partial P}=H|_{\partial P}$, fixing the bottom
and preserving every mark. Restrict to the original half.

\textbf{Restoring the initial boundary by a PL collar.}\quad
Choose an affine inner cube $P_{\rm in}$ and a fixed PL
collar
\[
 c:\partial P\times I\longrightarrow
                    P\setminus\operatorname{int}P_{\rm in},
\]
whose inner edge is the corresponding affine copy of
$\partial P$ and whose outer edge is $\partial P$.
It can be made by triangulating the finite boundary
simplex prisms between homothetic cubes; no bilinear
radial formula is asserted to be PL.
Place a rescaled copy of $G$ on $P_{\rm in}$ and put
$H_{z,1-u}$ at $c(z,u)$. At the inner edge these agree
because $G_z=H_{z,1}$; at the outer edge the value is
$V_z$. The track was proved jointly PL, so this is
a jointly PL filling with exactly the initial boundary.
It retains the bottom and top constraints.

Apply the extension-only assertion separately to
$\bar r,\bar s$. It gives PL families $R,S$ with
boundary values $r,s$, respectively. Their quotient
$W=SR^{-1}$ is jointly PL, bottom-fixed, top-preserving,
and has boundary value $sr^{-1}=U$. This proves \eqref{equation-94}.

A based jointly PL $k$-cube nullhomotopic through based
definable concordances has a definable filling on a
$(k+1)$-cube with PL full boundary: the original cube
on one face and identity on the others. Formula \eqref{equation-94}
gives a PL nullhomotopy. Hence comparison is injective
on every positive homotopy group. Theorem~\ref{result-15} gives
surjectivity and Theorem~\ref{result-4} gives the component bijection.
PL translation handles all base components.
Empty, disconnected and zero-dimensional bases are
treated as in the preceding marked constructions.
\end{proof}\Needspace{15\baselineskip}
\section{Comparison over every real closed field}
\begin{theorem}[Higher surjectivity over every fixed o-minimal real closed field]\label{result-43}
Let $M$ be any fixed o-minimal expansion of a real closed field,
let $B$ be a closed compact PL manifold over $M$, and put
$N=B\times I$. Then
\begin{equation}
 \pi_k P_{\mathrm{PL}}(B)\longrightarrow
 \pi_k P_{\mathrm{def}}(B)
 \quad\hbox{is surjective for every }k\geq1.
 \label{equation-105}
\end{equation}
Together with Theorem~\ref{result-4} this gives the component bijection and
surjectivity in every positive degree. This statement does not
assert injectivity in positive degrees over an arbitrary
nonreal model.

The sufficient regularity criterion of Theorem~\ref{result-41} is valid over $M$ on a fixed standard PL model when the forward parameter Lipschitz bound is bounded by an ordinary finite integer. No metric bound on the inverse is required.
\end{theorem}

\begin{proof}
Throughout this proof a \emph{standard bound} means a bound by
an ordinary integer, viewed in $M$. It does not mean merely
the existence of an element of $M$ bounding a function.
All simplicial complexes, posets, lists, and inductions below
are finite in the ordinary sense.
All coordinate and Hausdorff estimates use the maximum norm.

\textbf{The standard spatial model.}\quad
Choose a finite triangulation of $B$ and replace it by the
standard realization of its abstract simplicial complex,
with the vertices taken among the standard basis vectors.
The vertex correspondence is a PL homeomorphism over $M$,
so conjugation by its product with $\mathrm{id}_I$ identifies
the PL and definable concordance problems. All subsequent
metric estimates are taken on this standard model.

The real-algebraic realization is a closed PL manifold as well.
Indeed, the finite PL sphere or ball charts for the links in
the given realization have finite triangulations and affine
descriptors. Fixing those descriptors makes their existence
a first-order statement in real closed fields; the statement
holds over $M$, hence over the real algebraic numbers. This
uses finite PL chart certificates, not a comparison of
parameterized PL homotopies. Fix such charts once and for all.
The product $N$, its double $X=B\times S^1$, the two halves,
and the seams $A_0,A_1$ therefore have fixed standard finite
PL models. Any finite Lipschitz constants of maps with these
fixed real-algebraic PL descriptors have standard bounds.

We also use the cubical PL model of Theorem~\ref{result-15}. Explicitly, the
barycentre of a face $\sigma$ in the barycentric subdivision
is sent to its characteristic vector $\mathbf1_\sigma$.
The resulting affine maps on chains triangulate
\[
 Q=\{x\in[0,1]^v:
          \operatorname{supp}(x)\in K,\ \max_i x_i=1\}.
\]
Thus the conjugating map is PL, has standard Lipschitz bounds
in both directions, and takes every marked subcomplex to a
union of coordinate faces.

\textbf{Definable contractions require no completeness theorem.}\quad
Let $T:[0,1]^a\to[0,1]^a$ be continuous and definable, with
\[
 \|T(t)-T(t')\|_\infty\le c\|t-t'\|_\infty,\qquad c<1.
\]
The definable continuous residual
$r(t)=\|T(t)-t\|_\infty$ attains its minimum $m$ by definable
compactness. If $m>0$ and $t_0$ realizes it, then
\[
 r(T(t_0))
 =\|T(T(t_0))-T(t_0)\|_\infty
 \le c\,r(t_0)<m,
\]
a contradiction. Hence $T$ has a fixed point; contraction
makes it unique. For fixed points of $T$ and $T'$,
\begin{equation}
 \|t-t'\|_\infty
 \le (1-c)^{-1}\sup_z\|T(z)-T'(z)\|_\infty.
 \label{equation-106}
\end{equation}
Consequently the fixed point is definable and varies
continuously in any jointly continuous definable compact
family with a common contraction constant. Definable uniform
continuity supplies the required local uniform convergence
in \eqref{equation-106}. No infinite iteration or ordinary Banach
completeness is used.

\textbf{Finite target choices with standard bounds.}\quad
Consider first one forward parameter-Lipschitz family $V$
on the standard model, with bound $L\in\mathbb N$.
The auxiliary choices in Theorem~\ref{result-41} can be made independently of
$V$ and entirely on the fixed standard target geometry.
A finite marked triangulation and its barycentric coordinate
functions give the Hudson sequence
\[
 0=f^0\le f^1\le\cdots\le f^{R-1}=1.
\]
Consecutive differences are supported in prescribed vertex
stars. Subdividing the unit rise into an ordinary finite
number of increments makes the common spatial Lipschitz
constant $c$ of the functions satisfy $cL<1$.
The needed number of increments is an ordinary integer
because the original geometric constants and $L$ have
standard bounds. This is the step that cannot be justified
for an unrestricted infinite element $L\in M$.

All subsequent finite target data in Theorem~\ref{result-41} depend only on
this sequence, the standard spatial model, and the parameter
dimension: the interval boxes, support-label poset, small
conic ball configurations, strict enlargements, marked
coarsenings, and the target coloring of Mn\"ev's Lemma~8
with its unprescribed parameter part. They can be fixed over
the real algebraic numbers and then interpreted in $M$.
One precise way to make this assertion is to perform the
finite pure PL construction over the reals, retain its
finite combinatorial types and affine descriptors, and
express all incidences, strict inclusions, ball charts, and
compatibilities by their finite semialgebraic certificates.
Real-closed-field completeness then supplies the same
certificates over the real algebraic numbers. This transfers
only auxiliary finite PL data, never the given definable
family or a purported PL homotopy of it.

The marked finite construction and its use are detailed next,
so that no non-Archimedean Lebesgue-number argument is hidden
in this choice of data.


\textbf{The graph deformation and finite support labels.}\quad
Extend $V$ to the marked double by reflection. The fixed PL
reflection changes its parameter bound only by a standard
factor, which we absorb into the integer $L$. Use the same
Hudson functions and the following formulas from Theorem~\ref{result-41} and
Theorem~\ref{result-42}:
\[
 \rho_{y,v}(q)=\min_{0\leq j<R}
       \max\{f^j(y),q-2jv\},\qquad
 T_v=1+2(R-1)v,\quad D=2R-1,
\]
\[
 \Lambda(v,p)=\max\{p,D(p-1)+T_v\}.
\]
For $(x,p,v)$, solve the finite fixed-point equations
\begin{equation}
 t_i=\rho_{V_t(x),v}(\Lambda(v,p_i)),\qquad
 H_{p,v}(x)=V_t(x).
 \label{equation-107}
\end{equation}
The right side is a contraction in $t$ with constant $cL<1$;
the preceding residual-minimum argument proves existence,
uniqueness, definability, and joint continuity over $M$.
The inverse formula is
\begin{equation}
 H_{p,v}^{-1}(y)=
 V^{-1}_{(\rho_{y,v}(\Lambda(v,p_i)))_i}(y).
 \label{equation-108}
\end{equation}
It verifies the homeomorphism assertion directly. The initial
map is $V$, and all spatial marks are preserved. If $V$ fixes
the bottom or is the identity on a parameter face, so does
this deformation. These conclusions use the endpoint values
of $\rho$, not a compactness assertion about ordinary topology.

At $v=1$, put $q_i=Dp_i$. The moving bands are
$[2j,2j+1]$, $0\le j<R-1$. On such a band,
\[
 \rho_{y,1}(q_i)=
 \min\{f^{j+1}(y),\max\{f^j(y),q_i-2j\}\},
\]
and on the intervening plateau it equals $f^{j+1}(y)$;
the last index contributes only the final constant plateau.
Thus varying $q_i$ within that band can change \eqref{equation-108} only
where $f^{j+1}-f^j$ is nonzero, in its prescribed vertex star.
Choose the finite standard interval cover used in Theorem~\ref{result-41} so
that each interval meets at most one moving band. For a
coordinate box or a nonempty intersection of these boxes,
its active support is the union of at most $k$ such stars.
Define its label by the bands actually met by that intersection,
not by taking the union of the original box labels. Consequently
inclusion of boxes or their intersections gives inclusion of
support labels in the required direction. This is a finite
poset, and every geometric choice in it is independent of $V$.

The small-ball separation construction of Mn\"ev,
Proposition~57, pp.~89--90, as used in Theorem~\ref{result-41}, replaces those
monotone star supports by compatible configurations of finitely
many conic PL balls. Retain the full finite indexing poset,
even when two old configuration labels coincide. The explicit
finite rank induction of Theorem~\ref{result-41} enlarges the balls to make
proper indexing relations strict. This is the finite conic
extension of the construction behind Proposition~35,
pp.~61--62; it is not an invocation of that proposition with
an unmet strictness hypothesis. All input complexes and the
finite support poset have already been fixed over the real
algebraic numbers. The resulting choices can therefore be
included in the finite target certificates described above.
No subdivision is being asked to beat an infinitesimal
Lebesgue number.

\textbf{The marked coarsening and the coloring.}\quad
All target complexes are chosen compatible with the two
halves and the seams $A_0,A_1$, and the common fine complex
refines their barycentric subdivisions. Proposition~36,
pp.~62--63, supplies the conic-star invariance and its
cell-incidence property. Since the seams are locally flat
subcomplexes, these properties imply that each consistent
ball is either disjoint from the seams or forms with its
intersection a standard ball/equatorial-disk pair. This is
an inference using local flatness, not an additional literal
assertion in Proposition~36. The incidence property requires
the center to lie in any seam met by the ball, so it cannot
meet both disjoint seams.
For a crossing ball replace the single interior cell of the
unmarked coarsening by its two half-ball interiors and the
middle-disk interior, retaining their boundary cells. This
is the split marked coarsening of Theorem~\ref{result-41}, rather than the
original inscribed-ball conclusion of Lemma~6.

The terminal inverse images of the full balls are constant
over the relevant parameter cells. Outside the balls the
inverse map is pointwise constant; continuity gives the same
assertion on their boundaries. Hence any chosen boundary
subcells also have constant inverse images. For a fixed
mark $A$ or half $N^{\pm}$, preservation gives the exact
identities
\begin{equation}
 H_p^{-1}(D\cap A)=H_p^{-1}(D)\cap A,\qquad
 H_p^{-1}(D\cap N^{\pm})=H_p^{-1}(D)\cap N^{\pm}.
 \label{equation-109}
\end{equation}
Thus split prismaticity follows directly. In the decreasing
rank induction, boundary triangulations are chosen compatible
with the equatorial disks and extended on the complement by
finite relative PL triangulation. For nested balls, intersecting
with the same fixed global marks preserves every required
incidence. This is exactly the finite marked refinement checked
in Theorem~\ref{result-41}; it involves finite PL subdivisions and no limiting
process.

Apply the proof of Mn\"ev's Lemma~8, pp.~67--68, with its
unprescribed parameter part, no additional maps, and the
fixed target coloring. Its choices depend only on that coloring
and the configuration presheaf. They therefore produce the
same finite terminal target coloring for any families using
the same Hudson data. The equations that verify prismaticity
are the support constancies just proved, so their validity
over $M$ follows from the explicit formulas. The finite target
coloring itself may be taken from the real-algebraic auxiliary
certificates. This is the only use of the real version of the
finite PL construction; it supplies fixed finite data, not a
comparison theorem for the input family.

\textbf{Finite prism straightening over the field.}\quad
The target prisms are the literal products $C\times\sigma$
of Mn\"ev's Proposition~4, p.~21. Equation~(11) and
Proposition~5, pp.~22--23, identify their inverse images as
$D\times\sigma$, with $D$ independent of the parameter on
the entire closed simplex. These are finite ball complexes;
the whole boundary of each product ball is a union of proper
prisms of smaller total dimension. A lower parameter face
may carry a spatial refinement, which does not affect this
boundary assertion.

Triangulate the finitely many definable spatial source cells
compatibly, together with every cell of a fixed target PL
complex containing all marks. The finite marked
Hauptvermutung, Shiota, Theorem~2.1, p.~5, applies over the
given real closed field. If $\tau$ is the compatible definable
triangulation and $a$ is the marked PL homeomorphism supplied
by that theorem, then $\phi=\tau a^{-1}$ preserves every
cell of the fixed target complex and makes all source cells
PL after taking inverse images. Each is a PL ball by the
same theorem applied to its pair with its boundary.
The cell-preserving map $\phi$ is definably isotopic to the
identity by the finite cellwise Alexander construction below.
Thus common conjugation by $\phi$ does not change the
concordance class or any equality on the marked subcomplexes.

The PL endpoint is constructed by induction on total prism
dimension. Prescribe the required map on the already assigned
subcomplex. For a remaining product $D\times\sigma$, all of
$(\partial D\times\sigma)\cup(D\times\partial\sigma)$ is
already assigned. Choose compatible finite boundary
subdivisions making this map simplicial, without changing
any previously assigned map. Take fixed PL charts of $D$
and its target cell to convex balls and cone the TOTAL
product boundary from centers having the same base coordinate
$b_\sigma$. The resulting map is PL and preserves projection:
source and target points on corresponding cone segments have
the same parameter $(1-u)b_\sigma+up$. This is the construction
in Mn\"ev's Lemma~3, proof Steps~2--3, pp.~24--25.
It is a finite affine construction over $M$.

For completeness, the homotopy from a definable cellular map
to its PL endpoint also uses only finite field operations and
definable compactness. After taking the quotient by the PL
endpoint, construct the homotopy to the identity in increasing
prism dimension. In one spatial ball, let the parameter ball
be $D_0=\sigma\times I$. The already prescribed spatial
boundary maps extend radially to a continuous definable
family $R_p$ of homeomorphisms of the whole spatial ball.
On $p\in\partial D_0$, divide the prescribed whole-ball map
by $R_p$ to obtain a family $Q_p$ fixing the spatial boundary.
In fixed cube coordinates write $p=ru$, with
$r=\|p\|_\infty$ and $u\in\partial D_0$. Fill this family
using
\begin{equation}
 A_r(Q_u)(x)=
 \begin{cases}
 rQ_u(x/r),&\|x\|_\infty\le r,\ r>0,\\
 x,&\|x\|_\infty\ge r,
 \end{cases}
 \qquad A_0=\mathrm{id}.
 \label{equation-110}
\end{equation}
The formulas agree on the overlap. Their displacement is at
most $2r$, proving joint continuity at the parameter center
in the field order topology. Multiplying by $R_p$ extends
all the prescribed data. Fixed PL changes of ball coordinates
handle general cells. Finite pasting gives the required joint
definable homotopy, relative to any already prescribed
subcomplex. This also proves the cellwise contraction used
for $\phi$. No assertion that this particular homotopy is
jointly PL is needed.

For a single based bottom-fixed input, prescribe the identity
on the bottom and on the parameter boundary before the prism
induction. This proves the standard-bound sufficient criterion
corresponding to Theorem~\ref{result-41}.

\textbf{Separate straightening of fixed factors.}\quad
The spatial-clock construction below gives factors that individually
fix the bottom and are the identity on the parameter boundary.
Consequently the single-family criterion applies to each factor
separately.  If their homotopies are $R_t,S_t$, then $S_tR_t^{-1}$
is a based bottom-fixed homotopy, and its terminal map is jointly PL.
No synchronization of their source cell complexes is required.

\textbf{Producing the required standard bounds.}\quad
Let $U$ be an arbitrary based definable $k$-cube of concordances
on the standard model. The spatial factorization of Theorem~\ref{result-15}
works over $M$ and gives
\begin{equation}
 U_p=s_pr_p^{-1}.
 \label{equation-111}
\end{equation}
Both factors preserve the spatial marks, fix the bottom pointwise,
and are the identity on the parameter boundary. They have a
common STANDARD spatial Lipschitz bound. To check the last
word, run the coordinate clocks on the unit cubical model.
Every evolving component remains in a fixed unit cube. Each
stage has an ordinary finite monotonicity count $N$, so its
variation bound and all mixed-coordinate amplification factors
are ordinary finite numerical expressions in $N$, the number
of coordinates, and the preceding bounds. The final bounds
are therefore below an ordinary integer. The fixed PL changes
of spatial model multiply them only by standard constants.

Form the disjoint tagged graph family
\begin{equation}
 \mathcal G_p=
  (\{0\}\times\Gamma r_p)\cup
  (\{3\}\times\Gamma s_p)
 \label{equation-112}
\end{equation}
on the unit cubical model. The tag separation exceeds the
diameter of either untagged graph space. Hence the Hausdorff
distance of these unions is the maximum of the two individual
graph distances. Divide all ambient coordinates by $3$ to put
the family in a unit cube; undoing this scaling costs only the
standard factor $3$. It is a continuous compact-set family and
is constant on the parameter boundary.

Apply Theorem~\ref{result-18} to obtain a compact resolution consisting of
finitely many integer-tagged unit cubes, with standard
projection and Hausdorff constants. Every coordinate of this
compact source is bounded by an ordinary integer. Append the
boundary copy used in Theorem~\ref{result-17} with an integer tag larger than
all existing tags and diameters. Because the boundary family
is constant, this preserves standard bounds for both the
projection and the pulled family.

The labelled join of Theorem~\ref{result-17} uses at most $k+2$ slots. If its
source coordinate bound is $M_0$, its near-pair estimate is
$m(1+M_0)$ times the original Lipschitz bound, and its far-pair
estimate uses $2m$ times the bounded output diameter. Here
$m,M_0$, and all those input bounds are standard. Thus the
projection and compact-set family on the join have standard
Lipschitz constants. The finite universal resolution and
relative homotopy-extension argument give a continuous
map $w$ from the parameter cube to this join, with its literal
affine boundary section. Apply the finite vector coordinate
clocks to $(\mathrm{id},w)$, after a fixed standard scaling of
the target coordinates if needed. The same bounded-variation
calculation gives a standard bound on $w\chi$. Every boundary
coordinate line is affine, so its variation is $ct$ and the
normalized clock $(t+ct)/(1+c)$ is exactly $t$. Hence
$\chi$ fixes the parameter boundary pointwise.

It follows that the map $g$ supplied by Theorem~\ref{result-17} in this
application fixes the boundary and has STANDARD graph
Hausdorff control for $\mathcal G\circ g$. Let $L_\Gamma$
be an ordinary integer bounding that control in the original
untagged graph coordinates, and let $K$ be the common standard
spatial bound of the factors. A point $(x,r_{g(p)}(x))$ can
be matched, by definable compactness, to
$(x',r_{g(q)}(x'))$ at graph distance at most
$L_\Gamma\|p-q\|_\infty$. Therefore
\begin{equation}
 \|r_{g(p)}(x)-r_{g(q)}(x)\|_\infty
 \le(1+K)L_\Gamma\|p-q\|_\infty.
 \label{equation-113}
\end{equation}
The same estimate holds for $s$. No estimate for $r^{-1}$ or
$s^{-1}$ was used. We have obtained a pair of factors with
standard forward parameter bounds and all the equality
conditions of the single-family criterion applied separately to both factors.

Apply that criterion to $r\circ g,s\circ g$. Their quotient
$U\circ g$ is definably homotopic through based concordance
families to a jointly PL family. The convex parameter homotopy
\begin{equation}
 U_{(1-t)p+t g(p)},\qquad 0\le t\le1,
 \label{equation-114}
\end{equation}
is based, fixes the bottom, and connects the original family
to $U\circ g$. This proves \eqref{equation-105}. Translation by PL
representatives handles other base components, and Theorem~\ref{result-4}
supplies the component bijection. Empty bases have only their
empty maps; the constructions also apply componentwise and
to zero-dimensional bases.

\textbf{The remaining relative distinction.}\quad
For a nonconstant prescribed PL parameter boundary over $M$,
its parameter Lipschitz constant can be larger than every
ordinary integer. Appending that boundary as an exact section
in Theorem~\ref{result-17} then need not preserve a standard bound. Neither
Theorem~\ref{result-18} nor the present proof erases this exact-boundary issue.
Consequently this theorem does not infer nonreal injectivity
from Theorem~\ref{result-42}. A separate parameter-preserving PL field-invariance
argument, or a direct exact-relative standard-bound reduction,
is still needed for that inference. A definable path of PL
coefficient data alone is not such an argument.
\end{proof}\section{The relative comparison and injectivity}
\begin{proof}[Proof of Theorem~\ref{result-44}]
Higher surjectivity is Theorem~\ref{result-43} and the component bijection is
Theorem~\ref{result-4}. We prove injectivity, retaining the exact parameter
projection throughout. Replace $B$ by its standard abstract
finite PL model as in Theorem~\ref{result-43} and write $F=\mathbb R_{\mathrm{alg}}$
for the ordered real closed subfield of real algebraic numbers
in $M$. The fixed model and its two end markings are defined
over $F$.

\textbf{Parameter-preserving field comparison.}\quad
We use Theorem~\ref{result-40}. This is not an
appeal to field invariance of ordinary, unparameterized PL
isotopies. Its proof simultaneously subdivides the one-way
chain
\[
 N_M\times P_M\xrightarrow{h}N_M\times P_M
                  \xrightarrow{\operatorname{pr}_P}P_M,
\]
retains fixed product-cell carriers, and transfers the finite
vertex configuration to $F$. The matching-vertex maps $u,v$
cover the SAME base map $\beta$. Thus
$(\mathrm{id}_N\times\beta^{-1})u$ and
$(\mathrm{id}_N\times\beta^{-1})v$ preserve the parameter
pointwise and preserve every fixed product cell.
The finite total-product cone construction contracts them
with their agreement on the bottom and parameter boundary
retained. The same construction removes the base conjugation
by $\beta$. This produces a jointly PL concordance
homotopy to an $F$-coefficient representative. No polygonal
approximation to a coefficient-space path is used.

Let $h$ be a based jointly PL $k$-cube over $M$, $k\ge1$,
which is nullhomotopic through based definable concordances.
By Theorem~\ref{result-40}, choose a based jointly PL homotopy from $h$ to
an $F$-PL cube $h_0$. Concatenate the reversed homotopy
with the given definable nullhomotopy of $h$. We obtain a
definable filling $U$ of the full boundary data on a
$(k+1)$-cube $P$: $h_0$ on the initial face and the identity
on the other faces. These full boundary data are jointly PL
over $F$ and agree exactly on their intersections.

We now prove that such a filling with $F$-PL boundary admits
a PL filling over $M$ with exactly the same boundary.

\textbf{Standard bounds for the boundary factors.}\quad
Use the fixed $F$-PL cubical model $Q\subset C=[0,1]^n$ of
Theorem~\ref{result-15}, and write $U^Q$ for the conjugated family.
Its boundary restriction has a PL vector extension
\[
 E_\partial:C\times\partial P\longrightarrow C
\]
over $F$: take a finite compatible triangulation over $F$,
leave the prescribed vertex values unchanged, assign a
fixed rational point of $C$ to the other vertices, and
extend affinely. This is a vector extension and assumes
no extension by homeomorphisms. Paste it with $U^Q$ on
$Q\times P$ and use the full-subcomplex extension of Theorem~\ref{result-15}
to obtain a continuous definable vector map $E$ on $C\times P$.

Run the PL-normalized coordinate clocks of Theorem~\ref{result-42}, starting
with $(\mathrm{id},E)$. At each stage all component values
lie in a fixed unit cube. If their number is $d$ and the
uniform monotonicity count is $N$, the raw clock height is
at most $1+dN$. Choose its trapezoid normalization parameter
$D$ to be a STANDARD INTEGER larger than $\max(3,1+dN)$.
The counts $d,N$ are ordinary finite integers even though
they may depend on the input over $M$.

On the parameter boundary the initial data are PL over $F$.
The finite maximum formula for line variation is semilinear
over $F$ when the input is PL over $F$; continuity makes
the resulting variation function PL over $F$. The
normalization $T_D$ of Theorem~\ref{result-42} has coefficients in $F$,
and the inverse of the resulting parameter-preserving PL
clock has coefficients in $F$ as well. Induction shows
that every boundary clock, and the two final boundary
factors, remain jointly PL over $F$.
Hence the factorization
\begin{equation}
 U_p=s_pr_p^{-1}
 \label{equation-116}
\end{equation}
has uniformly spatially Lipschitz factors with standard
bounds, and its boundary factors have STANDARD forward
parameter Lipschitz bounds. The latter assertion follows
from their finite PL descriptors over the Archimedean field
$F$; it is not a rounding of an arbitrary element of $M$.

The factors individually fix the bottom. Indeed the bottom
is a union of coordinate faces of $Q$, and $E=\mathrm{id}$
there. On each free coordinate line in such a face the
evolving vector is $(\mathrm{id},\mathrm{id})$, so its raw
clock is $3t$. The normalization of Theorem~\ref{result-42} is exactly the
identity at height $3$. Coordinates fixed by that face
remain fixed endpoints. Induction gives $r=s=\mathrm{id}$
on the bottom. Both factors preserve the top setwise.

\textbf{Exact relative graph regularization.}\quad
Form the two tagged graph family as in Theorem~\ref{result-43}. It has a
standard spatial bound and its parameter-boundary graph
family has a standard Hausdorff bound, by the preceding
$F$-PL assertion. Theorem~\ref{result-18} gives a standard compact Lipschitz
resolution. Add the exact boundary copy of Theorem~\ref{result-17}.
Its tag coordinates, its projection, and its nonconstant
pulled boundary family all have standard bounds.
The finite join estimate and bounded vector clocks
therefore give a map $g:P\to P$ fixing $\partial P$
pointwise and standard Hausdorff control for both pulled
graph families. As in \eqref{equation-113}, the common standard
spatial bound converts this to a standard forward
parameter bound for
\[
 \bar r=r\circ g,\qquad \bar s=s\circ g.
\]
Their boundary restrictions are still exactly $r,s$,
and their bottom restrictions are still the identity.

\textbf{The exact PL endpoint and its boundary collar.}\quad
Apply the extension-only relative argument of Theorem~\ref{result-42} to
each factor. Its only Archimedean use was the choice of
a finite Hudson grid against the forward parameter bound.
That bound is now standard, so Theorem~\ref{result-43} supplies the fixed
finite target data over $F$. The contraction is solved
over $M$ by the residual-minimum argument.

The PL roof normalization of Theorem~\ref{result-42} gives a jointly PL
homotopy track on the parameter boundary. Its inverse
formula evaluates only the given PL boundary inverse.
The target prisms and all boundary source prisms are PL;
the common marked ambient triangulation preserves these
boundary cells setwise. Thus the ORIGINAL boundary map
has the required correspondence after conjugation.
Prescribe it, together with the bottom identity, first.
The total-product cone induction extends it over the
remaining finite prisms. Lower parameter faces may carry
spatial subdivisions; their prisms lie in the already
assigned lower-dimensional product boundary. No PL filling
is assumed at this step.

A fixed PL collar of an inner parameter cube, carrying
the reversed PL boundary track, restores the initial
boundary exactly. All these operations are finite PL
operations over $M$ and require no further metric mesh.
Consequently we obtain jointly PL bottom-fixed families
$R,S$ with the boundary values of $r,s$, respectively.
Their quotient
\begin{equation}
 W_p=S_pR_p^{-1}
 \label{equation-117}
\end{equation}
is a PL filling of $U|_{\partial P}$ with exactly its
original end markings and parameter-boundary values.

In particular $h_0$ is PL nullhomotopic over $M$.
Concatenate this nullhomotopy with the PL homotopy
$h\simeq h_0$ supplied by Theorem~\ref{result-40}. The result is a PL
nullhomotopy of $h$, proving injectivity on $\pi_k$.
All concatenations use fixed affine subdivisions of the
time interval, so are jointly PL where asserted.
Together with Theorem~\ref{result-43} and Theorem~\ref{result-4} this proves \eqref{equation-115}.
Translation by PL representatives handles all base
components. Empty or disconnected bases are treated
componentwise as in the preceding results.
\end{proof}
\subsection{Effective models over every o-minimal real closed field}

The all-field comparison makes the same finite PL codes describe
the entire definable concordance space in every chosen o-minimal
expansion.

\begin{corollary}\label{cor:effective-whole-definable-concordances}
Let $M$ be an o-minimal expansion of a real closed field $R$,
let $B$ be a closed definably compact definable manifold, and
choose a definable triangulation
$\tau:|K|_R\longrightarrow B$ with $K$ a finite abstract PL
manifold. The effective simplicial group
$G^{\mathrm{PL}}(K)$ of
Theorem~\ref{thm:effective-countable-concordance-models}
has a specified weak equivalence
\begin{equation}\label{eq:effective-whole-definable-model}
 G^{\mathrm{PL}}(K)\longrightarrow P_{\mathrm{def},M}(B).
\end{equation}
It is real-algebraic coefficient extension, followed by PL
inclusion and conjugation by $\tau\times\operatorname{id}_I$.
All families preserve the whole parameter, fix the bottom
pointwise, and have free top.

Uniformly from the finite simplex list of $K$ and $r\geq0$,
one can enumerate a complete countable presentation of
$\pi_rP_{\mathrm{def},M}(B)$. In degree zero this is the
full component group. In positive degrees the basepoint
is the identity, and translation gives the other components;
one may instead enumerate an integer relation matrix.
For the real realization $|K|_{\mathbb R}$, the same model
and presentations describe ordinary TOP concordances if
every component has dimension at most two or at least five.

The algorithm takes $K$ alone, with its manifold condition
as a promise. The triangulation specifies the comparison
map and need not be computable from a description of $B$.
No decision procedure for the theory of $M$ is an input.
The presentations are recursively enumerable; finiteness
and a terminating word problem are not asserted.
\end{corollary}

\begin{proof}
Let $F$ be the real algebraic field, embedded canonically
as an ordered subfield of $R$. The code construction gives
$G^{\mathrm{PL}}(K)=P_{\mathrm{PL},F}(|K|_F)$, and
Theorem~\ref{thm:effective-countable-concordance-models}
proves that coefficient extension to
$P_{\mathrm{PL},R}(|K|_R)$ is a weak equivalence.
Theorem~\ref{result-44} gives the next weak equivalence,
to $P_{\mathrm{def},M}(|K|_R)$.

Put $T=\tau\times\operatorname{id}_I$. On every parameter
simplex, conjugation sends a family $h_p$ to
$T h_pT^{-1}$. It preserves joint definability, continuity
of the total map and its inverse, the fixed bottom and
the free top. It commutes with every face and degeneracy
because $T$ is independent of the parameter. Conjugation
by $T^{-1}$ is its inverse. It is therefore an isomorphism
of definable simplicial groups, proving
\eqref{eq:effective-whole-definable-model} by composition.

For completeness, the presentation constructed earlier
uses the decidable cycle set $Z_r$ and Moore subgroup
$N_{r+1}$ in the coded group $G$. Its relations are
\[
 x_1=1,\qquad x_ax_b=x_{ab}\quad(a,b\in Z_r),
 \qquad x_{d_0u}=1\quad(u\in N_{r+1}).
\]
The multiplication relations present all of $Z_r$ and
the boundary relations quotient by precisely $d_0N_{r+1}$.
The first-face Moore formula established in the preceding
theorem identifies this quotient with $\pi_r|G|$, including
$r=0$. Enumerating finite strings, and pairs of strings,
up to each successive length enumerates every required
generator and relation. Every stage uses only terminating
tests on the finite codes for $K$. The comparison with
$M$ changes the interpretation of this presentation,
not the algorithm. Positive-degree groups are abelian,
so their multiplication and boundary relations give the
stated integer matrix as before.

For a real base, Theorem~\ref{result-4} gives the ordinary
PL--TOP comparison in dimensions at least five, and
Theorem~\ref{thm:low-dimensional-real-concordance-comparison}
gives it in dimensions at most two. A bottom-fixed
concordance preserves every component of $|K|\times I$:
each component meets its own fixed bottom. The entire
group is consequently the finite product of the groups
of those components, and the comparisons apply to that
product, including when the component dimensions differ.
For $K=\varnothing$, all groups are the one-element group.
\end{proof}


\subsection{Field extension, equal-boundary factors and graph regularization}
We record several consequences with their full hypotheses. Their proofs
use the completed comparison theorem and do not reinsert an earlier
factor criterion into its prerequisites.

\begin{corollary}[Semialgebraic coefficient extension]\label{result-52}
Let $F\subset K$ be real closed fields and let $B$ be a closed compact
finite PL manifold defined over $F$. Coefficient extension induces a
weak equivalence
\[
          P_{{\rm SA},F}(B)\longrightarrow P_{{\rm SA},K}(B_K).
\]
The square with the PL coefficient-extension maps commutes. Hence the
PL-to-semialgebraic comparison of a finite abstract PL model is
independent of the real closed field, and can be compared to its real
realization by a zigzag through the real algebraic field. In particular
it is a component bijection and an isomorphism in every positive
homotopy degree, in every base dimension.
\end{corollary}
\begin{proof}
Use Theorem~\ref{result-44} for the semialgebraic structures on both
fields. The two vertical PL-to-semialgebraic maps are weak equivalences,
and Theorem~\ref{result-40} makes the upper PL coefficient-extension
map a weak equivalence. The square commutes because each map interprets
the same finite formulas and affine descriptors. On each based homotopy
group, including components, the lower map is therefore an isomorphism.
The zigzag and the final assertions follow exactly as in
Theorem~\ref{result-40}. This is a statement about finite parameter
families, without any map from nonreal families to ordinary TOP spaces.
\end{proof}

\begin{corollary}[The full equal-boundary pair criterion]\label{result-53}
Let $M$ be any fixed o-minimal real closed field, let $B$ be a closed
compact finite PL manifold, and put $N=B\times I$, $A_j=B\times\{j\}$,
$P=I^k$, with $k\ge1$. Suppose $r,s:N\times P\to N$ are jointly
continuous definable families of homeomorphisms, each preserving both
$A_0,A_1$ setwise, and satisfying
\[
 r_p|_{A_0}=s_p|_{A_0}\quad(p\in P),\qquad
 r_p=s_p\quad(p\in\partial P).
\]
Then $U_p=s_pr_p^{-1}$ is based and bottom-fixed, and is based
definably homotopic through such concordances to a jointly PL family.
In particular this proves the pair criterion with uniform forward
parameter Lipschitz bounds over the real field, and with ordinary
integer bounds on a standard PL model over any $M$. No inverse
parameter bound is required. In fact no metric bound on these factors
is needed for the conclusion after Theorem~\ref{result-44}.

Conversely every class in the PL comparison image has both a jointly
PL representative with a finite field-valued forward parameter bound,
and a representation by two factors with those bounds and the displayed
agreements. Consequently the two sufficient regularity criteria remain
valid whenever the theory of $M$ has a real model. Over the real field,
surjectivity is equivalently expressed by requiring every based
concordance cube with Hausdorff-Lipschitz graphs to be homotopic to such
a two-factor representative; both assertions hold. These statements
include the semialgebraic six-torus in degree two.
\end{corollary}
\begin{proof}
If $x\in A_0$, then $r_p^{-1}(x)\in A_0$, so the first agreement
gives $s_pr_p^{-1}(x)=x$. The second agreement makes the quotient
identity on the full parameter boundary. The quotient preserves the
other end as well and is jointly continuous and definable with a
continuous inverse, by the family conventions. Theorem~\ref{result-44}
gives a based definable homotopy from this cube to a jointly PL one.
This argument does not assume that either original factor individually
fixes the bottom or the parameter boundary.

A jointly PL representative $W$ on a compact finite PL product has a
finite Lipschitz constant in fixed coordinates. Use the factors
$r=\operatorname{id}$ and $s=W$. They have all the asserted agreements,
and both forward parameter bounds exist. Over the real field these
are ordinary finite constants. If standard bounds on a fixed standard
model are desired over $M$, first use Theorem~\ref{result-40} to choose
a real-algebraic PL representative; its finite descriptors have
ordinary integer bounds. Thus all stated sufficient and converse
criteria follow. The graph-restricted assertion follows by applying
the same construction to those representatives; conversely both it
and surjectivity have now been proved for every family, so they are
equivalent. No homeomorphic reparameterization of a prescribed
representative is asserted.
\end{proof}

\begin{corollary}[Relative regularization of compact graphs]\label{result-54}
Let $M$ be any o-minimal real closed field, let $X$ be a compact
definable set, and let $U_p:X\to X$, $p\in P=[0,1]^k$, be a jointly
continuous definable family of homeomorphisms which is the identity
on $\partial P$. There is a definable Lipschitz map $g:P\to P$,
equal to the identity on $\partial P$, such that the graphs of
$V_p=U_{g(p)}$ are uniformly Hausdorff-Lipschitz. The family
$U_{(1-t)p+tg(p)}$ is a based definable homotopy from $U$ to $V$ and
preserves every spatial condition imposed on all members of $U$.
More generally the same relative parameter conclusion holds for any
bounded Hausdorff-continuous family of nonempty compact sets whose
boundary family is Hausdorff-Lipschitz. The map $g$ need not be
injective, and no point matching between graph fibres is asserted.
\end{corollary}
\begin{proof}
For nonempty $X$, the graphs form a bounded Hausdorff-continuous family:
uniform continuity of the total family gives a matching at the same
spatial input, in both directions. Their boundary family is constant.
Translate and scale the bounded coordinates into a unit cube, apply
Theorem~\ref{result-18} to obtain a compact Lipschitz resolution, and
apply the exact relative metric consequence of Theorem~\ref{result-17}.
Rescaling changes the field-valued constants by only fixed finite
factors. The resulting $g$ has the asserted graph bound and fixes the
boundary. Convexity of $P$ verifies the displayed homotopy, and every
member is an original family member, proving preservation of all spatial
constraints. If $X$ is empty, use $g=\operatorname{id}$; the family
contains only the empty map and its graph distance is zero. The same
two theorems give the stated general compact-set version directly.
\end{proof}

\begin{corollary}[Ordinary topological comparison and its kernel]\label{result-56}
Let $B$ be a closed compact PL manifold. Over any fixed o-minimal
real closed field, every based definable circle family of concordances
is based definably homotopic to a family that is PL as a map of total
spaces. Inclusion is an isomorphism on components and is surjective
on every positive homotopy group; in fact it is a weak equivalence.
In particular this includes the semialgebraic category over every
real closed field, and every o-minimal theory having a real model.

Over an o-minimal expansion of $\mathbb R$, suppose every component
of $B$ has dimension at most two or at least five. Then both maps
\[
 P_{\rm PL}(B)\longrightarrow P_{\rm def}(B)
             \longrightarrow P_{\rm TOP}(B)
\]
are weak equivalences. Thus PL inclusion is injective on every
homotopy group, is an isomorphism on $\pi_1$, and for $k\ge2$
there is the natural splitting
\[
 \pi_kP_{\rm def}(B)\simeq\pi_kP_{\rm PL}(B)\oplus
 \ker\bigl(\pi_kP_{\rm def}(B)\to\pi_kP_{\rm TOP}(B)\bigr),
\]
where the displayed kernel is zero. No comparison with ordinary TOP
spaces is asserted here in base dimensions three or four.
\end{corollary}
\begin{proof}
Theorem~\ref{result-44} gives every assertion about PL and definable
families, including the circle representative and its based homotopy.
For the real TOP comparison in dimensions at least five, use
Theorem~\ref{result-4}. For base dimension $n\le2$, repeat its
relative structure-space argument on $B\times I$: the interior
dimension is $n+1\le3$ and the boundary dimension is $n\le2$,
so neither equals four. Kirby--Siebenmann's Theorems 5.8(I) and 5.9,
Essay V, printed p.~253, make both
$\mathrm{TOP}_n/\mathrm{PL}_n$ and
$\mathrm{TOP}_{n+1}/\mathrm{PL}_{n+1}$ contractible
\cite{KirbySiebenmann1977}; the zero-dimensional group is trivial.
Their stabilization is therefore a weak equivalence. The relative
lift space with its free upper boundary is contractible by exactly
the same fibre-product and cylinder contraction, and the explicit
bottom compression in that proof handles our pointwise boundary
condition. This proves PL-to-TOP equivalence in all the stated
dimensions. Components of different dimensions are handled by
finite products.

The composite PL-to-TOP map factors through definable concordances.
The first map and the composite are weak equivalences, so the second
is one as well. The kernel in the displayed splitting consequently
vanishes. The inverse of the induced PL inclusion, together with
the zero kernel, gives that splitting. This proves its precise
kernel assertion rather than only a PL-to-definable surjectivity
statement.
\end{proof}

\section{Further comparison questions}
The all-field comparison concerns finite jointly PL and finite jointly
definable families.  The relation with the whole ordinary topological
concordance space in exceptional dimensions is an additional question;
it is not supplied by the definable comparison.

One component question remains in dimension five. Let $M$ be a
compact connected PL five-manifold with $\pi_2M=0$, with concordances
fixing $M\times\{0\}\cup\partial M\times I$, and let $K\simeq M$
be a finite connected pointed model. Proposition
\ref{prop:five-dimensional-stable-surjectivity} proves that
positive suspension gives a surjective homomorphism
$\pi_0P_{\mathrm{PL}}(M)\to E(K)$; its proof does not require
$\pi_2M=0$. The remaining question is whether this homomorphism
is injective under the stated hypothesis.
Igusa states the corresponding geometric exact sequence in
\cite[Theorem~7.1, p.~23]{Igusa2025}. Its local framed-curve
relations give part of the argument, but the null-deformation
criterion in Theorem~7.5, p.~27, refers to further geometric work.
The missing step is completeness: a concordance with zero
stabilized obstruction must admit the required relative
two-parameter deformation. The corrected sequence of
\cite[Theorem~8.a.2, pp.~149--150]{Igusa1984} applies in dimension
at least six. The closed simply connected case in
Theorem~\ref{thm:simply-connected-five-components} is already
settled and does not require $\pi_2M=0$.

There is also a remaining algebraic question.  Let $A$ be a complete
separated local ring with finite residue field and let $A\subset C$
be a commutative unital inclusion.  Suppose that $e\in C$ is an
idempotent, that $Q_0=eC$ is nonreduced local Artinian, and that
$A_0=eA$ is a complete one-dimensional Noetherian Cohen--Macaulay
local quotient.  Suppose a finite full $A_0$-order
$A_0\subseteq O\subseteq Q_0$, with $Q(O)=Q_0$, has nonzero
conductor $\{x\in O:exO\subseteq A\}$.  Must the binary tensor
map for $\Omega(A,C)$ fail to be injective?
Corollary~\ref{cor:binary-selected-nonzero-conductor} answers this
when $A$ is Noetherian; Lemma~\ref{lem:binary-selected-factor-transfer}
answers the arbitrary-base case with a regular conductor.
Proposition~\ref{prop:nonnoetherian-nilpotent-conductor} gives a
family with non-Noetherian $A$ and a square-zero conductor.
In general a nonzero conductor need not contain a regular
parameter power times $O$.  Moreover, a nonzero tensor over a
Noetherian subring can vanish after extension to $A$, so the
Noetherian classification does not settle the question by descent.
The finite residue field cannot be omitted, by
Proposition~\ref{prop:infinite-residue-nilpotent-conductor}.

\section*{Method and records}
The mathematics and text were produced by GPT-6 Astra (ChatGPT Pro mode).
Checks in this preparation were performed by the same model; the
numeric effort setting was not exposed.  The separate records identify
the scopes of direct checks, separate reconstructions, and remaining
external-source comparisons.
An editorial reading by Claude Opus 5.5 at high effort identified
the freeness consequence of the finite-conductor argument and its
relation to Elliott's conjectures. The broader formulation in terms
of analytic irreducibility, the explicit witnesses, and the complete
freeness argument were then proved and checked here, including a
separate reconstruction from the theorem's statement.
During a separate reconstruction from the statements, the polynomial
lattice and orthogonal-component arguments in
Lemmas~\ref{result-45}--\ref{result-46} yielded the stronger binary
injectivity criterion in Theorem~\ref{result-39}. The existing
generic-nilpotent obstruction and positive all-arity arguments are
reused; the records distinguish the independent necessity reconstruction
from its subsequent source comparison and from the separately checked
sufficient implication.
No human has checked this work.  The text is dedicated to the public
domain under CC0 1.0.  The program \href{\workbenchurl}{\emph{Finite polynomial tensors}} checks finite-field
examples, the explicit series witnesses, and the displayed double-complex signs; its finite results
do not replace the proofs above. Its proposed repository, with the
accompanying records, is the
\href{\workbenchurl}{homotopy and categories workbench}; the
\href{\recordurl}{mathematics workbenches register} is its registry.
\bibliographystyle{alpha}
\bibliography{references}
\end{document}
